% 原书第62—69页：4.2及4.2.1—4.2.4。
\FloatBarrier
\section{结构识别方法}

下面介绍如何利用第4.1节得到的可识别性结果进行因果发现，即根据有限数据集估计一个图。这些都是具有挑战性的统计问题，可以采用许多不同的途径处理。这里着重介绍方法思路，并不声称所介绍的方法能够最有效地利用数据。今后的研究完全可能提出新颖而有效的方法。我们只介绍少数例子，主要选择那些我们对其性能已经积累了一定经验的方法。

\subsection{加性噪声模型}

对于利用定理4.5中加性噪声模型可识别性进行因果学习的方法，主要请参见讨论多元情形的第7.2节。这里概述两种方法，但不声称它们最优。第一种方法检验\textbf{残差的独立性}，它是“回归后进行独立性检验”（regression with subsequent independence test，RESIT）算法的一个特例，见第7.2节。
\begin{enumerate}
 \item 对 $Y$ 关于 $X$ 进行回归，即采用某种回归技术，将 $Y$ 写成 $X$ 的函数 $\widehat f_Y$ 与某个噪声之和。
 \item 检验 $Y-\widehat f_Y(X)$ 是否与 $X$ 独立。
 \item 交换 $X$ 和 $Y$ 的角色，重复上述过程。
 \item 如果一个方向接受独立性，而另一个方向拒绝独立性，就将前一个方向推断为因果方向。
\end{enumerate}

% 中文 TikZ 重排：坐标及曲线提取自原 PDF，未重新模拟数据。
\begin{figure}[htbp]
\centering
\definecolor{chfourfig5c0}{rgb}{0.13669,0.12195,0.12529}
\definecolor{chfourfig5c1}{rgb}{0.93042,0.1109,0.14145}
\definecolor{chfourfig5c2}{rgb}{0.18155,0.18994,0.57328}
\begin{tikzpicture}[x=0.039cm,y=-0.039cm]
\useasboundingbox (0,0) rectangle (358,310);
\path[fill=chfourfig5c0,fill opacity=0.5,draw=chfourfig5c0,line width=0.3985pt,draw opacity=0.5] (119.7146,49.5059) .. controls (119.7146,48.7851) and (119.1303,48.2008) ..(118.4095,48.2008) .. controls (117.6887,48.2008) and (117.1044,48.7851) ..(117.1044,49.5059) .. controls (117.1044,50.2267) and (117.6887,50.8110) ..(118.4095,50.8110) .. controls (119.1303,50.8110) and (119.7146,50.2267) ..(119.7146,49.5059);
\path[fill=chfourfig5c0,fill opacity=0.5,draw=chfourfig5c0,line width=0.3985pt,draw opacity=0.5] (104.7505,75.7777) .. controls (104.7505,75.0569) and (104.1662,74.4726) ..(103.4454,74.4726) .. controls (102.7246,74.4726) and (102.1403,75.0569) ..(102.1403,75.7777) .. controls (102.1403,76.4985) and (102.7246,77.0828) ..(103.4454,77.0828) .. controls (104.1662,77.0828) and (104.7505,76.4985) ..(104.7505,75.7777);
\path[fill=chfourfig5c0,fill opacity=0.5,draw=chfourfig5c0,line width=0.3985pt,draw opacity=0.5] (95.8837,63.4338) .. controls (95.8837,62.7130) and (95.2994,62.1287) ..(94.5786,62.1287) .. controls (93.8578,62.1287) and (93.2735,62.7130) ..(93.2735,63.4338) .. controls (93.2735,64.1546) and (93.8578,64.7390) ..(94.5786,64.7390) .. controls (95.2994,64.7390) and (95.8837,64.1546) ..(95.8837,63.4338);
\path[fill=chfourfig5c0,fill opacity=0.5,draw=chfourfig5c0,line width=0.3985pt,draw opacity=0.5] (60.5757,111.4145) .. controls (60.5757,110.6937) and (59.9913,110.1094) ..(59.2705,110.1094) .. controls (58.5497,110.1094) and (57.9654,110.6937) ..(57.9654,111.4145) .. controls (57.9654,112.1353) and (58.5497,112.7197) ..(59.2705,112.7197) .. controls (59.9913,112.7197) and (60.5757,112.1353) ..(60.5757,111.4145);
\path[fill=chfourfig5c0,fill opacity=0.5,draw=chfourfig5c0,line width=0.3985pt,draw opacity=0.5] (122.1854,46.2182) .. controls (122.1854,45.4974) and (121.6011,44.9131) ..(120.8803,44.9131) .. controls (120.1595,44.9131) and (119.5752,45.4974) ..(119.5752,46.2182) .. controls (119.5752,46.9390) and (120.1595,47.5233) ..(120.8803,47.5233) .. controls (121.6011,47.5233) and (122.1854,46.9390) ..(122.1854,46.2182);
\path[fill=chfourfig5c0,fill opacity=0.5,draw=chfourfig5c0,line width=0.3985pt,draw opacity=0.5] (140.0686,44.8632) .. controls (140.0686,44.1424) and (139.4843,43.5581) ..(138.7635,43.5581) .. controls (138.0426,43.5581) and (137.4583,44.1424) ..(137.4583,44.8632) .. controls (137.4583,45.5840) and (138.0426,46.1684) ..(138.7635,46.1684) .. controls (139.4843,46.1684) and (140.0686,45.5840) ..(140.0686,44.8632);
\path[fill=chfourfig5c0,fill opacity=0.5,draw=chfourfig5c0,line width=0.3985pt,draw opacity=0.5] (58.6528,111.1953) .. controls (58.6528,110.4745) and (58.0685,109.8902) ..(57.3477,109.8902) .. controls (56.6269,109.8902) and (56.0426,110.4745) ..(56.0426,111.1953) .. controls (56.0426,111.9162) and (56.6269,112.5005) ..(57.3477,112.5005) .. controls (58.0685,112.5005) and (58.6528,111.9162) ..(58.6528,111.1953);
\path[fill=chfourfig5c0,fill opacity=0.5,draw=chfourfig5c0,line width=0.3985pt,draw opacity=0.5] (155.0426,50.7711) .. controls (155.0426,50.0503) and (154.4583,49.4660) ..(153.7375,49.4660) .. controls (153.0167,49.4660) and (152.4324,50.0503) ..(152.4324,50.7711) .. controls (152.4324,51.4919) and (153.0167,52.0763) ..(153.7375,52.0763) .. controls (154.4583,52.0763) and (155.0426,51.4919) ..(155.0426,50.7711);
\path[fill=chfourfig5c0,fill opacity=0.5,draw=chfourfig5c0,line width=0.3985pt,draw opacity=0.5] (108.7655,67.3990) .. controls (108.7655,66.6782) and (108.1812,66.0939) ..(107.4604,66.0939) .. controls (106.7396,66.0939) and (106.1553,66.6782) ..(106.1553,67.3990) .. controls (106.1553,68.1198) and (106.7396,68.7041) ..(107.4604,68.7041) .. controls (108.1812,68.7041) and (108.7655,68.1198) ..(108.7655,67.3990);
\path[fill=chfourfig5c0,fill opacity=0.5,draw=chfourfig5c0,line width=0.3985pt,draw opacity=0.5] (155.9891,59.5185) .. controls (155.9891,58.7977) and (155.4048,58.2133) ..(154.6840,58.2133) .. controls (153.9632,58.2133) and (153.3788,58.7977) ..(153.3788,59.5185) .. controls (153.3788,60.2393) and (153.9632,60.8236) ..(154.6840,60.8236) .. controls (155.4048,60.8236) and (155.9891,60.2393) ..(155.9891,59.5185);
\path[fill=chfourfig5c0,fill opacity=0.5,draw=chfourfig5c0,line width=0.3985pt,draw opacity=0.5] (94.9472,81.0978) .. controls (94.9472,80.3770) and (94.3629,79.7927) ..(93.6421,79.7927) .. controls (92.9213,79.7927) and (92.3369,80.3770) ..(92.3369,81.0978) .. controls (92.3369,81.8186) and (92.9213,82.4029) ..(93.6421,82.4029) .. controls (94.3629,82.4029) and (94.9472,81.8186) ..(94.9472,81.0978);
\path[fill=chfourfig5c0,fill opacity=0.5,draw=chfourfig5c0,line width=0.3985pt,draw opacity=0.5] (80.9594,75.6183) .. controls (80.9594,74.8975) and (80.3752,74.3132) ..(79.6543,74.3132) .. controls (78.9335,74.3132) and (78.3492,74.8975) ..(78.3492,75.6183) .. controls (78.3492,76.3391) and (78.9335,76.9234) ..(79.6543,76.9234) .. controls (80.3752,76.9234) and (80.9594,76.3391) ..(80.9594,75.6183);
\path[fill=chfourfig5c0,fill opacity=0.5,draw=chfourfig5c0,line width=0.3985pt,draw opacity=0.5] (92.7354,92.5351) .. controls (92.7354,91.8143) and (92.1511,91.2300) ..(91.4303,91.2300) .. controls (90.7095,91.2300) and (90.1252,91.8143) ..(90.1252,92.5351) .. controls (90.1252,93.2559) and (90.7095,93.8402) ..(91.4303,93.8402) .. controls (92.1511,93.8402) and (92.7354,93.2559) ..(92.7354,92.5351);
\path[fill=chfourfig5c0,fill opacity=0.5,draw=chfourfig5c0,line width=0.3985pt,draw opacity=0.5] (81.1388,91.0506) .. controls (81.1388,90.3298) and (80.5545,89.7455) ..(79.8337,89.7455) .. controls (79.1129,89.7455) and (78.5286,90.3298) ..(78.5286,91.0506) .. controls (78.5286,91.7714) and (79.1129,92.3557) ..(79.8337,92.3557) .. controls (80.5545,92.3557) and (81.1388,91.7714) ..(81.1388,91.0506);
\path[fill=chfourfig5c0,fill opacity=0.5,draw=chfourfig5c0,line width=0.3985pt,draw opacity=0.5] (92.3071,65.1674) .. controls (92.3071,64.4466) and (91.7227,63.8622) ..(91.0019,63.8622) .. controls (90.2811,63.8622) and (89.6968,64.4466) ..(89.6968,65.1674) .. controls (89.6968,65.8882) and (90.2811,66.4725) ..(91.0019,66.4725) .. controls (91.7227,66.4725) and (92.3071,65.8882) ..(92.3071,65.1674);
\path[fill=chfourfig5c0,fill opacity=0.5,draw=chfourfig5c0,line width=0.3985pt,draw opacity=0.5] (87.3456,92.6447) .. controls (87.3456,91.9239) and (86.7613,91.3395) ..(86.0405,91.3395) .. controls (85.3197,91.3395) and (84.7354,91.9239) ..(84.7354,92.6447) .. controls (84.7354,93.3655) and (85.3197,93.9498) ..(86.0405,93.9498) .. controls (86.7613,93.9498) and (87.3456,93.3655) ..(87.3456,92.6447);
\path[fill=chfourfig5c0,fill opacity=0.5,draw=chfourfig5c0,line width=0.3985pt,draw opacity=0.5] (95.1763,80.3805) .. controls (95.1763,79.6597) and (94.5920,79.0754) ..(93.8712,79.0754) .. controls (93.1504,79.0754) and (92.5661,79.6597) ..(92.5661,80.3805) .. controls (92.5661,81.1013) and (93.1504,81.6856) ..(93.8712,81.6856) .. controls (94.5920,81.6856) and (95.1763,81.1013) ..(95.1763,80.3805);
\path[fill=chfourfig5c0,fill opacity=0.5,draw=chfourfig5c0,line width=0.3985pt,draw opacity=0.5] (130.9427,58.7912) .. controls (130.9427,58.0704) and (130.3584,57.4861) ..(129.6375,57.4861) .. controls (128.9167,57.4861) and (128.3324,58.0704) ..(128.3324,58.7912) .. controls (128.3324,59.5120) and (128.9167,60.0963) ..(129.6375,60.0963) .. controls (130.3584,60.0963) and (130.9427,59.5120) ..(130.9427,58.7912);
\path[fill=chfourfig5c0,fill opacity=0.5,draw=chfourfig5c0,line width=0.3985pt,draw opacity=0.5] (132.5068,62.3778) .. controls (132.5068,61.6570) and (131.9225,61.0727) ..(131.2017,61.0727) .. controls (130.4809,61.0727) and (129.8966,61.6570) ..(129.8966,62.3778) .. controls (129.8966,63.0986) and (130.4809,63.6829) ..(131.2017,63.6829) .. controls (131.9225,63.6829) and (132.5068,63.0986) ..(132.5068,62.3778);
\path[fill=chfourfig5c0,fill opacity=0.5,draw=chfourfig5c0,line width=0.3985pt,draw opacity=0.5] (83.4203,73.0379) .. controls (83.4203,72.3171) and (82.8360,71.7328) ..(82.1151,71.7328) .. controls (81.3943,71.7328) and (80.8100,72.3171) ..(80.8100,73.0379) .. controls (80.8100,73.7587) and (81.3943,74.3431) ..(82.1151,74.3431) .. controls (82.8360,74.3431) and (83.4203,73.7587) ..(83.4203,73.0379);
\path[fill=chfourfig5c0,fill opacity=0.5,draw=chfourfig5c0,line width=0.3985pt,draw opacity=0.5] (65.8260,90.0942) .. controls (65.8260,89.3734) and (65.2417,88.7891) ..(64.5209,88.7891) .. controls (63.8001,88.7891) and (63.2158,89.3734) ..(63.2158,90.0942) .. controls (63.2158,90.8150) and (63.8001,91.3993) ..(64.5209,91.3993) .. controls (65.2417,91.3993) and (65.8260,90.8150) ..(65.8260,90.0942);
\path[fill=chfourfig5c0,fill opacity=0.5,draw=chfourfig5c0,line width=0.3985pt,draw opacity=0.5] (70.9070,108.8142) .. controls (70.9070,108.0934) and (70.3227,107.5091) ..(69.6019,107.5091) .. controls (68.8811,107.5091) and (68.2968,108.0934) ..(68.2968,108.8142) .. controls (68.2968,109.5350) and (68.8811,110.1194) ..(69.6019,110.1194) .. controls (70.3227,110.1194) and (70.9070,109.5350) ..(70.9070,108.8142);
\path[fill=chfourfig5c0,fill opacity=0.5,draw=chfourfig5c0,line width=0.3985pt,draw opacity=0.5] (79.7540,102.9761) .. controls (79.7540,102.2553) and (79.1697,101.6709) ..(78.4489,101.6709) .. controls (77.7281,101.6709) and (77.1437,102.2553) ..(77.1437,102.9761) .. controls (77.1437,103.6969) and (77.7281,104.2812) ..(78.4489,104.2812) .. controls (79.1697,104.2812) and (79.7540,103.6969) ..(79.7540,102.9761);
\path[fill=chfourfig5c0,fill opacity=0.5,draw=chfourfig5c0,line width=0.3985pt,draw opacity=0.5] (116.9749,78.6470) .. controls (116.9749,77.9262) and (116.3906,77.3419) ..(115.6697,77.3419) .. controls (114.9489,77.3419) and (114.3646,77.9262) ..(114.3646,78.6470) .. controls (114.3646,79.3678) and (114.9489,79.9521) ..(115.6697,79.9521) .. controls (116.3906,79.9521) and (116.9749,79.3678) ..(116.9749,78.6470);
\path[fill=chfourfig5c0,fill opacity=0.5,draw=chfourfig5c0,line width=0.3985pt,draw opacity=0.5] (100.9747,71.0653) .. controls (100.9747,70.3445) and (100.3904,69.7602) ..(99.6696,69.7602) .. controls (98.9487,69.7602) and (98.3644,70.3445) ..(98.3644,71.0653) .. controls (98.3644,71.7861) and (98.9487,72.3704) ..(99.6696,72.3704) .. controls (100.3904,72.3704) and (100.9747,71.7861) ..(100.9747,71.0653);
\path[fill=chfourfig5c0,fill opacity=0.5,draw=chfourfig5c0,line width=0.3985pt,draw opacity=0.5] (67.5297,108.9936) .. controls (67.5297,108.2728) and (66.9453,107.6884) ..(66.2245,107.6884) .. controls (65.5037,107.6884) and (64.9194,108.2728) ..(64.9194,108.9936) .. controls (64.9194,109.7144) and (65.5037,110.2987) ..(66.2245,110.2987) .. controls (66.9453,110.2987) and (67.5297,109.7144) ..(67.5297,108.9936);
\path[fill=chfourfig5c0,fill opacity=0.5,draw=chfourfig5c0,line width=0.3985pt,draw opacity=0.5] (52.3165,81.9347) .. controls (52.3165,81.2139) and (51.7322,80.6296) ..(51.0114,80.6296) .. controls (50.2906,80.6296) and (49.7063,81.2139) ..(49.7063,81.9347) .. controls (49.7063,82.6555) and (50.2906,83.2398) ..(51.0114,83.2398) .. controls (51.7322,83.2398) and (52.3165,82.6555) ..(52.3165,81.9347);
\path[fill=chfourfig5c0,fill opacity=0.5,draw=chfourfig5c0,line width=0.3985pt,draw opacity=0.5] (58.9218,90.3931) .. controls (58.9218,89.6723) and (58.3375,89.0880) ..(57.6167,89.0880) .. controls (56.8959,89.0880) and (56.3116,89.6723) ..(56.3116,90.3931) .. controls (56.3116,91.1139) and (56.8959,91.6982) ..(57.6167,91.6982) .. controls (58.3375,91.6982) and (58.9218,91.1139) ..(58.9218,90.3931);
\path[fill=chfourfig5c0,fill opacity=0.5,draw=chfourfig5c0,line width=0.3985pt,draw opacity=0.5] (118.2401,55.5234) .. controls (118.2401,54.8026) and (117.6558,54.2183) ..(116.9350,54.2183) .. controls (116.2142,54.2183) and (115.6299,54.8026) ..(115.6299,55.5234) .. controls (115.6299,56.2442) and (116.2142,56.8285) ..(116.9350,56.8285) .. controls (117.6558,56.8285) and (118.2401,56.2442) ..(118.2401,55.5234);
\path[fill=chfourfig5c0,fill opacity=0.5,draw=chfourfig5c0,line width=0.3985pt,draw opacity=0.5] (140.2977,59.7077) .. controls (140.2977,58.9870) and (139.7134,58.4026) ..(138.9926,58.4026) .. controls (138.2718,58.4026) and (137.6875,58.9870) ..(137.6875,59.7077) .. controls (137.6875,60.4286) and (138.2718,61.0129) ..(138.9926,61.0129) .. controls (139.7134,61.0129) and (140.2977,60.4286) ..(140.2977,59.7077);
\path[fill=chfourfig5c0,fill opacity=0.5,draw=chfourfig5c0,line width=0.3985pt,draw opacity=0.5] (64.8796,79.7429) .. controls (64.8796,79.0221) and (64.2952,78.4378) ..(63.5744,78.4378) .. controls (62.8536,78.4378) and (62.2693,79.0221) ..(62.2693,79.7429) .. controls (62.2693,80.4637) and (62.8536,81.0480) ..(63.5744,81.0480) .. controls (64.2952,81.0480) and (64.8796,80.4637) ..(64.8796,79.7429);
\path[fill=chfourfig5c0,fill opacity=0.5,draw=chfourfig5c0,line width=0.3985pt,draw opacity=0.5] (130.3150,56.0215) .. controls (130.3150,55.3007) and (129.7307,54.7164) ..(129.0099,54.7164) .. controls (128.2891,54.7164) and (127.7048,55.3007) ..(127.7048,56.0215) .. controls (127.7048,56.7423) and (128.2891,57.3266) ..(129.0099,57.3266) .. controls (129.7307,57.3266) and (130.3150,56.7423) ..(130.3150,56.0215);
\path[fill=chfourfig5c0,fill opacity=0.5,draw=chfourfig5c0,line width=0.3985pt,draw opacity=0.5] (65.8160,107.0907) .. controls (65.8160,106.3699) and (65.2317,105.7856) ..(64.5109,105.7856) .. controls (63.7901,105.7856) and (63.2058,106.3699) ..(63.2058,107.0907) .. controls (63.2058,107.8115) and (63.7901,108.3958) ..(64.5109,108.3958) .. controls (65.2317,108.3958) and (65.8160,107.8115) ..(65.8160,107.0907);
\path[fill=chfourfig5c0,fill opacity=0.5,draw=chfourfig5c0,line width=0.3985pt,draw opacity=0.5] (119.9139,51.2095) .. controls (119.9139,50.4887) and (119.3296,49.9044) ..(118.6088,49.9044) .. controls (117.8880,49.9044) and (117.3036,50.4887) ..(117.3036,51.2095) .. controls (117.3036,51.9303) and (117.8880,52.5146) ..(118.6088,52.5146) .. controls (119.3296,52.5146) and (119.9139,51.9303) ..(119.9139,51.2095);
\path[fill=chfourfig5c0,fill opacity=0.5,draw=chfourfig5c0,line width=0.3985pt,draw opacity=0.5] (147.1122,59.8273) .. controls (147.1122,59.1065) and (146.5279,58.5222) ..(145.8071,58.5222) .. controls (145.0863,58.5222) and (144.5020,59.1065) ..(144.5020,59.8273) .. controls (144.5020,60.5481) and (145.0863,61.1324) ..(145.8071,61.1324) .. controls (146.5279,61.1324) and (147.1122,60.5481) ..(147.1122,59.8273);
\path[fill=chfourfig5c0,fill opacity=0.5,draw=chfourfig5c0,line width=0.3985pt,draw opacity=0.5] (153.8471,34.3127) .. controls (153.8471,33.5919) and (153.2628,33.0075) ..(152.5420,33.0075) .. controls (151.8212,33.0075) and (151.2368,33.5919) ..(151.2368,34.3127) .. controls (151.2368,35.0335) and (151.8212,35.6178) ..(152.5420,35.6178) .. controls (153.2628,35.6178) and (153.8471,35.0335) ..(153.8471,34.3127);
\path[fill=chfourfig5c0,fill opacity=0.5,draw=chfourfig5c0,line width=0.3985pt,draw opacity=0.5] (77.1836,92.9934) .. controls (77.1836,92.2726) and (76.5992,91.6882) ..(75.8784,91.6882) .. controls (75.1577,91.6882) and (74.5733,92.2726) ..(74.5733,92.9934) .. controls (74.5733,93.7142) and (75.1577,94.2985) ..(75.8784,94.2985) .. controls (76.5992,94.2985) and (77.1836,93.7142) ..(77.1836,92.9934);
\path[fill=chfourfig5c0,fill opacity=0.5,draw=chfourfig5c0,line width=0.3985pt,draw opacity=0.5] (132.3773,40.4796) .. controls (132.3773,39.7588) and (131.7930,39.1745) ..(131.0722,39.1745) .. controls (130.3514,39.1745) and (129.7671,39.7588) ..(129.7671,40.4796) .. controls (129.7671,41.2004) and (130.3514,41.7848) ..(131.0722,41.7848) .. controls (131.7930,41.7848) and (132.3773,41.2004) ..(132.3773,40.4796);
\path[fill=chfourfig5c0,fill opacity=0.5,draw=chfourfig5c0,line width=0.3985pt,draw opacity=0.5] (154.2057,47.1248) .. controls (154.2057,46.4040) and (153.6214,45.8197) ..(152.9006,45.8197) .. controls (152.1798,45.8197) and (151.5955,46.4040) ..(151.5955,47.1248) .. controls (151.5955,47.8456) and (152.1798,48.4299) ..(152.9006,48.4299) .. controls (153.6214,48.4299) and (154.2057,47.8456) ..(154.2057,47.1248);
\path[fill=chfourfig5c0,fill opacity=0.5,draw=chfourfig5c0,line width=0.3985pt,draw opacity=0.5] (122.9326,67.9669) .. controls (122.9326,67.2461) and (122.3483,66.6618) ..(121.6275,66.6618) .. controls (120.9067,66.6618) and (120.3224,67.2461) ..(120.3224,67.9669) .. controls (120.3224,68.6877) and (120.9067,69.2720) ..(121.6275,69.2720) .. controls (122.3483,69.2720) and (122.9326,68.6877) ..(122.9326,67.9669);
\path[fill=chfourfig5c0,fill opacity=0.5,draw=chfourfig5c0,line width=0.3985pt,draw opacity=0.5] (155.8596,41.1670) .. controls (155.8596,40.4463) and (155.2752,39.8619) ..(154.5544,39.8619) .. controls (153.8336,39.8619) and (153.2493,40.4463) ..(153.2493,41.1670) .. controls (153.2493,41.8878) and (153.8336,42.4722) ..(154.5544,42.4722) .. controls (155.2752,42.4722) and (155.8596,41.8878) ..(155.8596,41.1670);
\path[fill=chfourfig5c0,fill opacity=0.5,draw=chfourfig5c0,line width=0.3985pt,draw opacity=0.5] (121.4382,55.4636) .. controls (121.4382,54.7428) and (120.8539,54.1585) ..(120.1331,54.1585) .. controls (119.4123,54.1585) and (118.8279,54.7428) ..(118.8279,55.4636) .. controls (118.8279,56.1844) and (119.4123,56.7687) ..(120.1331,56.7687) .. controls (120.8539,56.7687) and (121.4382,56.1844) ..(121.4382,55.4636);
\path[fill=chfourfig5c0,fill opacity=0.5,draw=chfourfig5c0,line width=0.3985pt,draw opacity=0.5] (66.4337,84.6047) .. controls (66.4337,83.8839) and (65.8494,83.2996) ..(65.1286,83.2996) .. controls (64.4078,83.2996) and (63.8235,83.8839) ..(63.8235,84.6047) .. controls (63.8235,85.3255) and (64.4078,85.9098) ..(65.1286,85.9098) .. controls (65.8494,85.9098) and (66.4337,85.3255) ..(66.4337,84.6047);
\path[fill=chfourfig5c0,fill opacity=0.5,draw=chfourfig5c0,line width=0.3985pt,draw opacity=0.5] (74.3741,70.7963) .. controls (74.3741,70.0755) and (73.7897,69.4912) ..(73.0689,69.4912) .. controls (72.3481,69.4912) and (71.7638,70.0755) ..(71.7638,70.7963) .. controls (71.7638,71.5171) and (72.3481,72.1014) ..(73.0689,72.1014) .. controls (73.7897,72.1014) and (74.3741,71.5171) ..(74.3741,70.7963);
\path[fill=chfourfig5c0,fill opacity=0.5,draw=chfourfig5c0,line width=0.3985pt,draw opacity=0.5] (57.7462,98.3733) .. controls (57.7462,97.6525) and (57.1619,97.0681) ..(56.4411,97.0681) .. controls (55.7203,97.0681) and (55.1360,97.6525) ..(55.1360,98.3733) .. controls (55.1360,99.0941) and (55.7203,99.6784) ..(56.4411,99.6784) .. controls (57.1619,99.6784) and (57.7462,99.0941) ..(57.7462,98.3733);
\path[fill=chfourfig5c0,fill opacity=0.5,draw=chfourfig5c0,line width=0.3985pt,draw opacity=0.5] (95.2959,71.0852) .. controls (95.2959,70.3644) and (94.7115,69.7801) ..(93.9908,69.7801) .. controls (93.2700,69.7801) and (92.6856,70.3644) ..(92.6856,71.0852) .. controls (92.6856,71.8060) and (93.2700,72.3904) ..(93.9908,72.3904) .. controls (94.7115,72.3904) and (95.2959,71.8060) ..(95.2959,71.0852);
\path[fill=chfourfig5c0,fill opacity=0.5,draw=chfourfig5c0,line width=0.3985pt,draw opacity=0.5] (108.6859,87.1253) .. controls (108.6859,86.4045) and (108.1015,85.8202) ..(107.3807,85.8202) .. controls (106.6599,85.8202) and (106.0756,86.4045) ..(106.0756,87.1253) .. controls (106.0756,87.8461) and (106.6599,88.4304) ..(107.3807,88.4304) .. controls (108.1015,88.4304) and (108.6859,87.8461) ..(108.6859,87.1253);
\path[fill=chfourfig5c0,fill opacity=0.5,draw=chfourfig5c0,line width=0.3985pt,draw opacity=0.5] (69.5023,73.6556) .. controls (69.5023,72.9348) and (68.9180,72.3505) ..(68.1971,72.3505) .. controls (67.4763,72.3505) and (66.8920,72.9348) ..(66.8920,73.6556) .. controls (66.8920,74.3764) and (67.4763,74.9608) ..(68.1971,74.9608) .. controls (68.9180,74.9608) and (69.5023,74.3764) ..(69.5023,73.6556);
\path[fill=chfourfig5c0,fill opacity=0.5,draw=chfourfig5c0,line width=0.3985pt,draw opacity=0.5] (121.5876,60.1362) .. controls (121.5876,59.4154) and (121.0033,58.8310) ..(120.2825,58.8310) .. controls (119.5617,58.8310) and (118.9774,59.4154) ..(118.9774,60.1362) .. controls (118.9774,60.8570) and (119.5617,61.4413) ..(120.2825,61.4413) .. controls (121.0033,61.4413) and (121.5876,60.8570) ..(121.5876,60.1362);
\path[fill=chfourfig5c0,fill opacity=0.5,draw=chfourfig5c0,line width=0.3985pt,draw opacity=0.5] (144.0636,41.3763) .. controls (144.0636,40.6555) and (143.4793,40.0711) ..(142.7585,40.0711) .. controls (142.0377,40.0711) and (141.4534,40.6555) ..(141.4534,41.3763) .. controls (141.4534,42.0971) and (142.0377,42.6814) ..(142.7585,42.6814) .. controls (143.4793,42.6814) and (144.0636,42.0971) ..(144.0636,41.3763);
\path[fill=chfourfig5c0,fill opacity=0.5,draw=chfourfig5c0,line width=0.3985pt,draw opacity=0.5] (123.8591,43.5581) .. controls (123.8591,42.8373) and (123.2748,42.2530) ..(122.5540,42.2530) .. controls (121.8332,42.2530) and (121.2489,42.8373) ..(121.2489,43.5581) .. controls (121.2489,44.2789) and (121.8332,44.8632) ..(122.5540,44.8632) .. controls (123.2748,44.8632) and (123.8591,44.2789) ..(123.8591,43.5581);
\path[fill=chfourfig5c0,fill opacity=0.5,draw=chfourfig5c0,line width=0.3985pt,draw opacity=0.5] (53.4224,91.5787) .. controls (53.4224,90.8578) and (52.8381,90.2735) ..(52.1172,90.2735) .. controls (51.3965,90.2735) and (50.8121,90.8578) ..(50.8121,91.5787) .. controls (50.8121,92.2995) and (51.3965,92.8838) ..(52.1172,92.8838) .. controls (52.8381,92.8838) and (53.4224,92.2995) ..(53.4224,91.5787);
\path[fill=chfourfig5c0,fill opacity=0.5,draw=chfourfig5c0,line width=0.3985pt,draw opacity=0.5] (90.6034,78.7466) .. controls (90.6034,78.0258) and (90.0191,77.4415) ..(89.2983,77.4415) .. controls (88.5775,77.4415) and (87.9932,78.0258) ..(87.9932,78.7466) .. controls (87.9932,79.4674) and (88.5775,80.0517) ..(89.2983,80.0517) .. controls (90.0191,80.0517) and (90.6034,79.4674) ..(90.6034,78.7466);
\path[fill=chfourfig5c0,fill opacity=0.5,draw=chfourfig5c0,line width=0.3985pt,draw opacity=0.5] (97.4379,61.7402) .. controls (97.4379,61.0194) and (96.8536,60.4350) ..(96.1328,60.4350) .. controls (95.4120,60.4350) and (94.8277,61.0194) ..(94.8277,61.7402) .. controls (94.8277,62.4610) and (95.4120,63.0453) ..(96.1328,63.0453) .. controls (96.8536,63.0453) and (97.4379,62.4610) ..(97.4379,61.7402);
\path[fill=chfourfig5c0,fill opacity=0.5,draw=chfourfig5c0,line width=0.3985pt,draw opacity=0.5] (130.3449,74.4327) .. controls (130.3449,73.7119) and (129.7606,73.1276) ..(129.0398,73.1276) .. controls (128.3190,73.1276) and (127.7347,73.7119) ..(127.7347,74.4327) .. controls (127.7347,75.1535) and (128.3190,75.7378) ..(129.0398,75.7378) .. controls (129.7606,75.7378) and (130.3449,75.1535) ..(130.3449,74.4327);
\path[fill=chfourfig5c0,fill opacity=0.5,draw=chfourfig5c0,line width=0.3985pt,draw opacity=0.5] (136.1133,50.7911) .. controls (136.1133,50.0703) and (135.5290,49.4860) ..(134.8082,49.4860) .. controls (134.0874,49.4860) and (133.5031,50.0703) ..(133.5031,50.7911) .. controls (133.5031,51.5119) and (134.0874,52.0962) ..(134.8082,52.0962) .. controls (135.5290,52.0962) and (136.1133,51.5119) ..(136.1133,50.7911);
\path[fill=chfourfig5c0,fill opacity=0.5,draw=chfourfig5c0,line width=0.3985pt,draw opacity=0.5] (69.1436,106.8615) .. controls (69.1436,106.1407) and (68.5593,105.5564) ..(67.8385,105.5564) .. controls (67.1177,105.5564) and (66.5334,106.1407) ..(66.5334,106.8615) .. controls (66.5334,107.5823) and (67.1177,108.1667) ..(67.8385,108.1667) .. controls (68.5593,108.1667) and (69.1436,107.5823) ..(69.1436,106.8615);
\path[fill=chfourfig5c0,fill opacity=0.5,draw=chfourfig5c0,line width=0.3985pt,draw opacity=0.5] (104.4716,80.4203) .. controls (104.4716,79.6996) and (103.8873,79.1152) ..(103.1665,79.1152) .. controls (102.4457,79.1152) and (101.8614,79.6996) ..(101.8614,80.4203) .. controls (101.8614,81.1412) and (102.4457,81.7255) ..(103.1665,81.7255) .. controls (103.8873,81.7255) and (104.4716,81.1412) ..(104.4716,80.4203);
\path[fill=chfourfig5c0,fill opacity=0.5,draw=chfourfig5c0,line width=0.3985pt,draw opacity=0.5] (143.8843,36.7037) .. controls (143.8843,35.9829) and (143.3000,35.3986) ..(142.5792,35.3986) .. controls (141.8584,35.3986) and (141.2741,35.9829) ..(141.2741,36.7037) .. controls (141.2741,37.4245) and (141.8584,38.0089) ..(142.5792,38.0089) .. controls (143.3000,38.0089) and (143.8843,37.4245) ..(143.8843,36.7037);
\path[fill=chfourfig5c0,fill opacity=0.5,draw=chfourfig5c0,line width=0.3985pt,draw opacity=0.5] (54.7972,89.8451) .. controls (54.7972,89.1243) and (54.2129,88.5400) ..(53.4921,88.5400) .. controls (52.7713,88.5400) and (52.1870,89.1243) ..(52.1870,89.8451) .. controls (52.1870,90.5659) and (52.7713,91.1503) ..(53.4921,91.1503) .. controls (54.2129,91.1503) and (54.7972,90.5659) ..(54.7972,89.8451);
\path[fill=chfourfig5c0,fill opacity=0.5,draw=chfourfig5c0,line width=0.3985pt,draw opacity=0.5] (100.4666,58.4026) .. controls (100.4666,57.6818) and (99.8822,57.0975) ..(99.1614,57.0975) .. controls (98.4406,57.0975) and (97.8563,57.6818) ..(97.8563,58.4026) .. controls (97.8563,59.1234) and (98.4406,59.7077) ..(99.1614,59.7077) .. controls (99.8822,59.7077) and (100.4666,59.1234) ..(100.4666,58.4026);
\path[fill=chfourfig5c0,fill opacity=0.5,draw=chfourfig5c0,line width=0.3985pt,draw opacity=0.5] (126.0111,67.2296) .. controls (126.0111,66.5089) and (125.4268,65.9245) ..(124.7060,65.9245) .. controls (123.9852,65.9245) and (123.4008,66.5089) ..(123.4008,67.2296) .. controls (123.4008,67.9504) and (123.9852,68.5348) ..(124.7060,68.5348) .. controls (125.4268,68.5348) and (126.0111,67.9504) ..(126.0111,67.2296);
\path[fill=chfourfig5c0,fill opacity=0.5,draw=chfourfig5c0,line width=0.3985pt,draw opacity=0.5] (95.6047,88.7492) .. controls (95.6047,88.0284) and (95.0204,87.4441) ..(94.2996,87.4441) .. controls (93.5788,87.4441) and (92.9945,88.0284) ..(92.9945,88.7492) .. controls (92.9945,89.4700) and (93.5788,90.0544) ..(94.2996,90.0544) .. controls (95.0204,90.0544) and (95.6047,89.4700) ..(95.6047,88.7492);
\path[fill=chfourfig5c0,fill opacity=0.5,draw=chfourfig5c0,line width=0.3985pt,draw opacity=0.5] (71.1960,85.8999) .. controls (71.1960,85.1791) and (70.6116,84.5948) ..(69.8908,84.5948) .. controls (69.1700,84.5948) and (68.5857,85.1791) ..(68.5857,85.8999) .. controls (68.5857,86.6207) and (69.1700,87.2050) ..(69.8908,87.2050) .. controls (70.6116,87.2050) and (71.1960,86.6207) ..(71.1960,85.8999);
\path[fill=chfourfig5c0,fill opacity=0.5,draw=chfourfig5c0,line width=0.3985pt,draw opacity=0.5] (55.5444,89.4566) .. controls (55.5444,88.7358) and (54.9601,88.1515) ..(54.2393,88.1515) .. controls (53.5185,88.1515) and (52.9342,88.7358) ..(52.9342,89.4566) .. controls (52.9342,90.1774) and (53.5185,90.7617) ..(54.2393,90.7617) .. controls (54.9601,90.7617) and (55.5444,90.1774) ..(55.5444,89.4566);
\path[fill=chfourfig5c0,fill opacity=0.5,draw=chfourfig5c0,line width=0.3985pt,draw opacity=0.5] (59.1908,110.0098) .. controls (59.1908,109.2890) and (58.6065,108.7047) ..(57.8857,108.7047) .. controls (57.1649,108.7047) and (56.5806,109.2890) ..(56.5806,110.0098) .. controls (56.5806,110.7306) and (57.1649,111.3149) ..(57.8857,111.3149) .. controls (58.6065,111.3149) and (59.1908,110.7306) ..(59.1908,110.0098);
\path[fill=chfourfig5c0,fill opacity=0.5,draw=chfourfig5c0,line width=0.3985pt,draw opacity=0.5] (77.5422,69.4214) .. controls (77.5422,68.7007) and (76.9579,68.1163) ..(76.2371,68.1163) .. controls (75.5163,68.1163) and (74.9320,68.7007) ..(74.9320,69.4214) .. controls (74.9320,70.1423) and (75.5163,70.7266) ..(76.2371,70.7266) .. controls (76.9579,70.7266) and (77.5422,70.1423) ..(77.5422,69.4214);
\path[fill=chfourfig5c0,fill opacity=0.5,draw=chfourfig5c0,line width=0.3985pt,draw opacity=0.5] (54.4585,88.1116) .. controls (54.4585,87.3908) and (53.8742,86.8065) ..(53.1534,86.8065) .. controls (52.4326,86.8065) and (51.8483,87.3908) ..(51.8483,88.1116) .. controls (51.8483,88.8324) and (52.4326,89.4167) ..(53.1534,89.4167) .. controls (53.8742,89.4167) and (54.4585,88.8324) ..(54.4585,88.1116);
\path[fill=chfourfig5c0,fill opacity=0.5,draw=chfourfig5c0,line width=0.3985pt,draw opacity=0.5] (146.3451,46.1185) .. controls (146.3451,45.3977) and (145.7608,44.8134) ..(145.0400,44.8134) .. controls (144.3192,44.8134) and (143.7349,45.3977) ..(143.7349,46.1185) .. controls (143.7349,46.8394) and (144.3192,47.4237) ..(145.0400,47.4237) .. controls (145.7608,47.4237) and (146.3451,46.8394) ..(146.3451,46.1185);
\path[fill=chfourfig5c0,fill opacity=0.5,draw=chfourfig5c0,line width=0.3985pt,draw opacity=0.5] (158.0514,43.9666) .. controls (158.0514,43.2458) and (157.4670,42.6615) ..(156.7462,42.6615) .. controls (156.0255,42.6615) and (155.4411,43.2458) ..(155.4411,43.9666) .. controls (155.4411,44.6874) and (156.0255,45.2717) ..(156.7462,45.2717) .. controls (157.4670,45.2717) and (158.0514,44.6874) ..(158.0514,43.9666);
\path[fill=chfourfig5c0,fill opacity=0.5,draw=chfourfig5c0,line width=0.3985pt,draw opacity=0.5] (66.0053,97.6659) .. controls (66.0053,96.9451) and (65.4210,96.3608) ..(64.7002,96.3608) .. controls (63.9794,96.3608) and (63.3951,96.9451) ..(63.3951,97.6659) .. controls (63.3951,98.3867) and (63.9794,98.9710) ..(64.7002,98.9710) .. controls (65.4210,98.9710) and (66.0053,98.3867) ..(66.0053,97.6659);
\path[fill=chfourfig5c0,fill opacity=0.5,draw=chfourfig5c0,line width=0.3985pt,draw opacity=0.5] (64.2918,111.6736) .. controls (64.2918,110.9528) and (63.7074,110.3684) ..(62.9866,110.3684) .. controls (62.2658,110.3684) and (61.6815,110.9528) ..(61.6815,111.6736) .. controls (61.6815,112.3944) and (62.2658,112.9787) ..(62.9866,112.9787) .. controls (63.7074,112.9787) and (64.2918,112.3944) ..(64.2918,111.6736);
\path[fill=chfourfig5c0,fill opacity=0.5,draw=chfourfig5c0,line width=0.3985pt,draw opacity=0.5] (124.1481,61.3018) .. controls (124.1481,60.5810) and (123.5637,59.9967) ..(122.8430,59.9967) .. controls (122.1221,59.9967) and (121.5378,60.5810) ..(121.5378,61.3018) .. controls (121.5378,62.0226) and (122.1221,62.6069) ..(122.8430,62.6069) .. controls (123.5637,62.6069) and (124.1481,62.0226) ..(124.1481,61.3018);
\path[fill=chfourfig5c0,fill opacity=0.5,draw=chfourfig5c0,line width=0.3985pt,draw opacity=0.5] (118.6586,53.6803) .. controls (118.6586,52.9595) and (118.0743,52.3752) ..(117.3535,52.3752) .. controls (116.6327,52.3752) and (116.0483,52.9595) ..(116.0483,53.6803) .. controls (116.0483,54.4011) and (116.6327,54.9854) ..(117.3535,54.9854) .. controls (118.0743,54.9854) and (118.6586,54.4011) ..(118.6586,53.6803);
\path[fill=chfourfig5c0,fill opacity=0.5,draw=chfourfig5c0,line width=0.3985pt,draw opacity=0.5] (121.6075,52.5943) .. controls (121.6075,51.8735) and (121.0232,51.2892) ..(120.3024,51.2892) .. controls (119.5816,51.2892) and (118.9973,51.8735) ..(118.9973,52.5943) .. controls (118.9973,53.3151) and (119.5816,53.8995) ..(120.3024,53.8995) .. controls (121.0232,53.8995) and (121.6075,53.3151) ..(121.6075,52.5943);
\path[fill=chfourfig5c0,fill opacity=0.5,draw=chfourfig5c0,line width=0.3985pt,draw opacity=0.5] (124.1580,48.3900) .. controls (124.1580,47.6693) and (123.5737,47.0849) ..(122.8529,47.0849) .. controls (122.1321,47.0849) and (121.5478,47.6693) ..(121.5478,48.3900) .. controls (121.5478,49.1108) and (122.1321,49.6952) ..(122.8529,49.6952) .. controls (123.5737,49.6952) and (124.1580,49.1108) ..(124.1580,48.3900);
\path[fill=chfourfig5c0,fill opacity=0.5,draw=chfourfig5c0,line width=0.3985pt,draw opacity=0.5] (149.6229,38.8258) .. controls (149.6229,38.1050) and (149.0385,37.5207) ..(148.3177,37.5207) .. controls (147.5969,37.5207) and (147.0126,38.1050) ..(147.0126,38.8258) .. controls (147.0126,39.5466) and (147.5969,40.1309) ..(148.3177,40.1309) .. controls (149.0385,40.1309) and (149.6229,39.5466) ..(149.6229,38.8258);
\path[fill=chfourfig5c0,fill opacity=0.5,draw=chfourfig5c0,line width=0.3985pt,draw opacity=0.5] (85.4128,76.9533) .. controls (85.4128,76.2325) and (84.8285,75.6482) ..(84.1077,75.6482) .. controls (83.3869,75.6482) and (82.8026,76.2325) ..(82.8026,76.9533) .. controls (82.8026,77.6741) and (83.3869,78.2584) ..(84.1077,78.2584) .. controls (84.8285,78.2584) and (85.4128,77.6741) ..(85.4128,76.9533);
\path[fill=chfourfig5c0,fill opacity=0.5,draw=chfourfig5c0,line width=0.3985pt,draw opacity=0.5] (111.0769,81.4465) .. controls (111.0769,80.7257) and (110.4926,80.1414) ..(109.7718,80.1414) .. controls (109.0510,80.1414) and (108.4667,80.7257) ..(108.4667,81.4465) .. controls (108.4667,82.1673) and (109.0510,82.7516) ..(109.7718,82.7516) .. controls (110.4926,82.7516) and (111.0769,82.1673) ..(111.0769,81.4465);
\path[fill=chfourfig5c0,fill opacity=0.5,draw=chfourfig5c0,line width=0.3985pt,draw opacity=0.5] (127.9738,63.3840) .. controls (127.9738,62.6632) and (127.3894,62.0789) ..(126.6686,62.0789) .. controls (125.9478,62.0789) and (125.3635,62.6632) ..(125.3635,63.3840) .. controls (125.3635,64.1048) and (125.9478,64.6891) ..(126.6686,64.6891) .. controls (127.3894,64.6891) and (127.9738,64.1048) ..(127.9738,63.3840);
\path[fill=chfourfig5c0,fill opacity=0.5,draw=chfourfig5c0,line width=0.3985pt,draw opacity=0.5] (128.1730,59.0403) .. controls (128.1730,58.3194) and (127.5887,57.7351) ..(126.8679,57.7351) .. controls (126.1471,57.7351) and (125.5628,58.3194) ..(125.5628,59.0403) .. controls (125.5628,59.7610) and (126.1471,60.3454) ..(126.8679,60.3454) .. controls (127.5887,60.3454) and (128.1730,59.7610) ..(128.1730,59.0403);
\path[fill=chfourfig5c0,fill opacity=0.5,draw=chfourfig5c0,line width=0.3985pt,draw opacity=0.5] (61.0140,93.6409) .. controls (61.0140,92.9201) and (60.4297,92.3358) ..(59.7089,92.3358) .. controls (58.9881,92.3358) and (58.4038,92.9201) ..(58.4038,93.6409) .. controls (58.4038,94.3617) and (58.9881,94.9461) ..(59.7089,94.9461) .. controls (60.4297,94.9461) and (61.0140,94.3617) ..(61.0140,93.6409);
\path[fill=chfourfig5c0,fill opacity=0.5,draw=chfourfig5c0,line width=0.3985pt,draw opacity=0.5] (110.7880,73.0579) .. controls (110.7880,72.3371) and (110.2037,71.7527) ..(109.4829,71.7527) .. controls (108.7621,71.7527) and (108.1777,72.3371) ..(108.1777,73.0579) .. controls (108.1777,73.7787) and (108.7621,74.3630) ..(109.4829,74.3630) .. controls (110.2037,74.3630) and (110.7880,73.7787) ..(110.7880,73.0579);
\path[fill=chfourfig5c0,fill opacity=0.5,draw=chfourfig5c0,line width=0.3985pt,draw opacity=0.5] (102.2599,59.5284) .. controls (102.2599,58.8076) and (101.6756,58.2233) ..(100.9547,58.2233) .. controls (100.2339,58.2233) and (99.6496,58.8076) ..(99.6496,59.5284) .. controls (99.6496,60.2492) and (100.2339,60.8335) ..(100.9547,60.8335) .. controls (101.6756,60.8335) and (102.2599,60.2492) ..(102.2599,59.5284);
\path[fill=chfourfig5c0,fill opacity=0.5,draw=chfourfig5c0,line width=0.3985pt,draw opacity=0.5] (147.3812,29.5704) .. controls (147.3812,28.8496) and (146.7969,28.2653) ..(146.0761,28.2653) .. controls (145.3553,28.2653) and (144.7710,28.8496) ..(144.7710,29.5704) .. controls (144.7710,30.2912) and (145.3553,30.8755) ..(146.0761,30.8755) .. controls (146.7969,30.8755) and (147.3812,30.2912) ..(147.3812,29.5704);
\path[fill=chfourfig5c0,fill opacity=0.5,draw=chfourfig5c0,line width=0.3985pt,draw opacity=0.5] (64.9991,111.3448) .. controls (64.9991,110.6240) and (64.4148,110.0397) ..(63.6940,110.0397) .. controls (62.9732,110.0397) and (62.3889,110.6240) ..(62.3889,111.3448) .. controls (62.3889,112.0656) and (62.9732,112.6499) ..(63.6940,112.6499) .. controls (64.4148,112.6499) and (64.9991,112.0656) ..(64.9991,111.3448);
\path[fill=chfourfig5c0,fill opacity=0.5,draw=chfourfig5c0,line width=0.3985pt,draw opacity=0.5] (138.6638,57.6355) .. controls (138.6638,56.9147) and (138.0795,56.3304) ..(137.3587,56.3304) .. controls (136.6379,56.3304) and (136.0536,56.9147) ..(136.0536,57.6355) .. controls (136.0536,58.3563) and (136.6379,58.9406) ..(137.3587,58.9406) .. controls (138.0795,58.9406) and (138.6638,58.3563) ..(138.6638,57.6355);
\path[fill=chfourfig5c0,fill opacity=0.5,draw=chfourfig5c0,line width=0.3985pt,draw opacity=0.5] (76.3865,98.2836) .. controls (76.3865,97.5628) and (75.8022,96.9785) ..(75.0814,96.9785) .. controls (74.3606,96.9785) and (73.7763,97.5628) ..(73.7763,98.2836) .. controls (73.7763,99.0044) and (74.3606,99.5887) ..(75.0814,99.5887) .. controls (75.8022,99.5887) and (76.3865,99.0044) ..(76.3865,98.2836);
\path[fill=chfourfig5c0,fill opacity=0.5,draw=chfourfig5c0,line width=0.3985pt,draw opacity=0.5] (108.3571,80.2410) .. controls (108.3571,79.5202) and (107.7728,78.9359) ..(107.0519,78.9359) .. controls (106.3311,78.9359) and (105.7468,79.5202) ..(105.7468,80.2410) .. controls (105.7468,80.9618) and (106.3311,81.5461) ..(107.0519,81.5461) .. controls (107.7728,81.5461) and (108.3571,80.9618) ..(108.3571,80.2410);
\path[fill=chfourfig5c0,fill opacity=0.5,draw=chfourfig5c0,line width=0.3985pt,draw opacity=0.5] (60.9044,106.7021) .. controls (60.9044,105.9813) and (60.3201,105.3970) ..(59.5993,105.3970) .. controls (58.8785,105.3970) and (58.2942,105.9813) ..(58.2942,106.7021) .. controls (58.2942,107.4229) and (58.8785,108.0072) ..(59.5993,108.0072) .. controls (60.3201,108.0072) and (60.9044,107.4229) ..(60.9044,106.7021);
\path[fill=chfourfig5c0,fill opacity=0.5,draw=chfourfig5c0,line width=0.3985pt,draw opacity=0.5] (89.0592,60.4251) .. controls (89.0592,59.7043) and (88.4749,59.1199) ..(87.7541,59.1199) .. controls (87.0333,59.1199) and (86.4489,59.7043) ..(86.4489,60.4251) .. controls (86.4489,61.1459) and (87.0333,61.7302) ..(87.7541,61.7302) .. controls (88.4749,61.7302) and (89.0592,61.1459) ..(89.0592,60.4251);
\path[fill=chfourfig5c0,fill opacity=0.5,draw=chfourfig5c0,line width=0.3985pt,draw opacity=0.5] (112.6510,58.5023) .. controls (112.6510,57.7815) and (112.0667,57.1971) ..(111.3459,57.1971) .. controls (110.6251,57.1971) and (110.0408,57.7815) ..(110.0408,58.5023) .. controls (110.0408,59.2231) and (110.6251,59.8074) ..(111.3459,59.8074) .. controls (112.0667,59.8074) and (112.6510,59.2231) ..(112.6510,58.5023);
\path[fill=chfourfig5c0,fill opacity=0.5,draw=chfourfig5c0,line width=0.3985pt,draw opacity=0.5] (88.3917,77.1625) .. controls (88.3917,76.4417) and (87.8074,75.8574) ..(87.0866,75.8574) .. controls (86.3658,75.8574) and (85.7814,76.4417) ..(85.7814,77.1625) .. controls (85.7814,77.8833) and (86.3658,78.4677) ..(87.0866,78.4677) .. controls (87.8074,78.4677) and (88.3917,77.8833) ..(88.3917,77.1625);
\path[fill=chfourfig5c0,fill opacity=0.5,draw=chfourfig5c0,line width=0.3985pt,draw opacity=0.5] (89.9658,83.5686) .. controls (89.9658,82.8478) and (89.3815,82.2635) ..(88.6607,82.2635) .. controls (87.9399,82.2635) and (87.3556,82.8478) ..(87.3556,83.5686) .. controls (87.3556,84.2894) and (87.9399,84.8737) ..(88.6607,84.8737) .. controls (89.3815,84.8737) and (89.9658,84.2894) ..(89.9658,83.5686);
\path[fill=chfourfig5c0,fill opacity=0.5,draw=chfourfig5c0,line width=0.3985pt,draw opacity=0.5] (86.2298,95.9324) .. controls (86.2298,95.2116) and (85.6455,94.6273) ..(84.9247,94.6273) .. controls (84.2039,94.6273) and (83.6195,95.2116) ..(83.6195,95.9324) .. controls (83.6195,96.6532) and (84.2039,97.2375) ..(84.9247,97.2375) .. controls (85.6455,97.2375) and (86.2298,96.6532) ..(86.2298,95.9324);
\path[fill=chfourfig5c0,fill opacity=0.5,draw=chfourfig5c0,line width=0.3985pt,draw opacity=0.5] (144.6415,42.8109) .. controls (144.6415,42.0901) and (144.0572,41.5058) ..(143.3364,41.5058) .. controls (142.6155,41.5058) and (142.0312,42.0901) ..(142.0312,42.8109) .. controls (142.0312,43.5317) and (142.6155,44.1160) ..(143.3364,44.1160) .. controls (144.0572,44.1160) and (144.6415,43.5317) ..(144.6415,42.8109);
\path[fill=chfourfig5c0,fill opacity=0.5,draw=chfourfig5c0,line width=0.3985pt,draw opacity=0.5] (110.3795,58.1835) .. controls (110.3795,57.4626) and (109.7952,56.8783) ..(109.0744,56.8783) .. controls (108.3536,56.8783) and (107.7693,57.4626) ..(107.7693,58.1835) .. controls (107.7693,58.9043) and (108.3536,59.4886) ..(109.0744,59.4886) .. controls (109.7952,59.4886) and (110.3795,58.9043) ..(110.3795,58.1835);
\path[fill=chfourfig5c0,fill opacity=0.5,draw=chfourfig5c0,line width=0.3985pt,draw opacity=0.5] (96.4416,74.5921) .. controls (96.4416,73.8713) and (95.8573,73.2870) ..(95.1365,73.2870) .. controls (94.4157,73.2870) and (93.8314,73.8713) ..(93.8314,74.5921) .. controls (93.8314,75.3129) and (94.4157,75.8972) ..(95.1365,75.8972) .. controls (95.8573,75.8972) and (96.4416,75.3129) ..(96.4416,74.5921);
\path[fill=chfourfig5c0,fill opacity=0.5,draw=chfourfig5c0,line width=0.3985pt,draw opacity=0.5] (84.6357,76.9135) .. controls (84.6357,76.1927) and (84.0514,75.6083) ..(83.3306,75.6083) .. controls (82.6098,75.6083) and (82.0255,76.1927) ..(82.0255,76.9135) .. controls (82.0255,77.6343) and (82.6098,78.2186) ..(83.3306,78.2186) .. controls (84.0514,78.2186) and (84.6357,77.6343) ..(84.6357,76.9135);
\path[fill=chfourfig5c0,fill opacity=0.5,draw=chfourfig5c0,line width=0.3985pt,draw opacity=0.5] (152.3427,45.4809) .. controls (152.3427,44.7601) and (151.7584,44.1758) ..(151.0376,44.1758) .. controls (150.3168,44.1758) and (149.7325,44.7601) ..(149.7325,45.4809) .. controls (149.7325,46.2017) and (150.3168,46.7860) ..(151.0376,46.7860) .. controls (151.7584,46.7860) and (152.3427,46.2017) ..(152.3427,45.4809);
\path[fill=chfourfig5c0,fill opacity=0.5,draw=chfourfig5c0,line width=0.3985pt,draw opacity=0.5] (54.6976,79.6632) .. controls (54.6976,78.9424) and (54.1133,78.3581) ..(53.3925,78.3581) .. controls (52.6717,78.3581) and (52.0874,78.9424) ..(52.0874,79.6632) .. controls (52.0874,80.3840) and (52.6717,80.9683) ..(53.3925,80.9683) .. controls (54.1133,80.9683) and (54.6976,80.3840) ..(54.6976,79.6632);
\path[fill=chfourfig5c0,fill opacity=0.5,draw=chfourfig5c0,line width=0.3985pt,draw opacity=0.5] (113.2588,51.5781) .. controls (113.2588,50.8573) and (112.6744,50.2730) ..(111.9536,50.2730) .. controls (111.2328,50.2730) and (110.6485,50.8573) ..(110.6485,51.5781) .. controls (110.6485,52.2990) and (111.2328,52.8833) ..(111.9536,52.8833) .. controls (112.6744,52.8833) and (113.2588,52.2990) ..(113.2588,51.5781);
\path[fill=chfourfig5c0,fill opacity=0.5,draw=chfourfig5c0,line width=0.3985pt,draw opacity=0.5] (96.9198,90.4927) .. controls (96.9198,89.7719) and (96.3355,89.1876) ..(95.6147,89.1876) .. controls (94.8939,89.1876) and (94.3096,89.7719) ..(94.3096,90.4927) .. controls (94.3096,91.2135) and (94.8939,91.7978) ..(95.6147,91.7978) .. controls (96.3355,91.7978) and (96.9198,91.2135) ..(96.9198,90.4927);
\path[fill=chfourfig5c0,fill opacity=0.5,draw=chfourfig5c0,line width=0.3985pt,draw opacity=0.5] (82.9022,92.3358) .. controls (82.9022,91.6150) and (82.3179,91.0307) ..(81.5971,91.0307) .. controls (80.8763,91.0307) and (80.2920,91.6150) ..(80.2920,92.3358) .. controls (80.2920,93.0566) and (80.8763,93.6409) ..(81.5971,93.6409) .. controls (82.3179,93.6409) and (82.9022,93.0566) ..(82.9022,92.3358);
\path[fill=chfourfig5c0,fill opacity=0.5,draw=chfourfig5c0,line width=0.3985pt,draw opacity=0.5] (149.1247,52.7936) .. controls (149.1247,52.0728) and (148.5404,51.4885) ..(147.8196,51.4885) .. controls (147.0988,51.4885) and (146.5145,52.0728) ..(146.5145,52.7936) .. controls (146.5145,53.5144) and (147.0988,54.0987) ..(147.8196,54.0987) .. controls (148.5404,54.0987) and (149.1247,53.5144) ..(149.1247,52.7936);
\path[fill=chfourfig5c0,fill opacity=0.5,draw=chfourfig5c0,line width=0.3985pt,draw opacity=0.5] (148.3476,38.8856) .. controls (148.3476,38.1648) and (147.7633,37.5804) ..(147.0425,37.5804) .. controls (146.3217,37.5804) and (145.7374,38.1648) ..(145.7374,38.8856) .. controls (145.7374,39.6064) and (146.3217,40.1907) ..(147.0425,40.1907) .. controls (147.7633,40.1907) and (148.3476,39.6064) ..(148.3476,38.8856);
\path[fill=chfourfig5c0,fill opacity=0.5,draw=chfourfig5c0,line width=0.3985pt,draw opacity=0.5] (82.7129,101.0831) .. controls (82.7129,100.3623) and (82.1286,99.7780) ..(81.4078,99.7780) .. controls (80.6870,99.7780) and (80.1027,100.3623) ..(80.1027,101.0831) .. controls (80.1027,101.8039) and (80.6870,102.3882) ..(81.4078,102.3882) .. controls (82.1286,102.3882) and (82.7129,101.8039) ..(82.7129,101.0831);
\path[fill=chfourfig5c0,fill opacity=0.5,draw=chfourfig5c0,line width=0.3985pt,draw opacity=0.5] (86.1600,61.4911) .. controls (86.1600,60.7703) and (85.5757,60.1860) ..(84.8549,60.1860) .. controls (84.1341,60.1860) and (83.5498,60.7703) ..(83.5498,61.4911) .. controls (83.5498,62.2119) and (84.1341,62.7962) ..(84.8549,62.7962) .. controls (85.5757,62.7962) and (86.1600,62.2119) ..(86.1600,61.4911);
\path[fill=chfourfig5c0,fill opacity=0.5,draw=chfourfig5c0,line width=0.3985pt,draw opacity=0.5] (140.4372,36.7137) .. controls (140.4372,35.9929) and (139.8529,35.4086) ..(139.1321,35.4086) .. controls (138.4113,35.4086) and (137.8270,35.9929) ..(137.8270,36.7137) .. controls (137.8270,37.4345) and (138.4113,38.0188) ..(139.1321,38.0188) .. controls (139.8529,38.0188) and (140.4372,37.4345) ..(140.4372,36.7137);
\path[fill=chfourfig5c0,fill opacity=0.5,draw=chfourfig5c0,line width=0.3985pt,draw opacity=0.5] (150.9380,39.2841) .. controls (150.9380,38.5633) and (150.3536,37.9789) ..(149.6328,37.9789) .. controls (148.9120,37.9789) and (148.3277,38.5633) ..(148.3277,39.2841) .. controls (148.3277,40.0049) and (148.9120,40.5892) ..(149.6328,40.5892) .. controls (150.3536,40.5892) and (150.9380,40.0049) ..(150.9380,39.2841);
\path[fill=chfourfig5c0,fill opacity=0.5,draw=chfourfig5c0,line width=0.3985pt,draw opacity=0.5] (128.3524,74.2733) .. controls (128.3524,73.5525) and (127.7680,72.9682) ..(127.0472,72.9682) .. controls (126.3264,72.9682) and (125.7421,73.5525) ..(125.7421,74.2733) .. controls (125.7421,74.9941) and (126.3264,75.5784) ..(127.0472,75.5784) .. controls (127.7680,75.5784) and (128.3524,74.9941) ..(128.3524,74.2733);
\path[fill=chfourfig5c0,fill opacity=0.5,draw=chfourfig5c0,line width=0.3985pt,draw opacity=0.5] (79.1363,68.2857) .. controls (79.1363,67.5649) and (78.5519,66.9806) ..(77.8311,66.9806) .. controls (77.1103,66.9806) and (76.5260,67.5649) ..(76.5260,68.2857) .. controls (76.5260,69.0065) and (77.1103,69.5908) ..(77.8311,69.5908) .. controls (78.5519,69.5908) and (79.1363,69.0065) ..(79.1363,68.2857);
\path[fill=chfourfig5c0,fill opacity=0.5,draw=chfourfig5c0,line width=0.3985pt,draw opacity=0.5] (126.9576,39.1845) .. controls (126.9576,38.4637) and (126.3732,37.8793) ..(125.6524,37.8793) .. controls (124.9316,37.8793) and (124.3473,38.4637) ..(124.3473,39.1845) .. controls (124.3473,39.9052) and (124.9316,40.4896) ..(125.6524,40.4896) .. controls (126.3732,40.4896) and (126.9576,39.9052) ..(126.9576,39.1845);
\path[fill=chfourfig5c0,fill opacity=0.5,draw=chfourfig5c0,line width=0.3985pt,draw opacity=0.5] (143.5854,65.3467) .. controls (143.5854,64.6259) and (143.0011,64.0416) ..(142.2803,64.0416) .. controls (141.5595,64.0416) and (140.9752,64.6259) ..(140.9752,65.3467) .. controls (140.9752,66.0675) and (141.5595,66.6518) ..(142.2803,66.6518) .. controls (143.0011,66.6518) and (143.5854,66.0675) ..(143.5854,65.3467);
\path[fill=chfourfig5c0,fill opacity=0.5,draw=chfourfig5c0,line width=0.3985pt,draw opacity=0.5] (70.6679,94.3184) .. controls (70.6679,93.5976) and (70.0836,93.0133) ..(69.3628,93.0133) .. controls (68.6420,93.0133) and (68.0577,93.5976) ..(68.0577,94.3184) .. controls (68.0577,95.0392) and (68.6420,95.6235) ..(69.3628,95.6235) .. controls (70.0836,95.6235) and (70.6679,95.0392) ..(70.6679,94.3184);
\path[fill=chfourfig5c0,fill opacity=0.5,draw=chfourfig5c0,line width=0.3985pt,draw opacity=0.5] (152.0139,48.7786) .. controls (152.0139,48.0578) and (151.4296,47.4735) ..(150.7088,47.4735) .. controls (149.9880,47.4735) and (149.4037,48.0578) ..(149.4037,48.7786) .. controls (149.4037,49.4994) and (149.9880,50.0837) ..(150.7088,50.0837) .. controls (151.4296,50.0837) and (152.0139,49.4994) ..(152.0139,48.7786);
\path[fill=chfourfig5c0,fill opacity=0.5,draw=chfourfig5c0,line width=0.3985pt,draw opacity=0.5] (99.8289,81.8251) .. controls (99.8289,81.1043) and (99.2446,80.5200) ..(98.5238,80.5200) .. controls (97.8030,80.5200) and (97.2187,81.1043) ..(97.2187,81.8251) .. controls (97.2187,82.5459) and (97.8030,83.1302) ..(98.5238,83.1302) .. controls (99.2446,83.1302) and (99.8289,82.5459) ..(99.8289,81.8251);
\path[fill=chfourfig5c0,fill opacity=0.5,draw=chfourfig5c0,line width=0.3985pt,draw opacity=0.5] (151.7051,26.3126) .. controls (151.7051,25.5918) and (151.1208,25.0074) ..(150.4000,25.0074) .. controls (149.6792,25.0074) and (149.0948,25.5918) ..(149.0948,26.3126) .. controls (149.0948,27.0334) and (149.6792,27.6177) ..(150.4000,27.6177) .. controls (151.1208,27.6177) and (151.7051,27.0334) ..(151.7051,26.3126);
\path[fill=chfourfig5c0,fill opacity=0.5,draw=chfourfig5c0,line width=0.3985pt,draw opacity=0.5] (132.6862,48.0414) .. controls (132.6862,47.3206) and (132.1018,46.7362) ..(131.3810,46.7362) .. controls (130.6602,46.7362) and (130.0759,47.3206) ..(130.0759,48.0414) .. controls (130.0759,48.7622) and (130.6602,49.3465) ..(131.3810,49.3465) .. controls (132.1018,49.3465) and (132.6862,48.7622) ..(132.6862,48.0414);
\path[fill=chfourfig5c0,fill opacity=0.5,draw=chfourfig5c0,line width=0.3985pt,draw opacity=0.5] (104.5513,62.0988) .. controls (104.5513,61.3780) and (103.9670,60.7937) ..(103.2462,60.7937) .. controls (102.5254,60.7937) and (101.9411,61.3780) ..(101.9411,62.0988) .. controls (101.9411,62.8196) and (102.5254,63.4039) ..(103.2462,63.4039) .. controls (103.9670,63.4039) and (104.5513,62.8196) ..(104.5513,62.0988);
\path[fill=chfourfig5c0,fill opacity=0.5,draw=chfourfig5c0,line width=0.3985pt,draw opacity=0.5] (87.2260,76.7341) .. controls (87.2260,76.0133) and (86.6417,75.4290) ..(85.9209,75.4290) .. controls (85.2001,75.4290) and (84.6158,76.0133) ..(84.6158,76.7341) .. controls (84.6158,77.4549) and (85.2001,78.0393) ..(85.9209,78.0393) .. controls (86.6417,78.0393) and (87.2260,77.4549) ..(87.2260,76.7341);
\path[fill=chfourfig5c0,fill opacity=0.5,draw=chfourfig5c0,line width=0.3985pt,draw opacity=0.5] (110.4991,79.3842) .. controls (110.4991,78.6634) and (109.9148,78.0791) ..(109.1939,78.0791) .. controls (108.4731,78.0791) and (107.8888,78.6634) ..(107.8888,79.3842) .. controls (107.8888,80.1050) and (108.4731,80.6893) ..(109.1939,80.6893) .. controls (109.9148,80.6893) and (110.4991,80.1050) ..(110.4991,79.3842);
\path[fill=chfourfig5c0,fill opacity=0.5,draw=chfourfig5c0,line width=0.3985pt,draw opacity=0.5] (83.2609,66.8411) .. controls (83.2609,66.1203) and (82.6765,65.5360) ..(81.9557,65.5360) .. controls (81.2350,65.5360) and (80.6506,66.1203) ..(80.6506,66.8411) .. controls (80.6506,67.5619) and (81.2350,68.1462) ..(81.9557,68.1462) .. controls (82.6765,68.1462) and (83.2609,67.5619) ..(83.2609,66.8411);
\path[fill=chfourfig5c0,fill opacity=0.5,draw=chfourfig5c0,line width=0.3985pt,draw opacity=0.5] (97.9858,78.2883) .. controls (97.9858,77.5675) and (97.4015,76.9832) ..(96.6807,76.9832) .. controls (95.9599,76.9832) and (95.3756,77.5675) ..(95.3756,78.2883) .. controls (95.3756,79.0091) and (95.9599,79.5934) ..(96.6807,79.5934) .. controls (97.4015,79.5934) and (97.9858,79.0091) ..(97.9858,78.2883);
\path[fill=chfourfig5c0,fill opacity=0.5,draw=chfourfig5c0,line width=0.3985pt,draw opacity=0.5] (95.0667,72.0715) .. controls (95.0667,71.3508) and (94.4824,70.7664) ..(93.7616,70.7664) .. controls (93.0408,70.7664) and (92.4565,71.3508) ..(92.4565,72.0715) .. controls (92.4565,72.7924) and (93.0408,73.3767) ..(93.7616,73.3767) .. controls (94.4824,73.3767) and (95.0667,72.7924) ..(95.0667,72.0715);
\path[fill=chfourfig5c0,fill opacity=0.5,draw=chfourfig5c0,line width=0.3985pt,draw opacity=0.5] (76.3766,69.7104) .. controls (76.3766,68.9896) and (75.7923,68.4053) ..(75.0715,68.4053) .. controls (74.3507,68.4053) and (73.7664,68.9896) ..(73.7664,69.7104) .. controls (73.7664,70.4312) and (74.3507,71.0155) ..(75.0715,71.0155) .. controls (75.7923,71.0155) and (76.3766,70.4312) ..(76.3766,69.7104);
\path[fill=chfourfig5c0,fill opacity=0.5,draw=chfourfig5c0,line width=0.3985pt,draw opacity=0.5] (112.6510,46.2680) .. controls (112.6510,45.5472) and (112.0667,44.9629) ..(111.3459,44.9629) .. controls (110.6251,44.9629) and (110.0408,45.5472) ..(110.0408,46.2680) .. controls (110.0408,46.9888) and (110.6251,47.5731) ..(111.3459,47.5731) .. controls (112.0667,47.5731) and (112.6510,46.9888) ..(112.6510,46.2680);
\path[fill=chfourfig5c0,fill opacity=0.5,draw=chfourfig5c0,line width=0.3985pt,draw opacity=0.5] (60.7450,93.7804) .. controls (60.7450,93.0596) and (60.1607,92.4753) ..(59.4399,92.4753) .. controls (58.7191,92.4753) and (58.1348,93.0596) ..(58.1348,93.7804) .. controls (58.1348,94.5012) and (58.7191,95.0856) ..(59.4399,95.0856) .. controls (60.1607,95.0856) and (60.7450,94.5012) ..(60.7450,93.7804);
\path[fill=chfourfig5c0,fill opacity=0.5,draw=chfourfig5c0,line width=0.3985pt,draw opacity=0.5] (127.7247,44.9131) .. controls (127.7247,44.1922) and (127.1404,43.6079) ..(126.4196,43.6079) .. controls (125.6988,43.6079) and (125.1144,44.1922) ..(125.1144,44.9131) .. controls (125.1144,45.6339) and (125.6988,46.2182) ..(126.4196,46.2182) .. controls (127.1404,46.2182) and (127.7247,45.6339) ..(127.7247,44.9131);
\path[fill=chfourfig5c0,fill opacity=0.5,draw=chfourfig5c0,line width=0.3985pt,draw opacity=0.5] (83.6594,99.6485) .. controls (83.6594,98.9277) and (83.0750,98.3434) ..(82.3542,98.3434) .. controls (81.6335,98.3434) and (81.0491,98.9277) ..(81.0491,99.6485) .. controls (81.0491,100.3693) and (81.6335,100.9536) ..(82.3542,100.9536) .. controls (83.0750,100.9536) and (83.6594,100.3693) ..(83.6594,99.6485);
\path[fill=chfourfig5c0,fill opacity=0.5,draw=chfourfig5c0,line width=0.3985pt,draw opacity=0.5] (111.2463,47.2643) .. controls (111.2463,46.5435) and (110.6620,45.9591) ..(109.9412,45.9591) .. controls (109.2204,45.9591) and (108.6360,46.5435) ..(108.6360,47.2643) .. controls (108.6360,47.9851) and (109.2204,48.5694) ..(109.9412,48.5694) .. controls (110.6620,48.5694) and (111.2463,47.9851) ..(111.2463,47.2643);
\path[fill=chfourfig5c0,fill opacity=0.5,draw=chfourfig5c0,line width=0.3985pt,draw opacity=0.5] (129.8069,54.9754) .. controls (129.8069,54.2546) and (129.2226,53.6703) ..(128.5018,53.6703) .. controls (127.7810,53.6703) and (127.1967,54.2546) ..(127.1967,54.9754) .. controls (127.1967,55.6962) and (127.7810,56.2806) ..(128.5018,56.2806) .. controls (129.2226,56.2806) and (129.8069,55.6962) ..(129.8069,54.9754);
\path[fill=chfourfig5c0,fill opacity=0.5,draw=chfourfig5c0,line width=0.3985pt,draw opacity=0.5] (126.8480,45.5507) .. controls (126.8480,44.8299) and (126.2637,44.2455) ..(125.5429,44.2455) .. controls (124.8221,44.2455) and (124.2377,44.8299) ..(124.2377,45.5507) .. controls (124.2377,46.2715) and (124.8221,46.8558) ..(125.5429,46.8558) .. controls (126.2637,46.8558) and (126.8480,46.2715) ..(126.8480,45.5507);
\path[fill=chfourfig5c0,fill opacity=0.5,draw=chfourfig5c0,line width=0.3985pt,draw opacity=0.5] (70.8672,78.2784) .. controls (70.8672,77.5576) and (70.2829,76.9732) ..(69.5621,76.9732) .. controls (68.8413,76.9732) and (68.2569,77.5576) ..(68.2569,78.2784) .. controls (68.2569,78.9992) and (68.8413,79.5835) ..(69.5621,79.5835) .. controls (70.2829,79.5835) and (70.8672,78.9992) ..(70.8672,78.2784);
\path[fill=chfourfig5c0,fill opacity=0.5,draw=chfourfig5c0,line width=0.3985pt,draw opacity=0.5] (53.2032,88.2312) .. controls (53.2032,87.5104) and (52.6189,86.9260) ..(51.8981,86.9260) .. controls (51.1772,86.9260) and (50.5929,87.5104) ..(50.5929,88.2312) .. controls (50.5929,88.9520) and (51.1772,89.5363) ..(51.8981,89.5363) .. controls (52.6189,89.5363) and (53.2032,88.9520) ..(53.2032,88.2312);
\path[fill=chfourfig5c0,fill opacity=0.5,draw=chfourfig5c0,line width=0.3985pt,draw opacity=0.5] (117.6324,54.3279) .. controls (117.6324,53.6071) and (117.0481,53.0227) ..(116.3273,53.0227) .. controls (115.6065,53.0227) and (115.0222,53.6071) ..(115.0222,54.3279) .. controls (115.0222,55.0487) and (115.6065,55.6330) ..(116.3273,55.6330) .. controls (117.0481,55.6330) and (117.6324,55.0487) ..(117.6324,54.3279);
\path[fill=chfourfig5c0,fill opacity=0.5,draw=chfourfig5c0,line width=0.3985pt,draw opacity=0.5] (68.2769,108.6648) .. controls (68.2769,107.9440) and (67.6925,107.3597) ..(66.9717,107.3597) .. controls (66.2509,107.3597) and (65.6666,107.9440) ..(65.6666,108.6648) .. controls (65.6666,109.3856) and (66.2509,109.9699) ..(66.9717,109.9699) .. controls (67.6925,109.9699) and (68.2769,109.3856) ..(68.2769,108.6648);
\path[fill=chfourfig5c0,fill opacity=0.5,draw=chfourfig5c0,line width=0.3985pt,draw opacity=0.5] (67.2806,84.1564) .. controls (67.2806,83.4356) and (66.6963,82.8513) ..(65.9755,82.8513) .. controls (65.2547,82.8513) and (64.6703,83.4356) ..(64.6703,84.1564) .. controls (64.6703,84.8772) and (65.2547,85.4615) ..(65.9755,85.4615) .. controls (66.6963,85.4615) and (67.2806,84.8772) ..(67.2806,84.1564);
\path[fill=chfourfig5c0,fill opacity=0.5,draw=chfourfig5c0,line width=0.3985pt,draw opacity=0.5] (117.4830,56.2009) .. controls (117.4830,55.4801) and (116.8986,54.8958) ..(116.1779,54.8958) .. controls (115.4570,54.8958) and (114.8727,55.4801) ..(114.8727,56.2009) .. controls (114.8727,56.9217) and (115.4570,57.5060) ..(116.1779,57.5060) .. controls (116.8986,57.5060) and (117.4830,56.9217) ..(117.4830,56.2009);
\path[fill=chfourfig5c0,fill opacity=0.5,draw=chfourfig5c0,line width=0.3985pt,draw opacity=0.5] (114.5240,83.8575) .. controls (114.5240,83.1367) and (113.9397,82.5524) ..(113.2189,82.5524) .. controls (112.4981,82.5524) and (111.9138,83.1367) ..(111.9138,83.8575) .. controls (111.9138,84.5783) and (112.4981,85.1626) ..(113.2189,85.1626) .. controls (113.9397,85.1626) and (114.5240,84.5783) ..(114.5240,83.8575);
\path[fill=chfourfig5c0,fill opacity=0.5,draw=chfourfig5c0,line width=0.3985pt,draw opacity=0.5] (119.5353,69.0130) .. controls (119.5353,68.2922) and (118.9510,67.7079) ..(118.2302,67.7079) .. controls (117.5094,67.7079) and (116.9250,68.2922) ..(116.9250,69.0130) .. controls (116.9250,69.7338) and (117.5094,70.3181) ..(118.2302,70.3181) .. controls (118.9510,70.3181) and (119.5353,69.7338) ..(119.5353,69.0130);
\path[fill=chfourfig5c0,fill opacity=0.5,draw=chfourfig5c0,line width=0.3985pt,draw opacity=0.5] (152.6217,33.4559) .. controls (152.6217,32.7351) and (152.0374,32.1508) ..(151.3165,32.1508) .. controls (150.5957,32.1508) and (150.0114,32.7351) ..(150.0114,33.4559) .. controls (150.0114,34.1767) and (150.5957,34.7610) ..(151.3165,34.7610) .. controls (152.0374,34.7610) and (152.6217,34.1767) ..(152.6217,33.4559);
\path[fill=chfourfig5c0,fill opacity=0.5,draw=chfourfig5c0,line width=0.3985pt,draw opacity=0.5] (140.6664,62.6866) .. controls (140.6664,61.9658) and (140.0820,61.3815) ..(139.3612,61.3815) .. controls (138.6404,61.3815) and (138.0561,61.9658) ..(138.0561,62.6866) .. controls (138.0561,63.4074) and (138.6404,63.9917) ..(139.3612,63.9917) .. controls (140.0820,63.9917) and (140.6664,63.4074) ..(140.6664,62.6866);
\path[fill=chfourfig5c0,fill opacity=0.5,draw=chfourfig5c0,line width=0.3985pt,draw opacity=0.5] (63.3154,112.6599) .. controls (63.3154,111.9391) and (62.7311,111.3548) ..(62.0103,111.3548) .. controls (61.2895,111.3548) and (60.7052,111.9391) ..(60.7052,112.6599) .. controls (60.7052,113.3807) and (61.2895,113.9650) ..(62.0103,113.9650) .. controls (62.7311,113.9650) and (63.3154,113.3807) ..(63.3154,112.6599);
\path[fill=chfourfig5c0,fill opacity=0.5,draw=chfourfig5c0,line width=0.3985pt,draw opacity=0.5] (68.1673,99.3098) .. controls (68.1673,98.5890) and (67.5829,98.0046) ..(66.8622,98.0046) .. controls (66.1414,98.0046) and (65.5570,98.5890) ..(65.5570,99.3098) .. controls (65.5570,100.0306) and (66.1414,100.6149) ..(66.8622,100.6149) .. controls (67.5829,100.6149) and (68.1673,100.0306) ..(68.1673,99.3098);
\path[fill=chfourfig5c0,fill opacity=0.5,draw=chfourfig5c0,line width=0.3985pt,draw opacity=0.5] (76.8050,72.2509) .. controls (76.8050,71.5301) and (76.2207,70.9458) ..(75.4999,70.9458) .. controls (74.7791,70.9458) and (74.1947,71.5301) ..(74.1947,72.2509) .. controls (74.1947,72.9717) and (74.7791,73.5560) ..(75.4999,73.5560) .. controls (76.2207,73.5560) and (76.8050,72.9717) ..(76.8050,72.2509);
\path[fill=chfourfig5c0,fill opacity=0.5,draw=chfourfig5c0,line width=0.3985pt,draw opacity=0.5] (52.2667,114.7620) .. controls (52.2667,114.0412) and (51.6824,113.4569) ..(50.9616,113.4569) .. controls (50.2408,113.4569) and (49.6565,114.0412) ..(49.6565,114.7620) .. controls (49.6565,115.4828) and (50.2408,116.0671) ..(50.9616,116.0671) .. controls (51.6824,116.0671) and (52.2667,115.4828) ..(52.2667,114.7620);
\path[fill=chfourfig5c0,fill opacity=0.5,draw=chfourfig5c0,line width=0.3985pt,draw opacity=0.5] (130.1457,40.2405) .. controls (130.1457,39.5197) and (129.5613,38.9354) ..(128.8405,38.9354) .. controls (128.1197,38.9354) and (127.5354,39.5197) ..(127.5354,40.2405) .. controls (127.5354,40.9613) and (128.1197,41.5456) ..(128.8405,41.5456) .. controls (129.5613,41.5456) and (130.1457,40.9613) ..(130.1457,40.2405);
\path[fill=chfourfig5c0,fill opacity=0.5,draw=chfourfig5c0,line width=0.3985pt,draw opacity=0.5] (139.2616,32.0810) .. controls (139.2616,31.3602) and (138.6773,30.7759) ..(137.9565,30.7759) .. controls (137.2357,30.7759) and (136.6513,31.3602) ..(136.6513,32.0810) .. controls (136.6513,32.8018) and (137.2357,33.3861) ..(137.9565,33.3861) .. controls (138.6773,33.3861) and (139.2616,32.8018) ..(139.2616,32.0810);
\path[fill=chfourfig5c0,fill opacity=0.5,draw=chfourfig5c0,line width=0.3985pt,draw opacity=0.5] (133.3337,59.4487) .. controls (133.3337,58.7279) and (132.7494,58.1436) ..(132.0286,58.1436) .. controls (131.3078,58.1436) and (130.7235,58.7279) ..(130.7235,59.4487) .. controls (130.7235,60.1695) and (131.3078,60.7538) ..(132.0286,60.7538) .. controls (132.7494,60.7538) and (133.3337,60.1695) ..(133.3337,59.4487);
\path[fill=chfourfig5c0,fill opacity=0.5,draw=chfourfig5c0,line width=0.3985pt,draw opacity=0.5] (68.2071,81.1875) .. controls (68.2071,80.4667) and (67.6228,79.8824) ..(66.9020,79.8824) .. controls (66.1812,79.8824) and (65.5969,80.4667) ..(65.5969,81.1875) .. controls (65.5969,81.9083) and (66.1812,82.4926) ..(66.9020,82.4926) .. controls (67.6228,82.4926) and (68.2071,81.9083) ..(68.2071,81.1875);
\path[fill=chfourfig5c0,fill opacity=0.5,draw=chfourfig5c0,line width=0.3985pt,draw opacity=0.5] (76.2371,88.0419) .. controls (76.2371,87.3211) and (75.6528,86.7367) ..(74.9320,86.7367) .. controls (74.2112,86.7367) and (73.6269,87.3211) ..(73.6269,88.0419) .. controls (73.6269,88.7627) and (74.2112,89.3470) ..(74.9320,89.3470) .. controls (75.6528,89.3470) and (76.2371,88.7627) ..(76.2371,88.0419);
\path[fill=chfourfig5c0,fill opacity=0.5,draw=chfourfig5c0,line width=0.3985pt,draw opacity=0.5] (108.4965,55.0053) .. controls (108.4965,54.2845) and (107.9122,53.7002) ..(107.1914,53.7002) .. controls (106.4706,53.7002) and (105.8863,54.2845) ..(105.8863,55.0053) .. controls (105.8863,55.7261) and (106.4706,56.3104) ..(107.1914,56.3104) .. controls (107.9122,56.3104) and (108.4965,55.7261) ..(108.4965,55.0053);
\path[fill=chfourfig5c0,fill opacity=0.5,draw=chfourfig5c0,line width=0.3985pt,draw opacity=0.5] (109.6921,69.8897) .. controls (109.6921,69.1689) and (109.1078,68.5846) ..(108.3870,68.5846) .. controls (107.6662,68.5846) and (107.0818,69.1689) ..(107.0818,69.8897) .. controls (107.0818,70.6105) and (107.6662,71.1948) ..(108.3870,71.1948) .. controls (109.1078,71.1948) and (109.6921,70.6105) ..(109.6921,69.8897);
\path[fill=chfourfig5c0,fill opacity=0.5,draw=chfourfig5c0,line width=0.3985pt,draw opacity=0.5] (123.7994,50.0638) .. controls (123.7994,49.3430) and (123.2151,48.7587) ..(122.4942,48.7587) .. controls (121.7735,48.7587) and (121.1891,49.3430) ..(121.1891,50.0638) .. controls (121.1891,50.7846) and (121.7735,51.3689) ..(122.4942,51.3689) .. controls (123.2151,51.3689) and (123.7994,50.7846) ..(123.7994,50.0638);
\path[fill=chfourfig5c0,fill opacity=0.5,draw=chfourfig5c0,line width=0.3985pt,draw opacity=0.5] (141.1844,40.8781) .. controls (141.1844,40.1573) and (140.6001,39.5730) ..(139.8793,39.5730) .. controls (139.1585,39.5730) and (138.5742,40.1573) ..(138.5742,40.8781) .. controls (138.5742,41.5989) and (139.1585,42.1833) ..(139.8793,42.1833) .. controls (140.6001,42.1833) and (141.1844,41.5989) ..(141.1844,40.8781);
\path[fill=chfourfig5c0,fill opacity=0.5,draw=chfourfig5c0,line width=0.3985pt,draw opacity=0.5] (115.9886,65.1873) .. controls (115.9886,64.4665) and (115.4042,63.8822) ..(114.6834,63.8822) .. controls (113.9626,63.8822) and (113.3783,64.4665) ..(113.3783,65.1873) .. controls (113.3783,65.9081) and (113.9626,66.4924) ..(114.6834,66.4924) .. controls (115.4042,66.4924) and (115.9886,65.9081) ..(115.9886,65.1873);
\path[fill=chfourfig5c0,fill opacity=0.5,draw=chfourfig5c0,line width=0.3985pt,draw opacity=0.5] (82.7029,78.1987) .. controls (82.7029,77.4778) and (82.1186,76.8935) ..(81.3978,76.8935) .. controls (80.6770,76.8935) and (80.0927,77.4778) ..(80.0927,78.1987) .. controls (80.0927,78.9194) and (80.6770,79.5038) ..(81.3978,79.5038) .. controls (82.1186,79.5038) and (82.7029,78.9194) ..(82.7029,78.1987);
\path[fill=chfourfig5c0,fill opacity=0.5,draw=chfourfig5c0,line width=0.3985pt,draw opacity=0.5] (87.4153,60.2657) .. controls (87.4153,59.5449) and (86.8310,58.9605) ..(86.1102,58.9605) .. controls (85.3894,58.9605) and (84.8051,59.5449) ..(84.8051,60.2657) .. controls (84.8051,60.9865) and (85.3894,61.5708) ..(86.1102,61.5708) .. controls (86.8310,61.5708) and (87.4153,60.9865) ..(87.4153,60.2657);
\path[fill=chfourfig5c0,fill opacity=0.5,draw=chfourfig5c0,line width=0.3985pt,draw opacity=0.5] (54.6976,82.7118) .. controls (54.6976,81.9910) and (54.1133,81.4067) ..(53.3925,81.4067) .. controls (52.6717,81.4067) and (52.0874,81.9910) ..(52.0874,82.7118) .. controls (52.0874,83.4326) and (52.6717,84.0169) ..(53.3925,84.0169) .. controls (54.1133,84.0169) and (54.6976,83.4326) ..(54.6976,82.7118);
\path[fill=chfourfig5c0,fill opacity=0.5,draw=chfourfig5c0,line width=0.3985pt,draw opacity=0.5] (57.7960,111.5939) .. controls (57.7960,110.8730) and (57.2117,110.2887) ..(56.4909,110.2887) .. controls (55.7701,110.2887) and (55.1858,110.8730) ..(55.1858,111.5939) .. controls (55.1858,112.3147) and (55.7701,112.8990) ..(56.4909,112.8990) .. controls (57.2117,112.8990) and (57.7960,112.3147) ..(57.7960,111.5939);
\path[fill=chfourfig5c0,fill opacity=0.5,draw=chfourfig5c0,line width=0.3985pt,draw opacity=0.5] (86.0106,80.0717) .. controls (86.0106,79.3508) and (85.4263,78.7665) ..(84.7055,78.7665) .. controls (83.9847,78.7665) and (83.4004,79.3508) ..(83.4004,80.0717) .. controls (83.4004,80.7924) and (83.9847,81.3768) ..(84.7055,81.3768) .. controls (85.4263,81.3768) and (86.0106,80.7924) ..(86.0106,80.0717);
\path[fill=chfourfig5c0,fill opacity=0.5,draw=chfourfig5c0,line width=0.3985pt,draw opacity=0.5] (148.5867,40.1508) .. controls (148.5867,39.4300) and (148.0024,38.8457) ..(147.2816,38.8457) .. controls (146.5608,38.8457) and (145.9765,39.4300) ..(145.9765,40.1508) .. controls (145.9765,40.8716) and (146.5608,41.4560) ..(147.2816,41.4560) .. controls (148.0024,41.4560) and (148.5867,40.8716) ..(148.5867,40.1508);
\path[fill=chfourfig5c0,fill opacity=0.5,draw=chfourfig5c0,line width=0.3985pt,draw opacity=0.5] (114.9823,54.7861) .. controls (114.9823,54.0653) and (114.3980,53.4810) ..(113.6772,53.4810) .. controls (112.9564,53.4810) and (112.3721,54.0653) ..(112.3721,54.7861) .. controls (112.3721,55.5069) and (112.9564,56.0913) ..(113.6772,56.0913) .. controls (114.3980,56.0913) and (114.9823,55.5069) ..(114.9823,54.7861);
\path[fill=chfourfig5c0,fill opacity=0.5,draw=chfourfig5c0,line width=0.3985pt,draw opacity=0.5] (127.7347,63.6032) .. controls (127.7347,62.8824) and (127.1504,62.2981) ..(126.4295,62.2981) .. controls (125.7087,62.2981) and (125.1244,62.8824) ..(125.1244,63.6032) .. controls (125.1244,64.3240) and (125.7087,64.9083) ..(126.4295,64.9083) .. controls (127.1504,64.9083) and (127.7347,64.3240) ..(127.7347,63.6032);
\path[fill=chfourfig5c0,fill opacity=0.5,draw=chfourfig5c0,line width=0.3985pt,draw opacity=0.5] (62.1497,84.6944) .. controls (62.1497,83.9736) and (61.5654,83.3893) ..(60.8446,83.3893) .. controls (60.1238,83.3893) and (59.5395,83.9736) ..(59.5395,84.6944) .. controls (59.5395,85.4152) and (60.1238,85.9995) ..(60.8446,85.9995) .. controls (61.5654,85.9995) and (62.1497,85.4152) ..(62.1497,84.6944);
\path[fill=chfourfig5c0,fill opacity=0.5,draw=chfourfig5c0,line width=0.3985pt,draw opacity=0.5] (113.6573,59.5085) .. controls (113.6573,58.7877) and (113.0730,58.2034) ..(112.3521,58.2034) .. controls (111.6314,58.2034) and (111.0470,58.7877) ..(111.0470,59.5085) .. controls (111.0470,60.2293) and (111.6314,60.8136) ..(112.3521,60.8136) .. controls (113.0730,60.8136) and (113.6573,60.2293) ..(113.6573,59.5085);
\path[fill=chfourfig5c0,fill opacity=0.5,draw=chfourfig5c0,line width=0.3985pt,draw opacity=0.5] (157.6429,43.7076) .. controls (157.6429,42.9868) and (157.0586,42.4024) ..(156.3378,42.4024) .. controls (155.6170,42.4024) and (155.0327,42.9868) ..(155.0327,43.7076) .. controls (155.0327,44.4283) and (155.6170,45.0127) ..(156.3378,45.0127) .. controls (157.0586,45.0127) and (157.6429,44.4283) ..(157.6429,43.7076);
\path[fill=chfourfig5c0,fill opacity=0.5,draw=chfourfig5c0,line width=0.3985pt,draw opacity=0.5] (78.7079,70.0192) .. controls (78.7079,69.2984) and (78.1236,68.7141) ..(77.4028,68.7141) .. controls (76.6819,68.7141) and (76.0976,69.2984) ..(76.0976,70.0192) .. controls (76.0976,70.7400) and (76.6819,71.3243) ..(77.4028,71.3243) .. controls (78.1236,71.3243) and (78.7079,70.7400) ..(78.7079,70.0192);
\path[fill=chfourfig5c0,fill opacity=0.5,draw=chfourfig5c0,line width=0.3985pt,draw opacity=0.5] (120.1131,43.0400) .. controls (120.1131,42.3192) and (119.5288,41.7349) ..(118.8080,41.7349) .. controls (118.0872,41.7349) and (117.5029,42.3192) ..(117.5029,43.0400) .. controls (117.5029,43.7608) and (118.0872,44.3452) ..(118.8080,44.3452) .. controls (119.5288,44.3452) and (120.1131,43.7608) ..(120.1131,43.0400);
\path[fill=chfourfig5c0,fill opacity=0.5,draw=chfourfig5c0,line width=0.3985pt,draw opacity=0.5] (105.0893,82.4029) .. controls (105.0893,81.6821) and (104.5050,81.0978) ..(103.7841,81.0978) .. controls (103.0634,81.0978) and (102.4790,81.6821) ..(102.4790,82.4029) .. controls (102.4790,83.1237) and (103.0634,83.7081) ..(103.7841,83.7081) .. controls (104.5050,83.7081) and (105.0893,83.1237) ..(105.0893,82.4029);
\path[fill=chfourfig5c0,fill opacity=0.5,draw=chfourfig5c0,line width=0.3985pt,draw opacity=0.5] (132.0784,60.3454) .. controls (132.0784,59.6246) and (131.4941,59.0403) ..(130.7733,59.0403) .. controls (130.0525,59.0403) and (129.4682,59.6246) ..(129.4682,60.3454) .. controls (129.4682,61.0662) and (130.0525,61.6505) ..(130.7733,61.6505) .. controls (131.4941,61.6505) and (132.0784,61.0662) ..(132.0784,60.3454);
\path[fill=chfourfig5c0,fill opacity=0.5,draw=chfourfig5c0,line width=0.3985pt,draw opacity=0.5] (125.5428,61.6206) .. controls (125.5428,60.8998) and (124.9585,60.3155) ..(124.2377,60.3155) .. controls (123.5169,60.3155) and (122.9326,60.8998) ..(122.9326,61.6206) .. controls (122.9326,62.3414) and (123.5169,62.9257) ..(124.2377,62.9257) .. controls (124.9585,62.9257) and (125.5428,62.3414) ..(125.5428,61.6206);
\path[fill=chfourfig5c0,fill opacity=0.5,draw=chfourfig5c0,line width=0.3985pt,draw opacity=0.5] (111.1467,83.3096) .. controls (111.1467,82.5887) and (110.5623,82.0044) ..(109.8415,82.0044) .. controls (109.1207,82.0044) and (108.5364,82.5887) ..(108.5364,83.3096) .. controls (108.5364,84.0303) and (109.1207,84.6147) ..(109.8415,84.6147) .. controls (110.5623,84.6147) and (111.1467,84.0303) ..(111.1467,83.3096);
\path[fill=chfourfig5c0,fill opacity=0.5,draw=chfourfig5c0,line width=0.3985pt,draw opacity=0.5] (128.9103,63.3043) .. controls (128.9103,62.5835) and (128.3260,61.9992) ..(127.6051,61.9992) .. controls (126.8843,61.9992) and (126.3000,62.5835) ..(126.3000,63.3043) .. controls (126.3000,64.0251) and (126.8843,64.6094) ..(127.6051,64.6094) .. controls (128.3260,64.6094) and (128.9103,64.0251) ..(128.9103,63.3043);
\path[fill=chfourfig5c0,fill opacity=0.5,draw=chfourfig5c0,line width=0.3985pt,draw opacity=0.5] (93.7516,74.3829) .. controls (93.7516,73.6621) and (93.1674,73.0778) ..(92.4465,73.0778) .. controls (91.7257,73.0778) and (91.1414,73.6621) ..(91.1414,74.3829) .. controls (91.1414,75.1037) and (91.7257,75.6880) ..(92.4465,75.6880) .. controls (93.1674,75.6880) and (93.7516,75.1037) ..(93.7516,74.3829);
\path[fill=chfourfig5c0,fill opacity=0.5,draw=chfourfig5c0,line width=0.3985pt,draw opacity=0.5] (101.0643,65.3168) .. controls (101.0643,64.5960) and (100.4800,64.0117) ..(99.7592,64.0117) .. controls (99.0384,64.0117) and (98.4541,64.5960) ..(98.4541,65.3168) .. controls (98.4541,66.0376) and (99.0384,66.6219) ..(99.7592,66.6219) .. controls (100.4800,66.6219) and (101.0643,66.0376) ..(101.0643,65.3168);
\path[fill=chfourfig5c0,fill opacity=0.5,draw=chfourfig5c0,line width=0.3985pt,draw opacity=0.5] (133.3537,49.9642) .. controls (133.3537,49.2434) and (132.7693,48.6591) ..(132.0486,48.6591) .. controls (131.3278,48.6591) and (130.7434,49.2434) ..(130.7434,49.9642) .. controls (130.7434,50.6850) and (131.3278,51.2693) ..(132.0486,51.2693) .. controls (132.7693,51.2693) and (133.3537,50.6850) ..(133.3537,49.9642);
\path[fill=chfourfig5c0,fill opacity=0.5,draw=chfourfig5c0,line width=0.3985pt,draw opacity=0.5] (108.0681,70.9657) .. controls (108.0681,70.2449) and (107.4838,69.6606) ..(106.7630,69.6606) .. controls (106.0422,69.6606) and (105.4579,70.2449) ..(105.4579,70.9657) .. controls (105.4579,71.6865) and (106.0422,72.2708) ..(106.7630,72.2708) .. controls (107.4838,72.2708) and (108.0681,71.6865) ..(108.0681,70.9657);
\path[fill=chfourfig5c0,fill opacity=0.5,draw=chfourfig5c0,line width=0.3985pt,draw opacity=0.5] (60.6753,107.5589) .. controls (60.6753,106.8381) and (60.0909,106.2538) ..(59.3701,106.2538) .. controls (58.6493,106.2538) and (58.0650,106.8381) ..(58.0650,107.5589) .. controls (58.0650,108.2797) and (58.6493,108.8640) ..(59.3701,108.8640) .. controls (60.0909,108.8640) and (60.6753,108.2797) ..(60.6753,107.5589);
\path[fill=chfourfig5c0,fill opacity=0.5,draw=chfourfig5c0,line width=0.3985pt,draw opacity=0.5] (98.6334,86.4578) .. controls (98.6334,85.7370) and (98.0491,85.1527) ..(97.3283,85.1527) .. controls (96.6075,85.1527) and (96.0232,85.7370) ..(96.0232,86.4578) .. controls (96.0232,87.1786) and (96.6075,87.7629) ..(97.3283,87.7629) .. controls (98.0491,87.7629) and (98.6334,87.1786) ..(98.6334,86.4578);
\path[fill=chfourfig5c0,fill opacity=0.5,draw=chfourfig5c0,line width=0.3985pt,draw opacity=0.5] (156.6665,27.7870) .. controls (156.6665,27.0663) and (156.0822,26.4819) ..(155.3614,26.4819) .. controls (154.6406,26.4819) and (154.0563,27.0663) ..(154.0563,27.7870) .. controls (154.0563,28.5078) and (154.6406,29.0922) ..(155.3614,29.0922) .. controls (156.0822,29.0922) and (156.6665,28.5078) ..(156.6665,27.7870);
\path[fill=chfourfig5c0,fill opacity=0.5,draw=chfourfig5c0,line width=0.3985pt,draw opacity=0.5] (138.7734,51.6279) .. controls (138.7734,50.9072) and (138.1891,50.3228) ..(137.4683,50.3228) .. controls (136.7475,50.3228) and (136.1631,50.9072) ..(136.1631,51.6279) .. controls (136.1631,52.3488) and (136.7475,52.9331) ..(137.4683,52.9331) .. controls (138.1891,52.9331) and (138.7734,52.3488) ..(138.7734,51.6279);
\path[fill=chfourfig5c0,fill opacity=0.5,draw=chfourfig5c0,line width=0.3985pt,draw opacity=0.5] (140.0088,31.3637) .. controls (140.0088,30.6429) and (139.4245,30.0586) ..(138.7037,30.0586) .. controls (137.9829,30.0586) and (137.3986,30.6429) ..(137.3986,31.3637) .. controls (137.3986,32.0845) and (137.9829,32.6688) ..(138.7037,32.6688) .. controls (139.4245,32.6688) and (140.0088,32.0845) ..(140.0088,31.3637);
\path[fill=chfourfig5c0,fill opacity=0.5,draw=chfourfig5c0,line width=0.3985pt,draw opacity=0.5] (74.7427,89.5462) .. controls (74.7427,88.8255) and (74.1584,88.2411) ..(73.4376,88.2411) .. controls (72.7168,88.2411) and (72.1324,88.8255) ..(72.1324,89.5462) .. controls (72.1324,90.2671) and (72.7168,90.8514) ..(73.4376,90.8514) .. controls (74.1584,90.8514) and (74.7427,90.2671) ..(74.7427,89.5462);
\path[fill=chfourfig5c0,fill opacity=0.5,draw=chfourfig5c0,line width=0.3985pt,draw opacity=0.5] (114.8329,77.3020) .. controls (114.8329,76.5812) and (114.2486,75.9969) ..(113.5277,75.9969) .. controls (112.8069,75.9969) and (112.2226,76.5812) ..(112.2226,77.3020) .. controls (112.2226,78.0228) and (112.8069,78.6071) ..(113.5277,78.6071) .. controls (114.2486,78.6071) and (114.8329,78.0228) ..(114.8329,77.3020);
\path[fill=chfourfig5c0,fill opacity=0.5,draw=chfourfig5c0,line width=0.3985pt,draw opacity=0.5] (141.3438,54.9057) .. controls (141.3438,54.1849) and (140.7595,53.6006) ..(140.0387,53.6006) .. controls (139.3179,53.6006) and (138.7336,54.1849) ..(138.7336,54.9057) .. controls (138.7336,55.6265) and (139.3179,56.2108) ..(140.0387,56.2108) .. controls (140.7595,56.2108) and (141.3438,55.6265) ..(141.3438,54.9057);
\path[fill=chfourfig5c0,fill opacity=0.5,draw=chfourfig5c0,line width=0.3985pt,draw opacity=0.5] (74.4438,68.4252) .. controls (74.4438,67.7044) and (73.8595,67.1201) ..(73.1387,67.1201) .. controls (72.4179,67.1201) and (71.8336,67.7044) ..(71.8336,68.4252) .. controls (71.8336,69.1460) and (72.4179,69.7303) ..(73.1387,69.7303) .. controls (73.8595,69.7303) and (74.4438,69.1460) ..(74.4438,68.4252);
\path[fill=chfourfig5c0,fill opacity=0.5,draw=chfourfig5c0,line width=0.3985pt,draw opacity=0.5] (137.8469,48.0115) .. controls (137.8469,47.2906) and (137.2625,46.7063) ..(136.5417,46.7063) .. controls (135.8210,46.7063) and (135.2366,47.2906) ..(135.2366,48.0115) .. controls (135.2366,48.7323) and (135.8210,49.3166) ..(136.5417,49.3166) .. controls (137.2625,49.3166) and (137.8469,48.7323) ..(137.8469,48.0115);
\path[fill=chfourfig5c0,fill opacity=0.5,draw=chfourfig5c0,line width=0.3985pt,draw opacity=0.5] (140.0486,58.3428) .. controls (140.0486,57.6221) and (139.4643,57.0377) ..(138.7435,57.0377) .. controls (138.0227,57.0377) and (137.4384,57.6221) ..(137.4384,58.3428) .. controls (137.4384,59.0637) and (138.0227,59.6480) ..(138.7435,59.6480) .. controls (139.4643,59.6480) and (140.0486,59.0637) ..(140.0486,58.3428);
\path[fill=chfourfig5c0,fill opacity=0.5,draw=chfourfig5c0,line width=0.3985pt,draw opacity=0.5] (88.5411,74.1040) .. controls (88.5411,73.3832) and (87.9568,72.7988) ..(87.2360,72.7988) .. controls (86.5152,72.7988) and (85.9309,73.3832) ..(85.9309,74.1040) .. controls (85.9309,74.8248) and (86.5152,75.4091) ..(87.2360,75.4091) .. controls (87.9568,75.4091) and (88.5411,74.8248) ..(88.5411,74.1040);
\path[fill=chfourfig5c0,fill opacity=0.5,draw=chfourfig5c0,line width=0.3985pt,draw opacity=0.5] (70.6082,69.9893) .. controls (70.6082,69.2685) and (70.0238,68.6842) ..(69.3030,68.6842) .. controls (68.5822,68.6842) and (67.9979,69.2685) ..(67.9979,69.9893) .. controls (67.9979,70.7101) and (68.5822,71.2944) ..(69.3030,71.2944) .. controls (70.0238,71.2944) and (70.6082,70.7101) ..(70.6082,69.9893);
\path[fill=chfourfig5c0,fill opacity=0.5,draw=chfourfig5c0,line width=0.3985pt,draw opacity=0.5] (71.7240,97.9947) .. controls (71.7240,97.2739) and (71.1396,96.6896) ..(70.4189,96.6896) .. controls (69.6981,96.6896) and (69.1137,97.2739) ..(69.1137,97.9947) .. controls (69.1137,98.7155) and (69.6981,99.2998) ..(70.4189,99.2998) .. controls (71.1396,99.2998) and (71.7240,98.7155) ..(71.7240,97.9947);
\path[fill=chfourfig5c0,fill opacity=0.5,draw=chfourfig5c0,line width=0.3985pt,draw opacity=0.5] (62.7176,100.8938) .. controls (62.7176,100.1730) and (62.1333,99.5887) ..(61.4125,99.5887) .. controls (60.6917,99.5887) and (60.1074,100.1730) ..(60.1074,100.8938) .. controls (60.1074,101.6147) and (60.6917,102.1990) ..(61.4125,102.1990) .. controls (62.1333,102.1990) and (62.7176,101.6147) ..(62.7176,100.8938);
\path[fill=chfourfig5c0,fill opacity=0.5,draw=chfourfig5c0,line width=0.3985pt,draw opacity=0.5] (111.9536,78.4178) .. controls (111.9536,77.6970) and (111.3693,77.1127) ..(110.6485,77.1127) .. controls (109.9277,77.1127) and (109.3434,77.6970) ..(109.3434,78.4178) .. controls (109.3434,79.1386) and (109.9277,79.7229) ..(110.6485,79.7229) .. controls (111.3693,79.7229) and (111.9536,79.1386) ..(111.9536,78.4178);
\path[fill=chfourfig5c0,fill opacity=0.5,draw=chfourfig5c0,line width=0.3985pt,draw opacity=0.5] (147.3812,42.0836) .. controls (147.3812,41.3628) and (146.7969,40.7785) ..(146.0761,40.7785) .. controls (145.3553,40.7785) and (144.7710,41.3628) ..(144.7710,42.0836) .. controls (144.7710,42.8044) and (145.3553,43.3887) ..(146.0761,43.3887) .. controls (146.7969,43.3887) and (147.3812,42.8044) ..(147.3812,42.0836);
\path[fill=chfourfig5c0,fill opacity=0.5,draw=chfourfig5c0,line width=0.3985pt,draw opacity=0.5] (95.4055,72.0317) .. controls (95.4055,71.3109) and (94.8212,70.7266) ..(94.1003,70.7266) .. controls (93.3796,70.7266) and (92.7952,71.3109) ..(92.7952,72.0317) .. controls (92.7952,72.7525) and (93.3796,73.3368) ..(94.1003,73.3368) .. controls (94.8212,73.3368) and (95.4055,72.7525) ..(95.4055,72.0317);
\path[fill=chfourfig5c0,fill opacity=0.5,draw=chfourfig5c0,line width=0.3985pt,draw opacity=0.5] (154.2854,45.4411) .. controls (154.2854,44.7203) and (153.7011,44.1360) ..(152.9803,44.1360) .. controls (152.2595,44.1360) and (151.6752,44.7203) ..(151.6752,45.4411) .. controls (151.6752,46.1619) and (152.2595,46.7462) ..(152.9803,46.7462) .. controls (153.7011,46.7462) and (154.2854,46.1619) ..(154.2854,45.4411);
\path[fill=chfourfig5c0,fill opacity=0.5,draw=chfourfig5c0,line width=0.3985pt,draw opacity=0.5] (119.4855,57.2469) .. controls (119.4855,56.5261) and (118.9012,55.9418) ..(118.1804,55.9418) .. controls (117.4596,55.9418) and (116.8752,56.5261) ..(116.8752,57.2469) .. controls (116.8752,57.9677) and (117.4596,58.5521) ..(118.1804,58.5521) .. controls (118.9012,58.5521) and (119.4855,57.9677) ..(119.4855,57.2469);
\path[fill=chfourfig5c0,fill opacity=0.5,draw=chfourfig5c0,line width=0.3985pt,draw opacity=0.5] (123.5304,50.0738) .. controls (123.5304,49.3530) and (122.9460,48.7686) ..(122.2253,48.7686) .. controls (121.5045,48.7686) and (120.9201,49.3530) ..(120.9201,50.0738) .. controls (120.9201,50.7946) and (121.5045,51.3789) ..(122.2253,51.3789) .. controls (122.9460,51.3789) and (123.5304,50.7946) ..(123.5304,50.0738);
\path[fill=chfourfig5c0,fill opacity=0.5,draw=chfourfig5c0,line width=0.3985pt,draw opacity=0.5] (93.6421,82.6420) .. controls (93.6421,81.9212) and (93.0577,81.3369) ..(92.3369,81.3369) .. controls (91.6161,81.3369) and (91.0318,81.9212) ..(91.0318,82.6420) .. controls (91.0318,83.3629) and (91.6161,83.9472) ..(92.3369,83.9472) .. controls (93.0577,83.9472) and (93.6421,83.3629) ..(93.6421,82.6420);
\path[fill=chfourfig5c0,fill opacity=0.5,draw=chfourfig5c0,line width=0.3985pt,draw opacity=0.5] (145.9665,61.2719) .. controls (145.9665,60.5511) and (145.3822,59.9668) ..(144.6614,59.9668) .. controls (143.9406,59.9668) and (143.3563,60.5511) ..(143.3563,61.2719) .. controls (143.3563,61.9927) and (143.9406,62.5770) ..(144.6614,62.5770) .. controls (145.3822,62.5770) and (145.9665,61.9927) ..(145.9665,61.2719);
\path[fill=chfourfig5c0,fill opacity=0.5,draw=chfourfig5c0,line width=0.3985pt,draw opacity=0.5] (87.5249,80.4702) .. controls (87.5249,79.7494) and (86.9406,79.1651) ..(86.2198,79.1651) .. controls (85.4990,79.1651) and (84.9147,79.7494) ..(84.9147,80.4702) .. controls (84.9147,81.1910) and (85.4990,81.7753) ..(86.2198,81.7753) .. controls (86.9406,81.7753) and (87.5249,81.1910) ..(87.5249,80.4702);
\path[fill=chfourfig5c0,fill opacity=0.5,draw=chfourfig5c0,line width=0.3985pt,draw opacity=0.5] (96.0929,86.8165) .. controls (96.0929,86.0957) and (95.5086,85.5113) ..(94.7878,85.5113) .. controls (94.0670,85.5113) and (93.4827,86.0957) ..(93.4827,86.8165) .. controls (93.4827,87.5372) and (94.0670,88.1216) ..(94.7878,88.1216) .. controls (95.5086,88.1216) and (96.0929,87.5372) ..(96.0929,86.8165);
\path[fill=chfourfig5c0,fill opacity=0.5,draw=chfourfig5c0,line width=0.3985pt,draw opacity=0.5] (155.5109,32.8581) .. controls (155.5109,32.1373) and (154.9265,31.5530) ..(154.2057,31.5530) .. controls (153.4849,31.5530) and (152.9006,32.1373) ..(152.9006,32.8581) .. controls (152.9006,33.5789) and (153.4849,34.1632) ..(154.2057,34.1632) .. controls (154.9265,34.1632) and (155.5109,33.5789) ..(155.5109,32.8581);
\path[fill=chfourfig5c0,fill opacity=0.5,draw=chfourfig5c0,line width=0.3985pt,draw opacity=0.5] (87.0467,99.5489) .. controls (87.0467,98.8281) and (86.4624,98.2437) ..(85.7416,98.2437) .. controls (85.0208,98.2437) and (84.4365,98.8281) ..(84.4365,99.5489) .. controls (84.4365,100.2697) and (85.0208,100.8540) ..(85.7416,100.8540) .. controls (86.4624,100.8540) and (87.0467,100.2697) ..(87.0467,99.5489);
\path[fill=chfourfig5c0,fill opacity=0.5,draw=chfourfig5c0,line width=0.3985pt,draw opacity=0.5] (89.4577,69.5012) .. controls (89.4577,68.7803) and (88.8734,68.1960) ..(88.1526,68.1960) .. controls (87.4318,68.1960) and (86.8474,68.7803) ..(86.8474,69.5012) .. controls (86.8474,70.2220) and (87.4318,70.8063) ..(88.1526,70.8063) .. controls (88.8734,70.8063) and (89.4577,70.2220) ..(89.4577,69.5012);
\path[fill=chfourfig5c0,fill opacity=0.5,draw=chfourfig5c0,line width=0.3985pt,draw opacity=0.5] (130.6637,65.4961) .. controls (130.6637,64.7753) and (130.0794,64.1910) ..(129.3586,64.1910) .. controls (128.6378,64.1910) and (128.0535,64.7753) ..(128.0535,65.4961) .. controls (128.0535,66.2169) and (128.6378,66.8012) ..(129.3586,66.8012) .. controls (130.0794,66.8012) and (130.6637,66.2169) ..(130.6637,65.4961);
\path[fill=chfourfig5c0,fill opacity=0.5,draw=chfourfig5c0,line width=0.3985pt,draw opacity=0.5] (70.7078,93.9199) .. controls (70.7078,93.1991) and (70.1234,92.6148) ..(69.4026,92.6148) .. controls (68.6819,92.6148) and (68.0975,93.1991) ..(68.0975,93.9199) .. controls (68.0975,94.6407) and (68.6819,95.2250) ..(69.4026,95.2250) .. controls (70.1234,95.2250) and (70.7078,94.6407) ..(70.7078,93.9199);
\path[fill=chfourfig5c0,fill opacity=0.5,draw=chfourfig5c0,line width=0.3985pt,draw opacity=0.5] (120.7209,53.0825) .. controls (120.7209,52.3617) and (120.1365,51.7774) ..(119.4157,51.7774) .. controls (118.6949,51.7774) and (118.1106,52.3617) ..(118.1106,53.0825) .. controls (118.1106,53.8033) and (118.6949,54.3876) ..(119.4157,54.3876) .. controls (120.1365,54.3876) and (120.7209,53.8033) ..(120.7209,53.0825);
\path[fill=chfourfig5c0,fill opacity=0.5,draw=chfourfig5c0,line width=0.3985pt,draw opacity=0.5] (80.9993,95.3645) .. controls (80.9993,94.6437) and (80.4150,94.0594) ..(79.6942,94.0594) .. controls (78.9734,94.0594) and (78.3891,94.6437) ..(78.3891,95.3645) .. controls (78.3891,96.0853) and (78.9734,96.6696) ..(79.6942,96.6696) .. controls (80.4150,96.6696) and (80.9993,96.0853) ..(80.9993,95.3645);
\path[fill=chfourfig5c0,fill opacity=0.5,draw=chfourfig5c0,line width=0.3985pt,draw opacity=0.5] (144.3227,33.5953) .. controls (144.3227,32.8745) and (143.7383,32.2902) ..(143.0175,32.2902) .. controls (142.2968,32.2902) and (141.7124,32.8745) ..(141.7124,33.5953) .. controls (141.7124,34.3161) and (142.2968,34.9005) ..(143.0175,34.9005) .. controls (143.7383,34.9005) and (144.3227,34.3161) ..(144.3227,33.5953);
\path[fill=chfourfig5c0,fill opacity=0.5,draw=chfourfig5c0,line width=0.3985pt,draw opacity=0.5] (72.3616,76.3555) .. controls (72.3616,75.6348) and (71.7773,75.0504) ..(71.0565,75.0504) .. controls (70.3357,75.0504) and (69.7513,75.6348) ..(69.7513,76.3555) .. controls (69.7513,77.0763) and (70.3357,77.6607) ..(71.0565,77.6607) .. controls (71.7773,77.6607) and (72.3616,77.0763) ..(72.3616,76.3555);
\path[fill=chfourfig5c0,fill opacity=0.5,draw=chfourfig5c0,line width=0.3985pt,draw opacity=0.5] (102.6683,73.5261) .. controls (102.6683,72.8053) and (102.0840,72.2210) ..(101.3632,72.2210) .. controls (100.6424,72.2210) and (100.0581,72.8053) ..(100.0581,73.5261) .. controls (100.0581,74.2469) and (100.6424,74.8312) ..(101.3632,74.8312) .. controls (102.0840,74.8312) and (102.6683,74.2469) ..(102.6683,73.5261);
\path[fill=chfourfig5c0,fill opacity=0.5,draw=chfourfig5c0,line width=0.3985pt,draw opacity=0.5] (113.9462,51.2892) .. controls (113.9462,50.5684) and (113.3619,49.9841) ..(112.6411,49.9841) .. controls (111.9203,49.9841) and (111.3360,50.5684) ..(111.3360,51.2892) .. controls (111.3360,52.0100) and (111.9203,52.5943) ..(112.6411,52.5943) .. controls (113.3619,52.5943) and (113.9462,52.0100) ..(113.9462,51.2892);
\path[fill=chfourfig5c0,fill opacity=0.5,draw=chfourfig5c0,line width=0.3985pt,draw opacity=0.5] (67.0315,73.1675) .. controls (67.0315,72.4467) and (66.4472,71.8623) ..(65.7264,71.8623) .. controls (65.0056,71.8623) and (64.4213,72.4467) ..(64.4213,73.1675) .. controls (64.4213,73.8883) and (65.0056,74.4726) ..(65.7264,74.4726) .. controls (66.4472,74.4726) and (67.0315,73.8883) ..(67.0315,73.1675);
\path[fill=chfourfig5c0,fill opacity=0.5,draw=chfourfig5c0,line width=0.3985pt,draw opacity=0.5] (151.5258,53.2818) .. controls (151.5258,52.5610) and (150.9414,51.9767) ..(150.2206,51.9767) .. controls (149.4998,51.9767) and (148.9155,52.5610) ..(148.9155,53.2818) .. controls (148.9155,54.0026) and (149.4998,54.5869) ..(150.2206,54.5869) .. controls (150.9414,54.5869) and (151.5258,54.0026) ..(151.5258,53.2818);
\path[fill=chfourfig5c0,fill opacity=0.5,draw=chfourfig5c0,line width=0.3985pt,draw opacity=0.5] (64.3515,80.3606) .. controls (64.3515,79.6398) and (63.7672,79.0555) ..(63.0464,79.0555) .. controls (62.3256,79.0555) and (61.7413,79.6398) ..(61.7413,80.3606) .. controls (61.7413,81.0814) and (62.3256,81.6657) ..(63.0464,81.6657) .. controls (63.7672,81.6657) and (64.3515,81.0814) ..(64.3515,80.3606);
\path[fill=chfourfig5c0,fill opacity=0.5,draw=chfourfig5c0,line width=0.3985pt,draw opacity=0.5] (107.1317,66.8511) .. controls (107.1317,66.1302) and (106.5473,65.5459) ..(105.8265,65.5459) .. controls (105.1057,65.5459) and (104.5214,66.1302) ..(104.5214,66.8511) .. controls (104.5214,67.5719) and (105.1057,68.1562) ..(105.8265,68.1562) .. controls (106.5473,68.1562) and (107.1317,67.5719) ..(107.1317,66.8511);
\path[fill=chfourfig5c0,fill opacity=0.5,draw=chfourfig5c0,line width=0.3985pt,draw opacity=0.5] (83.8387,77.3717) .. controls (83.8387,76.6510) and (83.2544,76.0666) ..(82.5336,76.0666) .. controls (81.8128,76.0666) and (81.2285,76.6510) ..(81.2285,77.3717) .. controls (81.2285,78.0926) and (81.8128,78.6769) ..(82.5336,78.6769) .. controls (83.2544,78.6769) and (83.8387,78.0926) ..(83.8387,77.3717);
\path[fill=chfourfig5c0,fill opacity=0.5,draw=chfourfig5c0,line width=0.3985pt,draw opacity=0.5] (97.7667,57.2868) .. controls (97.7667,56.5660) and (97.1823,55.9817) ..(96.4615,55.9817) .. controls (95.7407,55.9817) and (95.1564,56.5660) ..(95.1564,57.2868) .. controls (95.1564,58.0076) and (95.7407,58.5919) ..(96.4615,58.5919) .. controls (97.1823,58.5919) and (97.7667,58.0076) ..(97.7667,57.2868);
\path[fill=chfourfig5c0,fill opacity=0.5,draw=chfourfig5c0,line width=0.3985pt,draw opacity=0.5] (76.3567,80.6694) .. controls (76.3567,79.9486) and (75.7723,79.3643) ..(75.0515,79.3643) .. controls (74.3307,79.3643) and (73.7464,79.9486) ..(73.7464,80.6694) .. controls (73.7464,81.3902) and (74.3307,81.9745) ..(75.0515,81.9745) .. controls (75.7723,81.9745) and (76.3567,81.3902) ..(76.3567,80.6694);
\path[fill=chfourfig5c0,fill opacity=0.5,draw=chfourfig5c0,line width=0.3985pt,draw opacity=0.5] (92.9048,80.4104) .. controls (92.9048,79.6896) and (92.3205,79.1053) ..(91.5997,79.1053) .. controls (90.8789,79.1053) and (90.2946,79.6896) ..(90.2946,80.4104) .. controls (90.2946,81.1312) and (90.8789,81.7155) ..(91.5997,81.7155) .. controls (92.3205,81.7155) and (92.9048,81.1312) ..(92.9048,80.4104);
\path[fill=chfourfig5c0,fill opacity=0.5,draw=chfourfig5c0,line width=0.3985pt,draw opacity=0.5] (157.1647,31.8220) .. controls (157.1647,31.1012) and (156.5804,30.5168) ..(155.8596,30.5168) .. controls (155.1388,30.5168) and (154.5544,31.1012) ..(154.5544,31.8220) .. controls (154.5544,32.5428) and (155.1388,33.1271) ..(155.8596,33.1271) .. controls (156.5804,33.1271) and (157.1647,32.5428) ..(157.1647,31.8220);
\path[fill=chfourfig5c0,fill opacity=0.5,draw=chfourfig5c0,line width=0.3985pt,draw opacity=0.5] (92.3369,66.7514) .. controls (92.3369,66.0306) and (91.7526,65.4463) ..(91.0318,65.4463) .. controls (90.3110,65.4463) and (89.7267,66.0306) ..(89.7267,66.7514) .. controls (89.7267,67.4722) and (90.3110,68.0565) ..(91.0318,68.0565) .. controls (91.7526,68.0565) and (92.3369,67.4722) ..(92.3369,66.7514);
\path[fill=chfourfig5c0,fill opacity=0.5,draw=chfourfig5c0,line width=0.3985pt,draw opacity=0.5] (117.4630,69.9495) .. controls (117.4630,69.2287) and (116.8787,68.6444) ..(116.1579,68.6444) .. controls (115.4371,68.6444) and (114.8528,69.2287) ..(114.8528,69.9495) .. controls (114.8528,70.6703) and (115.4371,71.2546) ..(116.1579,71.2546) .. controls (116.8787,71.2546) and (117.4630,70.6703) ..(117.4630,69.9495);
\path[fill=chfourfig5c0,fill opacity=0.5,draw=chfourfig5c0,line width=0.3985pt,draw opacity=0.5] (143.1869,32.6987) .. controls (143.1869,31.9779) and (142.6026,31.3936) ..(141.8818,31.3936) .. controls (141.1610,31.3936) and (140.5767,31.9779) ..(140.5767,32.6987) .. controls (140.5767,33.4195) and (141.1610,34.0038) ..(141.8818,34.0038) .. controls (142.6026,34.0038) and (143.1869,33.4195) ..(143.1869,32.6987);
\path[fill=chfourfig5c0,fill opacity=0.5,draw=chfourfig5c0,line width=0.3985pt,draw opacity=0.5] (90.7827,93.0233) .. controls (90.7827,92.3024) and (90.1984,91.7181) ..(89.4776,91.7181) .. controls (88.7568,91.7181) and (88.1725,92.3024) ..(88.1725,93.0233) .. controls (88.1725,93.7440) and (88.7568,94.3284) ..(89.4776,94.3284) .. controls (90.1984,94.3284) and (90.7827,93.7440) ..(90.7827,93.0233);
\path[fill=chfourfig5c0,fill opacity=0.5,draw=chfourfig5c0,line width=0.3985pt,draw opacity=0.5] (102.2001,65.7452) .. controls (102.2001,65.0244) and (101.6158,64.4401) ..(100.8950,64.4401) .. controls (100.1742,64.4401) and (99.5898,65.0244) ..(99.5898,65.7452) .. controls (99.5898,66.4660) and (100.1742,67.0503) ..(100.8950,67.0503) .. controls (101.6158,67.0503) and (102.2001,66.4660) ..(102.2001,65.7452);
\path[fill=chfourfig5c0,fill opacity=0.5,draw=chfourfig5c0,line width=0.3985pt,draw opacity=0.5] (70.4288,90.7119) .. controls (70.4288,89.9911) and (69.8445,89.4068) ..(69.1237,89.4068) .. controls (68.4029,89.4068) and (67.8186,89.9911) ..(67.8186,90.7119) .. controls (67.8186,91.4327) and (68.4029,92.0170) ..(69.1237,92.0170) .. controls (69.8445,92.0170) and (70.4288,91.4327) ..(70.4288,90.7119);
\path[fill=chfourfig5c0,fill opacity=0.5,draw=chfourfig5c0,line width=0.3985pt,draw opacity=0.5] (53.7213,89.3171) .. controls (53.7213,88.5963) and (53.1369,88.0120) ..(52.4161,88.0120) .. controls (51.6953,88.0120) and (51.1110,88.5963) ..(51.1110,89.3171) .. controls (51.1110,90.0379) and (51.6953,90.6222) ..(52.4161,90.6222) .. controls (53.1369,90.6222) and (53.7213,90.0379) ..(53.7213,89.3171);
\path[fill=chfourfig5c0,fill opacity=0.5,draw=chfourfig5c0,line width=0.3985pt,draw opacity=0.5] (88.3120,78.8562) .. controls (88.3120,78.1354) and (87.7277,77.5511) ..(87.0069,77.5511) .. controls (86.2861,77.5511) and (85.7017,78.1354) ..(85.7017,78.8562) .. controls (85.7017,79.5770) and (86.2861,80.1613) ..(87.0069,80.1613) .. controls (87.7277,80.1613) and (88.3120,79.5770) ..(88.3120,78.8562);
\path[fill=chfourfig5c0,fill opacity=0.5,draw=chfourfig5c0,line width=0.3985pt,draw opacity=0.5] (135.2964,44.3950) .. controls (135.2964,43.6742) and (134.7121,43.0898) ..(133.9913,43.0898) .. controls (133.2705,43.0898) and (132.6862,43.6742) ..(132.6862,44.3950) .. controls (132.6862,45.1158) and (133.2705,45.7001) ..(133.9913,45.7001) .. controls (134.7121,45.7001) and (135.2964,45.1158) ..(135.2964,44.3950);
\path[fill=chfourfig5c0,fill opacity=0.5,draw=chfourfig5c0,line width=0.3985pt,draw opacity=0.5] (64.7401,108.1069) .. controls (64.7401,107.3861) and (64.1557,106.8018) ..(63.4349,106.8018) .. controls (62.7141,106.8018) and (62.1298,107.3861) ..(62.1298,108.1069) .. controls (62.1298,108.8277) and (62.7141,109.4120) ..(63.4349,109.4120) .. controls (64.1557,109.4120) and (64.7401,108.8277) ..(64.7401,108.1069);
\path[fill=chfourfig5c0,fill opacity=0.5,draw=chfourfig5c0,line width=0.3985pt,draw opacity=0.5] (72.4114,79.0056) .. controls (72.4114,78.2848) and (71.8271,77.7005) ..(71.1063,77.7005) .. controls (70.3855,77.7005) and (69.8012,78.2848) ..(69.8012,79.0056) .. controls (69.8012,79.7264) and (70.3855,80.3108) ..(71.1063,80.3108) .. controls (71.8271,80.3108) and (72.4114,79.7264) ..(72.4114,79.0056);
\path[fill=chfourfig5c0,fill opacity=0.5,draw=chfourfig5c0,line width=0.3985pt,draw opacity=0.5] (97.2386,87.3943) .. controls (97.2386,86.6735) and (96.6543,86.0892) ..(95.9335,86.0892) .. controls (95.2127,86.0892) and (94.6284,86.6735) ..(94.6284,87.3943) .. controls (94.6284,88.1151) and (95.2127,88.6994) ..(95.9335,88.6994) .. controls (96.6543,88.6994) and (97.2386,88.1151) ..(97.2386,87.3943);
\path[fill=chfourfig5c0,fill opacity=0.5,draw=chfourfig5c0,line width=0.3985pt,draw opacity=0.5] (84.9047,85.8600) .. controls (84.9047,85.1392) and (84.3204,84.5549) ..(83.5996,84.5549) .. controls (82.8788,84.5549) and (82.2945,85.1392) ..(82.2945,85.8600) .. controls (82.2945,86.5808) and (82.8788,87.1651) ..(83.5996,87.1651) .. controls (84.3204,87.1651) and (84.9047,86.5808) ..(84.9047,85.8600);
\path[fill=chfourfig5c0,fill opacity=0.5,draw=chfourfig5c0,line width=0.3985pt,draw opacity=0.5] (59.7089,93.4915) .. controls (59.7089,92.7707) and (59.1245,92.1864) ..(58.4038,92.1864) .. controls (57.6830,92.1864) and (57.0986,92.7707) ..(57.0986,93.4915) .. controls (57.0986,94.2123) and (57.6830,94.7966) ..(58.4038,94.7966) .. controls (59.1245,94.7966) and (59.7089,94.2123) ..(59.7089,93.4915);
\path[fill=chfourfig5c0,fill opacity=0.5,draw=chfourfig5c0,line width=0.3985pt,draw opacity=0.5] (108.7755,52.1560) .. controls (108.7755,51.4352) and (108.1912,50.8508) ..(107.4704,50.8508) .. controls (106.7496,50.8508) and (106.1653,51.4352) ..(106.1653,52.1560) .. controls (106.1653,52.8768) and (106.7496,53.4611) ..(107.4704,53.4611) .. controls (108.1912,53.4611) and (108.7755,52.8768) ..(108.7755,52.1560);
\path[fill=chfourfig5c0,fill opacity=0.5,draw=chfourfig5c0,line width=0.3985pt,draw opacity=0.5] (57.2879,112.6798) .. controls (57.2879,111.9590) and (56.7036,111.3747) ..(55.9828,111.3747) .. controls (55.2620,111.3747) and (54.6777,111.9590) ..(54.6777,112.6798) .. controls (54.6777,113.4006) and (55.2620,113.9849) ..(55.9828,113.9849) .. controls (56.7036,113.9849) and (57.2879,113.4006) ..(57.2879,112.6798);
\path[fill=chfourfig5c0,fill opacity=0.5,draw=chfourfig5c0,line width=0.3985pt,draw opacity=0.5] (99.8289,74.3829) .. controls (99.8289,73.6621) and (99.2446,73.0778) ..(98.5238,73.0778) .. controls (97.8030,73.0778) and (97.2187,73.6621) ..(97.2187,74.3829) .. controls (97.2187,75.1037) and (97.8030,75.6880) ..(98.5238,75.6880) .. controls (99.2446,75.6880) and (99.8289,75.1037) ..(99.8289,74.3829);
\path[fill=chfourfig5c0,fill opacity=0.5,draw=chfourfig5c0,line width=0.3985pt,draw opacity=0.5] (59.3602,101.5115) .. controls (59.3602,100.7907) and (58.7759,100.2064) ..(58.0551,100.2064) .. controls (57.3343,100.2064) and (56.7500,100.7907) ..(56.7500,101.5115) .. controls (56.7500,102.2323) and (57.3343,102.8167) ..(58.0551,102.8167) .. controls (58.7759,102.8167) and (59.3602,102.2323) ..(59.3602,101.5115);
\path[fill=chfourfig5c0,fill opacity=0.5,draw=chfourfig5c0,line width=0.3985pt,draw opacity=0.5] (100.6359,70.9557) .. controls (100.6359,70.2349) and (100.0516,69.6506) ..(99.3308,69.6506) .. controls (98.6100,69.6506) and (98.0257,70.2349) ..(98.0257,70.9557) .. controls (98.0257,71.6765) and (98.6100,72.2608) ..(99.3308,72.2608) .. controls (100.0516,72.2608) and (100.6359,71.6765) ..(100.6359,70.9557);
\path[fill=chfourfig5c0,fill opacity=0.5,draw=chfourfig5c0,line width=0.3985pt,draw opacity=0.5] (63.2855,89.2872) .. controls (63.2855,88.5664) and (62.7012,87.9821) ..(61.9804,87.9821) .. controls (61.2596,87.9821) and (60.6753,88.5664) ..(60.6753,89.2872) .. controls (60.6753,90.0080) and (61.2596,90.5923) ..(61.9804,90.5923) .. controls (62.7012,90.5923) and (63.2855,90.0080) ..(63.2855,89.2872);
\path[fill=chfourfig5c0,fill opacity=0.5,draw=chfourfig5c0,line width=0.3985pt,draw opacity=0.5] (62.2793,100.3060) .. controls (62.2793,99.5852) and (61.6949,99.0009) ..(60.9741,99.0009) .. controls (60.2534,99.0009) and (59.6690,99.5852) ..(59.6690,100.3060) .. controls (59.6690,101.0268) and (60.2534,101.6112) ..(60.9741,101.6112) .. controls (61.6949,101.6112) and (62.2793,101.0268) ..(62.2793,100.3060);
\path[fill=chfourfig5c0,fill opacity=0.5,draw=chfourfig5c0,line width=0.3985pt,draw opacity=0.5] (150.4398,61.7900) .. controls (150.4398,61.0692) and (149.8555,60.4849) ..(149.1347,60.4849) .. controls (148.4139,60.4849) and (147.8296,61.0692) ..(147.8296,61.7900) .. controls (147.8296,62.5108) and (148.4139,63.0951) ..(149.1347,63.0951) .. controls (149.8555,63.0951) and (150.4398,62.5108) ..(150.4398,61.7900);
\path[fill=chfourfig5c0,fill opacity=0.5,draw=chfourfig5c0,line width=0.3985pt,draw opacity=0.5] (149.9616,44.0463) .. controls (149.9616,43.3255) and (149.3773,42.7411) ..(148.6565,42.7411) .. controls (147.9357,42.7411) and (147.3513,43.3255) ..(147.3513,44.0463) .. controls (147.3513,44.7671) and (147.9357,45.3514) ..(148.6565,45.3514) .. controls (149.3773,45.3514) and (149.9616,44.7671) ..(149.9616,44.0463);
\path[fill=chfourfig5c0,fill opacity=0.5,draw=chfourfig5c0,line width=0.3985pt,draw opacity=0.5] (76.3467,71.0255) .. controls (76.3467,70.3047) and (75.7624,69.7203) ..(75.0416,69.7203) .. controls (74.3208,69.7203) and (73.7365,70.3047) ..(73.7365,71.0255) .. controls (73.7365,71.7463) and (74.3208,72.3306) ..(75.0416,72.3306) .. controls (75.7624,72.3306) and (76.3467,71.7463) ..(76.3467,71.0255);
\path[fill=chfourfig5c0,fill opacity=0.5,draw=chfourfig5c0,line width=0.3985pt,draw opacity=0.5] (118.4892,45.6403) .. controls (118.4892,44.9195) and (117.9049,44.3352) ..(117.1841,44.3352) .. controls (116.4633,44.3352) and (115.8790,44.9195) ..(115.8790,45.6403) .. controls (115.8790,46.3611) and (116.4633,46.9455) ..(117.1841,46.9455) .. controls (117.9049,46.9455) and (118.4892,46.3611) ..(118.4892,45.6403);
\path[fill=chfourfig5c0,fill opacity=0.5,draw=chfourfig5c0,line width=0.3985pt,draw opacity=0.5] (150.9579,52.4150) .. controls (150.9579,51.6942) and (150.3736,51.1099) ..(149.6528,51.1099) .. controls (148.9319,51.1099) and (148.3476,51.6942) ..(148.3476,52.4150) .. controls (148.3476,53.1358) and (148.9319,53.7201) ..(149.6528,53.7201) .. controls (150.3736,53.7201) and (150.9579,53.1358) ..(150.9579,52.4150);
\path[fill=chfourfig5c0,fill opacity=0.5,draw=chfourfig5c0,line width=0.3985pt,draw opacity=0.5] (61.0140,101.4318) .. controls (61.0140,100.7110) and (60.4297,100.1267) ..(59.7089,100.1267) .. controls (58.9881,100.1267) and (58.4038,100.7110) ..(58.4038,101.4318) .. controls (58.4038,102.1526) and (58.9881,102.7370) ..(59.7089,102.7370) .. controls (60.4297,102.7370) and (61.0140,102.1526) ..(61.0140,101.4318);
\path[fill=chfourfig5c0,fill opacity=0.5,draw=chfourfig5c0,line width=0.3985pt,draw opacity=0.5] (128.9501,52.6541) .. controls (128.9501,51.9333) and (128.3658,51.3490) ..(127.6450,51.3490) .. controls (126.9242,51.3490) and (126.3399,51.9333) ..(126.3399,52.6541) .. controls (126.3399,53.3749) and (126.9242,53.9592) ..(127.6450,53.9592) .. controls (128.3658,53.9592) and (128.9501,53.3749) ..(128.9501,52.6541);
\path[fill=chfourfig5c0,fill opacity=0.5,draw=chfourfig5c0,line width=0.3985pt,draw opacity=0.5] (82.6232,88.0020) .. controls (82.6232,87.2812) and (82.0389,86.6969) ..(81.3181,86.6969) .. controls (80.5973,86.6969) and (80.0130,87.2812) ..(80.0130,88.0020) .. controls (80.0130,88.7228) and (80.5973,89.3071) ..(81.3181,89.3071) .. controls (82.0389,89.3071) and (82.6232,88.7228) ..(82.6232,88.0020);
\path[fill=chfourfig5c0,fill opacity=0.5,draw=chfourfig5c0,line width=0.3985pt,draw opacity=0.5] (108.8552,68.3554) .. controls (108.8552,67.6346) and (108.2709,67.0503) ..(107.5501,67.0503) .. controls (106.8293,67.0503) and (106.2450,67.6346) ..(106.2450,68.3554) .. controls (106.2450,69.0762) and (106.8293,69.6606) ..(107.5501,69.6606) .. controls (108.2709,69.6606) and (108.8552,69.0762) ..(108.8552,68.3554);
\path[fill=chfourfig5c0,fill opacity=0.5,draw=chfourfig5c0,line width=0.3985pt,draw opacity=0.5] (116.0882,62.5372) .. controls (116.0882,61.8164) and (115.5039,61.2321) ..(114.7831,61.2321) .. controls (114.0623,61.2321) and (113.4779,61.8164) ..(113.4779,62.5372) .. controls (113.4779,63.2580) and (114.0623,63.8423) ..(114.7831,63.8423) .. controls (115.5039,63.8423) and (116.0882,63.2580) ..(116.0882,62.5372);
\path[fill=chfourfig5c0,fill opacity=0.5,draw=chfourfig5c0,line width=0.3985pt,draw opacity=0.5] (87.9832,69.2919) .. controls (87.9832,68.5711) and (87.3989,67.9868) ..(86.6781,67.9868) .. controls (85.9573,67.9868) and (85.3730,68.5711) ..(85.3730,69.2919) .. controls (85.3730,70.0127) and (85.9573,70.5971) ..(86.6781,70.5971) .. controls (87.3989,70.5971) and (87.9832,70.0127) ..(87.9832,69.2919);
\path[fill=chfourfig5c0,fill opacity=0.5,draw=chfourfig5c0,line width=0.3985pt,draw opacity=0.5] (153.2593,35.2990) .. controls (153.2593,34.5782) and (152.6750,33.9939) ..(151.9542,33.9939) .. controls (151.2334,33.9939) and (150.6490,34.5782) ..(150.6490,35.2990) .. controls (150.6490,36.0198) and (151.2334,36.6041) ..(151.9542,36.6041) .. controls (152.6750,36.6041) and (153.2593,36.0198) ..(153.2593,35.2990);
\path[fill=chfourfig5c0,fill opacity=0.5,draw=chfourfig5c0,line width=0.3985pt,draw opacity=0.5] (87.0168,71.3243) .. controls (87.0168,70.6035) and (86.4325,70.0192) ..(85.7117,70.0192) .. controls (84.9909,70.0192) and (84.4066,70.6035) ..(84.4066,71.3243) .. controls (84.4066,72.0452) and (84.9909,72.6295) ..(85.7117,72.6295) .. controls (86.4325,72.6295) and (87.0168,72.0452) ..(87.0168,71.3243);
\path[fill=chfourfig5c0,fill opacity=0.5,draw=chfourfig5c0,line width=0.3985pt,draw opacity=0.5] (155.4810,40.4896) .. controls (155.4810,39.7688) and (154.8966,39.1845) ..(154.1758,39.1845) .. controls (153.4550,39.1845) and (152.8707,39.7688) ..(152.8707,40.4896) .. controls (152.8707,41.2104) and (153.4550,41.7947) ..(154.1758,41.7947) .. controls (154.8966,41.7947) and (155.4810,41.2104) ..(155.4810,40.4896);
\path[fill=chfourfig5c0,fill opacity=0.5,draw=chfourfig5c0,line width=0.3985pt,draw opacity=0.5] (62.7176,103.9624) .. controls (62.7176,103.2416) and (62.1333,102.6572) ..(61.4125,102.6572) .. controls (60.6917,102.6572) and (60.1074,103.2416) ..(60.1074,103.9624) .. controls (60.1074,104.6832) and (60.6917,105.2675) ..(61.4125,105.2675) .. controls (62.1333,105.2675) and (62.7176,104.6832) ..(62.7176,103.9624);
\path[fill=chfourfig5c0,fill opacity=0.5,draw=chfourfig5c0,line width=0.3985pt,draw opacity=0.5] (54.7574,116.9239) .. controls (54.7574,116.2031) and (54.1731,115.6188) ..(53.4523,115.6188) .. controls (52.7315,115.6188) and (52.1471,116.2031) ..(52.1471,116.9239) .. controls (52.1471,117.6447) and (52.7315,118.2291) ..(53.4523,118.2291) .. controls (54.1731,118.2291) and (54.7574,117.6447) ..(54.7574,116.9239);
\path[fill=chfourfig5c0,fill opacity=0.5,draw=chfourfig5c0,line width=0.3985pt,draw opacity=0.5] (112.7506,70.4177) .. controls (112.7506,69.6969) and (112.1663,69.1126) ..(111.4455,69.1126) .. controls (110.7247,69.1126) and (110.1404,69.6969) ..(110.1404,70.4177) .. controls (110.1404,71.1385) and (110.7247,71.7229) ..(111.4455,71.7229) .. controls (112.1663,71.7229) and (112.7506,71.1385) ..(112.7506,70.4177);
\path[fill=chfourfig5c0,fill opacity=0.5,draw=chfourfig5c0,line width=0.3985pt,draw opacity=0.5] (70.9568,86.8065) .. controls (70.9568,86.0857) and (70.3725,85.5014) ..(69.6517,85.5014) .. controls (68.9309,85.5014) and (68.3466,86.0857) ..(68.3466,86.8065) .. controls (68.3466,87.5273) and (68.9309,88.1116) ..(69.6517,88.1116) .. controls (70.3725,88.1116) and (70.9568,87.5273) ..(70.9568,86.8065);
\path[fill=chfourfig5c0,fill opacity=0.5,draw=chfourfig5c0,line width=0.3985pt,draw opacity=0.5] (123.4606,60.8535) .. controls (123.4606,60.1327) and (122.8763,59.5484) ..(122.1555,59.5484) .. controls (121.4347,59.5484) and (120.8504,60.1327) ..(120.8504,60.8535) .. controls (120.8504,61.5743) and (121.4347,62.1586) ..(122.1555,62.1586) .. controls (122.8763,62.1586) and (123.4606,61.5743) ..(123.4606,60.8535);
\path[fill=chfourfig5c0,fill opacity=0.5,draw=chfourfig5c0,line width=0.3985pt,draw opacity=0.5] (100.6459,74.1239) .. controls (100.6459,73.4031) and (100.0616,72.8188) ..(99.3408,72.8188) .. controls (98.6200,72.8188) and (98.0356,73.4031) ..(98.0356,74.1239) .. controls (98.0356,74.8447) and (98.6200,75.4290) ..(99.3408,75.4290) .. controls (100.0616,75.4290) and (100.6459,74.8447) ..(100.6459,74.1239);
\path[fill=chfourfig5c0,fill opacity=0.5,draw=chfourfig5c0,line width=0.3985pt,draw opacity=0.5] (106.4243,80.4602) .. controls (106.4243,79.7394) and (105.8400,79.1551) ..(105.1192,79.1551) .. controls (104.3984,79.1551) and (103.8141,79.7394) ..(103.8141,80.4602) .. controls (103.8141,81.1810) and (104.3984,81.7653) ..(105.1192,81.7653) .. controls (105.8400,81.7653) and (106.4243,81.1810) ..(106.4243,80.4602);
\path[fill=chfourfig5c0,fill opacity=0.5,draw=chfourfig5c0,line width=0.3985pt,draw opacity=0.5] (156.5271,30.8855) .. controls (156.5271,30.1647) and (155.9427,29.5804) ..(155.2220,29.5804) .. controls (154.5011,29.5804) and (153.9168,30.1647) ..(153.9168,30.8855) .. controls (153.9168,31.6063) and (154.5011,32.1906) ..(155.2220,32.1906) .. controls (155.9427,32.1906) and (156.5271,31.6063) ..(156.5271,30.8855);
\path[fill=chfourfig5c0,fill opacity=0.5,draw=chfourfig5c0,line width=0.3985pt,draw opacity=0.5] (132.9352,57.8945) .. controls (132.9352,57.1737) and (132.3509,56.5894) ..(131.6301,56.5894) .. controls (130.9093,56.5894) and (130.3250,57.1737) ..(130.3250,57.8945) .. controls (130.3250,58.6153) and (130.9093,59.1997) ..(131.6301,59.1997) .. controls (132.3509,59.1997) and (132.9352,58.6153) ..(132.9352,57.8945);
\path[fill=chfourfig5c0,fill opacity=0.5,draw=chfourfig5c0,line width=0.3985pt,draw opacity=0.5] (143.2168,48.5993) .. controls (143.2168,47.8785) and (142.6325,47.2942) ..(141.9117,47.2942) .. controls (141.1909,47.2942) and (140.6066,47.8785) ..(140.6066,48.5993) .. controls (140.6066,49.3201) and (141.1909,49.9044) ..(141.9117,49.9044) .. controls (142.6325,49.9044) and (143.2168,49.3201) ..(143.2168,48.5993);
\path[fill=chfourfig5c0,fill opacity=0.5,draw=chfourfig5c0,line width=0.3985pt,draw opacity=0.5] (125.0547,44.9728) .. controls (125.0547,44.2520) and (124.4704,43.6677) ..(123.7496,43.6677) .. controls (123.0288,43.6677) and (122.4444,44.2520) ..(122.4444,44.9728) .. controls (122.4444,45.6936) and (123.0288,46.2779) ..(123.7496,46.2779) .. controls (124.4704,46.2779) and (125.0547,45.6936) ..(125.0547,44.9728);
\path[fill=chfourfig5c0,fill opacity=0.5,draw=chfourfig5c0,line width=0.3985pt,draw opacity=0.5] (88.7205,63.2844) .. controls (88.7205,62.5636) and (88.1361,61.9793) ..(87.4153,61.9793) .. controls (86.6945,61.9793) and (86.1102,62.5636) ..(86.1102,63.2844) .. controls (86.1102,64.0052) and (86.6945,64.5895) ..(87.4153,64.5895) .. controls (88.1361,64.5895) and (88.7205,64.0052) ..(88.7205,63.2844);
\path[fill=chfourfig5c0,fill opacity=0.5,draw=chfourfig5c0,line width=0.3985pt,draw opacity=0.5] (158.5196,34.2130) .. controls (158.5196,33.4922) and (157.9353,32.9079) ..(157.2145,32.9079) .. controls (156.4937,32.9079) and (155.9094,33.4922) ..(155.9094,34.2130) .. controls (155.9094,34.9338) and (156.4937,35.5182) ..(157.2145,35.5182) .. controls (157.9353,35.5182) and (158.5196,34.9338) ..(158.5196,34.2130);
\path[fill=chfourfig5c0,fill opacity=0.5,draw=chfourfig5c0,line width=0.3985pt,draw opacity=0.5] (126.8878,53.2618) .. controls (126.8878,52.5410) and (126.3035,51.9567) ..(125.5827,51.9567) .. controls (124.8619,51.9567) and (124.2776,52.5410) ..(124.2776,53.2618) .. controls (124.2776,53.9826) and (124.8619,54.5670) ..(125.5827,54.5670) .. controls (126.3035,54.5670) and (126.8878,53.9826) ..(126.8878,53.2618);
\path[fill=chfourfig5c0,fill opacity=0.5,draw=chfourfig5c0,line width=0.3985pt,draw opacity=0.5] (84.0778,68.6543) .. controls (84.0778,67.9335) and (83.4935,67.3492) ..(82.7727,67.3492) .. controls (82.0519,67.3492) and (81.4676,67.9335) ..(81.4676,68.6543) .. controls (81.4676,69.3751) and (82.0519,69.9594) ..(82.7727,69.9594) .. controls (83.4935,69.9594) and (84.0778,69.3751) ..(84.0778,68.6543);
\path[fill=chfourfig5c0,fill opacity=0.5,draw=chfourfig5c0,line width=0.3985pt,draw opacity=0.5] (78.7278,73.3767) .. controls (78.7278,72.6559) and (78.1435,72.0715) ..(77.4227,72.0715) .. controls (76.7019,72.0715) and (76.1176,72.6559) ..(76.1176,73.3767) .. controls (76.1176,74.0975) and (76.7019,74.6818) ..(77.4227,74.6818) .. controls (78.1435,74.6818) and (78.7278,74.0975) ..(78.7278,73.3767);
\path[fill=chfourfig5c0,fill opacity=0.5,draw=chfourfig5c0,line width=0.3985pt,draw opacity=0.5] (81.1487,98.6124) .. controls (81.1487,97.8916) and (80.5644,97.3073) ..(79.8436,97.3073) .. controls (79.1228,97.3073) and (78.5385,97.8916) ..(78.5385,98.6124) .. controls (78.5385,99.3332) and (79.1228,99.9175) ..(79.8436,99.9175) .. controls (80.5644,99.9175) and (81.1487,99.3332) ..(81.1487,98.6124);
\path[fill=chfourfig5c0,fill opacity=0.5,draw=chfourfig5c0,line width=0.3985pt,draw opacity=0.5] (78.0204,93.8801) .. controls (78.0204,93.1593) and (77.4361,92.5749) ..(76.7153,92.5749) .. controls (75.9945,92.5749) and (75.4102,93.1593) ..(75.4102,93.8801) .. controls (75.4102,94.6008) and (75.9945,95.1852) ..(76.7153,95.1852) .. controls (77.4361,95.1852) and (78.0204,94.6008) ..(78.0204,93.8801);
\path[fill=chfourfig5c0,fill opacity=0.5,draw=chfourfig5c0,line width=0.3985pt,draw opacity=0.5] (72.0129,72.1413) .. controls (72.0129,71.4205) and (71.4286,70.8362) ..(70.7078,70.8362) .. controls (69.9870,70.8362) and (69.4026,71.4205) ..(69.4026,72.1413) .. controls (69.4026,72.8621) and (69.9870,73.4464) ..(70.7078,73.4464) .. controls (71.4286,73.4464) and (72.0129,72.8621) ..(72.0129,72.1413);
\path[fill=chfourfig5c0,fill opacity=0.5,draw=chfourfig5c0,line width=0.3985pt,draw opacity=0.5] (139.3214,46.8857) .. controls (139.3214,46.1649) and (138.7370,45.5806) ..(138.0162,45.5806) .. controls (137.2954,45.5806) and (136.7111,46.1649) ..(136.7111,46.8857) .. controls (136.7111,47.6065) and (137.2954,48.1908) ..(138.0162,48.1908) .. controls (138.7370,48.1908) and (139.3214,47.6065) ..(139.3214,46.8857);
\path[fill=chfourfig5c0,fill opacity=0.5,draw=chfourfig5c0,line width=0.3985pt,draw opacity=0.5] (54.5282,91.8277) .. controls (54.5282,91.1069) and (53.9439,90.5226) ..(53.2231,90.5226) .. controls (52.5023,90.5226) and (51.9180,91.1069) ..(51.9180,91.8277) .. controls (51.9180,92.5485) and (52.5023,93.1328) ..(53.2231,93.1328) .. controls (53.9439,93.1328) and (54.5282,92.5485) ..(54.5282,91.8277);
\path[fill=chfourfig5c0,fill opacity=0.5,draw=chfourfig5c0,line width=0.3985pt,draw opacity=0.5] (71.4849,84.2759) .. controls (71.4849,83.5551) and (70.9005,82.9708) ..(70.1797,82.9708) .. controls (69.4590,82.9708) and (68.8746,83.5551) ..(68.8746,84.2759) .. controls (68.8746,84.9967) and (69.4590,85.5811) ..(70.1797,85.5811) .. controls (70.9005,85.5811) and (71.4849,84.9967) ..(71.4849,84.2759);
\path[fill=chfourfig5c0,fill opacity=0.5,draw=chfourfig5c0,line width=0.3985pt,draw opacity=0.5] (82.4140,84.2460) .. controls (82.4140,83.5253) and (81.8297,82.9409) ..(81.1089,82.9409) .. controls (80.3881,82.9409) and (79.8038,83.5253) ..(79.8038,84.2460) .. controls (79.8038,84.9669) and (80.3881,85.5512) ..(81.1089,85.5512) .. controls (81.8297,85.5512) and (82.4140,84.9669) ..(82.4140,84.2460);
\path[fill=chfourfig5c0,fill opacity=0.5,draw=chfourfig5c0,line width=0.3985pt,draw opacity=0.5] (53.2729,102.6174) .. controls (53.2729,101.8966) and (52.6886,101.3123) ..(51.9678,101.3123) .. controls (51.2470,101.3123) and (50.6627,101.8966) ..(50.6627,102.6174) .. controls (50.6627,103.3382) and (51.2470,103.9225) ..(51.9678,103.9225) .. controls (52.6886,103.9225) and (53.2729,103.3382) ..(53.2729,102.6174);
\path[fill=chfourfig5c0,fill opacity=0.5,draw=chfourfig5c0,line width=0.3985pt,draw opacity=0.5] (153.3987,53.5009) .. controls (153.3987,52.7802) and (152.8144,52.1958) ..(152.0936,52.1958) .. controls (151.3728,52.1958) and (150.7885,52.7802) ..(150.7885,53.5009) .. controls (150.7885,54.2218) and (151.3728,54.8061) ..(152.0936,54.8061) .. controls (152.8144,54.8061) and (153.3987,54.2218) ..(153.3987,53.5009);
\path[fill=chfourfig5c0,fill opacity=0.5,draw=chfourfig5c0,line width=0.3985pt,draw opacity=0.5] (122.4245,41.9342) .. controls (122.4245,41.2134) and (121.8402,40.6291) ..(121.1194,40.6291) .. controls (120.3986,40.6291) and (119.8143,41.2134) ..(119.8143,41.9342) .. controls (119.8143,42.6550) and (120.3986,43.2393) ..(121.1194,43.2393) .. controls (121.8402,43.2393) and (122.4245,42.6550) ..(122.4245,41.9342);
\path[fill=chfourfig5c0,fill opacity=0.5,draw=chfourfig5c0,line width=0.3985pt,draw opacity=0.5] (125.6624,57.5558) .. controls (125.6624,56.8350) and (125.0781,56.2507) ..(124.3573,56.2507) .. controls (123.6365,56.2507) and (123.0522,56.8350) ..(123.0522,57.5558) .. controls (123.0522,58.2766) and (123.6365,58.8609) ..(124.3573,58.8609) .. controls (125.0781,58.8609) and (125.6624,58.2766) ..(125.6624,57.5558);
\path[fill=chfourfig5c0,fill opacity=0.5,draw=chfourfig5c0,line width=0.3985pt,draw opacity=0.5] (68.5259,103.3746) .. controls (68.5259,102.6538) and (67.9416,102.0694) ..(67.2208,102.0694) .. controls (66.5000,102.0694) and (65.9157,102.6538) ..(65.9157,103.3746) .. controls (65.9157,104.0954) and (66.5000,104.6797) ..(67.2208,104.6797) .. controls (67.9416,104.6797) and (68.5259,104.0954) ..(68.5259,103.3746);
\path[fill=chfourfig5c0,fill opacity=0.5,draw=chfourfig5c0,line width=0.3985pt,draw opacity=0.5] (115.7395,80.6595) .. controls (115.7395,79.9386) and (115.1552,79.3543) ..(114.4344,79.3543) .. controls (113.7136,79.3543) and (113.1292,79.9386) ..(113.1292,80.6595) .. controls (113.1292,81.3802) and (113.7136,81.9646) ..(114.4344,81.9646) .. controls (115.1552,81.9646) and (115.7395,81.3802) ..(115.7395,80.6595);
\path[fill=chfourfig5c0,fill opacity=0.5,draw=chfourfig5c0,line width=0.3985pt,draw opacity=0.5] (65.9256,94.1391) .. controls (65.9256,93.4183) and (65.3413,92.8340) ..(64.6205,92.8340) .. controls (63.8997,92.8340) and (63.3154,93.4183) ..(63.3154,94.1391) .. controls (63.3154,94.8599) and (63.8997,95.4442) ..(64.6205,95.4442) .. controls (65.3413,95.4442) and (65.9256,94.8599) ..(65.9256,94.1391);
\path[fill=chfourfig5c0,fill opacity=0.5,draw=chfourfig5c0,line width=0.3985pt,draw opacity=0.5] (86.8076,74.8811) .. controls (86.8076,74.1603) and (86.2233,73.5759) ..(85.5025,73.5759) .. controls (84.7817,73.5759) and (84.1974,74.1603) ..(84.1974,74.8811) .. controls (84.1974,75.6019) and (84.7817,76.1862) ..(85.5025,76.1862) .. controls (86.2233,76.1862) and (86.8076,75.6019) ..(86.8076,74.8811);
\path[fill=chfourfig5c0,fill opacity=0.5,draw=chfourfig5c0,line width=0.3985pt,draw opacity=0.5] (102.1204,60.8037) .. controls (102.1204,60.0829) and (101.5361,59.4985) ..(100.8153,59.4985) .. controls (100.0945,59.4985) and (99.5101,60.0829) ..(99.5101,60.8037) .. controls (99.5101,61.5245) and (100.0945,62.1088) ..(100.8153,62.1088) .. controls (101.5361,62.1088) and (102.1204,61.5245) ..(102.1204,60.8037);
\path[fill=chfourfig5c0,fill opacity=0.5,draw=chfourfig5c0,line width=0.3985pt,draw opacity=0.5] (88.5013,95.1254) .. controls (88.5013,94.4046) and (87.9169,93.8203) ..(87.1962,93.8203) .. controls (86.4754,93.8203) and (85.8911,94.4046) ..(85.8911,95.1254) .. controls (85.8911,95.8462) and (86.4754,96.4305) ..(87.1962,96.4305) .. controls (87.9169,96.4305) and (88.5013,95.8462) ..(88.5013,95.1254);
\path[fill=chfourfig5c0,fill opacity=0.5,draw=chfourfig5c0,line width=0.3985pt,draw opacity=0.5] (158.0613,21.5703) .. controls (158.0613,20.8495) and (157.4770,20.2652) ..(156.7562,20.2652) .. controls (156.0354,20.2652) and (155.4511,20.8495) ..(155.4511,21.5703) .. controls (155.4511,22.2911) and (156.0354,22.8754) ..(156.7562,22.8754) .. controls (157.4770,22.8754) and (158.0613,22.2911) ..(158.0613,21.5703);
\path[fill=chfourfig5c0,fill opacity=0.5,draw=chfourfig5c0,line width=0.3985pt,draw opacity=0.5] (124.8156,76.7142) .. controls (124.8156,75.9934) and (124.2313,75.4091) ..(123.5105,75.4091) .. controls (122.7896,75.4091) and (122.2053,75.9934) ..(122.2053,76.7142) .. controls (122.2053,77.4350) and (122.7896,78.0193) ..(123.5105,78.0193) .. controls (124.2313,78.0193) and (124.8156,77.4350) ..(124.8156,76.7142);
\path[fill=chfourfig5c0,fill opacity=0.5,draw=chfourfig5c0,line width=0.3985pt,draw opacity=0.5] (54.7275,110.3286) .. controls (54.7275,109.6078) and (54.1432,109.0235) ..(53.4224,109.0235) .. controls (52.7016,109.0235) and (52.1172,109.6078) ..(52.1172,110.3286) .. controls (52.1172,111.0494) and (52.7016,111.6337) ..(53.4224,111.6337) .. controls (54.1432,111.6337) and (54.7275,111.0494) ..(54.7275,110.3286);
\path[fill=chfourfig5c0,fill opacity=0.5,draw=chfourfig5c0,line width=0.3985pt,draw opacity=0.5] (130.5840,43.8669) .. controls (130.5840,43.1461) and (129.9997,42.5618) ..(129.2789,42.5618) .. controls (128.5581,42.5618) and (127.9738,43.1461) ..(127.9738,43.8669) .. controls (127.9738,44.5878) and (128.5581,45.1721) ..(129.2789,45.1721) .. controls (129.9997,45.1721) and (130.5840,44.5878) ..(130.5840,43.8669);
\path[fill=chfourfig5c0,fill opacity=0.5,draw=chfourfig5c0,line width=0.3985pt,draw opacity=0.5] (100.7555,90.7119) .. controls (100.7555,89.9911) and (100.1712,89.4068) ..(99.4504,89.4068) .. controls (98.7296,89.4068) and (98.1452,89.9911) ..(98.1452,90.7119) .. controls (98.1452,91.4327) and (98.7296,92.0170) ..(99.4504,92.0170) .. controls (100.1712,92.0170) and (100.7555,91.4327) ..(100.7555,90.7119);
\path[fill=chfourfig5c0,fill opacity=0.5,draw=chfourfig5c0,line width=0.3985pt,draw opacity=0.5] (102.2997,66.9806) .. controls (102.2997,66.2598) and (101.7154,65.6755) ..(100.9946,65.6755) .. controls (100.2738,65.6755) and (99.6895,66.2598) ..(99.6895,66.9806) .. controls (99.6895,67.7014) and (100.2738,68.2857) ..(100.9946,68.2857) .. controls (101.7154,68.2857) and (102.2997,67.7014) ..(102.2997,66.9806);
\path[fill=chfourfig5c0,fill opacity=0.5,draw=chfourfig5c0,line width=0.3985pt,draw opacity=0.5] (63.3353,93.3122) .. controls (63.3353,92.5914) and (62.7510,92.0071) ..(62.0302,92.0071) .. controls (61.3094,92.0071) and (60.7251,92.5914) ..(60.7251,93.3122) .. controls (60.7251,94.0330) and (61.3094,94.6173) ..(62.0302,94.6173) .. controls (62.7510,94.6173) and (63.3353,94.0330) ..(63.3353,93.3122);
\path[fill=chfourfig5c0,fill opacity=0.5,draw=chfourfig5c0,line width=0.3985pt,draw opacity=0.5] (92.6657,93.7107) .. controls (92.6657,92.9899) and (92.0814,92.4056) ..(91.3606,92.4056) .. controls (90.6398,92.4056) and (90.0555,92.9899) ..(90.0555,93.7107) .. controls (90.0555,94.4315) and (90.6398,95.0158) ..(91.3606,95.0158) .. controls (92.0814,95.0158) and (92.6657,94.4315) ..(92.6657,93.7107);
\path[fill=chfourfig5c0,fill opacity=0.5,draw=chfourfig5c0,line width=0.3985pt,draw opacity=0.5] (95.1465,81.2572) .. controls (95.1465,80.5364) and (94.5621,79.9521) ..(93.8413,79.9521) .. controls (93.1205,79.9521) and (92.5362,80.5364) ..(92.5362,81.2572) .. controls (92.5362,81.9780) and (93.1205,82.5623) ..(93.8413,82.5623) .. controls (94.5621,82.5623) and (95.1465,81.9780) ..(95.1465,81.2572);
\path[fill=chfourfig5c0,fill opacity=0.5,draw=chfourfig5c0,line width=0.3985pt,draw opacity=0.5] (75.3903,88.3308) .. controls (75.3903,87.6100) and (74.8059,87.0257) ..(74.0851,87.0257) .. controls (73.3643,87.0257) and (72.7800,87.6100) ..(72.7800,88.3308) .. controls (72.7800,89.0516) and (73.3643,89.6359) ..(74.0851,89.6359) .. controls (74.8059,89.6359) and (75.3903,89.0516) ..(75.3903,88.3308);
\path[fill=chfourfig5c0,fill opacity=0.5,draw=chfourfig5c0,line width=0.3985pt,draw opacity=0.5] (67.7090,105.9350) .. controls (67.7090,105.2142) and (67.1247,104.6299) ..(66.4039,104.6299) .. controls (65.6831,104.6299) and (65.0988,105.2142) ..(65.0988,105.9350) .. controls (65.0988,106.6558) and (65.6831,107.2401) ..(66.4039,107.2401) .. controls (67.1247,107.2401) and (67.7090,106.6558) ..(67.7090,105.9350);
\path[fill=chfourfig5c0,fill opacity=0.5,draw=chfourfig5c0,line width=0.3985pt,draw opacity=0.5] (105.7966,85.7803) .. controls (105.7966,85.0595) and (105.2123,84.4752) ..(104.4915,84.4752) .. controls (103.7707,84.4752) and (103.1864,85.0595) ..(103.1864,85.7803) .. controls (103.1864,86.5011) and (103.7707,87.0854) ..(104.4915,87.0854) .. controls (105.2123,87.0854) and (105.7966,86.5011) ..(105.7966,85.7803);
\path[fill=chfourfig5c0,fill opacity=0.5,draw=chfourfig5c0,line width=0.3985pt,draw opacity=0.5] (91.7591,82.5026) .. controls (91.7591,81.7818) and (91.1748,81.1974) ..(90.4540,81.1974) .. controls (89.7332,81.1974) and (89.1489,81.7818) ..(89.1489,82.5026) .. controls (89.1489,83.2234) and (89.7332,83.8077) ..(90.4540,83.8077) .. controls (91.1748,83.8077) and (91.7591,83.2234) ..(91.7591,82.5026);
\path[fill=chfourfig5c0,fill opacity=0.5,draw=chfourfig5c0,line width=0.3985pt,draw opacity=0.5] (155.6603,34.7311) .. controls (155.6603,34.0103) and (155.0760,33.4260) ..(154.3552,33.4260) .. controls (153.6344,33.4260) and (153.0501,34.0103) ..(153.0501,34.7311) .. controls (153.0501,35.4519) and (153.6344,36.0362) ..(154.3552,36.0362) .. controls (155.0760,36.0362) and (155.6603,35.4519) ..(155.6603,34.7311);
\path[fill=chfourfig5c0,fill opacity=0.5,draw=chfourfig5c0,line width=0.3985pt,draw opacity=0.5] (157.7525,57.3466) .. controls (157.7525,56.6258) and (157.1682,56.0415) ..(156.4474,56.0415) .. controls (155.7266,56.0415) and (155.1422,56.6258) ..(155.1422,57.3466) .. controls (155.1422,58.0674) and (155.7266,58.6517) ..(156.4474,58.6517) .. controls (157.1682,58.6517) and (157.7525,58.0674) ..(157.7525,57.3466);
\path[fill=chfourfig5c0,fill opacity=0.5,draw=chfourfig5c0,line width=0.3985pt,draw opacity=0.5] (101.9311,67.7377) .. controls (101.9311,67.0169) and (101.3468,66.4326) ..(100.6260,66.4326) .. controls (99.9052,66.4326) and (99.3208,67.0169) ..(99.3208,67.7377) .. controls (99.3208,68.4585) and (99.9052,69.0429) ..(100.6260,69.0429) .. controls (101.3468,69.0429) and (101.9311,68.4585) ..(101.9311,67.7377);
\path[fill=chfourfig5c0,fill opacity=0.5,draw=chfourfig5c0,line width=0.3985pt,draw opacity=0.5] (96.8600,57.8945) .. controls (96.8600,57.1737) and (96.2757,56.5894) ..(95.5549,56.5894) .. controls (94.8341,56.5894) and (94.2498,57.1737) ..(94.2498,57.8945) .. controls (94.2498,58.6153) and (94.8341,59.1997) ..(95.5549,59.1997) .. controls (96.2757,59.1997) and (96.8600,58.6153) ..(96.8600,57.8945);
\path[fill=chfourfig5c0,fill opacity=0.5,draw=chfourfig5c0,line width=0.3985pt,draw opacity=0.5] (62.3391,93.0531) .. controls (62.3391,92.3323) and (61.7547,91.7480) ..(61.0339,91.7480) .. controls (60.3131,91.7480) and (59.7288,92.3323) ..(59.7288,93.0531) .. controls (59.7288,93.7739) and (60.3131,94.3583) ..(61.0339,94.3583) .. controls (61.7547,94.3583) and (62.3391,93.7739) ..(62.3391,93.0531);
\path[fill=chfourfig5c0,fill opacity=0.5,draw=chfourfig5c0,line width=0.3985pt,draw opacity=0.5] (154.6342,35.9167) .. controls (154.6342,35.1959) and (154.0498,34.6115) ..(153.3290,34.6115) .. controls (152.6082,34.6115) and (152.0239,35.1959) ..(152.0239,35.9167) .. controls (152.0239,36.6375) and (152.6082,37.2218) ..(153.3290,37.2218) .. controls (154.0498,37.2218) and (154.6342,36.6375) ..(154.6342,35.9167);
\path[fill=chfourfig5c0,fill opacity=0.5,draw=chfourfig5c0,line width=0.3985pt,draw opacity=0.5] (59.5893,104.0919) .. controls (59.5893,103.3711) and (59.0050,102.7868) ..(58.2842,102.7868) .. controls (57.5634,102.7868) and (56.9791,103.3711) ..(56.9791,104.0919) .. controls (56.9791,104.8127) and (57.5634,105.3970) ..(58.2842,105.3970) .. controls (59.0050,105.3970) and (59.5893,104.8127) ..(59.5893,104.0919);
\path[fill=chfourfig5c0,fill opacity=0.5,draw=chfourfig5c0,line width=0.3985pt,draw opacity=0.5] (119.4257,42.7910) .. controls (119.4257,42.0702) and (118.8414,41.4858) ..(118.1206,41.4858) .. controls (117.3998,41.4858) and (116.8155,42.0702) ..(116.8155,42.7910) .. controls (116.8155,43.5118) and (117.3998,44.0961) ..(118.1206,44.0961) .. controls (118.8414,44.0961) and (119.4257,43.5118) ..(119.4257,42.7910);
\path[fill=chfourfig5c0,fill opacity=0.5,draw=chfourfig5c0,line width=0.3985pt,draw opacity=0.5] (94.0007,74.6021) .. controls (94.0007,73.8813) and (93.4164,73.2970) ..(92.6956,73.2970) .. controls (91.9748,73.2970) and (91.3905,73.8813) ..(91.3905,74.6021) .. controls (91.3905,75.3229) and (91.9748,75.9072) ..(92.6956,75.9072) .. controls (93.4164,75.9072) and (94.0007,75.3229) ..(94.0007,74.6021);
\path[fill=chfourfig5c0,fill opacity=0.5,draw=chfourfig5c0,line width=0.3985pt,draw opacity=0.5] (130.9028,43.3887) .. controls (130.9028,42.6679) and (130.3185,42.0836) ..(129.5977,42.0836) .. controls (128.8769,42.0836) and (128.2926,42.6679) ..(128.2926,43.3887) .. controls (128.2926,44.1096) and (128.8769,44.6939) ..(129.5977,44.6939) .. controls (130.3185,44.6939) and (130.9028,44.1096) ..(130.9028,43.3887);
\path[fill=chfourfig5c0,fill opacity=0.5,draw=chfourfig5c0,line width=0.3985pt,draw opacity=0.5] (59.9380,97.8153) .. controls (59.9380,97.0945) and (59.3537,96.5102) ..(58.6329,96.5102) .. controls (57.9121,96.5102) and (57.3278,97.0945) ..(57.3278,97.8153) .. controls (57.3278,98.5361) and (57.9121,99.1205) ..(58.6329,99.1205) .. controls (59.3537,99.1205) and (59.9380,98.5361) ..(59.9380,97.8153);
\path[fill=chfourfig5c0,fill opacity=0.5,draw=chfourfig5c0,line width=0.3985pt,draw opacity=0.5] (68.7750,91.1403) .. controls (68.7750,90.4195) and (68.1907,89.8352) ..(67.4699,89.8352) .. controls (66.7491,89.8352) and (66.1648,90.4195) ..(66.1648,91.1403) .. controls (66.1648,91.8611) and (66.7491,92.4454) ..(67.4699,92.4454) .. controls (68.1907,92.4454) and (68.7750,91.8611) ..(68.7750,91.1403);
\path[fill=chfourfig5c0,fill opacity=0.5,draw=chfourfig5c0,line width=0.3985pt,draw opacity=0.5] (126.4295,43.0002) .. controls (126.4295,42.2794) and (125.8452,41.6951) ..(125.1244,41.6951) .. controls (124.4036,41.6951) and (123.8193,42.2794) ..(123.8193,43.0002) .. controls (123.8193,43.7210) and (124.4036,44.3053) ..(125.1244,44.3053) .. controls (125.8452,44.3053) and (126.4295,43.7210) ..(126.4295,43.0002);
\path[fill=chfourfig5c0,fill opacity=0.5,draw=chfourfig5c0,line width=0.3985pt,draw opacity=0.5] (83.1812,101.2126) .. controls (83.1812,100.4919) and (82.5968,99.9075) ..(81.8761,99.9075) .. controls (81.1552,99.9075) and (80.5709,100.4919) ..(80.5709,101.2126) .. controls (80.5709,101.9335) and (81.1552,102.5178) ..(81.8761,102.5178) .. controls (82.5968,102.5178) and (83.1812,101.9335) ..(83.1812,101.2126);
\path[fill=chfourfig5c0,fill opacity=0.5,draw=chfourfig5c0,line width=0.3985pt,draw opacity=0.5] (158.1709,20.5939) .. controls (158.1709,19.8731) and (157.5866,19.2888) ..(156.8658,19.2888) .. controls (156.1450,19.2888) and (155.5607,19.8731) ..(155.5607,20.5939) .. controls (155.5607,21.3147) and (156.1450,21.8990) ..(156.8658,21.8990) .. controls (157.5866,21.8990) and (158.1709,21.3147) ..(158.1709,20.5939);
\path[fill=chfourfig5c0,fill opacity=0.5,draw=chfourfig5c0,line width=0.3985pt,draw opacity=0.5] (142.8282,52.2755) .. controls (142.8282,51.5547) and (142.2439,50.9704) ..(141.5231,50.9704) .. controls (140.8023,50.9704) and (140.2180,51.5547) ..(140.2180,52.2755) .. controls (140.2180,52.9963) and (140.8023,53.5807) ..(141.5231,53.5807) .. controls (142.2439,53.5807) and (142.8282,52.9963) ..(142.8282,52.2755);
\path[fill=chfourfig5c0,fill opacity=0.5,draw=chfourfig5c0,line width=0.3985pt,draw opacity=0.5] (98.2747,85.6010) .. controls (98.2747,84.8802) and (97.6904,84.2959) ..(96.9696,84.2959) .. controls (96.2488,84.2959) and (95.6645,84.8802) ..(95.6645,85.6010) .. controls (95.6645,86.3218) and (96.2488,86.9061) ..(96.9696,86.9061) .. controls (97.6904,86.9061) and (98.2747,86.3218) ..(98.2747,85.6010);
\path[fill=chfourfig5c0,fill opacity=0.5,draw=chfourfig5c0,line width=0.3985pt,draw opacity=0.5] (86.8973,68.5846) .. controls (86.8973,67.8638) and (86.3130,67.2795) ..(85.5922,67.2795) .. controls (84.8713,67.2795) and (84.2870,67.8638) ..(84.2870,68.5846) .. controls (84.2870,69.3054) and (84.8713,69.8897) ..(85.5922,69.8897) .. controls (86.3130,69.8897) and (86.8973,69.3054) ..(86.8973,68.5846);
\path[fill=chfourfig5c0,fill opacity=0.5,draw=chfourfig5c0,line width=0.3985pt,draw opacity=0.5] (106.8726,79.1750) .. controls (106.8726,78.4542) and (106.2883,77.8699) ..(105.5675,77.8699) .. controls (104.8467,77.8699) and (104.2624,78.4542) ..(104.2624,79.1750) .. controls (104.2624,79.8958) and (104.8467,80.4801) ..(105.5675,80.4801) .. controls (106.2883,80.4801) and (106.8726,79.8958) ..(106.8726,79.1750);
\path[fill=chfourfig5c0,fill opacity=0.5,draw=chfourfig5c0,line width=0.3985pt,draw opacity=0.5] (70.4189,72.1014) .. controls (70.4189,71.3806) and (69.8345,70.7963) ..(69.1137,70.7963) .. controls (68.3929,70.7963) and (67.8086,71.3806) ..(67.8086,72.1014) .. controls (67.8086,72.8222) and (68.3929,73.4066) ..(69.1137,73.4066) .. controls (69.8345,73.4066) and (70.4189,72.8222) ..(70.4189,72.1014);
\path[fill=chfourfig5c0,fill opacity=0.5,draw=chfourfig5c0,line width=0.3985pt,draw opacity=0.5] (77.8212,65.2072) .. controls (77.8212,64.4864) and (77.2369,63.9021) ..(76.5161,63.9021) .. controls (75.7953,63.9021) and (75.2110,64.4864) ..(75.2110,65.2072) .. controls (75.2110,65.9280) and (75.7953,66.5123) ..(76.5161,66.5123) .. controls (77.2369,66.5123) and (77.8212,65.9280) ..(77.8212,65.2072);
\path[fill=chfourfig5c0,fill opacity=0.5,draw=chfourfig5c0,line width=0.3985pt,draw opacity=0.5] (75.1113,71.0454) .. controls (75.1113,70.3246) and (74.5270,69.7403) ..(73.8062,69.7403) .. controls (73.0854,69.7403) and (72.5011,70.3246) ..(72.5011,71.0454) .. controls (72.5011,71.7662) and (73.0854,72.3505) ..(73.8062,72.3505) .. controls (74.5270,72.3505) and (75.1113,71.7662) ..(75.1113,71.0454);
\path[fill=chfourfig5c0,fill opacity=0.5,draw=chfourfig5c0,line width=0.3985pt,draw opacity=0.5] (95.0369,90.0444) .. controls (95.0369,89.3236) and (94.4525,88.7393) ..(93.7317,88.7393) .. controls (93.0109,88.7393) and (92.4266,89.3236) ..(92.4266,90.0444) .. controls (92.4266,90.7652) and (93.0109,91.3495) ..(93.7317,91.3495) .. controls (94.4525,91.3495) and (95.0369,90.7652) ..(95.0369,90.0444);
\path[fill=chfourfig5c0,fill opacity=0.5,draw=chfourfig5c0,line width=0.3985pt,draw opacity=0.5] (68.5259,84.4353) .. controls (68.5259,83.7146) and (67.9416,83.1302) ..(67.2208,83.1302) .. controls (66.5000,83.1302) and (65.9157,83.7146) ..(65.9157,84.4353) .. controls (65.9157,85.1562) and (66.5000,85.7405) ..(67.2208,85.7405) .. controls (67.9416,85.7405) and (68.5259,85.1562) ..(68.5259,84.4353);
\path[fill=chfourfig5c0,fill opacity=0.5,draw=chfourfig5c0,line width=0.3985pt,draw opacity=0.5] (99.2810,92.9137) .. controls (99.2810,92.1929) and (98.6967,91.6085) ..(97.9759,91.6085) .. controls (97.2551,91.6085) and (96.6707,92.1929) ..(96.6707,92.9137) .. controls (96.6707,93.6345) and (97.2551,94.2188) ..(97.9759,94.2188) .. controls (98.6967,94.2188) and (99.2810,93.6345) ..(99.2810,92.9137);
\path[fill=chfourfig5c0,fill opacity=0.5,draw=chfourfig5c0,line width=0.3985pt,draw opacity=0.5] (118.3099,66.7116) .. controls (118.3099,65.9908) and (117.7256,65.4065) ..(117.0048,65.4065) .. controls (116.2840,65.4065) and (115.6996,65.9908) ..(115.6996,66.7116) .. controls (115.6996,67.4324) and (116.2840,68.0167) ..(117.0048,68.0167) .. controls (117.7256,68.0167) and (118.3099,67.4324) ..(118.3099,66.7116);
\path[fill=chfourfig5c0,fill opacity=0.5,draw=chfourfig5c0,line width=0.3985pt,draw opacity=0.5] (60.3864,101.0831) .. controls (60.3864,100.3623) and (59.8020,99.7780) ..(59.0812,99.7780) .. controls (58.3604,99.7780) and (57.7761,100.3623) ..(57.7761,101.0831) .. controls (57.7761,101.8039) and (58.3604,102.3882) ..(59.0812,102.3882) .. controls (59.8020,102.3882) and (60.3864,101.8039) ..(60.3864,101.0831);
\path[fill=chfourfig5c0,fill opacity=0.5,draw=chfourfig5c0,line width=0.3985pt,draw opacity=0.5] (151.5357,36.5842) .. controls (151.5357,35.8634) and (150.9514,35.2791) ..(150.2306,35.2791) .. controls (149.5098,35.2791) and (148.9255,35.8634) ..(148.9255,36.5842) .. controls (148.9255,37.3050) and (149.5098,37.8893) ..(150.2306,37.8893) .. controls (150.9514,37.8893) and (151.5357,37.3050) ..(151.5357,36.5842);
\path[fill=chfourfig5c0,fill opacity=0.5,draw=chfourfig5c0,line width=0.3985pt,draw opacity=0.5] (114.3646,81.6956) .. controls (114.3646,80.9748) and (113.7803,80.3905) ..(113.0595,80.3905) .. controls (112.3387,80.3905) and (111.7544,80.9748) ..(111.7544,81.6956) .. controls (111.7544,82.4164) and (112.3387,83.0007) ..(113.0595,83.0007) .. controls (113.7803,83.0007) and (114.3646,82.4164) ..(114.3646,81.6956);
\path[fill=chfourfig5c0,fill opacity=0.5,draw=chfourfig5c0,line width=0.3985pt,draw opacity=0.5] (76.1176,73.5660) .. controls (76.1176,72.8452) and (75.5332,72.2608) ..(74.8124,72.2608) .. controls (74.0916,72.2608) and (73.5073,72.8452) ..(73.5073,73.5660) .. controls (73.5073,74.2868) and (74.0916,74.8711) ..(74.8124,74.8711) .. controls (75.5332,74.8711) and (76.1176,74.2868) ..(76.1176,73.5660);
\path[fill=chfourfig5c0,fill opacity=0.5,draw=chfourfig5c0,line width=0.3985pt,draw opacity=0.5] (100.0581,65.1773) .. controls (100.0581,64.4565) and (99.4738,63.8722) ..(98.7530,63.8722) .. controls (98.0322,63.8722) and (97.4478,64.4565) ..(97.4478,65.1773) .. controls (97.4478,65.8981) and (98.0322,66.4824) ..(98.7530,66.4824) .. controls (99.4738,66.4824) and (100.0581,65.8981) ..(100.0581,65.1773);
\path[fill=chfourfig5c0,fill opacity=0.5,draw=chfourfig5c0,line width=0.3985pt,draw opacity=0.5] (86.9770,74.0741) .. controls (86.9770,73.3533) and (86.3927,72.7689) ..(85.6718,72.7689) .. controls (84.9510,72.7689) and (84.3667,73.3533) ..(84.3667,74.0741) .. controls (84.3667,74.7949) and (84.9510,75.3792) ..(85.6718,75.3792) .. controls (86.3927,75.3792) and (86.9770,74.7949) ..(86.9770,74.0741);
\path[fill=chfourfig5c0,fill opacity=0.5,draw=chfourfig5c0,line width=0.3985pt,draw opacity=0.5] (75.2906,105.8154) .. controls (75.2906,105.0946) and (74.7063,104.5103) ..(73.9855,104.5103) .. controls (73.2647,104.5103) and (72.6804,105.0946) ..(72.6804,105.8154) .. controls (72.6804,106.5362) and (73.2647,107.1206) ..(73.9855,107.1206) .. controls (74.7063,107.1206) and (75.2906,106.5362) ..(75.2906,105.8154);
\path[fill=chfourfig5c0,fill opacity=0.5,draw=chfourfig5c0,line width=0.3985pt,draw opacity=0.5] (105.8365,53.4511) .. controls (105.8365,52.7303) and (105.2522,52.1460) ..(104.5314,52.1460) .. controls (103.8106,52.1460) and (103.2263,52.7303) ..(103.2263,53.4511) .. controls (103.2263,54.1719) and (103.8106,54.7563) ..(104.5314,54.7563) .. controls (105.2522,54.7563) and (105.8365,54.1719) ..(105.8365,53.4511);
\path[fill=chfourfig5c0,fill opacity=0.5,draw=chfourfig5c0,line width=0.3985pt,draw opacity=0.5] (117.9512,54.9655) .. controls (117.9512,54.2447) and (117.3669,53.6604) ..(116.6461,53.6604) .. controls (115.9253,53.6604) and (115.3410,54.2447) ..(115.3410,54.9655) .. controls (115.3410,55.6863) and (115.9253,56.2706) ..(116.6461,56.2706) .. controls (117.3669,56.2706) and (117.9512,55.6863) ..(117.9512,54.9655);
\path[fill=chfourfig5c0,fill opacity=0.5,draw=chfourfig5c0,line width=0.3985pt,draw opacity=0.5] (59.0613,115.2701) .. controls (59.0613,114.5493) and (58.4770,113.9650) ..(57.7562,113.9650) .. controls (57.0354,113.9650) and (56.4510,114.5493) ..(56.4510,115.2701) .. controls (56.4510,115.9909) and (57.0354,116.5752) ..(57.7562,116.5752) .. controls (58.4770,116.5752) and (59.0613,115.9909) ..(59.0613,115.2701);
\path[fill=chfourfig5c0,fill opacity=0.5,draw=chfourfig5c0,line width=0.3985pt,draw opacity=0.5] (107.0320,54.5570) .. controls (107.0320,53.8362) and (106.4477,53.2519) ..(105.7269,53.2519) .. controls (105.0061,53.2519) and (104.4218,53.8362) ..(104.4218,54.5570) .. controls (104.4218,55.2778) and (105.0061,55.8621) ..(105.7269,55.8621) .. controls (106.4477,55.8621) and (107.0320,55.2778) ..(107.0320,54.5570);
\path[fill=chfourfig5c0,fill opacity=0.5,draw=chfourfig5c0,line width=0.3985pt,draw opacity=0.5] (75.1412,79.2049) .. controls (75.1412,78.4841) and (74.5569,77.8998) ..(73.8361,77.8998) .. controls (73.1153,77.8998) and (72.5310,78.4841) ..(72.5310,79.2049) .. controls (72.5310,79.9257) and (73.1153,80.5100) ..(73.8361,80.5100) .. controls (74.5569,80.5100) and (75.1412,79.9257) ..(75.1412,79.2049);
\path[fill=chfourfig5c0,fill opacity=0.5,draw=chfourfig5c0,line width=0.3985pt,draw opacity=0.5] (94.0605,83.0505) .. controls (94.0605,82.3297) and (93.4762,81.7454) ..(92.7554,81.7454) .. controls (92.0346,81.7454) and (91.4503,82.3297) ..(91.4503,83.0505) .. controls (91.4503,83.7713) and (92.0346,84.3557) ..(92.7554,84.3557) .. controls (93.4762,84.3557) and (94.0605,83.7713) ..(94.0605,83.0505);
\path[fill=chfourfig5c0,fill opacity=0.5,draw=chfourfig5c0,line width=0.3985pt,draw opacity=0.5] (145.5979,54.8658) .. controls (145.5979,54.1451) and (145.0136,53.5607) ..(144.2928,53.5607) .. controls (143.5720,53.5607) and (142.9877,54.1451) ..(142.9877,54.8658) .. controls (142.9877,55.5867) and (143.5720,56.1710) ..(144.2928,56.1710) .. controls (145.0136,56.1710) and (145.5979,55.5867) ..(145.5979,54.8658);
\path[fill=chfourfig5c0,fill opacity=0.5,draw=chfourfig5c0,line width=0.3985pt,draw opacity=0.5] (153.2095,44.1459) .. controls (153.2095,43.4251) and (152.6252,42.8408) ..(151.9043,42.8408) .. controls (151.1835,42.8408) and (150.5992,43.4251) ..(150.5992,44.1459) .. controls (150.5992,44.8667) and (151.1835,45.4510) ..(151.9043,45.4510) .. controls (152.6252,45.4510) and (153.2095,44.8667) ..(153.2095,44.1459);
\path[fill=chfourfig5c0,fill opacity=0.5,draw=chfourfig5c0,line width=0.3985pt,draw opacity=0.5] (148.0089,64.1312) .. controls (148.0089,63.4104) and (147.4246,62.8261) ..(146.7038,62.8261) .. controls (145.9830,62.8261) and (145.3987,63.4104) ..(145.3987,64.1312) .. controls (145.3987,64.8520) and (145.9830,65.4363) ..(146.7038,65.4363) .. controls (147.4246,65.4363) and (148.0089,64.8520) ..(148.0089,64.1312);
\path[fill=chfourfig5c0,fill opacity=0.5,draw=chfourfig5c0,line width=0.3985pt,draw opacity=0.5] (125.3137,47.0251) .. controls (125.3137,46.3044) and (124.7294,45.7200) ..(124.0086,45.7200) .. controls (123.2878,45.7200) and (122.7035,46.3044) ..(122.7035,47.0251) .. controls (122.7035,47.7459) and (123.2878,48.3303) ..(124.0086,48.3303) .. controls (124.7294,48.3303) and (125.3137,47.7459) ..(125.3137,47.0251);
\path[fill=chfourfig5c0,fill opacity=0.5,draw=chfourfig5c0,line width=0.3985pt,draw opacity=0.5] (108.4866,87.6732) .. controls (108.4866,86.9525) and (107.9023,86.3681) ..(107.1815,86.3681) .. controls (106.4607,86.3681) and (105.8763,86.9525) ..(105.8763,87.6732) .. controls (105.8763,88.3941) and (106.4607,88.9784) ..(107.1815,88.9784) .. controls (107.9023,88.9784) and (108.4866,88.3941) ..(108.4866,87.6732);
\path[fill=chfourfig5c0,fill opacity=0.5,draw=chfourfig5c0,line width=0.3985pt,draw opacity=0.5] (75.8187,94.4280) .. controls (75.8187,93.7072) and (75.2344,93.1229) ..(74.5135,93.1229) .. controls (73.7928,93.1229) and (73.2084,93.7072) ..(73.2084,94.4280) .. controls (73.2084,95.1488) and (73.7928,95.7331) ..(74.5135,95.7331) .. controls (75.2344,95.7331) and (75.8187,95.1488) ..(75.8187,94.4280);
\path[fill=chfourfig5c0,fill opacity=0.5,draw=chfourfig5c0,line width=0.3985pt,draw opacity=0.5] (145.5082,49.5158) .. controls (145.5082,48.7950) and (144.9239,48.2107) ..(144.2031,48.2107) .. controls (143.4823,48.2107) and (142.8980,48.7950) ..(142.8980,49.5158) .. controls (142.8980,50.2366) and (143.4823,50.8210) ..(144.2031,50.8210) .. controls (144.9239,50.8210) and (145.5082,50.2366) ..(145.5082,49.5158);
\path[fill=chfourfig5c0,fill opacity=0.5,draw=chfourfig5c0,line width=0.3985pt,draw opacity=0.5] (141.5829,37.1421) .. controls (141.5829,36.4213) and (140.9986,35.8370) ..(140.2778,35.8370) .. controls (139.5570,35.8370) and (138.9727,36.4213) ..(138.9727,37.1421) .. controls (138.9727,37.8629) and (139.5570,38.4472) ..(140.2778,38.4472) .. controls (140.9986,38.4472) and (141.5829,37.8629) ..(141.5829,37.1421);
\path[fill=chfourfig5c0,fill opacity=0.5,draw=chfourfig5c0,line width=0.3985pt,draw opacity=0.5] (117.4132,49.5756) .. controls (117.4132,48.8548) and (116.8289,48.2705) ..(116.1081,48.2705) .. controls (115.3873,48.2705) and (114.8030,48.8548) ..(114.8030,49.5756) .. controls (114.8030,50.2964) and (115.3873,50.8807) ..(116.1081,50.8807) .. controls (116.8289,50.8807) and (117.4132,50.2964) ..(117.4132,49.5756);
\path[fill=chfourfig5c0,fill opacity=0.5,draw=chfourfig5c0,line width=0.3985pt,draw opacity=0.5] (105.5276,75.5685) .. controls (105.5276,74.8477) and (104.9433,74.2634) ..(104.2225,74.2634) .. controls (103.5017,74.2634) and (102.9174,74.8477) ..(102.9174,75.5685) .. controls (102.9174,76.2893) and (103.5017,76.8736) ..(104.2225,76.8736) .. controls (104.9433,76.8736) and (105.5276,76.2893) ..(105.5276,75.5685);
\path[fill=chfourfig5c0,fill opacity=0.5,draw=chfourfig5c0,line width=0.3985pt,draw opacity=0.5] (115.7395,70.3380) .. controls (115.7395,69.6172) and (115.1552,69.0329) ..(114.4344,69.0329) .. controls (113.7136,69.0329) and (113.1292,69.6172) ..(113.1292,70.3380) .. controls (113.1292,71.0588) and (113.7136,71.6431) ..(114.4344,71.6431) .. controls (115.1552,71.6431) and (115.7395,71.0588) ..(115.7395,70.3380);
\path[fill=chfourfig5c0,fill opacity=0.5,draw=chfourfig5c0,line width=0.3985pt,draw opacity=0.5] (125.9713,70.5373) .. controls (125.9713,69.8165) and (125.3869,69.2322) ..(124.6661,69.2322) .. controls (123.9453,69.2322) and (123.3610,69.8165) ..(123.3610,70.5373) .. controls (123.3610,71.2581) and (123.9453,71.8424) ..(124.6661,71.8424) .. controls (125.3869,71.8424) and (125.9713,71.2581) ..(125.9713,70.5373);
\path[fill=chfourfig5c0,fill opacity=0.5,draw=chfourfig5c0,line width=0.3985pt,draw opacity=0.5] (129.1195,42.5120) .. controls (129.1195,41.7912) and (128.5352,41.2069) ..(127.8144,41.2069) .. controls (127.0936,41.2069) and (126.5093,41.7912) ..(126.5093,42.5120) .. controls (126.5093,43.2328) and (127.0936,43.8171) ..(127.8144,43.8171) .. controls (128.5352,43.8171) and (129.1195,43.2328) ..(129.1195,42.5120);
\path[fill=chfourfig5c0,fill opacity=0.5,draw=chfourfig5c0,line width=0.3985pt,draw opacity=0.5] (142.7784,39.4833) .. controls (142.7784,38.7625) and (142.1941,38.1782) ..(141.4733,38.1782) .. controls (140.7525,38.1782) and (140.1682,38.7625) ..(140.1682,39.4833) .. controls (140.1682,40.2041) and (140.7525,40.7885) ..(141.4733,40.7885) .. controls (142.1941,40.7885) and (142.7784,40.2041) ..(142.7784,39.4833);
\path[fill=chfourfig5c0,fill opacity=0.5,draw=chfourfig5c0,line width=0.3985pt,draw opacity=0.5] (62.5782,86.9061) .. controls (62.5782,86.1853) and (61.9938,85.6010) ..(61.2730,85.6010) .. controls (60.5522,85.6010) and (59.9679,86.1853) ..(59.9679,86.9061) .. controls (59.9679,87.6269) and (60.5522,88.2112) ..(61.2730,88.2112) .. controls (61.9938,88.2112) and (62.5782,87.6269) ..(62.5782,86.9061);
\path[fill=chfourfig5c0,fill opacity=0.5,draw=chfourfig5c0,line width=0.3985pt,draw opacity=0.5] (140.4272,36.0562) .. controls (140.4272,35.3354) and (139.8429,34.7510) ..(139.1221,34.7510) .. controls (138.4013,34.7510) and (137.8170,35.3354) ..(137.8170,36.0562) .. controls (137.8170,36.7769) and (138.4013,37.3613) ..(139.1221,37.3613) .. controls (139.8429,37.3613) and (140.4272,36.7769) ..(140.4272,36.0562);
\path[fill=chfourfig5c0,fill opacity=0.5,draw=chfourfig5c0,line width=0.3985pt,draw opacity=0.5] (158.0314,22.8754) .. controls (158.0314,22.1546) and (157.4471,21.5703) ..(156.7263,21.5703) .. controls (156.0055,21.5703) and (155.4212,22.1546) ..(155.4212,22.8754) .. controls (155.4212,23.5962) and (156.0055,24.1805) ..(156.7263,24.1805) .. controls (157.4471,24.1805) and (158.0314,23.5962) ..(158.0314,22.8754);
\path[fill=chfourfig5c0,fill opacity=0.5,draw=chfourfig5c0,line width=0.3985pt,draw opacity=0.5] (93.1140,79.4440) .. controls (93.1140,78.7232) and (92.5297,78.1389) ..(91.8089,78.1389) .. controls (91.0881,78.1389) and (90.5038,78.7232) ..(90.5038,79.4440) .. controls (90.5038,80.1648) and (91.0881,80.7491) ..(91.8089,80.7491) .. controls (92.5297,80.7491) and (93.1140,80.1648) ..(93.1140,79.4440);
\path[fill=chfourfig5c0,fill opacity=0.5,draw=chfourfig5c0,line width=0.3985pt,draw opacity=0.5] (150.6988,54.8658) .. controls (150.6988,54.1451) and (150.1145,53.5607) ..(149.3937,53.5607) .. controls (148.6729,53.5607) and (148.0886,54.1451) ..(148.0886,54.8658) .. controls (148.0886,55.5867) and (148.6729,56.1710) ..(149.3937,56.1710) .. controls (150.1145,56.1710) and (150.6988,55.5867) ..(150.6988,54.8658);
\path[fill=chfourfig5c0,fill opacity=0.5,draw=chfourfig5c0,line width=0.3985pt,draw opacity=0.5] (141.1047,69.1126) .. controls (141.1047,68.3918) and (140.5204,67.8075) ..(139.7996,67.8075) .. controls (139.0788,67.8075) and (138.4944,68.3918) ..(138.4944,69.1126) .. controls (138.4944,69.8334) and (139.0788,70.4177) ..(139.7996,70.4177) .. controls (140.5204,70.4177) and (141.1047,69.8334) ..(141.1047,69.1126);
\path[fill=chfourfig5c0,fill opacity=0.5,draw=chfourfig5c0,line width=0.3985pt,draw opacity=0.5] (149.7125,38.8358) .. controls (149.7125,38.1150) and (149.1282,37.5306) ..(148.4074,37.5306) .. controls (147.6866,37.5306) and (147.1023,38.1150) ..(147.1023,38.8358) .. controls (147.1023,39.5565) and (147.6866,40.1409) ..(148.4074,40.1409) .. controls (149.1282,40.1409) and (149.7125,39.5565) ..(149.7125,38.8358);
\path[fill=chfourfig5c0,fill opacity=0.5,draw=chfourfig5c0,line width=0.3985pt,draw opacity=0.5] (133.8319,52.7737) .. controls (133.8319,52.0529) and (133.2476,51.4686) ..(132.5268,51.4686) .. controls (131.8060,51.4686) and (131.2216,52.0529) ..(131.2216,52.7737) .. controls (131.2216,53.4945) and (131.8060,54.0788) ..(132.5268,54.0788) .. controls (133.2476,54.0788) and (133.8319,53.4945) ..(133.8319,52.7737);
\path[fill=chfourfig5c0,fill opacity=0.5,draw=chfourfig5c0,line width=0.3985pt,draw opacity=0.5] (105.6073,68.0167) .. controls (105.6073,67.2959) and (105.0230,66.7116) ..(104.3022,66.7116) .. controls (103.5814,66.7116) and (102.9971,67.2959) ..(102.9971,68.0167) .. controls (102.9971,68.7375) and (103.5814,69.3218) ..(104.3022,69.3218) .. controls (105.0230,69.3218) and (105.6073,68.7375) ..(105.6073,68.0167);
\path[fill=chfourfig5c0,fill opacity=0.5,draw=chfourfig5c0,line width=0.3985pt,draw opacity=0.5] (124.0883,70.6070) .. controls (124.0883,69.8862) and (123.5040,69.3019) ..(122.7832,69.3019) .. controls (122.0624,69.3019) and (121.4781,69.8862) ..(121.4781,70.6070) .. controls (121.4781,71.3278) and (122.0624,71.9121) ..(122.7832,71.9121) .. controls (123.5040,71.9121) and (124.0883,71.3278) ..(124.0883,70.6070);
\path[fill=chfourfig5c0,fill opacity=0.5,draw=chfourfig5c0,line width=0.3985pt,draw opacity=0.5] (98.8426,82.0642) .. controls (98.8426,81.3434) and (98.2583,80.7591) ..(97.5375,80.7591) .. controls (96.8167,80.7591) and (96.2324,81.3434) ..(96.2324,82.0642) .. controls (96.2324,82.7850) and (96.8167,83.3693) ..(97.5375,83.3693) .. controls (98.2583,83.3693) and (98.8426,82.7850) ..(98.8426,82.0642);
\path[fill=chfourfig5c0,fill opacity=0.5,draw=chfourfig5c0,line width=0.3985pt,draw opacity=0.5] (133.5230,73.7453) .. controls (133.5230,73.0245) and (132.9387,72.4402) ..(132.2179,72.4402) .. controls (131.4971,72.4402) and (130.9128,73.0245) ..(130.9128,73.7453) .. controls (130.9128,74.4661) and (131.4971,75.0504) ..(132.2179,75.0504) .. controls (132.9387,75.0504) and (133.5230,74.4661) ..(133.5230,73.7453);
\path[fill=chfourfig5c0,fill opacity=0.5,draw=chfourfig5c0,line width=0.3985pt,draw opacity=0.5] (112.9599,62.5571) .. controls (112.9599,61.8363) and (112.3755,61.2520) ..(111.6548,61.2520) .. controls (110.9340,61.2520) and (110.3496,61.8363) ..(110.3496,62.5571) .. controls (110.3496,63.2779) and (110.9340,63.8622) ..(111.6548,63.8622) .. controls (112.3755,63.8622) and (112.9599,63.2779) ..(112.9599,62.5571);
\path[fill=chfourfig5c0,fill opacity=0.5,draw=chfourfig5c0,line width=0.3985pt,draw opacity=0.5] (135.4558,71.9022) .. controls (135.4558,71.1814) and (134.8715,70.5971) ..(134.1507,70.5971) .. controls (133.4299,70.5971) and (132.8456,71.1814) ..(132.8456,71.9022) .. controls (132.8456,72.6230) and (133.4299,73.2073) ..(134.1507,73.2073) .. controls (134.8715,73.2073) and (135.4558,72.6230) ..(135.4558,71.9022);
\path[fill=chfourfig5c0,fill opacity=0.5,draw=chfourfig5c0,line width=0.3985pt,draw opacity=0.5] (115.7893,47.7425) .. controls (115.7893,47.0217) and (115.2050,46.4373) ..(114.4842,46.4373) .. controls (113.7634,46.4373) and (113.1791,47.0217) ..(113.1791,47.7425) .. controls (113.1791,48.4633) and (113.7634,49.0476) ..(114.4842,49.0476) .. controls (115.2050,49.0476) and (115.7893,48.4633) ..(115.7893,47.7425);
\path[fill=chfourfig5c0,fill opacity=0.5,draw=chfourfig5c0,line width=0.3985pt,draw opacity=0.5] (98.2349,62.0092) .. controls (98.2349,61.2883) and (97.6506,60.7040) ..(96.9298,60.7040) .. controls (96.2090,60.7040) and (95.6247,61.2883) ..(95.6247,62.0092) .. controls (95.6247,62.7299) and (96.2090,63.3143) ..(96.9298,63.3143) .. controls (97.6506,63.3143) and (98.2349,62.7299) ..(98.2349,62.0092);
\path[fill=chfourfig5c0,fill opacity=0.5,draw=chfourfig5c0,line width=0.3985pt,draw opacity=0.5] (83.8188,89.6857) .. controls (83.8188,88.9649) and (83.2345,88.3806) ..(82.5137,88.3806) .. controls (81.7929,88.3806) and (81.2085,88.9649) ..(81.2085,89.6857) .. controls (81.2085,90.4065) and (81.7929,90.9909) ..(82.5137,90.9909) .. controls (83.2345,90.9909) and (83.8188,90.4065) ..(83.8188,89.6857);
\path[fill=chfourfig5c0,fill opacity=0.5,draw=chfourfig5c0,line width=0.3985pt,draw opacity=0.5] (144.8009,32.8382) .. controls (144.8009,32.1173) and (144.2166,31.5331) ..(143.4958,31.5331) .. controls (142.7750,31.5331) and (142.1906,32.1173) ..(142.1906,32.8382) .. controls (142.1906,33.5590) and (142.7750,34.1433) ..(143.4958,34.1433) .. controls (144.2166,34.1433) and (144.8009,33.5590) ..(144.8009,32.8382);
\path[fill=chfourfig5c0,fill opacity=0.5,draw=chfourfig5c0,line width=0.3985pt,draw opacity=0.5] (96.3420,73.2472) .. controls (96.3420,72.5264) and (95.7577,71.9420) ..(95.0369,71.9420) .. controls (94.3160,71.9420) and (93.7317,72.5264) ..(93.7317,73.2472) .. controls (93.7317,73.9680) and (94.3160,74.5523) ..(95.0369,74.5523) .. controls (95.7577,74.5523) and (96.3420,73.9680) ..(96.3420,73.2472);
\path[fill=chfourfig5c0,fill opacity=0.5,draw=chfourfig5c0,line width=0.3985pt,draw opacity=0.5] (58.6429,78.3680) .. controls (58.6429,77.6472) and (58.0585,77.0629) ..(57.3378,77.0629) .. controls (56.6169,77.0629) and (56.0326,77.6472) ..(56.0326,78.3680) .. controls (56.0326,79.0888) and (56.6169,79.6731) ..(57.3378,79.6731) .. controls (58.0585,79.6731) and (58.6429,79.0888) ..(58.6429,78.3680);
\path[fill=chfourfig5c0,fill opacity=0.5,draw=chfourfig5c0,line width=0.3985pt,draw opacity=0.5] (67.5894,102.2787) .. controls (67.5894,101.5579) and (67.0051,100.9736) ..(66.2843,100.9736) .. controls (65.5635,100.9736) and (64.9792,101.5579) ..(64.9792,102.2787) .. controls (64.9792,102.9995) and (65.5635,103.5838) ..(66.2843,103.5838) .. controls (67.0051,103.5838) and (67.5894,102.9995) ..(67.5894,102.2787);
\path[fill=chfourfig5c0,fill opacity=0.5,draw=chfourfig5c0,line width=0.3985pt,draw opacity=0.5] (144.2529,63.2744) .. controls (144.2529,62.5536) and (143.6686,61.9693) ..(142.9478,61.9693) .. controls (142.2270,61.9693) and (141.6427,62.5536) ..(141.6427,63.2744) .. controls (141.6427,63.9952) and (142.2270,64.5795) ..(142.9478,64.5795) .. controls (143.6686,64.5795) and (144.2529,63.9952) ..(144.2529,63.2744);
\path[fill=chfourfig5c0,fill opacity=0.5,draw=chfourfig5c0,line width=0.3985pt,draw opacity=0.5] (110.6983,69.0728) .. controls (110.6983,68.3520) and (110.1140,67.7676) ..(109.3932,67.7676) .. controls (108.6724,67.7676) and (108.0881,68.3520) ..(108.0881,69.0728) .. controls (108.0881,69.7936) and (108.6724,70.3779) ..(109.3932,70.3779) .. controls (110.1140,70.3779) and (110.6983,69.7936) ..(110.6983,69.0728);
\path[fill=chfourfig5c0,fill opacity=0.5,draw=chfourfig5c0,line width=0.3985pt,draw opacity=0.5] (112.0832,58.9207) .. controls (112.0832,58.1999) and (111.4988,57.6156) ..(110.7780,57.6156) .. controls (110.0572,57.6156) and (109.4729,58.1999) ..(109.4729,58.9207) .. controls (109.4729,59.6415) and (110.0572,60.2258) ..(110.7780,60.2258) .. controls (111.4988,60.2258) and (112.0832,59.6415) ..(112.0832,58.9207);
\path[fill=chfourfig5c0,fill opacity=0.5,draw=chfourfig5c0,line width=0.3985pt,draw opacity=0.5] (63.1261,113.1779) .. controls (63.1261,112.4571) and (62.5418,111.8728) ..(61.8210,111.8728) .. controls (61.1002,111.8728) and (60.5159,112.4571) ..(60.5159,113.1779) .. controls (60.5159,113.8987) and (61.1002,114.4830) ..(61.8210,114.4830) .. controls (62.5418,114.4830) and (63.1261,113.8987) ..(63.1261,113.1779);
\path[fill=chfourfig5c0,fill opacity=0.5,draw=chfourfig5c0,line width=0.3985pt,draw opacity=0.5] (149.2941,33.4160) .. controls (149.2941,32.6952) and (148.7098,32.1109) ..(147.9890,32.1109) .. controls (147.2682,32.1109) and (146.6839,32.6952) ..(146.6839,33.4160) .. controls (146.6839,34.1368) and (147.2682,34.7211) ..(147.9890,34.7211) .. controls (148.7098,34.7211) and (149.2941,34.1368) ..(149.2941,33.4160);
\path[fill=chfourfig5c0,fill opacity=0.5,draw=chfourfig5c0,line width=0.3985pt,draw opacity=0.5] (67.4300,83.8874) .. controls (67.4300,83.1666) and (66.8457,82.5823) ..(66.1249,82.5823) .. controls (65.4041,82.5823) and (64.8198,83.1666) ..(64.8198,83.8874) .. controls (64.8198,84.6082) and (65.4041,85.1925) ..(66.1249,85.1925) .. controls (66.8457,85.1925) and (67.4300,84.6082) ..(67.4300,83.8874);
\path[fill=chfourfig5c0,fill opacity=0.5,draw=chfourfig5c0,line width=0.3985pt,draw opacity=0.5] (71.5945,72.5298) .. controls (71.5945,71.8090) and (71.0101,71.2247) ..(70.2894,71.2247) .. controls (69.5685,71.2247) and (68.9842,71.8090) ..(68.9842,72.5298) .. controls (68.9842,73.2506) and (69.5685,73.8350) ..(70.2894,73.8350) .. controls (71.0101,73.8350) and (71.5945,73.2506) ..(71.5945,72.5298);
\path[fill=chfourfig5c0,fill opacity=0.5,draw=chfourfig5c0,line width=0.3985pt,draw opacity=0.5] (75.2607,75.2397) .. controls (75.2607,74.5189) and (74.6764,73.9346) ..(73.9556,73.9346) .. controls (73.2348,73.9346) and (72.6505,74.5189) ..(72.6505,75.2397) .. controls (72.6505,75.9605) and (73.2348,76.5448) ..(73.9556,76.5448) .. controls (74.6764,76.5448) and (75.2607,75.9605) ..(75.2607,75.2397);
\path[fill=chfourfig5c0,fill opacity=0.5,draw=chfourfig5c0,line width=0.3985pt,draw opacity=0.5] (107.0619,52.3951) .. controls (107.0619,51.6743) and (106.4776,51.0900) ..(105.7568,51.0900) .. controls (105.0360,51.0900) and (104.4517,51.6743) ..(104.4517,52.3951) .. controls (104.4517,53.1159) and (105.0360,53.7002) ..(105.7568,53.7002) .. controls (106.4776,53.7002) and (107.0619,53.1159) ..(107.0619,52.3951);
\path[fill=chfourfig5c0,fill opacity=0.5,draw=chfourfig5c0,line width=0.3985pt,draw opacity=0.5] (92.8749,92.2362) .. controls (92.8749,91.5154) and (92.2906,90.9311) ..(91.5698,90.9311) .. controls (90.8490,90.9311) and (90.2647,91.5154) ..(90.2647,92.2362) .. controls (90.2647,92.9570) and (90.8490,93.5413) ..(91.5698,93.5413) .. controls (92.2906,93.5413) and (92.8749,92.9570) ..(92.8749,92.2362);
\path[fill=chfourfig5c0,fill opacity=0.5,draw=chfourfig5c0,line width=0.3985pt,draw opacity=0.5] (105.5476,81.2373) .. controls (105.5476,80.5165) and (104.9633,79.9322) ..(104.2425,79.9322) .. controls (103.5217,79.9322) and (102.9373,80.5165) ..(102.9373,81.2373) .. controls (102.9373,81.9581) and (103.5217,82.5424) ..(104.2425,82.5424) .. controls (104.9633,82.5424) and (105.5476,81.9581) ..(105.5476,81.2373);
\path[fill=chfourfig5c0,fill opacity=0.5,draw=chfourfig5c0,line width=0.3985pt,draw opacity=0.5] (83.9184,85.6109) .. controls (83.9184,84.8902) and (83.3341,84.3058) ..(82.6133,84.3058) .. controls (81.8925,84.3058) and (81.3082,84.8902) ..(81.3082,85.6109) .. controls (81.3082,86.3318) and (81.8925,86.9161) ..(82.6133,86.9161) .. controls (83.3341,86.9161) and (83.9184,86.3318) ..(83.9184,85.6109);
\path[fill=chfourfig5c0,fill opacity=0.5,draw=chfourfig5c0,line width=0.3985pt,draw opacity=0.5] (61.8011,100.0470) .. controls (61.8011,99.3262) and (61.2167,98.7419) ..(60.4959,98.7419) .. controls (59.7751,98.7419) and (59.1908,99.3262) ..(59.1908,100.0470) .. controls (59.1908,100.7678) and (59.7751,101.3521) ..(60.4959,101.3521) .. controls (61.2167,101.3521) and (61.8011,100.7678) ..(61.8011,100.0470);
\path[fill=chfourfig5c0,fill opacity=0.5,draw=chfourfig5c0,line width=0.3985pt,draw opacity=0.5] (116.9250,75.9570) .. controls (116.9250,75.2362) and (116.3407,74.6519) ..(115.6199,74.6519) .. controls (114.8991,74.6519) and (114.3148,75.2362) ..(114.3148,75.9570) .. controls (114.3148,76.6778) and (114.8991,77.2621) ..(115.6199,77.2621) .. controls (116.3407,77.2621) and (116.9250,76.6778) ..(116.9250,75.9570);
\path[fill=chfourfig5c0,fill opacity=0.5,draw=chfourfig5c0,line width=0.3985pt,draw opacity=0.5] (122.1256,41.5456) .. controls (122.1256,40.8248) and (121.5413,40.2405) ..(120.8205,40.2405) .. controls (120.0997,40.2405) and (119.5154,40.8248) ..(119.5154,41.5456) .. controls (119.5154,42.2664) and (120.0997,42.8508) ..(120.8205,42.8508) .. controls (121.5413,42.8508) and (122.1256,42.2664) ..(122.1256,41.5456);
\path[fill=chfourfig5c0,fill opacity=0.5,draw=chfourfig5c0,line width=0.3985pt,draw opacity=0.5] (109.4629,52.0962) .. controls (109.4629,51.3754) and (108.8786,50.7911) ..(108.1578,50.7911) .. controls (107.4370,50.7911) and (106.8527,51.3754) ..(106.8527,52.0962) .. controls (106.8527,52.8170) and (107.4370,53.4013) ..(108.1578,53.4013) .. controls (108.8786,53.4013) and (109.4629,52.8170) ..(109.4629,52.0962);
\path[fill=chfourfig5c0,fill opacity=0.5,draw=chfourfig5c0,line width=0.3985pt,draw opacity=0.5] (67.5595,100.6348) .. controls (67.5595,99.9140) and (66.9752,99.3297) ..(66.2544,99.3297) .. controls (65.5336,99.3297) and (64.9493,99.9140) ..(64.9493,100.6348) .. controls (64.9493,101.3556) and (65.5336,101.9399) ..(66.2544,101.9399) .. controls (66.9752,101.9399) and (67.5595,101.3556) ..(67.5595,100.6348);
\path[fill=chfourfig5c0,fill opacity=0.5,draw=chfourfig5c0,line width=0.3985pt,draw opacity=0.5] (76.6257,102.3484) .. controls (76.6257,101.6276) and (76.0413,101.0433) ..(75.3205,101.0433) .. controls (74.5997,101.0433) and (74.0154,101.6276) ..(74.0154,102.3484) .. controls (74.0154,103.0692) and (74.5997,103.6535) ..(75.3205,103.6535) .. controls (76.0413,103.6535) and (76.6257,103.0692) ..(76.6257,102.3484);
\path[fill=chfourfig5c0,fill opacity=0.5,draw=chfourfig5c0,line width=0.3985pt,draw opacity=0.5] (115.5303,62.0390) .. controls (115.5303,61.3183) and (114.9460,60.7339) ..(114.2252,60.7339) .. controls (113.5044,60.7339) and (112.9200,61.3183) ..(112.9200,62.0390) .. controls (112.9200,62.7599) and (113.5044,63.3442) ..(114.2252,63.3442) .. controls (114.9460,63.3442) and (115.5303,62.7599) ..(115.5303,62.0390);
\path[fill=chfourfig5c0,fill opacity=0.5,draw=chfourfig5c0,line width=0.3985pt,draw opacity=0.5] (148.3974,53.7799) .. controls (148.3974,53.0591) and (147.8131,52.4748) ..(147.0923,52.4748) .. controls (146.3715,52.4748) and (145.7872,53.0591) ..(145.7872,53.7799) .. controls (145.7872,54.5007) and (146.3715,55.0850) ..(147.0923,55.0850) .. controls (147.8131,55.0850) and (148.3974,54.5007) ..(148.3974,53.7799);
\path[fill=chfourfig5c0,fill opacity=0.5,draw=chfourfig5c0,line width=0.3985pt,draw opacity=0.5] (125.8218,65.1076) .. controls (125.8218,64.3868) and (125.2375,63.8024) ..(124.5167,63.8024) .. controls (123.7959,63.8024) and (123.2115,64.3868) ..(123.2115,65.1076) .. controls (123.2115,65.8284) and (123.7959,66.4127) ..(124.5167,66.4127) .. controls (125.2375,66.4127) and (125.8218,65.8284) ..(125.8218,65.1076);
\path[fill=chfourfig5c0,fill opacity=0.5,draw=chfourfig5c0,line width=0.3985pt,draw opacity=0.5] (120.3323,47.1547) .. controls (120.3323,46.4339) and (119.7480,45.8495) ..(119.0272,45.8495) .. controls (118.3064,45.8495) and (117.7221,46.4339) ..(117.7221,47.1547) .. controls (117.7221,47.8755) and (118.3064,48.4598) ..(119.0272,48.4598) .. controls (119.7480,48.4598) and (120.3323,47.8755) ..(120.3323,47.1547);
\path[fill=chfourfig5c0,fill opacity=0.5,draw=chfourfig5c0,line width=0.3985pt,draw opacity=0.5] (91.7790,97.0283) .. controls (91.7790,96.3075) and (91.1947,95.7232) ..(90.4739,95.7232) .. controls (89.7531,95.7232) and (89.1688,96.3075) ..(89.1688,97.0283) .. controls (89.1688,97.7491) and (89.7531,98.3334) ..(90.4739,98.3334) .. controls (91.1947,98.3334) and (91.7790,97.7491) ..(91.7790,97.0283);
\path[fill=chfourfig5c0,fill opacity=0.5,draw=chfourfig5c0,line width=0.3985pt,draw opacity=0.5] (87.9733,64.1810) .. controls (87.9733,63.4602) and (87.3889,62.8759) ..(86.6681,62.8759) .. controls (85.9473,62.8759) and (85.3630,63.4602) ..(85.3630,64.1810) .. controls (85.3630,64.9018) and (85.9473,65.4862) ..(86.6681,65.4862) .. controls (87.3889,65.4862) and (87.9733,64.9018) ..(87.9733,64.1810);
\path[fill=chfourfig5c0,fill opacity=0.5,draw=chfourfig5c0,line width=0.3985pt,draw opacity=0.5] (113.3285,62.1586) .. controls (113.3285,61.4378) and (112.7442,60.8535) ..(112.0234,60.8535) .. controls (111.3026,60.8535) and (110.7182,61.4378) ..(110.7182,62.1586) .. controls (110.7182,62.8794) and (111.3026,63.4637) ..(112.0234,63.4637) .. controls (112.7442,63.4637) and (113.3285,62.8794) ..(113.3285,62.1586);
\path[fill=chfourfig5c0,fill opacity=0.5,draw=chfourfig5c0,line width=0.3985pt,draw opacity=0.5] (70.0204,95.1752) .. controls (70.0204,94.4544) and (69.4360,93.8701) ..(68.7152,93.8701) .. controls (67.9944,93.8701) and (67.4101,94.4544) ..(67.4101,95.1752) .. controls (67.4101,95.8960) and (67.9944,96.4803) ..(68.7152,96.4803) .. controls (69.4360,96.4803) and (70.0204,95.8960) ..(70.0204,95.1752);
\path[fill=chfourfig5c0,fill opacity=0.5,draw=chfourfig5c0,line width=0.3985pt,draw opacity=0.5] (140.1084,48.4897) .. controls (140.1084,47.7689) and (139.5241,47.1845) ..(138.8033,47.1845) .. controls (138.0825,47.1845) and (137.4982,47.7689) ..(137.4982,48.4897) .. controls (137.4982,49.2105) and (138.0825,49.7948) ..(138.8033,49.7948) .. controls (139.5241,49.7948) and (140.1084,49.2105) ..(140.1084,48.4897);
\path[fill=chfourfig5c0,fill opacity=0.5,draw=chfourfig5c0,line width=0.3985pt,draw opacity=0.5] (63.6442,83.4590) .. controls (63.6442,82.7382) and (63.0598,82.1539) ..(62.3391,82.1539) .. controls (61.6183,82.1539) and (61.0339,82.7382) ..(61.0339,83.4590) .. controls (61.0339,84.1798) and (61.6183,84.7641) ..(62.3391,84.7641) .. controls (63.0598,84.7641) and (63.6442,84.1798) ..(63.6442,83.4590);
\path[fill=chfourfig5c0,fill opacity=0.5,draw=chfourfig5c0,line width=0.3985pt,draw opacity=0.5] (107.9586,50.6118) .. controls (107.9586,49.8909) and (107.3743,49.3066) ..(106.6534,49.3066) .. controls (105.9326,49.3066) and (105.3483,49.8909) ..(105.3483,50.6118) .. controls (105.3483,51.3326) and (105.9326,51.9169) ..(106.6534,51.9169) .. controls (107.3743,51.9169) and (107.9586,51.3326) ..(107.9586,50.6118);
\path[fill=chfourfig5c0,fill opacity=0.5,draw=chfourfig5c0,line width=0.3985pt,draw opacity=0.5] (98.1950,73.5560) .. controls (98.1950,72.8352) and (97.6107,72.2509) ..(96.8899,72.2509) .. controls (96.1691,72.2509) and (95.5848,72.8352) ..(95.5848,73.5560) .. controls (95.5848,74.2768) and (96.1691,74.8611) ..(96.8899,74.8611) .. controls (97.6107,74.8611) and (98.1950,74.2768) ..(98.1950,73.5560);
\path[fill=chfourfig5c0,fill opacity=0.5,draw=chfourfig5c0,line width=0.3985pt,draw opacity=0.5] (134.0511,39.2343) .. controls (134.0511,38.5135) and (133.4667,37.9292) ..(132.7459,37.9292) .. controls (132.0251,37.9292) and (131.4408,38.5135) ..(131.4408,39.2343) .. controls (131.4408,39.9551) and (132.0251,40.5394) ..(132.7459,40.5394) .. controls (133.4667,40.5394) and (134.0511,39.9551) ..(134.0511,39.2343);
\path[fill=chfourfig5c0,fill opacity=0.5,draw=chfourfig5c0,line width=0.3985pt,draw opacity=0.5] (116.0981,63.0851) .. controls (116.0981,62.3643) and (115.5138,61.7800) ..(114.7930,61.7800) .. controls (114.0722,61.7800) and (113.4879,62.3643) ..(113.4879,63.0851) .. controls (113.4879,63.8059) and (114.0722,64.3902) ..(114.7930,64.3902) .. controls (115.5138,64.3902) and (116.0981,63.8059) ..(116.0981,63.0851);
\path[fill=chfourfig5c0,fill opacity=0.5,draw=chfourfig5c0,line width=0.3985pt,draw opacity=0.5] (96.8501,69.4414) .. controls (96.8501,68.7206) and (96.2658,68.1363) ..(95.5450,68.1363) .. controls (94.8242,68.1363) and (94.2399,68.7206) ..(94.2399,69.4414) .. controls (94.2399,70.1622) and (94.8242,70.7465) ..(95.5450,70.7465) .. controls (96.2658,70.7465) and (96.8501,70.1622) ..(96.8501,69.4414);
\path[fill=chfourfig5c0,fill opacity=0.5,draw=chfourfig5c0,line width=0.3985pt,draw opacity=0.5] (121.4581,62.0092) .. controls (121.4581,61.2883) and (120.8738,60.7040) ..(120.1530,60.7040) .. controls (119.4322,60.7040) and (118.8479,61.2883) ..(118.8479,62.0092) .. controls (118.8479,62.7299) and (119.4322,63.3143) ..(120.1530,63.3143) .. controls (120.8738,63.3143) and (121.4581,62.7299) ..(121.4581,62.0092);
\path[fill=chfourfig5c0,fill opacity=0.5,draw=chfourfig5c0,line width=0.3985pt,draw opacity=0.5] (135.5355,48.0115) .. controls (135.5355,47.2906) and (134.9512,46.7063) ..(134.2304,46.7063) .. controls (133.5096,46.7063) and (132.9253,47.2906) ..(132.9253,48.0115) .. controls (132.9253,48.7323) and (133.5096,49.3166) ..(134.2304,49.3166) .. controls (134.9512,49.3166) and (135.5355,48.7323) ..(135.5355,48.0115);
\path[fill=chfourfig5c0,fill opacity=0.5,draw=chfourfig5c0,line width=0.3985pt,draw opacity=0.5] (149.5033,44.6839) .. controls (149.5033,43.9631) and (148.9190,43.3788) ..(148.1982,43.3788) .. controls (147.4774,43.3788) and (146.8931,43.9631) ..(146.8931,44.6839) .. controls (146.8931,45.4047) and (147.4774,45.9890) ..(148.1982,45.9890) .. controls (148.9190,45.9890) and (149.5033,45.4047) ..(149.5033,44.6839);
\path[fill=chfourfig5c0,fill opacity=0.5,draw=chfourfig5c0,line width=0.3985pt,draw opacity=0.5] (124.4270,61.1723) .. controls (124.4270,60.4515) and (123.8427,59.8672) ..(123.1219,59.8672) .. controls (122.4011,59.8672) and (121.8168,60.4515) ..(121.8168,61.1723) .. controls (121.8168,61.8931) and (122.4011,62.4774) ..(123.1219,62.4774) .. controls (123.8427,62.4774) and (124.4270,61.8931) ..(124.4270,61.1723);
\path[fill=chfourfig5c0,fill opacity=0.5,draw=chfourfig5c0,line width=0.3985pt,draw opacity=0.5] (79.2758,79.7329) .. controls (79.2758,79.0121) and (78.6915,78.4278) ..(77.9706,78.4278) .. controls (77.2498,78.4278) and (76.6655,79.0121) ..(76.6655,79.7329) .. controls (76.6655,80.4537) and (77.2498,81.0380) ..(77.9706,81.0380) .. controls (78.6915,81.0380) and (79.2758,80.4537) ..(79.2758,79.7329);
\path[fill=chfourfig5c0,fill opacity=0.5,draw=chfourfig5c0,line width=0.3985pt,draw opacity=0.5] (82.0454,88.1116) .. controls (82.0454,87.3908) and (81.4611,86.8065) ..(80.7403,86.8065) .. controls (80.0195,86.8065) and (79.4352,87.3908) ..(79.4352,88.1116) .. controls (79.4352,88.8324) and (80.0195,89.4167) ..(80.7403,89.4167) .. controls (81.4611,89.4167) and (82.0454,88.8324) ..(82.0454,88.1116);
\path[fill=chfourfig5c0,fill opacity=0.5,draw=chfourfig5c0,line width=0.3985pt,draw opacity=0.5] (95.7841,94.5974) .. controls (95.7841,93.8766) and (95.1998,93.2923) ..(94.4789,93.2923) .. controls (93.7581,93.2923) and (93.1738,93.8766) ..(93.1738,94.5974) .. controls (93.1738,95.3182) and (93.7581,95.9025) ..(94.4789,95.9025) .. controls (95.1998,95.9025) and (95.7841,95.3182) ..(95.7841,94.5974);
\path[fill=chfourfig5c0,fill opacity=0.5,draw=chfourfig5c0,line width=0.3985pt,draw opacity=0.5] (127.8741,62.6567) .. controls (127.8741,61.9359) and (127.2898,61.3516) ..(126.5690,61.3516) .. controls (125.8482,61.3516) and (125.2639,61.9359) ..(125.2639,62.6567) .. controls (125.2639,63.3775) and (125.8482,63.9619) ..(126.5690,63.9619) .. controls (127.2898,63.9619) and (127.8741,63.3775) ..(127.8741,62.6567);
\path[fill=chfourfig5c0,fill opacity=0.5,draw=chfourfig5c0,line width=0.3985pt,draw opacity=0.5] (112.8802,65.8747) .. controls (112.8802,65.1539) and (112.2959,64.5696) ..(111.5750,64.5696) .. controls (110.8542,64.5696) and (110.2699,65.1539) ..(110.2699,65.8747) .. controls (110.2699,66.5955) and (110.8542,67.1798) ..(111.5750,67.1798) .. controls (112.2959,67.1798) and (112.8802,66.5955) ..(112.8802,65.8747);
\path[fill=chfourfig5c0,fill opacity=0.5,draw=chfourfig5c0,line width=0.3985pt,draw opacity=0.5] (137.4583,37.0624) .. controls (137.4583,36.3416) and (136.8740,35.7573) ..(136.1532,35.7573) .. controls (135.4324,35.7573) and (134.8481,36.3416) ..(134.8481,37.0624) .. controls (134.8481,37.7832) and (135.4324,38.3675) ..(136.1532,38.3675) .. controls (136.8740,38.3675) and (137.4583,37.7832) ..(137.4583,37.0624);
\path[fill=chfourfig5c0,fill opacity=0.5,draw=chfourfig5c0,line width=0.3985pt,draw opacity=0.5] (99.0120,54.4673) .. controls (99.0120,53.7465) and (98.4277,53.1622) ..(97.7069,53.1622) .. controls (96.9861,53.1622) and (96.4018,53.7465) ..(96.4018,54.4673) .. controls (96.4018,55.1881) and (96.9861,55.7725) ..(97.7069,55.7725) .. controls (98.4277,55.7725) and (99.0120,55.1881) ..(99.0120,54.4673);
\path[fill=chfourfig5c0,fill opacity=0.5,draw=chfourfig5c0,line width=0.3985pt,draw opacity=0.5] (122.6337,77.6407) .. controls (122.6337,76.9199) and (122.0494,76.3356) ..(121.3286,76.3356) .. controls (120.6078,76.3356) and (120.0235,76.9199) ..(120.0235,77.6407) .. controls (120.0235,78.3615) and (120.6078,78.9459) ..(121.3286,78.9459) .. controls (122.0494,78.9459) and (122.6337,78.3615) ..(122.6337,77.6407);
\path[fill=chfourfig5c0,fill opacity=0.5,draw=chfourfig5c0,line width=0.3985pt,draw opacity=0.5] (129.2988,75.9471) .. controls (129.2988,75.2263) and (128.7145,74.6420) ..(127.9937,74.6420) .. controls (127.2729,74.6420) and (126.6886,75.2263) ..(126.6886,75.9471) .. controls (126.6886,76.6679) and (127.2729,77.2522) ..(127.9937,77.2522) .. controls (128.7145,77.2522) and (129.2988,76.6679) ..(129.2988,75.9471);
\path[fill=chfourfig5c0,fill opacity=0.5,draw=chfourfig5c0,line width=0.3985pt,draw opacity=0.5] (136.6215,61.5608) .. controls (136.6215,60.8400) and (136.0371,60.2557) ..(135.3163,60.2557) .. controls (134.5955,60.2557) and (134.0112,60.8400) ..(134.0112,61.5608) .. controls (134.0112,62.2816) and (134.5955,62.8660) ..(135.3163,62.8660) .. controls (136.0371,62.8660) and (136.6215,62.2816) ..(136.6215,61.5608);
\path[fill=chfourfig5c0,fill opacity=0.5,draw=chfourfig5c0,line width=0.3985pt,draw opacity=0.5] (120.5216,52.1759) .. controls (120.5216,51.4551) and (119.9373,50.8708) ..(119.2165,50.8708) .. controls (118.4957,50.8708) and (117.9114,51.4551) ..(117.9114,52.1759) .. controls (117.9114,52.8967) and (118.4957,53.4810) ..(119.2165,53.4810) .. controls (119.9373,53.4810) and (120.5216,52.8967) ..(120.5216,52.1759);
\path[fill=chfourfig5c0,fill opacity=0.5,draw=chfourfig5c0,line width=0.3985pt,draw opacity=0.5] (77.6319,91.5189) .. controls (77.6319,90.7981) and (77.0476,90.2138) ..(76.3268,90.2138) .. controls (75.6060,90.2138) and (75.0217,90.7981) ..(75.0217,91.5189) .. controls (75.0217,92.2397) and (75.6060,92.8240) ..(76.3268,92.8240) .. controls (77.0476,92.8240) and (77.6319,92.2397) ..(77.6319,91.5189);
\path[fill=chfourfig5c0,fill opacity=0.5,draw=chfourfig5c0,line width=0.3985pt,draw opacity=0.5] (76.9743,93.6808) .. controls (76.9743,92.9600) and (76.3900,92.3757) ..(75.6692,92.3757) .. controls (74.9484,92.3757) and (74.3641,92.9600) ..(74.3641,93.6808) .. controls (74.3641,94.4016) and (74.9484,94.9859) ..(75.6692,94.9859) .. controls (76.3900,94.9859) and (76.9743,94.4016) ..(76.9743,93.6808);
\path[fill=chfourfig5c0,fill opacity=0.5,draw=chfourfig5c0,line width=0.3985pt,draw opacity=0.5] (148.5569,45.3315) .. controls (148.5569,44.6107) and (147.9725,44.0264) ..(147.2517,44.0264) .. controls (146.5309,44.0264) and (145.9466,44.6107) ..(145.9466,45.3315) .. controls (145.9466,46.0523) and (146.5309,46.6366) ..(147.2517,46.6366) .. controls (147.9725,46.6366) and (148.5569,46.0523) ..(148.5569,45.3315);
\path[fill=chfourfig5c0,fill opacity=0.5,draw=chfourfig5c0,line width=0.3985pt,draw opacity=0.5] (55.4249,112.8492) .. controls (55.4249,112.1284) and (54.8406,111.5441) ..(54.1198,111.5441) .. controls (53.3990,111.5441) and (52.8146,112.1284) ..(52.8146,112.8492) .. controls (52.8146,113.5700) and (53.3990,114.1543) ..(54.1198,114.1543) .. controls (54.8406,114.1543) and (55.4249,113.5700) ..(55.4249,112.8492);
\path[fill=chfourfig5c0,fill opacity=0.5,draw=chfourfig5c0,line width=0.3985pt,draw opacity=0.5] (130.6139,55.9418) .. controls (130.6139,55.2210) and (130.0296,54.6367) ..(129.3088,54.6367) .. controls (128.5880,54.6367) and (128.0037,55.2210) ..(128.0037,55.9418) .. controls (128.0037,56.6626) and (128.5880,57.2469) ..(129.3088,57.2469) .. controls (130.0296,57.2469) and (130.6139,56.6626) ..(130.6139,55.9418);
\path[fill=chfourfig5c0,fill opacity=0.5,draw=chfourfig5c0,line width=0.3985pt,draw opacity=0.5] (84.6158,82.9609) .. controls (84.6158,82.2401) and (84.0315,81.6557) ..(83.3107,81.6557) .. controls (82.5899,81.6557) and (82.0056,82.2401) ..(82.0056,82.9609) .. controls (82.0056,83.6817) and (82.5899,84.2660) ..(83.3107,84.2660) .. controls (84.0315,84.2660) and (84.6158,83.6817) ..(84.6158,82.9609);
\path[fill=chfourfig5c0,fill opacity=0.5,draw=chfourfig5c0,line width=0.3985pt,draw opacity=0.5] (100.9946,55.5632) .. controls (100.9946,54.8425) and (100.4103,54.2581) ..(99.6895,54.2581) .. controls (98.9687,54.2581) and (98.3843,54.8425) ..(98.3843,55.5632) .. controls (98.3843,56.2841) and (98.9687,56.8684) ..(99.6895,56.8684) .. controls (100.4103,56.8684) and (100.9946,56.2841) ..(100.9946,55.5632);
\path[fill=chfourfig5c0,fill opacity=0.5,draw=chfourfig5c0,line width=0.3985pt,draw opacity=0.5] (113.9462,58.9307) .. controls (113.9462,58.2099) and (113.3619,57.6255) ..(112.6411,57.6255) .. controls (111.9203,57.6255) and (111.3360,58.2099) ..(111.3360,58.9307) .. controls (111.3360,59.6515) and (111.9203,60.2358) ..(112.6411,60.2358) .. controls (113.3619,60.2358) and (113.9462,59.6515) ..(113.9462,58.9307);
\path[fill=chfourfig5c0,fill opacity=0.5,draw=chfourfig5c0,line width=0.3985pt,draw opacity=0.5] (71.0465,73.9545) .. controls (71.0465,73.2337) and (70.4622,72.6494) ..(69.7414,72.6494) .. controls (69.0206,72.6494) and (68.4363,73.2337) ..(68.4363,73.9545) .. controls (68.4363,74.6753) and (69.0206,75.2596) ..(69.7414,75.2596) .. controls (70.4622,75.2596) and (71.0465,74.6753) ..(71.0465,73.9545);
\path[fill=chfourfig5c0,fill opacity=0.5,draw=chfourfig5c0,line width=0.3985pt,draw opacity=0.5] (115.6399,50.7612) .. controls (115.6399,50.0404) and (115.0555,49.4561) ..(114.3347,49.4561) .. controls (113.6140,49.4561) and (113.0296,50.0404) ..(113.0296,50.7612) .. controls (113.0296,51.4820) and (113.6140,52.0663) ..(114.3347,52.0663) .. controls (115.0555,52.0663) and (115.6399,51.4820) ..(115.6399,50.7612);
\path[fill=chfourfig5c0,fill opacity=0.5,draw=chfourfig5c0,line width=0.3985pt,draw opacity=0.5] (56.8894,109.0235) .. controls (56.8894,108.3027) and (56.3051,107.7183) ..(55.5843,107.7183) .. controls (54.8635,107.7183) and (54.2792,108.3027) ..(54.2792,109.0235) .. controls (54.2792,109.7443) and (54.8635,110.3286) ..(55.5843,110.3286) .. controls (56.3051,110.3286) and (56.8894,109.7443) ..(56.8894,109.0235);
\path[fill=chfourfig5c0,fill opacity=0.5,draw=chfourfig5c0,line width=0.3985pt,draw opacity=0.5] (94.4889,81.6557) .. controls (94.4889,80.9349) and (93.9046,80.3506) ..(93.1838,80.3506) .. controls (92.4630,80.3506) and (91.8787,80.9349) ..(91.8787,81.6557) .. controls (91.8787,82.3765) and (92.4630,82.9608) ..(93.1838,82.9608) .. controls (93.9046,82.9608) and (94.4889,82.3765) ..(94.4889,81.6557);
\path[fill=chfourfig5c0,fill opacity=0.5,draw=chfourfig5c0,line width=0.3985pt,draw opacity=0.5] (71.2458,101.5813) .. controls (71.2458,100.8605) and (70.6614,100.2762) ..(69.9406,100.2762) .. controls (69.2198,100.2762) and (68.6355,100.8605) ..(68.6355,101.5813) .. controls (68.6355,102.3021) and (69.2198,102.8864) ..(69.9406,102.8864) .. controls (70.6614,102.8864) and (71.2458,102.3021) ..(71.2458,101.5813);
\path[fill=chfourfig5c0,fill opacity=0.5,draw=chfourfig5c0,line width=0.3985pt,draw opacity=0.5] (83.3107,78.7167) .. controls (83.3107,77.9959) and (82.7263,77.4116) ..(82.0056,77.4116) .. controls (81.2848,77.4116) and (80.7004,77.9959) ..(80.7004,78.7167) .. controls (80.7004,79.4375) and (81.2848,80.0218) ..(82.0056,80.0218) .. controls (82.7263,80.0218) and (83.3107,79.4375) ..(83.3107,78.7167);
\path[fill=chfourfig5c0,fill opacity=0.5,draw=chfourfig5c0,line width=0.3985pt,draw opacity=0.5] (147.4211,33.8843) .. controls (147.4211,33.1635) and (146.8368,32.5791) ..(146.1160,32.5791) .. controls (145.3952,32.5791) and (144.8109,33.1635) ..(144.8109,33.8843) .. controls (144.8109,34.6050) and (145.3952,35.1894) ..(146.1160,35.1894) .. controls (146.8368,35.1894) and (147.4211,34.6050) ..(147.4211,33.8843);
\path[fill=chfourfig5c0,fill opacity=0.5,draw=chfourfig5c0,line width=0.3985pt,draw opacity=0.5] (151.4958,59.7177) .. controls (151.4958,58.9969) and (150.9115,58.4126) ..(150.1907,58.4126) .. controls (149.4699,58.4126) and (148.8856,58.9969) ..(148.8856,59.7177) .. controls (148.8856,60.4385) and (149.4699,61.0228) ..(150.1907,61.0228) .. controls (150.9115,61.0228) and (151.4958,60.4385) ..(151.4958,59.7177);
\path[fill=chfourfig5c0,fill opacity=0.5,draw=chfourfig5c0,line width=0.3985pt,draw opacity=0.5] (144.0935,54.1186) .. controls (144.0935,53.3978) and (143.5092,52.8135) ..(142.7884,52.8135) .. controls (142.0676,52.8135) and (141.4833,53.3978) ..(141.4833,54.1186) .. controls (141.4833,54.8394) and (142.0676,55.4238) ..(142.7884,55.4238) .. controls (143.5092,55.4238) and (144.0935,54.8394) ..(144.0935,54.1186);
\path[fill=chfourfig5c0,fill opacity=0.5,draw=chfourfig5c0,line width=0.3985pt,draw opacity=0.5] (124.3473,69.5410) .. controls (124.3473,68.8202) and (123.7630,68.2359) ..(123.0422,68.2359) .. controls (122.3214,68.2359) and (121.7371,68.8202) ..(121.7371,69.5410) .. controls (121.7371,70.2618) and (122.3214,70.8461) ..(123.0422,70.8461) .. controls (123.7630,70.8461) and (124.3473,70.2618) ..(124.3473,69.5410);
\path[fill=chfourfig5c0,fill opacity=0.5,draw=chfourfig5c0,line width=0.3985pt,draw opacity=0.5] (51.7187,104.1417) .. controls (51.7187,103.4209) and (51.1344,102.8366) ..(50.4136,102.8366) .. controls (49.6928,102.8366) and (49.1085,103.4209) ..(49.1085,104.1417) .. controls (49.1085,104.8625) and (49.6928,105.4468) ..(50.4136,105.4468) .. controls (51.1344,105.4468) and (51.7187,104.8625) ..(51.7187,104.1417);
\path[fill=chfourfig5c0,fill opacity=0.5,draw=chfourfig5c0,line width=0.3985pt,draw opacity=0.5] (61.7213,97.9449) .. controls (61.7213,97.2241) and (61.1370,96.6397) ..(60.4162,96.6397) .. controls (59.6954,96.6397) and (59.1111,97.2241) ..(59.1111,97.9449) .. controls (59.1111,98.6657) and (59.6954,99.2500) ..(60.4162,99.2500) .. controls (61.1370,99.2500) and (61.7213,98.6657) ..(61.7213,97.9449);
\path[fill=chfourfig5c0,fill opacity=0.5,draw=chfourfig5c0,line width=0.3985pt,draw opacity=0.5] (113.1791,84.2660) .. controls (113.1791,83.5452) and (112.5948,82.9609) ..(111.8739,82.9609) .. controls (111.1532,82.9609) and (110.5688,83.5452) ..(110.5688,84.2660) .. controls (110.5688,84.9868) and (111.1532,85.5711) ..(111.8739,85.5711) .. controls (112.5948,85.5711) and (113.1791,84.9868) ..(113.1791,84.2660);
\path[fill=chfourfig5c0,fill opacity=0.5,draw=chfourfig5c0,line width=0.3985pt,draw opacity=0.5] (56.8894,79.7529) .. controls (56.8894,79.0320) and (56.3051,78.4477) ..(55.5843,78.4477) .. controls (54.8635,78.4477) and (54.2792,79.0320) ..(54.2792,79.7529) .. controls (54.2792,80.4736) and (54.8635,81.0580) ..(55.5843,81.0580) .. controls (56.3051,81.0580) and (56.8894,80.4736) ..(56.8894,79.7529);
\path[fill=chfourfig5c0,fill opacity=0.5,draw=chfourfig5c0,line width=0.3985pt,draw opacity=0.5] (89.8762,95.9025) .. controls (89.8762,95.1817) and (89.2918,94.5974) ..(88.5710,94.5974) .. controls (87.8502,94.5974) and (87.2659,95.1817) ..(87.2659,95.9025) .. controls (87.2659,96.6233) and (87.8502,97.2076) ..(88.5710,97.2076) .. controls (89.2918,97.2076) and (89.8762,96.6233) ..(89.8762,95.9025);
\path[fill=chfourfig5c0,fill opacity=0.5,draw=chfourfig5c0,line width=0.3985pt,draw opacity=0.5] (99.0817,61.4313) .. controls (99.0817,60.7105) and (98.4974,60.1262) ..(97.7766,60.1262) .. controls (97.0558,60.1262) and (96.4715,60.7105) ..(96.4715,61.4313) .. controls (96.4715,62.1521) and (97.0558,62.7364) ..(97.7766,62.7364) .. controls (98.4974,62.7364) and (99.0817,62.1521) ..(99.0817,61.4313);
\path[fill=chfourfig5c0,fill opacity=0.5,draw=chfourfig5c0,line width=0.3985pt,draw opacity=0.5] (113.8266,46.8159) .. controls (113.8266,46.0951) and (113.2423,45.5108) ..(112.5215,45.5108) .. controls (111.8007,45.5108) and (111.2164,46.0951) ..(111.2164,46.8159) .. controls (111.2164,47.5367) and (111.8007,48.1210) ..(112.5215,48.1210) .. controls (113.2423,48.1210) and (113.8266,47.5367) ..(113.8266,46.8159);
\path[fill=chfourfig5c0,fill opacity=0.5,draw=chfourfig5c0,line width=0.3985pt,draw opacity=0.5] (75.7987,87.0356) .. controls (75.7987,86.3148) and (75.2144,85.7305) ..(74.4936,85.7305) .. controls (73.7728,85.7305) and (73.1885,86.3148) ..(73.1885,87.0356) .. controls (73.1885,87.7564) and (73.7728,88.3408) ..(74.4936,88.3408) .. controls (75.2144,88.3408) and (75.7987,87.7564) ..(75.7987,87.0356);
\path[fill=chfourfig5c0,fill opacity=0.5,draw=chfourfig5c0,line width=0.3985pt,draw opacity=0.5] (143.3264,38.3675) .. controls (143.3264,37.6467) and (142.7421,37.0624) ..(142.0213,37.0624) .. controls (141.3005,37.0624) and (140.7162,37.6467) ..(140.7162,38.3675) .. controls (140.7162,39.0883) and (141.3005,39.6726) ..(142.0213,39.6726) .. controls (142.7421,39.6726) and (143.3264,39.0883) ..(143.3264,38.3675);
\path[fill=chfourfig5c0,fill opacity=0.5,draw=chfourfig5c0,line width=0.3985pt,draw opacity=0.5] (137.4185,44.0562) .. controls (137.4185,43.3354) and (136.8341,42.7511) ..(136.1133,42.7511) .. controls (135.3925,42.7511) and (134.8082,43.3354) ..(134.8082,44.0562) .. controls (134.8082,44.7771) and (135.3925,45.3614) ..(136.1133,45.3614) .. controls (136.8341,45.3614) and (137.4185,44.7771) ..(137.4185,44.0562);
\path[fill=chfourfig5c0,fill opacity=0.5,draw=chfourfig5c0,line width=0.3985pt,draw opacity=0.5] (142.9179,50.6416) .. controls (142.9179,49.9208) and (142.3336,49.3365) ..(141.6128,49.3365) .. controls (140.8920,49.3365) and (140.3077,49.9208) ..(140.3077,50.6416) .. controls (140.3077,51.3624) and (140.8920,51.9468) ..(141.6128,51.9468) .. controls (142.3336,51.9468) and (142.9179,51.3624) ..(142.9179,50.6416);
\path[fill=chfourfig5c0,fill opacity=0.5,draw=chfourfig5c0,line width=0.3985pt,draw opacity=0.5] (107.4903,70.6768) .. controls (107.4903,69.9560) and (106.9060,69.3716) ..(106.1852,69.3716) .. controls (105.4644,69.3716) and (104.8801,69.9560) ..(104.8801,70.6768) .. controls (104.8801,71.3976) and (105.4644,71.9819) ..(106.1852,71.9819) .. controls (106.9060,71.9819) and (107.4903,71.3976) ..(107.4903,70.6768);
\path[fill=chfourfig5c0,fill opacity=0.5,draw=chfourfig5c0,line width=0.3985pt,draw opacity=0.5] (117.7719,43.3090) .. controls (117.7719,42.5883) and (117.1876,42.0039) ..(116.4668,42.0039) .. controls (115.7460,42.0039) and (115.1617,42.5883) ..(115.1617,43.3090) .. controls (115.1617,44.0298) and (115.7460,44.6142) ..(116.4668,44.6142) .. controls (117.1876,44.6142) and (117.7719,44.0298) ..(117.7719,43.3090);
\path[fill=chfourfig5c0,fill opacity=0.5,draw=chfourfig5c0,line width=0.3985pt,draw opacity=0.5] (106.7033,62.8261) .. controls (106.7033,62.1053) and (106.1190,61.5210) ..(105.3981,61.5210) .. controls (104.6774,61.5210) and (104.0930,62.1053) ..(104.0930,62.8261) .. controls (104.0930,63.5469) and (104.6774,64.1312) ..(105.3981,64.1312) .. controls (106.1190,64.1312) and (106.7033,63.5469) ..(106.7033,62.8261);
\path[fill=chfourfig5c0,fill opacity=0.5,draw=chfourfig5c0,line width=0.3985pt,draw opacity=0.5] (94.0107,65.1674) .. controls (94.0107,64.4466) and (93.4264,63.8622) ..(92.7056,63.8622) .. controls (91.9848,63.8622) and (91.4004,64.4466) ..(91.4004,65.1674) .. controls (91.4004,65.8882) and (91.9848,66.4725) ..(92.7056,66.4725) .. controls (93.4264,66.4725) and (94.0107,65.8882) ..(94.0107,65.1674);
\path[fill=chfourfig5c0,fill opacity=0.5,draw=chfourfig5c0,line width=0.3985pt,draw opacity=0.5] (105.0992,56.3105) .. controls (105.0992,55.5897) and (104.5149,55.0053) ..(103.7941,55.0053) .. controls (103.0733,55.0053) and (102.4890,55.5897) ..(102.4890,56.3105) .. controls (102.4890,57.0313) and (103.0733,57.6156) ..(103.7941,57.6156) .. controls (104.5149,57.6156) and (105.0992,57.0313) ..(105.0992,56.3105);
\path[fill=chfourfig5c0,fill opacity=0.5,draw=chfourfig5c0,line width=0.3985pt,draw opacity=0.5] (77.7415,91.6384) .. controls (77.7415,90.9176) and (77.1572,90.3333) ..(76.4364,90.3333) .. controls (75.7156,90.3333) and (75.1313,90.9176) ..(75.1313,91.6384) .. controls (75.1313,92.3592) and (75.7156,92.9435) ..(76.4364,92.9435) .. controls (77.1572,92.9435) and (77.7415,92.3592) ..(77.7415,91.6384);
\path[fill=chfourfig5c0,fill opacity=0.5,draw=chfourfig5c0,line width=0.3985pt,draw opacity=0.5] (149.4635,50.8907) .. controls (149.4635,50.1699) and (148.8792,49.5856) ..(148.1584,49.5856) .. controls (147.4375,49.5856) and (146.8532,50.1699) ..(146.8532,50.8907) .. controls (146.8532,51.6115) and (147.4375,52.1958) ..(148.1584,52.1958) .. controls (148.8792,52.1958) and (149.4635,51.6115) ..(149.4635,50.8907);
\path[fill=chfourfig5c0,fill opacity=0.5,draw=chfourfig5c0,line width=0.3985pt,draw opacity=0.5] (147.1820,55.4935) .. controls (147.1820,54.7727) and (146.5977,54.1884) ..(145.8769,54.1884) .. controls (145.1561,54.1884) and (144.5717,54.7727) ..(144.5717,55.4935) .. controls (144.5717,56.2143) and (145.1561,56.7986) ..(145.8769,56.7986) .. controls (146.5977,56.7986) and (147.1820,56.2143) ..(147.1820,55.4935);
\path[fill=chfourfig5c0,fill opacity=0.5,draw=chfourfig5c0,line width=0.3985pt,draw opacity=0.5] (62.1497,96.2313) .. controls (62.1497,95.5105) and (61.5654,94.9261) ..(60.8446,94.9261) .. controls (60.1238,94.9261) and (59.5395,95.5105) ..(59.5395,96.2313) .. controls (59.5395,96.9521) and (60.1238,97.5364) ..(60.8446,97.5364) .. controls (61.5654,97.5364) and (62.1497,96.9521) ..(62.1497,96.2313);
\path[fill=chfourfig5c0,fill opacity=0.5,draw=chfourfig5c0,line width=0.3985pt,draw opacity=0.5] (76.6356,89.0182) .. controls (76.6356,88.2974) and (76.0513,87.7131) ..(75.3305,87.7131) .. controls (74.6097,87.7131) and (74.0254,88.2974) ..(74.0254,89.0182) .. controls (74.0254,89.7390) and (74.6097,90.3233) ..(75.3305,90.3233) .. controls (76.0513,90.3233) and (76.6356,89.7390) ..(76.6356,89.0182);
\path[fill=chfourfig5c0,fill opacity=0.5,draw=chfourfig5c0,line width=0.3985pt,draw opacity=0.5] (52.9641,96.2412) .. controls (52.9641,95.5204) and (52.3798,94.9361) ..(51.6590,94.9361) .. controls (50.9382,94.9361) and (50.3538,95.5204) ..(50.3538,96.2412) .. controls (50.3538,96.9620) and (50.9382,97.5463) ..(51.6590,97.5463) .. controls (52.3798,97.5463) and (52.9641,96.9620) ..(52.9641,96.2412);
\path[fill=chfourfig5c0,fill opacity=0.5,draw=chfourfig5c0,line width=0.3985pt,draw opacity=0.5] (57.6466,94.9959) .. controls (57.6466,94.2751) and (57.0623,93.6908) ..(56.3415,93.6908) .. controls (55.6207,93.6908) and (55.0363,94.2751) ..(55.0363,94.9959) .. controls (55.0363,95.7167) and (55.6207,96.3010) ..(56.3415,96.3010) .. controls (57.0623,96.3010) and (57.6466,95.7167) ..(57.6466,94.9959);
\path[fill=chfourfig5c0,fill opacity=0.5,draw=chfourfig5c0,line width=0.3985pt,draw opacity=0.5] (101.9012,76.8935) .. controls (101.9012,76.1727) and (101.3169,75.5884) ..(100.5961,75.5884) .. controls (99.8753,75.5884) and (99.2910,76.1727) ..(99.2910,76.8935) .. controls (99.2910,77.6143) and (99.8753,78.1987) ..(100.5961,78.1987) .. controls (101.3169,78.1987) and (101.9012,77.6143) ..(101.9012,76.8935);
\path[fill=chfourfig5c0,fill opacity=0.5,draw=chfourfig5c0,line width=0.3985pt,draw opacity=0.5] (63.5446,108.1965) .. controls (63.5446,107.4757) and (62.9602,106.8914) ..(62.2394,106.8914) .. controls (61.5186,106.8914) and (60.9343,107.4757) ..(60.9343,108.1965) .. controls (60.9343,108.9173) and (61.5186,109.5017) ..(62.2394,109.5017) .. controls (62.9602,109.5017) and (63.5446,108.9173) ..(63.5446,108.1965);
\path[fill=chfourfig5c0,fill opacity=0.5,draw=chfourfig5c0,line width=0.3985pt,draw opacity=0.5] (59.0613,99.2301) .. controls (59.0613,98.5092) and (58.4770,97.9249) ..(57.7562,97.9249) .. controls (57.0354,97.9249) and (56.4510,98.5092) ..(56.4510,99.2301) .. controls (56.4510,99.9509) and (57.0354,100.5352) ..(57.7562,100.5352) .. controls (58.4770,100.5352) and (59.0613,99.9509) ..(59.0613,99.2301);
\path[fill=chfourfig5c0,fill opacity=0.5,draw=chfourfig5c0,line width=0.3985pt,draw opacity=0.5] (143.2267,62.8759) .. controls (143.2267,62.1551) and (142.6425,61.5708) ..(141.9216,61.5708) .. controls (141.2008,61.5708) and (140.6165,62.1551) ..(140.6165,62.8759) .. controls (140.6165,63.5967) and (141.2008,64.1810) ..(141.9216,64.1810) .. controls (142.6425,64.1810) and (143.2267,63.5967) ..(143.2267,62.8759);
\path[fill=chfourfig5c0,fill opacity=0.5,draw=chfourfig5c0,line width=0.3985pt,draw opacity=0.5] (104.7007,71.9819) .. controls (104.7007,71.2611) and (104.1164,70.6768) ..(103.3956,70.6768) .. controls (102.6748,70.6768) and (102.0905,71.2611) ..(102.0905,71.9819) .. controls (102.0905,72.7027) and (102.6748,73.2870) ..(103.3956,73.2870) .. controls (104.1164,73.2870) and (104.7007,72.7027) ..(104.7007,71.9819);
\path[fill=chfourfig5c0,fill opacity=0.5,draw=chfourfig5c0,line width=0.3985pt,draw opacity=0.5] (130.4146,52.6541) .. controls (130.4146,51.9333) and (129.8303,51.3490) ..(129.1095,51.3490) .. controls (128.3887,51.3490) and (127.8044,51.9333) ..(127.8044,52.6541) .. controls (127.8044,53.3749) and (128.3887,53.9592) ..(129.1095,53.9592) .. controls (129.8303,53.9592) and (130.4146,53.3749) ..(130.4146,52.6541);
\path[fill=chfourfig5c0,fill opacity=0.5,draw=chfourfig5c0,line width=0.3985pt,draw opacity=0.5] (83.0118,80.8587) .. controls (83.0118,80.1379) and (82.4275,79.5536) ..(81.7067,79.5536) .. controls (80.9859,79.5536) and (80.4016,80.1379) ..(80.4016,80.8587) .. controls (80.4016,81.5795) and (80.9859,82.1638) ..(81.7067,82.1638) .. controls (82.4275,82.1638) and (83.0118,81.5795) ..(83.0118,80.8587);
\path[fill=chfourfig5c0,fill opacity=0.5,draw=chfourfig5c0,line width=0.3985pt,draw opacity=0.5] (95.9236,88.4304) .. controls (95.9236,87.7096) and (95.3392,87.1253) ..(94.6184,87.1253) .. controls (93.8976,87.1253) and (93.3133,87.7096) ..(93.3133,88.4304) .. controls (93.3133,89.1512) and (93.8976,89.7355) ..(94.6184,89.7355) .. controls (95.3392,89.7355) and (95.9236,89.1512) ..(95.9236,88.4304);
\path[fill=chfourfig5c0,fill opacity=0.5,draw=chfourfig5c0,line width=0.3985pt,draw opacity=0.5] (127.6350,60.7837) .. controls (127.6350,60.0629) and (127.0507,59.4786) ..(126.3299,59.4786) .. controls (125.6091,59.4786) and (125.0248,60.0629) ..(125.0248,60.7837) .. controls (125.0248,61.5045) and (125.6091,62.0889) ..(126.3299,62.0889) .. controls (127.0507,62.0889) and (127.6350,61.5045) ..(127.6350,60.7837);
\path[fill=chfourfig5c0,fill opacity=0.5,draw=chfourfig5c0,line width=0.3985pt,draw opacity=0.5] (98.3345,78.8761) .. controls (98.3345,78.1553) and (97.7502,77.5710) ..(97.0294,77.5710) .. controls (96.3086,77.5710) and (95.7243,78.1553) ..(95.7243,78.8761) .. controls (95.7243,79.5969) and (96.3086,80.1812) ..(97.0294,80.1812) .. controls (97.7502,80.1812) and (98.3345,79.5969) ..(98.3345,78.8761);
\path[fill=chfourfig5c0,fill opacity=0.5,draw=chfourfig5c0,line width=0.3985pt,draw opacity=0.5] (88.8002,70.3380) .. controls (88.8002,69.6172) and (88.2158,69.0329) ..(87.4951,69.0329) .. controls (86.7742,69.0329) and (86.1899,69.6172) ..(86.1899,70.3380) .. controls (86.1899,71.0588) and (86.7742,71.6431) ..(87.4951,71.6431) .. controls (88.2158,71.6431) and (88.8002,71.0588) ..(88.8002,70.3380);
\path[fill=chfourfig5c0,fill opacity=0.5,draw=chfourfig5c0,line width=0.3985pt,draw opacity=0.5] (68.9543,101.9997) .. controls (68.9543,101.2789) and (68.3700,100.6946) ..(67.6492,100.6946) .. controls (66.9284,100.6946) and (66.3441,101.2789) ..(66.3441,101.9997) .. controls (66.3441,102.7205) and (66.9284,103.3048) ..(67.6492,103.3048) .. controls (68.3700,103.3048) and (68.9543,102.7205) ..(68.9543,101.9997);
\path[fill=chfourfig5c0,fill opacity=0.5,draw=chfourfig5c0,line width=0.3985pt,draw opacity=0.5] (155.7798,30.2080) .. controls (155.7798,29.4872) and (155.1955,28.9029) ..(154.4747,28.9029) .. controls (153.7539,28.9029) and (153.1696,29.4872) ..(153.1696,30.2080) .. controls (153.1696,30.9288) and (153.7539,31.5131) ..(154.4747,31.5131) .. controls (155.1955,31.5131) and (155.7798,30.9288) ..(155.7798,30.2080);
\path[fill=chfourfig5c0,fill opacity=0.5,draw=chfourfig5c0,line width=0.3985pt,draw opacity=0.5] (73.8062,102.1193) .. controls (73.8062,101.3985) and (73.2219,100.8141) ..(72.5011,100.8141) .. controls (71.7803,100.8141) and (71.1960,101.3985) ..(71.1960,102.1193) .. controls (71.1960,102.8401) and (71.7803,103.4244) ..(72.5011,103.4244) .. controls (73.2219,103.4244) and (73.8062,102.8401) ..(73.8062,102.1193);
\path[fill=chfourfig5c0,fill opacity=0.5,draw=chfourfig5c0,line width=0.3985pt,draw opacity=0.5] (131.8692,62.7564) .. controls (131.8692,62.0356) and (131.2849,61.4512) ..(130.5641,61.4512) .. controls (129.8433,61.4512) and (129.2590,62.0356) ..(129.2590,62.7564) .. controls (129.2590,63.4772) and (129.8433,64.0615) ..(130.5641,64.0615) .. controls (131.2849,64.0615) and (131.8692,63.4772) ..(131.8692,62.7564);
\path[fill=chfourfig5c0,fill opacity=0.5,draw=chfourfig5c0,line width=0.3985pt,draw opacity=0.5] (81.9258,73.1575) .. controls (81.9258,72.4367) and (81.3415,71.8524) ..(80.6207,71.8524) .. controls (79.8999,71.8524) and (79.3156,72.4367) ..(79.3156,73.1575) .. controls (79.3156,73.8783) and (79.8999,74.4626) ..(80.6207,74.4626) .. controls (81.3415,74.4626) and (81.9258,73.8783) ..(81.9258,73.1575);
\path[fill=chfourfig5c0,fill opacity=0.5,draw=chfourfig5c0,line width=0.3985pt,draw opacity=0.5] (122.6736,63.0353) .. controls (122.6736,62.3145) and (122.0893,61.7302) ..(121.3685,61.7302) .. controls (120.6476,61.7302) and (120.0633,62.3145) ..(120.0633,63.0353) .. controls (120.0633,63.7561) and (120.6476,64.3405) ..(121.3685,64.3405) .. controls (122.0893,64.3405) and (122.6736,63.7561) ..(122.6736,63.0353);
\path[fill=chfourfig5c0,fill opacity=0.5,draw=chfourfig5c0,line width=0.3985pt,draw opacity=0.5] (87.2260,81.9745) .. controls (87.2260,81.2537) and (86.6417,80.6694) ..(85.9209,80.6694) .. controls (85.2001,80.6694) and (84.6158,81.2537) ..(84.6158,81.9745) .. controls (84.6158,82.6953) and (85.2001,83.2797) ..(85.9209,83.2797) .. controls (86.6417,83.2797) and (87.2260,82.6953) ..(87.2260,81.9745);
\path[fill=chfourfig5c0,fill opacity=0.5,draw=chfourfig5c0,line width=0.3985pt,draw opacity=0.5] (117.2239,77.2821) .. controls (117.2239,76.5613) and (116.6396,75.9770) ..(115.9188,75.9770) .. controls (115.1980,75.9770) and (114.6137,76.5613) ..(114.6137,77.2821) .. controls (114.6137,78.0029) and (115.1980,78.5872) ..(115.9188,78.5872) .. controls (116.6396,78.5872) and (117.2239,78.0029) ..(117.2239,77.2821);
\path[fill=chfourfig5c0,fill opacity=0.5,draw=chfourfig5c0,line width=0.3985pt,draw opacity=0.5] (65.7164,90.8813) .. controls (65.7164,90.1605) and (65.1321,89.5761) ..(64.4113,89.5761) .. controls (63.6905,89.5761) and (63.1062,90.1605) ..(63.1062,90.8813) .. controls (63.1062,91.6021) and (63.6905,92.1864) ..(64.4113,92.1864) .. controls (65.1321,92.1864) and (65.7164,91.6021) ..(65.7164,90.8813);
\path[fill=chfourfig5c0,fill opacity=0.5,draw=chfourfig5c0,line width=0.3985pt,draw opacity=0.5] (98.6732,84.5649) .. controls (98.6732,83.8441) and (98.0889,83.2598) ..(97.3681,83.2598) .. controls (96.6473,83.2598) and (96.0630,83.8441) ..(96.0630,84.5649) .. controls (96.0630,85.2857) and (96.6473,85.8700) ..(97.3681,85.8700) .. controls (98.0889,85.8700) and (98.6732,85.2857) ..(98.6732,84.5649);
\path[fill=chfourfig5c0,fill opacity=0.5,draw=chfourfig5c0,line width=0.3985pt,draw opacity=0.5] (82.3542,76.5150) .. controls (82.3542,75.7941) and (81.7699,75.2098) ..(81.0491,75.2098) .. controls (80.3283,75.2098) and (79.7440,75.7941) ..(79.7440,76.5150) .. controls (79.7440,77.2357) and (80.3283,77.8201) ..(81.0491,77.8201) .. controls (81.7699,77.8201) and (82.3542,77.2357) ..(82.3542,76.5150);
\path[fill=chfourfig5c0,fill opacity=0.5,draw=chfourfig5c0,line width=0.3985pt,draw opacity=0.5] (83.2310,88.0219) .. controls (83.2310,87.3011) and (82.6467,86.7168) ..(81.9258,86.7168) .. controls (81.2050,86.7168) and (80.6207,87.3011) ..(80.6207,88.0219) .. controls (80.6207,88.7427) and (81.2050,89.3271) ..(81.9258,89.3271) .. controls (82.6467,89.3271) and (83.2310,88.7427) ..(83.2310,88.0219);
\path[fill=chfourfig5c0,fill opacity=0.5,draw=chfourfig5c0,line width=0.3985pt,draw opacity=0.5] (154.1659,47.0550) .. controls (154.1659,46.3343) and (153.5816,45.7499) ..(152.8608,45.7499) .. controls (152.1400,45.7499) and (151.5557,46.3343) ..(151.5557,47.0550) .. controls (151.5557,47.7758) and (152.1400,48.3602) ..(152.8608,48.3602) .. controls (153.5816,48.3602) and (154.1659,47.7758) ..(154.1659,47.0550);
\path[fill=chfourfig5c0,fill opacity=0.5,draw=chfourfig5c0,line width=0.3985pt,draw opacity=0.5] (106.0357,54.3976) .. controls (106.0357,53.6768) and (105.4514,53.0925) ..(104.7306,53.0925) .. controls (104.0098,53.0925) and (103.4255,53.6768) ..(103.4255,54.3976) .. controls (103.4255,55.1184) and (104.0098,55.7027) ..(104.7306,55.7027) .. controls (105.4514,55.7027) and (106.0357,55.1184) ..(106.0357,54.3976);
\path[fill=chfourfig5c0,fill opacity=0.5,draw=chfourfig5c0,line width=0.3985pt,draw opacity=0.5] (123.0422,54.2482) .. controls (123.0422,53.5274) and (122.4579,52.9430) ..(121.7371,52.9430) .. controls (121.0163,52.9430) and (120.4319,53.5274) ..(120.4319,54.2482) .. controls (120.4319,54.9690) and (121.0163,55.5533) ..(121.7371,55.5533) .. controls (122.4579,55.5533) and (123.0422,54.9690) ..(123.0422,54.2482);
\path[fill=chfourfig5c0,fill opacity=0.5,draw=chfourfig5c0,line width=0.3985pt,draw opacity=0.5] (61.5221,74.7914) .. controls (61.5221,74.0706) and (60.9378,73.4863) ..(60.2170,73.4863) .. controls (59.4962,73.4863) and (58.9119,74.0706) ..(58.9119,74.7914) .. controls (58.9119,75.5122) and (59.4962,76.0965) ..(60.2170,76.0965) .. controls (60.9378,76.0965) and (61.5221,75.5122) ..(61.5221,74.7914);
\path[fill=chfourfig5c0,fill opacity=0.5,draw=chfourfig5c0,line width=0.3985pt,draw opacity=0.5] (58.0949,97.9947) .. controls (58.0949,97.2739) and (57.5106,96.6896) ..(56.7898,96.6896) .. controls (56.0690,96.6896) and (55.4847,97.2739) ..(55.4847,97.9947) .. controls (55.4847,98.7155) and (56.0690,99.2998) ..(56.7898,99.2998) .. controls (57.5106,99.2998) and (58.0949,98.7155) ..(58.0949,97.9947);
\path[fill=chfourfig5c0,fill opacity=0.5,draw=chfourfig5c0,line width=0.3985pt,draw opacity=0.5] (54.7972,105.5564) .. controls (54.7972,104.8356) and (54.2129,104.2513) ..(53.4921,104.2513) .. controls (52.7713,104.2513) and (52.1870,104.8356) ..(52.1870,105.5564) .. controls (52.1870,106.2772) and (52.7713,106.8615) ..(53.4921,106.8615) .. controls (54.2129,106.8615) and (54.7972,106.2772) ..(54.7972,105.5564);
\path[fill=chfourfig5c0,fill opacity=0.5,draw=chfourfig5c0,line width=0.3985pt,draw opacity=0.5] (118.4992,58.7314) .. controls (118.4992,58.0106) and (117.9149,57.4263) ..(117.1940,57.4263) .. controls (116.4733,57.4263) and (115.8889,58.0106) ..(115.8889,58.7314) .. controls (115.8889,59.4522) and (116.4733,60.0365) ..(117.1940,60.0365) .. controls (117.9149,60.0365) and (118.4992,59.4522) ..(118.4992,58.7314);
\path[fill=chfourfig5c0,fill opacity=0.5,draw=chfourfig5c0,line width=0.3985pt,draw opacity=0.5] (76.5260,68.8835) .. controls (76.5260,68.1627) and (75.9417,67.5784) ..(75.2209,67.5784) .. controls (74.5001,67.5784) and (73.9158,68.1627) ..(73.9158,68.8835) .. controls (73.9158,69.6043) and (74.5001,70.1886) ..(75.2209,70.1886) .. controls (75.9417,70.1886) and (76.5260,69.6043) ..(76.5260,68.8835);
\path[fill=chfourfig5c0,fill opacity=0.5,draw=chfourfig5c0,line width=0.3985pt,draw opacity=0.5] (82.8225,102.0196) .. controls (82.8225,101.2988) and (82.2382,100.7145) ..(81.5174,100.7145) .. controls (80.7966,100.7145) and (80.2122,101.2988) ..(80.2122,102.0196) .. controls (80.2122,102.7404) and (80.7966,103.3248) ..(81.5174,103.3248) .. controls (82.2382,103.3248) and (82.8225,102.7404) ..(82.8225,102.0196);
\path[fill=none,draw=chfourfig5c0,line width=0.3985pt,draw opacity=1.0] (50.3837,120.7696) --(157.3540,120.7696);
\path[fill=none,draw=chfourfig5c0,line width=0.3985pt,draw opacity=1.0] (50.3837,120.7696) --(50.3837,125.7310);
\path[fill=none,draw=chfourfig5c0,line width=0.3985pt,draw opacity=1.0] (86.0405,120.7696) --(86.0405,125.7310);
\path[fill=none,draw=chfourfig5c0,line width=0.3985pt,draw opacity=1.0] (121.6972,120.7696) --(121.6972,125.7310);
\path[fill=none,draw=chfourfig5c0,line width=0.3985pt,draw opacity=1.0] (157.3540,120.7696) --(157.3540,125.7310);
\path[fill=none,draw=chfourfig5c0,line width=0.3985pt,draw opacity=1.0] (38.4085,119.7733) --(38.4085,39.6228);
\path[fill=none,draw=chfourfig5c0,line width=0.3985pt,draw opacity=1.0] (38.4085,119.7733) --(33.4470,119.7733);
\path[fill=none,draw=chfourfig5c0,line width=0.3985pt,draw opacity=1.0] (38.4085,79.6931) --(33.4470,79.6931);
\path[fill=none,draw=chfourfig5c0,line width=0.3985pt,draw opacity=1.0] (38.4085,39.6228) --(33.4470,39.6228);
\path[fill=none,draw=chfourfig5c0,line width=0.3985pt,draw opacity=1.0] (38.4085,120.7696)--(38.4085,16.7483)--(169.3392,16.7483)--(169.3392,120.7696)--cycle;
\path[fill=none,draw=chfourfig5c1,line width=1.1955pt,draw opacity=1.0] (50.4136,101.1529) --(50.9616,100.8341) --(51.0114,100.8042) --(51.6590,100.4156) --(51.8981,100.2762) --(51.9678,100.2363) --(52.1172,100.1467) --(52.4161,99.9773) --(53.1534,99.5389) --(53.2231,99.4891) --(53.3925,99.3995) --(53.4224,99.3795) --(53.4523,99.3596) --(53.4921,99.3397) --(54.1198,98.9711) --(54.2393,98.8914) --(55.5843,98.1043) --(55.9828,97.8652) --(56.3415,97.6560) --(56.4411,97.5962) --(56.4909,97.5663) --(56.7898,97.3870) --(57.3378,97.0682) --(57.3477,97.0582) --(57.6167,96.8988) --(57.7562,96.8191) --(57.8857,96.7394) --(58.0551,96.6398) --(58.2842,96.5103) --(58.4038,96.4405) --(58.6329,96.3010) --(59.0812,96.0320) --(59.2705,95.9225) --(59.3701,95.8627) --(59.4399,95.8228) --(59.5993,95.7232) --(59.7089,95.6634) --(60.2170,95.3645) --(60.4162,95.2450) --(60.4959,95.2051) --(60.8447,94.9959) --(60.8447,94.9859) --(60.9742,94.9162) --(61.0339,94.8863) --(61.2730,94.7369) --(61.4125,94.6572) --(61.8210,94.4181) --(61.9804,94.3185) --(62.0103,94.3085) --(62.0302,94.2886) --(62.2394,94.1690) --(62.3391,94.1092) --(62.9866,93.7307) --(63.0464,93.6908) --(63.4350,93.4616) --(63.5744,93.3819) --(63.6940,93.3122) --(64.4113,92.8838) --(64.5109,92.8340) --(64.5209,92.8240) --(64.6205,92.7643) --(64.7002,92.7144) --(65.1286,92.4654) --(65.7264,92.1067) --(65.9755,91.9573) --(66.1249,91.8776) --(66.2245,91.8178) --(66.2544,91.7979) --(66.2843,91.7779) --(66.4039,91.7082) --(66.8622,91.4392) --(66.9020,91.4093) --(66.9717,91.3695) --(67.2208,91.2200) --(67.4699,91.0806) --(67.6492,90.9709) --(67.8385,90.8614) --(68.1972,90.6521) --(68.7152,90.3433) --(69.1137,90.1042) --(69.1237,90.1042) --(69.3030,89.9946) --(69.3628,89.9647) --(69.4027,89.9348) --(69.5621,89.8452) --(69.6019,89.8153) --(69.6517,89.7953) --(69.7414,89.7356) --(69.8908,89.6459) --(69.9406,89.6160) --(70.1798,89.4765) --(70.2894,89.4168) --(70.4189,89.3371) --(70.7078,89.1677) --(71.0565,88.9585) --(71.1063,88.9286) --(72.5011,88.1116) --(73.0690,87.7729) --(73.1387,87.7330) --(73.4376,87.5537) --(73.8062,87.3345) --(73.8361,87.3146) --(73.9557,87.2449) --(73.9855,87.2250) --(74.0852,87.1751) --(74.4936,86.9261) --(74.5136,86.9161) --(74.8124,86.7368) --(74.9320,86.6670) --(75.0416,86.6073) --(75.0515,86.6073) --(75.0715,86.5873) --(75.0814,86.5873) --(75.2209,86.4977) --(75.3205,86.4379) --(75.3305,86.4379) --(75.4999,86.3383) --(75.6692,86.2386) --(75.8785,86.1091) --(76.2371,85.8999) --(76.3268,85.8501) --(76.4364,85.7803) --(76.5161,85.7305) --(76.7153,85.6210) --(77.4028,85.2125) --(77.4227,85.2025) --(77.8311,84.9534) --(77.9706,84.8737) --(78.4489,84.5948) --(79.6544,83.8775) --(79.6942,83.8575) --(79.8337,83.7778) --(79.8436,83.7679) --(80.6207,83.3096) --(80.7403,83.2399) --(81.0491,83.0605) --(81.1089,83.0207) --(81.3181,82.9011) --(81.3978,82.8513) --(81.4078,82.8513) --(81.5174,82.7816) --(81.5971,82.7317) --(81.7067,82.6720) --(81.8761,82.5723) --(81.9259,82.5325) --(81.9557,82.5225) --(82.0056,82.4926) --(82.1151,82.4329) --(82.3543,82.2834) --(82.5137,82.1938) --(82.5336,82.1838) --(82.6133,82.1340) --(82.7727,82.0443) --(83.3107,81.7255) --(83.3306,81.7056) --(83.5996,81.5561) --(84.1077,81.2572) --(84.7055,80.8986) --(84.8549,80.8089) --(84.9247,80.7691) --(85.5025,80.4303) --(85.5922,80.3706) --(85.6719,80.3208) --(85.7117,80.3008) --(85.7416,80.2809) --(85.9209,80.1813) --(86.0405,80.1115) --(86.1102,80.0617) --(86.2198,80.0020) --(86.6681,79.7429) --(86.6781,79.7330) --(87.0069,79.5337) --(87.0866,79.4938) --(87.1962,79.4241) --(87.2360,79.4042) --(87.4153,79.2946) --(87.4951,79.2448) --(87.7541,79.0953) --(88.1526,78.8562) --(88.5710,78.6171) --(88.6607,78.5573) --(89.2983,78.1887) --(89.4776,78.0791) --(90.4540,77.5013) --(90.4739,77.4913) --(91.0019,77.1725) --(91.0318,77.1626) --(91.3606,76.9733) --(91.4303,76.9234) --(91.5698,76.8438) --(91.5997,76.8238) --(91.8089,76.7043) --(92.3370,76.3954) --(92.4465,76.3257) --(92.6956,76.1762) --(92.7056,76.1762) --(92.7554,76.1364) --(93.1838,75.8873) --(93.6421,75.6183) --(93.7317,75.5685) --(93.7616,75.5486) --(93.8413,75.4988) --(93.8712,75.4888) --(93.9908,75.4091) --(94.1004,75.3493) --(94.2996,75.2298) --(94.4789,75.1302) --(94.5786,75.0604) --(94.6184,75.0405) --(94.7878,74.9409) --(95.0369,74.7914) --(95.1365,74.7416) --(95.5450,74.4925) --(95.5549,74.4925) --(95.6147,74.4527) --(95.9335,74.2634) --(96.1328,74.1438) --(96.4615,73.9545) --(96.6807,73.8250) --(96.8899,73.7055) --(96.9298,73.6756) --(96.9696,73.6557) --(97.0294,73.6158) --(97.3283,73.4365) --(97.3681,73.4166) --(97.5375,73.3169) --(97.7069,73.2173) --(97.7766,73.1774) --(97.9759,73.0579) --(98.5238,72.7391) --(98.5238,72.7291) --(98.7530,72.5996) --(99.1615,72.3605) --(99.3308,72.2609) --(99.3408,72.2509) --(99.4504,72.1812) --(99.6696,72.0616) --(99.6895,72.0417) --(99.7592,72.0018) --(100.5961,71.5137) --(100.6260,71.4937) --(100.8153,71.3842) --(100.8950,71.3343) --(100.9547,71.2945) --(100.9946,71.2746) --(101.3632,71.0554) --(103.1665,69.9894) --(103.2462,69.9495) --(103.3956,69.8598) --(103.4454,69.8300) --(103.7842,69.6307) --(103.7941,69.6207) --(104.2225,69.3617) --(104.2425,69.3517) --(104.3022,69.3219) --(104.4915,69.2123) --(104.5314,69.1824) --(104.7306,69.0628) --(105.1192,68.8436) --(105.3981,68.6743) --(105.5675,68.5747) --(105.7269,68.4850) --(105.7568,68.4651) --(105.8265,68.4153) --(106.1852,68.2060) --(106.6535,67.9370) --(106.7630,67.8673) --(107.0520,67.6979) --(107.1815,67.6182) --(107.1914,67.6083) --(107.3807,67.4987) --(107.4604,67.4588) --(107.4704,67.4489) --(107.5501,67.3990) --(108.1578,67.0404) --(108.3870,66.9009) --(109.0744,66.5024) --(109.1940,66.4326) --(109.3932,66.3131) --(109.4829,66.2633) --(109.7718,66.0939) --(109.8415,66.0441) --(109.9412,65.9843) --(110.6485,65.5659) --(110.7780,65.4962) --(111.3459,65.1574) --(111.4455,65.0976) --(111.5751,65.0279) --(111.6548,64.9781) --(111.8739,64.8486) --(111.9536,64.7988) --(112.0234,64.7589) --(112.3522,64.5597) --(112.5215,64.4600) --(112.6411,64.4003) --(112.6411,64.3903) --(113.0595,64.1512) --(113.2189,64.0515) --(113.5278,63.8722) --(113.6772,63.7826) --(114.2252,63.4637) --(114.3347,63.3940) --(114.4344,63.3342) --(114.4842,63.3043) --(114.6834,63.1848) --(114.7831,63.1250) --(114.7930,63.1250) --(115.6199,62.6368) --(115.6698,62.6069) --(115.9188,62.4575) --(116.1081,62.3479) --(116.1579,62.3180) --(116.1779,62.3081) --(116.3273,62.2184) --(116.4668,62.1387) --(116.6461,62.0291) --(116.9350,61.8597) --(117.0048,61.8199) --(117.1841,61.7103) --(117.1941,61.7003) --(117.3535,61.6107) --(118.1206,61.1624) --(118.1804,61.1225) --(118.2302,61.0926) --(118.4095,60.9930) --(118.6088,60.8734) --(118.8080,60.7539) --(119.0272,60.6244) --(119.2165,60.5148) --(119.4158,60.3952) --(120.1331,59.9668) --(120.1530,59.9568) --(120.2825,59.8772) --(120.3024,59.8672) --(120.8205,59.5583) --(120.8803,59.5285) --(121.1194,59.3890) --(121.3286,59.2595) --(121.3685,59.2395) --(121.6275,59.0901) --(121.7371,59.0204) --(122.1555,58.7812) --(122.2253,58.7314) --(122.4942,58.5720) --(122.5540,58.5421) --(122.7832,58.4027) --(122.8430,58.3628) --(122.8529,58.3628) --(123.0422,58.2532) --(123.1219,58.2034) --(123.5105,57.9743) --(123.7496,57.8348) --(124.0086,57.6754) --(124.2377,57.5459) --(124.3573,57.4761) --(124.5167,57.3765) --(124.6661,57.2968) --(124.7060,57.2669) --(125.1244,57.0178) --(125.5429,56.7787) --(125.5827,56.7488) --(125.6525,56.7090) --(126.3299,56.3105) --(126.4196,56.2507) --(126.4296,56.2507) --(126.5690,56.1710) --(126.6687,56.1112) --(126.8679,55.9917) --(127.0472,55.8821) --(127.6051,55.5533) --(127.6450,55.5334) --(127.8144,55.4338) --(127.9937,55.3242) --(128.5018,55.0253) --(128.8405,54.8260) --(129.0099,54.7264) --(129.0398,54.7065) --(129.1095,54.6666) --(129.2789,54.5670) --(129.3088,54.5471) --(129.3586,54.5172) --(129.5977,54.3777) --(129.6376,54.3578) --(130.5641,53.7999) --(130.7733,53.6903) --(131.0722,53.5109) --(131.2017,53.4312) --(131.3810,53.3217) --(131.6301,53.1722) --(132.0286,52.9431) --(132.0486,52.9331) --(132.2179,52.8335) --(132.5268,52.6442) --(132.7459,52.5147) --(133.9913,51.7874) --(134.1507,51.6878) --(134.2304,51.6380) --(134.8082,51.2992) --(135.3163,51.0003) --(136.1134,50.5321) --(136.1532,50.5022) --(136.5418,50.2730) --(137.3587,49.7948) --(137.4683,49.7251) --(137.9565,49.4362) --(138.0163,49.4063) --(138.7037,48.9978) --(138.7435,48.9779) --(138.7635,48.9679) --(138.8033,48.9380) --(138.9926,48.8284) --(139.1221,48.7487) --(139.1321,48.7487) --(139.3612,48.6093) --(139.7996,48.3502) --(139.8793,48.3004) --(140.0387,48.2108) --(140.2778,48.0713) --(141.4733,47.3639) --(141.5231,47.3340) --(141.6128,47.2842) --(141.8818,47.1248) --(141.9117,47.1049) --(141.9216,47.0949) --(142.0213,47.0352) --(142.2803,46.8857) --(142.5792,46.7064) --(142.7585,46.6068) --(142.7884,46.5868) --(142.9478,46.4972) --(143.0176,46.4573) --(143.3364,46.2680) --(143.4958,46.1684) --(144.2031,45.7500) --(144.2928,45.7001) --(144.6614,45.4810) --(145.0400,45.2518) --(145.8071,44.8035) --(145.8769,44.7636) --(146.0761,44.6441) --(146.1160,44.6142) --(146.7038,44.2655) --(147.0425,44.0662) --(147.0923,44.0463) --(147.2517,43.9467) --(147.2816,43.9268) --(147.8196,43.6179) --(147.9890,43.5183) --(148.1584,43.4187) --(148.1982,43.3888) --(148.3177,43.3190) --(148.4074,43.2692) --(148.6565,43.1198) --(149.1347,42.8408) --(149.3937,42.6814) --(149.6328,42.5419) --(149.6528,42.5320) --(150.1908,42.2132) --(150.2206,42.1932) --(150.2306,42.1833) --(150.4000,42.0836) --(150.7088,41.9043) --(151.0376,41.7150) --(151.3165,41.5457) --(151.9043,41.2069) --(151.9542,41.1770) --(152.0936,41.0874) --(152.5420,40.8184) --(152.8608,40.6390) --(152.9006,40.6092) --(152.9803,40.5594) --(153.3290,40.3601) --(153.7375,40.1210) --(154.1759,39.8620) --(154.2057,39.8420) --(154.3552,39.7523) --(154.4747,39.6826) --(154.5544,39.6328) --(154.6840,39.5531) --(155.2220,39.2443) --(155.3614,39.1546) --(155.8596,38.8657) --(156.3378,38.5768) --(156.4474,38.5170) --(156.7263,38.3576) --(156.7462,38.3476) --(156.7562,38.3277) --(156.8658,38.2679) --(157.2145,38.0686);
\path[fill=chfourfig5c0,fill opacity=0.5,draw=chfourfig5c0,line width=0.3985pt,draw opacity=0.5] (297.2811,49.5059) .. controls (297.2811,48.7851) and (296.6967,48.2008) ..(295.9759,48.2008) .. controls (295.2552,48.2008) and (294.6708,48.7851) ..(294.6708,49.5059) .. controls (294.6708,50.2267) and (295.2552,50.8110) ..(295.9759,50.8110) .. controls (296.6967,50.8110) and (297.2811,50.2267) ..(297.2811,49.5059);
\path[fill=chfourfig5c0,fill opacity=0.5,draw=chfourfig5c0,line width=0.3985pt,draw opacity=0.5] (282.3269,75.7777) .. controls (282.3269,75.0569) and (281.7426,74.4726) ..(281.0219,74.4726) .. controls (280.3010,74.4726) and (279.7167,75.0569) ..(279.7167,75.7777) .. controls (279.7167,76.4985) and (280.3010,77.0829) ..(281.0219,77.0829) .. controls (281.7426,77.0829) and (282.3269,76.4985) ..(282.3269,75.7777);
\path[fill=chfourfig5c0,fill opacity=0.5,draw=chfourfig5c0,line width=0.3985pt,draw opacity=0.5] (273.4601,63.4339) .. controls (273.4601,62.7131) and (272.8758,62.1287) ..(272.1550,62.1287) .. controls (271.4341,62.1287) and (270.8498,62.7131) ..(270.8498,63.4339) .. controls (270.8498,64.1547) and (271.4341,64.7390) ..(272.1550,64.7390) .. controls (272.8758,64.7390) and (273.4601,64.1547) ..(273.4601,63.4339);
\path[fill=chfourfig5c0,fill opacity=0.5,draw=chfourfig5c0,line width=0.3985pt,draw opacity=0.5] (238.1520,111.4146) .. controls (238.1520,110.6938) and (237.5677,110.1094) ..(236.8469,110.1094) .. controls (236.1261,110.1094) and (235.5418,110.6938) ..(235.5418,111.4146) .. controls (235.5418,112.1354) and (236.1261,112.7197) ..(236.8469,112.7197) .. controls (237.5677,112.7197) and (238.1520,112.1354) ..(238.1520,111.4146);
\path[fill=chfourfig5c0,fill opacity=0.5,draw=chfourfig5c0,line width=0.3985pt,draw opacity=0.5] (299.7618,46.2182) .. controls (299.7618,45.4974) and (299.1775,44.9131) ..(298.4567,44.9131) .. controls (297.7359,44.9131) and (297.1516,45.4974) ..(297.1516,46.2182) .. controls (297.1516,46.9390) and (297.7359,47.5233) ..(298.4567,47.5233) .. controls (299.1775,47.5233) and (299.7618,46.9390) ..(299.7618,46.2182);
\path[fill=chfourfig5c0,fill opacity=0.5,draw=chfourfig5c0,line width=0.3985pt,draw opacity=0.5] (317.6350,44.8633) .. controls (317.6350,44.1425) and (317.0507,43.5581) ..(316.3299,43.5581) .. controls (315.6091,43.5581) and (315.0248,44.1425) ..(315.0248,44.8633) .. controls (315.0248,45.5841) and (315.6091,46.1684) ..(316.3299,46.1684) .. controls (317.0507,46.1684) and (317.6350,45.5841) ..(317.6350,44.8633);
\path[fill=chfourfig5c0,fill opacity=0.5,draw=chfourfig5c0,line width=0.3985pt,draw opacity=0.5] (236.2292,111.1954) .. controls (236.2292,110.4746) and (235.6449,109.8903) ..(234.9241,109.8903) .. controls (234.2033,109.8903) and (233.6190,110.4746) ..(233.6190,111.1954) .. controls (233.6190,111.9162) and (234.2033,112.5005) ..(234.9241,112.5005) .. controls (235.6449,112.5005) and (236.2292,111.9162) ..(236.2292,111.1954);
\path[fill=chfourfig5c0,fill opacity=0.5,draw=chfourfig5c0,line width=0.3985pt,draw opacity=0.5] (332.6090,50.7712) .. controls (332.6090,50.0504) and (332.0247,49.4661) ..(331.3039,49.4661) .. controls (330.5831,49.4661) and (329.9988,50.0504) ..(329.9988,50.7712) .. controls (329.9988,51.4920) and (330.5831,52.0763) ..(331.3039,52.0763) .. controls (332.0247,52.0763) and (332.6090,51.4920) ..(332.6090,50.7712);
\path[fill=chfourfig5c0,fill opacity=0.5,draw=chfourfig5c0,line width=0.3985pt,draw opacity=0.5] (286.3419,67.3990) .. controls (286.3419,66.6782) and (285.7576,66.0939) ..(285.0368,66.0939) .. controls (284.3160,66.0939) and (283.7317,66.6782) ..(283.7317,67.3990) .. controls (283.7317,68.1198) and (284.3160,68.7042) ..(285.0368,68.7042) .. controls (285.7576,68.7042) and (286.3419,68.1198) ..(286.3419,67.3990);
\path[fill=chfourfig5c0,fill opacity=0.5,draw=chfourfig5c0,line width=0.3985pt,draw opacity=0.5] (333.5655,59.5185) .. controls (333.5655,58.7977) and (332.9811,58.2134) ..(332.2603,58.2134) .. controls (331.5396,58.2134) and (330.9552,58.7977) ..(330.9552,59.5185) .. controls (330.9552,60.2393) and (331.5396,60.8236) ..(332.2603,60.8236) .. controls (332.9811,60.8236) and (333.5655,60.2393) ..(333.5655,59.5185);
\path[fill=chfourfig5c0,fill opacity=0.5,draw=chfourfig5c0,line width=0.3985pt,draw opacity=0.5] (272.5236,81.0979) .. controls (272.5236,80.3770) and (271.9393,79.7927) ..(271.2184,79.7927) .. controls (270.4977,79.7927) and (269.9133,80.3770) ..(269.9133,81.0979) .. controls (269.9133,81.8186) and (270.4977,82.4030) ..(271.2184,82.4030) .. controls (271.9393,82.4030) and (272.5236,81.8186) ..(272.5236,81.0979);
\path[fill=chfourfig5c0,fill opacity=0.5,draw=chfourfig5c0,line width=0.3985pt,draw opacity=0.5] (258.5359,75.6183) .. controls (258.5359,74.8975) and (257.9515,74.3132) ..(257.2307,74.3132) .. controls (256.5099,74.3132) and (255.9256,74.8975) ..(255.9256,75.6183) .. controls (255.9256,76.3391) and (256.5099,76.9234) ..(257.2307,76.9234) .. controls (257.9515,76.9234) and (258.5359,76.3391) ..(258.5359,75.6183);
\path[fill=chfourfig5c0,fill opacity=0.5,draw=chfourfig5c0,line width=0.3985pt,draw opacity=0.5] (270.3119,92.5351) .. controls (270.3119,91.8143) and (269.7275,91.2300) ..(269.0067,91.2300) .. controls (268.2859,91.2300) and (267.7016,91.8143) ..(267.7016,92.5351) .. controls (267.7016,93.2559) and (268.2859,93.8402) ..(269.0067,93.8402) .. controls (269.7275,93.8402) and (270.3119,93.2559) ..(270.3119,92.5351);
\path[fill=chfourfig5c0,fill opacity=0.5,draw=chfourfig5c0,line width=0.3985pt,draw opacity=0.5] (258.7052,91.0507) .. controls (258.7052,90.3299) and (258.1209,89.7455) ..(257.4001,89.7455) .. controls (256.6793,89.7455) and (256.0950,90.3299) ..(256.0950,91.0507) .. controls (256.0950,91.7715) and (256.6793,92.3558) ..(257.4001,92.3558) .. controls (258.1209,92.3558) and (258.7052,91.7715) ..(258.7052,91.0507);
\path[fill=chfourfig5c0,fill opacity=0.5,draw=chfourfig5c0,line width=0.3985pt,draw opacity=0.5] (269.8835,65.1674) .. controls (269.8835,64.4466) and (269.2991,63.8623) ..(268.5783,63.8623) .. controls (267.8575,63.8623) and (267.2732,64.4466) ..(267.2732,65.1674) .. controls (267.2732,65.8882) and (267.8575,66.4725) ..(268.5783,66.4725) .. controls (269.2991,66.4725) and (269.8835,65.8882) ..(269.8835,65.1674);
\path[fill=chfourfig5c0,fill opacity=0.5,draw=chfourfig5c0,line width=0.3985pt,draw opacity=0.5] (264.9120,92.6447) .. controls (264.9120,91.9239) and (264.3277,91.3396) ..(263.6069,91.3396) .. controls (262.8861,91.3396) and (262.3018,91.9239) ..(262.3018,92.6447) .. controls (262.3018,93.3655) and (262.8861,93.9498) ..(263.6069,93.9498) .. controls (264.3277,93.9498) and (264.9120,93.3655) ..(264.9120,92.6447);
\path[fill=chfourfig5c0,fill opacity=0.5,draw=chfourfig5c0,line width=0.3985pt,draw opacity=0.5] (272.7428,80.3805) .. controls (272.7428,79.6597) and (272.1584,79.0754) ..(271.4376,79.0754) .. controls (270.7169,79.0754) and (270.1325,79.6597) ..(270.1325,80.3805) .. controls (270.1325,81.1013) and (270.7169,81.6857) ..(271.4376,81.6857) .. controls (272.1584,81.6857) and (272.7428,81.1013) ..(272.7428,80.3805);
\path[fill=chfourfig5c0,fill opacity=0.5,draw=chfourfig5c0,line width=0.3985pt,draw opacity=0.5] (308.5091,58.7912) .. controls (308.5091,58.0704) and (307.9248,57.4861) ..(307.2040,57.4861) .. controls (306.4832,57.4861) and (305.8989,58.0704) ..(305.8989,58.7912) .. controls (305.8989,59.5120) and (306.4832,60.0963) ..(307.2040,60.0963) .. controls (307.9248,60.0963) and (308.5091,59.5120) ..(308.5091,58.7912);
\path[fill=chfourfig5c0,fill opacity=0.5,draw=chfourfig5c0,line width=0.3985pt,draw opacity=0.5] (310.0832,62.3778) .. controls (310.0832,61.6570) and (309.4989,61.0727) ..(308.7781,61.0727) .. controls (308.0573,61.0727) and (307.4730,61.6570) ..(307.4730,62.3778) .. controls (307.4730,63.0986) and (308.0573,63.6829) ..(308.7781,63.6829) .. controls (309.4989,63.6829) and (310.0832,63.0986) ..(310.0832,62.3778);
\path[fill=chfourfig5c0,fill opacity=0.5,draw=chfourfig5c0,line width=0.3985pt,draw opacity=0.5] (260.9867,73.0380) .. controls (260.9867,72.3172) and (260.4024,71.7328) ..(259.6816,71.7328) .. controls (258.9608,71.7328) and (258.3765,72.3172) ..(258.3765,73.0380) .. controls (258.3765,73.7588) and (258.9608,74.3431) ..(259.6816,74.3431) .. controls (260.4024,74.3431) and (260.9867,73.7588) ..(260.9867,73.0380);
\path[fill=chfourfig5c0,fill opacity=0.5,draw=chfourfig5c0,line width=0.3985pt,draw opacity=0.5] (243.4024,90.0942) .. controls (243.4024,89.3734) and (242.8181,88.7891) ..(242.0973,88.7891) .. controls (241.3765,88.7891) and (240.7922,89.3734) ..(240.7922,90.0942) .. controls (240.7922,90.8150) and (241.3765,91.3994) ..(242.0973,91.3994) .. controls (242.8181,91.3994) and (243.4024,90.8150) ..(243.4024,90.0942);
\path[fill=chfourfig5c0,fill opacity=0.5,draw=chfourfig5c0,line width=0.3985pt,draw opacity=0.5] (248.4834,108.8143) .. controls (248.4834,108.0935) and (247.8991,107.5092) ..(247.1783,107.5092) .. controls (246.4575,107.5092) and (245.8732,108.0935) ..(245.8732,108.8143) .. controls (245.8732,109.5351) and (246.4575,110.1194) ..(247.1783,110.1194) .. controls (247.8991,110.1194) and (248.4834,109.5351) ..(248.4834,108.8143);
\path[fill=chfourfig5c0,fill opacity=0.5,draw=chfourfig5c0,line width=0.3985pt,draw opacity=0.5] (257.3204,102.9761) .. controls (257.3204,102.2553) and (256.7361,101.6710) ..(256.0153,101.6710) .. controls (255.2945,101.6710) and (254.7101,102.2553) ..(254.7101,102.9761) .. controls (254.7101,103.6969) and (255.2945,104.2812) ..(256.0153,104.2812) .. controls (256.7361,104.2812) and (257.3204,103.6969) ..(257.3204,102.9761);
\path[fill=chfourfig5c0,fill opacity=0.5,draw=chfourfig5c0,line width=0.3985pt,draw opacity=0.5] (294.5513,78.6470) .. controls (294.5513,77.9262) and (293.9669,77.3419) ..(293.2461,77.3419) .. controls (292.5253,77.3419) and (291.9410,77.9262) ..(291.9410,78.6470) .. controls (291.9410,79.3678) and (292.5253,79.9521) ..(293.2461,79.9521) .. controls (293.9669,79.9521) and (294.5513,79.3678) ..(294.5513,78.6470);
\path[fill=chfourfig5c0,fill opacity=0.5,draw=chfourfig5c0,line width=0.3985pt,draw opacity=0.5] (278.5411,71.0654) .. controls (278.5411,70.3445) and (277.9568,69.7602) ..(277.2360,69.7602) .. controls (276.5152,69.7602) and (275.9309,70.3445) ..(275.9309,71.0654) .. controls (275.9309,71.7861) and (276.5152,72.3705) ..(277.2360,72.3705) .. controls (277.9568,72.3705) and (278.5411,71.7861) ..(278.5411,71.0654);
\path[fill=chfourfig5c0,fill opacity=0.5,draw=chfourfig5c0,line width=0.3985pt,draw opacity=0.5] (245.1060,108.9936) .. controls (245.1060,108.2728) and (244.5217,107.6885) ..(243.8009,107.6885) .. controls (243.0801,107.6885) and (242.4958,108.2728) ..(242.4958,108.9936) .. controls (242.4958,109.7144) and (243.0801,110.2987) ..(243.8009,110.2987) .. controls (244.5217,110.2987) and (245.1060,109.7144) ..(245.1060,108.9936);
\path[fill=chfourfig5c0,fill opacity=0.5,draw=chfourfig5c0,line width=0.3985pt,draw opacity=0.5] (229.8929,81.9347) .. controls (229.8929,81.2139) and (229.3086,80.6296) ..(228.5878,80.6296) .. controls (227.8670,80.6296) and (227.2827,81.2139) ..(227.2827,81.9347) .. controls (227.2827,82.6555) and (227.8670,83.2398) ..(228.5878,83.2398) .. controls (229.3086,83.2398) and (229.8929,82.6555) ..(229.8929,81.9347);
\path[fill=chfourfig5c0,fill opacity=0.5,draw=chfourfig5c0,line width=0.3985pt,draw opacity=0.5] (236.4982,90.3931) .. controls (236.4982,89.6723) and (235.9139,89.0880) ..(235.1931,89.0880) .. controls (234.4723,89.0880) and (233.8880,89.6723) ..(233.8880,90.3931) .. controls (233.8880,91.1139) and (234.4723,91.6982) ..(235.1931,91.6982) .. controls (235.9139,91.6982) and (236.4982,91.1139) ..(236.4982,90.3931);
\path[fill=chfourfig5c0,fill opacity=0.5,draw=chfourfig5c0,line width=0.3985pt,draw opacity=0.5] (295.8165,55.5234) .. controls (295.8165,54.8026) and (295.2322,54.2183) ..(294.5114,54.2183) .. controls (293.7906,54.2183) and (293.2063,54.8026) ..(293.2063,55.5234) .. controls (293.2063,56.2442) and (293.7906,56.8285) ..(294.5114,56.8285) .. controls (295.2322,56.8285) and (295.8165,56.2442) ..(295.8165,55.5234);
\path[fill=chfourfig5c0,fill opacity=0.5,draw=chfourfig5c0,line width=0.3985pt,draw opacity=0.5] (317.8641,59.7078) .. controls (317.8641,58.9870) and (317.2798,58.4027) ..(316.5590,58.4027) .. controls (315.8382,58.4027) and (315.2539,58.9870) ..(315.2539,59.7078) .. controls (315.2539,60.4286) and (315.8382,61.0129) ..(316.5590,61.0129) .. controls (317.2798,61.0129) and (317.8641,60.4286) ..(317.8641,59.7078);
\path[fill=chfourfig5c0,fill opacity=0.5,draw=chfourfig5c0,line width=0.3985pt,draw opacity=0.5] (242.4559,79.7429) .. controls (242.4559,79.0221) and (241.8716,78.4378) ..(241.1508,78.4378) .. controls (240.4300,78.4378) and (239.8457,79.0221) ..(239.8457,79.7429) .. controls (239.8457,80.4637) and (240.4300,81.0480) ..(241.1508,81.0480) .. controls (241.8716,81.0480) and (242.4559,80.4637) ..(242.4559,79.7429);
\path[fill=chfourfig5c0,fill opacity=0.5,draw=chfourfig5c0,line width=0.3985pt,draw opacity=0.5] (307.8815,56.0216) .. controls (307.8815,55.3008) and (307.2971,54.7164) ..(306.5763,54.7164) .. controls (305.8555,54.7164) and (305.2712,55.3008) ..(305.2712,56.0216) .. controls (305.2712,56.7424) and (305.8555,57.3267) ..(306.5763,57.3267) .. controls (307.2971,57.3267) and (307.8815,56.7424) ..(307.8815,56.0216);
\path[fill=chfourfig5c0,fill opacity=0.5,draw=chfourfig5c0,line width=0.3985pt,draw opacity=0.5] (243.3825,107.0907) .. controls (243.3825,106.3699) and (242.7982,105.7856) ..(242.0774,105.7856) .. controls (241.3566,105.7856) and (240.7722,106.3699) ..(240.7722,107.0907) .. controls (240.7722,107.8115) and (241.3566,108.3958) ..(242.0774,108.3958) .. controls (242.7982,108.3958) and (243.3825,107.8115) ..(243.3825,107.0907);
\path[fill=chfourfig5c0,fill opacity=0.5,draw=chfourfig5c0,line width=0.3985pt,draw opacity=0.5] (297.4903,51.2095) .. controls (297.4903,50.4887) and (296.9060,49.9044) ..(296.1852,49.9044) .. controls (295.4644,49.9044) and (294.8800,50.4887) ..(294.8800,51.2095) .. controls (294.8800,51.9303) and (295.4644,52.5147) ..(296.1852,52.5147) .. controls (296.9060,52.5147) and (297.4903,51.9303) ..(297.4903,51.2095);
\path[fill=chfourfig5c0,fill opacity=0.5,draw=chfourfig5c0,line width=0.3985pt,draw opacity=0.5] (324.6787,59.8273) .. controls (324.6787,59.1065) and (324.0944,58.5222) ..(323.3735,58.5222) .. controls (322.6528,58.5222) and (322.0685,59.1065) ..(322.0685,59.8273) .. controls (322.0685,60.5481) and (322.6528,61.1325) ..(323.3735,61.1325) .. controls (324.0944,61.1325) and (324.6787,60.5481) ..(324.6787,59.8273);
\path[fill=chfourfig5c0,fill opacity=0.5,draw=chfourfig5c0,line width=0.3985pt,draw opacity=0.5] (331.4235,34.3127) .. controls (331.4235,33.5919) and (330.8391,33.0076) ..(330.1183,33.0076) .. controls (329.3976,33.0076) and (328.8132,33.5919) ..(328.8132,34.3127) .. controls (328.8132,35.0335) and (329.3976,35.6178) ..(330.1183,35.6178) .. controls (330.8391,35.6178) and (331.4235,35.0335) ..(331.4235,34.3127);
\path[fill=chfourfig5c0,fill opacity=0.5,draw=chfourfig5c0,line width=0.3985pt,draw opacity=0.5] (254.7500,92.9934) .. controls (254.7500,92.2726) and (254.1657,91.6883) ..(253.4449,91.6883) .. controls (252.7241,91.6883) and (252.1398,92.2726) ..(252.1398,92.9934) .. controls (252.1398,93.7142) and (252.7241,94.2985) ..(253.4449,94.2985) .. controls (254.1657,94.2985) and (254.7500,93.7142) ..(254.7500,92.9934);
\path[fill=chfourfig5c0,fill opacity=0.5,draw=chfourfig5c0,line width=0.3985pt,draw opacity=0.5] (309.9437,40.4796) .. controls (309.9437,39.7589) and (309.3594,39.1745) ..(308.6386,39.1745) .. controls (307.9178,39.1745) and (307.3335,39.7589) ..(307.3335,40.4796) .. controls (307.3335,41.2005) and (307.9178,41.7848) ..(308.6386,41.7848) .. controls (309.3594,41.7848) and (309.9437,41.2005) ..(309.9437,40.4796);
\path[fill=chfourfig5c0,fill opacity=0.5,draw=chfourfig5c0,line width=0.3985pt,draw opacity=0.5] (331.7821,47.1248) .. controls (331.7821,46.4040) and (331.1978,45.8197) ..(330.4770,45.8197) .. controls (329.7562,45.8197) and (329.1719,46.4040) ..(329.1719,47.1248) .. controls (329.1719,47.8456) and (329.7562,48.4299) ..(330.4770,48.4299) .. controls (331.1978,48.4299) and (331.7821,47.8456) ..(331.7821,47.1248);
\path[fill=chfourfig5c0,fill opacity=0.5,draw=chfourfig5c0,line width=0.3985pt,draw opacity=0.5] (300.5090,67.9669) .. controls (300.5090,67.2461) and (299.9247,66.6618) ..(299.2039,66.6618) .. controls (298.4831,66.6618) and (297.8987,67.2461) ..(297.8987,67.9669) .. controls (297.8987,68.6877) and (298.4831,69.2720) ..(299.2039,69.2720) .. controls (299.9247,69.2720) and (300.5090,68.6877) ..(300.5090,67.9669);
\path[fill=chfourfig5c0,fill opacity=0.5,draw=chfourfig5c0,line width=0.3985pt,draw opacity=0.5] (333.4359,41.1671) .. controls (333.4359,40.4463) and (332.8516,39.8620) ..(332.1308,39.8620) .. controls (331.4100,39.8620) and (330.8257,40.4463) ..(330.8257,41.1671) .. controls (330.8257,41.8879) and (331.4100,42.4722) ..(332.1308,42.4722) .. controls (332.8516,42.4722) and (333.4359,41.8879) ..(333.4359,41.1671);
\path[fill=chfourfig5c0,fill opacity=0.5,draw=chfourfig5c0,line width=0.3985pt,draw opacity=0.5] (299.0146,55.4637) .. controls (299.0146,54.7428) and (298.4303,54.1585) ..(297.7095,54.1585) .. controls (296.9887,54.1585) and (296.4044,54.7428) ..(296.4044,55.4637) .. controls (296.4044,56.1844) and (296.9887,56.7688) ..(297.7095,56.7688) .. controls (298.4303,56.7688) and (299.0146,56.1844) ..(299.0146,55.4637);
\path[fill=chfourfig5c0,fill opacity=0.5,draw=chfourfig5c0,line width=0.3985pt,draw opacity=0.5] (244.0002,84.6048) .. controls (244.0002,83.8839) and (243.4159,83.2996) ..(242.6951,83.2996) .. controls (241.9742,83.2996) and (241.3900,83.8839) ..(241.3900,84.6048) .. controls (241.3900,85.3255) and (241.9742,85.9099) ..(242.6951,85.9099) .. controls (243.4159,85.9099) and (244.0002,85.3255) ..(244.0002,84.6048);
\path[fill=chfourfig5c0,fill opacity=0.5,draw=chfourfig5c0,line width=0.3985pt,draw opacity=0.5] (251.9505,70.7963) .. controls (251.9505,70.0755) and (251.3661,69.4912) ..(250.6453,69.4912) .. controls (249.9246,69.4912) and (249.3402,70.0755) ..(249.3402,70.7963) .. controls (249.3402,71.5172) and (249.9246,72.1015) ..(250.6453,72.1015) .. controls (251.3661,72.1015) and (251.9505,71.5172) ..(251.9505,70.7963);
\path[fill=chfourfig5c0,fill opacity=0.5,draw=chfourfig5c0,line width=0.3985pt,draw opacity=0.5] (235.3127,98.3733) .. controls (235.3127,97.6525) and (234.7283,97.0682) ..(234.0075,97.0682) .. controls (233.2867,97.0682) and (232.7024,97.6525) ..(232.7024,98.3733) .. controls (232.7024,99.0941) and (233.2867,99.6784) ..(234.0075,99.6784) .. controls (234.7283,99.6784) and (235.3127,99.0941) ..(235.3127,98.3733);
\path[fill=chfourfig5c0,fill opacity=0.5,draw=chfourfig5c0,line width=0.3985pt,draw opacity=0.5] (272.8723,71.0853) .. controls (272.8723,70.3645) and (272.2880,69.7801) ..(271.5672,69.7801) .. controls (270.8463,69.7801) and (270.2620,70.3645) ..(270.2620,71.0853) .. controls (270.2620,71.8061) and (270.8463,72.3904) ..(271.5672,72.3904) .. controls (272.2880,72.3904) and (272.8723,71.8061) ..(272.8723,71.0853);
\path[fill=chfourfig5c0,fill opacity=0.5,draw=chfourfig5c0,line width=0.3985pt,draw opacity=0.5] (286.2622,87.1253) .. controls (286.2622,86.4045) and (285.6779,85.8202) ..(284.9571,85.8202) .. controls (284.2363,85.8202) and (283.6520,86.4045) ..(283.6520,87.1253) .. controls (283.6520,87.8461) and (284.2363,88.4305) ..(284.9571,88.4305) .. controls (285.6779,88.4305) and (286.2622,87.8461) ..(286.2622,87.1253);
\path[fill=chfourfig5c0,fill opacity=0.5,draw=chfourfig5c0,line width=0.3985pt,draw opacity=0.5] (247.0687,73.6557) .. controls (247.0687,72.9349) and (246.4844,72.3505) ..(245.7636,72.3505) .. controls (245.0428,72.3505) and (244.4585,72.9349) ..(244.4585,73.6557) .. controls (244.4585,74.3765) and (245.0428,74.9608) ..(245.7636,74.9608) .. controls (246.4844,74.9608) and (247.0687,74.3765) ..(247.0687,73.6557);
\path[fill=chfourfig5c0,fill opacity=0.5,draw=chfourfig5c0,line width=0.3985pt,draw opacity=0.5] (299.1541,60.1362) .. controls (299.1541,59.4154) and (298.5697,58.8311) ..(297.8490,58.8311) .. controls (297.1281,58.8311) and (296.5438,59.4154) ..(296.5438,60.1362) .. controls (296.5438,60.8570) and (297.1281,61.4413) ..(297.8490,61.4413) .. controls (298.5697,61.4413) and (299.1541,60.8570) ..(299.1541,60.1362);
\path[fill=chfourfig5c0,fill opacity=0.5,draw=chfourfig5c0,line width=0.3985pt,draw opacity=0.5] (321.6301,41.3763) .. controls (321.6301,40.6555) and (321.0457,40.0712) ..(320.3250,40.0712) .. controls (319.6042,40.0712) and (319.0198,40.6555) ..(319.0198,41.3763) .. controls (319.0198,42.0971) and (319.6042,42.6814) ..(320.3250,42.6814) .. controls (321.0457,42.6814) and (321.6301,42.0971) ..(321.6301,41.3763);
\path[fill=chfourfig5c0,fill opacity=0.5,draw=chfourfig5c0,line width=0.3985pt,draw opacity=0.5] (301.4355,43.5581) .. controls (301.4355,42.8373) and (300.8512,42.2530) ..(300.1304,42.2530) .. controls (299.4096,42.2530) and (298.8253,42.8373) ..(298.8253,43.5581) .. controls (298.8253,44.2789) and (299.4096,44.8633) ..(300.1304,44.8633) .. controls (300.8512,44.8633) and (301.4355,44.2789) ..(301.4355,43.5581);
\path[fill=chfourfig5c0,fill opacity=0.5,draw=chfourfig5c0,line width=0.3985pt,draw opacity=0.5] (230.9888,91.5787) .. controls (230.9888,90.8579) and (230.4045,90.2736) ..(229.6837,90.2736) .. controls (228.9629,90.2736) and (228.3786,90.8579) ..(228.3786,91.5787) .. controls (228.3786,92.2995) and (228.9629,92.8838) ..(229.6837,92.8838) .. controls (230.4045,92.8838) and (230.9888,92.2995) ..(230.9888,91.5787);
\path[fill=chfourfig5c0,fill opacity=0.5,draw=chfourfig5c0,line width=0.3985pt,draw opacity=0.5] (268.1699,78.7466) .. controls (268.1699,78.0258) and (267.5855,77.4415) ..(266.8647,77.4415) .. controls (266.1440,77.4415) and (265.5596,78.0258) ..(265.5596,78.7466) .. controls (265.5596,79.4674) and (266.1440,80.0518) ..(266.8647,80.0518) .. controls (267.5855,80.0518) and (268.1699,79.4674) ..(268.1699,78.7466);
\path[fill=chfourfig5c0,fill opacity=0.5,draw=chfourfig5c0,line width=0.3985pt,draw opacity=0.5] (275.0143,61.7402) .. controls (275.0143,61.0194) and (274.4299,60.4351) ..(273.7092,60.4351) .. controls (272.9883,60.4351) and (272.4040,61.0194) ..(272.4040,61.7402) .. controls (272.4040,62.4610) and (272.9883,63.0453) ..(273.7092,63.0453) .. controls (274.4299,63.0453) and (275.0143,62.4610) ..(275.0143,61.7402);
\path[fill=chfourfig5c0,fill opacity=0.5,draw=chfourfig5c0,line width=0.3985pt,draw opacity=0.5] (307.9213,74.4328) .. controls (307.9213,73.7119) and (307.3370,73.1276) ..(306.6162,73.1276) .. controls (305.8954,73.1276) and (305.3111,73.7119) ..(305.3111,74.4328) .. controls (305.3111,75.1535) and (305.8954,75.7379) ..(306.6162,75.7379) .. controls (307.3370,75.7379) and (307.9213,75.1535) ..(307.9213,74.4328);
\path[fill=chfourfig5c0,fill opacity=0.5,draw=chfourfig5c0,line width=0.3985pt,draw opacity=0.5] (313.6898,50.7911) .. controls (313.6898,50.0703) and (313.1054,49.4860) ..(312.3846,49.4860) .. controls (311.6638,49.4860) and (311.0795,50.0703) ..(311.0795,50.7911) .. controls (311.0795,51.5119) and (311.6638,52.0962) ..(312.3846,52.0962) .. controls (313.1054,52.0962) and (313.6898,51.5119) ..(313.6898,50.7911);
\path[fill=chfourfig5c0,fill opacity=0.5,draw=chfourfig5c0,line width=0.3985pt,draw opacity=0.5] (246.7200,106.8616) .. controls (246.7200,106.1408) and (246.1357,105.5564) ..(245.4149,105.5564) .. controls (244.6941,105.5564) and (244.1098,106.1408) ..(244.1098,106.8616) .. controls (244.1098,107.5824) and (244.6941,108.1667) ..(245.4149,108.1667) .. controls (246.1357,108.1667) and (246.7200,107.5824) ..(246.7200,106.8616);
\path[fill=chfourfig5c0,fill opacity=0.5,draw=chfourfig5c0,line width=0.3985pt,draw opacity=0.5] (282.0380,80.4204) .. controls (282.0380,79.6996) and (281.4537,79.1153) ..(280.7329,79.1153) .. controls (280.0121,79.1153) and (279.4278,79.6996) ..(279.4278,80.4204) .. controls (279.4278,81.1412) and (280.0121,81.7255) ..(280.7329,81.7255) .. controls (281.4537,81.7255) and (282.0380,81.1412) ..(282.0380,80.4204);
\path[fill=chfourfig5c0,fill opacity=0.5,draw=chfourfig5c0,line width=0.3985pt,draw opacity=0.5] (321.4607,36.7038) .. controls (321.4607,35.9829) and (320.8764,35.3986) ..(320.1556,35.3986) .. controls (319.4348,35.3986) and (318.8505,35.9829) ..(318.8505,36.7038) .. controls (318.8505,37.4246) and (319.4348,38.0089) ..(320.1556,38.0089) .. controls (320.8764,38.0089) and (321.4607,37.4246) ..(321.4607,36.7038);
\path[fill=chfourfig5c0,fill opacity=0.5,draw=chfourfig5c0,line width=0.3985pt,draw opacity=0.5] (232.3736,89.8452) .. controls (232.3736,89.1244) and (231.7893,88.5400) ..(231.0685,88.5400) .. controls (230.3477,88.5400) and (229.7634,89.1244) ..(229.7634,89.8452) .. controls (229.7634,90.5660) and (230.3477,91.1503) ..(231.0685,91.1503) .. controls (231.7893,91.1503) and (232.3736,90.5660) ..(232.3736,89.8452);
\path[fill=chfourfig5c0,fill opacity=0.5,draw=chfourfig5c0,line width=0.3985pt,draw opacity=0.5] (278.0330,58.4027) .. controls (278.0330,57.6819) and (277.4487,57.0975) ..(276.7279,57.0975) .. controls (276.0071,57.0975) and (275.4228,57.6819) ..(275.4228,58.4027) .. controls (275.4228,59.1235) and (276.0071,59.7078) ..(276.7279,59.7078) .. controls (277.4487,59.7078) and (278.0330,59.1235) ..(278.0330,58.4027);
\path[fill=chfourfig5c0,fill opacity=0.5,draw=chfourfig5c0,line width=0.3985pt,draw opacity=0.5] (303.5775,67.2297) .. controls (303.5775,66.5089) and (302.9932,65.9246) ..(302.2724,65.9246) .. controls (301.5516,65.9246) and (300.9673,66.5089) ..(300.9673,67.2297) .. controls (300.9673,67.9505) and (301.5516,68.5348) ..(302.2724,68.5348) .. controls (302.9932,68.5348) and (303.5775,67.9505) ..(303.5775,67.2297);
\path[fill=chfourfig5c0,fill opacity=0.5,draw=chfourfig5c0,line width=0.3985pt,draw opacity=0.5] (273.1712,88.7493) .. controls (273.1712,88.0285) and (272.5869,87.4441) ..(271.8661,87.4441) .. controls (271.1452,87.4441) and (270.5609,88.0285) ..(270.5609,88.7493) .. controls (270.5609,89.4701) and (271.1452,90.0544) ..(271.8661,90.0544) .. controls (272.5869,90.0544) and (273.1712,89.4701) ..(273.1712,88.7493);
\path[fill=chfourfig5c0,fill opacity=0.5,draw=chfourfig5c0,line width=0.3985pt,draw opacity=0.5] (248.7624,85.8999) .. controls (248.7624,85.1791) and (248.1781,84.5948) ..(247.4573,84.5948) .. controls (246.7365,84.5948) and (246.1522,85.1791) ..(246.1522,85.8999) .. controls (246.1522,86.6207) and (246.7365,87.2050) ..(247.4573,87.2050) .. controls (248.1781,87.2050) and (248.7624,86.6207) ..(248.7624,85.8999);
\path[fill=chfourfig5c0,fill opacity=0.5,draw=chfourfig5c0,line width=0.3985pt,draw opacity=0.5] (233.1208,89.4566) .. controls (233.1208,88.7358) and (232.5365,88.1515) ..(231.8157,88.1515) .. controls (231.0949,88.1515) and (230.5106,88.7358) ..(230.5106,89.4566) .. controls (230.5106,90.1774) and (231.0949,90.7617) ..(231.8157,90.7617) .. controls (232.5365,90.7617) and (233.1208,90.1774) ..(233.1208,89.4566);
\path[fill=chfourfig5c0,fill opacity=0.5,draw=chfourfig5c0,line width=0.3985pt,draw opacity=0.5] (236.7573,110.0098) .. controls (236.7573,109.2890) and (236.1729,108.7047) ..(235.4521,108.7047) .. controls (234.7313,108.7047) and (234.1470,109.2890) ..(234.1470,110.0098) .. controls (234.1470,110.7306) and (234.7313,111.3149) ..(235.4521,111.3149) .. controls (236.1729,111.3149) and (236.7573,110.7306) ..(236.7573,110.0098);
\path[fill=chfourfig5c0,fill opacity=0.5,draw=chfourfig5c0,line width=0.3985pt,draw opacity=0.5] (255.1187,69.4215) .. controls (255.1187,68.7007) and (254.5343,68.1164) ..(253.8135,68.1164) .. controls (253.0927,68.1164) and (252.5084,68.7007) ..(252.5084,69.4215) .. controls (252.5084,70.1423) and (253.0927,70.7266) ..(253.8135,70.7266) .. controls (254.5343,70.7266) and (255.1187,70.1423) ..(255.1187,69.4215);
\path[fill=chfourfig5c0,fill opacity=0.5,draw=chfourfig5c0,line width=0.3985pt,draw opacity=0.5] (232.0349,88.1116) .. controls (232.0349,87.3908) and (231.4506,86.8065) ..(230.7298,86.8065) .. controls (230.0090,86.8065) and (229.4247,87.3908) ..(229.4247,88.1116) .. controls (229.4247,88.8324) and (230.0090,89.4168) ..(230.7298,89.4168) .. controls (231.4506,89.4168) and (232.0349,88.8324) ..(232.0349,88.1116);
\path[fill=chfourfig5c0,fill opacity=0.5,draw=chfourfig5c0,line width=0.3985pt,draw opacity=0.5] (323.9215,46.1186) .. controls (323.9215,45.3978) and (323.3372,44.8134) ..(322.6164,44.8134) .. controls (321.8956,44.8134) and (321.3113,45.3978) ..(321.3113,46.1186) .. controls (321.3113,46.8394) and (321.8956,47.4237) ..(322.6164,47.4237) .. controls (323.3372,47.4237) and (323.9215,46.8394) ..(323.9215,46.1186);
\path[fill=chfourfig5c0,fill opacity=0.5,draw=chfourfig5c0,line width=0.3985pt,draw opacity=0.5] (335.6178,43.9666) .. controls (335.6178,43.2458) and (335.0335,42.6615) ..(334.3127,42.6615) .. controls (333.5919,42.6615) and (333.0076,43.2458) ..(333.0076,43.9666) .. controls (333.0076,44.6874) and (333.5919,45.2718) ..(334.3127,45.2718) .. controls (335.0335,45.2718) and (335.6178,44.6874) ..(335.6178,43.9666);
\path[fill=chfourfig5c0,fill opacity=0.5,draw=chfourfig5c0,line width=0.3985pt,draw opacity=0.5] (243.5817,97.6659) .. controls (243.5817,96.9451) and (242.9974,96.3608) ..(242.2766,96.3608) .. controls (241.5558,96.3608) and (240.9715,96.9451) ..(240.9715,97.6659) .. controls (240.9715,98.3867) and (241.5558,98.9711) ..(242.2766,98.9711) .. controls (242.9974,98.9711) and (243.5817,98.3867) ..(243.5817,97.6659);
\path[fill=chfourfig5c0,fill opacity=0.5,draw=chfourfig5c0,line width=0.3985pt,draw opacity=0.5] (241.8582,111.6736) .. controls (241.8582,110.9528) and (241.2739,110.3685) ..(240.5530,110.3685) .. controls (239.8323,110.3685) and (239.2479,110.9528) ..(239.2479,111.6736) .. controls (239.2479,112.3944) and (239.8323,112.9787) ..(240.5530,112.9787) .. controls (241.2739,112.9787) and (241.8582,112.3944) ..(241.8582,111.6736);
\path[fill=chfourfig5c0,fill opacity=0.5,draw=chfourfig5c0,line width=0.3985pt,draw opacity=0.5] (301.7245,61.3018) .. controls (301.7245,60.5810) and (301.1402,59.9967) ..(300.4193,59.9967) .. controls (299.6985,59.9967) and (299.1142,60.5810) ..(299.1142,61.3018) .. controls (299.1142,62.0226) and (299.6985,62.6069) ..(300.4193,62.6069) .. controls (301.1402,62.6069) and (301.7245,62.0226) ..(301.7245,61.3018);
\path[fill=chfourfig5c0,fill opacity=0.5,draw=chfourfig5c0,line width=0.3985pt,draw opacity=0.5] (296.2350,53.6803) .. controls (296.2350,52.9595) and (295.6506,52.3752) ..(294.9299,52.3752) .. controls (294.2090,52.3752) and (293.6247,52.9595) ..(293.6247,53.6803) .. controls (293.6247,54.4011) and (294.2090,54.9854) ..(294.9299,54.9854) .. controls (295.6506,54.9854) and (296.2350,54.4011) ..(296.2350,53.6803);
\path[fill=chfourfig5c0,fill opacity=0.5,draw=chfourfig5c0,line width=0.3985pt,draw opacity=0.5] (299.1840,52.5944) .. controls (299.1840,51.8736) and (298.5996,51.2892) ..(297.8788,51.2892) .. controls (297.1581,51.2892) and (296.5737,51.8736) ..(296.5737,52.5944) .. controls (296.5737,53.3152) and (297.1581,53.8995) ..(297.8788,53.8995) .. controls (298.5996,53.8995) and (299.1840,53.3152) ..(299.1840,52.5944);
\path[fill=chfourfig5c0,fill opacity=0.5,draw=chfourfig5c0,line width=0.3985pt,draw opacity=0.5] (301.7245,48.3901) .. controls (301.7245,47.6693) and (301.1402,47.0850) ..(300.4193,47.0850) .. controls (299.6985,47.0850) and (299.1142,47.6693) ..(299.1142,48.3901) .. controls (299.1142,49.1109) and (299.6985,49.6952) ..(300.4193,49.6952) .. controls (301.1402,49.6952) and (301.7245,49.1109) ..(301.7245,48.3901);
\path[fill=chfourfig5c0,fill opacity=0.5,draw=chfourfig5c0,line width=0.3985pt,draw opacity=0.5] (327.1992,38.8258) .. controls (327.1992,38.1050) and (326.6149,37.5207) ..(325.8941,37.5207) .. controls (325.1733,37.5207) and (324.5890,38.1050) ..(324.5890,38.8258) .. controls (324.5890,39.5466) and (325.1733,40.1310) ..(325.8941,40.1310) .. controls (326.6149,40.1310) and (327.1992,39.5466) ..(327.1992,38.8258);
\path[fill=chfourfig5c0,fill opacity=0.5,draw=chfourfig5c0,line width=0.3985pt,draw opacity=0.5] (262.9793,76.9533) .. controls (262.9793,76.2325) and (262.3950,75.6482) ..(261.6741,75.6482) .. controls (260.9533,75.6482) and (260.3690,76.2325) ..(260.3690,76.9533) .. controls (260.3690,77.6741) and (260.9533,78.2585) ..(261.6741,78.2585) .. controls (262.3950,78.2585) and (262.9793,77.6741) ..(262.9793,76.9533);
\path[fill=chfourfig5c0,fill opacity=0.5,draw=chfourfig5c0,line width=0.3985pt,draw opacity=0.5] (288.6433,81.4465) .. controls (288.6433,80.7258) and (288.0590,80.1414) ..(287.3382,80.1414) .. controls (286.6174,80.1414) and (286.0331,80.7258) ..(286.0331,81.4465) .. controls (286.0331,82.1674) and (286.6174,82.7517) ..(287.3382,82.7517) .. controls (288.0590,82.7517) and (288.6433,82.1674) ..(288.6433,81.4465);
\path[fill=chfourfig5c0,fill opacity=0.5,draw=chfourfig5c0,line width=0.3985pt,draw opacity=0.5] (305.5402,63.3840) .. controls (305.5402,62.6632) and (304.9559,62.0789) ..(304.2351,62.0789) .. controls (303.5143,62.0789) and (302.9299,62.6632) ..(302.9299,63.3840) .. controls (302.9299,64.1048) and (303.5143,64.6892) ..(304.2351,64.6892) .. controls (304.9559,64.6892) and (305.5402,64.1048) ..(305.5402,63.3840);
\path[fill=chfourfig5c0,fill opacity=0.5,draw=chfourfig5c0,line width=0.3985pt,draw opacity=0.5] (305.7494,59.0403) .. controls (305.7494,58.3195) and (305.1651,57.7352) ..(304.4443,57.7352) .. controls (303.7235,57.7352) and (303.1392,58.3195) ..(303.1392,59.0403) .. controls (303.1392,59.7611) and (303.7235,60.3454) ..(304.4443,60.3454) .. controls (305.1651,60.3454) and (305.7494,59.7611) ..(305.7494,59.0403);
\path[fill=chfourfig5c0,fill opacity=0.5,draw=chfourfig5c0,line width=0.3985pt,draw opacity=0.5] (238.5904,93.6410) .. controls (238.5904,92.9202) and (238.0061,92.3358) ..(237.2853,92.3358) .. controls (236.5645,92.3358) and (235.9802,92.9202) ..(235.9802,93.6410) .. controls (235.9802,94.3618) and (236.5645,94.9461) ..(237.2853,94.9461) .. controls (238.0061,94.9461) and (238.5904,94.3618) ..(238.5904,93.6410);
\path[fill=chfourfig5c0,fill opacity=0.5,draw=chfourfig5c0,line width=0.3985pt,draw opacity=0.5] (288.3544,73.0579) .. controls (288.3544,72.3371) and (287.7701,71.7528) ..(287.0493,71.7528) .. controls (286.3285,71.7528) and (285.7442,72.3371) ..(285.7442,73.0579) .. controls (285.7442,73.7787) and (286.3285,74.3630) ..(287.0493,74.3630) .. controls (287.7701,74.3630) and (288.3544,73.7787) ..(288.3544,73.0579);
\path[fill=chfourfig5c0,fill opacity=0.5,draw=chfourfig5c0,line width=0.3985pt,draw opacity=0.5] (279.8362,59.5285) .. controls (279.8362,58.8077) and (279.2519,58.2233) ..(278.5312,58.2233) .. controls (277.8103,58.2233) and (277.2260,58.8077) ..(277.2260,59.5285) .. controls (277.2260,60.2493) and (277.8103,60.8336) ..(278.5312,60.8336) .. controls (279.2519,60.8336) and (279.8362,60.2493) ..(279.8362,59.5285);
\path[fill=chfourfig5c0,fill opacity=0.5,draw=chfourfig5c0,line width=0.3985pt,draw opacity=0.5] (324.9576,29.5704) .. controls (324.9576,28.8496) and (324.3733,28.2653) ..(323.6525,28.2653) .. controls (322.9317,28.2653) and (322.3474,28.8496) ..(322.3474,29.5704) .. controls (322.3474,30.2912) and (322.9317,30.8755) ..(323.6525,30.8755) .. controls (324.3733,30.8755) and (324.9576,30.2912) ..(324.9576,29.5704);
\path[fill=chfourfig5c0,fill opacity=0.5,draw=chfourfig5c0,line width=0.3985pt,draw opacity=0.5] (242.5755,111.3448) .. controls (242.5755,110.6240) and (241.9912,110.0397) ..(241.2704,110.0397) .. controls (240.5496,110.0397) and (239.9653,110.6240) ..(239.9653,111.3448) .. controls (239.9653,112.0656) and (240.5496,112.6499) ..(241.2704,112.6499) .. controls (241.9912,112.6499) and (242.5755,112.0656) ..(242.5755,111.3448);
\path[fill=chfourfig5c0,fill opacity=0.5,draw=chfourfig5c0,line width=0.3985pt,draw opacity=0.5] (316.2402,57.6355) .. controls (316.2402,56.9147) and (315.6559,56.3304) ..(314.9351,56.3304) .. controls (314.2143,56.3304) and (313.6300,56.9147) ..(313.6300,57.6355) .. controls (313.6300,58.3563) and (314.2143,58.9406) ..(314.9351,58.9406) .. controls (315.6559,58.9406) and (316.2402,58.3563) ..(316.2402,57.6355);
\path[fill=chfourfig5c0,fill opacity=0.5,draw=chfourfig5c0,line width=0.3985pt,draw opacity=0.5] (253.9630,98.2836) .. controls (253.9630,97.5628) and (253.3786,96.9785) ..(252.6578,96.9785) .. controls (251.9370,96.9785) and (251.3527,97.5628) ..(251.3527,98.2836) .. controls (251.3527,99.0044) and (251.9370,99.5887) ..(252.6578,99.5887) .. controls (253.3786,99.5887) and (253.9630,99.0044) ..(253.9630,98.2836);
\path[fill=chfourfig5c0,fill opacity=0.5,draw=chfourfig5c0,line width=0.3985pt,draw opacity=0.5] (285.9235,80.2410) .. controls (285.9235,79.5202) and (285.3392,78.9359) ..(284.6184,78.9359) .. controls (283.8976,78.9359) and (283.3133,79.5202) ..(283.3133,80.2410) .. controls (283.3133,80.9619) and (283.8976,81.5462) ..(284.6184,81.5462) .. controls (285.3392,81.5462) and (285.9235,80.9619) ..(285.9235,80.2410);
\path[fill=chfourfig5c0,fill opacity=0.5,draw=chfourfig5c0,line width=0.3985pt,draw opacity=0.5] (238.4808,106.7022) .. controls (238.4808,105.9814) and (237.8965,105.3970) ..(237.1757,105.3970) .. controls (236.4549,105.3970) and (235.8705,105.9814) ..(235.8705,106.7022) .. controls (235.8705,107.4230) and (236.4549,108.0073) ..(237.1757,108.0073) .. controls (237.8965,108.0073) and (238.4808,107.4230) ..(238.4808,106.7022);
\path[fill=chfourfig5c0,fill opacity=0.5,draw=chfourfig5c0,line width=0.3985pt,draw opacity=0.5] (266.6356,60.4251) .. controls (266.6356,59.7043) and (266.0513,59.1200) ..(265.3305,59.1200) .. controls (264.6097,59.1200) and (264.0254,59.7043) ..(264.0254,60.4251) .. controls (264.0254,61.1459) and (264.6097,61.7302) ..(265.3305,61.7302) .. controls (266.0513,61.7302) and (266.6356,61.1459) ..(266.6356,60.4251);
\path[fill=chfourfig5c0,fill opacity=0.5,draw=chfourfig5c0,line width=0.3985pt,draw opacity=0.5] (290.2175,58.5023) .. controls (290.2175,57.7815) and (289.6331,57.1972) ..(288.9123,57.1972) .. controls (288.1916,57.1972) and (287.6072,57.7815) ..(287.6072,58.5023) .. controls (287.6072,59.2231) and (288.1916,59.8074) ..(288.9123,59.8074) .. controls (289.6331,59.8074) and (290.2175,59.2231) ..(290.2175,58.5023);
\path[fill=chfourfig5c0,fill opacity=0.5,draw=chfourfig5c0,line width=0.3985pt,draw opacity=0.5] (265.9581,77.1626) .. controls (265.9581,76.4418) and (265.3738,75.8574) ..(264.6530,75.8574) .. controls (263.9322,75.8574) and (263.3479,76.4418) ..(263.3479,77.1626) .. controls (263.3479,77.8834) and (263.9322,78.4677) ..(264.6530,78.4677) .. controls (265.3738,78.4677) and (265.9581,77.8834) ..(265.9581,77.1626);
\path[fill=chfourfig5c0,fill opacity=0.5,draw=chfourfig5c0,line width=0.3985pt,draw opacity=0.5] (267.5422,83.5686) .. controls (267.5422,82.8478) and (266.9579,82.2635) ..(266.2371,82.2635) .. controls (265.5163,82.2635) and (264.9319,82.8478) ..(264.9319,83.5686) .. controls (264.9319,84.2894) and (265.5163,84.8737) ..(266.2371,84.8737) .. controls (266.9579,84.8737) and (267.5422,84.2894) ..(267.5422,83.5686);
\path[fill=chfourfig5c0,fill opacity=0.5,draw=chfourfig5c0,line width=0.3985pt,draw opacity=0.5] (263.7962,95.9324) .. controls (263.7962,95.2116) and (263.2119,94.6273) ..(262.4911,94.6273) .. controls (261.7703,94.6273) and (261.1859,95.2116) ..(261.1859,95.9324) .. controls (261.1859,96.6532) and (261.7703,97.2375) ..(262.4911,97.2375) .. controls (263.2119,97.2375) and (263.7962,96.6532) ..(263.7962,95.9324);
\path[fill=chfourfig5c0,fill opacity=0.5,draw=chfourfig5c0,line width=0.3985pt,draw opacity=0.5] (322.2079,42.8109) .. controls (322.2079,42.0901) and (321.6236,41.5058) ..(320.9028,41.5058) .. controls (320.1820,41.5058) and (319.5977,42.0901) ..(319.5977,42.8109) .. controls (319.5977,43.5317) and (320.1820,44.1161) ..(320.9028,44.1161) .. controls (321.6236,44.1161) and (322.2079,43.5317) ..(322.2079,42.8109);
\path[fill=chfourfig5c0,fill opacity=0.5,draw=chfourfig5c0,line width=0.3985pt,draw opacity=0.5] (287.9460,58.1835) .. controls (287.9460,57.4627) and (287.3616,56.8784) ..(286.6408,56.8784) .. controls (285.9200,56.8784) and (285.3357,57.4627) ..(285.3357,58.1835) .. controls (285.3357,58.9043) and (285.9200,59.4886) ..(286.6408,59.4886) .. controls (287.3616,59.4886) and (287.9460,58.9043) ..(287.9460,58.1835);
\path[fill=chfourfig5c0,fill opacity=0.5,draw=chfourfig5c0,line width=0.3985pt,draw opacity=0.5] (274.0080,74.5922) .. controls (274.0080,73.8714) and (273.4237,73.2870) ..(272.7029,73.2870) .. controls (271.9821,73.2870) and (271.3978,73.8714) ..(271.3978,74.5922) .. controls (271.3978,75.3130) and (271.9821,75.8973) ..(272.7029,75.8973) .. controls (273.4237,75.8973) and (274.0080,75.3130) ..(274.0080,74.5922);
\path[fill=chfourfig5c0,fill opacity=0.5,draw=chfourfig5c0,line width=0.3985pt,draw opacity=0.5] (262.2121,76.9135) .. controls (262.2121,76.1927) and (261.6278,75.6084) ..(260.9070,75.6084) .. controls (260.1862,75.6084) and (259.6019,76.1927) ..(259.6019,76.9135) .. controls (259.6019,77.6343) and (260.1862,78.2186) ..(260.9070,78.2186) .. controls (261.6278,78.2186) and (262.2121,77.6343) ..(262.2121,76.9135);
\path[fill=chfourfig5c0,fill opacity=0.5,draw=chfourfig5c0,line width=0.3985pt,draw opacity=0.5] (329.9091,45.4810) .. controls (329.9091,44.7601) and (329.3248,44.1758) ..(328.6040,44.1758) .. controls (327.8832,44.1758) and (327.2989,44.7601) ..(327.2989,45.4810) .. controls (327.2989,46.2018) and (327.8832,46.7861) ..(328.6040,46.7861) .. controls (329.3248,46.7861) and (329.9091,46.2018) ..(329.9091,45.4810);
\path[fill=chfourfig5c0,fill opacity=0.5,draw=chfourfig5c0,line width=0.3985pt,draw opacity=0.5] (232.2640,79.6632) .. controls (232.2640,78.9424) and (231.6797,78.3581) ..(230.9589,78.3581) .. controls (230.2381,78.3581) and (229.6538,78.9424) ..(229.6538,79.6632) .. controls (229.6538,80.3840) and (230.2381,80.9683) ..(230.9589,80.9683) .. controls (231.6797,80.9683) and (232.2640,80.3840) ..(232.2640,79.6632);
\path[fill=chfourfig5c0,fill opacity=0.5,draw=chfourfig5c0,line width=0.3985pt,draw opacity=0.5] (290.8351,51.5782) .. controls (290.8351,50.8574) and (290.2508,50.2730) ..(289.5300,50.2730) .. controls (288.8092,50.2730) and (288.2249,50.8574) ..(288.2249,51.5782) .. controls (288.2249,52.2990) and (288.8092,52.8833) ..(289.5300,52.8833) .. controls (290.2508,52.8833) and (290.8351,52.2990) ..(290.8351,51.5782);
\path[fill=chfourfig5c0,fill opacity=0.5,draw=chfourfig5c0,line width=0.3985pt,draw opacity=0.5] (274.4962,90.4928) .. controls (274.4962,89.7719) and (273.9119,89.1876) ..(273.1911,89.1876) .. controls (272.4703,89.1876) and (271.8860,89.7719) ..(271.8860,90.4928) .. controls (271.8860,91.2135) and (272.4703,91.7979) ..(273.1911,91.7979) .. controls (273.9119,91.7979) and (274.4962,91.2135) ..(274.4962,90.4928);
\path[fill=chfourfig5c0,fill opacity=0.5,draw=chfourfig5c0,line width=0.3985pt,draw opacity=0.5] (260.4786,92.3358) .. controls (260.4786,91.6151) and (259.8943,91.0307) ..(259.1735,91.0307) .. controls (258.4527,91.0307) and (257.8683,91.6151) ..(257.8683,92.3358) .. controls (257.8683,93.0567) and (258.4527,93.6410) ..(259.1735,93.6410) .. controls (259.8943,93.6410) and (260.4786,93.0567) ..(260.4786,92.3358);
\path[fill=chfourfig5c0,fill opacity=0.5,draw=chfourfig5c0,line width=0.3985pt,draw opacity=0.5] (326.7011,52.7936) .. controls (326.7011,52.0728) and (326.1168,51.4885) ..(325.3960,51.4885) .. controls (324.6752,51.4885) and (324.0909,52.0728) ..(324.0909,52.7936) .. controls (324.0909,53.5144) and (324.6752,54.0987) ..(325.3960,54.0987) .. controls (326.1168,54.0987) and (326.7011,53.5144) ..(326.7011,52.7936);
\path[fill=chfourfig5c0,fill opacity=0.5,draw=chfourfig5c0,line width=0.3985pt,draw opacity=0.5] (325.9240,38.8856) .. controls (325.9240,38.1648) and (325.3397,37.5805) ..(324.6189,37.5805) .. controls (323.8981,37.5805) and (323.3138,38.1648) ..(323.3138,38.8856) .. controls (323.3138,39.6064) and (323.8981,40.1907) ..(324.6189,40.1907) .. controls (325.3397,40.1907) and (325.9240,39.6064) ..(325.9240,38.8856);
\path[fill=chfourfig5c0,fill opacity=0.5,draw=chfourfig5c0,line width=0.3985pt,draw opacity=0.5] (260.2794,101.0832) .. controls (260.2794,100.3624) and (259.6950,99.7780) ..(258.9742,99.7780) .. controls (258.2534,99.7780) and (257.6691,100.3624) ..(257.6691,101.0832) .. controls (257.6691,101.8040) and (258.2534,102.3883) ..(258.9742,102.3883) .. controls (259.6950,102.3883) and (260.2794,101.8040) ..(260.2794,101.0832);
\path[fill=chfourfig5c0,fill opacity=0.5,draw=chfourfig5c0,line width=0.3985pt,draw opacity=0.5] (263.7265,61.4911) .. controls (263.7265,60.7703) and (263.1422,60.1860) ..(262.4213,60.1860) .. controls (261.7006,60.1860) and (261.1162,60.7703) ..(261.1162,61.4911) .. controls (261.1162,62.2119) and (261.7006,62.7962) ..(262.4213,62.7962) .. controls (263.1422,62.7962) and (263.7265,62.2119) ..(263.7265,61.4911);
\path[fill=chfourfig5c0,fill opacity=0.5,draw=chfourfig5c0,line width=0.3985pt,draw opacity=0.5] (318.0136,36.7137) .. controls (318.0136,35.9929) and (317.4293,35.4086) ..(316.7085,35.4086) .. controls (315.9877,35.4086) and (315.4034,35.9929) ..(315.4034,36.7137) .. controls (315.4034,37.4345) and (315.9877,38.0188) ..(316.7085,38.0188) .. controls (317.4293,38.0188) and (318.0136,37.4345) ..(318.0136,36.7137);
\path[fill=chfourfig5c0,fill opacity=0.5,draw=chfourfig5c0,line width=0.3985pt,draw opacity=0.5] (328.5044,39.2841) .. controls (328.5044,38.5633) and (327.9200,37.9790) ..(327.1992,37.9790) .. controls (326.4785,37.9790) and (325.8941,38.5633) ..(325.8941,39.2841) .. controls (325.8941,40.0049) and (326.4785,40.5892) ..(327.1992,40.5892) .. controls (327.9200,40.5892) and (328.5044,40.0049) ..(328.5044,39.2841);
\path[fill=chfourfig5c0,fill opacity=0.5,draw=chfourfig5c0,line width=0.3985pt,draw opacity=0.5] (305.9188,74.2734) .. controls (305.9188,73.5526) and (305.3344,72.9682) ..(304.6137,72.9682) .. controls (303.8929,72.9682) and (303.3085,73.5526) ..(303.3085,74.2734) .. controls (303.3085,74.9942) and (303.8929,75.5785) ..(304.6137,75.5785) .. controls (305.3344,75.5785) and (305.9188,74.9942) ..(305.9188,74.2734);
\path[fill=chfourfig5c0,fill opacity=0.5,draw=chfourfig5c0,line width=0.3985pt,draw opacity=0.5] (256.7126,68.2857) .. controls (256.7126,67.5649) and (256.1284,66.9806) ..(255.4075,66.9806) .. controls (254.6868,66.9806) and (254.1024,67.5649) ..(254.1024,68.2857) .. controls (254.1024,69.0065) and (254.6868,69.5909) ..(255.4075,69.5909) .. controls (256.1284,69.5909) and (256.7126,69.0065) ..(256.7126,68.2857);
\path[fill=chfourfig5c0,fill opacity=0.5,draw=chfourfig5c0,line width=0.3985pt,draw opacity=0.5] (304.5339,39.1845) .. controls (304.5339,38.4637) and (303.9496,37.8794) ..(303.2289,37.8794) .. controls (302.5080,37.8794) and (301.9237,38.4637) ..(301.9237,39.1845) .. controls (301.9237,39.9053) and (302.5080,40.4896) ..(303.2289,40.4896) .. controls (303.9496,40.4896) and (304.5339,39.9053) ..(304.5339,39.1845);
\path[fill=chfourfig5c0,fill opacity=0.5,draw=chfourfig5c0,line width=0.3985pt,draw opacity=0.5] (321.1519,65.3467) .. controls (321.1519,64.6259) and (320.5675,64.0416) ..(319.8467,64.0416) .. controls (319.1259,64.0416) and (318.5416,64.6259) ..(318.5416,65.3467) .. controls (318.5416,66.0675) and (319.1259,66.6518) ..(319.8467,66.6518) .. controls (320.5675,66.6518) and (321.1519,66.0675) ..(321.1519,65.3467);
\path[fill=chfourfig5c0,fill opacity=0.5,draw=chfourfig5c0,line width=0.3985pt,draw opacity=0.5] (248.2443,94.3185) .. controls (248.2443,93.5976) and (247.6600,93.0133) ..(246.9392,93.0133) .. controls (246.2184,93.0133) and (245.6341,93.5976) ..(245.6341,94.3185) .. controls (245.6341,95.0392) and (246.2184,95.6236) ..(246.9392,95.6236) .. controls (247.6600,95.6236) and (248.2443,95.0392) ..(248.2443,94.3185);
\path[fill=chfourfig5c0,fill opacity=0.5,draw=chfourfig5c0,line width=0.3985pt,draw opacity=0.5] (329.5903,48.7786) .. controls (329.5903,48.0578) and (329.0060,47.4735) ..(328.2852,47.4735) .. controls (327.5644,47.4735) and (326.9801,48.0578) ..(326.9801,48.7786) .. controls (326.9801,49.4994) and (327.5644,50.0837) ..(328.2852,50.0837) .. controls (329.0060,50.0837) and (329.5903,49.4994) ..(329.5903,48.7786);
\path[fill=chfourfig5c0,fill opacity=0.5,draw=chfourfig5c0,line width=0.3985pt,draw opacity=0.5] (277.4053,81.8251) .. controls (277.4053,81.1043) and (276.8210,80.5200) ..(276.1002,80.5200) .. controls (275.3794,80.5200) and (274.7951,81.1043) ..(274.7951,81.8251) .. controls (274.7951,82.5459) and (275.3794,83.1303) ..(276.1002,83.1303) .. controls (276.8210,83.1303) and (277.4053,82.5459) ..(277.4053,81.8251);
\path[fill=chfourfig5c0,fill opacity=0.5,draw=chfourfig5c0,line width=0.3985pt,draw opacity=0.5] (329.2815,26.3126) .. controls (329.2815,25.5918) and (328.6971,25.0075) ..(327.9763,25.0075) .. controls (327.2556,25.0075) and (326.6712,25.5918) ..(326.6712,26.3126) .. controls (326.6712,27.0334) and (327.2556,27.6177) ..(327.9763,27.6177) .. controls (328.6971,27.6177) and (329.2815,27.0334) ..(329.2815,26.3126);
\path[fill=chfourfig5c0,fill opacity=0.5,draw=chfourfig5c0,line width=0.3985pt,draw opacity=0.5] (310.2526,48.0414) .. controls (310.2526,47.3206) and (309.6683,46.7363) ..(308.9475,46.7363) .. controls (308.2267,46.7363) and (307.6423,47.3206) ..(307.6423,48.0414) .. controls (307.6423,48.7622) and (308.2267,49.3465) ..(308.9475,49.3465) .. controls (309.6683,49.3465) and (310.2526,48.7622) ..(310.2526,48.0414);
\path[fill=chfourfig5c0,fill opacity=0.5,draw=chfourfig5c0,line width=0.3985pt,draw opacity=0.5] (282.1177,62.0989) .. controls (282.1177,61.3781) and (281.5334,60.7937) ..(280.8126,60.7937) .. controls (280.0918,60.7937) and (279.5075,61.3781) ..(279.5075,62.0989) .. controls (279.5075,62.8197) and (280.0918,63.4040) ..(280.8126,63.4040) .. controls (281.5334,63.4040) and (282.1177,62.8197) ..(282.1177,62.0989);
\path[fill=chfourfig5c0,fill opacity=0.5,draw=chfourfig5c0,line width=0.3985pt,draw opacity=0.5] (264.8024,76.7342) .. controls (264.8024,76.0134) and (264.2181,75.4290) ..(263.4973,75.4290) .. controls (262.7765,75.4290) and (262.1922,76.0134) ..(262.1922,76.7342) .. controls (262.1922,77.4550) and (262.7765,78.0393) ..(263.4973,78.0393) .. controls (264.2181,78.0393) and (264.8024,77.4550) ..(264.8024,76.7342);
\path[fill=chfourfig5c0,fill opacity=0.5,draw=chfourfig5c0,line width=0.3985pt,draw opacity=0.5] (288.0655,79.3842) .. controls (288.0655,78.6635) and (287.4812,78.0791) ..(286.7604,78.0791) .. controls (286.0396,78.0791) and (285.4553,78.6635) ..(285.4553,79.3842) .. controls (285.4553,80.1051) and (286.0396,80.6894) ..(286.7604,80.6894) .. controls (287.4812,80.6894) and (288.0655,80.1051) ..(288.0655,79.3842);
\path[fill=chfourfig5c0,fill opacity=0.5,draw=chfourfig5c0,line width=0.3985pt,draw opacity=0.5] (260.8372,66.8411) .. controls (260.8372,66.1203) and (260.2529,65.5360) ..(259.5321,65.5360) .. controls (258.8113,65.5360) and (258.2270,66.1203) ..(258.2270,66.8411) .. controls (258.2270,67.5619) and (258.8113,68.1463) ..(259.5321,68.1463) .. controls (260.2529,68.1463) and (260.8372,67.5619) ..(260.8372,66.8411);
\path[fill=chfourfig5c0,fill opacity=0.5,draw=chfourfig5c0,line width=0.3985pt,draw opacity=0.5] (275.5522,78.2884) .. controls (275.5522,77.5676) and (274.9679,76.9832) ..(274.2472,76.9832) .. controls (273.5264,76.9832) and (272.9420,77.5676) ..(272.9420,78.2884) .. controls (272.9420,79.0092) and (273.5264,79.5935) ..(274.2472,79.5935) .. controls (274.9679,79.5935) and (275.5522,79.0092) ..(275.5522,78.2884);
\path[fill=chfourfig5c0,fill opacity=0.5,draw=chfourfig5c0,line width=0.3985pt,draw opacity=0.5] (272.6432,72.0716) .. controls (272.6432,71.3508) and (272.0588,70.7665) ..(271.3380,70.7665) .. controls (270.6172,70.7665) and (270.0329,71.3508) ..(270.0329,72.0716) .. controls (270.0329,72.7924) and (270.6172,73.3767) ..(271.3380,73.3767) .. controls (272.0588,73.3767) and (272.6432,72.7924) ..(272.6432,72.0716);
\path[fill=chfourfig5c0,fill opacity=0.5,draw=chfourfig5c0,line width=0.3985pt,draw opacity=0.5] (253.9430,69.7104) .. controls (253.9430,68.9896) and (253.3587,68.4053) ..(252.6379,68.4053) .. controls (251.9171,68.4053) and (251.3328,68.9896) ..(251.3328,69.7104) .. controls (251.3328,70.4312) and (251.9171,71.0155) ..(252.6379,71.0155) .. controls (253.3587,71.0155) and (253.9430,70.4312) ..(253.9430,69.7104);
\path[fill=chfourfig5c0,fill opacity=0.5,draw=chfourfig5c0,line width=0.3985pt,draw opacity=0.5] (290.2274,46.2680) .. controls (290.2274,45.5472) and (289.6431,44.9629) ..(288.9223,44.9629) .. controls (288.2015,44.9629) and (287.6172,45.5472) ..(287.6172,46.2680) .. controls (287.6172,46.9888) and (288.2015,47.5731) ..(288.9223,47.5731) .. controls (289.6431,47.5731) and (290.2274,46.9888) ..(290.2274,46.2680);
\path[fill=chfourfig5c0,fill opacity=0.5,draw=chfourfig5c0,line width=0.3985pt,draw opacity=0.5] (238.3114,93.7805) .. controls (238.3114,93.0597) and (237.7271,92.4753) ..(237.0063,92.4753) .. controls (236.2855,92.4753) and (235.7012,93.0597) ..(235.7012,93.7805) .. controls (235.7012,94.5013) and (236.2855,95.0856) ..(237.0063,95.0856) .. controls (237.7271,95.0856) and (238.3114,94.5013) ..(238.3114,93.7805);
\path[fill=chfourfig5c0,fill opacity=0.5,draw=chfourfig5c0,line width=0.3985pt,draw opacity=0.5] (305.2911,44.9131) .. controls (305.2911,44.1923) and (304.7068,43.6079) ..(303.9860,43.6079) .. controls (303.2652,43.6079) and (302.6809,44.1923) ..(302.6809,44.9131) .. controls (302.6809,45.6339) and (303.2652,46.2182) ..(303.9860,46.2182) .. controls (304.7068,46.2182) and (305.2911,45.6339) ..(305.2911,44.9131);
\path[fill=chfourfig5c0,fill opacity=0.5,draw=chfourfig5c0,line width=0.3985pt,draw opacity=0.5] (261.2358,99.6485) .. controls (261.2358,98.9277) and (260.6515,98.3434) ..(259.9306,98.3434) .. controls (259.2099,98.3434) and (258.6255,98.9277) ..(258.6255,99.6485) .. controls (258.6255,100.3693) and (259.2099,100.9537) ..(259.9306,100.9537) .. controls (260.6515,100.9537) and (261.2358,100.3693) ..(261.2358,99.6485);
\path[fill=chfourfig5c0,fill opacity=0.5,draw=chfourfig5c0,line width=0.3985pt,draw opacity=0.5] (288.8227,47.2643) .. controls (288.8227,46.5435) and (288.2383,45.9592) ..(287.5175,45.9592) .. controls (286.7967,45.9592) and (286.2124,46.5435) ..(286.2124,47.2643) .. controls (286.2124,47.9851) and (286.7967,48.5694) ..(287.5175,48.5694) .. controls (288.2383,48.5694) and (288.8227,47.9851) ..(288.8227,47.2643);
\path[fill=chfourfig5c0,fill opacity=0.5,draw=chfourfig5c0,line width=0.3985pt,draw opacity=0.5] (307.3833,54.9755) .. controls (307.3833,54.2547) and (306.7990,53.6703) ..(306.0782,53.6703) .. controls (305.3574,53.6703) and (304.7731,54.2547) ..(304.7731,54.9755) .. controls (304.7731,55.6963) and (305.3574,56.2806) ..(306.0782,56.2806) .. controls (306.7990,56.2806) and (307.3833,55.6963) ..(307.3833,54.9755);
\path[fill=chfourfig5c0,fill opacity=0.5,draw=chfourfig5c0,line width=0.3985pt,draw opacity=0.5] (304.4144,45.5507) .. controls (304.4144,44.8299) and (303.8301,44.2456) ..(303.1093,44.2456) .. controls (302.3885,44.2456) and (301.8042,44.8299) ..(301.8042,45.5507) .. controls (301.8042,46.2715) and (302.3885,46.8558) ..(303.1093,46.8558) .. controls (303.8301,46.8558) and (304.4144,46.2715) ..(304.4144,45.5507);
\path[fill=chfourfig5c0,fill opacity=0.5,draw=chfourfig5c0,line width=0.3985pt,draw opacity=0.5] (248.4436,78.2784) .. controls (248.4436,77.5576) and (247.8593,76.9733) ..(247.1384,76.9733) .. controls (246.4177,76.9733) and (245.8333,77.5576) ..(245.8333,78.2784) .. controls (245.8333,78.9992) and (246.4177,79.5835) ..(247.1384,79.5835) .. controls (247.8593,79.5835) and (248.4436,78.9992) ..(248.4436,78.2784);
\path[fill=chfourfig5c0,fill opacity=0.5,draw=chfourfig5c0,line width=0.3985pt,draw opacity=0.5] (230.7696,88.2312) .. controls (230.7696,87.5104) and (230.1853,86.9261) ..(229.4645,86.9261) .. controls (228.7437,86.9261) and (228.1594,87.5104) ..(228.1594,88.2312) .. controls (228.1594,88.9520) and (228.7437,89.5363) ..(229.4645,89.5363) .. controls (230.1853,89.5363) and (230.7696,88.9520) ..(230.7696,88.2312);
\path[fill=chfourfig5c0,fill opacity=0.5,draw=chfourfig5c0,line width=0.3985pt,draw opacity=0.5] (295.1989,54.3279) .. controls (295.1989,53.6071) and (294.6145,53.0228) ..(293.8937,53.0228) .. controls (293.1729,53.0228) and (292.5886,53.6071) ..(292.5886,54.3279) .. controls (292.5886,55.0487) and (293.1729,55.6330) ..(293.8937,55.6330) .. controls (294.6145,55.6330) and (295.1989,55.0487) ..(295.1989,54.3279);
\path[fill=chfourfig5c0,fill opacity=0.5,draw=chfourfig5c0,line width=0.3985pt,draw opacity=0.5] (245.8533,108.6648) .. controls (245.8533,107.9440) and (245.2689,107.3597) ..(244.5482,107.3597) .. controls (243.8273,107.3597) and (243.2430,107.9440) ..(243.2430,108.6648) .. controls (243.2430,109.3856) and (243.8273,109.9699) ..(244.5482,109.9699) .. controls (245.2689,109.9699) and (245.8533,109.3856) ..(245.8533,108.6648);
\path[fill=chfourfig5c0,fill opacity=0.5,draw=chfourfig5c0,line width=0.3985pt,draw opacity=0.5] (244.8570,84.1564) .. controls (244.8570,83.4356) and (244.2726,82.8513) ..(243.5519,82.8513) .. controls (242.8311,82.8513) and (242.2467,83.4356) ..(242.2467,84.1564) .. controls (242.2467,84.8772) and (242.8311,85.4615) ..(243.5519,85.4615) .. controls (244.2726,85.4615) and (244.8570,84.8772) ..(244.8570,84.1564);
\path[fill=chfourfig5c0,fill opacity=0.5,draw=chfourfig5c0,line width=0.3985pt,draw opacity=0.5] (295.0494,56.2009) .. controls (295.0494,55.4801) and (294.4651,54.8958) ..(293.7443,54.8958) .. controls (293.0235,54.8958) and (292.4392,55.4801) ..(292.4392,56.2009) .. controls (292.4392,56.9217) and (293.0235,57.5060) ..(293.7443,57.5060) .. controls (294.4651,57.5060) and (295.0494,56.9217) ..(295.0494,56.2009);
\path[fill=chfourfig5c0,fill opacity=0.5,draw=chfourfig5c0,line width=0.3985pt,draw opacity=0.5] (292.1004,83.8575) .. controls (292.1004,83.1367) and (291.5161,82.5524) ..(290.7953,82.5524) .. controls (290.0745,82.5524) and (289.4902,83.1367) ..(289.4902,83.8575) .. controls (289.4902,84.5783) and (290.0745,85.1627) ..(290.7953,85.1627) .. controls (291.5161,85.1627) and (292.1004,84.5783) ..(292.1004,83.8575);
\path[fill=chfourfig5c0,fill opacity=0.5,draw=chfourfig5c0,line width=0.3985pt,draw opacity=0.5] (297.1117,69.0130) .. controls (297.1117,68.2922) and (296.5274,67.7079) ..(295.8066,67.7079) .. controls (295.0858,67.7079) and (294.5014,68.2922) ..(294.5014,69.0130) .. controls (294.5014,69.7338) and (295.0858,70.3181) ..(295.8066,70.3181) .. controls (296.5274,70.3181) and (297.1117,69.7338) ..(297.1117,69.0130);
\path[fill=chfourfig5c0,fill opacity=0.5,draw=chfourfig5c0,line width=0.3985pt,draw opacity=0.5] (330.1981,33.4559) .. controls (330.1981,32.7351) and (329.6137,32.1508) ..(328.8929,32.1508) .. controls (328.1721,32.1508) and (327.5878,32.7351) ..(327.5878,33.4559) .. controls (327.5878,34.1767) and (328.1721,34.7610) ..(328.8929,34.7610) .. controls (329.6137,34.7610) and (330.1981,34.1767) ..(330.1981,33.4559);
\path[fill=chfourfig5c0,fill opacity=0.5,draw=chfourfig5c0,line width=0.3985pt,draw opacity=0.5] (318.2427,62.6867) .. controls (318.2427,61.9659) and (317.6584,61.3815) ..(316.9376,61.3815) .. controls (316.2168,61.3815) and (315.6325,61.9659) ..(315.6325,62.6867) .. controls (315.6325,63.4075) and (316.2168,63.9918) ..(316.9376,63.9918) .. controls (317.6584,63.9918) and (318.2427,63.4075) ..(318.2427,62.6867);
\path[fill=chfourfig5c0,fill opacity=0.5,draw=chfourfig5c0,line width=0.3985pt,draw opacity=0.5] (240.8818,112.6599) .. controls (240.8818,111.9391) and (240.2975,111.3548) ..(239.5767,111.3548) .. controls (238.8559,111.3548) and (238.2716,111.9391) ..(238.2716,112.6599) .. controls (238.2716,113.3807) and (238.8559,113.9650) ..(239.5767,113.9650) .. controls (240.2975,113.9650) and (240.8818,113.3807) ..(240.8818,112.6599);
\path[fill=chfourfig5c0,fill opacity=0.5,draw=chfourfig5c0,line width=0.3985pt,draw opacity=0.5] (245.7337,99.3098) .. controls (245.7337,98.5890) and (245.1494,98.0047) ..(244.4286,98.0047) .. controls (243.7078,98.0047) and (243.1235,98.5890) ..(243.1235,99.3098) .. controls (243.1235,100.0306) and (243.7078,100.6149) ..(244.4286,100.6149) .. controls (245.1494,100.6149) and (245.7337,100.0306) ..(245.7337,99.3098);
\path[fill=chfourfig5c0,fill opacity=0.5,draw=chfourfig5c0,line width=0.3985pt,draw opacity=0.5] (254.3714,72.2509) .. controls (254.3714,71.5301) and (253.7871,70.9458) ..(253.0663,70.9458) .. controls (252.3455,70.9458) and (251.7612,71.5301) ..(251.7612,72.2509) .. controls (251.7612,72.9717) and (252.3455,73.5560) ..(253.0663,73.5560) .. controls (253.7871,73.5560) and (254.3714,72.9717) ..(254.3714,72.2509);
\path[fill=chfourfig5c0,fill opacity=0.5,draw=chfourfig5c0,line width=0.3985pt,draw opacity=0.5] (229.8431,114.7621) .. controls (229.8431,114.0412) and (229.2588,113.4569) ..(228.5380,113.4569) .. controls (227.8171,113.4569) and (227.2328,114.0412) ..(227.2328,114.7621) .. controls (227.2328,115.4828) and (227.8171,116.0672) ..(228.5380,116.0672) .. controls (229.2588,116.0672) and (229.8431,115.4828) ..(229.8431,114.7621);
\path[fill=chfourfig5c0,fill opacity=0.5,draw=chfourfig5c0,line width=0.3985pt,draw opacity=0.5] (307.7220,40.2405) .. controls (307.7220,39.5197) and (307.1377,38.9354) ..(306.4169,38.9354) .. controls (305.6961,38.9354) and (305.1118,39.5197) ..(305.1118,40.2405) .. controls (305.1118,40.9613) and (305.6961,41.5457) ..(306.4169,41.5457) .. controls (307.1377,41.5457) and (307.7220,40.9613) ..(307.7220,40.2405);
\path[fill=chfourfig5c0,fill opacity=0.5,draw=chfourfig5c0,line width=0.3985pt,draw opacity=0.5] (316.8380,32.0810) .. controls (316.8380,31.3602) and (316.2537,30.7759) ..(315.5329,30.7759) .. controls (314.8120,30.7759) and (314.2277,31.3602) ..(314.2277,32.0810) .. controls (314.2277,32.8018) and (314.8120,33.3862) ..(315.5329,33.3862) .. controls (316.2537,33.3862) and (316.8380,32.8018) ..(316.8380,32.0810);
\path[fill=chfourfig5c0,fill opacity=0.5,draw=chfourfig5c0,line width=0.3985pt,draw opacity=0.5] (310.9101,59.4488) .. controls (310.9101,58.7280) and (310.3258,58.1436) ..(309.6050,58.1436) .. controls (308.8842,58.1436) and (308.2999,58.7280) ..(308.2999,59.4488) .. controls (308.2999,60.1696) and (308.8842,60.7539) ..(309.6050,60.7539) .. controls (310.3258,60.7539) and (310.9101,60.1696) ..(310.9101,59.4488);
\path[fill=chfourfig5c0,fill opacity=0.5,draw=chfourfig5c0,line width=0.3985pt,draw opacity=0.5] (245.7835,81.1875) .. controls (245.7835,80.4667) and (245.1992,79.8824) ..(244.4784,79.8824) .. controls (243.7576,79.8824) and (243.1733,80.4667) ..(243.1733,81.1875) .. controls (243.1733,81.9083) and (243.7576,82.4926) ..(244.4784,82.4926) .. controls (245.1992,82.4926) and (245.7835,81.9083) ..(245.7835,81.1875);
\path[fill=chfourfig5c0,fill opacity=0.5,draw=chfourfig5c0,line width=0.3985pt,draw opacity=0.5] (253.8035,88.0419) .. controls (253.8035,87.3211) and (253.2192,86.7368) ..(252.4984,86.7368) .. controls (251.7776,86.7368) and (251.1933,87.3211) ..(251.1933,88.0419) .. controls (251.1933,88.7627) and (251.7776,89.3470) ..(252.4984,89.3470) .. controls (253.2192,89.3470) and (253.8035,88.7627) ..(253.8035,88.0419);
\path[fill=chfourfig5c0,fill opacity=0.5,draw=chfourfig5c0,line width=0.3985pt,draw opacity=0.5] (286.0730,55.0054) .. controls (286.0730,54.2845) and (285.4886,53.7002) ..(284.7678,53.7002) .. controls (284.0470,53.7002) and (283.4627,54.2845) ..(283.4627,55.0054) .. controls (283.4627,55.7262) and (284.0470,56.3105) ..(284.7678,56.3105) .. controls (285.4886,56.3105) and (286.0730,55.7262) ..(286.0730,55.0054);
\path[fill=chfourfig5c0,fill opacity=0.5,draw=chfourfig5c0,line width=0.3985pt,draw opacity=0.5] (287.2685,69.8897) .. controls (287.2685,69.1689) and (286.6841,68.5846) ..(285.9634,68.5846) .. controls (285.2426,68.5846) and (284.6582,69.1689) ..(284.6582,69.8897) .. controls (284.6582,70.6105) and (285.2426,71.1949) ..(285.9634,71.1949) .. controls (286.6841,71.1949) and (287.2685,70.6105) ..(287.2685,69.8897);
\path[fill=chfourfig5c0,fill opacity=0.5,draw=chfourfig5c0,line width=0.3985pt,draw opacity=0.5] (301.3658,50.0638) .. controls (301.3658,49.3430) and (300.7815,48.7587) ..(300.0607,48.7587) .. controls (299.3399,48.7587) and (298.7556,49.3430) ..(298.7556,50.0638) .. controls (298.7556,50.7846) and (299.3399,51.3689) ..(300.0607,51.3689) .. controls (300.7815,51.3689) and (301.3658,50.7846) ..(301.3658,50.0638);
\path[fill=chfourfig5c0,fill opacity=0.5,draw=chfourfig5c0,line width=0.3985pt,draw opacity=0.5] (318.7608,40.8781) .. controls (318.7608,40.1573) and (318.1765,39.5730) ..(317.4557,39.5730) .. controls (316.7349,39.5730) and (316.1505,40.1573) ..(316.1505,40.8781) .. controls (316.1505,41.5990) and (316.7349,42.1833) ..(317.4557,42.1833) .. controls (318.1765,42.1833) and (318.7608,41.5990) ..(318.7608,40.8781);
\path[fill=chfourfig5c0,fill opacity=0.5,draw=chfourfig5c0,line width=0.3985pt,draw opacity=0.5] (293.5650,65.1873) .. controls (293.5650,64.4665) and (292.9807,63.8822) ..(292.2598,63.8822) .. controls (291.5390,63.8822) and (290.9547,64.4665) ..(290.9547,65.1873) .. controls (290.9547,65.9081) and (291.5390,66.4924) ..(292.2598,66.4924) .. controls (292.9807,66.4924) and (293.5650,65.9081) ..(293.5650,65.1873);
\path[fill=chfourfig5c0,fill opacity=0.5,draw=chfourfig5c0,line width=0.3985pt,draw opacity=0.5] (260.2694,78.1987) .. controls (260.2694,77.4779) and (259.6851,76.8936) ..(258.9643,76.8936) .. controls (258.2435,76.8936) and (257.6591,77.4779) ..(257.6591,78.1987) .. controls (257.6591,78.9195) and (258.2435,79.5038) ..(258.9643,79.5038) .. controls (259.6851,79.5038) and (260.2694,78.9195) ..(260.2694,78.1987);
\path[fill=chfourfig5c0,fill opacity=0.5,draw=chfourfig5c0,line width=0.3985pt,draw opacity=0.5] (264.9917,60.2657) .. controls (264.9917,59.5449) and (264.4074,58.9606) ..(263.6866,58.9606) .. controls (262.9658,58.9606) and (262.3815,59.5449) ..(262.3815,60.2657) .. controls (262.3815,60.9865) and (262.9658,61.5708) ..(263.6866,61.5708) .. controls (264.4074,61.5708) and (264.9917,60.9865) ..(264.9917,60.2657);
\path[fill=chfourfig5c0,fill opacity=0.5,draw=chfourfig5c0,line width=0.3985pt,draw opacity=0.5] (232.2740,82.7118) .. controls (232.2740,81.9910) and (231.6897,81.4067) ..(230.9689,81.4067) .. controls (230.2481,81.4067) and (229.6638,81.9910) ..(229.6638,82.7118) .. controls (229.6638,83.4326) and (230.2481,84.0170) ..(230.9689,84.0170) .. controls (231.6897,84.0170) and (232.2740,83.4326) ..(232.2740,82.7118);
\path[fill=chfourfig5c0,fill opacity=0.5,draw=chfourfig5c0,line width=0.3985pt,draw opacity=0.5] (235.3625,111.5939) .. controls (235.3625,110.8731) and (234.7781,110.2888) ..(234.0573,110.2888) .. controls (233.3365,110.2888) and (232.7522,110.8731) ..(232.7522,111.5939) .. controls (232.7522,112.3147) and (233.3365,112.8990) ..(234.0573,112.8990) .. controls (234.7781,112.8990) and (235.3625,112.3147) ..(235.3625,111.5939);
\path[fill=chfourfig5c0,fill opacity=0.5,draw=chfourfig5c0,line width=0.3985pt,draw opacity=0.5] (263.5770,80.0717) .. controls (263.5770,79.3509) and (262.9927,78.7666) ..(262.2719,78.7666) .. controls (261.5511,78.7666) and (260.9668,79.3509) ..(260.9668,80.0717) .. controls (260.9668,80.7925) and (261.5511,81.3768) ..(262.2719,81.3768) .. controls (262.9927,81.3768) and (263.5770,80.7925) ..(263.5770,80.0717);
\path[fill=chfourfig5c0,fill opacity=0.5,draw=chfourfig5c0,line width=0.3985pt,draw opacity=0.5] (326.1631,40.1509) .. controls (326.1631,39.4301) and (325.5788,38.8458) ..(324.8580,38.8458) .. controls (324.1372,38.8458) and (323.5529,39.4301) ..(323.5529,40.1509) .. controls (323.5529,40.8717) and (324.1372,41.4560) ..(324.8580,41.4560) .. controls (325.5788,41.4560) and (326.1631,40.8717) ..(326.1631,40.1509);
\path[fill=chfourfig5c0,fill opacity=0.5,draw=chfourfig5c0,line width=0.3985pt,draw opacity=0.5] (292.5487,54.7862) .. controls (292.5487,54.0654) and (291.9644,53.4810) ..(291.2436,53.4810) .. controls (290.5228,53.4810) and (289.9385,54.0654) ..(289.9385,54.7862) .. controls (289.9385,55.5070) and (290.5228,56.0913) ..(291.2436,56.0913) .. controls (291.9644,56.0913) and (292.5487,55.5070) ..(292.5487,54.7862);
\path[fill=chfourfig5c0,fill opacity=0.5,draw=chfourfig5c0,line width=0.3985pt,draw opacity=0.5] (305.3011,63.6032) .. controls (305.3011,62.8824) and (304.7168,62.2981) ..(303.9960,62.2981) .. controls (303.2752,62.2981) and (302.6909,62.8824) ..(302.6909,63.6032) .. controls (302.6909,64.3240) and (303.2752,64.9084) ..(303.9960,64.9084) .. controls (304.7168,64.9084) and (305.3011,64.3240) ..(305.3011,63.6032);
\path[fill=chfourfig5c0,fill opacity=0.5,draw=chfourfig5c0,line width=0.3985pt,draw opacity=0.5] (239.7261,84.6944) .. controls (239.7261,83.9736) and (239.1418,83.3893) ..(238.4210,83.3893) .. controls (237.7002,83.3893) and (237.1159,83.9736) ..(237.1159,84.6944) .. controls (237.1159,85.4152) and (237.7002,85.9995) ..(238.4210,85.9995) .. controls (239.1418,85.9995) and (239.7261,85.4152) ..(239.7261,84.6944);
\path[fill=chfourfig5c0,fill opacity=0.5,draw=chfourfig5c0,line width=0.3985pt,draw opacity=0.5] (291.2337,59.5085) .. controls (291.2337,58.7877) and (290.6494,58.2034) ..(289.9286,58.2034) .. controls (289.2077,58.2034) and (288.6234,58.7877) ..(288.6234,59.5085) .. controls (288.6234,60.2293) and (289.2077,60.8136) ..(289.9286,60.8136) .. controls (290.6494,60.8136) and (291.2337,60.2293) ..(291.2337,59.5085);
\path[fill=chfourfig5c0,fill opacity=0.5,draw=chfourfig5c0,line width=0.3985pt,draw opacity=0.5] (335.2193,43.7076) .. controls (335.2193,42.9868) and (334.6350,42.4025) ..(333.9142,42.4025) .. controls (333.1934,42.4025) and (332.6090,42.9868) ..(332.6090,43.7076) .. controls (332.6090,44.4284) and (333.1934,45.0127) ..(333.9142,45.0127) .. controls (334.6350,45.0127) and (335.2193,44.4284) ..(335.2193,43.7076);
\path[fill=chfourfig5c0,fill opacity=0.5,draw=chfourfig5c0,line width=0.3985pt,draw opacity=0.5] (256.2843,70.0193) .. controls (256.2843,69.2984) and (255.7000,68.7141) ..(254.9792,68.7141) .. controls (254.2584,68.7141) and (253.6740,69.2984) ..(253.6740,70.0193) .. controls (253.6740,70.7401) and (254.2584,71.3244) ..(254.9792,71.3244) .. controls (255.7000,71.3244) and (256.2843,70.7401) ..(256.2843,70.0193);
\path[fill=chfourfig5c0,fill opacity=0.5,draw=chfourfig5c0,line width=0.3985pt,draw opacity=0.5] (297.6895,43.0401) .. controls (297.6895,42.3193) and (297.1052,41.7350) ..(296.3844,41.7350) .. controls (295.6636,41.7350) and (295.0793,42.3193) ..(295.0793,43.0401) .. controls (295.0793,43.7609) and (295.6636,44.3452) ..(296.3844,44.3452) .. controls (297.1052,44.3452) and (297.6895,43.7609) ..(297.6895,43.0401);
\path[fill=chfourfig5c0,fill opacity=0.5,draw=chfourfig5c0,line width=0.3985pt,draw opacity=0.5] (282.6657,82.4030) .. controls (282.6657,81.6822) and (282.0814,81.0979) ..(281.3606,81.0979) .. controls (280.6398,81.0979) and (280.0555,81.6822) ..(280.0555,82.4030) .. controls (280.0555,83.1238) and (280.6398,83.7081) ..(281.3606,83.7081) .. controls (282.0814,83.7081) and (282.6657,83.1238) ..(282.6657,82.4030);
\path[fill=chfourfig5c0,fill opacity=0.5,draw=chfourfig5c0,line width=0.3985pt,draw opacity=0.5] (309.6449,60.3454) .. controls (309.6449,59.6246) and (309.0605,59.0403) ..(308.3397,59.0403) .. controls (307.6190,59.0403) and (307.0346,59.6246) ..(307.0346,60.3454) .. controls (307.0346,61.0662) and (307.6190,61.6505) ..(308.3397,61.6505) .. controls (309.0605,61.6505) and (309.6449,61.0662) ..(309.6449,60.3454);
\path[fill=chfourfig5c0,fill opacity=0.5,draw=chfourfig5c0,line width=0.3985pt,draw opacity=0.5] (303.1093,61.6206) .. controls (303.1093,60.8998) and (302.5250,60.3155) ..(301.8042,60.3155) .. controls (301.0833,60.3155) and (300.4991,60.8998) ..(300.4991,61.6206) .. controls (300.4991,62.3414) and (301.0833,62.9258) ..(301.8042,62.9258) .. controls (302.5250,62.9258) and (303.1093,62.3414) ..(303.1093,61.6206);
\path[fill=chfourfig5c0,fill opacity=0.5,draw=chfourfig5c0,line width=0.3985pt,draw opacity=0.5] (288.7231,83.3096) .. controls (288.7231,82.5888) and (288.1387,82.0045) ..(287.4179,82.0045) .. controls (286.6971,82.0045) and (286.1128,82.5888) ..(286.1128,83.3096) .. controls (286.1128,84.0304) and (286.6971,84.6147) ..(287.4179,84.6147) .. controls (288.1387,84.6147) and (288.7231,84.0304) ..(288.7231,83.3096);
\path[fill=chfourfig5c0,fill opacity=0.5,draw=chfourfig5c0,line width=0.3985pt,draw opacity=0.5] (306.4867,63.3043) .. controls (306.4867,62.5835) and (305.9023,61.9992) ..(305.1815,61.9992) .. controls (304.4607,61.9992) and (303.8764,62.5835) ..(303.8764,63.3043) .. controls (303.8764,64.0251) and (304.4607,64.6095) ..(305.1815,64.6095) .. controls (305.9023,64.6095) and (306.4867,64.0251) ..(306.4867,63.3043);
\path[fill=chfourfig5c0,fill opacity=0.5,draw=chfourfig5c0,line width=0.3985pt,draw opacity=0.5] (271.3280,74.3829) .. controls (271.3280,73.6622) and (270.7437,73.0778) ..(270.0229,73.0778) .. controls (269.3021,73.0778) and (268.7178,73.6622) ..(268.7178,74.3829) .. controls (268.7178,75.1038) and (269.3021,75.6881) ..(270.0229,75.6881) .. controls (270.7437,75.6881) and (271.3280,75.1038) ..(271.3280,74.3829);
\path[fill=chfourfig5c0,fill opacity=0.5,draw=chfourfig5c0,line width=0.3985pt,draw opacity=0.5] (278.6308,65.3168) .. controls (278.6308,64.5960) and (278.0464,64.0117) ..(277.3256,64.0117) .. controls (276.6049,64.0117) and (276.0205,64.5960) ..(276.0205,65.3168) .. controls (276.0205,66.0376) and (276.6049,66.6219) ..(277.3256,66.6219) .. controls (278.0464,66.6219) and (278.6308,66.0376) ..(278.6308,65.3168);
\path[fill=chfourfig5c0,fill opacity=0.5,draw=chfourfig5c0,line width=0.3985pt,draw opacity=0.5] (310.9201,49.9642) .. controls (310.9201,49.2434) and (310.3358,48.6591) ..(309.6150,48.6591) .. controls (308.8942,48.6591) and (308.3098,49.2434) ..(308.3098,49.9642) .. controls (308.3098,50.6850) and (308.8942,51.2693) ..(309.6150,51.2693) .. controls (310.3358,51.2693) and (310.9201,50.6850) ..(310.9201,49.9642);
\path[fill=chfourfig5c0,fill opacity=0.5,draw=chfourfig5c0,line width=0.3985pt,draw opacity=0.5] (285.6446,70.9657) .. controls (285.6446,70.2449) and (285.0602,69.6606) ..(284.3394,69.6606) .. controls (283.6186,69.6606) and (283.0343,70.2449) ..(283.0343,70.9657) .. controls (283.0343,71.6865) and (283.6186,72.2708) ..(284.3394,72.2708) .. controls (285.0602,72.2708) and (285.6446,71.6865) ..(285.6446,70.9657);
\path[fill=chfourfig5c0,fill opacity=0.5,draw=chfourfig5c0,line width=0.3985pt,draw opacity=0.5] (238.2417,107.5590) .. controls (238.2417,106.8382) and (237.6574,106.2538) ..(236.9366,106.2538) .. controls (236.2158,106.2538) and (235.6315,106.8382) ..(235.6315,107.5590) .. controls (235.6315,108.2798) and (236.2158,108.8641) ..(236.9366,108.8641) .. controls (237.6574,108.8641) and (238.2417,108.2798) ..(238.2417,107.5590);
\path[fill=chfourfig5c0,fill opacity=0.5,draw=chfourfig5c0,line width=0.3985pt,draw opacity=0.5] (276.1999,86.4578) .. controls (276.1999,85.7370) and (275.6155,85.1527) ..(274.8947,85.1527) .. controls (274.1739,85.1527) and (273.5896,85.7370) ..(273.5896,86.4578) .. controls (273.5896,87.1786) and (274.1739,87.7629) ..(274.8947,87.7629) .. controls (275.6155,87.7629) and (276.1999,87.1786) ..(276.1999,86.4578);
\path[fill=chfourfig5c0,fill opacity=0.5,draw=chfourfig5c0,line width=0.3985pt,draw opacity=0.5] (334.2430,27.7871) .. controls (334.2430,27.0663) and (333.6586,26.4820) ..(332.9378,26.4820) .. controls (332.2170,26.4820) and (331.6327,27.0663) ..(331.6327,27.7871) .. controls (331.6327,28.5079) and (332.2170,29.0922) ..(332.9378,29.0922) .. controls (333.6586,29.0922) and (334.2430,28.5079) ..(334.2430,27.7871);
\path[fill=chfourfig5c0,fill opacity=0.5,draw=chfourfig5c0,line width=0.3985pt,draw opacity=0.5] (316.3398,51.6280) .. controls (316.3398,50.9072) and (315.7555,50.3229) ..(315.0347,50.3229) .. controls (314.3139,50.3229) and (313.7296,50.9072) ..(313.7296,51.6280) .. controls (313.7296,52.3488) and (314.3139,52.9331) ..(315.0347,52.9331) .. controls (315.7555,52.9331) and (316.3398,52.3488) ..(316.3398,51.6280);
\path[fill=chfourfig5c0,fill opacity=0.5,draw=chfourfig5c0,line width=0.3985pt,draw opacity=0.5] (317.5752,31.3637) .. controls (317.5752,30.6429) and (316.9909,30.0586) ..(316.2701,30.0586) .. controls (315.5493,30.0586) and (314.9650,30.6429) ..(314.9650,31.3637) .. controls (314.9650,32.0845) and (315.5493,32.6688) ..(316.2701,32.6688) .. controls (316.9909,32.6688) and (317.5752,32.0845) ..(317.5752,31.3637);
\path[fill=chfourfig5c0,fill opacity=0.5,draw=chfourfig5c0,line width=0.3985pt,draw opacity=0.5] (252.3091,89.5463) .. controls (252.3091,88.8255) and (251.7248,88.2412) ..(251.0040,88.2412) .. controls (250.2832,88.2412) and (249.6989,88.8255) ..(249.6989,89.5463) .. controls (249.6989,90.2671) and (250.2832,90.8514) ..(251.0040,90.8514) .. controls (251.7248,90.8514) and (252.3091,90.2671) ..(252.3091,89.5463);
\path[fill=chfourfig5c0,fill opacity=0.5,draw=chfourfig5c0,line width=0.3985pt,draw opacity=0.5] (292.4093,77.3020) .. controls (292.4093,76.5812) and (291.8250,75.9969) ..(291.1042,75.9969) .. controls (290.3834,75.9969) and (289.7990,76.5812) ..(289.7990,77.3020) .. controls (289.7990,78.0228) and (290.3834,78.6072) ..(291.1042,78.6072) .. controls (291.8250,78.6072) and (292.4093,78.0228) ..(292.4093,77.3020);
\path[fill=chfourfig5c0,fill opacity=0.5,draw=chfourfig5c0,line width=0.3985pt,draw opacity=0.5] (318.9202,54.9057) .. controls (318.9202,54.1849) and (318.3359,53.6006) ..(317.6151,53.6006) .. controls (316.8943,53.6006) and (316.3099,54.1849) ..(316.3099,54.9057) .. controls (316.3099,55.6265) and (316.8943,56.2109) ..(317.6151,56.2109) .. controls (318.3359,56.2109) and (318.9202,55.6265) ..(318.9202,54.9057);
\path[fill=chfourfig5c0,fill opacity=0.5,draw=chfourfig5c0,line width=0.3985pt,draw opacity=0.5] (252.0103,68.4252) .. controls (252.0103,67.7044) and (251.4259,67.1201) ..(250.7051,67.1201) .. controls (249.9843,67.1201) and (249.4000,67.7044) ..(249.4000,68.4252) .. controls (249.4000,69.1460) and (249.9843,69.7303) ..(250.7051,69.7303) .. controls (251.4259,69.7303) and (252.0103,69.1460) ..(252.0103,68.4252);
\path[fill=chfourfig5c0,fill opacity=0.5,draw=chfourfig5c0,line width=0.3985pt,draw opacity=0.5] (315.4233,48.0115) .. controls (315.4233,47.2907) and (314.8390,46.7064) ..(314.1181,46.7064) .. controls (313.3974,46.7064) and (312.8130,47.2907) ..(312.8130,48.0115) .. controls (312.8130,48.7323) and (313.3974,49.3166) ..(314.1181,49.3166) .. controls (314.8390,49.3166) and (315.4233,48.7323) ..(315.4233,48.0115);
\path[fill=chfourfig5c0,fill opacity=0.5,draw=chfourfig5c0,line width=0.3985pt,draw opacity=0.5] (317.6250,58.3429) .. controls (317.6250,57.6221) and (317.0407,57.0378) ..(316.3199,57.0378) .. controls (315.5991,57.0378) and (315.0148,57.6221) ..(315.0148,58.3429) .. controls (315.0148,59.0637) and (315.5991,59.6480) ..(316.3199,59.6480) .. controls (317.0407,59.6480) and (317.6250,59.0637) ..(317.6250,58.3429);
\path[fill=chfourfig5c0,fill opacity=0.5,draw=chfourfig5c0,line width=0.3985pt,draw opacity=0.5] (266.1075,74.1040) .. controls (266.1075,73.3832) and (265.5233,72.7989) ..(264.8025,72.7989) .. controls (264.0816,72.7989) and (263.4973,73.3832) ..(263.4973,74.1040) .. controls (263.4973,74.8248) and (264.0816,75.4091) ..(264.8025,75.4091) .. controls (265.5233,75.4091) and (266.1075,74.8248) ..(266.1075,74.1040);
\path[fill=chfourfig5c0,fill opacity=0.5,draw=chfourfig5c0,line width=0.3985pt,draw opacity=0.5] (248.1845,69.9894) .. controls (248.1845,69.2686) and (247.6002,68.6842) ..(246.8794,68.6842) .. controls (246.1586,68.6842) and (245.5743,69.2686) ..(245.5743,69.9894) .. controls (245.5743,70.7102) and (246.1586,71.2945) ..(246.8794,71.2945) .. controls (247.6002,71.2945) and (248.1845,70.7102) ..(248.1845,69.9894);
\path[fill=chfourfig5c0,fill opacity=0.5,draw=chfourfig5c0,line width=0.3985pt,draw opacity=0.5] (249.2904,97.9947) .. controls (249.2904,97.2739) and (248.7061,96.6896) ..(247.9853,96.6896) .. controls (247.2645,96.6896) and (246.6802,97.2739) ..(246.6802,97.9947) .. controls (246.6802,98.7155) and (247.2645,99.2998) ..(247.9853,99.2998) .. controls (248.7061,99.2998) and (249.2904,98.7155) ..(249.2904,97.9947);
\path[fill=chfourfig5c0,fill opacity=0.5,draw=chfourfig5c0,line width=0.3985pt,draw opacity=0.5] (240.2841,100.8939) .. controls (240.2841,100.1731) and (239.6997,99.5888) ..(238.9789,99.5888) .. controls (238.2581,99.5888) and (237.6738,100.1731) ..(237.6738,100.8939) .. controls (237.6738,101.6147) and (238.2581,102.1990) ..(238.9789,102.1990) .. controls (239.6997,102.1990) and (240.2841,101.6147) ..(240.2841,100.8939);
\path[fill=chfourfig5c0,fill opacity=0.5,draw=chfourfig5c0,line width=0.3985pt,draw opacity=0.5] (289.5300,78.4179) .. controls (289.5300,77.6971) and (288.9457,77.1127) ..(288.2249,77.1127) .. controls (287.5041,77.1127) and (286.9198,77.6971) ..(286.9198,78.4179) .. controls (286.9198,79.1387) and (287.5041,79.7230) ..(288.2249,79.7230) .. controls (288.9457,79.7230) and (289.5300,79.1387) ..(289.5300,78.4179);
\path[fill=chfourfig5c0,fill opacity=0.5,draw=chfourfig5c0,line width=0.3985pt,draw opacity=0.5] (324.9576,42.0836) .. controls (324.9576,41.3628) and (324.3733,40.7785) ..(323.6525,40.7785) .. controls (322.9317,40.7785) and (322.3474,41.3628) ..(322.3474,42.0836) .. controls (322.3474,42.8044) and (322.9317,43.3888) ..(323.6525,43.3888) .. controls (324.3733,43.3888) and (324.9576,42.8044) ..(324.9576,42.0836);
\path[fill=chfourfig5c0,fill opacity=0.5,draw=chfourfig5c0,line width=0.3985pt,draw opacity=0.5] (272.9719,72.0317) .. controls (272.9719,71.3109) and (272.3876,70.7266) ..(271.6668,70.7266) .. controls (270.9460,70.7266) and (270.3617,71.3109) ..(270.3617,72.0317) .. controls (270.3617,72.7525) and (270.9460,73.3369) ..(271.6668,73.3369) .. controls (272.3876,73.3369) and (272.9719,72.7525) ..(272.9719,72.0317);
\path[fill=chfourfig5c0,fill opacity=0.5,draw=chfourfig5c0,line width=0.3985pt,draw opacity=0.5] (331.8618,45.4411) .. controls (331.8618,44.7203) and (331.2775,44.1360) ..(330.5567,44.1360) .. controls (329.8359,44.1360) and (329.2516,44.7203) ..(329.2516,45.4411) .. controls (329.2516,46.1619) and (329.8359,46.7462) ..(330.5567,46.7462) .. controls (331.2775,46.7462) and (331.8618,46.1619) ..(331.8618,45.4411);
\path[fill=chfourfig5c0,fill opacity=0.5,draw=chfourfig5c0,line width=0.3985pt,draw opacity=0.5] (297.0519,57.2470) .. controls (297.0519,56.5262) and (296.4676,55.9419) ..(295.7468,55.9419) .. controls (295.0260,55.9419) and (294.4417,56.5262) ..(294.4417,57.2470) .. controls (294.4417,57.9678) and (295.0260,58.5521) ..(295.7468,58.5521) .. controls (296.4676,58.5521) and (297.0519,57.9678) ..(297.0519,57.2470);
\path[fill=chfourfig5c0,fill opacity=0.5,draw=chfourfig5c0,line width=0.3985pt,draw opacity=0.5] (301.1068,50.0738) .. controls (301.1068,49.3530) and (300.5225,48.7687) ..(299.8016,48.7687) .. controls (299.0808,48.7687) and (298.4965,49.3530) ..(298.4965,50.0738) .. controls (298.4965,50.7946) and (299.0808,51.3789) ..(299.8016,51.3789) .. controls (300.5225,51.3789) and (301.1068,50.7946) ..(301.1068,50.0738);
\path[fill=chfourfig5c0,fill opacity=0.5,draw=chfourfig5c0,line width=0.3985pt,draw opacity=0.5] (271.2085,82.6421) .. controls (271.2085,81.9213) and (270.6242,81.3370) ..(269.9034,81.3370) .. controls (269.1826,81.3370) and (268.5982,81.9213) ..(268.5982,82.6421) .. controls (268.5982,83.3629) and (269.1826,83.9472) ..(269.9034,83.9472) .. controls (270.6242,83.9472) and (271.2085,83.3629) ..(271.2085,82.6421);
\path[fill=chfourfig5c0,fill opacity=0.5,draw=chfourfig5c0,line width=0.3985pt,draw opacity=0.5] (323.5429,61.2719) .. controls (323.5429,60.5511) and (322.9586,59.9668) ..(322.2378,59.9668) .. controls (321.5170,59.9668) and (320.9327,60.5511) ..(320.9327,61.2719) .. controls (320.9327,61.9928) and (321.5170,62.5771) ..(322.2378,62.5771) .. controls (322.9586,62.5771) and (323.5429,61.9928) ..(323.5429,61.2719);
\path[fill=chfourfig5c0,fill opacity=0.5,draw=chfourfig5c0,line width=0.3985pt,draw opacity=0.5] (265.1013,80.4702) .. controls (265.1013,79.7494) and (264.5170,79.1651) ..(263.7962,79.1651) .. controls (263.0754,79.1651) and (262.4911,79.7494) ..(262.4911,80.4702) .. controls (262.4911,81.1910) and (263.0754,81.7753) ..(263.7962,81.7753) .. controls (264.5170,81.7753) and (265.1013,81.1910) ..(265.1013,80.4702);
\path[fill=chfourfig5c0,fill opacity=0.5,draw=chfourfig5c0,line width=0.3985pt,draw opacity=0.5] (273.6593,86.8165) .. controls (273.6593,86.0957) and (273.0750,85.5114) ..(272.3542,85.5114) .. controls (271.6334,85.5114) and (271.0491,86.0957) ..(271.0491,86.8165) .. controls (271.0491,87.5373) and (271.6334,88.1216) ..(272.3542,88.1216) .. controls (273.0750,88.1216) and (273.6593,87.5373) ..(273.6593,86.8165);
\path[fill=chfourfig5c0,fill opacity=0.5,draw=chfourfig5c0,line width=0.3985pt,draw opacity=0.5] (333.0872,32.8581) .. controls (333.0872,32.1373) and (332.5029,31.5530) ..(331.7821,31.5530) .. controls (331.0613,31.5530) and (330.4770,32.1373) ..(330.4770,32.8581) .. controls (330.4770,33.5789) and (331.0613,34.1632) ..(331.7821,34.1632) .. controls (332.5029,34.1632) and (333.0872,33.5789) ..(333.0872,32.8581);
\path[fill=chfourfig5c0,fill opacity=0.5,draw=chfourfig5c0,line width=0.3985pt,draw opacity=0.5] (264.6231,99.5489) .. controls (264.6231,98.8281) and (264.0388,98.2438) ..(263.3180,98.2438) .. controls (262.5972,98.2438) and (262.0129,98.8281) ..(262.0129,99.5489) .. controls (262.0129,100.2697) and (262.5972,100.8540) ..(263.3180,100.8540) .. controls (264.0388,100.8540) and (264.6231,100.2697) ..(264.6231,99.5489);
\path[fill=chfourfig5c0,fill opacity=0.5,draw=chfourfig5c0,line width=0.3985pt,draw opacity=0.5] (267.0341,69.5012) .. controls (267.0341,68.7804) and (266.4498,68.1961) ..(265.7290,68.1961) .. controls (265.0082,68.1961) and (264.4239,68.7804) ..(264.4239,69.5012) .. controls (264.4239,70.2220) and (265.0082,70.8063) ..(265.7290,70.8063) .. controls (266.4498,70.8063) and (267.0341,70.2220) ..(267.0341,69.5012);
\path[fill=chfourfig5c0,fill opacity=0.5,draw=chfourfig5c0,line width=0.3985pt,draw opacity=0.5] (308.2401,65.4962) .. controls (308.2401,64.7753) and (307.6558,64.1910) ..(306.9350,64.1910) .. controls (306.2142,64.1910) and (305.6299,64.7753) ..(305.6299,65.4962) .. controls (305.6299,66.2169) and (306.2142,66.8013) ..(306.9350,66.8013) .. controls (307.6558,66.8013) and (308.2401,66.2169) ..(308.2401,65.4962);
\path[fill=chfourfig5c0,fill opacity=0.5,draw=chfourfig5c0,line width=0.3985pt,draw opacity=0.5] (248.2842,93.9199) .. controls (248.2842,93.1991) and (247.6999,92.6148) ..(246.9791,92.6148) .. controls (246.2582,92.6148) and (245.6740,93.1991) ..(245.6740,93.9199) .. controls (245.6740,94.6407) and (246.2582,95.2251) ..(246.9791,95.2251) .. controls (247.6999,95.2251) and (248.2842,94.6407) ..(248.2842,93.9199);
\path[fill=chfourfig5c0,fill opacity=0.5,draw=chfourfig5c0,line width=0.3985pt,draw opacity=0.5] (298.2973,53.0826) .. controls (298.2973,52.3617) and (297.7130,51.7774) ..(296.9921,51.7774) .. controls (296.2713,51.7774) and (295.6870,52.3617) ..(295.6870,53.0826) .. controls (295.6870,53.8033) and (296.2713,54.3877) ..(296.9921,54.3877) .. controls (297.7130,54.3877) and (298.2973,53.8033) ..(298.2973,53.0826);
\path[fill=chfourfig5c0,fill opacity=0.5,draw=chfourfig5c0,line width=0.3985pt,draw opacity=0.5] (258.5657,95.3645) .. controls (258.5657,94.6437) and (257.9814,94.0594) ..(257.2606,94.0594) .. controls (256.5398,94.0594) and (255.9555,94.6437) ..(255.9555,95.3645) .. controls (255.9555,96.0853) and (256.5398,96.6697) ..(257.2606,96.6697) .. controls (257.9814,96.6697) and (258.5657,96.0853) ..(258.5657,95.3645);
\path[fill=chfourfig5c0,fill opacity=0.5,draw=chfourfig5c0,line width=0.3985pt,draw opacity=0.5] (321.8891,33.5954) .. controls (321.8891,32.8746) and (321.3048,32.2903) ..(320.5840,32.2903) .. controls (319.8632,32.2903) and (319.2789,32.8746) ..(319.2789,33.5954) .. controls (319.2789,34.3162) and (319.8632,34.9005) ..(320.5840,34.9005) .. controls (321.3048,34.9005) and (321.8891,34.3162) ..(321.8891,33.5954);
\path[fill=chfourfig5c0,fill opacity=0.5,draw=chfourfig5c0,line width=0.3985pt,draw opacity=0.5] (249.9280,76.3556) .. controls (249.9280,75.6348) and (249.3437,75.0505) ..(248.6229,75.0505) .. controls (247.9021,75.0505) and (247.3178,75.6348) ..(247.3178,76.3556) .. controls (247.3178,77.0764) and (247.9021,77.6607) ..(248.6229,77.6607) .. controls (249.3437,77.6607) and (249.9280,77.0764) ..(249.9280,76.3556);
\path[fill=chfourfig5c0,fill opacity=0.5,draw=chfourfig5c0,line width=0.3985pt,draw opacity=0.5] (280.2348,73.5262) .. controls (280.2348,72.8053) and (279.6505,72.2210) ..(278.9296,72.2210) .. controls (278.2088,72.2210) and (277.6245,72.8053) ..(277.6245,73.5262) .. controls (277.6245,74.2469) and (278.2088,74.8313) ..(278.9296,74.8313) .. controls (279.6505,74.8313) and (280.2348,74.2469) ..(280.2348,73.5262);
\path[fill=chfourfig5c0,fill opacity=0.5,draw=chfourfig5c0,line width=0.3985pt,draw opacity=0.5] (291.5126,51.2892) .. controls (291.5126,50.5685) and (290.9283,49.9841) ..(290.2075,49.9841) .. controls (289.4867,49.9841) and (288.9024,50.5685) ..(288.9024,51.2892) .. controls (288.9024,52.0100) and (289.4867,52.5944) ..(290.2075,52.5944) .. controls (290.9283,52.5944) and (291.5126,52.0100) ..(291.5126,51.2892);
\path[fill=chfourfig5c0,fill opacity=0.5,draw=chfourfig5c0,line width=0.3985pt,draw opacity=0.5] (244.5979,73.1675) .. controls (244.5979,72.4467) and (244.0136,71.8624) ..(243.2928,71.8624) .. controls (242.5720,71.8624) and (241.9877,72.4467) ..(241.9877,73.1675) .. controls (241.9877,73.8883) and (242.5720,74.4726) ..(243.2928,74.4726) .. controls (244.0136,74.4726) and (244.5979,73.8883) ..(244.5979,73.1675);
\path[fill=chfourfig5c0,fill opacity=0.5,draw=chfourfig5c0,line width=0.3985pt,draw opacity=0.5] (329.0922,53.2818) .. controls (329.0922,52.5610) and (328.5079,51.9767) ..(327.7871,51.9767) .. controls (327.0663,51.9767) and (326.4819,52.5610) ..(326.4819,53.2818) .. controls (326.4819,54.0026) and (327.0663,54.5869) ..(327.7871,54.5869) .. controls (328.5079,54.5869) and (329.0922,54.0026) ..(329.0922,53.2818);
\path[fill=chfourfig5c0,fill opacity=0.5,draw=chfourfig5c0,line width=0.3985pt,draw opacity=0.5] (241.9279,80.3606) .. controls (241.9279,79.6398) and (241.3436,79.0555) ..(240.6228,79.0555) .. controls (239.9020,79.0555) and (239.3177,79.6398) ..(239.3177,80.3606) .. controls (239.3177,81.0814) and (239.9020,81.6657) ..(240.6228,81.6657) .. controls (241.3436,81.6657) and (241.9279,81.0814) ..(241.9279,80.3606);
\path[fill=chfourfig5c0,fill opacity=0.5,draw=chfourfig5c0,line width=0.3985pt,draw opacity=0.5] (284.7081,66.8511) .. controls (284.7081,66.1303) and (284.1237,65.5460) ..(283.4029,65.5460) .. controls (282.6821,65.5460) and (282.0978,66.1303) ..(282.0978,66.8511) .. controls (282.0978,67.5719) and (282.6821,68.1562) ..(283.4029,68.1562) .. controls (284.1237,68.1562) and (284.7081,67.5719) ..(284.7081,66.8511);
\path[fill=chfourfig5c0,fill opacity=0.5,draw=chfourfig5c0,line width=0.3985pt,draw opacity=0.5] (261.4151,77.3718) .. controls (261.4151,76.6510) and (260.8308,76.0667) ..(260.1100,76.0667) .. controls (259.3892,76.0667) and (258.8049,76.6510) ..(258.8049,77.3718) .. controls (258.8049,78.0926) and (259.3892,78.6769) ..(260.1100,78.6769) .. controls (260.8308,78.6769) and (261.4151,78.0926) ..(261.4151,77.3718);
\path[fill=chfourfig5c0,fill opacity=0.5,draw=chfourfig5c0,line width=0.3985pt,draw opacity=0.5] (275.3331,57.2868) .. controls (275.3331,56.5660) and (274.7488,55.9817) ..(274.0280,55.9817) .. controls (273.3072,55.9817) and (272.7228,56.5660) ..(272.7228,57.2868) .. controls (272.7228,58.0076) and (273.3072,58.5919) ..(274.0280,58.5919) .. controls (274.7488,58.5919) and (275.3331,58.0076) ..(275.3331,57.2868);
\path[fill=chfourfig5c0,fill opacity=0.5,draw=chfourfig5c0,line width=0.3985pt,draw opacity=0.5] (253.9231,80.6694) .. controls (253.9231,79.9487) and (253.3388,79.3643) ..(252.6180,79.3643) .. controls (251.8972,79.3643) and (251.3129,79.9487) ..(251.3129,80.6694) .. controls (251.3129,81.3903) and (251.8972,81.9746) ..(252.6180,81.9746) .. controls (253.3388,81.9746) and (253.9231,81.3903) ..(253.9231,80.6694);
\path[fill=chfourfig5c0,fill opacity=0.5,draw=chfourfig5c0,line width=0.3985pt,draw opacity=0.5] (270.4713,80.4104) .. controls (270.4713,79.6896) and (269.8869,79.1053) ..(269.1661,79.1053) .. controls (268.4453,79.1053) and (267.8610,79.6896) ..(267.8610,80.4104) .. controls (267.8610,81.1312) and (268.4453,81.7155) ..(269.1661,81.7155) .. controls (269.8869,81.7155) and (270.4713,81.1312) ..(270.4713,80.4104);
\path[fill=chfourfig5c0,fill opacity=0.5,draw=chfourfig5c0,line width=0.3985pt,draw opacity=0.5] (334.7311,31.8220) .. controls (334.7311,31.1012) and (334.1468,30.5169) ..(333.4260,30.5169) .. controls (332.7052,30.5169) and (332.1209,31.1012) ..(332.1209,31.8220) .. controls (332.1209,32.5428) and (332.7052,33.1271) ..(333.4260,33.1271) .. controls (334.1468,33.1271) and (334.7311,32.5428) ..(334.7311,31.8220);
\path[fill=chfourfig5c0,fill opacity=0.5,draw=chfourfig5c0,line width=0.3985pt,draw opacity=0.5] (269.9034,66.7515) .. controls (269.9034,66.0307) and (269.3191,65.4464) ..(268.5982,65.4464) .. controls (267.8775,65.4464) and (267.2932,66.0307) ..(267.2932,66.7515) .. controls (267.2932,67.4723) and (267.8775,68.0566) ..(268.5982,68.0566) .. controls (269.3191,68.0566) and (269.9034,67.4723) ..(269.9034,66.7515);
\path[fill=chfourfig5c0,fill opacity=0.5,draw=chfourfig5c0,line width=0.3985pt,draw opacity=0.5] (295.0395,69.9495) .. controls (295.0395,69.2287) and (294.4551,68.6444) ..(293.7343,68.6444) .. controls (293.0135,68.6444) and (292.4292,69.2287) ..(292.4292,69.9495) .. controls (292.4292,70.6703) and (293.0135,71.2546) ..(293.7343,71.2546) .. controls (294.4551,71.2546) and (295.0395,70.6703) ..(295.0395,69.9495);
\path[fill=chfourfig5c0,fill opacity=0.5,draw=chfourfig5c0,line width=0.3985pt,draw opacity=0.5] (320.7633,32.6987) .. controls (320.7633,31.9779) and (320.1790,31.3936) ..(319.4582,31.3936) .. controls (318.7374,31.3936) and (318.1531,31.9779) ..(318.1531,32.6987) .. controls (318.1531,33.4195) and (318.7374,34.0038) ..(319.4582,34.0038) .. controls (320.1790,34.0038) and (320.7633,33.4195) ..(320.7633,32.6987);
\path[fill=chfourfig5c0,fill opacity=0.5,draw=chfourfig5c0,line width=0.3985pt,draw opacity=0.5] (268.3592,93.0233) .. controls (268.3592,92.3025) and (267.7748,91.7182) ..(267.0540,91.7182) .. controls (266.3332,91.7182) and (265.7489,92.3025) ..(265.7489,93.0233) .. controls (265.7489,93.7441) and (266.3332,94.3284) ..(267.0540,94.3284) .. controls (267.7748,94.3284) and (268.3592,93.7441) ..(268.3592,93.0233);
\path[fill=chfourfig5c0,fill opacity=0.5,draw=chfourfig5c0,line width=0.3985pt,draw opacity=0.5] (279.7665,65.7452) .. controls (279.7665,65.0244) and (279.1822,64.4401) ..(278.4614,64.4401) .. controls (277.7406,64.4401) and (277.1563,65.0244) ..(277.1563,65.7452) .. controls (277.1563,66.4660) and (277.7406,67.0503) ..(278.4614,67.0503) .. controls (279.1822,67.0503) and (279.7665,66.4660) ..(279.7665,65.7452);
\path[fill=chfourfig5c0,fill opacity=0.5,draw=chfourfig5c0,line width=0.3985pt,draw opacity=0.5] (248.0052,90.7119) .. controls (248.0052,89.9911) and (247.4209,89.4068) ..(246.7001,89.4068) .. controls (245.9793,89.4068) and (245.3950,89.9911) ..(245.3950,90.7119) .. controls (245.3950,91.4327) and (245.9793,92.0170) ..(246.7001,92.0170) .. controls (247.4209,92.0170) and (248.0052,91.4327) ..(248.0052,90.7119);
\path[fill=chfourfig5c0,fill opacity=0.5,draw=chfourfig5c0,line width=0.3985pt,draw opacity=0.5] (231.2877,89.3171) .. controls (231.2877,88.5963) and (230.7034,88.0120) ..(229.9825,88.0120) .. controls (229.2617,88.0120) and (228.6774,88.5963) ..(228.6774,89.3171) .. controls (228.6774,90.0379) and (229.2617,90.6223) ..(229.9825,90.6223) .. controls (230.7034,90.6223) and (231.2877,90.0379) ..(231.2877,89.3171);
\path[fill=chfourfig5c0,fill opacity=0.5,draw=chfourfig5c0,line width=0.3985pt,draw opacity=0.5] (265.8884,78.8562) .. controls (265.8884,78.1354) and (265.3041,77.5511) ..(264.5833,77.5511) .. controls (263.8625,77.5511) and (263.2781,78.1354) ..(263.2781,78.8562) .. controls (263.2781,79.5770) and (263.8625,80.1613) ..(264.5833,80.1613) .. controls (265.3041,80.1613) and (265.8884,79.5770) ..(265.8884,78.8562);
\path[fill=chfourfig5c0,fill opacity=0.5,draw=chfourfig5c0,line width=0.3985pt,draw opacity=0.5] (312.8628,44.3950) .. controls (312.8628,43.6742) and (312.2785,43.0899) ..(311.5577,43.0899) .. controls (310.8369,43.0899) and (310.2526,43.6742) ..(310.2526,44.3950) .. controls (310.2526,45.1158) and (310.8369,45.7001) ..(311.5577,45.7001) .. controls (312.2785,45.7001) and (312.8628,45.1158) ..(312.8628,44.3950);
\path[fill=chfourfig5c0,fill opacity=0.5,draw=chfourfig5c0,line width=0.3985pt,draw opacity=0.5] (242.3065,108.1069) .. controls (242.3065,107.3861) and (241.7222,106.8018) ..(241.0014,106.8018) .. controls (240.2806,106.8018) and (239.6963,107.3861) ..(239.6963,108.1069) .. controls (239.6963,108.8277) and (240.2806,109.4120) ..(241.0014,109.4120) .. controls (241.7222,109.4120) and (242.3065,108.8277) ..(242.3065,108.1069);
\path[fill=chfourfig5c0,fill opacity=0.5,draw=chfourfig5c0,line width=0.3985pt,draw opacity=0.5] (249.9778,79.0057) .. controls (249.9778,78.2849) and (249.3935,77.7006) ..(248.6727,77.7006) .. controls (247.9519,77.7006) and (247.3676,78.2849) ..(247.3676,79.0057) .. controls (247.3676,79.7265) and (247.9519,80.3108) ..(248.6727,80.3108) .. controls (249.3935,80.3108) and (249.9778,79.7265) ..(249.9778,79.0057);
\path[fill=chfourfig5c0,fill opacity=0.5,draw=chfourfig5c0,line width=0.3985pt,draw opacity=0.5] (274.8150,87.3943) .. controls (274.8150,86.6735) and (274.2307,86.0892) ..(273.5099,86.0892) .. controls (272.7891,86.0892) and (272.2048,86.6735) ..(272.2048,87.3943) .. controls (272.2048,88.1151) and (272.7891,88.6994) ..(273.5099,88.6994) .. controls (274.2307,88.6994) and (274.8150,88.1151) ..(274.8150,87.3943);
\path[fill=chfourfig5c0,fill opacity=0.5,draw=chfourfig5c0,line width=0.3985pt,draw opacity=0.5] (262.4712,85.8600) .. controls (262.4712,85.1393) and (261.8868,84.5549) ..(261.1660,84.5549) .. controls (260.4452,84.5549) and (259.8609,85.1393) ..(259.8609,85.8600) .. controls (259.8609,86.5809) and (260.4452,87.1652) ..(261.1660,87.1652) .. controls (261.8868,87.1652) and (262.4712,86.5809) ..(262.4712,85.8600);
\path[fill=chfourfig5c0,fill opacity=0.5,draw=chfourfig5c0,line width=0.3985pt,draw opacity=0.5] (237.2753,93.4915) .. controls (237.2753,92.7707) and (236.6910,92.1864) ..(235.9702,92.1864) .. controls (235.2494,92.1864) and (234.6651,92.7707) ..(234.6651,93.4915) .. controls (234.6651,94.2123) and (235.2494,94.7967) ..(235.9702,94.7967) .. controls (236.6910,94.7967) and (237.2753,94.2123) ..(237.2753,93.4915);
\path[fill=chfourfig5c0,fill opacity=0.5,draw=chfourfig5c0,line width=0.3985pt,draw opacity=0.5] (286.3519,52.1560) .. controls (286.3519,51.4352) and (285.7676,50.8509) ..(285.0468,50.8509) .. controls (284.3260,50.8509) and (283.7417,51.4352) ..(283.7417,52.1560) .. controls (283.7417,52.8768) and (284.3260,53.4611) ..(285.0468,53.4611) .. controls (285.7676,53.4611) and (286.3519,52.8768) ..(286.3519,52.1560);
\path[fill=chfourfig5c0,fill opacity=0.5,draw=chfourfig5c0,line width=0.3985pt,draw opacity=0.5] (234.8643,112.6798) .. controls (234.8643,111.9590) and (234.2800,111.3747) ..(233.5592,111.3747) .. controls (232.8384,111.3747) and (232.2541,111.9590) ..(232.2541,112.6798) .. controls (232.2541,113.4006) and (232.8384,113.9849) ..(233.5592,113.9849) .. controls (234.2800,113.9849) and (234.8643,113.4006) ..(234.8643,112.6798);
\path[fill=chfourfig5c0,fill opacity=0.5,draw=chfourfig5c0,line width=0.3985pt,draw opacity=0.5] (277.3954,74.3829) .. controls (277.3954,73.6622) and (276.8110,73.0778) ..(276.0903,73.0778) .. controls (275.3694,73.0778) and (274.7851,73.6622) ..(274.7851,74.3829) .. controls (274.7851,75.1038) and (275.3694,75.6881) ..(276.0903,75.6881) .. controls (276.8110,75.6881) and (277.3954,75.1038) ..(277.3954,74.3829);
\path[fill=chfourfig5c0,fill opacity=0.5,draw=chfourfig5c0,line width=0.3985pt,draw opacity=0.5] (236.9366,101.5116) .. controls (236.9366,100.7908) and (236.3523,100.2065) ..(235.6315,100.2065) .. controls (234.9106,100.2065) and (234.3264,100.7908) ..(234.3264,101.5116) .. controls (234.3264,102.2324) and (234.9106,102.8167) ..(235.6315,102.8167) .. controls (236.3523,102.8167) and (236.9366,102.2324) ..(236.9366,101.5116);
\path[fill=chfourfig5c0,fill opacity=0.5,draw=chfourfig5c0,line width=0.3985pt,draw opacity=0.5] (278.2024,70.9557) .. controls (278.2024,70.2350) and (277.6180,69.6506) ..(276.8972,69.6506) .. controls (276.1764,69.6506) and (275.5921,70.2350) ..(275.5921,70.9557) .. controls (275.5921,71.6766) and (276.1764,72.2609) ..(276.8972,72.2609) .. controls (277.6180,72.2609) and (278.2024,71.6766) ..(278.2024,70.9557);
\path[fill=chfourfig5c0,fill opacity=0.5,draw=chfourfig5c0,line width=0.3985pt,draw opacity=0.5] (240.8619,89.2872) .. controls (240.8619,88.5664) and (240.2776,87.9821) ..(239.5568,87.9821) .. controls (238.8360,87.9821) and (238.2516,88.5664) ..(238.2516,89.2872) .. controls (238.2516,90.0080) and (238.8360,90.5924) ..(239.5568,90.5924) .. controls (240.2776,90.5924) and (240.8619,90.0080) ..(240.8619,89.2872);
\path[fill=chfourfig5c0,fill opacity=0.5,draw=chfourfig5c0,line width=0.3985pt,draw opacity=0.5] (239.8457,100.3061) .. controls (239.8457,99.5853) and (239.2614,99.0009) ..(238.5406,99.0009) .. controls (237.8198,99.0009) and (237.2355,99.5853) ..(237.2355,100.3061) .. controls (237.2355,101.0269) and (237.8198,101.6112) ..(238.5406,101.6112) .. controls (239.2614,101.6112) and (239.8457,101.0269) ..(239.8457,100.3061);
\path[fill=chfourfig5c0,fill opacity=0.5,draw=chfourfig5c0,line width=0.3985pt,draw opacity=0.5] (328.0062,61.7900) .. controls (328.0062,61.0692) and (327.4219,60.4849) ..(326.7011,60.4849) .. controls (325.9803,60.4849) and (325.3960,61.0692) ..(325.3960,61.7900) .. controls (325.3960,62.5108) and (325.9803,63.0951) ..(326.7011,63.0951) .. controls (327.4219,63.0951) and (328.0062,62.5108) ..(328.0062,61.7900);
\path[fill=chfourfig5c0,fill opacity=0.5,draw=chfourfig5c0,line width=0.3985pt,draw opacity=0.5] (327.5380,44.0463) .. controls (327.5380,43.3255) and (326.9537,42.7412) ..(326.2329,42.7412) .. controls (325.5121,42.7412) and (324.9277,43.3255) ..(324.9277,44.0463) .. controls (324.9277,44.7671) and (325.5121,45.3514) ..(326.2329,45.3514) .. controls (326.9537,45.3514) and (327.5380,44.7671) ..(327.5380,44.0463);
\path[fill=chfourfig5c0,fill opacity=0.5,draw=chfourfig5c0,line width=0.3985pt,draw opacity=0.5] (253.9231,71.0255) .. controls (253.9231,70.3047) and (253.3388,69.7204) ..(252.6180,69.7204) .. controls (251.8972,69.7204) and (251.3129,70.3047) ..(251.3129,71.0255) .. controls (251.3129,71.7463) and (251.8972,72.3306) ..(252.6180,72.3306) .. controls (253.3388,72.3306) and (253.9231,71.7463) ..(253.9231,71.0255);
\path[fill=chfourfig5c0,fill opacity=0.5,draw=chfourfig5c0,line width=0.3985pt,draw opacity=0.5] (296.0556,45.6404) .. controls (296.0556,44.9196) and (295.4713,44.3352) ..(294.7505,44.3352) .. controls (294.0297,44.3352) and (293.4454,44.9196) ..(293.4454,45.6404) .. controls (293.4454,46.3612) and (294.0297,46.9455) ..(294.7505,46.9455) .. controls (295.4713,46.9455) and (296.0556,46.3612) ..(296.0556,45.6404);
\path[fill=chfourfig5c0,fill opacity=0.5,draw=chfourfig5c0,line width=0.3985pt,draw opacity=0.5] (328.5243,52.4150) .. controls (328.5243,51.6942) and (327.9400,51.1099) ..(327.2192,51.1099) .. controls (326.4984,51.1099) and (325.9141,51.6942) ..(325.9141,52.4150) .. controls (325.9141,53.1358) and (326.4984,53.7202) ..(327.2192,53.7202) .. controls (327.9400,53.7202) and (328.5243,53.1358) ..(328.5243,52.4150);
\path[fill=chfourfig5c0,fill opacity=0.5,draw=chfourfig5c0,line width=0.3985pt,draw opacity=0.5] (238.5904,101.4319) .. controls (238.5904,100.7111) and (238.0061,100.1267) ..(237.2853,100.1267) .. controls (236.5645,100.1267) and (235.9802,100.7111) ..(235.9802,101.4319) .. controls (235.9802,102.1527) and (236.5645,102.7370) ..(237.2853,102.7370) .. controls (238.0061,102.7370) and (238.5904,102.1527) ..(238.5904,101.4319);
\path[fill=chfourfig5c0,fill opacity=0.5,draw=chfourfig5c0,line width=0.3985pt,draw opacity=0.5] (306.5166,52.6541) .. controls (306.5166,51.9333) and (305.9323,51.3490) ..(305.2114,51.3490) .. controls (304.4906,51.3490) and (303.9063,51.9333) ..(303.9063,52.6541) .. controls (303.9063,53.3749) and (304.4906,53.9593) ..(305.2114,53.9593) .. controls (305.9323,53.9593) and (306.5166,53.3749) ..(306.5166,52.6541);
\path[fill=chfourfig5c0,fill opacity=0.5,draw=chfourfig5c0,line width=0.3985pt,draw opacity=0.5] (260.1996,88.0021) .. controls (260.1996,87.2812) and (259.6153,86.6969) ..(258.8945,86.6969) .. controls (258.1737,86.6969) and (257.5894,87.2812) ..(257.5894,88.0021) .. controls (257.5894,88.7229) and (258.1737,89.3072) ..(258.8945,89.3072) .. controls (259.6153,89.3072) and (260.1996,88.7229) ..(260.1996,88.0021);
\path[fill=chfourfig5c0,fill opacity=0.5,draw=chfourfig5c0,line width=0.3985pt,draw opacity=0.5] (286.4217,68.3555) .. controls (286.4217,67.6347) and (285.8373,67.0504) ..(285.1165,67.0504) .. controls (284.3957,67.0504) and (283.8114,67.6347) ..(283.8114,68.3555) .. controls (283.8114,69.0763) and (284.3957,69.6606) ..(285.1165,69.6606) .. controls (285.8373,69.6606) and (286.4217,69.0763) ..(286.4217,68.3555);
\path[fill=chfourfig5c0,fill opacity=0.5,draw=chfourfig5c0,line width=0.3985pt,draw opacity=0.5] (293.6646,62.5372) .. controls (293.6646,61.8164) and (293.0803,61.2321) ..(292.3594,61.2321) .. controls (291.6387,61.2321) and (291.0544,61.8164) ..(291.0544,62.5372) .. controls (291.0544,63.2580) and (291.6387,63.8423) ..(292.3594,63.8423) .. controls (293.0803,63.8423) and (293.6646,63.2580) ..(293.6646,62.5372);
\path[fill=chfourfig5c0,fill opacity=0.5,draw=chfourfig5c0,line width=0.3985pt,draw opacity=0.5] (265.5496,69.2920) .. controls (265.5496,68.5712) and (264.9653,67.9869) ..(264.2445,67.9869) .. controls (263.5237,67.9869) and (262.9394,68.5712) ..(262.9394,69.2920) .. controls (262.9394,70.0128) and (263.5237,70.5971) ..(264.2445,70.5971) .. controls (264.9653,70.5971) and (265.5496,70.0128) ..(265.5496,69.2920);
\path[fill=chfourfig5c0,fill opacity=0.5,draw=chfourfig5c0,line width=0.3985pt,draw opacity=0.5] (330.8257,35.2990) .. controls (330.8257,34.5782) and (330.2414,33.9939) ..(329.5206,33.9939) .. controls (328.7998,33.9939) and (328.2155,34.5782) ..(328.2155,35.2990) .. controls (328.2155,36.0198) and (328.7998,36.6041) ..(329.5206,36.6041) .. controls (330.2414,36.6041) and (330.8257,36.0198) ..(330.8257,35.2990);
\path[fill=chfourfig5c0,fill opacity=0.5,draw=chfourfig5c0,line width=0.3985pt,draw opacity=0.5] (264.5833,71.3244) .. controls (264.5833,70.6036) and (263.9989,70.0193) ..(263.2781,70.0193) .. controls (262.5573,70.0193) and (261.9730,70.6036) ..(261.9730,71.3244) .. controls (261.9730,72.0452) and (262.5573,72.6295) ..(263.2781,72.6295) .. controls (263.9989,72.6295) and (264.5833,72.0452) ..(264.5833,71.3244);
\path[fill=chfourfig5c0,fill opacity=0.5,draw=chfourfig5c0,line width=0.3985pt,draw opacity=0.5] (333.0474,40.4896) .. controls (333.0474,39.7688) and (332.4631,39.1845) ..(331.7423,39.1845) .. controls (331.0215,39.1845) and (330.4371,39.7688) ..(330.4371,40.4896) .. controls (330.4371,41.2104) and (331.0215,41.7947) ..(331.7423,41.7947) .. controls (332.4631,41.7947) and (333.0474,41.2104) ..(333.0474,40.4896);
\path[fill=chfourfig5c0,fill opacity=0.5,draw=chfourfig5c0,line width=0.3985pt,draw opacity=0.5] (240.2940,103.9624) .. controls (240.2940,103.2416) and (239.7097,102.6573) ..(238.9889,102.6573) .. controls (238.2681,102.6573) and (237.6838,103.2416) ..(237.6838,103.9624) .. controls (237.6838,104.6832) and (238.2681,105.2675) ..(238.9889,105.2675) .. controls (239.7097,105.2675) and (240.2940,104.6832) ..(240.2940,103.9624);
\path[fill=chfourfig5c0,fill opacity=0.5,draw=chfourfig5c0,line width=0.3985pt,draw opacity=0.5] (232.3338,116.9240) .. controls (232.3338,116.2032) and (231.7495,115.6189) ..(231.0287,115.6189) .. controls (230.3079,115.6189) and (229.7235,116.2032) ..(229.7235,116.9240) .. controls (229.7235,117.6448) and (230.3079,118.2291) ..(231.0287,118.2291) .. controls (231.7495,118.2291) and (232.3338,117.6448) ..(232.3338,116.9240);
\path[fill=chfourfig5c0,fill opacity=0.5,draw=chfourfig5c0,line width=0.3985pt,draw opacity=0.5] (290.3271,70.4178) .. controls (290.3271,69.6970) and (289.7427,69.1126) ..(289.0219,69.1126) .. controls (288.3011,69.1126) and (287.7168,69.6970) ..(287.7168,70.4178) .. controls (287.7168,71.1386) and (288.3011,71.7229) ..(289.0219,71.7229) .. controls (289.7427,71.7229) and (290.3271,71.1386) ..(290.3271,70.4178);
\path[fill=chfourfig5c0,fill opacity=0.5,draw=chfourfig5c0,line width=0.3985pt,draw opacity=0.5] (248.5233,86.8065) .. controls (248.5233,86.0857) and (247.9390,85.5014) ..(247.2181,85.5014) .. controls (246.4973,85.5014) and (245.9130,86.0857) ..(245.9130,86.8065) .. controls (245.9130,87.5273) and (246.4973,88.1116) ..(247.2181,88.1116) .. controls (247.9390,88.1116) and (248.5233,87.5273) ..(248.5233,86.8065);
\path[fill=chfourfig5c0,fill opacity=0.5,draw=chfourfig5c0,line width=0.3985pt,draw opacity=0.5] (301.0271,60.8535) .. controls (301.0271,60.1327) and (300.4427,59.5484) ..(299.7219,59.5484) .. controls (299.0011,59.5484) and (298.4168,60.1327) ..(298.4168,60.8535) .. controls (298.4168,61.5743) and (299.0011,62.1586) ..(299.7219,62.1586) .. controls (300.4427,62.1586) and (301.0271,61.5743) ..(301.0271,60.8535);
\path[fill=chfourfig5c0,fill opacity=0.5,draw=chfourfig5c0,line width=0.3985pt,draw opacity=0.5] (278.2123,74.1239) .. controls (278.2123,73.4031) and (277.6280,72.8188) ..(276.9072,72.8188) .. controls (276.1864,72.8188) and (275.6021,73.4031) ..(275.6021,74.1239) .. controls (275.6021,74.8447) and (276.1864,75.4290) ..(276.9072,75.4290) .. controls (277.6280,75.4290) and (278.2123,74.8447) ..(278.2123,74.1239);
\path[fill=chfourfig5c0,fill opacity=0.5,draw=chfourfig5c0,line width=0.3985pt,draw opacity=0.5] (283.9907,80.4602) .. controls (283.9907,79.7394) and (283.4064,79.1551) ..(282.6856,79.1551) .. controls (281.9648,79.1551) and (281.3805,79.7394) ..(281.3805,80.4602) .. controls (281.3805,81.1810) and (281.9648,81.7654) ..(282.6856,81.7654) .. controls (283.4064,81.7654) and (283.9907,81.1810) ..(283.9907,80.4602);
\path[fill=chfourfig5c0,fill opacity=0.5,draw=chfourfig5c0,line width=0.3985pt,draw opacity=0.5] (334.0935,30.8855) .. controls (334.0935,30.1647) and (333.5092,29.5804) ..(332.7884,29.5804) .. controls (332.0676,29.5804) and (331.4832,30.1647) ..(331.4832,30.8855) .. controls (331.4832,31.6063) and (332.0676,32.1906) ..(332.7884,32.1906) .. controls (333.5092,32.1906) and (334.0935,31.6063) ..(334.0935,30.8855);
\path[fill=chfourfig5c0,fill opacity=0.5,draw=chfourfig5c0,line width=0.3985pt,draw opacity=0.5] (310.5116,57.8946) .. controls (310.5116,57.1738) and (309.9273,56.5894) ..(309.2065,56.5894) .. controls (308.4857,56.5894) and (307.9014,57.1738) ..(307.9014,57.8946) .. controls (307.9014,58.6154) and (308.4857,59.1997) ..(309.2065,59.1997) .. controls (309.9273,59.1997) and (310.5116,58.6154) ..(310.5116,57.8946);
\path[fill=chfourfig5c0,fill opacity=0.5,draw=chfourfig5c0,line width=0.3985pt,draw opacity=0.5] (320.7932,48.5993) .. controls (320.7932,47.8785) and (320.2089,47.2942) ..(319.4881,47.2942) .. controls (318.7673,47.2942) and (318.1830,47.8785) ..(318.1830,48.5993) .. controls (318.1830,49.3201) and (318.7673,49.9044) ..(319.4881,49.9044) .. controls (320.2089,49.9044) and (320.7932,49.3201) ..(320.7932,48.5993);
\path[fill=chfourfig5c0,fill opacity=0.5,draw=chfourfig5c0,line width=0.3985pt,draw opacity=0.5] (302.6211,44.9729) .. controls (302.6211,44.2521) and (302.0368,43.6677) ..(301.3160,43.6677) .. controls (300.5952,43.6677) and (300.0109,44.2521) ..(300.0109,44.9729) .. controls (300.0109,45.6936) and (300.5952,46.2780) ..(301.3160,46.2780) .. controls (302.0368,46.2780) and (302.6211,45.6936) ..(302.6211,44.9729);
\path[fill=chfourfig5c0,fill opacity=0.5,draw=chfourfig5c0,line width=0.3985pt,draw opacity=0.5] (266.2968,63.2844) .. controls (266.2968,62.5636) and (265.7125,61.9793) ..(264.9917,61.9793) .. controls (264.2709,61.9793) and (263.6866,62.5636) ..(263.6866,63.2844) .. controls (263.6866,64.0052) and (264.2709,64.5896) ..(264.9917,64.5896) .. controls (265.7125,64.5896) and (266.2968,64.0052) ..(266.2968,63.2844);
\path[fill=chfourfig5c0,fill opacity=0.5,draw=chfourfig5c0,line width=0.3985pt,draw opacity=0.5] (336.0861,34.2130) .. controls (336.0861,33.4923) and (335.5017,32.9080) ..(334.7809,32.9080) .. controls (334.0601,32.9080) and (333.4758,33.4923) ..(333.4758,34.2130) .. controls (333.4758,34.9339) and (334.0601,35.5182) ..(334.7809,35.5182) .. controls (335.5017,35.5182) and (336.0861,34.9339) ..(336.0861,34.2130);
\path[fill=chfourfig5c0,fill opacity=0.5,draw=chfourfig5c0,line width=0.3985pt,draw opacity=0.5] (304.4543,53.2619) .. controls (304.4543,52.5411) and (303.8699,51.9568) ..(303.1491,51.9568) .. controls (302.4283,51.9568) and (301.8440,52.5411) ..(301.8440,53.2619) .. controls (301.8440,53.9827) and (302.4283,54.5670) ..(303.1491,54.5670) .. controls (303.8699,54.5670) and (304.4543,53.9827) ..(304.4543,53.2619);
\path[fill=chfourfig5c0,fill opacity=0.5,draw=chfourfig5c0,line width=0.3985pt,draw opacity=0.5] (261.6442,68.6544) .. controls (261.6442,67.9336) and (261.0599,67.3492) ..(260.3391,67.3492) .. controls (259.6183,67.3492) and (259.0340,67.9336) ..(259.0340,68.6544) .. controls (259.0340,69.3752) and (259.6183,69.9595) ..(260.3391,69.9595) .. controls (261.0599,69.9595) and (261.6442,69.3752) ..(261.6442,68.6544);
\path[fill=chfourfig5c0,fill opacity=0.5,draw=chfourfig5c0,line width=0.3985pt,draw opacity=0.5] (256.2943,73.3767) .. controls (256.2943,72.6559) and (255.7099,72.0716) ..(254.9891,72.0716) .. controls (254.2683,72.0716) and (253.6840,72.6559) ..(253.6840,73.3767) .. controls (253.6840,74.0975) and (254.2683,74.6818) ..(254.9891,74.6818) .. controls (255.7099,74.6818) and (256.2943,74.0975) ..(256.2943,73.3767);
\path[fill=chfourfig5c0,fill opacity=0.5,draw=chfourfig5c0,line width=0.3985pt,draw opacity=0.5] (258.7152,98.6124) .. controls (258.7152,97.8916) and (258.1309,97.3073) ..(257.4101,97.3073) .. controls (256.6893,97.3073) and (256.1049,97.8916) ..(256.1049,98.6124) .. controls (256.1049,99.3332) and (256.6893,99.9175) ..(257.4101,99.9175) .. controls (258.1309,99.9175) and (258.7152,99.3332) ..(258.7152,98.6124);
\path[fill=chfourfig5c0,fill opacity=0.5,draw=chfourfig5c0,line width=0.3985pt,draw opacity=0.5] (255.5969,93.8801) .. controls (255.5969,93.1593) and (255.0125,92.5750) ..(254.2917,92.5750) .. controls (253.5709,92.5750) and (252.9866,93.1593) ..(252.9866,93.8801) .. controls (252.9866,94.6009) and (253.5709,95.1852) ..(254.2917,95.1852) .. controls (255.0125,95.1852) and (255.5969,94.6009) ..(255.5969,93.8801);
\path[fill=chfourfig5c0,fill opacity=0.5,draw=chfourfig5c0,line width=0.3985pt,draw opacity=0.5] (249.5893,72.1413) .. controls (249.5893,71.4205) and (249.0050,70.8362) ..(248.2842,70.8362) .. controls (247.5634,70.8362) and (246.9791,71.4205) ..(246.9791,72.1413) .. controls (246.9791,72.8621) and (247.5634,73.4464) ..(248.2842,73.4464) .. controls (249.0050,73.4464) and (249.5893,72.8621) ..(249.5893,72.1413);
\path[fill=chfourfig5c0,fill opacity=0.5,draw=chfourfig5c0,line width=0.3985pt,draw opacity=0.5] (316.8878,46.8857) .. controls (316.8878,46.1649) and (316.3035,45.5806) ..(315.5827,45.5806) .. controls (314.8619,45.5806) and (314.2776,46.1649) ..(314.2776,46.8857) .. controls (314.2776,47.6065) and (314.8619,48.1908) ..(315.5827,48.1908) .. controls (316.3035,48.1908) and (316.8878,47.6065) ..(316.8878,46.8857);
\path[fill=chfourfig5c0,fill opacity=0.5,draw=chfourfig5c0,line width=0.3985pt,draw opacity=0.5] (232.1046,91.8278) .. controls (232.1046,91.1069) and (231.5203,90.5226) ..(230.7995,90.5226) .. controls (230.0787,90.5226) and (229.4944,91.1069) ..(229.4944,91.8278) .. controls (229.4944,92.5486) and (230.0787,93.1329) ..(230.7995,93.1329) .. controls (231.5203,93.1329) and (232.1046,92.5486) ..(232.1046,91.8278);
\path[fill=chfourfig5c0,fill opacity=0.5,draw=chfourfig5c0,line width=0.3985pt,draw opacity=0.5] (249.0613,84.2760) .. controls (249.0613,83.5552) and (248.4769,82.9709) ..(247.7562,82.9709) .. controls (247.0353,82.9709) and (246.4510,83.5552) ..(246.4510,84.2760) .. controls (246.4510,84.9968) and (247.0353,85.5811) ..(247.7562,85.5811) .. controls (248.4769,85.5811) and (249.0613,84.9968) ..(249.0613,84.2760);
\path[fill=chfourfig5c0,fill opacity=0.5,draw=chfourfig5c0,line width=0.3985pt,draw opacity=0.5] (259.9805,84.2461) .. controls (259.9805,83.5253) and (259.3961,82.9410) ..(258.6754,82.9410) .. controls (257.9545,82.9410) and (257.3702,83.5253) ..(257.3702,84.2461) .. controls (257.3702,84.9669) and (257.9545,85.5512) ..(258.6754,85.5512) .. controls (259.3961,85.5512) and (259.9805,84.9669) ..(259.9805,84.2461);
\path[fill=chfourfig5c0,fill opacity=0.5,draw=chfourfig5c0,line width=0.3985pt,draw opacity=0.5] (230.8394,102.6174) .. controls (230.8394,101.8966) and (230.2550,101.3123) ..(229.5342,101.3123) .. controls (228.8134,101.3123) and (228.2291,101.8966) ..(228.2291,102.6174) .. controls (228.2291,103.3382) and (228.8134,103.9225) ..(229.5342,103.9225) .. controls (230.2550,103.9225) and (230.8394,103.3382) ..(230.8394,102.6174);
\path[fill=chfourfig5c0,fill opacity=0.5,draw=chfourfig5c0,line width=0.3985pt,draw opacity=0.5] (330.9752,53.5010) .. controls (330.9752,52.7802) and (330.3908,52.1959) ..(329.6700,52.1959) .. controls (328.9492,52.1959) and (328.3649,52.7802) ..(328.3649,53.5010) .. controls (328.3649,54.2218) and (328.9492,54.8061) ..(329.6700,54.8061) .. controls (330.3908,54.8061) and (330.9752,54.2218) ..(330.9752,53.5010);
\path[fill=chfourfig5c0,fill opacity=0.5,draw=chfourfig5c0,line width=0.3985pt,draw opacity=0.5] (299.9909,41.9342) .. controls (299.9909,41.2134) and (299.4066,40.6291) ..(298.6858,40.6291) .. controls (297.9650,40.6291) and (297.3807,41.2134) ..(297.3807,41.9342) .. controls (297.3807,42.6550) and (297.9650,43.2393) ..(298.6858,43.2393) .. controls (299.4066,43.2393) and (299.9909,42.6550) ..(299.9909,41.9342);
\path[fill=chfourfig5c0,fill opacity=0.5,draw=chfourfig5c0,line width=0.3985pt,draw opacity=0.5] (303.2388,57.5558) .. controls (303.2388,56.8350) and (302.6545,56.2507) ..(301.9337,56.2507) .. controls (301.2129,56.2507) and (300.6286,56.8350) ..(300.6286,57.5558) .. controls (300.6286,58.2766) and (301.2129,58.8610) ..(301.9337,58.8610) .. controls (302.6545,58.8610) and (303.2388,58.2766) ..(303.2388,57.5558);
\path[fill=chfourfig5c0,fill opacity=0.5,draw=chfourfig5c0,line width=0.3985pt,draw opacity=0.5] (246.1023,103.3746) .. controls (246.1023,102.6538) and (245.5180,102.0695) ..(244.7972,102.0695) .. controls (244.0764,102.0695) and (243.4921,102.6538) ..(243.4921,103.3746) .. controls (243.4921,104.0954) and (244.0764,104.6797) ..(244.7972,104.6797) .. controls (245.5180,104.6797) and (246.1023,104.0954) ..(246.1023,103.3746);
\path[fill=chfourfig5c0,fill opacity=0.5,draw=chfourfig5c0,line width=0.3985pt,draw opacity=0.5] (293.3159,80.6595) .. controls (293.3159,79.9387) and (292.7316,79.3544) ..(292.0107,79.3544) .. controls (291.2900,79.3544) and (290.7057,79.9387) ..(290.7057,80.6595) .. controls (290.7057,81.3803) and (291.2900,81.9646) ..(292.0107,81.9646) .. controls (292.7316,81.9646) and (293.3159,81.3803) ..(293.3159,80.6595);
\path[fill=chfourfig5c0,fill opacity=0.5,draw=chfourfig5c0,line width=0.3985pt,draw opacity=0.5] (243.5020,94.1391) .. controls (243.5020,93.4183) and (242.9177,92.8340) ..(242.1969,92.8340) .. controls (241.4761,92.8340) and (240.8918,93.4183) ..(240.8918,94.1391) .. controls (240.8918,94.8599) and (241.4761,95.4442) ..(242.1969,95.4442) .. controls (242.9177,95.4442) and (243.5020,94.8599) ..(243.5020,94.1391);
\path[fill=chfourfig5c0,fill opacity=0.5,draw=chfourfig5c0,line width=0.3985pt,draw opacity=0.5] (264.3740,74.8811) .. controls (264.3740,74.1603) and (263.7897,73.5760) ..(263.0689,73.5760) .. controls (262.3481,73.5760) and (261.7638,74.1603) ..(261.7638,74.8811) .. controls (261.7638,75.6019) and (262.3481,76.1862) ..(263.0689,76.1862) .. controls (263.7897,76.1862) and (264.3740,75.6019) ..(264.3740,74.8811);
\path[fill=chfourfig5c0,fill opacity=0.5,draw=chfourfig5c0,line width=0.3985pt,draw opacity=0.5] (279.6868,60.8037) .. controls (279.6868,60.0829) and (279.1025,59.4986) ..(278.3817,59.4986) .. controls (277.6609,59.4986) and (277.0766,60.0829) ..(277.0766,60.8037) .. controls (277.0766,61.5245) and (277.6609,62.1088) ..(278.3817,62.1088) .. controls (279.1025,62.1088) and (279.6868,61.5245) ..(279.6868,60.8037);
\path[fill=chfourfig5c0,fill opacity=0.5,draw=chfourfig5c0,line width=0.3985pt,draw opacity=0.5] (266.0777,95.1254) .. controls (266.0777,94.4046) and (265.4934,93.8203) ..(264.7726,93.8203) .. controls (264.0518,93.8203) and (263.4674,94.4046) ..(263.4674,95.1254) .. controls (263.4674,95.8462) and (264.0518,96.4305) ..(264.7726,96.4305) .. controls (265.4934,96.4305) and (266.0777,95.8462) ..(266.0777,95.1254);
\path[fill=chfourfig5c0,fill opacity=0.5,draw=chfourfig5c0,line width=0.3985pt,draw opacity=0.5] (335.6377,21.5703) .. controls (335.6377,20.8495) and (335.0534,20.2652) ..(334.3326,20.2652) .. controls (333.6118,20.2652) and (333.0275,20.8495) ..(333.0275,21.5703) .. controls (333.0275,22.2911) and (333.6118,22.8754) ..(334.3326,22.8754) .. controls (335.0534,22.8754) and (335.6377,22.2911) ..(335.6377,21.5703);
\path[fill=chfourfig5c0,fill opacity=0.5,draw=chfourfig5c0,line width=0.3985pt,draw opacity=0.5] (302.3920,76.7142) .. controls (302.3920,75.9934) and (301.8076,75.4091) ..(301.0869,75.4091) .. controls (300.3661,75.4091) and (299.7817,75.9934) ..(299.7817,76.7142) .. controls (299.7817,77.4350) and (300.3661,78.0193) ..(301.0869,78.0193) .. controls (301.8076,78.0193) and (302.3920,77.4350) ..(302.3920,76.7142);
\path[fill=chfourfig5c0,fill opacity=0.5,draw=chfourfig5c0,line width=0.3985pt,draw opacity=0.5] (232.2939,110.3286) .. controls (232.2939,109.6078) and (231.7096,109.0235) ..(230.9888,109.0235) .. controls (230.2680,109.0235) and (229.6837,109.6078) ..(229.6837,110.3286) .. controls (229.6837,111.0494) and (230.2680,111.6337) ..(230.9888,111.6337) .. controls (231.7096,111.6337) and (232.2939,111.0494) ..(232.2939,110.3286);
\path[fill=chfourfig5c0,fill opacity=0.5,draw=chfourfig5c0,line width=0.3985pt,draw opacity=0.5] (308.1505,43.8670) .. controls (308.1505,43.1462) and (307.5661,42.5619) ..(306.8453,42.5619) .. controls (306.1245,42.5619) and (305.5402,43.1462) ..(305.5402,43.8670) .. controls (305.5402,44.5878) and (306.1245,45.1721) ..(306.8453,45.1721) .. controls (307.5661,45.1721) and (308.1505,44.5878) ..(308.1505,43.8670);
\path[fill=chfourfig5c0,fill opacity=0.5,draw=chfourfig5c0,line width=0.3985pt,draw opacity=0.5] (278.3319,90.7119) .. controls (278.3319,89.9911) and (277.7476,89.4068) ..(277.0267,89.4068) .. controls (276.3060,89.4068) and (275.7216,89.9911) ..(275.7216,90.7119) .. controls (275.7216,91.4327) and (276.3060,92.0170) ..(277.0267,92.0170) .. controls (277.7476,92.0170) and (278.3319,91.4327) ..(278.3319,90.7119);
\path[fill=chfourfig5c0,fill opacity=0.5,draw=chfourfig5c0,line width=0.3985pt,draw opacity=0.5] (279.8761,66.9806) .. controls (279.8761,66.2598) and (279.2918,65.6755) ..(278.5710,65.6755) .. controls (277.8502,65.6755) and (277.2659,66.2598) ..(277.2659,66.9806) .. controls (277.2659,67.7014) and (277.8502,68.2857) ..(278.5710,68.2857) .. controls (279.2918,68.2857) and (279.8761,67.7014) ..(279.8761,66.9806);
\path[fill=chfourfig5c0,fill opacity=0.5,draw=chfourfig5c0,line width=0.3985pt,draw opacity=0.5] (240.9117,93.3122) .. controls (240.9117,92.5914) and (240.3274,92.0071) ..(239.6066,92.0071) .. controls (238.8858,92.0071) and (238.3015,92.5914) ..(238.3015,93.3122) .. controls (238.3015,94.0330) and (238.8858,94.6173) ..(239.6066,94.6173) .. controls (240.3274,94.6173) and (240.9117,94.0330) ..(240.9117,93.3122);
\path[fill=chfourfig5c0,fill opacity=0.5,draw=chfourfig5c0,line width=0.3985pt,draw opacity=0.5] (270.2321,93.7107) .. controls (270.2321,92.9899) and (269.6478,92.4056) ..(268.9270,92.4056) .. controls (268.2062,92.4056) and (267.6219,92.9899) ..(267.6219,93.7107) .. controls (267.6219,94.4315) and (268.2062,95.0158) ..(268.9270,95.0158) .. controls (269.6478,95.0158) and (270.2321,94.4315) ..(270.2321,93.7107);
\path[fill=chfourfig5c0,fill opacity=0.5,draw=chfourfig5c0,line width=0.3985pt,draw opacity=0.5] (272.7228,81.2572) .. controls (272.7228,80.5365) and (272.1385,79.9521) ..(271.4177,79.9521) .. controls (270.6969,79.9521) and (270.1126,80.5365) ..(270.1126,81.2572) .. controls (270.1126,81.9781) and (270.6969,82.5624) ..(271.4177,82.5624) .. controls (272.1385,82.5624) and (272.7228,81.9781) ..(272.7228,81.2572);
\path[fill=chfourfig5c0,fill opacity=0.5,draw=chfourfig5c0,line width=0.3985pt,draw opacity=0.5] (252.9567,88.3308) .. controls (252.9567,87.6100) and (252.3724,87.0257) ..(251.6516,87.0257) .. controls (250.9308,87.0257) and (250.3465,87.6100) ..(250.3465,88.3308) .. controls (250.3465,89.0516) and (250.9308,89.6359) ..(251.6516,89.6359) .. controls (252.3724,89.6359) and (252.9567,89.0516) ..(252.9567,88.3308);
\path[fill=chfourfig5c0,fill opacity=0.5,draw=chfourfig5c0,line width=0.3985pt,draw opacity=0.5] (245.2854,105.9350) .. controls (245.2854,105.2142) and (244.7010,104.6299) ..(243.9803,104.6299) .. controls (243.2595,104.6299) and (242.6751,105.2142) ..(242.6751,105.9350) .. controls (242.6751,106.6558) and (243.2595,107.2402) ..(243.9803,107.2402) .. controls (244.7010,107.2402) and (245.2854,106.6558) ..(245.2854,105.9350);
\path[fill=chfourfig5c0,fill opacity=0.5,draw=chfourfig5c0,line width=0.3985pt,draw opacity=0.5] (283.3730,85.7803) .. controls (283.3730,85.0595) and (282.7887,84.4752) ..(282.0679,84.4752) .. controls (281.3471,84.4752) and (280.7628,85.0595) ..(280.7628,85.7803) .. controls (280.7628,86.5011) and (281.3471,87.0855) ..(282.0679,87.0855) .. controls (282.7887,87.0855) and (283.3730,86.5011) ..(283.3730,85.7803);
\path[fill=chfourfig5c0,fill opacity=0.5,draw=chfourfig5c0,line width=0.3985pt,draw opacity=0.5] (269.3355,82.5026) .. controls (269.3355,81.7818) and (268.7512,81.1975) ..(268.0304,81.1975) .. controls (267.3096,81.1975) and (266.7253,81.7818) ..(266.7253,82.5026) .. controls (266.7253,83.2234) and (267.3096,83.8077) ..(268.0304,83.8077) .. controls (268.7512,83.8077) and (269.3355,83.2234) ..(269.3355,82.5026);
\path[fill=chfourfig5c0,fill opacity=0.5,draw=chfourfig5c0,line width=0.3985pt,draw opacity=0.5] (333.2367,34.7311) .. controls (333.2367,34.0103) and (332.6524,33.4260) ..(331.9316,33.4260) .. controls (331.2108,33.4260) and (330.6264,34.0103) ..(330.6264,34.7311) .. controls (330.6264,35.4519) and (331.2108,36.0363) ..(331.9316,36.0363) .. controls (332.6524,36.0363) and (333.2367,35.4519) ..(333.2367,34.7311);
\path[fill=chfourfig5c0,fill opacity=0.5,draw=chfourfig5c0,line width=0.3985pt,draw opacity=0.5] (335.3189,57.3466) .. controls (335.3189,56.6258) and (334.7346,56.0415) ..(334.0138,56.0415) .. controls (333.2930,56.0415) and (332.7087,56.6258) ..(332.7087,57.3466) .. controls (332.7087,58.0674) and (333.2930,58.6517) ..(334.0138,58.6517) .. controls (334.7346,58.6517) and (335.3189,58.0674) ..(335.3189,57.3466);
\path[fill=chfourfig5c0,fill opacity=0.5,draw=chfourfig5c0,line width=0.3985pt,draw opacity=0.5] (279.5075,67.7378) .. controls (279.5075,67.0170) and (278.9232,66.4326) ..(278.2024,66.4326) .. controls (277.4816,66.4326) and (276.8972,67.0170) ..(276.8972,67.7378) .. controls (276.8972,68.4586) and (277.4816,69.0429) ..(278.2024,69.0429) .. controls (278.9232,69.0429) and (279.5075,68.4586) ..(279.5075,67.7378);
\path[fill=chfourfig5c0,fill opacity=0.5,draw=chfourfig5c0,line width=0.3985pt,draw opacity=0.5] (274.4265,57.8946) .. controls (274.4265,57.1738) and (273.8421,56.5894) ..(273.1214,56.5894) .. controls (272.4005,56.5894) and (271.8162,57.1738) ..(271.8162,57.8946) .. controls (271.8162,58.6154) and (272.4005,59.1997) ..(273.1214,59.1997) .. controls (273.8421,59.1997) and (274.4265,58.6154) ..(274.4265,57.8946);
\path[fill=chfourfig5c0,fill opacity=0.5,draw=chfourfig5c0,line width=0.3985pt,draw opacity=0.5] (239.9055,93.0532) .. controls (239.9055,92.3324) and (239.3212,91.7480) ..(238.6003,91.7480) .. controls (237.8796,91.7480) and (237.2952,92.3324) ..(237.2952,93.0532) .. controls (237.2952,93.7740) and (237.8796,94.3583) ..(238.6003,94.3583) .. controls (239.3212,94.3583) and (239.9055,93.7740) ..(239.9055,93.0532);
\path[fill=chfourfig5c0,fill opacity=0.5,draw=chfourfig5c0,line width=0.3985pt,draw opacity=0.5] (332.2006,35.9167) .. controls (332.2006,35.1959) and (331.6163,34.6116) ..(330.8954,34.6116) .. controls (330.1747,34.6116) and (329.5903,35.1959) ..(329.5903,35.9167) .. controls (329.5903,36.6375) and (330.1747,37.2218) ..(330.8954,37.2218) .. controls (331.6163,37.2218) and (332.2006,36.6375) ..(332.2006,35.9167);
\path[fill=chfourfig5c0,fill opacity=0.5,draw=chfourfig5c0,line width=0.3985pt,draw opacity=0.5] (237.1558,104.0919) .. controls (237.1558,103.3711) and (236.5714,102.7868) ..(235.8506,102.7868) .. controls (235.1298,102.7868) and (234.5455,103.3711) ..(234.5455,104.0919) .. controls (234.5455,104.8127) and (235.1298,105.3970) ..(235.8506,105.3970) .. controls (236.5714,105.3970) and (237.1558,104.8127) ..(237.1558,104.0919);
\path[fill=chfourfig5c0,fill opacity=0.5,draw=chfourfig5c0,line width=0.3985pt,draw opacity=0.5] (296.9921,42.7910) .. controls (296.9921,42.0702) and (296.4078,41.4859) ..(295.6870,41.4859) .. controls (294.9662,41.4859) and (294.3819,42.0702) ..(294.3819,42.7910) .. controls (294.3819,43.5118) and (294.9662,44.0961) ..(295.6870,44.0961) .. controls (296.4078,44.0961) and (296.9921,43.5118) ..(296.9921,42.7910);
\path[fill=chfourfig5c0,fill opacity=0.5,draw=chfourfig5c0,line width=0.3985pt,draw opacity=0.5] (271.5672,74.6021) .. controls (271.5672,73.8813) and (270.9828,73.2970) ..(270.2620,73.2970) .. controls (269.5413,73.2970) and (268.9569,73.8813) ..(268.9569,74.6021) .. controls (268.9569,75.3229) and (269.5413,75.9073) ..(270.2620,75.9073) .. controls (270.9828,75.9073) and (271.5672,75.3229) ..(271.5672,74.6021);
\path[fill=chfourfig5c0,fill opacity=0.5,draw=chfourfig5c0,line width=0.3985pt,draw opacity=0.5] (308.4693,43.3888) .. controls (308.4693,42.6680) and (307.8849,42.0837) ..(307.1641,42.0837) .. controls (306.4433,42.0837) and (305.8590,42.6680) ..(305.8590,43.3888) .. controls (305.8590,44.1096) and (306.4433,44.6939) ..(307.1641,44.6939) .. controls (307.8849,44.6939) and (308.4693,44.1096) ..(308.4693,43.3888);
\path[fill=chfourfig5c0,fill opacity=0.5,draw=chfourfig5c0,line width=0.3985pt,draw opacity=0.5] (237.5144,97.8154) .. controls (237.5144,97.0946) and (236.9301,96.5103) ..(236.2093,96.5103) .. controls (235.4885,96.5103) and (234.9042,97.0946) ..(234.9042,97.8154) .. controls (234.9042,98.5362) and (235.4885,99.1205) ..(236.2093,99.1205) .. controls (236.9301,99.1205) and (237.5144,98.5362) ..(237.5144,97.8154);
\path[fill=chfourfig5c0,fill opacity=0.5,draw=chfourfig5c0,line width=0.3985pt,draw opacity=0.5] (246.3414,91.1403) .. controls (246.3414,90.4195) and (245.7571,89.8352) ..(245.0363,89.8352) .. controls (244.3155,89.8352) and (243.7312,90.4195) ..(243.7312,91.1403) .. controls (243.7312,91.8611) and (244.3155,92.4454) ..(245.0363,92.4454) .. controls (245.7571,92.4454) and (246.3414,91.8611) ..(246.3414,91.1403);
\path[fill=chfourfig5c0,fill opacity=0.5,draw=chfourfig5c0,line width=0.3985pt,draw opacity=0.5] (303.9960,43.0002) .. controls (303.9960,42.2794) and (303.4117,41.6951) ..(302.6909,41.6951) .. controls (301.9700,41.6951) and (301.3857,42.2794) ..(301.3857,43.0002) .. controls (301.3857,43.7210) and (301.9700,44.3054) ..(302.6909,44.3054) .. controls (303.4117,44.3054) and (303.9960,43.7210) ..(303.9960,43.0002);
\path[fill=chfourfig5c0,fill opacity=0.5,draw=chfourfig5c0,line width=0.3985pt,draw opacity=0.5] (260.7476,101.2127) .. controls (260.7476,100.4919) and (260.1633,99.9076) ..(259.4425,99.9076) .. controls (258.7217,99.9076) and (258.1373,100.4919) ..(258.1373,101.2127) .. controls (258.1373,101.9335) and (258.7217,102.5178) ..(259.4425,102.5178) .. controls (260.1633,102.5178) and (260.7476,101.9335) ..(260.7476,101.2127);
\path[fill=chfourfig5c0,fill opacity=0.5,draw=chfourfig5c0,line width=0.3985pt,draw opacity=0.5] (335.7373,20.5940) .. controls (335.7373,19.8732) and (335.1530,19.2888) ..(334.4323,19.2888) .. controls (333.7114,19.2888) and (333.1271,19.8732) ..(333.1271,20.5940) .. controls (333.1271,21.3148) and (333.7114,21.8991) ..(334.4323,21.8991) .. controls (335.1530,21.8991) and (335.7373,21.3148) ..(335.7373,20.5940);
\path[fill=chfourfig5c0,fill opacity=0.5,draw=chfourfig5c0,line width=0.3985pt,draw opacity=0.5] (320.3947,52.2756) .. controls (320.3947,51.5548) and (319.8104,50.9705) ..(319.0896,50.9705) .. controls (318.3687,50.9705) and (317.7844,51.5548) ..(317.7844,52.2756) .. controls (317.7844,52.9964) and (318.3687,53.5807) ..(319.0896,53.5807) .. controls (319.8104,53.5807) and (320.3947,52.9964) ..(320.3947,52.2756);
\path[fill=chfourfig5c0,fill opacity=0.5,draw=chfourfig5c0,line width=0.3985pt,draw opacity=0.5] (275.8412,85.6010) .. controls (275.8412,84.8802) and (275.2569,84.2959) ..(274.5361,84.2959) .. controls (273.8153,84.2959) and (273.2309,84.8802) ..(273.2309,85.6010) .. controls (273.2309,86.3218) and (273.8153,86.9061) ..(274.5361,86.9061) .. controls (275.2569,86.9061) and (275.8412,86.3218) ..(275.8412,85.6010);
\path[fill=chfourfig5c0,fill opacity=0.5,draw=chfourfig5c0,line width=0.3985pt,draw opacity=0.5] (264.4737,68.5846) .. controls (264.4737,67.8638) and (263.8893,67.2795) ..(263.1685,67.2795) .. controls (262.4478,67.2795) and (261.8634,67.8638) ..(261.8634,68.5846) .. controls (261.8634,69.3054) and (262.4478,69.8897) ..(263.1685,69.8897) .. controls (263.8893,69.8897) and (264.4737,69.3054) ..(264.4737,68.5846);
\path[fill=chfourfig5c0,fill opacity=0.5,draw=chfourfig5c0,line width=0.3985pt,draw opacity=0.5] (284.4490,79.1750) .. controls (284.4490,78.4542) and (283.8647,77.8699) ..(283.1439,77.8699) .. controls (282.4231,77.8699) and (281.8388,78.4542) ..(281.8388,79.1750) .. controls (281.8388,79.8958) and (282.4231,80.4802) ..(283.1439,80.4802) .. controls (283.8647,80.4802) and (284.4490,79.8958) ..(284.4490,79.1750);
\path[fill=chfourfig5c0,fill opacity=0.5,draw=chfourfig5c0,line width=0.3985pt,draw opacity=0.5] (247.9952,72.1015) .. controls (247.9952,71.3807) and (247.4109,70.7963) ..(246.6901,70.7963) .. controls (245.9693,70.7963) and (245.3850,71.3807) ..(245.3850,72.1015) .. controls (245.3850,72.8223) and (245.9693,73.4066) ..(246.6901,73.4066) .. controls (247.4109,73.4066) and (247.9952,72.8223) ..(247.9952,72.1015);
\path[fill=chfourfig5c0,fill opacity=0.5,draw=chfourfig5c0,line width=0.3985pt,draw opacity=0.5] (255.3976,65.2072) .. controls (255.3976,64.4864) and (254.8133,63.9021) ..(254.0925,63.9021) .. controls (253.3717,63.9021) and (252.7874,64.4864) ..(252.7874,65.2072) .. controls (252.7874,65.9280) and (253.3717,66.5124) ..(254.0925,66.5124) .. controls (254.8133,66.5124) and (255.3976,65.9280) ..(255.3976,65.2072);
\path[fill=chfourfig5c0,fill opacity=0.5,draw=chfourfig5c0,line width=0.3985pt,draw opacity=0.5] (252.6877,71.0454) .. controls (252.6877,70.3246) and (252.1034,69.7403) ..(251.3826,69.7403) .. controls (250.6618,69.7403) and (250.0775,70.3246) ..(250.0775,71.0454) .. controls (250.0775,71.7662) and (250.6618,72.3505) ..(251.3826,72.3505) .. controls (252.1034,72.3505) and (252.6877,71.7662) ..(252.6877,71.0454);
\path[fill=chfourfig5c0,fill opacity=0.5,draw=chfourfig5c0,line width=0.3985pt,draw opacity=0.5] (272.6033,90.0444) .. controls (272.6033,89.3236) and (272.0190,88.7393) ..(271.2982,88.7393) .. controls (270.5774,88.7393) and (269.9930,89.3236) ..(269.9930,90.0444) .. controls (269.9930,90.7652) and (270.5774,91.3495) ..(271.2982,91.3495) .. controls (272.0190,91.3495) and (272.6033,90.7652) ..(272.6033,90.0444);
\path[fill=chfourfig5c0,fill opacity=0.5,draw=chfourfig5c0,line width=0.3985pt,draw opacity=0.5] (246.1023,84.4354) .. controls (246.1023,83.7146) and (245.5180,83.1303) ..(244.7972,83.1303) .. controls (244.0764,83.1303) and (243.4921,83.7146) ..(243.4921,84.4354) .. controls (243.4921,85.1562) and (244.0764,85.7405) ..(244.7972,85.7405) .. controls (245.5180,85.7405) and (246.1023,85.1562) ..(246.1023,84.4354);
\path[fill=chfourfig5c0,fill opacity=0.5,draw=chfourfig5c0,line width=0.3985pt,draw opacity=0.5] (276.8474,92.9137) .. controls (276.8474,92.1929) and (276.2631,91.6086) ..(275.5423,91.6086) .. controls (274.8215,91.6086) and (274.2372,92.1929) ..(274.2372,92.9137) .. controls (274.2372,93.6345) and (274.8215,94.2188) ..(275.5423,94.2188) .. controls (276.2631,94.2188) and (276.8474,93.6345) ..(276.8474,92.9137);
\path[fill=chfourfig5c0,fill opacity=0.5,draw=chfourfig5c0,line width=0.3985pt,draw opacity=0.5] (295.8863,66.7116) .. controls (295.8863,65.9908) and (295.3020,65.4065) ..(294.5811,65.4065) .. controls (293.8604,65.4065) and (293.2761,65.9908) ..(293.2761,66.7116) .. controls (293.2761,67.4324) and (293.8604,68.0167) ..(294.5811,68.0167) .. controls (295.3020,68.0167) and (295.8863,67.4324) ..(295.8863,66.7116);
\path[fill=chfourfig5c0,fill opacity=0.5,draw=chfourfig5c0,line width=0.3985pt,draw opacity=0.5] (237.9627,101.0832) .. controls (237.9627,100.3624) and (237.3784,99.7780) ..(236.6576,99.7780) .. controls (235.9368,99.7780) and (235.3525,100.3624) ..(235.3525,101.0832) .. controls (235.3525,101.8040) and (235.9368,102.3883) ..(236.6576,102.3883) .. controls (237.3784,102.3883) and (237.9627,101.8040) ..(237.9627,101.0832);
\path[fill=chfourfig5c0,fill opacity=0.5,draw=chfourfig5c0,line width=0.3985pt,draw opacity=0.5] (329.1121,36.5842) .. controls (329.1121,35.8634) and (328.5278,35.2791) ..(327.8070,35.2791) .. controls (327.0862,35.2791) and (326.5019,35.8634) ..(326.5019,36.5842) .. controls (326.5019,37.3050) and (327.0862,37.8893) ..(327.8070,37.8893) .. controls (328.5278,37.8893) and (329.1121,37.3050) ..(329.1121,36.5842);
\path[fill=chfourfig5c0,fill opacity=0.5,draw=chfourfig5c0,line width=0.3985pt,draw opacity=0.5] (291.9311,81.6956) .. controls (291.9311,80.9748) and (291.3467,80.3905) ..(290.6259,80.3905) .. controls (289.9052,80.3905) and (289.3208,80.9748) ..(289.3208,81.6956) .. controls (289.3208,82.4164) and (289.9052,83.0007) ..(290.6259,83.0007) .. controls (291.3467,83.0007) and (291.9311,82.4164) ..(291.9311,81.6956);
\path[fill=chfourfig5c0,fill opacity=0.5,draw=chfourfig5c0,line width=0.3985pt,draw opacity=0.5] (253.6940,73.5660) .. controls (253.6940,72.8452) and (253.1096,72.2609) ..(252.3888,72.2609) .. controls (251.6680,72.2609) and (251.0837,72.8452) ..(251.0837,73.5660) .. controls (251.0837,74.2868) and (251.6680,74.8711) ..(252.3888,74.8711) .. controls (253.1096,74.8711) and (253.6940,74.2868) ..(253.6940,73.5660);
\path[fill=chfourfig5c0,fill opacity=0.5,draw=chfourfig5c0,line width=0.3985pt,draw opacity=0.5] (277.6345,65.1774) .. controls (277.6345,64.4565) and (277.0502,63.8722) ..(276.3293,63.8722) .. controls (275.6086,63.8722) and (275.0243,64.4565) ..(275.0243,65.1774) .. controls (275.0243,65.8981) and (275.6086,66.4825) ..(276.3293,66.4825) .. controls (277.0502,66.4825) and (277.6345,65.8981) ..(277.6345,65.1774);
\path[fill=chfourfig5c0,fill opacity=0.5,draw=chfourfig5c0,line width=0.3985pt,draw opacity=0.5] (264.5533,74.0741) .. controls (264.5533,73.3533) and (263.9690,72.7690) ..(263.2483,72.7690) .. controls (262.5274,72.7690) and (261.9431,73.3533) ..(261.9431,74.0741) .. controls (261.9431,74.7949) and (262.5274,75.3792) ..(263.2483,75.3792) .. controls (263.9690,75.3792) and (264.5533,74.7949) ..(264.5533,74.0741);
\path[fill=chfourfig5c0,fill opacity=0.5,draw=chfourfig5c0,line width=0.3985pt,draw opacity=0.5] (252.8670,105.8155) .. controls (252.8670,105.0947) and (252.2827,104.5104) ..(251.5619,104.5104) .. controls (250.8411,104.5104) and (250.2568,105.0947) ..(250.2568,105.8155) .. controls (250.2568,106.5363) and (250.8411,107.1206) ..(251.5619,107.1206) .. controls (252.2827,107.1206) and (252.8670,106.5363) ..(252.8670,105.8155);
\path[fill=chfourfig5c0,fill opacity=0.5,draw=chfourfig5c0,line width=0.3985pt,draw opacity=0.5] (283.4129,53.4512) .. controls (283.4129,52.7304) and (282.8286,52.1460) ..(282.1078,52.1460) .. controls (281.3870,52.1460) and (280.8026,52.7304) ..(280.8026,53.4512) .. controls (280.8026,54.1720) and (281.3870,54.7563) ..(282.1078,54.7563) .. controls (282.8286,54.7563) and (283.4129,54.1720) ..(283.4129,53.4512);
\path[fill=chfourfig5c0,fill opacity=0.5,draw=chfourfig5c0,line width=0.3985pt,draw opacity=0.5] (295.5276,54.9655) .. controls (295.5276,54.2447) and (294.9433,53.6604) ..(294.2225,53.6604) .. controls (293.5017,53.6604) and (292.9174,54.2447) ..(292.9174,54.9655) .. controls (292.9174,55.6863) and (293.5017,56.2706) ..(294.2225,56.2706) .. controls (294.9433,56.2706) and (295.5276,55.6863) ..(295.5276,54.9655);
\path[fill=chfourfig5c0,fill opacity=0.5,draw=chfourfig5c0,line width=0.3985pt,draw opacity=0.5] (236.6277,115.2701) .. controls (236.6277,114.5493) and (236.0434,113.9650) ..(235.3226,113.9650) .. controls (234.6018,113.9650) and (234.0175,114.5493) ..(234.0175,115.2701) .. controls (234.0175,115.9910) and (234.6018,116.5753) ..(235.3226,116.5753) .. controls (236.0434,116.5753) and (236.6277,115.9910) ..(236.6277,115.2701);
\path[fill=chfourfig5c0,fill opacity=0.5,draw=chfourfig5c0,line width=0.3985pt,draw opacity=0.5] (284.5985,54.5570) .. controls (284.5985,53.8362) and (284.0142,53.2519) ..(283.2933,53.2519) .. controls (282.5725,53.2519) and (281.9882,53.8362) ..(281.9882,54.5570) .. controls (281.9882,55.2778) and (282.5725,55.8622) ..(283.2933,55.8622) .. controls (284.0142,55.8622) and (284.5985,55.2778) ..(284.5985,54.5570);
\path[fill=chfourfig5c0,fill opacity=0.5,draw=chfourfig5c0,line width=0.3985pt,draw opacity=0.5] (252.7176,79.2049) .. controls (252.7176,78.4841) and (252.1333,77.8998) ..(251.4125,77.8998) .. controls (250.6917,77.8998) and (250.1074,78.4841) ..(250.1074,79.2049) .. controls (250.1074,79.9257) and (250.6917,80.5100) ..(251.4125,80.5100) .. controls (252.1333,80.5100) and (252.7176,79.9257) ..(252.7176,79.2049);
\path[fill=chfourfig5c0,fill opacity=0.5,draw=chfourfig5c0,line width=0.3985pt,draw opacity=0.5] (271.6369,83.0506) .. controls (271.6369,82.3298) and (271.0526,81.7454) ..(270.3318,81.7454) .. controls (269.6110,81.7454) and (269.0266,82.3298) ..(269.0266,83.0506) .. controls (269.0266,83.7714) and (269.6110,84.3557) ..(270.3318,84.3557) .. controls (271.0526,84.3557) and (271.6369,83.7714) ..(271.6369,83.0506);
\path[fill=chfourfig5c0,fill opacity=0.5,draw=chfourfig5c0,line width=0.3985pt,draw opacity=0.5] (323.1743,54.8659) .. controls (323.1743,54.1451) and (322.5900,53.5608) ..(321.8692,53.5608) .. controls (321.1484,53.5608) and (320.5641,54.1451) ..(320.5641,54.8659) .. controls (320.5641,55.5867) and (321.1484,56.1710) ..(321.8692,56.1710) .. controls (322.5900,56.1710) and (323.1743,55.5867) ..(323.1743,54.8659);
\path[fill=chfourfig5c0,fill opacity=0.5,draw=chfourfig5c0,line width=0.3985pt,draw opacity=0.5] (330.7759,44.1459) .. controls (330.7759,43.4251) and (330.1916,42.8408) ..(329.4708,42.8408) .. controls (328.7500,42.8408) and (328.1656,43.4251) ..(328.1656,44.1459) .. controls (328.1656,44.8667) and (328.7500,45.4510) ..(329.4708,45.4510) .. controls (330.1916,45.4510) and (330.7759,44.8667) ..(330.7759,44.1459);
\path[fill=chfourfig5c0,fill opacity=0.5,draw=chfourfig5c0,line width=0.3985pt,draw opacity=0.5] (325.5853,64.1313) .. controls (325.5853,63.4104) and (325.0010,62.8261) ..(324.2802,62.8261) .. controls (323.5594,62.8261) and (322.9751,63.4104) ..(322.9751,64.1313) .. controls (322.9751,64.8521) and (323.5594,65.4364) ..(324.2802,65.4364) .. controls (325.0010,65.4364) and (325.5853,64.8521) ..(325.5853,64.1313);
\path[fill=chfourfig5c0,fill opacity=0.5,draw=chfourfig5c0,line width=0.3985pt,draw opacity=0.5] (302.8801,47.0252) .. controls (302.8801,46.3044) and (302.2958,45.7201) ..(301.5750,45.7201) .. controls (300.8542,45.7201) and (300.2699,46.3044) ..(300.2699,47.0252) .. controls (300.2699,47.7460) and (300.8542,48.3303) ..(301.5750,48.3303) .. controls (302.2958,48.3303) and (302.8801,47.7460) ..(302.8801,47.0252);
\path[fill=chfourfig5c0,fill opacity=0.5,draw=chfourfig5c0,line width=0.3985pt,draw opacity=0.5] (286.0530,87.6733) .. controls (286.0530,86.9525) and (285.4687,86.3682) ..(284.7479,86.3682) .. controls (284.0271,86.3682) and (283.4428,86.9525) ..(283.4428,87.6733) .. controls (283.4428,88.3941) and (284.0271,88.9784) ..(284.7479,88.9784) .. controls (285.4687,88.9784) and (286.0530,88.3941) ..(286.0530,87.6733);
\path[fill=chfourfig5c0,fill opacity=0.5,draw=chfourfig5c0,line width=0.3985pt,draw opacity=0.5] (253.3851,94.4280) .. controls (253.3851,93.7072) and (252.8008,93.1229) ..(252.0800,93.1229) .. controls (251.3592,93.1229) and (250.7749,93.7072) ..(250.7749,94.4280) .. controls (250.7749,95.1488) and (251.3592,95.7332) ..(252.0800,95.7332) .. controls (252.8008,95.7332) and (253.3851,95.1488) ..(253.3851,94.4280);
\path[fill=chfourfig5c0,fill opacity=0.5,draw=chfourfig5c0,line width=0.3985pt,draw opacity=0.5] (323.0846,49.5159) .. controls (323.0846,48.7951) and (322.5003,48.2108) ..(321.7795,48.2108) .. controls (321.0587,48.2108) and (320.4744,48.7951) ..(320.4744,49.5159) .. controls (320.4744,50.2367) and (321.0587,50.8210) ..(321.7795,50.8210) .. controls (322.5003,50.8210) and (323.0846,50.2367) ..(323.0846,49.5159);
\path[fill=chfourfig5c0,fill opacity=0.5,draw=chfourfig5c0,line width=0.3985pt,draw opacity=0.5] (319.1493,37.1421) .. controls (319.1493,36.4213) and (318.5650,35.8370) ..(317.8442,35.8370) .. controls (317.1234,35.8370) and (316.5391,36.4213) ..(316.5391,37.1421) .. controls (316.5391,37.8629) and (317.1234,38.4472) ..(317.8442,38.4472) .. controls (318.5650,38.4472) and (319.1493,37.8629) ..(319.1493,37.1421);
\path[fill=chfourfig5c0,fill opacity=0.5,draw=chfourfig5c0,line width=0.3985pt,draw opacity=0.5] (294.9796,49.5757) .. controls (294.9796,48.8549) and (294.3953,48.2705) ..(293.6746,48.2705) .. controls (292.9537,48.2705) and (292.3694,48.8549) ..(292.3694,49.5757) .. controls (292.3694,50.2964) and (292.9537,50.8808) ..(293.6746,50.8808) .. controls (294.3953,50.8808) and (294.9796,50.2964) ..(294.9796,49.5757);
\path[fill=chfourfig5c0,fill opacity=0.5,draw=chfourfig5c0,line width=0.3985pt,draw opacity=0.5] (283.1040,75.5685) .. controls (283.1040,74.8477) and (282.5197,74.2634) ..(281.7989,74.2634) .. controls (281.0781,74.2634) and (280.4938,74.8477) ..(280.4938,75.5685) .. controls (280.4938,76.2893) and (281.0781,76.8736) ..(281.7989,76.8736) .. controls (282.5197,76.8736) and (283.1040,76.2893) ..(283.1040,75.5685);
\path[fill=chfourfig5c0,fill opacity=0.5,draw=chfourfig5c0,line width=0.3985pt,draw opacity=0.5] (293.3159,70.3381) .. controls (293.3159,69.6172) and (292.7316,69.0329) ..(292.0107,69.0329) .. controls (291.2900,69.0329) and (290.7057,69.6172) ..(290.7057,70.3381) .. controls (290.7057,71.0589) and (291.2900,71.6432) ..(292.0107,71.6432) .. controls (292.7316,71.6432) and (293.3159,71.0589) ..(293.3159,70.3381);
\path[fill=chfourfig5c0,fill opacity=0.5,draw=chfourfig5c0,line width=0.3985pt,draw opacity=0.5] (303.5377,70.5373) .. controls (303.5377,69.8165) and (302.9534,69.2322) ..(302.2325,69.2322) .. controls (301.5118,69.2322) and (300.9275,69.8165) ..(300.9275,70.5373) .. controls (300.9275,71.2581) and (301.5118,71.8424) ..(302.2325,71.8424) .. controls (302.9534,71.8424) and (303.5377,71.2581) ..(303.5377,70.5373);
\path[fill=chfourfig5c0,fill opacity=0.5,draw=chfourfig5c0,line width=0.3985pt,draw opacity=0.5] (306.6859,42.5121) .. controls (306.6859,41.7913) and (306.1016,41.2069) ..(305.3808,41.2069) .. controls (304.6600,41.2069) and (304.0757,41.7913) ..(304.0757,42.5121) .. controls (304.0757,43.2328) and (304.6600,43.8172) ..(305.3808,43.8172) .. controls (306.1016,43.8172) and (306.6859,43.2328) ..(306.6859,42.5121);
\path[fill=chfourfig5c0,fill opacity=0.5,draw=chfourfig5c0,line width=0.3985pt,draw opacity=0.5] (320.3549,39.4834) .. controls (320.3549,38.7626) and (319.7705,38.1783) ..(319.0497,38.1783) .. controls (318.3289,38.1783) and (317.7446,38.7626) ..(317.7446,39.4834) .. controls (317.7446,40.2042) and (318.3289,40.7885) ..(319.0497,40.7885) .. controls (319.7705,40.7885) and (320.3549,40.2042) ..(320.3549,39.4834);
\path[fill=chfourfig5c0,fill opacity=0.5,draw=chfourfig5c0,line width=0.3985pt,draw opacity=0.5] (240.1545,86.9061) .. controls (240.1545,86.1853) and (239.5703,85.6010) ..(238.8494,85.6010) .. controls (238.1286,85.6010) and (237.5443,86.1853) ..(237.5443,86.9061) .. controls (237.5443,87.6269) and (238.1286,88.2113) ..(238.8494,88.2113) .. controls (239.5703,88.2113) and (240.1545,87.6269) ..(240.1545,86.9061);
\path[fill=chfourfig5c0,fill opacity=0.5,draw=chfourfig5c0,line width=0.3985pt,draw opacity=0.5] (318.0036,36.0562) .. controls (318.0036,35.3354) and (317.4193,34.7511) ..(316.6985,34.7511) .. controls (315.9777,34.7511) and (315.3934,35.3354) ..(315.3934,36.0562) .. controls (315.3934,36.7770) and (315.9777,37.3613) ..(316.6985,37.3613) .. controls (317.4193,37.3613) and (318.0036,36.7770) ..(318.0036,36.0562);
\path[fill=chfourfig5c0,fill opacity=0.5,draw=chfourfig5c0,line width=0.3985pt,draw opacity=0.5] (335.5979,22.8754) .. controls (335.5979,22.1546) and (335.0136,21.5703) ..(334.2928,21.5703) .. controls (333.5720,21.5703) and (332.9876,22.1546) ..(332.9876,22.8754) .. controls (332.9876,23.5962) and (333.5720,24.1805) ..(334.2928,24.1805) .. controls (335.0136,24.1805) and (335.5979,23.5962) ..(335.5979,22.8754);
\path[fill=chfourfig5c0,fill opacity=0.5,draw=chfourfig5c0,line width=0.3985pt,draw opacity=0.5] (270.6805,79.4440) .. controls (270.6805,78.7232) and (270.0962,78.1389) ..(269.3753,78.1389) .. controls (268.6546,78.1389) and (268.0703,78.7232) ..(268.0703,79.4440) .. controls (268.0703,80.1648) and (268.6546,80.7491) ..(269.3753,80.7491) .. controls (270.0962,80.7491) and (270.6805,80.1648) ..(270.6805,79.4440);
\path[fill=chfourfig5c0,fill opacity=0.5,draw=chfourfig5c0,line width=0.3985pt,draw opacity=0.5] (328.2653,54.8659) .. controls (328.2653,54.1451) and (327.6810,53.5608) ..(326.9601,53.5608) .. controls (326.2393,53.5608) and (325.6551,54.1451) ..(325.6551,54.8659) .. controls (325.6551,55.5867) and (326.2393,56.1710) ..(326.9601,56.1710) .. controls (327.6810,56.1710) and (328.2653,55.5867) ..(328.2653,54.8659);
\path[fill=chfourfig5c0,fill opacity=0.5,draw=chfourfig5c0,line width=0.3985pt,draw opacity=0.5] (318.6711,69.1126) .. controls (318.6711,68.3918) and (318.0868,67.8075) ..(317.3660,67.8075) .. controls (316.6452,67.8075) and (316.0609,68.3918) ..(316.0609,69.1126) .. controls (316.0609,69.8334) and (316.6452,70.4178) ..(317.3660,70.4178) .. controls (318.0868,70.4178) and (318.6711,69.8334) ..(318.6711,69.1126);
\path[fill=chfourfig5c0,fill opacity=0.5,draw=chfourfig5c0,line width=0.3985pt,draw opacity=0.5] (327.2889,38.8358) .. controls (327.2889,38.1150) and (326.7046,37.5307) ..(325.9838,37.5307) .. controls (325.2630,37.5307) and (324.6787,38.1150) ..(324.6787,38.8358) .. controls (324.6787,39.5566) and (325.2630,40.1409) ..(325.9838,40.1409) .. controls (326.7046,40.1409) and (327.2889,39.5566) ..(327.2889,38.8358);
\path[fill=chfourfig5c0,fill opacity=0.5,draw=chfourfig5c0,line width=0.3985pt,draw opacity=0.5] (311.4083,52.7737) .. controls (311.4083,52.0529) and (310.8239,51.4686) ..(310.1031,51.4686) .. controls (309.3824,51.4686) and (308.7980,52.0529) ..(308.7980,52.7737) .. controls (308.7980,53.4945) and (309.3824,54.0788) ..(310.1031,54.0788) .. controls (310.8239,54.0788) and (311.4083,53.4945) ..(311.4083,52.7737);
\path[fill=chfourfig5c0,fill opacity=0.5,draw=chfourfig5c0,line width=0.3985pt,draw opacity=0.5] (283.1837,68.0167) .. controls (283.1837,67.2959) and (282.5994,66.7116) ..(281.8786,66.7116) .. controls (281.1578,66.7116) and (280.5735,67.2959) ..(280.5735,68.0167) .. controls (280.5735,68.7375) and (281.1578,69.3219) ..(281.8786,69.3219) .. controls (282.5994,69.3219) and (283.1837,68.7375) ..(283.1837,68.0167);
\path[fill=chfourfig5c0,fill opacity=0.5,draw=chfourfig5c0,line width=0.3985pt,draw opacity=0.5] (301.6647,70.6071) .. controls (301.6647,69.8862) and (301.0804,69.3019) ..(300.3596,69.3019) .. controls (299.6387,69.3019) and (299.0544,69.8862) ..(299.0544,70.6071) .. controls (299.0544,71.3279) and (299.6387,71.9122) ..(300.3596,71.9122) .. controls (301.0804,71.9122) and (301.6647,71.3279) ..(301.6647,70.6071);
\path[fill=chfourfig5c0,fill opacity=0.5,draw=chfourfig5c0,line width=0.3985pt,draw opacity=0.5] (276.4091,82.0642) .. controls (276.4091,81.3434) and (275.8248,80.7591) ..(275.1039,80.7591) .. controls (274.3831,80.7591) and (273.7988,81.3434) ..(273.7988,82.0642) .. controls (273.7988,82.7850) and (274.3831,83.3694) ..(275.1039,83.3694) .. controls (275.8248,83.3694) and (276.4091,82.7850) ..(276.4091,82.0642);
\path[fill=chfourfig5c0,fill opacity=0.5,draw=chfourfig5c0,line width=0.3985pt,draw opacity=0.5] (311.0994,73.7453) .. controls (311.0994,73.0245) and (310.5151,72.4402) ..(309.7943,72.4402) .. controls (309.0735,72.4402) and (308.4892,73.0245) ..(308.4892,73.7453) .. controls (308.4892,74.4661) and (309.0735,75.0505) ..(309.7943,75.0505) .. controls (310.5151,75.0505) and (311.0994,74.4661) ..(311.0994,73.7453);
\path[fill=chfourfig5c0,fill opacity=0.5,draw=chfourfig5c0,line width=0.3985pt,draw opacity=0.5] (290.5263,62.5571) .. controls (290.5263,61.8363) and (289.9420,61.2520) ..(289.2212,61.2520) .. controls (288.5004,61.2520) and (287.9160,61.8363) ..(287.9160,62.5571) .. controls (287.9160,63.2779) and (288.5004,63.8623) ..(289.2212,63.8623) .. controls (289.9420,63.8623) and (290.5263,63.2779) ..(290.5263,62.5571);
\path[fill=chfourfig5c0,fill opacity=0.5,draw=chfourfig5c0,line width=0.3985pt,draw opacity=0.5] (313.0222,71.9022) .. controls (313.0222,71.1814) and (312.4379,70.5971) ..(311.7171,70.5971) .. controls (310.9963,70.5971) and (310.4120,71.1814) ..(310.4120,71.9022) .. controls (310.4120,72.6230) and (310.9963,73.2073) ..(311.7171,73.2073) .. controls (312.4379,73.2073) and (313.0222,72.6230) ..(313.0222,71.9022);
\path[fill=chfourfig5c0,fill opacity=0.5,draw=chfourfig5c0,line width=0.3985pt,draw opacity=0.5] (293.3657,47.7425) .. controls (293.3657,47.0217) and (292.7814,46.4374) ..(292.0606,46.4374) .. controls (291.3398,46.4374) and (290.7555,47.0217) ..(290.7555,47.7425) .. controls (290.7555,48.4633) and (291.3398,49.0476) ..(292.0606,49.0476) .. controls (292.7814,49.0476) and (293.3657,48.4633) ..(293.3657,47.7425);
\path[fill=chfourfig5c0,fill opacity=0.5,draw=chfourfig5c0,line width=0.3985pt,draw opacity=0.5] (275.8113,62.0092) .. controls (275.8113,61.2884) and (275.2270,60.7041) ..(274.5062,60.7041) .. controls (273.7854,60.7041) and (273.2010,61.2884) ..(273.2010,62.0092) .. controls (273.2010,62.7300) and (273.7854,63.3143) ..(274.5062,63.3143) .. controls (275.2270,63.3143) and (275.8113,62.7300) ..(275.8113,62.0092);
\path[fill=chfourfig5c0,fill opacity=0.5,draw=chfourfig5c0,line width=0.3985pt,draw opacity=0.5] (261.3952,89.6858) .. controls (261.3952,88.9650) and (260.8109,88.3806) ..(260.0901,88.3806) .. controls (259.3692,88.3806) and (258.7849,88.9650) ..(258.7849,89.6858) .. controls (258.7849,90.4066) and (259.3692,90.9909) ..(260.0901,90.9909) .. controls (260.8109,90.9909) and (261.3952,90.4066) ..(261.3952,89.6858);
\path[fill=chfourfig5c0,fill opacity=0.5,draw=chfourfig5c0,line width=0.3985pt,draw opacity=0.5] (322.3673,32.8382) .. controls (322.3673,32.1174) and (321.7830,31.5331) ..(321.0622,31.5331) .. controls (320.3414,31.5331) and (319.7571,32.1174) ..(319.7571,32.8382) .. controls (319.7571,33.5590) and (320.3414,34.1433) ..(321.0622,34.1433) .. controls (321.7830,34.1433) and (322.3673,33.5590) ..(322.3673,32.8382);
\path[fill=chfourfig5c0,fill opacity=0.5,draw=chfourfig5c0,line width=0.3985pt,draw opacity=0.5] (273.9184,73.2472) .. controls (273.9184,72.5264) and (273.3340,71.9421) ..(272.6133,71.9421) .. controls (271.8925,71.9421) and (271.3081,72.5264) ..(271.3081,73.2472) .. controls (271.3081,73.9680) and (271.8925,74.5523) ..(272.6133,74.5523) .. controls (273.3340,74.5523) and (273.9184,73.9680) ..(273.9184,73.2472);
\path[fill=chfourfig5c0,fill opacity=0.5,draw=chfourfig5c0,line width=0.3985pt,draw opacity=0.5] (236.2093,78.3681) .. controls (236.2093,77.6473) and (235.6250,77.0629) ..(234.9042,77.0629) .. controls (234.1833,77.0629) and (233.5991,77.6473) ..(233.5991,78.3681) .. controls (233.5991,79.0889) and (234.1833,79.6732) ..(234.9042,79.6732) .. controls (235.6250,79.6732) and (236.2093,79.0889) ..(236.2093,78.3681);
\path[fill=chfourfig5c0,fill opacity=0.5,draw=chfourfig5c0,line width=0.3985pt,draw opacity=0.5] (245.1658,102.2787) .. controls (245.1658,101.5579) and (244.5815,100.9736) ..(243.8607,100.9736) .. controls (243.1399,100.9736) and (242.5556,101.5579) ..(242.5556,102.2787) .. controls (242.5556,102.9995) and (243.1399,103.5838) ..(243.8607,103.5838) .. controls (244.5815,103.5838) and (245.1658,102.9995) ..(245.1658,102.2787);
\path[fill=chfourfig5c0,fill opacity=0.5,draw=chfourfig5c0,line width=0.3985pt,draw opacity=0.5] (321.8194,63.2745) .. controls (321.8194,62.5536) and (321.2350,61.9693) ..(320.5143,61.9693) .. controls (319.7934,61.9693) and (319.2091,62.5536) ..(319.2091,63.2745) .. controls (319.2091,63.9953) and (319.7934,64.5796) ..(320.5143,64.5796) .. controls (321.2350,64.5796) and (321.8194,63.9953) ..(321.8194,63.2745);
\path[fill=chfourfig5c0,fill opacity=0.5,draw=chfourfig5c0,line width=0.3985pt,draw opacity=0.5] (288.2648,69.0728) .. controls (288.2648,68.3520) and (287.6805,67.7677) ..(286.9596,67.7677) .. controls (286.2388,67.7677) and (285.6545,68.3520) ..(285.6545,69.0728) .. controls (285.6545,69.7936) and (286.2388,70.3779) ..(286.9596,70.3779) .. controls (287.6805,70.3779) and (288.2648,69.7936) ..(288.2648,69.0728);
\path[fill=chfourfig5c0,fill opacity=0.5,draw=chfourfig5c0,line width=0.3985pt,draw opacity=0.5] (289.6496,58.9207) .. controls (289.6496,58.1999) and (289.0653,57.6156) ..(288.3445,57.6156) .. controls (287.6237,57.6156) and (287.0393,58.1999) ..(287.0393,58.9207) .. controls (287.0393,59.6415) and (287.6237,60.2258) ..(288.3445,60.2258) .. controls (289.0653,60.2258) and (289.6496,59.6415) ..(289.6496,58.9207);
\path[fill=chfourfig5c0,fill opacity=0.5,draw=chfourfig5c0,line width=0.3985pt,draw opacity=0.5] (240.7025,113.1780) .. controls (240.7025,112.4572) and (240.1182,111.8728) ..(239.3974,111.8728) .. controls (238.6766,111.8728) and (238.0923,112.4572) ..(238.0923,113.1780) .. controls (238.0923,113.8988) and (238.6766,114.4831) ..(239.3974,114.4831) .. controls (240.1182,114.4831) and (240.7025,113.8988) ..(240.7025,113.1780);
\path[fill=chfourfig5c0,fill opacity=0.5,draw=chfourfig5c0,line width=0.3985pt,draw opacity=0.5] (326.8605,33.4160) .. controls (326.8605,32.6952) and (326.2762,32.1109) ..(325.5554,32.1109) .. controls (324.8346,32.1109) and (324.2503,32.6952) ..(324.2503,33.4160) .. controls (324.2503,34.1368) and (324.8346,34.7212) ..(325.5554,34.7212) .. controls (326.2762,34.7212) and (326.8605,34.1368) ..(326.8605,33.4160);
\path[fill=chfourfig5c0,fill opacity=0.5,draw=chfourfig5c0,line width=0.3985pt,draw opacity=0.5] (244.9965,83.8874) .. controls (244.9965,83.1666) and (244.4121,82.5823) ..(243.6913,82.5823) .. controls (242.9705,82.5823) and (242.3862,83.1666) ..(242.3862,83.8874) .. controls (242.3862,84.6082) and (242.9705,85.1926) ..(243.6913,85.1926) .. controls (244.4121,85.1926) and (244.9965,84.6082) ..(244.9965,83.8874);
\path[fill=chfourfig5c0,fill opacity=0.5,draw=chfourfig5c0,line width=0.3985pt,draw opacity=0.5] (249.1708,72.5299) .. controls (249.1708,71.8091) and (248.5865,71.2248) ..(247.8657,71.2248) .. controls (247.1450,71.2248) and (246.5606,71.8091) ..(246.5606,72.5299) .. controls (246.5606,73.2507) and (247.1450,73.8350) ..(247.8657,73.8350) .. controls (248.5865,73.8350) and (249.1708,73.2507) ..(249.1708,72.5299);
\path[fill=chfourfig5c0,fill opacity=0.5,draw=chfourfig5c0,line width=0.3985pt,draw opacity=0.5] (252.8372,75.2397) .. controls (252.8372,74.5190) and (252.2528,73.9346) ..(251.5320,73.9346) .. controls (250.8112,73.9346) and (250.2269,74.5190) ..(250.2269,75.2397) .. controls (250.2269,75.9606) and (250.8112,76.5449) ..(251.5320,76.5449) .. controls (252.2528,76.5449) and (252.8372,75.9606) ..(252.8372,75.2397);
\path[fill=chfourfig5c0,fill opacity=0.5,draw=chfourfig5c0,line width=0.3985pt,draw opacity=0.5] (284.6383,52.3951) .. controls (284.6383,51.6743) and (284.0540,51.0900) ..(283.3332,51.0900) .. controls (282.6124,51.0900) and (282.0280,51.6743) ..(282.0280,52.3951) .. controls (282.0280,53.1159) and (282.6124,53.7002) ..(283.3332,53.7002) .. controls (284.0540,53.7002) and (284.6383,53.1159) ..(284.6383,52.3951);
\path[fill=chfourfig5c0,fill opacity=0.5,draw=chfourfig5c0,line width=0.3985pt,draw opacity=0.5] (270.4513,92.2362) .. controls (270.4513,91.5154) and (269.8670,90.9311) ..(269.1462,90.9311) .. controls (268.4254,90.9311) and (267.8411,91.5154) ..(267.8411,92.2362) .. controls (267.8411,92.9570) and (268.4254,93.5414) ..(269.1462,93.5414) .. controls (269.8670,93.5414) and (270.4513,92.9570) ..(270.4513,92.2362);
\path[fill=chfourfig5c0,fill opacity=0.5,draw=chfourfig5c0,line width=0.3985pt,draw opacity=0.5] (283.1240,81.2373) .. controls (283.1240,80.5165) and (282.5396,79.9322) ..(281.8188,79.9322) .. controls (281.0980,79.9322) and (280.5137,80.5165) ..(280.5137,81.2373) .. controls (280.5137,81.9581) and (281.0980,82.5424) ..(281.8188,82.5424) .. controls (282.5396,82.5424) and (283.1240,81.9581) ..(283.1240,81.2373);
\path[fill=chfourfig5c0,fill opacity=0.5,draw=chfourfig5c0,line width=0.3985pt,draw opacity=0.5] (261.4848,85.6110) .. controls (261.4848,84.8902) and (260.9005,84.3059) ..(260.1797,84.3059) .. controls (259.4589,84.3059) and (258.8746,84.8902) ..(258.8746,85.6110) .. controls (258.8746,86.3318) and (259.4589,86.9161) ..(260.1797,86.9161) .. controls (260.9005,86.9161) and (261.4848,86.3318) ..(261.4848,85.6110);
\path[fill=chfourfig5c0,fill opacity=0.5,draw=chfourfig5c0,line width=0.3985pt,draw opacity=0.5] (239.3675,100.0470) .. controls (239.3675,99.3262) and (238.7832,98.7419) ..(238.0624,98.7419) .. controls (237.3416,98.7419) and (236.7573,99.3262) ..(236.7573,100.0470) .. controls (236.7573,100.7678) and (237.3416,101.3522) ..(238.0624,101.3522) .. controls (238.7832,101.3522) and (239.3675,100.7678) ..(239.3675,100.0470);
\path[fill=chfourfig5c0,fill opacity=0.5,draw=chfourfig5c0,line width=0.3985pt,draw opacity=0.5] (294.5014,75.9571) .. controls (294.5014,75.2363) and (293.9171,74.6519) ..(293.1964,74.6519) .. controls (292.4755,74.6519) and (291.8912,75.2363) ..(291.8912,75.9571) .. controls (291.8912,76.6779) and (292.4755,77.2622) ..(293.1964,77.2622) .. controls (293.9171,77.2622) and (294.5014,76.6779) ..(294.5014,75.9571);
\path[fill=chfourfig5c0,fill opacity=0.5,draw=chfourfig5c0,line width=0.3985pt,draw opacity=0.5] (299.7020,41.5457) .. controls (299.7020,40.8249) and (299.1177,40.2405) ..(298.3969,40.2405) .. controls (297.6761,40.2405) and (297.0918,40.8249) ..(297.0918,41.5457) .. controls (297.0918,42.2665) and (297.6761,42.8508) ..(298.3969,42.8508) .. controls (299.1177,42.8508) and (299.7020,42.2665) ..(299.7020,41.5457);
\path[fill=chfourfig5c0,fill opacity=0.5,draw=chfourfig5c0,line width=0.3985pt,draw opacity=0.5] (287.0393,52.0962) .. controls (287.0393,51.3754) and (286.4550,50.7911) ..(285.7342,50.7911) .. controls (285.0134,50.7911) and (284.4291,51.3754) ..(284.4291,52.0962) .. controls (284.4291,52.8170) and (285.0134,53.4014) ..(285.7342,53.4014) .. controls (286.4550,53.4014) and (287.0393,52.8170) ..(287.0393,52.0962);
\path[fill=chfourfig5c0,fill opacity=0.5,draw=chfourfig5c0,line width=0.3985pt,draw opacity=0.5] (245.1260,100.6348) .. controls (245.1260,99.9140) and (244.5416,99.3297) ..(243.8209,99.3297) .. controls (243.1000,99.3297) and (242.5157,99.9140) ..(242.5157,100.6348) .. controls (242.5157,101.3556) and (243.1000,101.9400) ..(243.8209,101.9400) .. controls (244.5416,101.9400) and (245.1260,101.3556) ..(245.1260,100.6348);
\path[fill=chfourfig5c0,fill opacity=0.5,draw=chfourfig5c0,line width=0.3985pt,draw opacity=0.5] (254.1921,102.3484) .. controls (254.1921,101.6276) and (253.6078,101.0433) ..(252.8870,101.0433) .. controls (252.1662,101.0433) and (251.5818,101.6276) ..(251.5818,102.3484) .. controls (251.5818,103.0692) and (252.1662,103.6535) ..(252.8870,103.6535) .. controls (253.6078,103.6535) and (254.1921,103.0692) ..(254.1921,102.3484);
\path[fill=chfourfig5c0,fill opacity=0.5,draw=chfourfig5c0,line width=0.3985pt,draw opacity=0.5] (293.0967,62.0391) .. controls (293.0967,61.3183) and (292.5124,60.7340) ..(291.7916,60.7340) .. controls (291.0708,60.7340) and (290.4865,61.3183) ..(290.4865,62.0391) .. controls (290.4865,62.7599) and (291.0708,63.3442) ..(291.7916,63.3442) .. controls (292.5124,63.3442) and (293.0967,62.7599) ..(293.0967,62.0391);
\path[fill=chfourfig5c0,fill opacity=0.5,draw=chfourfig5c0,line width=0.3985pt,draw opacity=0.5] (325.9738,53.7799) .. controls (325.9738,53.0591) and (325.3895,52.4748) ..(324.6687,52.4748) .. controls (323.9479,52.4748) and (323.3636,53.0591) ..(323.3636,53.7799) .. controls (323.3636,54.5007) and (323.9479,55.0851) ..(324.6687,55.0851) .. controls (325.3895,55.0851) and (325.9738,54.5007) ..(325.9738,53.7799);
\path[fill=chfourfig5c0,fill opacity=0.5,draw=chfourfig5c0,line width=0.3985pt,draw opacity=0.5] (303.3882,65.1076) .. controls (303.3882,64.3868) and (302.8039,63.8025) ..(302.0831,63.8025) .. controls (301.3623,63.8025) and (300.7780,64.3868) ..(300.7780,65.1076) .. controls (300.7780,65.8284) and (301.3623,66.4127) ..(302.0831,66.4127) .. controls (302.8039,66.4127) and (303.3882,65.8284) ..(303.3882,65.1076);
\path[fill=chfourfig5c0,fill opacity=0.5,draw=chfourfig5c0,line width=0.3985pt,draw opacity=0.5] (297.9087,47.1547) .. controls (297.9087,46.4339) and (297.3244,45.8496) ..(296.6036,45.8496) .. controls (295.8828,45.8496) and (295.2985,46.4339) ..(295.2985,47.1547) .. controls (295.2985,47.8755) and (295.8828,48.4598) ..(296.6036,48.4598) .. controls (297.3244,48.4598) and (297.9087,47.8755) ..(297.9087,47.1547);
\path[fill=chfourfig5c0,fill opacity=0.5,draw=chfourfig5c0,line width=0.3985pt,draw opacity=0.5] (269.3455,97.0283) .. controls (269.3455,96.3075) and (268.7612,95.7232) ..(268.0403,95.7232) .. controls (267.3195,95.7232) and (266.7352,96.3075) ..(266.7352,97.0283) .. controls (266.7352,97.7491) and (267.3195,98.3334) ..(268.0403,98.3334) .. controls (268.7612,98.3334) and (269.3455,97.7491) ..(269.3455,97.0283);
\path[fill=chfourfig5c0,fill opacity=0.5,draw=chfourfig5c0,line width=0.3985pt,draw opacity=0.5] (265.5397,64.1811) .. controls (265.5397,63.4603) and (264.9554,62.8759) ..(264.2346,62.8759) .. controls (263.5138,62.8759) and (262.9294,63.4603) ..(262.9294,64.1811) .. controls (262.9294,64.9019) and (263.5138,65.4862) ..(264.2346,65.4862) .. controls (264.9554,65.4862) and (265.5397,64.9019) ..(265.5397,64.1811);
\path[fill=chfourfig5c0,fill opacity=0.5,draw=chfourfig5c0,line width=0.3985pt,draw opacity=0.5] (290.8949,62.1586) .. controls (290.8949,61.4378) and (290.3106,60.8535) ..(289.5898,60.8535) .. controls (288.8690,60.8535) and (288.2847,61.4378) ..(288.2847,62.1586) .. controls (288.2847,62.8794) and (288.8690,63.4637) ..(289.5898,63.4637) .. controls (290.3106,63.4637) and (290.8949,62.8794) ..(290.8949,62.1586);
\path[fill=chfourfig5c0,fill opacity=0.5,draw=chfourfig5c0,line width=0.3985pt,draw opacity=0.5] (247.5868,95.1752) .. controls (247.5868,94.4544) and (247.0024,93.8701) ..(246.2816,93.8701) .. controls (245.5609,93.8701) and (244.9765,94.4544) ..(244.9765,95.1752) .. controls (244.9765,95.8960) and (245.5609,96.4804) ..(246.2816,96.4804) .. controls (247.0024,96.4804) and (247.5868,95.8960) ..(247.5868,95.1752);
\path[fill=chfourfig5c0,fill opacity=0.5,draw=chfourfig5c0,line width=0.3985pt,draw opacity=0.5] (317.6748,48.4897) .. controls (317.6748,47.7689) and (317.0905,47.1846) ..(316.3698,47.1846) .. controls (315.6489,47.1846) and (315.0646,47.7689) ..(315.0646,48.4897) .. controls (315.0646,49.2105) and (315.6489,49.7948) ..(316.3698,49.7948) .. controls (317.0905,49.7948) and (317.6748,49.2105) ..(317.6748,48.4897);
\path[fill=chfourfig5c0,fill opacity=0.5,draw=chfourfig5c0,line width=0.3985pt,draw opacity=0.5] (241.2206,83.4590) .. controls (241.2206,82.7382) and (240.6362,82.1539) ..(239.9154,82.1539) .. controls (239.1946,82.1539) and (238.6103,82.7382) ..(238.6103,83.4590) .. controls (238.6103,84.1798) and (239.1946,84.7642) ..(239.9154,84.7642) .. controls (240.6362,84.7642) and (241.2206,84.1798) ..(241.2206,83.4590);
\path[fill=chfourfig5c0,fill opacity=0.5,draw=chfourfig5c0,line width=0.3985pt,draw opacity=0.5] (285.5250,50.6118) .. controls (285.5250,49.8910) and (284.9407,49.3067) ..(284.2199,49.3067) .. controls (283.4991,49.3067) and (282.9147,49.8910) ..(282.9147,50.6118) .. controls (282.9147,51.3326) and (283.4991,51.9169) ..(284.2199,51.9169) .. controls (284.9407,51.9169) and (285.5250,51.3326) ..(285.5250,50.6118);
\path[fill=chfourfig5c0,fill opacity=0.5,draw=chfourfig5c0,line width=0.3985pt,draw opacity=0.5] (275.7615,73.5560) .. controls (275.7615,72.8352) and (275.1772,72.2509) ..(274.4564,72.2509) .. controls (273.7355,72.2509) and (273.1512,72.8352) ..(273.1512,73.5560) .. controls (273.1512,74.2768) and (273.7355,74.8612) ..(274.4564,74.8612) .. controls (275.1772,74.8612) and (275.7615,74.2768) ..(275.7615,73.5560);
\path[fill=chfourfig5c0,fill opacity=0.5,draw=chfourfig5c0,line width=0.3985pt,draw opacity=0.5] (311.6175,39.2343) .. controls (311.6175,38.5135) and (311.0332,37.9292) ..(310.3124,37.9292) .. controls (309.5916,37.9292) and (309.0072,38.5135) ..(309.0072,39.2343) .. controls (309.0072,39.9551) and (309.5916,40.5394) ..(310.3124,40.5394) .. controls (311.0332,40.5394) and (311.6175,39.9551) ..(311.6175,39.2343);
\path[fill=chfourfig5c0,fill opacity=0.5,draw=chfourfig5c0,line width=0.3985pt,draw opacity=0.5] (293.6646,63.0852) .. controls (293.6646,62.3644) and (293.0803,61.7800) ..(292.3594,61.7800) .. controls (291.6387,61.7800) and (291.0544,62.3644) ..(291.0544,63.0852) .. controls (291.0544,63.8060) and (291.6387,64.3903) ..(292.3594,64.3903) .. controls (293.0803,64.3903) and (293.6646,63.8060) ..(293.6646,63.0852);
\path[fill=chfourfig5c0,fill opacity=0.5,draw=chfourfig5c0,line width=0.3985pt,draw opacity=0.5] (274.4265,69.4414) .. controls (274.4265,68.7206) and (273.8421,68.1363) ..(273.1214,68.1363) .. controls (272.4005,68.1363) and (271.8162,68.7206) ..(271.8162,69.4414) .. controls (271.8162,70.1622) and (272.4005,70.7465) ..(273.1214,70.7465) .. controls (273.8421,70.7465) and (274.4265,70.1622) ..(274.4265,69.4414);
\path[fill=chfourfig5c0,fill opacity=0.5,draw=chfourfig5c0,line width=0.3985pt,draw opacity=0.5] (299.0345,62.0092) .. controls (299.0345,61.2884) and (298.4502,60.7041) ..(297.7294,60.7041) .. controls (297.0086,60.7041) and (296.4243,61.2884) ..(296.4243,62.0092) .. controls (296.4243,62.7300) and (297.0086,63.3143) ..(297.7294,63.3143) .. controls (298.4502,63.3143) and (299.0345,62.7300) ..(299.0345,62.0092);
\path[fill=chfourfig5c0,fill opacity=0.5,draw=chfourfig5c0,line width=0.3985pt,draw opacity=0.5] (313.1119,48.0115) .. controls (313.1119,47.2907) and (312.5276,46.7064) ..(311.8068,46.7064) .. controls (311.0860,46.7064) and (310.5017,47.2907) ..(310.5017,48.0115) .. controls (310.5017,48.7323) and (311.0860,49.3166) ..(311.8068,49.3166) .. controls (312.5276,49.3166) and (313.1119,48.7323) ..(313.1119,48.0115);
\path[fill=chfourfig5c0,fill opacity=0.5,draw=chfourfig5c0,line width=0.3985pt,draw opacity=0.5] (327.0797,44.6839) .. controls (327.0797,43.9631) and (326.4954,43.3788) ..(325.7746,43.3788) .. controls (325.0538,43.3788) and (324.4695,43.9631) ..(324.4695,44.6839) .. controls (324.4695,45.4047) and (325.0538,45.9890) ..(325.7746,45.9890) .. controls (326.4954,45.9890) and (327.0797,45.4047) ..(327.0797,44.6839);
\path[fill=chfourfig5c0,fill opacity=0.5,draw=chfourfig5c0,line width=0.3985pt,draw opacity=0.5] (302.0034,61.1723) .. controls (302.0034,60.4515) and (301.4191,59.8672) ..(300.6983,59.8672) .. controls (299.9775,59.8672) and (299.3932,60.4515) ..(299.3932,61.1723) .. controls (299.3932,61.8931) and (299.9775,62.4774) ..(300.6983,62.4774) .. controls (301.4191,62.4774) and (302.0034,61.8931) ..(302.0034,61.1723);
\path[fill=chfourfig5c0,fill opacity=0.5,draw=chfourfig5c0,line width=0.3985pt,draw opacity=0.5] (256.8422,79.7330) .. controls (256.8422,79.0121) and (256.2579,78.4278) ..(255.5370,78.4278) .. controls (254.8163,78.4278) and (254.2319,79.0121) ..(254.2319,79.7330) .. controls (254.2319,80.4538) and (254.8163,81.0381) ..(255.5370,81.0381) .. controls (256.2579,81.0381) and (256.8422,80.4538) ..(256.8422,79.7330);
\path[fill=chfourfig5c0,fill opacity=0.5,draw=chfourfig5c0,line width=0.3985pt,draw opacity=0.5] (259.6118,88.1116) .. controls (259.6118,87.3908) and (259.0275,86.8065) ..(258.3067,86.8065) .. controls (257.5859,86.8065) and (257.0016,87.3908) ..(257.0016,88.1116) .. controls (257.0016,88.8324) and (257.5859,89.4168) ..(258.3067,89.4168) .. controls (259.0275,89.4168) and (259.6118,88.8324) ..(259.6118,88.1116);
\path[fill=chfourfig5c0,fill opacity=0.5,draw=chfourfig5c0,line width=0.3985pt,draw opacity=0.5] (273.3505,94.5974) .. controls (273.3505,93.8766) and (272.7662,93.2923) ..(272.0454,93.2923) .. controls (271.3246,93.2923) and (270.7402,93.8766) ..(270.7402,94.5974) .. controls (270.7402,95.3182) and (271.3246,95.9025) ..(272.0454,95.9025) .. controls (272.7662,95.9025) and (273.3505,95.3182) ..(273.3505,94.5974);
\path[fill=chfourfig5c0,fill opacity=0.5,draw=chfourfig5c0,line width=0.3985pt,draw opacity=0.5] (305.4406,62.6568) .. controls (305.4406,61.9360) and (304.8562,61.3517) ..(304.1355,61.3517) .. controls (303.4146,61.3517) and (302.8303,61.9360) ..(302.8303,62.6568) .. controls (302.8303,63.3776) and (303.4146,63.9619) ..(304.1355,63.9619) .. controls (304.8562,63.9619) and (305.4406,63.3776) ..(305.4406,62.6568);
\path[fill=chfourfig5c0,fill opacity=0.5,draw=chfourfig5c0,line width=0.3985pt,draw opacity=0.5] (290.4565,65.8747) .. controls (290.4565,65.1539) and (289.8722,64.5696) ..(289.1515,64.5696) .. controls (288.4306,64.5696) and (287.8463,65.1539) ..(287.8463,65.8747) .. controls (287.8463,66.5955) and (288.4306,67.1799) ..(289.1515,67.1799) .. controls (289.8722,67.1799) and (290.4565,66.5955) ..(290.4565,65.8747);
\path[fill=chfourfig5c0,fill opacity=0.5,draw=chfourfig5c0,line width=0.3985pt,draw opacity=0.5] (315.0347,37.0624) .. controls (315.0347,36.3416) and (314.4504,35.7573) ..(313.7296,35.7573) .. controls (313.0088,35.7573) and (312.4245,36.3416) ..(312.4245,37.0624) .. controls (312.4245,37.7832) and (313.0088,38.3675) ..(313.7296,38.3675) .. controls (314.4504,38.3675) and (315.0347,37.7832) ..(315.0347,37.0624);
\path[fill=chfourfig5c0,fill opacity=0.5,draw=chfourfig5c0,line width=0.3985pt,draw opacity=0.5] (276.5884,54.4674) .. controls (276.5884,53.7466) and (276.0041,53.1622) ..(275.2833,53.1622) .. controls (274.5625,53.1622) and (273.9781,53.7466) ..(273.9781,54.4674) .. controls (273.9781,55.1882) and (274.5625,55.7725) ..(275.2833,55.7725) .. controls (276.0041,55.7725) and (276.5884,55.1882) ..(276.5884,54.4674);
\path[fill=chfourfig5c0,fill opacity=0.5,draw=chfourfig5c0,line width=0.3985pt,draw opacity=0.5] (300.2101,77.6408) .. controls (300.2101,76.9200) and (299.6258,76.3356) ..(298.9050,76.3356) .. controls (298.1842,76.3356) and (297.5999,76.9200) ..(297.5999,77.6408) .. controls (297.5999,78.3616) and (298.1842,78.9459) ..(298.9050,78.9459) .. controls (299.6258,78.9459) and (300.2101,78.3616) ..(300.2101,77.6408);
\path[fill=chfourfig5c0,fill opacity=0.5,draw=chfourfig5c0,line width=0.3985pt,draw opacity=0.5] (306.8652,75.9471) .. controls (306.8652,75.2263) and (306.2809,74.6420) ..(305.5601,74.6420) .. controls (304.8393,74.6420) and (304.2550,75.2263) ..(304.2550,75.9471) .. controls (304.2550,76.6679) and (304.8393,77.2522) ..(305.5601,77.2522) .. controls (306.2809,77.2522) and (306.8652,76.6679) ..(306.8652,75.9471);
\path[fill=chfourfig5c0,fill opacity=0.5,draw=chfourfig5c0,line width=0.3985pt,draw opacity=0.5] (314.1879,61.5609) .. controls (314.1879,60.8401) and (313.6035,60.2557) ..(312.8828,60.2557) .. controls (312.1620,60.2557) and (311.5776,60.8401) ..(311.5776,61.5609) .. controls (311.5776,62.2817) and (312.1620,62.8660) ..(312.8828,62.8660) .. controls (313.6035,62.8660) and (314.1879,62.2817) ..(314.1879,61.5609);
\path[fill=chfourfig5c0,fill opacity=0.5,draw=chfourfig5c0,line width=0.3985pt,draw opacity=0.5] (298.0880,52.1759) .. controls (298.0880,51.4551) and (297.5037,50.8708) ..(296.7829,50.8708) .. controls (296.0621,50.8708) and (295.4778,51.4551) ..(295.4778,52.1759) .. controls (295.4778,52.8967) and (296.0621,53.4810) ..(296.7829,53.4810) .. controls (297.5037,53.4810) and (298.0880,52.8967) ..(298.0880,52.1759);
\path[fill=chfourfig5c0,fill opacity=0.5,draw=chfourfig5c0,line width=0.3985pt,draw opacity=0.5] (255.1983,91.5189) .. controls (255.1983,90.7981) and (254.6140,90.2138) ..(253.8932,90.2138) .. controls (253.1724,90.2138) and (252.5881,90.7981) ..(252.5881,91.5189) .. controls (252.5881,92.2397) and (253.1724,92.8240) ..(253.8932,92.8240) .. controls (254.6140,92.8240) and (255.1983,92.2397) ..(255.1983,91.5189);
\path[fill=chfourfig5c0,fill opacity=0.5,draw=chfourfig5c0,line width=0.3985pt,draw opacity=0.5] (254.5408,93.6808) .. controls (254.5408,92.9600) and (253.9565,92.3757) ..(253.2357,92.3757) .. controls (252.5149,92.3757) and (251.9305,92.9600) ..(251.9305,93.6808) .. controls (251.9305,94.4016) and (252.5149,94.9859) ..(253.2357,94.9859) .. controls (253.9565,94.9859) and (254.5408,94.4016) ..(254.5408,93.6808);
\path[fill=chfourfig5c0,fill opacity=0.5,draw=chfourfig5c0,line width=0.3985pt,draw opacity=0.5] (326.1332,45.3315) .. controls (326.1332,44.6107) and (325.5490,44.0264) ..(324.8281,44.0264) .. controls (324.1073,44.0264) and (323.5230,44.6107) ..(323.5230,45.3315) .. controls (323.5230,46.0523) and (324.1073,46.6366) ..(324.8281,46.6366) .. controls (325.5490,46.6366) and (326.1332,46.0523) ..(326.1332,45.3315);
\path[fill=chfourfig5c0,fill opacity=0.5,draw=chfourfig5c0,line width=0.3985pt,draw opacity=0.5] (232.9913,112.8492) .. controls (232.9913,112.1284) and (232.4070,111.5441) ..(231.6862,111.5441) .. controls (230.9654,111.5441) and (230.3811,112.1284) ..(230.3811,112.8492) .. controls (230.3811,113.5700) and (230.9654,114.1543) ..(231.6862,114.1543) .. controls (232.4070,114.1543) and (232.9913,113.5700) ..(232.9913,112.8492);
\path[fill=chfourfig5c0,fill opacity=0.5,draw=chfourfig5c0,line width=0.3985pt,draw opacity=0.5] (308.1903,55.9419) .. controls (308.1903,55.2211) and (307.6060,54.6367) ..(306.8852,54.6367) .. controls (306.1644,54.6367) and (305.5800,55.2211) ..(305.5800,55.9419) .. controls (305.5800,56.6627) and (306.1644,57.2470) ..(306.8852,57.2470) .. controls (307.6060,57.2470) and (308.1903,56.6627) ..(308.1903,55.9419);
\path[fill=chfourfig5c0,fill opacity=0.5,draw=chfourfig5c0,line width=0.3985pt,draw opacity=0.5] (262.1823,82.9609) .. controls (262.1823,82.2401) and (261.5979,81.6558) ..(260.8771,81.6558) .. controls (260.1563,81.6558) and (259.5720,82.2401) ..(259.5720,82.9609) .. controls (259.5720,83.6817) and (260.1563,84.2660) ..(260.8771,84.2660) .. controls (261.5979,84.2660) and (262.1823,83.6817) ..(262.1823,82.9609);
\path[fill=chfourfig5c0,fill opacity=0.5,draw=chfourfig5c0,line width=0.3985pt,draw opacity=0.5] (278.5710,55.5633) .. controls (278.5710,54.8425) and (277.9867,54.2582) ..(277.2659,54.2582) .. controls (276.5450,54.2582) and (275.9607,54.8425) ..(275.9607,55.5633) .. controls (275.9607,56.2841) and (276.5450,56.8684) ..(277.2659,56.8684) .. controls (277.9867,56.8684) and (278.5710,56.2841) ..(278.5710,55.5633);
\path[fill=chfourfig5c0,fill opacity=0.5,draw=chfourfig5c0,line width=0.3985pt,draw opacity=0.5] (291.5226,58.9307) .. controls (291.5226,58.2099) and (290.9383,57.6256) ..(290.2175,57.6256) .. controls (289.4966,57.6256) and (288.9123,58.2099) ..(288.9123,58.9307) .. controls (288.9123,59.6515) and (289.4966,60.2358) ..(290.2175,60.2358) .. controls (290.9383,60.2358) and (291.5226,59.6515) ..(291.5226,58.9307);
\path[fill=chfourfig5c0,fill opacity=0.5,draw=chfourfig5c0,line width=0.3985pt,draw opacity=0.5] (248.6229,73.9545) .. controls (248.6229,73.2337) and (248.0386,72.6494) ..(247.3178,72.6494) .. controls (246.5970,72.6494) and (246.0127,73.2337) ..(246.0127,73.9545) .. controls (246.0127,74.6753) and (246.5970,75.2597) ..(247.3178,75.2597) .. controls (248.0386,75.2597) and (248.6229,74.6753) ..(248.6229,73.9545);
\path[fill=chfourfig5c0,fill opacity=0.5,draw=chfourfig5c0,line width=0.3985pt,draw opacity=0.5] (293.2162,50.7612) .. controls (293.2162,50.0404) and (292.6319,49.4561) ..(291.9111,49.4561) .. controls (291.1903,49.4561) and (290.6060,50.0404) ..(290.6060,50.7612) .. controls (290.6060,51.4820) and (291.1903,52.0663) ..(291.9111,52.0663) .. controls (292.6319,52.0663) and (293.2162,51.4820) ..(293.2162,50.7612);
\path[fill=chfourfig5c0,fill opacity=0.5,draw=chfourfig5c0,line width=0.3985pt,draw opacity=0.5] (234.4558,109.0235) .. controls (234.4558,108.3027) and (233.8715,107.7184) ..(233.1507,107.7184) .. controls (232.4299,107.7184) and (231.8456,108.3027) ..(231.8456,109.0235) .. controls (231.8456,109.7443) and (232.4299,110.3286) ..(233.1507,110.3286) .. controls (233.8715,110.3286) and (234.4558,109.7443) ..(234.4558,109.0235);
\path[fill=chfourfig5c0,fill opacity=0.5,draw=chfourfig5c0,line width=0.3985pt,draw opacity=0.5] (272.0653,81.6558) .. controls (272.0653,80.9350) and (271.4810,80.3506) ..(270.7602,80.3506) .. controls (270.0394,80.3506) and (269.4551,80.9350) ..(269.4551,81.6558) .. controls (269.4551,82.3766) and (270.0394,82.9609) ..(270.7602,82.9609) .. controls (271.4810,82.9609) and (272.0653,82.3766) ..(272.0653,81.6558);
\path[fill=chfourfig5c0,fill opacity=0.5,draw=chfourfig5c0,line width=0.3985pt,draw opacity=0.5] (248.8221,101.5813) .. controls (248.8221,100.8605) and (248.2379,100.2762) ..(247.5170,100.2762) .. controls (246.7962,100.2762) and (246.2119,100.8605) ..(246.2119,101.5813) .. controls (246.2119,102.3021) and (246.7962,102.8864) ..(247.5170,102.8864) .. controls (248.2379,102.8864) and (248.8221,102.3021) ..(248.8221,101.5813);
\path[fill=chfourfig5c0,fill opacity=0.5,draw=chfourfig5c0,line width=0.3985pt,draw opacity=0.5] (260.8771,78.7168) .. controls (260.8771,77.9959) and (260.2928,77.4116) ..(259.5720,77.4116) .. controls (258.8512,77.4116) and (258.2669,77.9959) ..(258.2669,78.7168) .. controls (258.2669,79.4375) and (258.8512,80.0219) ..(259.5720,80.0219) .. controls (260.2928,80.0219) and (260.8771,79.4375) ..(260.8771,78.7168);
\path[fill=chfourfig5c0,fill opacity=0.5,draw=chfourfig5c0,line width=0.3985pt,draw opacity=0.5] (324.9975,33.8843) .. controls (324.9975,33.1635) and (324.4132,32.5792) ..(323.6924,32.5792) .. controls (322.9716,32.5792) and (322.3873,33.1635) ..(322.3873,33.8843) .. controls (322.3873,34.6051) and (322.9716,35.1894) ..(323.6924,35.1894) .. controls (324.4132,35.1894) and (324.9975,34.6051) ..(324.9975,33.8843);
\path[fill=chfourfig5c0,fill opacity=0.5,draw=chfourfig5c0,line width=0.3985pt,draw opacity=0.5] (329.0623,59.7178) .. controls (329.0623,58.9969) and (328.4780,58.4126) ..(327.7572,58.4126) .. controls (327.0363,58.4126) and (326.4520,58.9969) ..(326.4520,59.7178) .. controls (326.4520,60.4386) and (327.0363,61.0229) ..(327.7572,61.0229) .. controls (328.4780,61.0229) and (329.0623,60.4386) ..(329.0623,59.7178);
\path[fill=chfourfig5c0,fill opacity=0.5,draw=chfourfig5c0,line width=0.3985pt,draw opacity=0.5] (321.6599,54.1187) .. controls (321.6599,53.3979) and (321.0756,52.8136) ..(320.3549,52.8136) .. controls (319.6340,52.8136) and (319.0497,53.3979) ..(319.0497,54.1187) .. controls (319.0497,54.8395) and (319.6340,55.4238) ..(320.3549,55.4238) .. controls (321.0756,55.4238) and (321.6599,54.8395) ..(321.6599,54.1187);
\path[fill=chfourfig5c0,fill opacity=0.5,draw=chfourfig5c0,line width=0.3985pt,draw opacity=0.5] (301.9138,69.5410) .. controls (301.9138,68.8202) and (301.3294,68.2359) ..(300.6086,68.2359) .. controls (299.8878,68.2359) and (299.3035,68.8202) ..(299.3035,69.5410) .. controls (299.3035,70.2618) and (299.8878,70.8462) ..(300.6086,70.8462) .. controls (301.3294,70.8462) and (301.9138,70.2618) ..(301.9138,69.5410);
\path[fill=chfourfig5c0,fill opacity=0.5,draw=chfourfig5c0,line width=0.3985pt,draw opacity=0.5] (229.2951,104.1417) .. controls (229.2951,103.4209) and (228.7108,102.8366) ..(227.9900,102.8366) .. controls (227.2692,102.8366) and (226.6849,103.4209) ..(226.6849,104.1417) .. controls (226.6849,104.8625) and (227.2692,105.4469) ..(227.9900,105.4469) .. controls (228.7108,105.4469) and (229.2951,104.8625) ..(229.2951,104.1417);
\path[fill=chfourfig5c0,fill opacity=0.5,draw=chfourfig5c0,line width=0.3985pt,draw opacity=0.5] (239.2878,97.9449) .. controls (239.2878,97.2241) and (238.7035,96.6398) ..(237.9827,96.6398) .. controls (237.2618,96.6398) and (236.6776,97.2241) ..(236.6776,97.9449) .. controls (236.6776,98.6657) and (237.2618,99.2500) ..(237.9827,99.2500) .. controls (238.7035,99.2500) and (239.2878,98.6657) ..(239.2878,97.9449);
\path[fill=chfourfig5c0,fill opacity=0.5,draw=chfourfig5c0,line width=0.3985pt,draw opacity=0.5] (290.7554,84.2660) .. controls (290.7554,83.5452) and (290.1711,82.9609) ..(289.4503,82.9609) .. controls (288.7295,82.9609) and (288.1452,83.5452) ..(288.1452,84.2660) .. controls (288.1452,84.9868) and (288.7295,85.5711) ..(289.4503,85.5711) .. controls (290.1711,85.5711) and (290.7554,84.9868) ..(290.7554,84.2660);
\path[fill=chfourfig5c0,fill opacity=0.5,draw=chfourfig5c0,line width=0.3985pt,draw opacity=0.5] (234.4558,79.7529) .. controls (234.4558,79.0321) and (233.8715,78.4478) ..(233.1507,78.4478) .. controls (232.4299,78.4478) and (231.8456,79.0321) ..(231.8456,79.7529) .. controls (231.8456,80.4737) and (232.4299,81.0580) ..(233.1507,81.0580) .. controls (233.8715,81.0580) and (234.4558,80.4737) ..(234.4558,79.7529);
\path[fill=chfourfig5c0,fill opacity=0.5,draw=chfourfig5c0,line width=0.3985pt,draw opacity=0.5] (267.4525,95.9025) .. controls (267.4525,95.1817) and (266.8682,94.5974) ..(266.1474,94.5974) .. controls (265.4266,94.5974) and (264.8423,95.1817) ..(264.8423,95.9025) .. controls (264.8423,96.6233) and (265.4266,97.2076) ..(266.1474,97.2076) .. controls (266.8682,97.2076) and (267.4525,96.6233) ..(267.4525,95.9025);
\path[fill=chfourfig5c0,fill opacity=0.5,draw=chfourfig5c0,line width=0.3985pt,draw opacity=0.5] (276.6581,61.4314) .. controls (276.6581,60.7106) and (276.0738,60.1262) ..(275.3530,60.1262) .. controls (274.6322,60.1262) and (274.0479,60.7106) ..(274.0479,61.4314) .. controls (274.0479,62.1522) and (274.6322,62.7365) ..(275.3530,62.7365) .. controls (276.0738,62.7365) and (276.6581,62.1522) ..(276.6581,61.4314);
\path[fill=chfourfig5c0,fill opacity=0.5,draw=chfourfig5c0,line width=0.3985pt,draw opacity=0.5] (291.4030,46.8160) .. controls (291.4030,46.0952) and (290.8187,45.5108) ..(290.0979,45.5108) .. controls (289.3771,45.5108) and (288.7928,46.0952) ..(288.7928,46.8160) .. controls (288.7928,47.5367) and (289.3771,48.1211) ..(290.0979,48.1211) .. controls (290.8187,48.1211) and (291.4030,47.5367) ..(291.4030,46.8160);
\path[fill=chfourfig5c0,fill opacity=0.5,draw=chfourfig5c0,line width=0.3985pt,draw opacity=0.5] (253.3752,87.0357) .. controls (253.3752,86.3149) and (252.7908,85.7305) ..(252.0700,85.7305) .. controls (251.3492,85.7305) and (250.7649,86.3149) ..(250.7649,87.0357) .. controls (250.7649,87.7565) and (251.3492,88.3408) ..(252.0700,88.3408) .. controls (252.7908,88.3408) and (253.3752,87.7565) ..(253.3752,87.0357);
\path[fill=chfourfig5c0,fill opacity=0.5,draw=chfourfig5c0,line width=0.3985pt,draw opacity=0.5] (320.9028,38.3675) .. controls (320.9028,37.6467) and (320.3185,37.0624) ..(319.5977,37.0624) .. controls (318.8769,37.0624) and (318.2926,37.6467) ..(318.2926,38.3675) .. controls (318.2926,39.0883) and (318.8769,39.6727) ..(319.5977,39.6727) .. controls (320.3185,39.6727) and (320.9028,39.0883) ..(320.9028,38.3675);
\path[fill=chfourfig5c0,fill opacity=0.5,draw=chfourfig5c0,line width=0.3985pt,draw opacity=0.5] (314.9849,44.0563) .. controls (314.9849,43.3355) and (314.4006,42.7512) ..(313.6798,42.7512) .. controls (312.9590,42.7512) and (312.3747,43.3355) ..(312.3747,44.0563) .. controls (312.3747,44.7771) and (312.9590,45.3614) ..(313.6798,45.3614) .. controls (314.4006,45.3614) and (314.9849,44.7771) ..(314.9849,44.0563);
\path[fill=chfourfig5c0,fill opacity=0.5,draw=chfourfig5c0,line width=0.3985pt,draw opacity=0.5] (320.4943,50.6417) .. controls (320.4943,49.9209) and (319.9100,49.3365) ..(319.1892,49.3365) .. controls (318.4684,49.3365) and (317.8841,49.9209) ..(317.8841,50.6417) .. controls (317.8841,51.3625) and (318.4684,51.9468) ..(319.1892,51.9468) .. controls (319.9100,51.9468) and (320.4943,51.3625) ..(320.4943,50.6417);
\path[fill=chfourfig5c0,fill opacity=0.5,draw=chfourfig5c0,line width=0.3985pt,draw opacity=0.5] (285.0667,70.6768) .. controls (285.0667,69.9560) and (284.4824,69.3717) ..(283.7616,69.3717) .. controls (283.0408,69.3717) and (282.4565,69.9560) ..(282.4565,70.6768) .. controls (282.4565,71.3976) and (283.0408,71.9819) ..(283.7616,71.9819) .. controls (284.4824,71.9819) and (285.0667,71.3976) ..(285.0667,70.6768);
\path[fill=chfourfig5c0,fill opacity=0.5,draw=chfourfig5c0,line width=0.3985pt,draw opacity=0.5] (295.3383,43.3091) .. controls (295.3383,42.5883) and (294.7540,42.0040) ..(294.0332,42.0040) .. controls (293.3124,42.0040) and (292.7281,42.5883) ..(292.7281,43.3091) .. controls (292.7281,44.0299) and (293.3124,44.6142) ..(294.0332,44.6142) .. controls (294.7540,44.6142) and (295.3383,44.0299) ..(295.3383,43.3091);
\path[fill=chfourfig5c0,fill opacity=0.5,draw=chfourfig5c0,line width=0.3985pt,draw opacity=0.5] (284.2697,62.8261) .. controls (284.2697,62.1053) and (283.6854,61.5210) ..(282.9646,61.5210) .. controls (282.2438,61.5210) and (281.6595,62.1053) ..(281.6595,62.8261) .. controls (281.6595,63.5469) and (282.2438,64.1313) ..(282.9646,64.1313) .. controls (283.6854,64.1313) and (284.2697,63.5469) ..(284.2697,62.8261);
\path[fill=chfourfig5c0,fill opacity=0.5,draw=chfourfig5c0,line width=0.3985pt,draw opacity=0.5] (271.5771,65.1674) .. controls (271.5771,64.4466) and (270.9928,63.8623) ..(270.2720,63.8623) .. controls (269.5512,63.8623) and (268.9669,64.4466) ..(268.9669,65.1674) .. controls (268.9669,65.8882) and (269.5512,66.4725) ..(270.2720,66.4725) .. controls (270.9928,66.4725) and (271.5771,65.8882) ..(271.5771,65.1674);
\path[fill=chfourfig5c0,fill opacity=0.5,draw=chfourfig5c0,line width=0.3985pt,draw opacity=0.5] (282.6657,56.3105) .. controls (282.6657,55.5897) and (282.0814,55.0054) ..(281.3606,55.0054) .. controls (280.6398,55.0054) and (280.0555,55.5897) ..(280.0555,56.3105) .. controls (280.0555,57.0313) and (280.6398,57.6156) ..(281.3606,57.6156) .. controls (282.0814,57.6156) and (282.6657,57.0313) ..(282.6657,56.3105);
\path[fill=chfourfig5c0,fill opacity=0.5,draw=chfourfig5c0,line width=0.3985pt,draw opacity=0.5] (255.3079,91.6385) .. controls (255.3079,90.9177) and (254.7236,90.3333) ..(254.0028,90.3333) .. controls (253.2820,90.3333) and (252.6977,90.9177) ..(252.6977,91.6385) .. controls (252.6977,92.3593) and (253.2820,92.9436) ..(254.0028,92.9436) .. controls (254.7236,92.9436) and (255.3079,92.3593) ..(255.3079,91.6385);
\path[fill=chfourfig5c0,fill opacity=0.5,draw=chfourfig5c0,line width=0.3985pt,draw opacity=0.5] (327.0299,50.8907) .. controls (327.0299,50.1699) and (326.4456,49.5856) ..(325.7248,49.5856) .. controls (325.0040,49.5856) and (324.4196,50.1699) ..(324.4196,50.8907) .. controls (324.4196,51.6115) and (325.0040,52.1959) ..(325.7248,52.1959) .. controls (326.4456,52.1959) and (327.0299,51.6115) ..(327.0299,50.8907);
\path[fill=chfourfig5c0,fill opacity=0.5,draw=chfourfig5c0,line width=0.3985pt,draw opacity=0.5] (324.7584,55.4935) .. controls (324.7584,54.7727) and (324.1740,54.1884) ..(323.4533,54.1884) .. controls (322.7325,54.1884) and (322.1481,54.7727) ..(322.1481,55.4935) .. controls (322.1481,56.2143) and (322.7325,56.7987) ..(323.4533,56.7987) .. controls (324.1740,56.7987) and (324.7584,56.2143) ..(324.7584,55.4935);
\path[fill=chfourfig5c0,fill opacity=0.5,draw=chfourfig5c0,line width=0.3985pt,draw opacity=0.5] (239.7261,96.2313) .. controls (239.7261,95.5105) and (239.1418,94.9262) ..(238.4210,94.9262) .. controls (237.7002,94.9262) and (237.1159,95.5105) ..(237.1159,96.2313) .. controls (237.1159,96.9521) and (237.7002,97.5364) ..(238.4210,97.5364) .. controls (239.1418,97.5364) and (239.7261,96.9521) ..(239.7261,96.2313);
\path[fill=chfourfig5c0,fill opacity=0.5,draw=chfourfig5c0,line width=0.3985pt,draw opacity=0.5] (254.2020,89.0182) .. controls (254.2020,88.2974) and (253.6177,87.7131) ..(252.8969,87.7131) .. controls (252.1761,87.7131) and (251.5918,88.2974) ..(251.5918,89.0182) .. controls (251.5918,89.7390) and (252.1761,90.3234) ..(252.8969,90.3234) .. controls (253.6177,90.3234) and (254.2020,89.7390) ..(254.2020,89.0182);
\path[fill=chfourfig5c0,fill opacity=0.5,draw=chfourfig5c0,line width=0.3985pt,draw opacity=0.5] (230.5405,96.2413) .. controls (230.5405,95.5204) and (229.9561,94.9361) ..(229.2354,94.9361) .. controls (228.5145,94.9361) and (227.9302,95.5204) ..(227.9302,96.2413) .. controls (227.9302,96.9621) and (228.5145,97.5464) ..(229.2354,97.5464) .. controls (229.9561,97.5464) and (230.5405,96.9621) ..(230.5405,96.2413);
\path[fill=chfourfig5c0,fill opacity=0.5,draw=chfourfig5c0,line width=0.3985pt,draw opacity=0.5] (235.2230,94.9959) .. controls (235.2230,94.2751) and (234.6387,93.6908) ..(233.9178,93.6908) .. controls (233.1971,93.6908) and (232.6127,94.2751) ..(232.6127,94.9959) .. controls (232.6127,95.7167) and (233.1971,96.3010) ..(233.9178,96.3010) .. controls (234.6387,96.3010) and (235.2230,95.7167) ..(235.2230,94.9959);
\path[fill=chfourfig5c0,fill opacity=0.5,draw=chfourfig5c0,line width=0.3985pt,draw opacity=0.5] (279.4676,76.8936) .. controls (279.4676,76.1728) and (278.8833,75.5884) ..(278.1625,75.5884) .. controls (277.4417,75.5884) and (276.8574,76.1728) ..(276.8574,76.8936) .. controls (276.8574,77.6144) and (277.4417,78.1987) ..(278.1625,78.1987) .. controls (278.8833,78.1987) and (279.4676,77.6144) ..(279.4676,76.8936);
\path[fill=chfourfig5c0,fill opacity=0.5,draw=chfourfig5c0,line width=0.3985pt,draw opacity=0.5] (241.1110,108.1966) .. controls (241.1110,107.4758) and (240.5267,106.8914) ..(239.8058,106.8914) .. controls (239.0851,106.8914) and (238.5007,107.4758) ..(238.5007,108.1966) .. controls (238.5007,108.9174) and (239.0851,109.5017) ..(239.8058,109.5017) .. controls (240.5267,109.5017) and (241.1110,108.9174) ..(241.1110,108.1966);
\path[fill=chfourfig5c0,fill opacity=0.5,draw=chfourfig5c0,line width=0.3985pt,draw opacity=0.5] (236.6277,99.2301) .. controls (236.6277,98.5093) and (236.0434,97.9250) ..(235.3226,97.9250) .. controls (234.6018,97.9250) and (234.0175,98.5093) ..(234.0175,99.2301) .. controls (234.0175,99.9509) and (234.6018,100.5352) ..(235.3226,100.5352) .. controls (236.0434,100.5352) and (236.6277,99.9509) ..(236.6277,99.2301);
\path[fill=chfourfig5c0,fill opacity=0.5,draw=chfourfig5c0,line width=0.3985pt,draw opacity=0.5] (320.7932,62.8759) .. controls (320.7932,62.1552) and (320.2089,61.5708) ..(319.4881,61.5708) .. controls (318.7673,61.5708) and (318.1830,62.1552) ..(318.1830,62.8759) .. controls (318.1830,63.5968) and (318.7673,64.1811) ..(319.4881,64.1811) .. controls (320.2089,64.1811) and (320.7932,63.5968) ..(320.7932,62.8759);
\path[fill=chfourfig5c0,fill opacity=0.5,draw=chfourfig5c0,line width=0.3985pt,draw opacity=0.5] (282.2672,71.9819) .. controls (282.2672,71.2611) and (281.6829,70.6768) ..(280.9620,70.6768) .. controls (280.2413,70.6768) and (279.6570,71.2611) ..(279.6570,71.9819) .. controls (279.6570,72.7027) and (280.2413,73.2870) ..(280.9620,73.2870) .. controls (281.6829,73.2870) and (282.2672,72.7027) ..(282.2672,71.9819);
\path[fill=chfourfig5c0,fill opacity=0.5,draw=chfourfig5c0,line width=0.3985pt,draw opacity=0.5] (307.9911,52.6541) .. controls (307.9911,51.9333) and (307.4067,51.3490) ..(306.6859,51.3490) .. controls (305.9651,51.3490) and (305.3808,51.9333) ..(305.3808,52.6541) .. controls (305.3808,53.3749) and (305.9651,53.9593) ..(306.6859,53.9593) .. controls (307.4067,53.9593) and (307.9911,53.3749) ..(307.9911,52.6541);
\path[fill=chfourfig5c0,fill opacity=0.5,draw=chfourfig5c0,line width=0.3985pt,draw opacity=0.5] (260.5882,80.8587) .. controls (260.5882,80.1380) and (260.0038,79.5536) ..(259.2831,79.5536) .. controls (258.5623,79.5536) and (257.9780,80.1380) ..(257.9780,80.8587) .. controls (257.9780,81.5796) and (258.5623,82.1639) ..(259.2831,82.1639) .. controls (260.0038,82.1639) and (260.5882,81.5796) ..(260.5882,80.8587);
\path[fill=chfourfig5c0,fill opacity=0.5,draw=chfourfig5c0,line width=0.3985pt,draw opacity=0.5] (273.4999,88.4305) .. controls (273.4999,87.7096) and (272.9156,87.1253) ..(272.1948,87.1253) .. controls (271.4740,87.1253) and (270.8897,87.7096) ..(270.8897,88.4305) .. controls (270.8897,89.1512) and (271.4740,89.7356) ..(272.1948,89.7356) .. controls (272.9156,89.7356) and (273.4999,89.1512) ..(273.4999,88.4305);
\path[fill=chfourfig5c0,fill opacity=0.5,draw=chfourfig5c0,line width=0.3985pt,draw opacity=0.5] (305.2114,60.7838) .. controls (305.2114,60.0630) and (304.6271,59.4786) ..(303.9063,59.4786) .. controls (303.1855,59.4786) and (302.6012,60.0630) ..(302.6012,60.7838) .. controls (302.6012,61.5046) and (303.1855,62.0889) ..(303.9063,62.0889) .. controls (304.6271,62.0889) and (305.2114,61.5046) ..(305.2114,60.7838);
\path[fill=chfourfig5c0,fill opacity=0.5,draw=chfourfig5c0,line width=0.3985pt,draw opacity=0.5] (275.9109,78.8762) .. controls (275.9109,78.1553) and (275.3266,77.5710) ..(274.6058,77.5710) .. controls (273.8850,77.5710) and (273.3007,78.1553) ..(273.3007,78.8762) .. controls (273.3007,79.5970) and (273.8850,80.1813) ..(274.6058,80.1813) .. controls (275.3266,80.1813) and (275.9109,79.5970) ..(275.9109,78.8762);
\path[fill=chfourfig5c0,fill opacity=0.5,draw=chfourfig5c0,line width=0.3985pt,draw opacity=0.5] (266.3766,70.3381) .. controls (266.3766,69.6172) and (265.7922,69.0329) ..(265.0714,69.0329) .. controls (264.3506,69.0329) and (263.7663,69.6172) ..(263.7663,70.3381) .. controls (263.7663,71.0589) and (264.3506,71.6432) ..(265.0714,71.6432) .. controls (265.7922,71.6432) and (266.3766,71.0589) ..(266.3766,70.3381);
\path[fill=chfourfig5c0,fill opacity=0.5,draw=chfourfig5c0,line width=0.3985pt,draw opacity=0.5] (246.5307,101.9997) .. controls (246.5307,101.2789) and (245.9464,100.6946) ..(245.2256,100.6946) .. controls (244.5048,100.6946) and (243.9205,101.2789) ..(243.9205,101.9997) .. controls (243.9205,102.7206) and (244.5048,103.3049) ..(245.2256,103.3049) .. controls (245.9464,103.3049) and (246.5307,102.7206) ..(246.5307,101.9997);
\path[fill=chfourfig5c0,fill opacity=0.5,draw=chfourfig5c0,line width=0.3985pt,draw opacity=0.5] (333.3562,30.2080) .. controls (333.3562,29.4872) and (332.7719,28.9029) ..(332.0511,28.9029) .. controls (331.3303,28.9029) and (330.7460,29.4872) ..(330.7460,30.2080) .. controls (330.7460,30.9288) and (331.3303,31.5131) ..(332.0511,31.5131) .. controls (332.7719,31.5131) and (333.3562,30.9288) ..(333.3562,30.2080);
\path[fill=chfourfig5c0,fill opacity=0.5,draw=chfourfig5c0,line width=0.3985pt,draw opacity=0.5] (251.3826,102.1193) .. controls (251.3826,101.3985) and (250.7983,100.8142) ..(250.0775,100.8142) .. controls (249.3567,100.8142) and (248.7723,101.3985) ..(248.7723,102.1193) .. controls (248.7723,102.8401) and (249.3567,103.4244) ..(250.0775,103.4244) .. controls (250.7983,103.4244) and (251.3826,102.8401) ..(251.3826,102.1193);
\path[fill=chfourfig5c0,fill opacity=0.5,draw=chfourfig5c0,line width=0.3985pt,draw opacity=0.5] (309.4456,62.7564) .. controls (309.4456,62.0356) and (308.8613,61.4513) ..(308.1405,61.4513) .. controls (307.4197,61.4513) and (306.8354,62.0356) ..(306.8354,62.7564) .. controls (306.8354,63.4772) and (307.4197,64.0615) ..(308.1405,64.0615) .. controls (308.8613,64.0615) and (309.4456,63.4772) ..(309.4456,62.7564);
\path[fill=chfourfig5c0,fill opacity=0.5,draw=chfourfig5c0,line width=0.3985pt,draw opacity=0.5] (259.4923,73.1575) .. controls (259.4923,72.4367) and (258.9080,71.8524) ..(258.1872,71.8524) .. controls (257.4664,71.8524) and (256.8820,72.4367) ..(256.8820,73.1575) .. controls (256.8820,73.8783) and (257.4664,74.4626) ..(258.1872,74.4626) .. controls (258.9080,74.4626) and (259.4923,73.8783) ..(259.4923,73.1575);
\path[fill=chfourfig5c0,fill opacity=0.5,draw=chfourfig5c0,line width=0.3985pt,draw opacity=0.5] (300.2500,63.0354) .. controls (300.2500,62.3146) and (299.6656,61.7302) ..(298.9448,61.7302) .. controls (298.2241,61.7302) and (297.6397,62.3146) ..(297.6397,63.0354) .. controls (297.6397,63.7562) and (298.2241,64.3405) ..(298.9448,64.3405) .. controls (299.6656,64.3405) and (300.2500,63.7562) ..(300.2500,63.0354);
\path[fill=chfourfig5c0,fill opacity=0.5,draw=chfourfig5c0,line width=0.3985pt,draw opacity=0.5] (264.7925,81.9746) .. controls (264.7925,81.2538) and (264.2082,80.6694) ..(263.4873,80.6694) .. controls (262.7666,80.6694) and (262.1823,81.2538) ..(262.1823,81.9746) .. controls (262.1823,82.6954) and (262.7666,83.2797) ..(263.4873,83.2797) .. controls (264.2082,83.2797) and (264.7925,82.6954) ..(264.7925,81.9746);
\path[fill=chfourfig5c0,fill opacity=0.5,draw=chfourfig5c0,line width=0.3985pt,draw opacity=0.5] (294.8003,77.2821) .. controls (294.8003,76.5613) and (294.2160,75.9770) ..(293.4952,75.9770) .. controls (292.7744,75.9770) and (292.1901,76.5613) ..(292.1901,77.2821) .. controls (292.1901,78.0029) and (292.7744,78.5872) ..(293.4952,78.5872) .. controls (294.2160,78.5872) and (294.8003,78.0029) ..(294.8003,77.2821);
\path[fill=chfourfig5c0,fill opacity=0.5,draw=chfourfig5c0,line width=0.3985pt,draw opacity=0.5] (243.2928,90.8813) .. controls (243.2928,90.1605) and (242.7085,89.5762) ..(241.9877,89.5762) .. controls (241.2669,89.5762) and (240.6826,90.1605) ..(240.6826,90.8813) .. controls (240.6826,91.6021) and (241.2669,92.1864) ..(241.9877,92.1864) .. controls (242.7085,92.1864) and (243.2928,91.6021) ..(243.2928,90.8813);
\path[fill=chfourfig5c0,fill opacity=0.5,draw=chfourfig5c0,line width=0.3985pt,draw opacity=0.5] (276.2397,84.5649) .. controls (276.2397,83.8441) and (275.6554,83.2598) ..(274.9346,83.2598) .. controls (274.2138,83.2598) and (273.6295,83.8441) ..(273.6295,84.5649) .. controls (273.6295,85.2857) and (274.2138,85.8700) ..(274.9346,85.8700) .. controls (275.6554,85.8700) and (276.2397,85.2857) ..(276.2397,84.5649);
\path[fill=chfourfig5c0,fill opacity=0.5,draw=chfourfig5c0,line width=0.3985pt,draw opacity=0.5] (259.9207,76.5150) .. controls (259.9207,75.7942) and (259.3364,75.2099) ..(258.6155,75.2099) .. controls (257.8948,75.2099) and (257.3104,75.7942) ..(257.3104,76.5150) .. controls (257.3104,77.2358) and (257.8948,77.8201) ..(258.6155,77.8201) .. controls (259.3364,77.8201) and (259.9207,77.2358) ..(259.9207,76.5150);
\path[fill=chfourfig5c0,fill opacity=0.5,draw=chfourfig5c0,line width=0.3985pt,draw opacity=0.5] (260.8074,88.0220) .. controls (260.8074,87.3012) and (260.2231,86.7168) ..(259.5023,86.7168) .. controls (258.7814,86.7168) and (258.1971,87.3012) ..(258.1971,88.0220) .. controls (258.1971,88.7428) and (258.7814,89.3271) ..(259.5023,89.3271) .. controls (260.2231,89.3271) and (260.8074,88.7428) ..(260.8074,88.0220);
\path[fill=chfourfig5c0,fill opacity=0.5,draw=chfourfig5c0,line width=0.3985pt,draw opacity=0.5] (331.7423,47.0551) .. controls (331.7423,46.3343) and (331.1580,45.7500) ..(330.4371,45.7500) .. controls (329.7164,45.7500) and (329.1320,46.3343) ..(329.1320,47.0551) .. controls (329.1320,47.7759) and (329.7164,48.3602) ..(330.4371,48.3602) .. controls (331.1580,48.3602) and (331.7423,47.7759) ..(331.7423,47.0551);
\path[fill=chfourfig5c0,fill opacity=0.5,draw=chfourfig5c0,line width=0.3985pt,draw opacity=0.5] (283.6121,54.3976) .. controls (283.6121,53.6768) and (283.0278,53.0925) ..(282.3070,53.0925) .. controls (281.5862,53.0925) and (281.0019,53.6768) ..(281.0019,54.3976) .. controls (281.0019,55.1184) and (281.5862,55.7027) ..(282.3070,55.7027) .. controls (283.0278,55.7027) and (283.6121,55.1184) ..(283.6121,54.3976);
\path[fill=chfourfig5c0,fill opacity=0.5,draw=chfourfig5c0,line width=0.3985pt,draw opacity=0.5] (300.6086,54.2482) .. controls (300.6086,53.5274) and (300.0243,52.9431) ..(299.3035,52.9431) .. controls (298.5827,52.9431) and (297.9984,53.5274) ..(297.9984,54.2482) .. controls (297.9984,54.9690) and (298.5827,55.5533) ..(299.3035,55.5533) .. controls (300.0243,55.5533) and (300.6086,54.9690) ..(300.6086,54.2482);
\path[fill=chfourfig5c0,fill opacity=0.5,draw=chfourfig5c0,line width=0.3985pt,draw opacity=0.5] (239.0985,74.7914) .. controls (239.0985,74.0706) and (238.5142,73.4863) ..(237.7934,73.4863) .. controls (237.0726,73.4863) and (236.4883,74.0706) ..(236.4883,74.7914) .. controls (236.4883,75.5122) and (237.0726,76.0965) ..(237.7934,76.0965) .. controls (238.5142,76.0965) and (239.0985,75.5122) ..(239.0985,74.7914);
\path[fill=chfourfig5c0,fill opacity=0.5,draw=chfourfig5c0,line width=0.3985pt,draw opacity=0.5] (235.6613,97.9947) .. controls (235.6613,97.2739) and (235.0770,96.6896) ..(234.3562,96.6896) .. controls (233.6354,96.6896) and (233.0511,97.2739) ..(233.0511,97.9947) .. controls (233.0511,98.7155) and (233.6354,99.2998) ..(234.3562,99.2998) .. controls (235.0770,99.2998) and (235.6613,98.7155) ..(235.6613,97.9947);
\path[fill=chfourfig5c0,fill opacity=0.5,draw=chfourfig5c0,line width=0.3985pt,draw opacity=0.5] (232.3636,105.5564) .. controls (232.3636,104.8356) and (231.7794,104.2513) ..(231.0585,104.2513) .. controls (230.3377,104.2513) and (229.7534,104.8356) ..(229.7534,105.5564) .. controls (229.7534,106.2773) and (230.3377,106.8616) ..(231.0585,106.8616) .. controls (231.7794,106.8616) and (232.3636,106.2773) ..(232.3636,105.5564);
\path[fill=chfourfig5c0,fill opacity=0.5,draw=chfourfig5c0,line width=0.3985pt,draw opacity=0.5] (296.0756,58.7314) .. controls (296.0756,58.0106) and (295.4912,57.4263) ..(294.7704,57.4263) .. controls (294.0496,57.4263) and (293.4653,58.0106) ..(293.4653,58.7314) .. controls (293.4653,59.4522) and (294.0496,60.0366) ..(294.7704,60.0366) .. controls (295.4912,60.0366) and (296.0756,59.4522) ..(296.0756,58.7314);
\path[fill=chfourfig5c0,fill opacity=0.5,draw=chfourfig5c0,line width=0.3985pt,draw opacity=0.5] (254.0925,68.8835) .. controls (254.0925,68.1627) and (253.5081,67.5784) ..(252.7874,67.5784) .. controls (252.0665,67.5784) and (251.4822,68.1627) ..(251.4822,68.8835) .. controls (251.4822,69.6043) and (252.0665,70.1886) ..(252.7874,70.1886) .. controls (253.5081,70.1886) and (254.0925,69.6043) ..(254.0925,68.8835);
\path[fill=chfourfig5c0,fill opacity=0.5,draw=chfourfig5c0,line width=0.3985pt,draw opacity=0.5] (260.3989,102.0197) .. controls (260.3989,101.2989) and (259.8146,100.7145) ..(259.0938,100.7145) .. controls (258.3730,100.7145) and (257.7886,101.2989) ..(257.7886,102.0197) .. controls (257.7886,102.7405) and (258.3730,103.3248) ..(259.0938,103.3248) .. controls (259.8146,103.3248) and (260.3989,102.7405) ..(260.3989,102.0197);
\path[fill=none,draw=chfourfig5c0,line width=0.3985pt,draw opacity=1.0] (227.9601,120.7696) --(334.9304,120.7696);
\path[fill=none,draw=chfourfig5c0,line width=0.3985pt,draw opacity=1.0] (227.9601,120.7696) --(227.9601,125.7310);
\path[fill=none,draw=chfourfig5c0,line width=0.3985pt,draw opacity=1.0] (263.6169,120.7696) --(263.6169,125.7310);
\path[fill=none,draw=chfourfig5c0,line width=0.3985pt,draw opacity=1.0] (299.2736,120.7696) --(299.2736,125.7310);
\path[fill=none,draw=chfourfig5c0,line width=0.3985pt,draw opacity=1.0] (334.9304,120.7696) --(334.9304,125.7310);
\path[fill=none,draw=chfourfig5c0,line width=0.3985pt,draw opacity=1.0] (215.9749,119.7733) --(215.9749,39.6228);
\path[fill=none,draw=chfourfig5c0,line width=0.3985pt,draw opacity=1.0] (215.9749,119.7733) --(211.0135,119.7733);
\path[fill=none,draw=chfourfig5c0,line width=0.3985pt,draw opacity=1.0] (215.9749,79.6931) --(211.0135,79.6931);
\path[fill=none,draw=chfourfig5c0,line width=0.3985pt,draw opacity=1.0] (215.9749,39.6228) --(211.0135,39.6228);
\path[fill=none,draw=chfourfig5c0,line width=0.3985pt,draw opacity=1.0] (215.9749,120.7696)--(215.9749,16.7483)--(346.9056,16.7483)--(346.9056,120.7696)--cycle;
\path[fill=none,draw=chfourfig5c2,line width=1.1955pt,draw opacity=1.0] (223.4370,116.9240) --(225.4495,115.2701) --(226.0672,114.7621) --(227.9900,113.1780) --(228.3885,112.8492) --(228.5977,112.6798) --(228.6177,112.6599) --(229.8132,111.6736) --(229.9228,111.5939) --(230.1320,111.4146) --(230.2217,111.3448) --(230.4010,111.1954) --(231.4471,110.3286) --(231.8456,110.0098) --(233.0411,109.0235) --(233.0810,108.9936) --(233.3002,108.8143) --(233.4795,108.6648) --(234.0474,108.1966) --(234.1570,108.1069) --(234.8145,107.5590) --(235.3923,107.0907) --(235.6713,106.8616) --(235.8606,106.7022) --(236.7971,105.9350) --(236.9366,105.8155) --(237.2554,105.5564) --(238.9789,104.1417) --(239.0387,104.0919) --(239.1981,103.9624) --(239.9154,103.3746) --(240.4036,102.9761) --(240.8320,102.6174) --(241.1608,102.3484) --(241.2505,102.2787) --(241.4398,102.1193) --(241.5593,102.0197) --(241.5892,101.9997) --(242.0973,101.5813) --(242.1770,101.5116) --(242.2766,101.4319) --(242.5456,101.2127) --(242.7050,101.0832) --(242.9242,100.8939) --(243.2430,100.6348) --(243.6415,100.3061) --(243.9603,100.0470) --(244.4385,99.6485) --(244.5681,99.5489) --(244.8470,99.3098) --(244.9566,99.2301) --(245.7038,98.6124) --(245.9927,98.3733) --(246.1023,98.2836) --(246.4610,97.9947) --(246.5108,97.9449) --(246.6802,97.8154) --(246.8595,97.6659) --(247.6366,97.0283) --(248.5831,96.2413) --(248.5930,96.2313) --(248.9517,95.9324) --(248.9915,95.9025) --(249.6590,95.3645) --(249.8782,95.1752) --(249.9380,95.1254) --(250.1074,94.9959) --(250.5856,94.5974) --(250.7948,94.4280) --(250.9243,94.3185) --(251.1335,94.1391) --(251.4125,93.9199) --(251.4523,93.8801) --(251.5719,93.7805) --(251.6715,93.7107) --(251.6914,93.6808) --(251.7512,93.6410) --(251.9305,93.4915) --(252.1497,93.3122) --(252.4685,93.0532) --(252.4984,93.0233) --(252.5283,92.9934) --(252.6279,92.9137) --(252.9667,92.6447) --(253.0962,92.5351) --(253.3353,92.3358) --(253.4548,92.2362) --(253.9530,91.8278) --(254.1821,91.6385) --(254.2519,91.5787) --(254.3316,91.5189) --(254.7899,91.1403) --(254.8994,91.0507) --(255.0987,90.8813) --(255.3079,90.7119) --(255.5769,90.4928) --(255.6965,90.3931) --(256.0551,90.0942) --(256.1249,90.0444) --(256.3640,89.8452) --(256.5632,89.6858) --(256.7227,89.5463) --(256.8322,89.4566) --(257.0016,89.3171) --(257.0514,89.2872) --(257.3702,89.0182) --(257.6990,88.7493) --(258.0776,88.4305) --(258.2071,88.3308) --(258.3267,88.2312) --(258.4661,88.1116) --(258.4761,88.1116) --(258.5558,88.0419) --(258.5757,88.0220) --(258.6056,88.0021) --(259.0041,87.6733) --(259.3428,87.3943) --(259.6716,87.1253) --(259.7812,87.0357) --(259.9406,86.9061) --(260.0502,86.8165) --(260.0602,86.8065) --(260.4886,86.4578) --(261.1561,85.8999) --(261.2159,85.8600) --(261.3155,85.7803) --(261.5147,85.6110) --(261.5247,85.6010) --(262.6306,84.6944) --(262.7401,84.6048) --(262.7900,84.5649) --(262.9394,84.4354) --(263.1387,84.2760) --(263.1486,84.2660) --(263.1685,84.2461) --(263.2781,84.1564) --(263.6069,83.8874) --(263.6468,83.8575) --(263.9955,83.5686) --(264.1349,83.4590) --(264.3143,83.3096) --(264.6331,83.0506) --(264.7327,82.9609) --(265.0415,82.7118) --(265.1312,82.6421) --(265.2906,82.5026) --(265.4102,82.4030) --(265.8186,82.0642) --(265.9282,81.9746) --(265.9880,81.9347) --(266.1275,81.8251) --(266.2769,81.6956) --(266.3268,81.6558) --(266.5758,81.4465) --(266.8050,81.2572) --(266.8348,81.2373) --(266.8946,81.1875) --(267.0042,81.0979) --(267.2932,80.8587) --(267.5223,80.6694) --(267.5322,80.6595) --(267.7614,80.4702) --(267.7813,80.4602) --(267.8311,80.4204) --(267.8411,80.4104) --(267.8710,80.3805) --(267.9009,80.3606) --(268.0503,80.2410) --(268.2495,80.0717) --(268.6381,79.7529) --(268.6481,79.7429) --(268.6680,79.7330) --(268.7477,79.6632) --(269.0067,79.4440) --(269.0864,79.3842) --(269.3056,79.2049) --(269.3455,79.1750) --(269.5447,79.0057) --(269.7041,78.8762) --(269.7241,78.8562) --(269.8635,78.7466) --(269.8934,78.7168) --(269.9731,78.6470) --(270.2521,78.4179) --(270.3218,78.3681) --(270.4115,78.2884) --(270.4314,78.2784) --(270.5310,78.1987) --(271.2085,77.6408) --(271.5273,77.3718) --(271.6170,77.3020) --(271.6369,77.2821) --(271.7863,77.1626) --(272.0354,76.9533) --(272.0852,76.9135) --(272.1151,76.8936) --(272.3144,76.7342) --(272.3343,76.7142) --(272.5734,76.5150) --(272.7726,76.3556) --(273.2608,75.9571) --(273.2708,75.9471) --(273.4700,75.7777) --(273.6693,75.6183) --(273.7291,75.5685) --(274.1276,75.2397) --(274.5659,74.8811) --(274.6755,74.7914) --(274.9047,74.6021) --(274.9146,74.5922) --(275.1139,74.4328) --(275.1637,74.3829) --(275.3032,74.2734) --(275.4825,74.1239) --(275.5124,74.1040) --(275.5522,74.0741) --(275.6917,73.9545) --(275.9508,73.7453) --(276.0504,73.6557) --(276.1600,73.5660) --(276.1799,73.5560) --(276.2098,73.5262) --(276.3891,73.3767) --(276.5485,73.2472) --(276.6482,73.1675) --(276.6681,73.1575) --(276.7777,73.0579) --(276.7976,73.0380) --(277.4253,72.5299) --(277.7541,72.2509) --(277.8935,72.1413) --(277.9433,72.1015) --(277.9732,72.0716) --(278.0230,72.0317) --(278.0828,71.9819) --(278.1924,71.9022) --(278.8898,71.3244) --(279.1787,71.0853) --(279.1986,71.0654) --(279.2285,71.0454) --(279.2484,71.0255) --(279.3182,70.9657) --(279.3381,70.9557) --(279.5274,70.7963) --(279.6768,70.6768) --(279.7665,70.6071) --(279.8462,70.5373) --(279.9857,70.4178) --(280.0853,70.3381) --(280.0953,70.3381) --(280.4838,70.0193) --(280.5137,69.9894) --(280.5635,69.9495) --(280.6333,69.8897) --(280.8525,69.7104) --(281.0617,69.5410) --(281.1015,69.5012) --(281.1812,69.4414) --(281.2111,69.4215) --(281.3606,69.2920) --(281.5797,69.1126) --(281.6296,69.0728) --(281.6993,69.0130) --(281.8587,68.8835) --(282.1277,68.6544) --(282.2173,68.5846) --(282.4166,68.4252) --(282.4963,68.3555) --(282.5760,68.2857) --(282.9147,68.0167) --(282.9646,67.9669) --(283.2435,67.7378) --(283.6520,67.3990) --(283.8712,67.2297) --(284.1800,66.9806) --(284.3295,66.8511) --(284.3394,66.8411) --(284.4391,66.7515) --(284.4988,66.7116) --(285.5250,65.8747) --(285.6744,65.7452) --(285.9733,65.4962) --(286.1626,65.3467) --(286.2025,65.3168) --(286.3320,65.2072) --(286.3519,65.1873) --(286.3619,65.1774) --(286.3718,65.1674) --(286.3818,65.1674) --(286.4515,65.1076) --(287.5674,64.1811) --(287.6371,64.1313) --(288.2747,63.6032) --(288.4740,63.4339) --(288.5437,63.3840) --(288.6433,63.3043) --(288.6732,63.2844) --(288.6832,63.2745) --(288.9024,63.0852) --(288.9721,63.0354) --(289.1614,62.8759) --(289.2212,62.8261) --(289.3109,62.7564) --(289.3905,62.6867) --(289.4304,62.6568) --(289.5599,62.5571) --(289.5699,62.5372) --(289.7691,62.3778) --(290.0381,62.1586) --(290.1079,62.0989) --(290.1876,62.0391) --(290.2075,62.0092) --(290.2175,62.0092) --(290.4865,61.7900) --(290.5363,61.7402) --(290.6857,61.6206) --(290.7654,61.5609) --(290.8551,61.4911) --(290.9149,61.4314) --(291.0743,61.3018) --(291.1141,61.2719) --(291.2337,61.1723) --(291.6122,60.8535) --(291.6820,60.8037) --(291.7119,60.7838) --(292.1403,60.4251) --(292.2399,60.3454) --(292.3395,60.2657) --(292.4890,60.1362) --(292.8676,59.8273) --(292.9971,59.7178) --(293.0170,59.7078) --(293.2262,59.5285) --(293.2461,59.5185) --(293.2561,59.5085) --(293.3358,59.4488) --(293.8240,59.0403) --(293.9635,58.9307) --(293.9734,58.9207) --(294.1328,58.7912) --(294.2026,58.7314) --(294.4915,58.5023) --(294.6011,58.4027) --(294.6708,58.3429) --(294.8701,58.1835) --(295.2188,57.8946) --(295.5376,57.6355) --(295.6372,57.5558) --(295.8863,57.3466) --(295.9560,57.2868) --(296.0058,57.2470) --(297.1416,56.3105) --(297.2811,56.2009) --(297.5002,56.0216) --(297.5999,55.9419) --(298.0581,55.5633) --(298.1080,55.5234) --(298.1379,55.4935) --(298.1777,55.4637) --(298.7356,55.0054) --(298.7655,54.9755) --(298.7855,54.9655) --(298.8552,54.9057) --(298.8951,54.8659) --(298.9050,54.8659) --(298.9947,54.7862) --(299.2736,54.5570) --(299.3832,54.4674) --(299.4629,54.3976) --(299.5526,54.3279) --(299.6522,54.2482) --(299.8016,54.1187) --(300.2201,53.7799) --(300.3397,53.6803) --(300.5588,53.5010) --(300.6186,53.4512) --(300.8378,53.2818) --(300.8577,53.2619) --(301.0769,53.0826) --(301.4256,52.7936) --(301.4455,52.7737) --(301.5949,52.6541) --(301.6747,52.5944) --(301.8839,52.4150) --(301.9038,52.3951) --(302.0632,52.2756) --(302.1728,52.1759) --(302.2027,52.1560) --(302.2824,52.0962) --(302.8503,51.6280) --(302.9001,51.5782) --(303.2488,51.2892) --(303.3484,51.2095) --(303.7369,50.8907) --(303.8665,50.7911) --(303.8864,50.7712) --(303.8864,50.7612) --(304.0358,50.6417) --(304.0856,50.6118) --(304.7332,50.0738) --(304.7432,50.0638) --(304.8727,49.9642) --(305.3409,49.5757) --(305.4107,49.5159) --(305.4206,49.5059) --(306.3173,48.7786) --(306.5265,48.5993) --(306.6560,48.4897) --(306.7756,48.3901) --(307.2040,48.0414) --(307.2438,48.0115) --(307.5726,47.7425) --(308.1505,47.2643) --(308.2800,47.1547) --(308.3198,47.1248) --(308.3995,47.0551) --(308.4394,47.0252) --(308.6087,46.8857) --(308.6884,46.8160) --(309.3559,46.2680) --(309.4157,46.2182) --(309.5452,46.1186) --(310.1231,45.6404) --(310.2327,45.5507) --(310.3223,45.4810) --(310.3622,45.4411) --(310.4917,45.3315) --(310.9400,44.9729) --(310.9998,44.9131) --(311.0596,44.8633) --(311.2788,44.6839) --(311.6374,44.3950) --(311.9463,44.1459) --(312.0558,44.0563) --(312.0658,44.0463) --(312.1555,43.9666) --(312.2850,43.8670) --(312.4743,43.7076) --(312.6536,43.5581) --(312.8628,43.3888) --(312.9525,43.3091) --(313.2813,43.0401) --(313.3410,43.0002) --(313.5602,42.8109) --(313.5901,42.7910) --(313.9288,42.5121) --(314.4469,42.0836) --(314.6263,41.9342) --(315.1045,41.5457) --(315.3037,41.3763) --(315.5627,41.1671) --(315.9114,40.8781) --(316.3897,40.4896) --(316.3897,40.4796) --(316.6985,40.2405) --(316.7981,40.1509) --(317.6051,39.4834) --(317.8542,39.2841) --(317.9140,39.2343) --(317.9737,39.1845) --(318.3324,38.8856) --(318.4021,38.8358) --(318.4021,38.8258) --(318.9700,38.3675) --(320.4644,37.1421) --(320.5541,37.0624) --(320.9725,36.7137) --(320.9925,36.7038) --(321.1319,36.5842) --(321.7696,36.0562) --(321.9489,35.9167) --(322.6961,35.2990) --(323.3935,34.7311) --(323.8916,34.3127) --(324.0211,34.2130) --(324.4196,33.8843) --(324.7683,33.5954) --(324.9477,33.4559) --(324.9875,33.4160) --(325.6650,32.8581) --(325.6949,32.8382) --(325.8643,32.6987) --(326.6115,32.0810) --(326.9303,31.8220) --(327.4882,31.3637) --(328.0660,30.8855) --(328.8929,30.2080) --(329.6700,29.5704) --(331.8320,27.7871) --(333.6252,26.3126) --(337.8096,22.8754) --(339.4036,21.5703) --(340.5793,20.5940);
\path[fill=chfourfig5c0,fill opacity=0.5,draw=chfourfig5c0,line width=0.3985pt,draw opacity=0.5] (121.7371,192.3720) .. controls (121.7371,191.6512) and (121.1528,191.0669) ..(120.4319,191.0669) .. controls (119.7112,191.0669) and (119.1268,191.6512) ..(119.1268,192.3720) .. controls (119.1268,193.0928) and (119.7112,193.6771) ..(120.4319,193.6771) .. controls (121.1528,193.6771) and (121.7371,193.0928) ..(121.7371,192.3720);
\path[fill=chfourfig5c0,fill opacity=0.5,draw=chfourfig5c0,line width=0.3985pt,draw opacity=0.5] (104.7605,232.8607) .. controls (104.7605,232.1399) and (104.1762,231.5556) ..(103.4554,231.5556) .. controls (102.7346,231.5556) and (102.1503,232.1399) ..(102.1503,232.8607) .. controls (102.1503,233.5815) and (102.7346,234.1658) ..(103.4554,234.1658) .. controls (104.1762,234.1658) and (104.7605,233.5815) ..(104.7605,232.8607);
\path[fill=chfourfig5c0,fill opacity=0.5,draw=chfourfig5c0,line width=0.3985pt,draw opacity=0.5] (94.6882,192.0333) .. controls (94.6882,191.3125) and (94.1038,190.7282) ..(93.3830,190.7282) .. controls (92.6622,190.7282) and (92.0779,191.3125) ..(92.0779,192.0333) .. controls (92.0779,192.7541) and (92.6622,193.3384) ..(93.3830,193.3384) .. controls (94.1038,193.3384) and (94.6882,192.7541) ..(94.6882,192.0333);
\path[fill=chfourfig5c0,fill opacity=0.5,draw=chfourfig5c0,line width=0.3985pt,draw opacity=0.5] (54.6179,254.9980) .. controls (54.6179,254.2772) and (54.0336,253.6929) ..(53.3128,253.6929) .. controls (52.5920,253.6929) and (52.0077,254.2772) ..(52.0077,254.9980) .. controls (52.0077,255.7188) and (52.5920,256.3031) ..(53.3128,256.3031) .. controls (54.0336,256.3031) and (54.6179,255.7188) ..(54.6179,254.9980);
\path[fill=chfourfig5c0,fill opacity=0.5,draw=chfourfig5c0,line width=0.3985pt,draw opacity=0.5] (124.5565,188.1378) .. controls (124.5565,187.4170) and (123.9722,186.8327) ..(123.2514,186.8327) .. controls (122.5306,186.8327) and (121.9463,187.4170) ..(121.9463,188.1378) .. controls (121.9463,188.8586) and (122.5306,189.4430) ..(123.2514,189.4430) .. controls (123.9722,189.4430) and (124.5565,188.8586) ..(124.5565,188.1378);
\path[fill=chfourfig5c0,fill opacity=0.5,draw=chfourfig5c0,line width=0.3985pt,draw opacity=0.5] (144.8407,209.5080) .. controls (144.8407,208.7872) and (144.2564,208.2029) ..(143.5356,208.2029) .. controls (142.8148,208.2029) and (142.2305,208.7872) ..(142.2305,209.5080) .. controls (142.2305,210.2288) and (142.8148,210.8131) ..(143.5356,210.8131) .. controls (144.2564,210.8131) and (144.8407,210.2288) ..(144.8407,209.5080);
\path[fill=chfourfig5c0,fill opacity=0.5,draw=chfourfig5c0,line width=0.3985pt,draw opacity=0.5] (52.4261,251.8497) .. controls (52.4261,251.1289) and (51.8418,250.5446) ..(51.1210,250.5446) .. controls (50.4002,250.5446) and (49.8159,251.1289) ..(49.8159,251.8497) .. controls (49.8159,252.5706) and (50.4002,253.1549) ..(51.1210,253.1549) .. controls (51.8418,253.1549) and (52.4261,252.5706) ..(52.4261,251.8497);
\path[fill=chfourfig5c0,fill opacity=0.5,draw=chfourfig5c0,line width=0.3985pt,draw opacity=0.5] (161.8472,243.7500) .. controls (161.8472,243.0292) and (161.2629,242.4449) ..(160.5421,242.4449) .. controls (159.8213,242.4449) and (159.2369,243.0292) ..(159.2369,243.7500) .. controls (159.2369,244.4708) and (159.8213,245.0551) ..(160.5421,245.0551) .. controls (161.2629,245.0551) and (161.8472,244.4708) ..(161.8472,243.7500);
\path[fill=chfourfig5c0,fill opacity=0.5,draw=chfourfig5c0,line width=0.3985pt,draw opacity=0.5] (109.3135,218.9128) .. controls (109.3135,218.1920) and (108.7292,217.6077) ..(108.0084,217.6077) .. controls (107.2876,217.6077) and (106.7033,218.1920) ..(106.7033,218.9128) .. controls (106.7033,219.6336) and (107.2876,220.2180) ..(108.0084,220.2180) .. controls (108.7292,220.2180) and (109.3135,219.6336) ..(109.3135,218.9128);
\path[fill=chfourfig5c0,fill opacity=0.5,draw=chfourfig5c0,line width=0.3985pt,draw opacity=0.5] (162.9232,265.3692) .. controls (162.9232,264.6484) and (162.3389,264.0641) ..(161.6180,264.0641) .. controls (160.8972,264.0641) and (160.3129,264.6484) ..(160.3129,265.3692) .. controls (160.3129,266.0900) and (160.8972,266.6743) ..(161.6180,266.6743) .. controls (162.3389,266.6743) and (162.9232,266.0900) ..(162.9232,265.3692);
\path[fill=chfourfig5c0,fill opacity=0.5,draw=chfourfig5c0,line width=0.3985pt,draw opacity=0.5] (93.6321,231.7548) .. controls (93.6321,231.0340) and (93.0478,230.4497) ..(92.3270,230.4497) .. controls (91.6062,230.4497) and (91.0219,231.0340) ..(91.0219,231.7548) .. controls (91.0219,232.4756) and (91.6062,233.0599) ..(92.3270,233.0599) .. controls (93.0478,233.0599) and (93.6321,232.4756) ..(93.6321,231.7548);
\path[fill=chfourfig5c0,fill opacity=0.5,draw=chfourfig5c0,line width=0.3985pt,draw opacity=0.5] (77.7514,199.8341) .. controls (77.7514,199.1133) and (77.1671,198.5290) ..(76.4463,198.5290) .. controls (75.7255,198.5290) and (75.1412,199.1133) ..(75.1412,199.8341) .. controls (75.1412,200.5549) and (75.7255,201.1393) ..(76.4463,201.1393) .. controls (77.1671,201.1393) and (77.7514,200.5549) ..(77.7514,199.8341);
\path[fill=chfourfig5c0,fill opacity=0.5,draw=chfourfig5c0,line width=0.3985pt,draw opacity=0.5] (91.1215,255.2769) .. controls (91.1215,254.5561) and (90.5372,253.9718) ..(89.8164,253.9718) .. controls (89.0956,253.9718) and (88.5113,254.5561) ..(88.5113,255.2769) .. controls (88.5113,255.9977) and (89.0956,256.5820) ..(89.8164,256.5820) .. controls (90.5372,256.5820) and (91.1215,255.9977) ..(91.1215,255.2769);
\path[fill=chfourfig5c0,fill opacity=0.5,draw=chfourfig5c0,line width=0.3985pt,draw opacity=0.5] (77.9507,235.9093) .. controls (77.9507,235.1885) and (77.3664,234.6042) ..(76.6456,234.6042) .. controls (75.9248,234.6042) and (75.3405,235.1885) ..(75.3405,235.9093) .. controls (75.3405,236.6301) and (75.9248,237.2144) ..(76.6456,237.2144) .. controls (77.3664,237.2144) and (77.9507,236.6301) ..(77.9507,235.9093);
\path[fill=chfourfig5c0,fill opacity=0.5,draw=chfourfig5c0,line width=0.3985pt,draw opacity=0.5] (90.6333,191.1466) .. controls (90.6333,190.4258) and (90.0490,189.8415) ..(89.3282,189.8415) .. controls (88.6074,189.8415) and (88.0231,190.4258) ..(88.0231,191.1466) .. controls (88.0231,191.8674) and (88.6074,192.4517) ..(89.3282,192.4517) .. controls (90.0490,192.4517) and (90.6333,191.8674) ..(90.6333,191.1466);
\path[fill=chfourfig5c0,fill opacity=0.5,draw=chfourfig5c0,line width=0.3985pt,draw opacity=0.5] (84.9944,248.1237) .. controls (84.9944,247.4029) and (84.4101,246.8185) ..(83.6893,246.8185) .. controls (82.9685,246.8185) and (82.3842,247.4029) ..(82.3842,248.1237) .. controls (82.3842,248.8445) and (82.9685,249.4288) ..(83.6893,249.4288) .. controls (84.4101,249.4288) and (84.9944,248.8445) ..(84.9944,248.1237);
\path[fill=chfourfig5c0,fill opacity=0.5,draw=chfourfig5c0,line width=0.3985pt,draw opacity=0.5] (93.8812,230.3900) .. controls (93.8812,229.6691) and (93.2969,229.0848) ..(92.5760,229.0848) .. controls (91.8553,229.0848) and (91.2709,229.6691) ..(91.2709,230.3900) .. controls (91.2709,231.1107) and (91.8553,231.6951) ..(92.5760,231.6951) .. controls (93.2969,231.6951) and (93.8812,231.1107) ..(93.8812,230.3900);
\path[fill=chfourfig5c0,fill opacity=0.5,draw=chfourfig5c0,line width=0.3985pt,draw opacity=0.5] (134.4894,229.3140) .. controls (134.4894,228.5931) and (133.9051,228.0089) ..(133.1843,228.0089) .. controls (132.4635,228.0089) and (131.8792,228.5931) ..(131.8792,229.3140) .. controls (131.8792,230.0348) and (132.4635,230.6191) ..(133.1843,230.6191) .. controls (133.9051,230.6191) and (134.4894,230.0348) ..(134.4894,229.3140);
\path[fill=chfourfig5c0,fill opacity=0.5,draw=chfourfig5c0,line width=0.3985pt,draw opacity=0.5] (136.2728,239.8047) .. controls (136.2728,239.0840) and (135.6884,238.4996) ..(134.9676,238.4996) .. controls (134.2468,238.4996) and (133.6625,239.0840) ..(133.6625,239.8047) .. controls (133.6625,240.5256) and (134.2468,241.1099) ..(134.9676,241.1099) .. controls (135.6884,241.1099) and (136.2728,240.5256) ..(136.2728,239.8047);
\path[fill=chfourfig5c0,fill opacity=0.5,draw=chfourfig5c0,line width=0.3985pt,draw opacity=0.5] (80.5410,197.2239) .. controls (80.5410,196.5031) and (79.9567,195.9188) ..(79.2359,195.9188) .. controls (78.5151,195.9188) and (77.9308,196.5031) ..(77.9308,197.2239) .. controls (77.9308,197.9447) and (78.5151,198.5290) ..(79.2359,198.5290) .. controls (79.9567,198.5290) and (80.5410,197.9447) ..(80.5410,197.2239);
\path[fill=chfourfig5c0,fill opacity=0.5,draw=chfourfig5c0,line width=0.3985pt,draw opacity=0.5] (60.5757,212.7060) .. controls (60.5757,211.9852) and (59.9913,211.4009) ..(59.2705,211.4009) .. controls (58.5497,211.4009) and (57.9654,211.9852) ..(57.9654,212.7060) .. controls (57.9654,213.4268) and (58.5497,214.0111) ..(59.2705,214.0111) .. controls (59.9913,214.0111) and (60.5757,213.4268) ..(60.5757,212.7060);
\path[fill=chfourfig5c0,fill opacity=0.5,draw=chfourfig5c0,line width=0.3985pt,draw opacity=0.5] (66.3441,263.1276) .. controls (66.3441,262.4068) and (65.7598,261.8225) ..(65.0390,261.8225) .. controls (64.3182,261.8225) and (63.7338,262.4068) ..(63.7338,263.1276) .. controls (63.7338,263.8484) and (64.3182,264.4327) ..(65.0390,264.4327) .. controls (65.7598,264.4327) and (66.3441,263.8484) ..(66.3441,263.1276);
\path[fill=chfourfig5c0,fill opacity=0.5,draw=chfourfig5c0,line width=0.3985pt,draw opacity=0.5] (76.3766,261.6930) .. controls (76.3766,260.9722) and (75.7923,260.3878) ..(75.0715,260.3878) .. controls (74.3507,260.3878) and (73.7664,260.9722) ..(73.7664,261.6930) .. controls (73.7664,262.4138) and (74.3507,262.9981) ..(75.0715,262.9981) .. controls (75.7923,262.9981) and (76.3766,262.4138) ..(76.3766,261.6930);
\path[fill=chfourfig5c0,fill opacity=0.5,draw=chfourfig5c0,line width=0.3985pt,draw opacity=0.5] (118.6387,256.2832) .. controls (118.6387,255.5624) and (118.0543,254.9781) ..(117.3336,254.9781) .. controls (116.6127,254.9781) and (116.0284,255.5624) ..(116.0284,256.2832) .. controls (116.0284,257.0040) and (116.6127,257.5883) ..(117.3336,257.5883) .. controls (118.0543,257.5883) and (118.6387,257.0040) ..(118.6387,256.2832);
\path[fill=chfourfig5c0,fill opacity=0.5,draw=chfourfig5c0,line width=0.3985pt,draw opacity=0.5] (100.4666,216.7210) .. controls (100.4666,216.0002) and (99.8822,215.4159) ..(99.1615,215.4159) .. controls (98.4406,215.4159) and (97.8563,216.0002) ..(97.8563,216.7210) .. controls (97.8563,217.4418) and (98.4406,218.0261) ..(99.1615,218.0261) .. controls (99.8822,218.0261) and (100.4666,217.4418) ..(100.4666,216.7210);
\path[fill=chfourfig5c0,fill opacity=0.5,draw=chfourfig5c0,line width=0.3985pt,draw opacity=0.5] (62.5084,258.9133) .. controls (62.5084,258.1926) and (61.9241,257.6082) ..(61.2033,257.6082) .. controls (60.4825,257.6082) and (59.8982,258.1926) ..(59.8982,258.9133) .. controls (59.8982,259.6342) and (60.4825,260.2185) ..(61.2033,260.2185) .. controls (61.9241,260.2185) and (62.5084,259.6342) ..(62.5084,258.9133);
\path[fill=chfourfig5c0,fill opacity=0.5,draw=chfourfig5c0,line width=0.3985pt,draw opacity=0.5] (45.2430,175.2161) .. controls (45.2430,174.4953) and (44.6586,173.9110) ..(43.9378,173.9110) .. controls (43.2170,173.9110) and (42.6327,174.4953) ..(42.6327,175.2161) .. controls (42.6327,175.9369) and (43.2170,176.5212) ..(43.9378,176.5212) .. controls (44.6586,176.5212) and (45.2430,175.9369) ..(45.2430,175.2161);
\path[fill=chfourfig5c0,fill opacity=0.5,draw=chfourfig5c0,line width=0.3985pt,draw opacity=0.5] (52.7349,203.9189) .. controls (52.7349,203.1981) and (52.1506,202.6137) ..(51.4298,202.6137) .. controls (50.7090,202.6137) and (50.1247,203.1981) ..(50.1247,203.9189) .. controls (50.1247,204.6396) and (50.7090,205.2240) ..(51.4298,205.2240) .. controls (52.1506,205.2240) and (52.7349,204.6396) ..(52.7349,203.9189);
\path[fill=chfourfig5c0,fill opacity=0.5,draw=chfourfig5c0,line width=0.3985pt,draw opacity=0.5] (120.0733,204.3174) .. controls (120.0733,203.5966) and (119.4890,203.0123) ..(118.7682,203.0123) .. controls (118.0474,203.0123) and (117.4631,203.5966) ..(117.4631,204.3174) .. controls (117.4631,205.0382) and (118.0474,205.6225) ..(118.7682,205.6225) .. controls (119.4890,205.6225) and (120.0733,205.0382) ..(120.0733,204.3174);
\path[fill=chfourfig5c0,fill opacity=0.5,draw=chfourfig5c0,line width=0.3985pt,draw opacity=0.5] (145.1097,244.2780) .. controls (145.1097,243.5573) and (144.5254,242.9729) ..(143.8046,242.9729) .. controls (143.0838,242.9729) and (142.4995,243.5573) ..(142.4995,244.2780) .. controls (142.4995,244.9988) and (143.0838,245.5832) ..(143.8046,245.5832) .. controls (144.5254,245.5832) and (145.1097,244.9988) ..(145.1097,244.2780);
\path[fill=chfourfig5c0,fill opacity=0.5,draw=chfourfig5c0,line width=0.3985pt,draw opacity=0.5] (59.4997,187.3807) .. controls (59.4997,186.6599) and (58.9153,186.0755) ..(58.1945,186.0755) .. controls (57.4737,186.0755) and (56.8894,186.6599) ..(56.8894,187.3807) .. controls (56.8894,188.1015) and (57.4737,188.6858) ..(58.1945,188.6858) .. controls (58.9153,188.6858) and (59.4997,188.1015) ..(59.4997,187.3807);
\path[fill=chfourfig5c0,fill opacity=0.5,draw=chfourfig5c0,line width=0.3985pt,draw opacity=0.5] (133.7721,222.0411) .. controls (133.7721,221.3203) and (133.1878,220.7360) ..(132.4670,220.7360) .. controls (131.7462,220.7360) and (131.1619,221.3203) ..(131.1619,222.0411) .. controls (131.1619,222.7619) and (131.7462,223.3463) ..(132.4670,223.3463) .. controls (133.1878,223.3463) and (133.7721,222.7619) ..(133.7721,222.0411);
\path[fill=chfourfig5c0,fill opacity=0.5,draw=chfourfig5c0,line width=0.3985pt,draw opacity=0.5] (60.5557,252.1287) .. controls (60.5557,251.4079) and (59.9714,250.8236) ..(59.2506,250.8236) .. controls (58.5298,250.8236) and (57.9455,251.4079) ..(57.9455,252.1287) .. controls (57.9455,252.8495) and (58.5298,253.4338) ..(59.2506,253.4338) .. controls (59.9714,253.4338) and (60.5557,252.8495) ..(60.5557,252.1287);
\path[fill=chfourfig5c0,fill opacity=0.5,draw=chfourfig5c0,line width=0.3985pt,draw opacity=0.5] (121.9762,196.6062) .. controls (121.9762,195.8854) and (121.3918,195.3011) ..(120.6711,195.3011) .. controls (119.9503,195.3011) and (119.3659,195.8854) ..(119.3659,196.6062) .. controls (119.3659,197.3270) and (119.9503,197.9113) ..(120.6711,197.9113) .. controls (121.3918,197.9113) and (121.9762,197.3270) ..(121.9762,196.6062);
\path[fill=chfourfig5c0,fill opacity=0.5,draw=chfourfig5c0,line width=0.3985pt,draw opacity=0.5] (152.8409,253.9220) .. controls (152.8409,253.2012) and (152.2565,252.6169) ..(151.5357,252.6169) .. controls (150.8149,252.6169) and (150.2306,253.2012) ..(150.2306,253.9220) .. controls (150.2306,254.6428) and (150.8149,255.2271) ..(151.5357,255.2271) .. controls (152.2565,255.2271) and (152.8409,254.6428) ..(152.8409,253.9220);
\path[fill=chfourfig5c0,fill opacity=0.5,draw=chfourfig5c0,line width=0.3985pt,draw opacity=0.5] (160.4922,203.9288) .. controls (160.4922,203.2080) and (159.9079,202.6237) ..(159.1871,202.6237) .. controls (158.4663,202.6237) and (157.8820,203.2080) ..(157.8820,203.9288) .. controls (157.8820,204.6496) and (158.4663,205.2339) ..(159.1871,205.2339) .. controls (159.9079,205.2339) and (160.4922,204.6496) ..(160.4922,203.9288);
\path[fill=chfourfig5c0,fill opacity=0.5,draw=chfourfig5c0,line width=0.3985pt,draw opacity=0.5] (73.4575,235.0127) .. controls (73.4575,234.2919) and (72.8732,233.7076) ..(72.1524,233.7076) .. controls (71.4316,233.7076) and (70.8473,234.2919) ..(70.8473,235.0127) .. controls (70.8473,235.7335) and (71.4316,236.3178) ..(72.1524,236.3178) .. controls (72.8732,236.3178) and (73.4575,235.7335) ..(73.4575,235.0127);
\path[fill=chfourfig5c0,fill opacity=0.5,draw=chfourfig5c0,line width=0.3985pt,draw opacity=0.5] (136.1134,188.7854) .. controls (136.1134,188.0646) and (135.5290,187.4803) ..(134.8082,187.4803) .. controls (134.0874,187.4803) and (133.5031,188.0646) ..(133.5031,188.7854) .. controls (133.5031,189.5062) and (134.0874,190.0905) ..(134.8082,190.0905) .. controls (135.5290,190.0905) and (136.1134,189.5062) ..(136.1134,188.7854);
\path[fill=chfourfig5c0,fill opacity=0.5,draw=chfourfig5c0,line width=0.3985pt,draw opacity=0.5] (160.9007,234.1459) .. controls (160.9007,233.4251) and (160.3164,232.8408) ..(159.5956,232.8408) .. controls (158.8748,232.8408) and (158.2905,233.4251) ..(158.2905,234.1459) .. controls (158.2905,234.8667) and (158.8748,235.4510) ..(159.5956,235.4510) .. controls (160.3164,235.4510) and (160.9007,234.8667) ..(160.9007,234.1459);
\path[fill=chfourfig5c0,fill opacity=0.5,draw=chfourfig5c0,line width=0.3985pt,draw opacity=0.5] (125.4034,239.6553) .. controls (125.4034,238.9345) and (124.8191,238.3502) ..(124.0982,238.3502) .. controls (123.3775,238.3502) and (122.7931,238.9345) ..(122.7931,239.6553) .. controls (122.7931,240.3761) and (123.3775,240.9604) ..(124.0982,240.9604) .. controls (124.8191,240.9604) and (125.4034,240.3761) ..(125.4034,239.6553);
\path[fill=chfourfig5c0,fill opacity=0.5,draw=chfourfig5c0,line width=0.3985pt,draw opacity=0.5] (162.7837,222.5891) .. controls (162.7837,221.8683) and (162.1994,221.2840) ..(161.4786,221.2840) .. controls (160.7578,221.2840) and (160.1734,221.8683) ..(160.1734,222.5891) .. controls (160.1734,223.3099) and (160.7578,223.8942) ..(161.4786,223.8942) .. controls (162.1994,223.8942) and (162.7837,223.3099) ..(162.7837,222.5891);
\path[fill=chfourfig5c0,fill opacity=0.5,draw=chfourfig5c0,line width=0.3985pt,draw opacity=0.5] (123.6998,208.5615) .. controls (123.6998,207.8407) and (123.1154,207.2564) ..(122.3946,207.2564) .. controls (121.6738,207.2564) and (121.0895,207.8407) ..(121.0895,208.5615) .. controls (121.0895,209.2823) and (121.6738,209.8666) ..(122.3946,209.8666) .. controls (123.1154,209.8666) and (123.6998,209.2823) ..(123.6998,208.5615);
\path[fill=chfourfig5c0,fill opacity=0.5,draw=chfourfig5c0,line width=0.3985pt,draw opacity=0.5] (61.2531,200.7806) .. controls (61.2531,200.0598) and (60.6688,199.4755) ..(59.9480,199.4755) .. controls (59.2272,199.4755) and (58.6429,200.0598) ..(58.6429,200.7806) .. controls (58.6429,201.5014) and (59.2272,202.0857) ..(59.9480,202.0857) .. controls (60.6688,202.0857) and (61.2531,201.5014) ..(61.2531,200.7806);
\path[fill=chfourfig5c0,fill opacity=0.5,draw=chfourfig5c0,line width=0.3985pt,draw opacity=0.5] (70.2794,179.6197) .. controls (70.2794,178.8989) and (69.6951,178.3145) ..(68.9743,178.3145) .. controls (68.2534,178.3145) and (67.6691,178.8989) ..(67.6691,179.6197) .. controls (67.6691,180.3405) and (68.2534,180.9248) ..(68.9743,180.9248) .. controls (69.6951,180.9248) and (70.2794,180.3405) ..(70.2794,179.6197);
\path[fill=chfourfig5c0,fill opacity=0.5,draw=chfourfig5c0,line width=0.3985pt,draw opacity=0.5] (51.3900,220.8356) .. controls (51.3900,220.1149) and (50.8056,219.5305) ..(50.0849,219.5305) .. controls (49.3641,219.5305) and (48.7797,220.1149) ..(48.7797,220.8356) .. controls (48.7797,221.5565) and (49.3641,222.1407) ..(50.0849,222.1407) .. controls (50.8056,222.1407) and (51.3900,221.5565) ..(51.3900,220.8356);
\path[fill=chfourfig5c0,fill opacity=0.5,draw=chfourfig5c0,line width=0.3985pt,draw opacity=0.5] (94.0306,208.9899) .. controls (94.0306,208.2691) and (93.4463,207.6848) ..(92.7255,207.6848) .. controls (92.0047,207.6848) and (91.4204,208.2691) ..(91.4204,208.9899) .. controls (91.4204,209.7107) and (92.0047,210.2950) ..(92.7255,210.2950) .. controls (93.4463,210.2950) and (94.0306,209.7107) ..(94.0306,208.9899);
\path[fill=chfourfig5c0,fill opacity=0.5,draw=chfourfig5c0,line width=0.3985pt,draw opacity=0.5] (109.2238,264.5822) .. controls (109.2238,263.8614) and (108.6395,263.2770) ..(107.9187,263.2770) .. controls (107.1979,263.2770) and (106.6136,263.8614) ..(106.6136,264.5822) .. controls (106.6136,265.3030) and (107.1979,265.8873) ..(107.9187,265.8873) .. controls (108.6395,265.8873) and (109.2238,265.3030) ..(109.2238,264.5822);
\path[fill=chfourfig5c0,fill opacity=0.5,draw=chfourfig5c0,line width=0.3985pt,draw opacity=0.5] (64.7401,179.5699) .. controls (64.7401,178.8490) and (64.1558,178.2647) ..(63.4350,178.2647) .. controls (62.7142,178.2647) and (62.1298,178.8490) ..(62.1298,179.5699) .. controls (62.1298,180.2906) and (62.7142,180.8750) ..(63.4350,180.8750) .. controls (64.1558,180.8750) and (64.7401,180.2906) ..(64.7401,179.5699);
\path[fill=chfourfig5c0,fill opacity=0.5,draw=chfourfig5c0,line width=0.3985pt,draw opacity=0.5] (123.8691,219.6301) .. controls (123.8691,218.9094) and (123.2848,218.3250) ..(122.5640,218.3250) .. controls (121.8432,218.3250) and (121.2589,218.9094) ..(121.2589,219.6301) .. controls (121.2589,220.3510) and (121.8432,220.9353) ..(122.5640,220.9353) .. controls (123.2848,220.9353) and (123.8691,220.3510) ..(123.8691,219.6301);
\path[fill=chfourfig5c0,fill opacity=0.5,draw=chfourfig5c0,line width=0.3985pt,draw opacity=0.5] (149.3838,206.8878) .. controls (149.3838,206.1670) and (148.7995,205.5826) ..(148.0786,205.5826) .. controls (147.3578,205.5826) and (146.7735,206.1670) ..(146.7735,206.8878) .. controls (146.7735,207.6086) and (147.3578,208.1929) ..(148.0786,208.1929) .. controls (148.7995,208.1929) and (149.3838,207.6086) ..(149.3838,206.8878);
\path[fill=chfourfig5c0,fill opacity=0.5,draw=chfourfig5c0,line width=0.3985pt,draw opacity=0.5] (126.4495,184.2424) .. controls (126.4495,183.5216) and (125.8651,182.9373) ..(125.1443,182.9373) .. controls (124.4236,182.9373) and (123.8392,183.5216) ..(123.8392,184.2424) .. controls (123.8392,184.9632) and (124.4236,185.5475) ..(125.1443,185.5475) .. controls (125.8651,185.5475) and (126.4495,184.9632) ..(126.4495,184.2424);
\path[fill=chfourfig5c0,fill opacity=0.5,draw=chfourfig5c0,line width=0.3985pt,draw opacity=0.5] (46.4883,199.1367) .. controls (46.4883,198.4159) and (45.9040,197.8316) ..(45.1832,197.8316) .. controls (44.4624,197.8316) and (43.8781,198.4159) ..(43.8781,199.1367) .. controls (43.8781,199.8575) and (44.4624,200.4419) ..(45.1832,200.4419) .. controls (45.9040,200.4419) and (46.4883,199.8575) ..(46.4883,199.1367);
\path[fill=chfourfig5c0,fill opacity=0.5,draw=chfourfig5c0,line width=0.3985pt,draw opacity=0.5] (88.6906,220.3275) .. controls (88.6906,219.6068) and (88.1063,219.0224) ..(87.3855,219.0224) .. controls (86.6647,219.0224) and (86.0803,219.6068) ..(86.0803,220.3275) .. controls (86.0803,221.0483) and (86.6647,221.6327) ..(87.3855,221.6327) .. controls (88.1063,221.6327) and (88.6906,221.0483) ..(88.6906,220.3275);
\path[fill=chfourfig5c0,fill opacity=0.5,draw=chfourfig5c0,line width=0.3985pt,draw opacity=0.5] (96.4615,190.2300) .. controls (96.4615,189.5092) and (95.8772,188.9249) ..(95.1564,188.9249) .. controls (94.4356,188.9249) and (93.8513,189.5092) ..(93.8513,190.2300) .. controls (93.8513,190.9508) and (94.4356,191.5352) ..(95.1564,191.5352) .. controls (95.8772,191.5352) and (96.4615,190.9508) ..(96.4615,190.2300);
\path[fill=chfourfig5c0,fill opacity=0.5,draw=chfourfig5c0,line width=0.3985pt,draw opacity=0.5] (133.8120,264.8213) .. controls (133.8120,264.1005) and (133.2276,263.5161) ..(132.5068,263.5161) .. controls (131.7860,263.5161) and (131.2017,264.1005) ..(131.2017,264.8213) .. controls (131.2017,265.5421) and (131.7860,266.1264) ..(132.5068,266.1264) .. controls (133.2276,266.1264) and (133.8120,265.5421) ..(133.8120,264.8213);
\path[fill=chfourfig5c0,fill opacity=0.5,draw=chfourfig5c0,line width=0.3985pt,draw opacity=0.5] (140.3575,217.8468) .. controls (140.3575,217.1260) and (139.7732,216.5417) ..(139.0524,216.5417) .. controls (138.3316,216.5417) and (137.7473,217.1260) ..(137.7473,217.8468) .. controls (137.7473,218.5676) and (138.3316,219.1519) ..(139.0524,219.1519) .. controls (139.7732,219.1519) and (140.3575,218.5676) ..(140.3575,217.8468);
\path[fill=chfourfig5c0,fill opacity=0.5,draw=chfourfig5c0,line width=0.3985pt,draw opacity=0.5] (64.3416,256.1835) .. controls (64.3416,255.4628) and (63.7572,254.8784) ..(63.0365,254.8784) .. controls (62.3157,254.8784) and (61.7313,255.4628) ..(61.7313,256.1835) .. controls (61.7313,256.9044) and (62.3157,257.4887) ..(63.0365,257.4887) .. controls (63.7572,257.4887) and (64.3416,256.9044) ..(64.3416,256.1835);
\path[fill=chfourfig5c0,fill opacity=0.5,draw=chfourfig5c0,line width=0.3985pt,draw opacity=0.5] (104.4318,243.2319) .. controls (104.4318,242.5112) and (103.8474,241.9268) ..(103.1266,241.9268) .. controls (102.4058,241.9268) and (101.8215,242.5112) ..(101.8215,243.2319) .. controls (101.8215,243.9528) and (102.4058,244.5370) ..(103.1266,244.5370) .. controls (103.8474,244.5370) and (104.4318,243.9528) ..(104.4318,243.2319);
\path[fill=chfourfig5c0,fill opacity=0.5,draw=chfourfig5c0,line width=0.3985pt,draw opacity=0.5] (149.1845,195.7992) .. controls (149.1845,195.0784) and (148.6002,194.4941) ..(147.8794,194.4941) .. controls (147.1586,194.4941) and (146.5743,195.0784) ..(146.5743,195.7992) .. controls (146.5743,196.5200) and (147.1586,197.1043) ..(147.8794,197.1043) .. controls (148.6002,197.1043) and (149.1845,196.5200) ..(149.1845,195.7992);
\path[fill=chfourfig5c0,fill opacity=0.5,draw=chfourfig5c0,line width=0.3985pt,draw opacity=0.5] (48.0524,196.9848) .. controls (48.0524,196.2640) and (47.4681,195.6797) ..(46.7473,195.6797) .. controls (46.0265,195.6797) and (45.4422,196.2640) ..(45.4422,196.9848) .. controls (45.4422,197.7056) and (46.0265,198.2899) ..(46.7473,198.2899) .. controls (47.4681,198.2899) and (48.0524,197.7056) ..(48.0524,196.9848);
\path[fill=chfourfig5c0,fill opacity=0.5,draw=chfourfig5c0,line width=0.3985pt,draw opacity=0.5] (99.8887,186.6135) .. controls (99.8887,185.8927) and (99.3044,185.3084) ..(98.5836,185.3084) .. controls (97.8628,185.3084) and (97.2785,185.8927) ..(97.2785,186.6135) .. controls (97.2785,187.3343) and (97.8628,187.9187) ..(98.5836,187.9187) .. controls (99.3044,187.9187) and (99.8887,187.3343) ..(99.8887,186.6135);
\path[fill=chfourfig5c0,fill opacity=0.5,draw=chfourfig5c0,line width=0.3985pt,draw opacity=0.5] (128.8904,242.1560) .. controls (128.8904,241.4352) and (128.3060,240.8508) ..(127.5852,240.8508) .. controls (126.8644,240.8508) and (126.2801,241.4352) ..(126.2801,242.1560) .. controls (126.2801,242.8768) and (126.8644,243.4611) ..(127.5852,243.4611) .. controls (128.3060,243.4611) and (128.8904,242.8768) ..(128.8904,242.1560);
\path[fill=chfourfig5c0,fill opacity=0.5,draw=chfourfig5c0,line width=0.3985pt,draw opacity=0.5] (94.3694,250.4151) .. controls (94.3694,249.6943) and (93.7850,249.1100) ..(93.0642,249.1100) .. controls (92.3434,249.1100) and (91.7591,249.6943) ..(91.7591,250.4151) .. controls (91.7591,251.1359) and (92.3434,251.7202) ..(93.0642,251.7202) .. controls (93.7850,251.7202) and (94.3694,251.1359) ..(94.3694,250.4151);
\path[fill=chfourfig5c0,fill opacity=0.5,draw=chfourfig5c0,line width=0.3985pt,draw opacity=0.5] (66.6629,210.3249) .. controls (66.6629,209.6041) and (66.0786,209.0198) ..(65.3578,209.0198) .. controls (64.6370,209.0198) and (64.0527,209.6041) ..(64.0527,210.3249) .. controls (64.0527,211.0457) and (64.6370,211.6301) ..(65.3578,211.6301) .. controls (66.0786,211.6301) and (66.6629,211.0457) ..(66.6629,210.3249);
\path[fill=chfourfig5c0,fill opacity=0.5,draw=chfourfig5c0,line width=0.3985pt,draw opacity=0.5] (48.8993,197.1243) .. controls (48.8993,196.4034) and (48.3150,195.8192) ..(47.5942,195.8192) .. controls (46.8734,195.8192) and (46.2890,196.4034) ..(46.2890,197.1243) .. controls (46.2890,197.8451) and (46.8734,198.4294) ..(47.5942,198.4294) .. controls (48.3150,198.4294) and (48.8993,197.8451) ..(48.8993,197.1243);
\path[fill=chfourfig5c0,fill opacity=0.5,draw=chfourfig5c0,line width=0.3985pt,draw opacity=0.5] (53.0338,249.8273) .. controls (53.0338,249.1065) and (52.4495,248.5222) ..(51.7287,248.5222) .. controls (51.0079,248.5222) and (50.4236,249.1065) ..(50.4236,249.8273) .. controls (50.4236,250.5481) and (51.0079,251.1324) ..(51.7287,251.1324) .. controls (52.4495,251.1324) and (53.0338,250.5481) ..(53.0338,249.8273);
\path[fill=chfourfig5c0,fill opacity=0.5,draw=chfourfig5c0,line width=0.3985pt,draw opacity=0.5] (73.8759,180.7554) .. controls (73.8759,180.0346) and (73.2916,179.4503) ..(72.5708,179.4503) .. controls (71.8500,179.4503) and (71.2657,180.0346) ..(71.2657,180.7554) .. controls (71.2657,181.4762) and (71.8500,182.0605) ..(72.5708,182.0605) .. controls (73.2916,182.0605) and (73.8759,181.4762) ..(73.8759,180.7554);
\path[fill=chfourfig5c0,fill opacity=0.5,draw=chfourfig5c0,line width=0.3985pt,draw opacity=0.5] (47.6639,192.5015) .. controls (47.6639,191.7807) and (47.0796,191.1964) ..(46.3588,191.1964) .. controls (45.6380,191.1964) and (45.0536,191.7807) ..(45.0536,192.5015) .. controls (45.0536,193.2224) and (45.6380,193.8066) ..(46.3588,193.8066) .. controls (47.0796,193.8066) and (47.6639,193.2224) ..(47.6639,192.5015);
\path[fill=chfourfig5c0,fill opacity=0.5,draw=chfourfig5c0,line width=0.3985pt,draw opacity=0.5] (151.9741,221.0249) .. controls (151.9741,220.3041) and (151.3898,219.7198) ..(150.6690,219.7198) .. controls (149.9482,219.7198) and (149.3638,220.3041) ..(149.3638,221.0249) .. controls (149.3638,221.7457) and (149.9482,222.3300) ..(150.6690,222.3300) .. controls (151.3898,222.3300) and (151.9741,221.7457) ..(151.9741,221.0249);
\path[fill=chfourfig5c0,fill opacity=0.5,draw=chfourfig5c0,line width=0.3985pt,draw opacity=0.5] (165.2545,232.0936) .. controls (165.2545,231.3727) and (164.6701,230.7885) ..(163.9493,230.7885) .. controls (163.2285,230.7885) and (162.6442,231.3727) ..(162.6442,232.0936) .. controls (162.6442,232.8144) and (163.2285,233.3987) ..(163.9493,233.3987) .. controls (164.6701,233.3987) and (165.2545,232.8144) ..(165.2545,232.0936);
\path[fill=chfourfig5c0,fill opacity=0.5,draw=chfourfig5c0,line width=0.3985pt,draw opacity=0.5] (60.7749,230.5195) .. controls (60.7749,229.7987) and (60.1906,229.2144) ..(59.4698,229.2144) .. controls (58.7490,229.2144) and (58.1647,229.7987) ..(58.1647,230.5195) .. controls (58.1647,231.2403) and (58.7490,231.8246) ..(59.4698,231.8246) .. controls (60.1906,231.8246) and (60.7749,231.2403) ..(60.7749,230.5195);
\path[fill=chfourfig5c0,fill opacity=0.5,draw=chfourfig5c0,line width=0.3985pt,draw opacity=0.5] (58.8222,260.6967) .. controls (58.8222,259.9759) and (58.2379,259.3915) ..(57.5171,259.3915) .. controls (56.7963,259.3915) and (56.2119,259.9759) ..(56.2119,260.6967) .. controls (56.2119,261.4175) and (56.7963,262.0018) ..(57.5171,262.0018) .. controls (58.2379,262.0018) and (58.8222,261.4175) ..(58.8222,260.6967);
\path[fill=chfourfig5c0,fill opacity=0.5,draw=chfourfig5c0,line width=0.3985pt,draw opacity=0.5] (126.7782,225.8469) .. controls (126.7782,225.1261) and (126.1939,224.5418) ..(125.4731,224.5418) .. controls (124.7523,224.5418) and (124.1680,225.1261) ..(124.1680,225.8469) .. controls (124.1680,226.5677) and (124.7523,227.1520) ..(125.4731,227.1520) .. controls (126.1939,227.1520) and (126.7782,226.5677) ..(126.7782,225.8469);
\path[fill=chfourfig5c0,fill opacity=0.5,draw=chfourfig5c0,line width=0.3985pt,draw opacity=0.5] (120.5515,200.6212) .. controls (120.5515,199.9004) and (119.9672,199.3160) ..(119.2464,199.3160) .. controls (118.5256,199.3160) and (117.9413,199.9004) ..(117.9413,200.6212) .. controls (117.9413,201.3420) and (118.5256,201.9263) ..(119.2464,201.9263) .. controls (119.9672,201.9263) and (120.5515,201.3420) ..(120.5515,200.6212);
\path[fill=chfourfig5c0,fill opacity=0.5,draw=chfourfig5c0,line width=0.3985pt,draw opacity=0.5] (123.8890,202.1255) .. controls (123.8890,201.4048) and (123.3047,200.8204) ..(122.5839,200.8204) .. controls (121.8631,200.8204) and (121.2788,201.4048) ..(121.2788,202.1255) .. controls (121.2788,202.8464) and (121.8631,203.4307) ..(122.5839,203.4307) .. controls (123.3047,203.4307) and (123.8890,202.8464) ..(123.8890,202.1255);
\path[fill=chfourfig5c0,fill opacity=0.5,draw=chfourfig5c0,line width=0.3985pt,draw opacity=0.5] (126.7882,195.8789) .. controls (126.7882,195.1581) and (126.2039,194.5738) ..(125.4831,194.5738) .. controls (124.7623,194.5738) and (124.1780,195.1581) ..(124.1780,195.8789) .. controls (124.1780,196.5997) and (124.7623,197.1840) ..(125.4831,197.1840) .. controls (126.2039,197.1840) and (126.7882,196.5997) ..(126.7882,195.8789);
\path[fill=chfourfig5c0,fill opacity=0.5,draw=chfourfig5c0,line width=0.3985pt,draw opacity=0.5] (155.7002,208.6113) .. controls (155.7002,207.8905) and (155.1158,207.3062) ..(154.3950,207.3062) .. controls (153.6742,207.3062) and (153.0899,207.8905) ..(153.0899,208.6113) .. controls (153.0899,209.3321) and (153.6742,209.9164) ..(154.3950,209.9164) .. controls (155.1158,209.9164) and (155.7002,209.3321) ..(155.7002,208.6113);
\path[fill=chfourfig5c0,fill opacity=0.5,draw=chfourfig5c0,line width=0.3985pt,draw opacity=0.5] (82.8026,209.0497) .. controls (82.8026,208.3289) and (82.2183,207.7446) ..(81.4975,207.7446) .. controls (80.7767,207.7446) and (80.1923,208.3289) ..(80.1923,209.0497) .. controls (80.1923,209.7705) and (80.7767,210.3548) ..(81.4975,210.3548) .. controls (82.2183,210.3548) and (82.8026,209.7705) ..(82.8026,209.0497);
\path[fill=chfourfig5c0,fill opacity=0.5,draw=chfourfig5c0,line width=0.3985pt,draw opacity=0.5] (111.9337,254.6792) .. controls (111.9337,253.9584) and (111.3494,253.3741) ..(110.6286,253.3741) .. controls (109.9078,253.3741) and (109.3235,253.9584) ..(109.3235,254.6792) .. controls (109.3235,255.4000) and (109.9078,255.9843) ..(110.6286,255.9843) .. controls (111.3494,255.9843) and (111.9337,255.4000) ..(111.9337,254.6792);
\path[fill=chfourfig5c0,fill opacity=0.5,draw=chfourfig5c0,line width=0.3985pt,draw opacity=0.5] (131.1120,235.9193) .. controls (131.1120,235.1985) and (130.5277,234.6142) ..(129.8069,234.6142) .. controls (129.0861,234.6142) and (128.5018,235.1985) ..(128.5018,235.9193) .. controls (128.5018,236.6401) and (129.0861,237.2244) ..(129.8069,237.2244) .. controls (130.5277,237.2244) and (131.1120,236.6401) ..(131.1120,235.9193);
\path[fill=chfourfig5c0,fill opacity=0.5,draw=chfourfig5c0,line width=0.3985pt,draw opacity=0.5] (131.3512,226.1159) .. controls (131.3512,225.3951) and (130.7668,224.8108) ..(130.0460,224.8108) .. controls (129.3252,224.8108) and (128.7409,225.3951) ..(128.7409,226.1159) .. controls (128.7409,226.8367) and (129.3252,227.4210) ..(130.0460,227.4210) .. controls (130.7668,227.4210) and (131.3512,226.8367) ..(131.3512,226.1159);
\path[fill=chfourfig5c0,fill opacity=0.5,draw=chfourfig5c0,line width=0.3985pt,draw opacity=0.5] (55.1061,214.3399) .. controls (55.1061,213.6191) and (54.5218,213.0348) ..(53.8010,213.0348) .. controls (53.0802,213.0348) and (52.4958,213.6191) ..(52.4958,214.3399) .. controls (52.4958,215.0607) and (53.0802,215.6450) ..(53.8010,215.6450) .. controls (54.5218,215.6450) and (55.1061,215.0607) ..(55.1061,214.3399);
\path[fill=chfourfig5c0,fill opacity=0.5,draw=chfourfig5c0,line width=0.3985pt,draw opacity=0.5] (111.6049,234.8034) .. controls (111.6049,234.0826) and (111.0206,233.4983) ..(110.2998,233.4983) .. controls (109.5790,233.4983) and (108.9947,234.0826) ..(108.9947,234.8034) .. controls (108.9947,235.5242) and (109.5790,236.1086) ..(110.2998,236.1086) .. controls (111.0206,236.1086) and (111.6049,235.5242) ..(111.6049,234.8034);
\path[fill=chfourfig5c0,fill opacity=0.5,draw=chfourfig5c0,line width=0.3985pt,draw opacity=0.5] (101.9311,191.7145) .. controls (101.9311,190.9937) and (101.3468,190.4093) ..(100.6260,190.4093) .. controls (99.9052,190.4093) and (99.3209,190.9937) ..(99.3209,191.7145) .. controls (99.3209,192.4353) and (99.9052,193.0196) ..(100.6260,193.0196) .. controls (101.3468,193.0196) and (101.9311,192.4353) ..(101.9311,191.7145);
\path[fill=chfourfig5c0,fill opacity=0.5,draw=chfourfig5c0,line width=0.3985pt,draw opacity=0.5] (153.1497,184.0332) .. controls (153.1497,183.3124) and (152.5654,182.7281) ..(151.8446,182.7281) .. controls (151.1238,182.7281) and (150.5394,183.3124) ..(150.5394,184.0332) .. controls (150.5394,184.7540) and (151.1238,185.3383) ..(151.8446,185.3383) .. controls (152.5654,185.3383) and (153.1497,184.7540) ..(153.1497,184.0332);
\path[fill=chfourfig5c0,fill opacity=0.5,draw=chfourfig5c0,line width=0.3985pt,draw opacity=0.5] (59.6391,260.9059) .. controls (59.6391,260.1851) and (59.0548,259.6008) ..(58.3340,259.6008) .. controls (57.6132,259.6008) and (57.0289,260.1851) ..(57.0289,260.9059) .. controls (57.0289,261.6267) and (57.6132,262.2110) ..(58.3340,262.2110) .. controls (59.0548,262.2110) and (59.6391,261.6267) ..(59.6391,260.9059);
\path[fill=chfourfig5c0,fill opacity=0.5,draw=chfourfig5c0,line width=0.3985pt,draw opacity=0.5] (143.2567,237.2343) .. controls (143.2567,236.5135) and (142.6723,235.9292) ..(141.9515,235.9292) .. controls (141.2307,235.9292) and (140.6464,236.5135) ..(140.6464,237.2343) .. controls (140.6464,237.9551) and (141.2307,238.5395) ..(141.9515,238.5395) .. controls (142.6723,238.5395) and (143.2567,237.9551) ..(143.2567,237.2343);
\path[fill=chfourfig5c0,fill opacity=0.5,draw=chfourfig5c0,line width=0.3985pt,draw opacity=0.5] (72.5609,246.1909) .. controls (72.5609,245.4701) and (71.9765,244.8858) ..(71.2557,244.8858) .. controls (70.5349,244.8858) and (69.9506,245.4701) ..(69.9506,246.1909) .. controls (69.9506,246.9117) and (70.5349,247.4960) ..(71.2557,247.4960) .. controls (71.9765,247.4960) and (72.5609,246.9117) ..(72.5609,246.1909);
\path[fill=chfourfig5c0,fill opacity=0.5,draw=chfourfig5c0,line width=0.3985pt,draw opacity=0.5] (108.8453,248.1436) .. controls (108.8453,247.4228) and (108.2609,246.8385) ..(107.5401,246.8385) .. controls (106.8193,246.8385) and (106.2350,247.4228) ..(106.2350,248.1436) .. controls (106.2350,248.8644) and (106.8193,249.4487) ..(107.5401,249.4487) .. controls (108.2609,249.4487) and (108.8453,248.8644) ..(108.8453,248.1436);
\path[fill=chfourfig5c0,fill opacity=0.5,draw=chfourfig5c0,line width=0.3985pt,draw opacity=0.5] (54.9865,244.5172) .. controls (54.9865,243.7964) and (54.4022,243.2120) ..(53.6814,243.2120) .. controls (52.9606,243.2120) and (52.3763,243.7964) ..(52.3763,244.5172) .. controls (52.3763,245.2379) and (52.9606,245.8223) ..(53.6814,245.8223) .. controls (54.4022,245.8223) and (54.9865,245.2379) ..(54.9865,244.5172);
\path[fill=chfourfig5c0,fill opacity=0.5,draw=chfourfig5c0,line width=0.3985pt,draw opacity=0.5] (86.9471,175.6844) .. controls (86.9471,174.9636) and (86.3628,174.3793) ..(85.6420,174.3793) .. controls (84.9212,174.3793) and (84.3369,174.9636) ..(84.3369,175.6844) .. controls (84.3369,176.4052) and (84.9212,176.9895) ..(85.6420,176.9895) .. controls (86.3628,176.9895) and (86.9471,176.4052) ..(86.9471,175.6844);
\path[fill=chfourfig5c0,fill opacity=0.5,draw=chfourfig5c0,line width=0.3985pt,draw opacity=0.5] (113.7270,203.5602) .. controls (113.7270,202.8394) and (113.1427,202.2551) ..(112.4219,202.2551) .. controls (111.7011,202.2551) and (111.1168,202.8394) ..(111.1168,203.5602) .. controls (111.1168,204.2810) and (111.7011,204.8653) ..(112.4219,204.8653) .. controls (113.1427,204.8653) and (113.7270,204.2810) ..(113.7270,203.5602);
\path[fill=chfourfig5c0,fill opacity=0.5,draw=chfourfig5c0,line width=0.3985pt,draw opacity=0.5] (86.1800,213.6326) .. controls (86.1800,212.9117) and (85.5956,212.3275) ..(84.8748,212.3275) .. controls (84.1540,212.3275) and (83.5697,212.9117) ..(83.5697,213.6326) .. controls (83.5697,214.3534) and (84.1540,214.9377) ..(84.8748,214.9377) .. controls (85.5956,214.9377) and (86.1800,214.3534) ..(86.1800,213.6326);
\path[fill=chfourfig5c0,fill opacity=0.5,draw=chfourfig5c0,line width=0.3985pt,draw opacity=0.5] (87.9733,230.6589) .. controls (87.9733,229.9381) and (87.3889,229.3538) ..(86.6681,229.3538) .. controls (85.9474,229.3538) and (85.3630,229.9381) ..(85.3630,230.6589) .. controls (85.3630,231.3798) and (85.9474,231.9641) ..(86.6681,231.9641) .. controls (87.3889,231.9641) and (87.9733,231.3798) ..(87.9733,230.6589);
\path[fill=chfourfig5c0,fill opacity=0.5,draw=chfourfig5c0,line width=0.3985pt,draw opacity=0.5] (83.7291,254.2408) .. controls (83.7291,253.5200) and (83.1448,252.9357) ..(82.4240,252.9357) .. controls (81.7032,252.9357) and (81.1189,253.5200) ..(81.1189,254.2408) .. controls (81.1189,254.9616) and (81.7032,255.5459) ..(82.4240,255.5459) .. controls (83.1448,255.5459) and (83.7291,254.9616) ..(83.7291,254.2408);
\path[fill=chfourfig5c0,fill opacity=0.5,draw=chfourfig5c0,line width=0.3985pt,draw opacity=0.5] (150.0313,211.0124) .. controls (150.0313,210.2915) and (149.4470,209.7072) ..(148.7262,209.7072) .. controls (148.0054,209.7072) and (147.4211,210.2915) ..(147.4211,211.0124) .. controls (147.4211,211.7332) and (148.0054,212.3175) ..(148.7262,212.3175) .. controls (149.4470,212.3175) and (150.0313,211.7332) ..(150.0313,211.0124);
\path[fill=chfourfig5c0,fill opacity=0.5,draw=chfourfig5c0,line width=0.3985pt,draw opacity=0.5] (111.1467,199.7046) .. controls (111.1467,198.9838) and (110.5623,198.3995) ..(109.8415,198.3995) .. controls (109.1207,198.3995) and (108.5364,198.9838) ..(108.5364,199.7046) .. controls (108.5364,200.4254) and (109.1207,201.0097) ..(109.8415,201.0097) .. controls (110.5623,201.0097) and (111.1467,200.4254) ..(111.1467,199.7046);
\path[fill=chfourfig5c0,fill opacity=0.5,draw=chfourfig5c0,line width=0.3985pt,draw opacity=0.5] (95.3258,218.6837) .. controls (95.3258,217.9629) and (94.7415,217.3786) ..(94.0207,217.3786) .. controls (93.2999,217.3786) and (92.7155,217.9629) ..(92.7155,218.6837) .. controls (92.7155,219.4045) and (93.2999,219.9888) ..(94.0207,219.9888) .. controls (94.7415,219.9888) and (95.3258,219.4045) ..(95.3258,218.6837);
\path[fill=chfourfig5c0,fill opacity=0.5,draw=chfourfig5c0,line width=0.3985pt,draw opacity=0.5] (81.9259,207.9040) .. controls (81.9259,207.1832) and (81.3415,206.5988) ..(80.6207,206.5988) .. controls (79.8999,206.5988) and (79.3156,207.1832) ..(79.3156,207.9040) .. controls (79.3156,208.6248) and (79.8999,209.2091) ..(80.6207,209.2091) .. controls (81.3415,209.2091) and (81.9259,208.6248) ..(81.9259,207.9040);
\path[fill=chfourfig5c0,fill opacity=0.5,draw=chfourfig5c0,line width=0.3985pt,draw opacity=0.5] (158.7787,227.7697) .. controls (158.7787,227.0489) and (158.1943,226.4646) ..(157.4735,226.4646) .. controls (156.7527,226.4646) and (156.1684,227.0489) ..(156.1684,227.7697) .. controls (156.1684,228.4905) and (156.7527,229.0748) ..(157.4735,229.0748) .. controls (158.1943,229.0748) and (158.7787,228.4905) ..(158.7787,227.7697);
\path[fill=chfourfig5c0,fill opacity=0.5,draw=chfourfig5c0,line width=0.3985pt,draw opacity=0.5] (47.9329,173.2136) .. controls (47.9329,172.4928) and (47.3486,171.9085) ..(46.6278,171.9085) .. controls (45.9070,171.9085) and (45.3226,172.4928) ..(45.3226,173.2136) .. controls (45.3226,173.9344) and (45.9070,174.5187) ..(46.6278,174.5187) .. controls (47.3486,174.5187) and (47.9329,173.9344) ..(47.9329,173.2136);
\path[fill=chfourfig5c0,fill opacity=0.5,draw=chfourfig5c0,line width=0.3985pt,draw opacity=0.5] (114.4244,188.3371) .. controls (114.4244,187.6163) and (113.8401,187.0320) ..(113.1193,187.0320) .. controls (112.3985,187.0320) and (111.8142,187.6163) ..(111.8142,188.3371) .. controls (111.8142,189.0579) and (112.3985,189.6422) ..(113.1193,189.6422) .. controls (113.8401,189.6422) and (114.4244,189.0579) ..(114.4244,188.3371);
\path[fill=chfourfig5c0,fill opacity=0.5,draw=chfourfig5c0,line width=0.3985pt,draw opacity=0.5] (95.8737,256.2732) .. controls (95.8737,255.5524) and (95.2894,254.9681) ..(94.5686,254.9681) .. controls (93.8478,254.9681) and (93.2635,255.5524) ..(93.2635,256.2732) .. controls (93.2635,256.9940) and (93.8478,257.5783) ..(94.5686,257.5783) .. controls (95.2894,257.5783) and (95.8737,256.9940) ..(95.8737,256.2732);
\path[fill=chfourfig5c0,fill opacity=0.5,draw=chfourfig5c0,line width=0.3985pt,draw opacity=0.5] (79.9532,241.3290) .. controls (79.9532,240.6083) and (79.3689,240.0240) ..(78.6481,240.0240) .. controls (77.9273,240.0240) and (77.3430,240.6083) ..(77.3430,241.3290) .. controls (77.3430,242.0499) and (77.9273,242.6342) ..(78.6481,242.6342) .. controls (79.3689,242.6342) and (79.9532,242.0499) ..(79.9532,241.3290);
\path[fill=chfourfig5c0,fill opacity=0.5,draw=chfourfig5c0,line width=0.3985pt,draw opacity=0.5] (155.1323,240.3427) .. controls (155.1323,239.6219) and (154.5480,239.0376) ..(153.8272,239.0376) .. controls (153.1064,239.0376) and (152.5220,239.6219) ..(152.5220,240.3427) .. controls (152.5220,241.0635) and (153.1064,241.6479) ..(153.8272,241.6479) .. controls (154.5480,241.6479) and (155.1323,241.0635) ..(155.1323,240.3427);
\path[fill=chfourfig5c0,fill opacity=0.5,draw=chfourfig5c0,line width=0.3985pt,draw opacity=0.5] (154.2556,206.9974) .. controls (154.2556,206.2766) and (153.6712,205.6923) ..(152.9504,205.6923) .. controls (152.2296,205.6923) and (151.6453,206.2766) ..(151.6453,206.9974) .. controls (151.6453,207.7182) and (152.2296,208.3025) ..(152.9504,208.3025) .. controls (153.6712,208.3025) and (154.2556,207.7182) ..(154.2556,206.9974);
\path[fill=chfourfig5c0,fill opacity=0.5,draw=chfourfig5c0,line width=0.3985pt,draw opacity=0.5] (79.7341,261.3642) .. controls (79.7341,260.6434) and (79.1497,260.0591) ..(78.4289,260.0591) .. controls (77.7081,260.0591) and (77.1238,260.6434) ..(77.1238,261.3642) .. controls (77.1238,262.0850) and (77.7081,262.6693) ..(78.4289,262.6693) .. controls (79.1497,262.6693) and (79.7341,262.0850) ..(79.7341,261.3642);
\path[fill=chfourfig5c0,fill opacity=0.5,draw=chfourfig5c0,line width=0.3985pt,draw opacity=0.5] (83.6494,174.1700) .. controls (83.6494,173.4492) and (83.0651,172.8649) ..(82.3443,172.8649) .. controls (81.6235,172.8649) and (81.0392,173.4492) ..(81.0392,174.1700) .. controls (81.0392,174.8908) and (81.6235,175.4751) ..(82.3443,175.4751) .. controls (83.0651,175.4751) and (83.6494,174.8908) ..(83.6494,174.1700);
\path[fill=chfourfig5c0,fill opacity=0.5,draw=chfourfig5c0,line width=0.3985pt,draw opacity=0.5] (145.2691,191.1068) .. controls (145.2691,190.3860) and (144.6848,189.8016) ..(143.9640,189.8016) .. controls (143.2432,189.8016) and (142.6589,190.3860) ..(142.6589,191.1068) .. controls (142.6589,191.8275) and (143.2432,192.4119) ..(143.9640,192.4119) .. controls (144.6848,192.4119) and (145.2691,191.8275) ..(145.2691,191.1068);
\path[fill=chfourfig5c0,fill opacity=0.5,draw=chfourfig5c0,line width=0.3985pt,draw opacity=0.5] (157.1846,211.4507) .. controls (157.1846,210.7299) and (156.6003,210.1456) ..(155.8795,210.1456) .. controls (155.1587,210.1456) and (154.5744,210.7299) ..(154.5744,211.4507) .. controls (154.5744,212.1715) and (155.1587,212.7558) ..(155.8795,212.7558) .. controls (156.6003,212.7558) and (157.1846,212.1715) ..(157.1846,211.4507);
\path[fill=chfourfig5c0,fill opacity=0.5,draw=chfourfig5c0,line width=0.3985pt,draw opacity=0.5] (131.5504,261.7228) .. controls (131.5504,261.0020) and (130.9661,260.4177) ..(130.2453,260.4177) .. controls (129.5245,260.4177) and (128.9402,261.0020) ..(128.9402,261.7228) .. controls (128.9402,262.4436) and (129.5245,263.0280) ..(130.2453,263.0280) .. controls (130.9661,263.0280) and (131.5504,262.4436) ..(131.5504,261.7228);
\path[fill=chfourfig5c0,fill opacity=0.5,draw=chfourfig5c0,line width=0.3985pt,draw opacity=0.5] (75.6792,180.3270) .. controls (75.6792,179.6062) and (75.0949,179.0219) ..(74.3741,179.0219) .. controls (73.6533,179.0219) and (73.0690,179.6062) ..(73.0690,180.3270) .. controls (73.0690,181.0478) and (73.6533,181.6321) ..(74.3741,181.6321) .. controls (75.0949,181.6321) and (75.6792,181.0478) ..(75.6792,180.3270);
\path[fill=chfourfig5c0,fill opacity=0.5,draw=chfourfig5c0,line width=0.3985pt,draw opacity=0.5] (129.9663,178.3444) .. controls (129.9663,177.6236) and (129.3820,177.0393) ..(128.6612,177.0393) .. controls (127.9404,177.0393) and (127.3561,177.6236) ..(127.3561,178.3444) .. controls (127.3561,179.0652) and (127.9404,179.6496) ..(128.6612,179.6496) .. controls (129.3820,179.6496) and (129.9663,179.0652) ..(129.9663,178.3444);
\path[fill=chfourfig5c0,fill opacity=0.5,draw=chfourfig5c0,line width=0.3985pt,draw opacity=0.5] (148.8358,261.8723) .. controls (148.8358,261.1515) and (148.2515,260.5671) ..(147.5307,260.5671) .. controls (146.8099,260.5671) and (146.2256,261.1515) ..(146.2256,261.8723) .. controls (146.2256,262.5931) and (146.8099,263.1774) ..(147.5307,263.1774) .. controls (148.2515,263.1774) and (148.8358,262.5931) ..(148.8358,261.8723);
\path[fill=chfourfig5c0,fill opacity=0.5,draw=chfourfig5c0,line width=0.3985pt,draw opacity=0.5] (66.0651,229.1545) .. controls (66.0651,228.4337) and (65.4808,227.8494) ..(64.7600,227.8494) .. controls (64.0392,227.8494) and (63.4549,228.4337) ..(63.4549,229.1545) .. controls (63.4549,229.8754) and (64.0392,230.4597) ..(64.7600,230.4597) .. controls (65.4808,230.4597) and (66.0651,229.8754) ..(66.0651,229.1545);
\path[fill=chfourfig5c0,fill opacity=0.5,draw=chfourfig5c0,line width=0.3985pt,draw opacity=0.5] (158.4100,234.9728) .. controls (158.4100,234.2520) and (157.8257,233.6677) ..(157.1049,233.6677) .. controls (156.3841,233.6677) and (155.7998,234.2520) ..(155.7998,234.9728) .. controls (155.7998,235.6936) and (156.3841,236.2779) ..(157.1049,236.2779) .. controls (157.8257,236.2779) and (158.4100,235.6936) ..(158.4100,234.9728);
\path[fill=chfourfig5c0,fill opacity=0.5,draw=chfourfig5c0,line width=0.3985pt,draw opacity=0.5] (99.1714,240.1335) .. controls (99.1714,239.4127) and (98.5871,238.8284) ..(97.8663,238.8284) .. controls (97.1455,238.8284) and (96.5612,239.4127) ..(96.5612,240.1335) .. controls (96.5612,240.8543) and (97.1455,241.4387) ..(97.8663,241.4387) .. controls (98.5871,241.4387) and (99.1714,240.8543) ..(99.1714,240.1335);
\path[fill=chfourfig5c0,fill opacity=0.5,draw=chfourfig5c0,line width=0.3985pt,draw opacity=0.5] (158.0613,182.4092) .. controls (158.0613,181.6884) and (157.4770,181.1041) ..(156.7562,181.1041) .. controls (156.0354,181.1041) and (155.4511,181.6884) ..(155.4511,182.4092) .. controls (155.4511,183.1301) and (156.0354,183.7144) ..(156.7562,183.7144) .. controls (157.4770,183.7144) and (158.0613,183.1301) ..(158.0613,182.4092);
\path[fill=chfourfig5c0,fill opacity=0.5,draw=chfourfig5c0,line width=0.3985pt,draw opacity=0.5] (136.4720,206.7682) .. controls (136.4720,206.0474) and (135.8877,205.4631) ..(135.1669,205.4631) .. controls (134.4461,205.4631) and (133.8618,206.0474) ..(133.8618,206.7682) .. controls (133.8618,207.4890) and (134.4461,208.0733) ..(135.1669,208.0733) .. controls (135.8877,208.0733) and (136.4720,207.4890) ..(136.4720,206.7682);
\path[fill=chfourfig5c0,fill opacity=0.5,draw=chfourfig5c0,line width=0.3985pt,draw opacity=0.5] (104.5314,200.8105) .. controls (104.5314,200.0897) and (103.9471,199.5054) ..(103.2263,199.5054) .. controls (102.5055,199.5054) and (101.9211,200.0897) ..(101.9211,200.8105) .. controls (101.9211,201.5313) and (102.5055,202.1156) ..(103.2263,202.1156) .. controls (103.9471,202.1156) and (104.5314,201.5313) ..(104.5314,200.8105);
\path[fill=chfourfig5c0,fill opacity=0.5,draw=chfourfig5c0,line width=0.3985pt,draw opacity=0.5] (84.8649,211.0223) .. controls (84.8649,210.3015) and (84.2805,209.7172) ..(83.5598,209.7172) .. controls (82.8389,209.7172) and (82.2546,210.3015) ..(82.2546,211.0223) .. controls (82.2546,211.7431) and (82.8389,212.3275) ..(83.5598,212.3275) .. controls (84.2805,212.3275) and (84.8649,211.7431) ..(84.8649,211.0223);
\path[fill=chfourfig5c0,fill opacity=0.5,draw=chfourfig5c0,line width=0.3985pt,draw opacity=0.5] (111.2762,249.1000) .. controls (111.2762,248.3792) and (110.6919,247.7949) ..(109.9711,247.7949) .. controls (109.2503,247.7949) and (108.6659,248.3792) ..(108.6659,249.1000) .. controls (108.6659,249.8208) and (109.2503,250.4052) ..(109.9711,250.4052) .. controls (110.6919,250.4052) and (111.2762,249.8208) ..(111.2762,249.1000);
\path[fill=chfourfig5c0,fill opacity=0.5,draw=chfourfig5c0,line width=0.3985pt,draw opacity=0.5] (80.3617,182.6185) .. controls (80.3617,181.8976) and (79.7774,181.3134) ..(79.0566,181.3134) .. controls (78.3358,181.3134) and (77.7515,181.8976) ..(77.7515,182.6185) .. controls (77.7515,183.3393) and (78.3358,183.9236) ..(79.0566,183.9236) .. controls (79.7774,183.9236) and (80.3617,183.3393) ..(80.3617,182.6185);
\path[fill=chfourfig5c0,fill opacity=0.5,draw=chfourfig5c0,line width=0.3985pt,draw opacity=0.5] (97.0792,229.4036) .. controls (97.0792,228.6828) and (96.4949,228.0985) ..(95.7741,228.0985) .. controls (95.0533,228.0985) and (94.4690,228.6828) ..(94.4690,229.4036) .. controls (94.4690,230.1244) and (95.0533,230.7087) ..(95.7741,230.7087) .. controls (96.4949,230.7087) and (97.0792,230.1244) ..(97.0792,229.4036);
\path[fill=chfourfig5c0,fill opacity=0.5,draw=chfourfig5c0,line width=0.3985pt,draw opacity=0.5] (93.7716,210.9625) .. controls (93.7716,210.2418) and (93.1873,209.6574) ..(92.4665,209.6574) .. controls (91.7457,209.6574) and (91.1613,210.2418) ..(91.1613,210.9625) .. controls (91.1613,211.6833) and (91.7457,212.2676) ..(92.4665,212.2676) .. controls (93.1873,212.2676) and (93.7716,211.6833) ..(93.7716,210.9625);
\path[fill=chfourfig5c0,fill opacity=0.5,draw=chfourfig5c0,line width=0.3985pt,draw opacity=0.5] (72.5409,179.8289) .. controls (72.5409,179.1081) and (71.9566,178.5238) ..(71.2358,178.5238) .. controls (70.5150,178.5238) and (69.9307,179.1081) ..(69.9307,179.8289) .. controls (69.9307,180.5497) and (70.5150,181.1340) ..(71.2358,181.1340) .. controls (71.9566,181.1340) and (72.5409,180.5497) ..(72.5409,179.8289);
\path[fill=chfourfig5c0,fill opacity=0.5,draw=chfourfig5c0,line width=0.3985pt,draw opacity=0.5] (113.7270,175.1663) .. controls (113.7270,174.4455) and (113.1427,173.8612) ..(112.4219,173.8612) .. controls (111.7011,173.8612) and (111.1168,174.4455) ..(111.1168,175.1663) .. controls (111.1168,175.8871) and (111.7011,176.4714) ..(112.4219,176.4714) .. controls (113.1427,176.4714) and (113.7270,175.8871) ..(113.7270,175.1663);
\path[fill=chfourfig5c0,fill opacity=0.5,draw=chfourfig5c0,line width=0.3985pt,draw opacity=0.5] (54.7972,214.2901) .. controls (54.7972,213.5693) and (54.2129,212.9850) ..(53.4921,212.9850) .. controls (52.7713,212.9850) and (52.1870,213.5693) ..(52.1870,214.2901) .. controls (52.1870,215.0109) and (52.7713,215.5952) ..(53.4921,215.5952) .. controls (54.2129,215.5952) and (54.7972,215.0109) ..(54.7972,214.2901);
\path[fill=chfourfig5c0,fill opacity=0.5,draw=chfourfig5c0,line width=0.3985pt,draw opacity=0.5] (130.8331,192.7008) .. controls (130.8331,191.9800) and (130.2488,191.3957) ..(129.5280,191.3957) .. controls (128.8072,191.3957) and (128.2229,191.9800) ..(128.2229,192.7008) .. controls (128.2229,193.4216) and (128.8072,194.0059) ..(129.5280,194.0059) .. controls (130.2488,194.0059) and (130.8331,193.4216) ..(130.8331,192.7008);
\path[fill=chfourfig5c0,fill opacity=0.5,draw=chfourfig5c0,line width=0.3985pt,draw opacity=0.5] (80.8200,259.3517) .. controls (80.8200,258.6309) and (80.2357,258.0466) ..(79.5149,258.0466) .. controls (78.7941,258.0466) and (78.2097,258.6309) ..(78.2097,259.3517) .. controls (78.2097,260.0725) and (78.7941,260.6568) ..(79.5149,260.6568) .. controls (80.2357,260.6568) and (80.8200,260.0725) ..(80.8200,259.3517);
\path[fill=chfourfig5c0,fill opacity=0.5,draw=chfourfig5c0,line width=0.3985pt,draw opacity=0.5] (112.1330,175.5549) .. controls (112.1330,174.8340) and (111.5487,174.2497) ..(110.8279,174.2497) .. controls (110.1071,174.2497) and (109.5227,174.8340) ..(109.5227,175.5549) .. controls (109.5227,176.2757) and (110.1071,176.8600) ..(110.8279,176.8600) .. controls (111.5487,176.8600) and (112.1330,176.2757) ..(112.1330,175.5549);
\path[fill=chfourfig5c0,fill opacity=0.5,draw=chfourfig5c0,line width=0.3985pt,draw opacity=0.5] (133.2042,218.9228) .. controls (133.2042,218.2020) and (132.6199,217.6177) ..(131.8991,217.6177) .. controls (131.1783,217.6177) and (130.5940,218.2020) ..(130.5940,218.9228) .. controls (130.5940,219.6436) and (131.1783,220.2279) ..(131.8991,220.2279) .. controls (132.6199,220.2279) and (133.2042,219.6436) ..(133.2042,218.9228);
\path[fill=chfourfig5c0,fill opacity=0.5,draw=chfourfig5c0,line width=0.3985pt,draw opacity=0.5] (129.8368,192.9698) .. controls (129.8368,192.2490) and (129.2525,191.6647) ..(128.5317,191.6647) .. controls (127.8109,191.6647) and (127.2266,192.2490) ..(127.2266,192.9698) .. controls (127.2266,193.6906) and (127.8109,194.2749) ..(128.5317,194.2749) .. controls (129.2525,194.2749) and (129.8368,193.6906) ..(129.8368,192.9698);
\path[fill=chfourfig5c0,fill opacity=0.5,draw=chfourfig5c0,line width=0.3985pt,draw opacity=0.5] (66.2943,192.1728) .. controls (66.2943,191.4520) and (65.7099,190.8676) ..(64.9892,190.8676) .. controls (64.2683,190.8676) and (63.6840,191.4520) ..(63.6840,192.1728) .. controls (63.6840,192.8936) and (64.2683,193.4779) ..(64.9892,193.4779) .. controls (65.7099,193.4779) and (66.2943,192.8936) ..(66.2943,192.1728);
\path[fill=chfourfig5c0,fill opacity=0.5,draw=chfourfig5c0,line width=0.3985pt,draw opacity=0.5] (46.2392,191.0470) .. controls (46.2392,190.3262) and (45.6549,189.7418) ..(44.9341,189.7418) .. controls (44.2133,189.7418) and (43.6290,190.3262) ..(43.6290,191.0470) .. controls (43.6290,191.7678) and (44.2133,192.3521) ..(44.9341,192.3521) .. controls (45.6549,192.3521) and (46.2392,191.7678) ..(46.2392,191.0470);
\path[fill=chfourfig5c0,fill opacity=0.5,draw=chfourfig5c0,line width=0.3985pt,draw opacity=0.5] (119.3759,200.7108) .. controls (119.3759,199.9901) and (118.7916,199.4057) ..(118.0708,199.4057) .. controls (117.3500,199.4057) and (116.7657,199.9901) ..(116.7657,200.7108) .. controls (116.7657,201.4316) and (117.3500,202.0160) ..(118.0708,202.0160) .. controls (118.7916,202.0160) and (119.3759,201.4316) ..(119.3759,200.7108);
\path[fill=chfourfig5c0,fill opacity=0.5,draw=chfourfig5c0,line width=0.3985pt,draw opacity=0.5] (63.3553,259.1823) .. controls (63.3553,258.4615) and (62.7709,257.8772) ..(62.0501,257.8772) .. controls (61.3293,257.8772) and (60.7450,258.4615) ..(60.7450,259.1823) .. controls (60.7450,259.9031) and (61.3293,260.4875) ..(62.0501,260.4875) .. controls (62.7709,260.4875) and (63.3553,259.9031) ..(63.3553,259.1823);
\path[fill=chfourfig5c0,fill opacity=0.5,draw=chfourfig5c0,line width=0.3985pt,draw opacity=0.5] (62.2295,200.9201) .. controls (62.2295,200.1993) and (61.6452,199.6149) ..(60.9243,199.6149) .. controls (60.2036,199.6149) and (59.6192,200.1993) ..(59.6192,200.9201) .. controls (59.6192,201.6409) and (60.2036,202.2252) ..(60.9243,202.2252) .. controls (61.6452,202.2252) and (62.2295,201.6409) ..(62.2295,200.9201);
\path[fill=chfourfig5c0,fill opacity=0.5,draw=chfourfig5c0,line width=0.3985pt,draw opacity=0.5] (119.2065,204.8553) .. controls (119.2065,204.1346) and (118.6222,203.5502) ..(117.9014,203.5502) .. controls (117.1806,203.5502) and (116.5963,204.1346) ..(116.5963,204.8553) .. controls (116.5963,205.5762) and (117.1806,206.1605) ..(117.9014,206.1605) .. controls (118.6222,206.1605) and (119.2065,205.5762) ..(119.2065,204.8553);
\path[fill=chfourfig5c0,fill opacity=0.5,draw=chfourfig5c0,line width=0.3985pt,draw opacity=0.5] (115.8491,265.0106) .. controls (115.8491,264.2898) and (115.2648,263.7054) ..(114.5440,263.7054) .. controls (113.8232,263.7054) and (113.2388,264.2898) ..(113.2388,265.0106) .. controls (113.2388,265.7314) and (113.8232,266.3157) ..(114.5440,266.3157) .. controls (115.2648,266.3157) and (115.8491,265.7314) ..(115.8491,265.0106);
\path[fill=chfourfig5c0,fill opacity=0.5,draw=chfourfig5c0,line width=0.3985pt,draw opacity=0.5] (121.5478,237.4236) .. controls (121.5478,236.7029) and (120.9635,236.1185) ..(120.2427,236.1185) .. controls (119.5219,236.1185) and (118.9375,236.7029) ..(118.9375,237.4236) .. controls (118.9375,238.1445) and (119.5219,238.7288) ..(120.2427,238.7288) .. controls (120.9635,238.7288) and (121.5478,238.1445) ..(121.5478,237.4236);
\path[fill=chfourfig5c0,fill opacity=0.5,draw=chfourfig5c0,line width=0.3985pt,draw opacity=0.5] (159.0975,200.2326) .. controls (159.0975,199.5118) and (158.5132,198.9275) ..(157.7923,198.9275) .. controls (157.0715,198.9275) and (156.4872,199.5118) ..(156.4872,200.2326) .. controls (156.4872,200.9534) and (157.0715,201.5378) ..(157.7923,201.5378) .. controls (158.5132,201.5378) and (159.0975,200.9534) ..(159.0975,200.2326);
\path[fill=chfourfig5c0,fill opacity=0.5,draw=chfourfig5c0,line width=0.3985pt,draw opacity=0.5] (145.5282,251.7103) .. controls (145.5282,250.9894) and (144.9439,250.4052) ..(144.2231,250.4052) .. controls (143.5023,250.4052) and (142.9179,250.9894) ..(142.9179,251.7103) .. controls (142.9179,252.4311) and (143.5023,253.0154) ..(144.2231,253.0154) .. controls (144.9439,253.0154) and (145.5282,252.4311) ..(145.5282,251.7103);
\path[fill=chfourfig5c0,fill opacity=0.5,draw=chfourfig5c0,line width=0.3985pt,draw opacity=0.5] (57.7163,261.6432) .. controls (57.7163,260.9223) and (57.1320,260.3380) ..(56.4112,260.3380) .. controls (55.6904,260.3380) and (55.1061,260.9223) ..(55.1061,261.6432) .. controls (55.1061,262.3639) and (55.6904,262.9482) ..(56.4112,262.9482) .. controls (57.1320,262.9482) and (57.7163,262.3639) ..(57.7163,261.6432);
\path[fill=chfourfig5c0,fill opacity=0.5,draw=chfourfig5c0,line width=0.3985pt,draw opacity=0.5] (63.2258,237.3041) .. controls (63.2258,236.5833) and (62.6414,235.9990) ..(61.9206,235.9990) .. controls (61.1998,235.9990) and (60.6155,236.5833) ..(60.6155,237.3041) .. controls (60.6155,238.0249) and (61.1998,238.6092) ..(61.9206,238.6092) .. controls (62.6414,238.6092) and (63.2258,238.0249) ..(63.2258,237.3041);
\path[fill=chfourfig5c0,fill opacity=0.5,draw=chfourfig5c0,line width=0.3985pt,draw opacity=0.5] (73.0291,186.3346) .. controls (73.0291,185.6138) and (72.4448,185.0294) ..(71.7240,185.0294) .. controls (71.0032,185.0294) and (70.4189,185.6138) ..(70.4189,186.3346) .. controls (70.4189,187.0554) and (71.0032,187.6397) ..(71.7240,187.6397) .. controls (72.4448,187.6397) and (73.0291,187.0554) ..(73.0291,186.3346);
\path[fill=chfourfig5c0,fill opacity=0.5,draw=chfourfig5c0,line width=0.3985pt,draw opacity=0.5] (45.1832,251.3715) .. controls (45.1832,250.6507) and (44.5989,250.0664) ..(43.8781,250.0664) .. controls (43.1573,250.0664) and (42.5729,250.6507) ..(42.5729,251.3715) .. controls (42.5729,252.0923) and (43.1573,252.6767) ..(43.8781,252.6767) .. controls (44.5989,252.6767) and (45.1832,252.0923) ..(45.1832,251.3715);
\path[fill=chfourfig5c0,fill opacity=0.5,draw=chfourfig5c0,line width=0.3985pt,draw opacity=0.5] (133.5928,185.1590) .. controls (133.5928,184.4382) and (133.0085,183.8538) ..(132.2877,183.8538) .. controls (131.5668,183.8538) and (130.9825,184.4382) ..(130.9825,185.1590) .. controls (130.9825,185.8798) and (131.5668,186.4641) ..(132.2877,186.4641) .. controls (133.0085,186.4641) and (133.5928,185.8798) ..(133.5928,185.1590);
\path[fill=chfourfig5c0,fill opacity=0.5,draw=chfourfig5c0,line width=0.3985pt,draw opacity=0.5] (143.9341,178.7230) .. controls (143.9341,178.0022) and (143.3498,177.4179) ..(142.6290,177.4179) .. controls (141.9082,177.4179) and (141.3239,178.0022) ..(141.3239,178.7230) .. controls (141.3239,179.4438) and (141.9082,180.0281) ..(142.6290,180.0281) .. controls (143.3498,180.0281) and (143.9341,179.4438) ..(143.9341,178.7230);
\path[fill=chfourfig5c0,fill opacity=0.5,draw=chfourfig5c0,line width=0.3985pt,draw opacity=0.5] (137.2093,234.1360) .. controls (137.2093,233.4151) and (136.6249,232.8308) ..(135.9041,232.8308) .. controls (135.1833,232.8308) and (134.5990,233.4151) ..(134.5990,234.1360) .. controls (134.5990,234.8568) and (135.1833,235.4410) ..(135.9041,235.4410) .. controls (136.6249,235.4410) and (137.2093,234.8568) ..(137.2093,234.1360);
\path[fill=chfourfig5c0,fill opacity=0.5,draw=chfourfig5c0,line width=0.3985pt,draw opacity=0.5] (63.2755,195.2911) .. controls (63.2755,194.5703) and (62.6912,193.9860) ..(61.9704,193.9860) .. controls (61.2496,193.9860) and (60.6653,194.5703) ..(60.6653,195.2911) .. controls (60.6653,196.0119) and (61.2496,196.5963) ..(61.9704,196.5963) .. controls (62.6912,196.5963) and (63.2755,196.0119) ..(63.2755,195.2911);
\path[fill=chfourfig5c0,fill opacity=0.5,draw=chfourfig5c0,line width=0.3985pt,draw opacity=0.5] (72.3915,222.2105) .. controls (72.3915,221.4897) and (71.8072,220.9054) ..(71.0864,220.9054) .. controls (70.3656,220.9054) and (69.7812,221.4897) ..(69.7812,222.2105) .. controls (69.7812,222.9313) and (70.3656,223.5156) ..(71.0864,223.5156) .. controls (71.8072,223.5156) and (72.3915,222.9313) ..(72.3915,222.2105);
\path[fill=chfourfig5c0,fill opacity=0.5,draw=chfourfig5c0,line width=0.3985pt,draw opacity=0.5] (109.0146,189.7518) .. controls (109.0146,189.0310) and (108.4303,188.4467) ..(107.7095,188.4467) .. controls (106.9887,188.4467) and (106.4044,189.0310) ..(106.4044,189.7518) .. controls (106.4044,190.4726) and (106.9887,191.0569) ..(107.7095,191.0569) .. controls (108.4303,191.0569) and (109.0146,190.4726) ..(109.0146,189.7518);
\path[fill=chfourfig5c0,fill opacity=0.5,draw=chfourfig5c0,line width=0.3985pt,draw opacity=0.5] (110.3696,225.9465) .. controls (110.3696,225.2257) and (109.7852,224.6414) ..(109.0644,224.6414) .. controls (108.3436,224.6414) and (107.7593,225.2257) ..(107.7593,225.9465) .. controls (107.7593,226.6674) and (108.3436,227.2516) ..(109.0644,227.2516) .. controls (109.7852,227.2516) and (110.3696,226.6674) ..(110.3696,225.9465);
\path[fill=chfourfig5c0,fill opacity=0.5,draw=chfourfig5c0,line width=0.3985pt,draw opacity=0.5] (126.3797,199.2663) .. controls (126.3797,198.5454) and (125.7954,197.9611) ..(125.0746,197.9611) .. controls (124.3538,197.9611) and (123.7695,198.5454) ..(123.7695,199.2663) .. controls (123.7695,199.9871) and (124.3538,200.5714) ..(125.0746,200.5714) .. controls (125.7954,200.5714) and (126.3797,199.9871) ..(126.3797,199.2663);
\path[fill=chfourfig5c0,fill opacity=0.5,draw=chfourfig5c0,line width=0.3985pt,draw opacity=0.5] (146.1160,201.7968) .. controls (146.1160,201.0760) and (145.5317,200.4917) ..(144.8109,200.4917) .. controls (144.0901,200.4917) and (143.5057,201.0760) ..(143.5057,201.7968) .. controls (143.5057,202.5176) and (144.0901,203.1019) ..(144.8109,203.1019) .. controls (145.5317,203.1019) and (146.1160,202.5176) ..(146.1160,201.7968);
\path[fill=chfourfig5c0,fill opacity=0.5,draw=chfourfig5c0,line width=0.3985pt,draw opacity=0.5] (117.5129,223.6750) .. controls (117.5129,222.9542) and (116.9285,222.3699) ..(116.2077,222.3699) .. controls (115.4869,222.3699) and (114.9026,222.9542) ..(114.9026,223.6750) .. controls (114.9026,224.3958) and (115.4869,224.9802) ..(116.2077,224.9802) .. controls (116.9285,224.9802) and (117.5129,224.3958) ..(117.5129,223.6750);
\path[fill=chfourfig5c0,fill opacity=0.5,draw=chfourfig5c0,line width=0.3985pt,draw opacity=0.5] (79.7241,208.2128) .. controls (79.7241,207.4920) and (79.1398,206.9077) ..(78.4190,206.9077) .. controls (77.6982,206.9077) and (77.1138,207.4920) ..(77.1138,208.2128) .. controls (77.1138,208.9336) and (77.6982,209.5179) ..(78.4190,209.5179) .. controls (79.1398,209.5179) and (79.7241,208.9336) ..(79.7241,208.2128);
\path[fill=chfourfig5c0,fill opacity=0.5,draw=chfourfig5c0,line width=0.3985pt,draw opacity=0.5] (85.0840,173.0443) .. controls (85.0840,172.3234) and (84.4997,171.7391) ..(83.7789,171.7391) .. controls (83.0581,171.7391) and (82.4738,172.3234) ..(82.4738,173.0443) .. controls (82.4738,173.7650) and (83.0581,174.3494) ..(83.7789,174.3494) .. controls (84.4997,174.3494) and (85.0840,173.7650) ..(85.0840,173.0443);
\path[fill=chfourfig5c0,fill opacity=0.5,draw=chfourfig5c0,line width=0.3985pt,draw opacity=0.5] (47.9429,180.2872) .. controls (47.9429,179.5664) and (47.3585,178.9821) ..(46.6377,178.9821) .. controls (45.9169,178.9821) and (45.3326,179.5664) ..(45.3326,180.2872) .. controls (45.3326,181.0080) and (45.9169,181.5923) ..(46.6377,181.5923) .. controls (47.3585,181.5923) and (47.9429,181.0080) ..(47.9429,180.2872);
\path[fill=chfourfig5c0,fill opacity=0.5,draw=chfourfig5c0,line width=0.3985pt,draw opacity=0.5] (51.4498,251.5907) .. controls (51.4498,250.8699) and (50.8654,250.2856) ..(50.1446,250.2856) .. controls (49.4238,250.2856) and (48.8395,250.8699) ..(48.8395,251.5907) .. controls (48.8395,252.3115) and (49.4238,252.8958) ..(50.1446,252.8958) .. controls (50.8654,252.8958) and (51.4498,252.3115) ..(51.4498,251.5907);
\path[fill=chfourfig5c0,fill opacity=0.5,draw=chfourfig5c0,line width=0.3985pt,draw opacity=0.5] (83.4800,217.0996) .. controls (83.4800,216.3788) and (82.8957,215.7945) ..(82.1749,215.7945) .. controls (81.4541,215.7945) and (80.8698,216.3788) ..(80.8698,217.0996) .. controls (80.8698,217.8204) and (81.4541,218.4047) ..(82.1749,218.4047) .. controls (82.8957,218.4047) and (83.4800,217.8204) ..(83.4800,217.0996);
\path[fill=chfourfig5c0,fill opacity=0.5,draw=chfourfig5c0,line width=0.3985pt,draw opacity=0.5] (154.5246,210.2651) .. controls (154.5246,209.5443) and (153.9402,208.9600) ..(153.2194,208.9600) .. controls (152.4986,208.9600) and (151.9143,209.5443) ..(151.9143,210.2651) .. controls (151.9143,210.9860) and (152.4986,211.5703) ..(153.2194,211.5703) .. controls (153.9402,211.5703) and (154.5246,210.9860) ..(154.5246,210.2651);
\path[fill=chfourfig5c0,fill opacity=0.5,draw=chfourfig5c0,line width=0.3985pt,draw opacity=0.5] (116.3672,198.1404) .. controls (116.3672,197.4197) and (115.7828,196.8353) ..(115.0620,196.8353) .. controls (114.3412,196.8353) and (113.7569,197.4197) ..(113.7569,198.1404) .. controls (113.7569,198.8613) and (114.3412,199.4456) ..(115.0620,199.4456) .. controls (115.7828,199.4456) and (116.3672,198.8613) ..(116.3672,198.1404);
\path[fill=chfourfig5c0,fill opacity=0.5,draw=chfourfig5c0,line width=0.3985pt,draw opacity=0.5] (130.8530,236.1086) .. controls (130.8530,235.3878) and (130.2687,234.8034) ..(129.5479,234.8034) .. controls (128.8271,234.8034) and (128.2428,235.3878) ..(128.2428,236.1086) .. controls (128.2428,236.8294) and (128.8271,237.4137) ..(129.5479,237.4137) .. controls (130.2687,237.4137) and (130.8530,236.8294) ..(130.8530,236.1086);
\path[fill=chfourfig5c0,fill opacity=0.5,draw=chfourfig5c0,line width=0.3985pt,draw opacity=0.5] (56.4012,195.1118) .. controls (56.4012,194.3910) and (55.8169,193.8066) ..(55.0961,193.8066) .. controls (54.3753,193.8066) and (53.7910,194.3910) ..(53.7910,195.1118) .. controls (53.7910,195.8326) and (54.3753,196.4169) ..(55.0961,196.4169) .. controls (55.8169,196.4169) and (56.4012,195.8326) ..(56.4012,195.1118);
\path[fill=chfourfig5c0,fill opacity=0.5,draw=chfourfig5c0,line width=0.3985pt,draw opacity=0.5] (114.8727,207.2863) .. controls (114.8727,206.5655) and (114.2884,205.9811) ..(113.5676,205.9811) .. controls (112.8468,205.9811) and (112.2625,206.5655) ..(112.2625,207.2863) .. controls (112.2625,208.0071) and (112.8468,208.5914) ..(113.5676,208.5914) .. controls (114.2884,208.5914) and (114.8727,208.0071) ..(114.8727,207.2863);
\path[fill=chfourfig5c0,fill opacity=0.5,draw=chfourfig5c0,line width=0.3985pt,draw opacity=0.5] (164.8061,230.9379) .. controls (164.8061,230.2171) and (164.2218,229.6328) ..(163.5010,229.6328) .. controls (162.7802,229.6328) and (162.1959,230.2171) ..(162.1959,230.9379) .. controls (162.1959,231.6587) and (162.7802,232.2430) ..(163.5010,232.2430) .. controls (164.2218,232.2430) and (164.8061,231.6587) ..(164.8061,230.9379);
\path[fill=chfourfig5c0,fill opacity=0.5,draw=chfourfig5c0,line width=0.3985pt,draw opacity=0.5] (75.1910,183.7442) .. controls (75.1910,183.0235) and (74.6067,182.4391) ..(73.8859,182.4391) .. controls (73.1651,182.4391) and (72.5808,183.0235) ..(72.5808,183.7442) .. controls (72.5808,184.4651) and (73.1651,185.0494) ..(73.8859,185.0494) .. controls (74.6067,185.0494) and (75.1910,184.4651) ..(75.1910,183.7442);
\path[fill=chfourfig5c0,fill opacity=0.5,draw=chfourfig5c0,line width=0.3985pt,draw opacity=0.5] (122.1954,177.9160) .. controls (122.1954,177.1952) and (121.6110,176.6109) ..(120.8902,176.6109) .. controls (120.1694,176.6109) and (119.5851,177.1952) ..(119.5851,177.9160) .. controls (119.5851,178.6368) and (120.1694,179.2212) ..(120.8902,179.2212) .. controls (121.6110,179.2212) and (122.1954,178.6368) ..(122.1954,177.9160);
\path[fill=chfourfig5c0,fill opacity=0.5,draw=chfourfig5c0,line width=0.3985pt,draw opacity=0.5] (105.1391,248.7015) .. controls (105.1391,247.9807) and (104.5548,247.3964) ..(103.8340,247.3964) .. controls (103.1132,247.3964) and (102.5289,247.9807) ..(102.5289,248.7015) .. controls (102.5289,249.4223) and (103.1132,250.0066) ..(103.8340,250.0066) .. controls (104.5548,250.0066) and (105.1391,249.4223) ..(105.1391,248.7015);
\path[fill=chfourfig5c0,fill opacity=0.5,draw=chfourfig5c0,line width=0.3985pt,draw opacity=0.5] (135.7746,234.4846) .. controls (135.7746,233.7638) and (135.1903,233.1795) ..(134.4695,233.1795) .. controls (133.7487,233.1795) and (133.1644,233.7638) ..(133.1644,234.4846) .. controls (133.1644,235.2054) and (133.7487,235.7898) ..(134.4695,235.7898) .. controls (135.1903,235.7898) and (135.7746,235.2054) ..(135.7746,234.4846);
\path[fill=chfourfig5c0,fill opacity=0.5,draw=chfourfig5c0,line width=0.3985pt,draw opacity=0.5] (128.3623,228.4871) .. controls (128.3623,227.7663) and (127.7780,227.1819) ..(127.0572,227.1819) .. controls (126.3364,227.1819) and (125.7521,227.7663) ..(125.7521,228.4871) .. controls (125.7521,229.2079) and (126.3364,229.7922) ..(127.0572,229.7922) .. controls (127.7780,229.7922) and (128.3623,229.2079) ..(128.3623,228.4871);
\path[fill=chfourfig5c0,fill opacity=0.5,draw=chfourfig5c0,line width=0.3985pt,draw opacity=0.5] (112.0234,259.1126) .. controls (112.0234,258.3918) and (111.4391,257.8075) ..(110.7183,257.8075) .. controls (109.9975,257.8075) and (109.4131,258.3918) ..(109.4131,259.1126) .. controls (109.4131,259.8334) and (109.9975,260.4177) ..(110.7183,260.4177) .. controls (111.4391,260.4177) and (112.0234,259.8334) ..(112.0234,259.1126);
\path[fill=chfourfig5c0,fill opacity=0.5,draw=chfourfig5c0,line width=0.3985pt,draw opacity=0.5] (132.1880,237.0251) .. controls (132.1880,236.3044) and (131.6037,235.7200) ..(130.8829,235.7200) .. controls (130.1621,235.7200) and (129.5778,236.3044) ..(129.5778,237.0251) .. controls (129.5778,237.7460) and (130.1621,238.3303) ..(130.8829,238.3303) .. controls (131.6037,238.3303) and (132.1880,237.7460) ..(132.1880,237.0251);
\path[fill=chfourfig5c0,fill opacity=0.5,draw=chfourfig5c0,line width=0.3985pt,draw opacity=0.5] (92.2772,214.5292) .. controls (92.2772,213.8084) and (91.6929,213.2241) ..(90.9720,213.2241) .. controls (90.2513,213.2241) and (89.6669,213.8084) ..(89.6669,214.5292) .. controls (89.6669,215.2500) and (90.2513,215.8344) ..(90.9720,215.8344) .. controls (91.6929,215.8344) and (92.2772,215.2500) ..(92.2772,214.5292);
\path[fill=chfourfig5c0,fill opacity=0.5,draw=chfourfig5c0,line width=0.3985pt,draw opacity=0.5] (100.5662,203.4905) .. controls (100.5662,202.7697) and (99.9819,202.1854) ..(99.2611,202.1854) .. controls (98.5403,202.1854) and (97.9560,202.7697) ..(97.9560,203.4905) .. controls (97.9560,204.2113) and (98.5403,204.7956) ..(99.2611,204.7956) .. controls (99.9819,204.7956) and (100.5662,204.2113) ..(100.5662,203.4905);
\path[fill=chfourfig5c0,fill opacity=0.5,draw=chfourfig5c0,line width=0.3985pt,draw opacity=0.5] (137.2292,212.1282) .. controls (137.2292,211.4074) and (136.6449,210.8231) ..(135.9241,210.8231) .. controls (135.2033,210.8231) and (134.6190,211.4074) ..(134.6190,212.1282) .. controls (134.6190,212.8490) and (135.2033,213.4333) ..(135.9241,213.4333) .. controls (136.6449,213.4333) and (137.2292,212.8490) ..(137.2292,212.1282);
\path[fill=chfourfig5c0,fill opacity=0.5,draw=chfourfig5c0,line width=0.3985pt,draw opacity=0.5] (108.5264,226.2355) .. controls (108.5264,225.5146) and (107.9421,224.9304) ..(107.2213,224.9304) .. controls (106.5005,224.9304) and (105.9162,225.5146) ..(105.9162,226.2355) .. controls (105.9162,226.9563) and (106.5005,227.5406) ..(107.2213,227.5406) .. controls (107.9421,227.5406) and (108.5264,226.9563) ..(108.5264,226.2355);
\path[fill=chfourfig5c0,fill opacity=0.5,draw=chfourfig5c0,line width=0.3985pt,draw opacity=0.5] (54.7275,246.1909) .. controls (54.7275,245.4701) and (54.1432,244.8858) ..(53.4224,244.8858) .. controls (52.7016,244.8858) and (52.1172,245.4701) ..(52.1172,246.1909) .. controls (52.1172,246.9117) and (52.7016,247.4960) ..(53.4224,247.4960) .. controls (54.1432,247.4960) and (54.7275,246.9117) ..(54.7275,246.1909);
\path[fill=chfourfig5c0,fill opacity=0.5,draw=chfourfig5c0,line width=0.3985pt,draw opacity=0.5] (97.8165,249.2395) .. controls (97.8165,248.5187) and (97.2321,247.9344) ..(96.5114,247.9344) .. controls (95.7906,247.9344) and (95.2062,248.5187) ..(95.2062,249.2395) .. controls (95.2062,249.9603) and (95.7906,250.5446) ..(96.5114,250.5446) .. controls (97.2321,250.5446) and (97.8165,249.9603) ..(97.8165,249.2395);
\path[fill=chfourfig5c0,fill opacity=0.5,draw=chfourfig5c0,line width=0.3985pt,draw opacity=0.5] (163.7003,192.6410) .. controls (163.7003,191.9202) and (163.1159,191.3359) ..(162.3951,191.3359) .. controls (161.6743,191.3359) and (161.0900,191.9202) ..(161.0900,192.6410) .. controls (161.0900,193.3618) and (161.6743,193.9461) ..(162.3951,193.9461) .. controls (163.1159,193.9461) and (163.7003,193.3618) ..(163.7003,192.6410);
\path[fill=chfourfig5c0,fill opacity=0.5,draw=chfourfig5c0,line width=0.3985pt,draw opacity=0.5] (143.3762,223.4260) .. controls (143.3762,222.7051) and (142.7919,222.1208) ..(142.0711,222.1208) .. controls (141.3503,222.1208) and (140.7660,222.7051) ..(140.7660,223.4260) .. controls (140.7660,224.1468) and (141.3503,224.7311) ..(142.0711,224.7311) .. controls (142.7919,224.7311) and (143.3762,224.1468) ..(143.3762,223.4260);
\path[fill=chfourfig5c0,fill opacity=0.5,draw=chfourfig5c0,line width=0.3985pt,draw opacity=0.5] (144.7810,178.0655) .. controls (144.7810,177.3447) and (144.1967,176.7604) ..(143.4758,176.7604) .. controls (142.7551,176.7604) and (142.1707,177.3447) ..(142.1707,178.0655) .. controls (142.1707,178.7863) and (142.7551,179.3706) ..(143.4758,179.3706) .. controls (144.1967,179.3706) and (144.7810,178.7863) ..(144.7810,178.0655);
\path[fill=chfourfig5c0,fill opacity=0.5,draw=chfourfig5c0,line width=0.3985pt,draw opacity=0.5] (70.6878,223.6551) .. controls (70.6878,222.9343) and (70.1035,222.3500) ..(69.3827,222.3500) .. controls (68.6619,222.3500) and (68.0776,222.9343) ..(68.0776,223.6551) .. controls (68.0776,224.3759) and (68.6619,224.9602) ..(69.3827,224.9602) .. controls (70.1035,224.9602) and (70.6878,224.3759) ..(70.6878,223.6551);
\path[fill=chfourfig5c0,fill opacity=0.5,draw=chfourfig5c0,line width=0.3985pt,draw opacity=0.5] (116.2077,250.2258) .. controls (116.2077,249.5050) and (115.6234,248.9207) ..(114.9026,248.9207) .. controls (114.1818,248.9207) and (113.5975,249.5050) ..(113.5975,250.2258) .. controls (113.5975,250.9466) and (114.1818,251.5309) ..(114.9026,251.5309) .. controls (115.6234,251.5309) and (116.2077,250.9466) ..(116.2077,250.2258);
\path[fill=chfourfig5c0,fill opacity=0.5,draw=chfourfig5c0,line width=0.3985pt,draw opacity=0.5] (146.2953,234.5843) .. controls (146.2953,233.8634) and (145.7110,233.2791) ..(144.9902,233.2791) .. controls (144.2694,233.2791) and (143.6851,233.8634) ..(143.6851,234.5843) .. controls (143.6851,235.3051) and (144.2694,235.8894) ..(144.9902,235.8894) .. controls (145.7110,235.8894) and (146.2953,235.3051) ..(146.2953,234.5843);
\path[fill=chfourfig5c0,fill opacity=0.5,draw=chfourfig5c0,line width=0.3985pt,draw opacity=0.5] (70.3491,174.1999) .. controls (70.3491,173.4791) and (69.7648,172.8948) ..(69.0440,172.8948) .. controls (68.3232,172.8948) and (67.7389,173.4791) ..(67.7389,174.1999) .. controls (67.7389,174.9207) and (68.3232,175.5050) ..(69.0440,175.5050) .. controls (69.7648,175.5050) and (70.3491,174.9207) ..(70.3491,174.1999);
\path[fill=chfourfig5c0,fill opacity=0.5,draw=chfourfig5c0,line width=0.3985pt,draw opacity=0.5] (142.3301,213.7720) .. controls (142.3301,213.0512) and (141.7458,212.4669) ..(141.0250,212.4669) .. controls (140.3042,212.4669) and (139.7199,213.0512) ..(139.7199,213.7720) .. controls (139.7199,214.4929) and (140.3042,215.0771) ..(141.0250,215.0771) .. controls (141.7458,215.0771) and (142.3301,214.4929) ..(142.3301,213.7720);
\path[fill=chfourfig5c0,fill opacity=0.5,draw=chfourfig5c0,line width=0.3985pt,draw opacity=0.5] (144.8308,240.7811) .. controls (144.8308,240.0603) and (144.2465,239.4760) ..(143.5257,239.4760) .. controls (142.8049,239.4760) and (142.2205,240.0603) ..(142.2205,240.7811) .. controls (142.2205,241.5019) and (142.8049,242.0862) ..(143.5257,242.0862) .. controls (144.2465,242.0862) and (144.8308,241.5019) ..(144.8308,240.7811);
\path[fill=chfourfig5c0,fill opacity=0.5,draw=chfourfig5c0,line width=0.3985pt,draw opacity=0.5] (86.3493,206.7184) .. controls (86.3493,205.9976) and (85.7650,205.4133) ..(85.0442,205.4133) .. controls (84.3234,205.4133) and (83.7391,205.9976) ..(83.7391,206.7184) .. controls (83.7391,207.4392) and (84.3234,208.0235) ..(85.0442,208.0235) .. controls (85.7650,208.0235) and (86.3493,207.4392) ..(86.3493,206.7184);
\path[fill=chfourfig5c0,fill opacity=0.5,draw=chfourfig5c0,line width=0.3985pt,draw opacity=0.5] (65.9954,172.5660) .. controls (65.9954,171.8452) and (65.4111,171.2609) ..(64.6903,171.2609) .. controls (63.9695,171.2609) and (63.3852,171.8452) ..(63.3852,172.5660) .. controls (63.3852,173.2868) and (63.9695,173.8711) ..(64.6903,173.8711) .. controls (65.4111,173.8711) and (65.9954,173.2868) ..(65.9954,172.5660);
\path[fill=chfourfig5c0,fill opacity=0.5,draw=chfourfig5c0,line width=0.3985pt,draw opacity=0.5] (67.2607,239.1173) .. controls (67.2607,238.3965) and (66.6763,237.8122) ..(65.9555,237.8122) .. controls (65.2347,237.8122) and (64.6504,238.3965) ..(64.6504,239.1173) .. controls (64.6504,239.8381) and (65.2347,240.4224) ..(65.9555,240.4224) .. controls (66.6763,240.4224) and (67.2607,239.8381) ..(67.2607,239.1173);
\path[fill=chfourfig5c0,fill opacity=0.5,draw=chfourfig5c0,line width=0.3985pt,draw opacity=0.5] (57.0389,233.5083) .. controls (57.0389,232.7875) and (56.4545,232.2032) ..(55.7337,232.2032) .. controls (55.0129,232.2032) and (54.4286,232.7875) ..(54.4286,233.5083) .. controls (54.4286,234.2291) and (55.0129,234.8134) ..(55.7337,234.8134) .. controls (56.4545,234.8134) and (57.0389,234.2291) ..(57.0389,233.5083);
\path[fill=chfourfig5c0,fill opacity=0.5,draw=chfourfig5c0,line width=0.3985pt,draw opacity=0.5] (112.9400,248.8709) .. controls (112.9400,248.1501) and (112.3556,247.5658) ..(111.6349,247.5658) .. controls (110.9140,247.5658) and (110.3297,248.1501) ..(110.3297,248.8709) .. controls (110.3297,249.5917) and (110.9140,250.1760) ..(111.6349,250.1760) .. controls (112.3556,250.1760) and (112.9400,249.5917) ..(112.9400,248.8709);
\path[fill=chfourfig5c0,fill opacity=0.5,draw=chfourfig5c0,line width=0.3985pt,draw opacity=0.5] (153.1497,213.0946) .. controls (153.1497,212.3738) and (152.5654,211.7895) ..(151.8446,211.7895) .. controls (151.1238,211.7895) and (150.5394,212.3738) ..(150.5394,213.0946) .. controls (150.5394,213.8154) and (151.1238,214.3997) ..(151.8446,214.3997) .. controls (152.5654,214.3997) and (153.1497,213.8154) ..(153.1497,213.0946);
\path[fill=chfourfig5c0,fill opacity=0.5,draw=chfourfig5c0,line width=0.3985pt,draw opacity=0.5] (94.1502,211.3411) .. controls (94.1502,210.6203) and (93.5659,210.0360) ..(92.8451,210.0360) .. controls (92.1243,210.0360) and (91.5399,210.6203) ..(91.5399,211.3411) .. controls (91.5399,212.0620) and (92.1243,212.6462) ..(92.8451,212.6462) .. controls (93.5659,212.6462) and (94.1502,212.0620) ..(94.1502,211.3411);
\path[fill=chfourfig5c0,fill opacity=0.5,draw=chfourfig5c0,line width=0.3985pt,draw opacity=0.5] (160.9904,230.3600) .. controls (160.9904,229.6393) and (160.4061,229.0549) ..(159.6853,229.0549) .. controls (158.9645,229.0549) and (158.3801,229.6393) ..(158.3801,230.3600) .. controls (158.3801,231.0809) and (158.9645,231.6652) ..(159.6853,231.6652) .. controls (160.4061,231.6652) and (160.9904,231.0809) ..(160.9904,230.3600);
\path[fill=chfourfig5c0,fill opacity=0.5,draw=chfourfig5c0,line width=0.3985pt,draw opacity=0.5] (121.4880,210.0360) .. controls (121.4880,209.3152) and (120.9037,208.7309) ..(120.1829,208.7309) .. controls (119.4621,208.7309) and (118.8778,209.3152) ..(118.8778,210.0360) .. controls (118.8778,210.7568) and (119.4621,211.3411) ..(120.1829,211.3411) .. controls (120.9037,211.3411) and (121.4880,210.7568) ..(121.4880,210.0360);
\path[fill=chfourfig5c0,fill opacity=0.5,draw=chfourfig5c0,line width=0.3985pt,draw opacity=0.5] (126.0809,198.9175) .. controls (126.0809,198.1968) and (125.4965,197.6124) ..(124.7757,197.6124) .. controls (124.0549,197.6124) and (123.4706,198.1968) ..(123.4706,198.9175) .. controls (123.4706,199.6384) and (124.0549,200.2227) ..(124.7757,200.2227) .. controls (125.4965,200.2227) and (126.0809,199.6384) ..(126.0809,198.9175);
\path[fill=chfourfig5c0,fill opacity=0.5,draw=chfourfig5c0,line width=0.3985pt,draw opacity=0.5] (92.1377,233.5382) .. controls (92.1377,232.8174) and (91.5534,232.2331) ..(90.8326,232.2331) .. controls (90.1118,232.2331) and (89.5275,232.8174) ..(89.5275,233.5382) .. controls (89.5275,234.2590) and (90.1118,234.8433) ..(90.8326,234.8433) .. controls (91.5534,234.8433) and (92.1377,234.2590) ..(92.1377,233.5382);
\path[fill=chfourfig5c0,fill opacity=0.5,draw=chfourfig5c0,line width=0.3985pt,draw opacity=0.5] (151.5457,255.6954) .. controls (151.5457,254.9745) and (150.9614,254.3903) ..(150.2406,254.3903) .. controls (149.5198,254.3903) and (148.9354,254.9745) ..(148.9354,255.6954) .. controls (148.9354,256.4162) and (149.5198,257.0005) ..(150.2406,257.0005) .. controls (150.9614,257.0005) and (151.5457,256.4162) ..(151.5457,255.6954);
\path[fill=chfourfig5c0,fill opacity=0.5,draw=chfourfig5c0,line width=0.3985pt,draw opacity=0.5] (85.2036,220.1283) .. controls (85.2036,219.4075) and (84.6193,218.8232) ..(83.8985,218.8232) .. controls (83.1777,218.8232) and (82.5934,219.4075) ..(82.5934,220.1283) .. controls (82.5934,220.8491) and (83.1777,221.4334) ..(83.8985,221.4334) .. controls (84.6193,221.4334) and (85.2036,220.8491) ..(85.2036,220.1283);
\path[fill=chfourfig5c0,fill opacity=0.5,draw=chfourfig5c0,line width=0.3985pt,draw opacity=0.5] (94.9273,246.5894) .. controls (94.9273,245.8686) and (94.3430,245.2843) ..(93.6221,245.2843) .. controls (92.9014,245.2843) and (92.3170,245.8686) ..(92.3170,246.5894) .. controls (92.3170,247.3102) and (92.9014,247.8945) ..(93.6221,247.8945) .. controls (94.3430,247.8945) and (94.9273,247.3102) ..(94.9273,246.5894);
\path[fill=chfourfig5c0,fill opacity=0.5,draw=chfourfig5c0,line width=0.3985pt,draw opacity=0.5] (162.3852,202.8230) .. controls (162.3852,202.1022) and (161.8009,201.5178) ..(161.0801,201.5178) .. controls (160.3593,201.5178) and (159.7749,202.1022) ..(159.7749,202.8230) .. controls (159.7749,203.5438) and (160.3593,204.1281) ..(161.0801,204.1281) .. controls (161.8009,204.1281) and (162.3852,203.5438) ..(162.3852,202.8230);
\path[fill=chfourfig5c0,fill opacity=0.5,draw=chfourfig5c0,line width=0.3985pt,draw opacity=0.5] (84.6656,263.7453) .. controls (84.6656,263.0245) and (84.0813,262.4402) ..(83.3605,262.4402) .. controls (82.6397,262.4402) and (82.0554,263.0245) ..(82.0554,263.7453) .. controls (82.0554,264.4661) and (82.6397,265.0504) ..(83.3605,265.0504) .. controls (84.0813,265.0504) and (84.6656,264.4661) ..(84.6656,263.7453);
\path[fill=chfourfig5c0,fill opacity=0.5,draw=chfourfig5c0,line width=0.3985pt,draw opacity=0.5] (87.3954,197.3036) .. controls (87.3954,196.5828) and (86.8111,195.9984) ..(86.0903,195.9984) .. controls (85.3695,195.9984) and (84.7852,196.5828) ..(84.7852,197.3036) .. controls (84.7852,198.0244) and (85.3695,198.6087) ..(86.0903,198.6087) .. controls (86.8111,198.6087) and (87.3954,198.0244) ..(87.3954,197.3036);
\path[fill=chfourfig5c0,fill opacity=0.5,draw=chfourfig5c0,line width=0.3985pt,draw opacity=0.5] (134.1806,244.5172) .. controls (134.1806,243.7964) and (133.5963,243.2120) ..(132.8755,243.2120) .. controls (132.1546,243.2120) and (131.5703,243.7964) ..(131.5703,244.5172) .. controls (131.5703,245.2379) and (132.1546,245.8223) ..(132.8755,245.8223) .. controls (133.5963,245.8223) and (134.1806,245.2379) ..(134.1806,244.5172);
\path[fill=chfourfig5c0,fill opacity=0.5,draw=chfourfig5c0,line width=0.3985pt,draw opacity=0.5] (66.1149,228.2778) .. controls (66.1149,227.5570) and (65.5306,226.9727) ..(64.8098,226.9727) .. controls (64.0890,226.9727) and (63.5047,227.5570) ..(63.5047,228.2778) .. controls (63.5047,228.9987) and (64.0890,229.5829) ..(64.8098,229.5829) .. controls (65.5306,229.5829) and (66.1149,228.9987) ..(66.1149,228.2778);
\path[fill=chfourfig5c0,fill opacity=0.5,draw=chfourfig5c0,line width=0.3985pt,draw opacity=0.5] (122.8928,202.0558) .. controls (122.8928,201.3350) and (122.3085,200.7507) ..(121.5876,200.7507) .. controls (120.8669,200.7507) and (120.2825,201.3350) ..(120.2825,202.0558) .. controls (120.2825,202.7766) and (120.8669,203.3610) ..(121.5876,203.3610) .. controls (122.3085,203.3610) and (122.8928,202.7766) ..(122.8928,202.0558);
\path[fill=chfourfig5c0,fill opacity=0.5,draw=chfourfig5c0,line width=0.3985pt,draw opacity=0.5] (77.7913,245.7326) .. controls (77.7913,245.0118) and (77.2070,244.4275) ..(76.4862,244.4275) .. controls (75.7654,244.4275) and (75.1811,245.0118) ..(75.1811,245.7326) .. controls (75.1811,246.4534) and (75.7654,247.0377) ..(76.4862,247.0377) .. controls (77.2070,247.0377) and (77.7913,246.4534) ..(77.7913,245.7326);
\path[fill=chfourfig5c0,fill opacity=0.5,draw=chfourfig5c0,line width=0.3985pt,draw opacity=0.5] (149.6727,189.1740) .. controls (149.6727,188.4532) and (149.0884,187.8689) ..(148.3676,187.8689) .. controls (147.6468,187.8689) and (147.0624,188.4532) ..(147.0624,189.1740) .. controls (147.0624,189.8948) and (147.6468,190.4791) ..(148.3676,190.4791) .. controls (149.0884,190.4791) and (149.6727,189.8948) ..(149.6727,189.1740);
\path[fill=chfourfig5c0,fill opacity=0.5,draw=chfourfig5c0,line width=0.3985pt,draw opacity=0.5] (67.9879,189.7618) .. controls (67.9879,189.0410) and (67.4036,188.4566) ..(66.6828,188.4566) .. controls (65.9620,188.4566) and (65.3777,189.0410) ..(65.3777,189.7618) .. controls (65.3777,190.4826) and (65.9620,191.0669) ..(66.6828,191.0669) .. controls (67.4036,191.0669) and (67.9879,190.4826) ..(67.9879,189.7618);
\path[fill=chfourfig5c0,fill opacity=0.5,draw=chfourfig5c0,line width=0.3985pt,draw opacity=0.5] (102.3894,224.7510) .. controls (102.3894,224.0302) and (101.8051,223.4459) ..(101.0843,223.4459) .. controls (100.3635,223.4459) and (99.7791,224.0302) ..(99.7791,224.7510) .. controls (99.7791,225.4718) and (100.3635,226.0562) ..(101.0843,226.0562) .. controls (101.8051,226.0562) and (102.3894,225.4718) ..(102.3894,224.7510);
\path[fill=chfourfig5c0,fill opacity=0.5,draw=chfourfig5c0,line width=0.3985pt,draw opacity=0.5] (115.1915,188.5961) .. controls (115.1915,187.8753) and (114.6072,187.2910) ..(113.8864,187.2910) .. controls (113.1656,187.2910) and (112.5813,187.8753) ..(112.5813,188.5961) .. controls (112.5813,189.3169) and (113.1656,189.9012) ..(113.8864,189.9012) .. controls (114.6072,189.9012) and (115.1915,189.3169) ..(115.1915,188.5961);
\path[fill=chfourfig5c0,fill opacity=0.5,draw=chfourfig5c0,line width=0.3985pt,draw opacity=0.5] (61.9405,175.0368) .. controls (61.9405,174.3160) and (61.3562,173.7317) ..(60.6354,173.7317) .. controls (59.9146,173.7317) and (59.3303,174.3160) ..(59.3303,175.0368) .. controls (59.3303,175.7576) and (59.9146,176.3419) ..(60.6354,176.3419) .. controls (61.3562,176.3419) and (61.9405,175.7576) ..(61.9405,175.0368);
\path[fill=chfourfig5c0,fill opacity=0.5,draw=chfourfig5c0,line width=0.3985pt,draw opacity=0.5] (157.8521,244.7563) .. controls (157.8521,244.0355) and (157.2678,243.4511) ..(156.5470,243.4511) .. controls (155.8262,243.4511) and (155.2419,244.0355) ..(155.2419,244.7563) .. controls (155.2419,245.4771) and (155.8262,246.0614) ..(156.5470,246.0614) .. controls (157.2678,246.0614) and (157.8521,245.4771) ..(157.8521,244.7563);
\path[fill=chfourfig5c0,fill opacity=0.5,draw=chfourfig5c0,line width=0.3985pt,draw opacity=0.5] (58.9019,188.0781) .. controls (58.9019,187.3573) and (58.3176,186.7729) ..(57.5968,186.7729) .. controls (56.8760,186.7729) and (56.2917,187.3573) ..(56.2917,188.0781) .. controls (56.2917,188.7989) and (56.8760,189.3832) ..(57.5968,189.3832) .. controls (58.3176,189.3832) and (58.9019,188.7989) ..(58.9019,188.0781);
\path[fill=chfourfig5c0,fill opacity=0.5,draw=chfourfig5c0,line width=0.3985pt,draw opacity=0.5] (107.4604,215.3860) .. controls (107.4604,214.6652) and (106.8761,214.0809) ..(106.1553,214.0809) .. controls (105.4345,214.0809) and (104.8502,214.6652) ..(104.8502,215.3860) .. controls (104.8502,216.1068) and (105.4345,216.6911) ..(106.1553,216.6911) .. controls (106.8761,216.6911) and (107.4604,216.1068) ..(107.4604,215.3860);
\path[fill=chfourfig5c0,fill opacity=0.5,draw=chfourfig5c0,line width=0.3985pt,draw opacity=0.5] (81.0193,207.8741) .. controls (81.0193,207.1533) and (80.4349,206.5690) ..(79.7141,206.5690) .. controls (78.9933,206.5690) and (78.4090,207.1533) ..(78.4090,207.8741) .. controls (78.4090,208.5949) and (78.9933,209.1792) ..(79.7141,209.1792) .. controls (80.4349,209.1792) and (81.0193,208.5949) ..(81.0193,207.8741);
\path[fill=chfourfig5c0,fill opacity=0.5,draw=chfourfig5c0,line width=0.3985pt,draw opacity=0.5] (96.8302,180.3270) .. controls (96.8302,179.6062) and (96.2458,179.0219) ..(95.5250,179.0219) .. controls (94.8042,179.0219) and (94.2199,179.6062) ..(94.2199,180.3270) .. controls (94.2199,181.0478) and (94.8042,181.6321) ..(95.5250,181.6321) .. controls (96.2458,181.6321) and (96.8302,181.0478) ..(96.8302,180.3270);
\path[fill=chfourfig5c0,fill opacity=0.5,draw=chfourfig5c0,line width=0.3985pt,draw opacity=0.5] (72.5210,205.2439) .. controls (72.5210,204.5231) and (71.9367,203.9388) ..(71.2159,203.9388) .. controls (70.4951,203.9388) and (69.9107,204.5231) ..(69.9107,205.2439) .. controls (69.9107,205.9647) and (70.4951,206.5490) ..(71.2159,206.5490) .. controls (71.9367,206.5490) and (72.5210,205.9647) ..(72.5210,205.2439);
\path[fill=chfourfig5c0,fill opacity=0.5,draw=chfourfig5c0,line width=0.3985pt,draw opacity=0.5] (91.3108,227.3513) .. controls (91.3108,226.6305) and (90.7265,226.0462) ..(90.0056,226.0462) .. controls (89.2849,226.0462) and (88.7005,226.6305) ..(88.7005,227.3513) .. controls (88.7005,228.0721) and (89.2849,228.6564) ..(90.0056,228.6564) .. controls (90.7265,228.6564) and (91.3108,228.0721) ..(91.3108,227.3513);
\path[fill=chfourfig5c0,fill opacity=0.5,draw=chfourfig5c0,line width=0.3985pt,draw opacity=0.5] (164.2582,202.6835) .. controls (164.2582,201.9627) and (163.6739,201.3784) ..(162.9530,201.3784) .. controls (162.2323,201.3784) and (161.6479,201.9627) ..(161.6479,202.6835) .. controls (161.6479,203.4043) and (162.2323,203.9886) ..(162.9530,203.9886) .. controls (163.6739,203.9886) and (164.2582,203.4043) ..(164.2582,202.6835);
\path[fill=chfourfig5c0,fill opacity=0.5,draw=chfourfig5c0,line width=0.3985pt,draw opacity=0.5] (90.6632,194.8627) .. controls (90.6632,194.1419) and (90.0789,193.5576) ..(89.3581,193.5576) .. controls (88.6373,193.5576) and (88.0530,194.1419) ..(88.0530,194.8627) .. controls (88.0530,195.5835) and (88.6373,196.1678) ..(89.3581,196.1678) .. controls (90.0789,196.1678) and (90.6632,195.5835) ..(90.6632,194.8627);
\path[fill=chfourfig5c0,fill opacity=0.5,draw=chfourfig5c0,line width=0.3985pt,draw opacity=0.5] (119.1866,236.7462) .. controls (119.1866,236.0254) and (118.6023,235.4411) ..(117.8815,235.4411) .. controls (117.1607,235.4411) and (116.5764,236.0254) ..(116.5764,236.7462) .. controls (116.5764,237.4670) and (117.1607,238.0513) ..(117.8815,238.0513) .. controls (118.6023,238.0513) and (119.1866,237.4670) ..(119.1866,236.7462);
\path[fill=chfourfig5c0,fill opacity=0.5,draw=chfourfig5c0,line width=0.3985pt,draw opacity=0.5] (148.3875,185.5376) .. controls (148.3875,184.8168) and (147.8032,184.2324) ..(147.0824,184.2324) .. controls (146.3616,184.2324) and (145.7773,184.8168) ..(145.7773,185.5376) .. controls (145.7773,186.2583) and (146.3616,186.8427) ..(147.0824,186.8427) .. controls (147.8032,186.8427) and (148.3875,186.2583) ..(148.3875,185.5376);
\path[fill=chfourfig5c0,fill opacity=0.5,draw=chfourfig5c0,line width=0.3985pt,draw opacity=0.5] (88.9098,253.7227) .. controls (88.9098,253.0019) and (88.3254,252.4176) ..(87.6046,252.4176) .. controls (86.8838,252.4176) and (86.2995,253.0019) ..(86.2995,253.7227) .. controls (86.2995,254.4435) and (86.8838,255.0279) ..(87.6046,255.0279) .. controls (88.3254,255.0279) and (88.9098,254.4435) ..(88.9098,253.7227);
\path[fill=chfourfig5c0,fill opacity=0.5,draw=chfourfig5c0,line width=0.3985pt,draw opacity=0.5] (101.8614,206.0509) .. controls (101.8614,205.3301) and (101.2770,204.7458) ..(100.5562,204.7458) .. controls (99.8354,204.7458) and (99.2511,205.3301) ..(99.2511,206.0509) .. controls (99.2511,206.7717) and (99.8354,207.3560) ..(100.5562,207.3560) .. controls (101.2770,207.3560) and (101.8614,206.7717) ..(101.8614,206.0509);
\path[fill=chfourfig5c0,fill opacity=0.5,draw=chfourfig5c0,line width=0.3985pt,draw opacity=0.5] (65.7961,220.4571) .. controls (65.7961,219.7363) and (65.2118,219.1519) ..(64.4910,219.1519) .. controls (63.7702,219.1519) and (63.1859,219.7363) ..(63.1859,220.4571) .. controls (63.1859,221.1779) and (63.7702,221.7622) ..(64.4910,221.7622) .. controls (65.2118,221.7622) and (65.7961,221.1779) ..(65.7961,220.4571);
\path[fill=chfourfig5c0,fill opacity=0.5,draw=chfourfig5c0,line width=0.3985pt,draw opacity=0.5] (46.8270,194.2849) .. controls (46.8270,193.5641) and (46.2427,192.9797) ..(45.5219,192.9797) .. controls (44.8011,192.9797) and (44.2168,193.5641) ..(44.2168,194.2849) .. controls (44.2168,195.0057) and (44.8011,195.5900) ..(45.5219,195.5900) .. controls (46.2427,195.5900) and (46.8270,195.0057) ..(46.8270,194.2849);
\path[fill=chfourfig5c0,fill opacity=0.5,draw=chfourfig5c0,line width=0.3985pt,draw opacity=0.5] (86.1003,217.4583) .. controls (86.1003,216.7375) and (85.5159,216.1531) ..(84.7951,216.1531) .. controls (84.0743,216.1531) and (83.4900,216.7375) ..(83.4900,217.4583) .. controls (83.4900,218.1791) and (84.0743,218.7634) ..(84.7951,218.7634) .. controls (85.5159,218.7634) and (86.1003,218.1791) ..(86.1003,217.4583);
\path[fill=chfourfig5c0,fill opacity=0.5,draw=chfourfig5c0,line width=0.3985pt,draw opacity=0.5] (139.4310,201.8665) .. controls (139.4310,201.1457) and (138.8466,200.5614) ..(138.1258,200.5614) .. controls (137.4050,200.5614) and (136.8207,201.1457) ..(136.8207,201.8665) .. controls (136.8207,202.5873) and (137.4050,203.1716) ..(138.1258,203.1716) .. controls (138.8466,203.1716) and (139.4310,202.5873) ..(139.4310,201.8665);
\path[fill=chfourfig5c0,fill opacity=0.5,draw=chfourfig5c0,line width=0.3985pt,draw opacity=0.5] (59.3403,253.0253) .. controls (59.3403,252.3046) and (58.7559,251.7202) ..(58.0351,251.7202) .. controls (57.3143,251.7202) and (56.7300,252.3046) ..(56.7300,253.0253) .. controls (56.7300,253.7462) and (57.3143,254.3305) ..(58.0351,254.3305) .. controls (58.7559,254.3305) and (59.3403,253.7462) ..(59.3403,253.0253);
\path[fill=chfourfig5c0,fill opacity=0.5,draw=chfourfig5c0,line width=0.3985pt,draw opacity=0.5] (68.0378,195.9885) .. controls (68.0378,195.2677) and (67.4534,194.6834) ..(66.7326,194.6834) .. controls (66.0118,194.6834) and (65.4275,195.2677) ..(65.4275,195.9885) .. controls (65.4275,196.7093) and (66.0118,197.2936) ..(66.7326,197.2936) .. controls (67.4534,197.2936) and (68.0378,196.7093) ..(68.0378,195.9885);
\path[fill=chfourfig5c0,fill opacity=0.5,draw=chfourfig5c0,line width=0.3985pt,draw opacity=0.5] (96.2324,249.5284) .. controls (96.2324,248.8076) and (95.6481,248.2233) ..(94.9273,248.2233) .. controls (94.2065,248.2233) and (93.6221,248.8076) ..(93.6221,249.5284) .. controls (93.6221,250.2492) and (94.2065,250.8335) ..(94.9273,250.8335) .. controls (95.6481,250.8335) and (96.2324,250.2492) ..(96.2324,249.5284);
\path[fill=chfourfig5c0,fill opacity=0.5,draw=chfourfig5c0,line width=0.3985pt,draw opacity=0.5] (82.2247,229.0250) .. controls (82.2247,228.3043) and (81.6404,227.7199) ..(80.9196,227.7199) .. controls (80.1988,227.7199) and (79.6145,228.3043) ..(79.6145,229.0250) .. controls (79.6145,229.7458) and (80.1988,230.3302) ..(80.9196,230.3302) .. controls (81.6404,230.3302) and (82.2247,229.7458) ..(82.2247,229.0250);
\path[fill=chfourfig5c0,fill opacity=0.5,draw=chfourfig5c0,line width=0.3985pt,draw opacity=0.5] (53.6216,212.1979) .. controls (53.6216,211.4771) and (53.0373,210.8928) ..(52.3165,210.8928) .. controls (51.5957,210.8928) and (51.0114,211.4771) ..(51.0114,212.1979) .. controls (51.0114,212.9187) and (51.5957,213.5031) ..(52.3165,213.5031) .. controls (53.0373,213.5031) and (53.6216,212.9187) ..(53.6216,212.1979);
\path[fill=chfourfig5c0,fill opacity=0.5,draw=chfourfig5c0,line width=0.3985pt,draw opacity=0.5] (109.3235,183.5151) .. controls (109.3235,182.7943) and (108.7391,182.2100) ..(108.0183,182.2100) .. controls (107.2975,182.2100) and (106.7132,182.7943) ..(106.7132,183.5151) .. controls (106.7132,184.2359) and (107.2975,184.8202) ..(108.0183,184.8202) .. controls (108.7391,184.8202) and (109.3235,184.2359) ..(109.3235,183.5151);
\path[fill=chfourfig5c0,fill opacity=0.5,draw=chfourfig5c0,line width=0.3985pt,draw opacity=0.5] (50.8819,253.4239) .. controls (50.8819,252.7030) and (50.2976,252.1187) ..(49.5768,252.1187) .. controls (48.8560,252.1187) and (48.2716,252.7030) ..(48.2716,253.4239) .. controls (48.2716,254.1447) and (48.8560,254.7290) ..(49.5768,254.7290) .. controls (50.2976,254.7290) and (50.8819,254.1447) ..(50.8819,253.4239);
\path[fill=chfourfig5c0,fill opacity=0.5,draw=chfourfig5c0,line width=0.3985pt,draw opacity=0.5] (99.1615,222.8581) .. controls (99.1615,222.1373) and (98.5771,221.5530) ..(97.8563,221.5530) .. controls (97.1355,221.5530) and (96.5512,222.1373) ..(96.5512,222.8581) .. controls (96.5512,223.5789) and (97.1355,224.1632) ..(97.8563,224.1632) .. controls (98.5771,224.1632) and (99.1615,223.5789) ..(99.1615,222.8581);
\path[fill=chfourfig5c0,fill opacity=0.5,draw=chfourfig5c0,line width=0.3985pt,draw opacity=0.5] (53.2331,230.3401) .. controls (53.2331,229.6193) and (52.6488,229.0350) ..(51.9280,229.0350) .. controls (51.2072,229.0350) and (50.6228,229.6193) ..(50.6228,230.3401) .. controls (50.6228,231.0609) and (51.2072,231.6453) ..(51.9280,231.6453) .. controls (52.6488,231.6453) and (53.2331,231.0609) ..(53.2331,230.3401);
\path[fill=chfourfig5c0,fill opacity=0.5,draw=chfourfig5c0,line width=0.3985pt,draw opacity=0.5] (100.0780,216.0037) .. controls (100.0780,215.2829) and (99.4937,214.6986) ..(98.7729,214.6986) .. controls (98.0521,214.6986) and (97.4678,215.2829) ..(97.4678,216.0037) .. controls (97.4678,216.7245) and (98.0521,217.3088) ..(98.7729,217.3088) .. controls (99.4937,217.3088) and (100.0780,216.7245) ..(100.0780,216.0037);
\path[fill=chfourfig5c0,fill opacity=0.5,draw=chfourfig5c0,line width=0.3985pt,draw opacity=0.5] (57.6864,207.3261) .. controls (57.6864,206.6053) and (57.1021,206.0210) ..(56.3813,206.0210) .. controls (55.6605,206.0210) and (55.0762,206.6053) ..(55.0762,207.3261) .. controls (55.0762,208.0469) and (55.6605,208.6313) ..(56.3813,208.6313) .. controls (57.1021,208.6313) and (57.6864,208.0469) ..(57.6864,207.3261);
\path[fill=chfourfig5c0,fill opacity=0.5,draw=chfourfig5c0,line width=0.3985pt,draw opacity=0.5] (56.5407,231.5456) .. controls (56.5407,230.8248) and (55.9564,230.2405) ..(55.2356,230.2405) .. controls (54.5148,230.2405) and (53.9305,230.8248) ..(53.9305,231.5456) .. controls (53.9305,232.2664) and (54.5148,232.8507) ..(55.2356,232.8507) .. controls (55.9564,232.8507) and (56.5407,232.2664) ..(56.5407,231.5456);
\path[fill=chfourfig5c0,fill opacity=0.5,draw=chfourfig5c0,line width=0.3985pt,draw opacity=0.5] (156.6167,263.0280) .. controls (156.6167,262.3072) and (156.0324,261.7228) ..(155.3116,261.7228) .. controls (154.5908,261.7228) and (154.0065,262.3072) ..(154.0065,263.0280) .. controls (154.0065,263.7488) and (154.5908,264.3331) ..(155.3116,264.3331) .. controls (156.0324,264.3331) and (156.6167,263.7488) ..(156.6167,263.0280);
\path[fill=chfourfig5c0,fill opacity=0.5,draw=chfourfig5c0,line width=0.3985pt,draw opacity=0.5] (156.0788,221.1843) .. controls (156.0788,220.4635) and (155.4944,219.8792) ..(154.7736,219.8792) .. controls (154.0528,219.8792) and (153.4685,220.4635) ..(153.4685,221.1843) .. controls (153.4685,221.9052) and (154.0528,222.4894) ..(154.7736,222.4894) .. controls (155.4944,222.4894) and (156.0788,221.9052) ..(156.0788,221.1843);
\path[fill=chfourfig5c0,fill opacity=0.5,draw=chfourfig5c0,line width=0.3985pt,draw opacity=0.5] (72.5110,182.8575) .. controls (72.5110,182.1368) and (71.9267,181.5524) ..(71.2059,181.5524) .. controls (70.4851,181.5524) and (69.9008,182.1368) ..(69.9008,182.8575) .. controls (69.9008,183.5784) and (70.4851,184.1627) ..(71.2059,184.1627) .. controls (71.9267,184.1627) and (72.5110,183.5784) ..(72.5110,182.8575);
\path[fill=chfourfig5c0,fill opacity=0.5,draw=chfourfig5c0,line width=0.3985pt,draw opacity=0.5] (120.3523,181.7019) .. controls (120.3523,180.9811) and (119.7679,180.3968) ..(119.0471,180.3968) .. controls (118.3263,180.3968) and (117.7420,180.9811) ..(117.7420,181.7019) .. controls (117.7420,182.4227) and (118.3263,183.0070) ..(119.0471,183.0070) .. controls (119.7679,183.0070) and (120.3523,182.4227) ..(120.3523,181.7019);
\path[fill=chfourfig5c0,fill opacity=0.5,draw=chfourfig5c0,line width=0.3985pt,draw opacity=0.5] (157.2145,241.9866) .. controls (157.2145,241.2658) and (156.6302,240.6815) ..(155.9094,240.6815) .. controls (155.1886,240.6815) and (154.6043,241.2658) ..(154.6043,241.9866) .. controls (154.6043,242.7074) and (155.1886,243.2917) ..(155.9094,243.2917) .. controls (156.6302,243.2917) and (157.2145,242.7074) ..(157.2145,241.9866);
\path[fill=chfourfig5c0,fill opacity=0.5,draw=chfourfig5c0,line width=0.3985pt,draw opacity=0.5] (55.1161,232.4323) .. controls (55.1161,231.7115) and (54.5317,231.1272) ..(53.8109,231.1272) .. controls (53.0901,231.1272) and (52.5058,231.7115) ..(52.5058,232.4323) .. controls (52.5058,233.1531) and (53.0901,233.7374) ..(53.8109,233.7374) .. controls (54.5317,233.7374) and (55.1161,233.1531) ..(55.1161,232.4323);
\path[fill=chfourfig5c0,fill opacity=0.5,draw=chfourfig5c0,line width=0.3985pt,draw opacity=0.5] (132.2279,212.3374) .. controls (132.2279,211.6166) and (131.6436,211.0323) ..(130.9228,211.0323) .. controls (130.2020,211.0323) and (129.6176,211.6166) ..(129.6176,212.3374) .. controls (129.6176,213.0582) and (130.2020,213.6425) ..(130.9228,213.6425) .. controls (131.6436,213.6425) and (132.2279,213.0582) ..(132.2279,212.3374);
\path[fill=chfourfig5c0,fill opacity=0.5,draw=chfourfig5c0,line width=0.3985pt,draw opacity=0.5] (79.6444,230.8881) .. controls (79.6444,230.1673) and (79.0601,229.5829) ..(78.3393,229.5829) .. controls (77.6185,229.5829) and (77.0341,230.1673) ..(77.0341,230.8881) .. controls (77.0341,231.6089) and (77.6185,232.1932) ..(78.3393,232.1932) .. controls (79.0601,232.1932) and (79.6444,231.6089) ..(79.6444,230.8881);
\path[fill=chfourfig5c0,fill opacity=0.5,draw=chfourfig5c0,line width=0.3985pt,draw opacity=0.5] (109.4131,221.2441) .. controls (109.4131,220.5233) and (108.8288,219.9390) ..(108.1080,219.9390) .. controls (107.3872,219.9390) and (106.8029,220.5233) ..(106.8029,221.2441) .. controls (106.8029,221.9649) and (107.3872,222.5492) ..(108.1080,222.5492) .. controls (108.8288,222.5492) and (109.4131,221.9649) ..(109.4131,221.2441);
\path[fill=chfourfig5c0,fill opacity=0.5,draw=chfourfig5c0,line width=0.3985pt,draw opacity=0.5] (117.6324,217.6675) .. controls (117.6324,216.9467) and (117.0481,216.3624) ..(116.3273,216.3624) .. controls (115.6065,216.3624) and (115.0222,216.9467) ..(115.0222,217.6675) .. controls (115.0222,218.3883) and (115.6065,218.9726) ..(116.3273,218.9726) .. controls (117.0481,218.9726) and (117.6324,218.3883) ..(117.6324,217.6675);
\path[fill=chfourfig5c0,fill opacity=0.5,draw=chfourfig5c0,line width=0.3985pt,draw opacity=0.5] (85.7217,194.7830) .. controls (85.7217,194.0622) and (85.1373,193.4779) ..(84.4165,193.4779) .. controls (83.6957,193.4779) and (83.1114,194.0622) ..(83.1114,194.7830) .. controls (83.1114,195.5038) and (83.6957,196.0881) ..(84.4165,196.0881) .. controls (85.1373,196.0881) and (85.7217,195.5038) ..(85.7217,194.7830);
\path[fill=chfourfig5c0,fill opacity=0.5,draw=chfourfig5c0,line width=0.3985pt,draw opacity=0.5] (159.8247,205.3934) .. controls (159.8247,204.6725) and (159.2404,204.0883) ..(158.5196,204.0883) .. controls (157.7988,204.0883) and (157.2145,204.6725) ..(157.2145,205.3934) .. controls (157.2145,206.1142) and (157.7988,206.6985) ..(158.5196,206.6985) .. controls (159.2404,206.6985) and (159.8247,206.1142) ..(159.8247,205.3934);
\path[fill=chfourfig5c0,fill opacity=0.5,draw=chfourfig5c0,line width=0.3985pt,draw opacity=0.5] (84.6258,198.1803) .. controls (84.6258,197.4595) and (84.0415,196.8752) ..(83.3206,196.8752) .. controls (82.5999,196.8752) and (82.0155,197.4595) ..(82.0155,198.1803) .. controls (82.0155,198.9011) and (82.5999,199.4854) ..(83.3206,199.4854) .. controls (84.0415,199.4854) and (84.6258,198.9011) ..(84.6258,198.1803);
\path[fill=chfourfig5c0,fill opacity=0.5,draw=chfourfig5c0,line width=0.3985pt,draw opacity=0.5] (162.3453,220.4869) .. controls (162.3453,219.7661) and (161.7610,219.1818) ..(161.0402,219.1818) .. controls (160.3194,219.1818) and (159.7351,219.7661) ..(159.7351,220.4869) .. controls (159.7351,221.2078) and (160.3194,221.7921) ..(161.0402,221.7921) .. controls (161.7610,221.7921) and (162.3453,221.2078) ..(162.3453,220.4869);
\path[fill=chfourfig5c0,fill opacity=0.5,draw=chfourfig5c0,line width=0.3985pt,draw opacity=0.5] (57.0488,240.6217) .. controls (57.0488,239.9009) and (56.4645,239.3166) ..(55.7437,239.3166) .. controls (55.0229,239.3166) and (54.4386,239.9009) ..(54.4386,240.6217) .. controls (54.4386,241.3425) and (55.0229,241.9268) ..(55.7437,241.9268) .. controls (56.4645,241.9268) and (57.0488,241.3425) ..(57.0488,240.6217);
\path[fill=chfourfig5c0,fill opacity=0.5,draw=chfourfig5c0,line width=0.3985pt,draw opacity=0.5] (48.0126,259.8100) .. controls (48.0126,259.0892) and (47.4283,258.5049) ..(46.7075,258.5049) .. controls (45.9867,258.5049) and (45.4024,259.0892) ..(45.4024,259.8100) .. controls (45.4024,260.5308) and (45.9867,261.1151) ..(46.7075,261.1151) .. controls (47.4283,261.1151) and (48.0126,260.5308) ..(48.0126,259.8100);
\path[fill=chfourfig5c0,fill opacity=0.5,draw=chfourfig5c0,line width=0.3985pt,draw opacity=0.5] (113.8366,231.3763) .. controls (113.8366,230.6555) and (113.2523,230.0711) ..(112.5315,230.0711) .. controls (111.8107,230.0711) and (111.2264,230.6555) ..(111.2264,231.3763) .. controls (111.2264,232.0970) and (111.8107,232.6814) ..(112.5315,232.6814) .. controls (113.2523,232.6814) and (113.8366,232.0970) ..(113.8366,231.3763);
\path[fill=chfourfig5c0,fill opacity=0.5,draw=chfourfig5c0,line width=0.3985pt,draw opacity=0.5] (66.3939,212.0883) .. controls (66.3939,211.3676) and (65.8096,210.7832) ..(65.0888,210.7832) .. controls (64.3680,210.7832) and (63.7837,211.3676) ..(63.7837,212.0883) .. controls (63.7837,212.8091) and (64.3680,213.3935) ..(65.0888,213.3935) .. controls (65.8096,213.3935) and (66.3939,212.8091) ..(66.3939,212.0883);
\path[fill=chfourfig5c0,fill opacity=0.5,draw=chfourfig5c0,line width=0.3985pt,draw opacity=0.5] (125.9912,223.8643) .. controls (125.9912,223.1435) and (125.4069,222.5592) ..(124.6861,222.5592) .. controls (123.9653,222.5592) and (123.3810,223.1435) ..(123.3810,223.8643) .. controls (123.3810,224.5851) and (123.9653,225.1694) ..(124.6861,225.1694) .. controls (125.4069,225.1694) and (125.9912,224.5851) ..(125.9912,223.8643);
\path[fill=chfourfig5c0,fill opacity=0.5,draw=chfourfig5c0,line width=0.3985pt,draw opacity=0.5] (100.0979,223.3762) .. controls (100.0979,222.6553) and (99.5136,222.0710) ..(98.7928,222.0710) .. controls (98.0720,222.0710) and (97.4877,222.6553) ..(97.4877,223.3762) .. controls (97.4877,224.0969) and (98.0720,224.6812) ..(98.7928,224.6812) .. controls (99.5136,224.6812) and (100.0979,224.0969) ..(100.0979,223.3762);
\path[fill=chfourfig5c0,fill opacity=0.5,draw=chfourfig5c0,line width=0.3985pt,draw opacity=0.5] (106.6535,246.0116) .. controls (106.6535,245.2908) and (106.0691,244.7064) ..(105.3483,244.7064) .. controls (104.6275,244.7064) and (104.0432,245.2908) ..(104.0432,246.0116) .. controls (104.0432,246.7324) and (104.6275,247.3167) ..(105.3483,247.3167) .. controls (106.0691,247.3167) and (106.6535,246.7324) ..(106.6535,246.0116);
\path[fill=chfourfig5c0,fill opacity=0.5,draw=chfourfig5c0,line width=0.3985pt,draw opacity=0.5] (163.5309,199.6249) .. controls (163.5309,198.9041) and (162.9466,198.3198) ..(162.2258,198.3198) .. controls (161.5050,198.3198) and (160.9207,198.9041) ..(160.9207,199.6249) .. controls (160.9207,200.3457) and (161.5050,200.9300) ..(162.2258,200.9300) .. controls (162.9466,200.9300) and (163.5309,200.3457) ..(163.5309,199.6249);
\path[fill=chfourfig5c0,fill opacity=0.5,draw=chfourfig5c0,line width=0.3985pt,draw opacity=0.5] (136.7510,229.9914) .. controls (136.7510,229.2706) and (136.1667,228.6863) ..(135.4458,228.6863) .. controls (134.7251,228.6863) and (134.1407,229.2706) ..(134.1407,229.9914) .. controls (134.1407,230.7122) and (134.7251,231.2965) ..(135.4458,231.2965) .. controls (136.1667,231.2965) and (136.7510,230.7122) ..(136.7510,229.9914);
\path[fill=chfourfig5c0,fill opacity=0.5,draw=chfourfig5c0,line width=0.3985pt,draw opacity=0.5] (148.4273,222.5094) .. controls (148.4273,221.7886) and (147.8430,221.2043) ..(147.1222,221.2043) .. controls (146.4014,221.2043) and (145.8171,221.7886) ..(145.8171,222.5094) .. controls (145.8171,223.2302) and (146.4014,223.8145) ..(147.1222,223.8145) .. controls (147.8430,223.8145) and (148.4273,223.2302) ..(148.4273,222.5094);
\path[fill=chfourfig5c0,fill opacity=0.5,draw=chfourfig5c0,line width=0.3985pt,draw opacity=0.5] (127.8044,189.1540) .. controls (127.8044,188.4332) and (127.2201,187.8489) ..(126.4993,187.8489) .. controls (125.7785,187.8489) and (125.1942,188.4332) ..(125.1942,189.1540) .. controls (125.1942,189.8748) and (125.7785,190.4592) ..(126.4993,190.4592) .. controls (127.2201,190.4592) and (127.8044,189.8748) ..(127.8044,189.1540);
\path[fill=chfourfig5c0,fill opacity=0.5,draw=chfourfig5c0,line width=0.3985pt,draw opacity=0.5] (86.5585,181.8414) .. controls (86.5585,181.1206) and (85.9742,180.5363) ..(85.2534,180.5363) .. controls (84.5326,180.5363) and (83.9483,181.1206) ..(83.9483,181.8414) .. controls (83.9483,182.5622) and (84.5326,183.1465) ..(85.2534,183.1465) .. controls (85.9742,183.1465) and (86.5585,182.5622) ..(86.5585,181.8414);
\path[fill=chfourfig5c0,fill opacity=0.5,draw=chfourfig5c0,line width=0.3985pt,draw opacity=0.5] (165.7924,210.0858) .. controls (165.7924,209.3650) and (165.2081,208.7807) ..(164.4873,208.7807) .. controls (163.7665,208.7807) and (163.1822,209.3650) ..(163.1822,210.0858) .. controls (163.1822,210.8066) and (163.7665,211.3909) ..(164.4873,211.3909) .. controls (165.2081,211.3909) and (165.7924,210.8066) ..(165.7924,210.0858);
\path[fill=chfourfig5c0,fill opacity=0.5,draw=chfourfig5c0,line width=0.3985pt,draw opacity=0.5] (129.8866,210.9227) .. controls (129.8866,210.2019) and (129.3023,209.6176) ..(128.5815,209.6176) .. controls (127.8607,209.6176) and (127.2764,210.2019) ..(127.2764,210.9227) .. controls (127.2764,211.6435) and (127.8607,212.2278) ..(128.5815,212.2278) .. controls (129.3023,212.2278) and (129.8866,211.6435) ..(129.8866,210.9227);
\path[fill=chfourfig5c0,fill opacity=0.5,draw=chfourfig5c0,line width=0.3985pt,draw opacity=0.5] (81.2882,187.9585) .. controls (81.2882,187.2377) and (80.7039,186.6534) ..(79.9831,186.6534) .. controls (79.2623,186.6534) and (78.6780,187.2377) ..(78.6780,187.9585) .. controls (78.6780,188.6793) and (79.2623,189.2636) ..(79.9831,189.2636) .. controls (80.7039,189.2636) and (81.2882,188.6793) ..(81.2882,187.9585);
\path[fill=chfourfig5c0,fill opacity=0.5,draw=chfourfig5c0,line width=0.3985pt,draw opacity=0.5] (75.2110,191.5750) .. controls (75.2110,190.8542) and (74.6266,190.2699) ..(73.9058,190.2699) .. controls (73.1850,190.2699) and (72.6007,190.8542) ..(72.6007,191.5750) .. controls (72.6007,192.2958) and (73.1850,192.8801) ..(73.9058,192.8801) .. controls (74.6266,192.8801) and (75.2110,192.2958) ..(75.2110,191.5750);
\path[fill=chfourfig5c0,fill opacity=0.5,draw=chfourfig5c0,line width=0.3985pt,draw opacity=0.5] (77.9607,253.4836) .. controls (77.9607,252.7628) and (77.3764,252.1785) ..(76.6556,252.1785) .. controls (75.9348,252.1785) and (75.3504,252.7628) ..(75.3504,253.4836) .. controls (75.3504,254.2044) and (75.9348,254.7888) ..(76.6556,254.7888) .. controls (77.3764,254.7888) and (77.9607,254.2044) ..(77.9607,253.4836);
\path[fill=chfourfig5c0,fill opacity=0.5,draw=chfourfig5c0,line width=0.3985pt,draw opacity=0.5] (74.4139,238.2207) .. controls (74.4139,237.4999) and (73.8296,236.9156) ..(73.1088,236.9156) .. controls (72.3880,236.9156) and (71.8037,237.4999) ..(71.8037,238.2207) .. controls (71.8037,238.9415) and (72.3880,239.5258) ..(73.1088,239.5258) .. controls (73.8296,239.5258) and (74.4139,238.9415) ..(74.4139,238.2207);
\path[fill=chfourfig5c0,fill opacity=0.5,draw=chfourfig5c0,line width=0.3985pt,draw opacity=0.5] (67.5994,179.5101) .. controls (67.5994,178.7893) and (67.0151,178.2050) ..(66.2943,178.2050) .. controls (65.5735,178.2050) and (64.9892,178.7893) ..(64.9892,179.5101) .. controls (64.9892,180.2309) and (65.5735,180.8152) ..(66.2943,180.8152) .. controls (67.0151,180.8152) and (67.5994,180.2309) ..(67.5994,179.5101);
\path[fill=chfourfig5c0,fill opacity=0.5,draw=chfourfig5c0,line width=0.3985pt,draw opacity=0.5] (143.9939,213.1643) .. controls (143.9939,212.4435) and (143.4096,211.8592) ..(142.6888,211.8592) .. controls (141.9680,211.8592) and (141.3837,212.4435) ..(141.3837,213.1643) .. controls (141.3837,213.8851) and (141.9680,214.4694) ..(142.6888,214.4694) .. controls (143.4096,214.4694) and (143.9939,213.8851) ..(143.9939,213.1643);
\path[fill=chfourfig5c0,fill opacity=0.5,draw=chfourfig5c0,line width=0.3985pt,draw opacity=0.5] (47.7536,201.2289) .. controls (47.7536,200.5081) and (47.1692,199.9238) ..(46.4484,199.9238) .. controls (45.7276,199.9238) and (45.1433,200.5081) ..(45.1433,201.2289) .. controls (45.1433,201.9497) and (45.7276,202.5340) ..(46.4484,202.5340) .. controls (47.1692,202.5340) and (47.7536,201.9497) ..(47.7536,201.2289);
\path[fill=chfourfig5c0,fill opacity=0.5,draw=chfourfig5c0,line width=0.3985pt,draw opacity=0.5] (67.0016,206.9475) .. controls (67.0016,206.2267) and (66.4173,205.6424) ..(65.6965,205.6424) .. controls (64.9757,205.6424) and (64.3914,206.2267) ..(64.3914,206.9475) .. controls (64.3914,207.6684) and (64.9757,208.2527) ..(65.6965,208.2527) .. controls (66.4173,208.2527) and (67.0016,207.6684) ..(67.0016,206.9475);
\path[fill=chfourfig5c0,fill opacity=0.5,draw=chfourfig5c0,line width=0.3985pt,draw opacity=0.5] (79.3953,221.8718) .. controls (79.3953,221.1510) and (78.8110,220.5667) ..(78.0902,220.5667) .. controls (77.3694,220.5667) and (76.7851,221.1510) ..(76.7851,221.8718) .. controls (76.7851,222.5926) and (77.3694,223.1769) ..(78.0902,223.1769) .. controls (78.8110,223.1769) and (79.3953,222.5926) ..(79.3953,221.8718);
\path[fill=chfourfig5c0,fill opacity=0.5,draw=chfourfig5c0,line width=0.3985pt,draw opacity=0.5] (46.3189,224.5518) .. controls (46.3189,223.8310) and (45.7346,223.2466) ..(45.0138,223.2466) .. controls (44.2930,223.2466) and (43.7087,223.8310) ..(43.7087,224.5518) .. controls (43.7087,225.2726) and (44.2930,225.8569) ..(45.0138,225.8569) .. controls (45.7346,225.8569) and (46.3189,225.2726) ..(46.3189,224.5518);
\path[fill=chfourfig5c0,fill opacity=0.5,draw=chfourfig5c0,line width=0.3985pt,draw opacity=0.5] (159.9841,247.8547) .. controls (159.9841,247.1339) and (159.3998,246.5496) ..(158.6790,246.5496) .. controls (157.9582,246.5496) and (157.3739,247.1339) ..(157.3739,247.8547) .. controls (157.3739,248.5755) and (157.9582,249.1598) ..(158.6790,249.1598) .. controls (159.3998,249.1598) and (159.9841,248.5755) ..(159.9841,247.8547);
\path[fill=chfourfig5c0,fill opacity=0.5,draw=chfourfig5c0,line width=0.3985pt,draw opacity=0.5] (124.8156,178.5038) .. controls (124.8156,177.7830) and (124.2313,177.1987) ..(123.5105,177.1987) .. controls (122.7897,177.1987) and (122.2053,177.7830) ..(122.2053,178.5038) .. controls (122.2053,179.2246) and (122.7897,179.8090) ..(123.5105,179.8090) .. controls (124.2313,179.8090) and (124.8156,179.2246) ..(124.8156,178.5038);
\path[fill=chfourfig5c0,fill opacity=0.5,draw=chfourfig5c0,line width=0.3985pt,draw opacity=0.5] (128.4918,219.2117) .. controls (128.4918,218.4909) and (127.9075,217.9066) ..(127.1867,217.9066) .. controls (126.4659,217.9066) and (125.8816,218.4909) ..(125.8816,219.2117) .. controls (125.8816,219.9325) and (126.4659,220.5168) ..(127.1867,220.5168) .. controls (127.9075,220.5168) and (128.4918,219.9325) ..(128.4918,219.2117);
\path[fill=chfourfig5c0,fill opacity=0.5,draw=chfourfig5c0,line width=0.3985pt,draw opacity=0.5] (63.6342,247.2270) .. controls (63.6342,246.5062) and (63.0499,245.9219) ..(62.3291,245.9219) .. controls (61.6083,245.9219) and (61.0240,246.5062) ..(61.0240,247.2270) .. controls (61.0240,247.9478) and (61.6083,248.5321) ..(62.3291,248.5321) .. controls (63.0499,248.5321) and (63.6342,247.9478) ..(63.6342,247.2270);
\path[fill=chfourfig5c0,fill opacity=0.5,draw=chfourfig5c0,line width=0.3985pt,draw opacity=0.5] (117.2339,259.2620) .. controls (117.2339,258.5413) and (116.6496,257.9569) ..(115.9288,257.9569) .. controls (115.2080,257.9569) and (114.6237,258.5413) ..(114.6237,259.2620) .. controls (114.6237,259.9828) and (115.2080,260.5671) ..(115.9288,260.5671) .. controls (116.6496,260.5671) and (117.2339,259.9828) ..(117.2339,259.2620);
\path[fill=chfourfig5c0,fill opacity=0.5,draw=chfourfig5c0,line width=0.3985pt,draw opacity=0.5] (60.6852,222.2404) .. controls (60.6852,221.5196) and (60.1009,220.9353) ..(59.3801,220.9353) .. controls (58.6593,220.9353) and (58.0750,221.5196) ..(58.0750,222.2404) .. controls (58.0750,222.9612) and (58.6593,223.5455) ..(59.3801,223.5455) .. controls (60.1009,223.5455) and (60.6852,222.9612) ..(60.6852,222.2404);
\path[fill=chfourfig5c0,fill opacity=0.5,draw=chfourfig5c0,line width=0.3985pt,draw opacity=0.5] (84.3867,206.1406) .. controls (84.3867,205.4198) and (83.8023,204.8354) ..(83.0815,204.8354) .. controls (82.3607,204.8354) and (81.7764,205.4198) ..(81.7764,206.1406) .. controls (81.7764,206.8614) and (82.3607,207.4457) ..(83.0815,207.4457) .. controls (83.8023,207.4457) and (84.3867,206.8614) ..(84.3867,206.1406);
\path[fill=chfourfig5c0,fill opacity=0.5,draw=chfourfig5c0,line width=0.3985pt,draw opacity=0.5] (101.7717,194.4742) .. controls (101.7717,193.7534) and (101.1874,193.1690) ..(100.4666,193.1690) .. controls (99.7458,193.1690) and (99.1615,193.7534) ..(99.1615,194.4742) .. controls (99.1615,195.1949) and (99.7458,195.7793) ..(100.4666,195.7793) .. controls (101.1874,195.7793) and (101.7717,195.1949) ..(101.7717,194.4742);
\path[fill=chfourfig5c0,fill opacity=0.5,draw=chfourfig5c0,line width=0.3985pt,draw opacity=0.5] (86.3095,255.4862) .. controls (86.3095,254.7654) and (85.7252,254.1810) ..(85.0044,254.1810) .. controls (84.2836,254.1810) and (83.6992,254.7654) ..(83.6992,255.4862) .. controls (83.6992,256.2070) and (84.2836,256.7913) ..(85.0044,256.7913) .. controls (85.7252,256.7913) and (86.3095,256.2070) ..(86.3095,255.4862);
\path[fill=chfourfig5c0,fill opacity=0.5,draw=chfourfig5c0,line width=0.3985pt,draw opacity=0.5] (165.2744,180.0979) .. controls (165.2744,179.3771) and (164.6901,178.7928) ..(163.9693,178.7928) .. controls (163.2485,178.7928) and (162.6641,179.3771) ..(162.6641,180.0979) .. controls (162.6641,180.8187) and (163.2485,181.4030) ..(163.9693,181.4030) .. controls (164.6901,181.4030) and (165.2744,180.8187) ..(165.2744,180.0979);
\path[fill=chfourfig5c0,fill opacity=0.5,draw=chfourfig5c0,line width=0.3985pt,draw opacity=0.5] (127.5354,262.5298) .. controls (127.5354,261.8090) and (126.9511,261.2247) ..(126.2303,261.2247) .. controls (125.5095,261.2247) and (124.9252,261.8090) ..(124.9252,262.5298) .. controls (124.9252,263.2506) and (125.5095,263.8350) ..(126.2303,263.8350) .. controls (126.9511,263.8350) and (127.5354,263.2506) ..(127.5354,262.5298);
\path[fill=chfourfig5c0,fill opacity=0.5,draw=chfourfig5c0,line width=0.3985pt,draw opacity=0.5] (47.9628,244.4574) .. controls (47.9628,243.7365) and (47.3785,243.1522) ..(46.6577,243.1522) .. controls (45.9369,243.1522) and (45.3525,243.7365) ..(45.3525,244.4574) .. controls (45.3525,245.1782) and (45.9369,245.7625) ..(46.6577,245.7625) .. controls (47.3785,245.7625) and (47.9628,245.1782) ..(47.9628,244.4574);
\path[fill=chfourfig5c0,fill opacity=0.5,draw=chfourfig5c0,line width=0.3985pt,draw opacity=0.5] (134.0809,194.1753) .. controls (134.0809,193.4545) and (133.4966,192.8701) ..(132.7758,192.8701) .. controls (132.0550,192.8701) and (131.4707,193.4545) ..(131.4707,194.1753) .. controls (131.4707,194.8961) and (132.0550,195.4804) ..(132.7758,195.4804) .. controls (133.4966,195.4804) and (134.0809,194.8961) ..(134.0809,194.1753);
\path[fill=chfourfig5c0,fill opacity=0.5,draw=chfourfig5c0,line width=0.3985pt,draw opacity=0.5] (100.2275,262.0417) .. controls (100.2275,261.3209) and (99.6432,260.7365) ..(98.9223,260.7365) .. controls (98.2016,260.7365) and (97.6172,261.3209) ..(97.6172,262.0417) .. controls (97.6172,262.7625) and (98.2016,263.3468) ..(98.9223,263.3468) .. controls (99.6432,263.3468) and (100.2275,262.7625) ..(100.2275,262.0417);
\path[fill=chfourfig5c0,fill opacity=0.5,draw=chfourfig5c0,line width=0.3985pt,draw opacity=0.5] (101.9809,209.0497) .. controls (101.9809,208.3289) and (101.3966,207.7446) ..(100.6758,207.7446) .. controls (99.9550,207.7446) and (99.3707,208.3289) ..(99.3707,209.0497) .. controls (99.3707,209.7705) and (99.9550,210.3548) ..(100.6758,210.3548) .. controls (101.3966,210.3548) and (101.9809,209.7705) ..(101.9809,209.0497);
\path[fill=chfourfig5c0,fill opacity=0.5,draw=chfourfig5c0,line width=0.3985pt,draw opacity=0.5] (57.7462,216.7609) .. controls (57.7462,216.0401) and (57.1619,215.4557) ..(56.4411,215.4557) .. controls (55.7203,215.4557) and (55.1360,216.0401) ..(55.1360,216.7609) .. controls (55.1360,217.4817) and (55.7203,218.0660) ..(56.4411,218.0660) .. controls (57.1619,218.0660) and (57.7462,217.4817) ..(57.7462,216.7609);
\path[fill=chfourfig5c0,fill opacity=0.5,draw=chfourfig5c0,line width=0.3985pt,draw opacity=0.5] (91.0318,257.8872) .. controls (91.0318,257.1664) and (90.4475,256.5820) ..(89.7267,256.5820) .. controls (89.0059,256.5820) and (88.4216,257.1664) ..(88.4216,257.8872) .. controls (88.4216,258.6080) and (89.0059,259.1923) ..(89.7267,259.1923) .. controls (90.4475,259.1923) and (91.0318,258.6080) ..(91.0318,257.8872);
\path[fill=chfourfig5c0,fill opacity=0.5,draw=chfourfig5c0,line width=0.3985pt,draw opacity=0.5] (93.8613,232.3925) .. controls (93.8613,231.6716) and (93.2769,231.0873) ..(92.5561,231.0873) .. controls (91.8354,231.0873) and (91.2510,231.6716) ..(91.2510,232.3925) .. controls (91.2510,233.1133) and (91.8354,233.6976) ..(92.5561,233.6976) .. controls (93.2769,233.6976) and (93.8613,233.1133) ..(93.8613,232.3925);
\path[fill=chfourfig5c0,fill opacity=0.5,draw=chfourfig5c0,line width=0.3985pt,draw opacity=0.5] (71.4251,221.7224) .. controls (71.4251,221.0015) and (70.8408,220.4172) ..(70.1200,220.4172) .. controls (69.3992,220.4172) and (68.8149,221.0015) ..(68.8149,221.7224) .. controls (68.8149,222.4431) and (69.3992,223.0275) ..(70.1200,223.0275) .. controls (70.8408,223.0275) and (71.4251,222.4431) ..(71.4251,221.7224);
\path[fill=chfourfig5c0,fill opacity=0.5,draw=chfourfig5c0,line width=0.3985pt,draw opacity=0.5] (62.7077,252.0590) .. controls (62.7077,251.3382) and (62.1234,250.7538) ..(61.4026,250.7538) .. controls (60.6818,250.7538) and (60.0974,251.3382) ..(60.0974,252.0590) .. controls (60.0974,252.7798) and (60.6818,253.3641) ..(61.4026,253.3641) .. controls (62.1234,253.3641) and (62.7077,252.7798) ..(62.7077,252.0590);
\path[fill=chfourfig5c0,fill opacity=0.5,draw=chfourfig5c0,line width=0.3985pt,draw opacity=0.5] (105.9461,257.4986) .. controls (105.9461,256.7778) and (105.3618,256.1935) ..(104.6410,256.1935) .. controls (103.9202,256.1935) and (103.3358,256.7778) ..(103.3358,257.4986) .. controls (103.3358,258.2194) and (103.9202,258.8037) ..(104.6410,258.8037) .. controls (105.3618,258.8037) and (105.9461,258.2194) ..(105.9461,257.4986);
\path[fill=chfourfig5c0,fill opacity=0.5,draw=chfourfig5c0,line width=0.3985pt,draw opacity=0.5] (90.0156,230.6589) .. controls (90.0156,229.9381) and (89.4313,229.3538) ..(88.7105,229.3538) .. controls (87.9897,229.3538) and (87.4054,229.9381) ..(87.4054,230.6589) .. controls (87.4054,231.3798) and (87.9897,231.9641) ..(88.7105,231.9641) .. controls (89.4313,231.9641) and (90.0156,231.3798) ..(90.0156,230.6589);
\path[fill=chfourfig5c0,fill opacity=0.5,draw=chfourfig5c0,line width=0.3985pt,draw opacity=0.5] (162.5546,207.3660) .. controls (162.5546,206.6452) and (161.9702,206.0609) ..(161.2494,206.0609) .. controls (160.5286,206.0609) and (159.9443,206.6452) ..(159.9443,207.3660) .. controls (159.9443,208.0868) and (160.5286,208.6711) ..(161.2494,208.6711) .. controls (161.9702,208.6711) and (162.5546,208.0868) ..(162.5546,207.3660);
\path[fill=chfourfig5c0,fill opacity=0.5,draw=chfourfig5c0,line width=0.3985pt,draw opacity=0.5] (164.9257,262.7490) .. controls (164.9257,262.0282) and (164.3414,261.4439) ..(163.6206,261.4439) .. controls (162.8997,261.4439) and (162.3154,262.0282) ..(162.3154,262.7490) .. controls (162.3154,263.4698) and (162.8997,264.0541) ..(163.6206,264.0541) .. controls (164.3414,264.0541) and (164.9257,263.4698) ..(164.9257,262.7490);
\path[fill=chfourfig5c0,fill opacity=0.5,draw=chfourfig5c0,line width=0.3985pt,draw opacity=0.5] (101.5625,210.3149) .. controls (101.5625,209.5941) and (100.9782,209.0098) ..(100.2574,209.0098) .. controls (99.5366,209.0098) and (98.9522,209.5941) ..(98.9522,210.3149) .. controls (98.9522,211.0358) and (99.5366,211.6201) ..(100.2574,211.6201) .. controls (100.9782,211.6201) and (101.5625,211.0358) ..(101.5625,210.3149);
\path[fill=chfourfig5c0,fill opacity=0.5,draw=chfourfig5c0,line width=0.3985pt,draw opacity=0.5] (95.8040,180.5063) .. controls (95.8040,179.7856) and (95.2197,179.2012) ..(94.4989,179.2012) .. controls (93.7781,179.2012) and (93.1938,179.7856) ..(93.1938,180.5063) .. controls (93.1938,181.2272) and (93.7781,181.8115) ..(94.4989,181.8115) .. controls (95.2197,181.8115) and (95.8040,181.2272) ..(95.8040,180.5063);
\path[fill=chfourfig5c0,fill opacity=0.5,draw=chfourfig5c0,line width=0.3985pt,draw opacity=0.5] (56.6105,214.7783) .. controls (56.6105,214.0575) and (56.0262,213.4731) ..(55.3053,213.4731) .. controls (54.5845,213.4731) and (54.0002,214.0575) ..(54.0002,214.7783) .. controls (54.0002,215.4991) and (54.5845,216.0834) ..(55.3053,216.0834) .. controls (56.0262,216.0834) and (56.6105,215.4991) ..(56.6105,214.7783);
\path[fill=chfourfig5c0,fill opacity=0.5,draw=chfourfig5c0,line width=0.3985pt,draw opacity=0.5] (161.3789,208.7209) .. controls (161.3789,208.0001) and (160.7946,207.4158) ..(160.0738,207.4158) .. controls (159.3530,207.4158) and (158.7687,208.0001) ..(158.7687,208.7209) .. controls (158.7687,209.4417) and (159.3530,210.0260) ..(160.0738,210.0260) .. controls (160.7946,210.0260) and (161.3789,209.4417) ..(161.3789,208.7209);
\path[fill=chfourfig5c0,fill opacity=0.5,draw=chfourfig5c0,line width=0.3985pt,draw opacity=0.5] (53.4921,236.6466) .. controls (53.4921,235.9258) and (52.9078,235.3414) ..(52.1870,235.3414) .. controls (51.4662,235.3414) and (50.8819,235.9258) ..(50.8819,236.6466) .. controls (50.8819,237.3674) and (51.4662,237.9517) ..(52.1870,237.9517) .. controls (52.9078,237.9517) and (53.4921,237.3674) ..(53.4921,236.6466);
\path[fill=chfourfig5c0,fill opacity=0.5,draw=chfourfig5c0,line width=0.3985pt,draw opacity=0.5] (121.4083,176.3718) .. controls (121.4083,175.6510) and (120.8240,175.0667) ..(120.1032,175.0667) .. controls (119.3824,175.0667) and (118.7981,175.6510) ..(118.7981,176.3718) .. controls (118.7981,177.0926) and (119.3824,177.6769) ..(120.1032,177.6769) .. controls (120.8240,177.6769) and (121.4083,177.0926) ..(121.4083,176.3718);
\path[fill=chfourfig5c0,fill opacity=0.5,draw=chfourfig5c0,line width=0.3985pt,draw opacity=0.5] (92.5462,215.3561) .. controls (92.5462,214.6353) and (91.9619,214.0510) ..(91.2410,214.0510) .. controls (90.5202,214.0510) and (89.9359,214.6353) ..(89.9359,215.3561) .. controls (89.9359,216.0769) and (90.5202,216.6612) ..(91.2410,216.6612) .. controls (91.9619,216.6612) and (92.5462,216.0769) ..(92.5462,215.3561);
\path[fill=chfourfig5c0,fill opacity=0.5,draw=chfourfig5c0,line width=0.3985pt,draw opacity=0.5] (134.4396,193.5177) .. controls (134.4396,192.7969) and (133.8553,192.2126) ..(133.1345,192.2126) .. controls (132.4137,192.2126) and (131.8294,192.7969) ..(131.8294,193.5177) .. controls (131.8294,194.2385) and (132.4137,194.8228) ..(133.1345,194.8228) .. controls (133.8553,194.8228) and (134.4396,194.2385) ..(134.4396,193.5177);
\path[fill=chfourfig5c0,fill opacity=0.5,draw=chfourfig5c0,line width=0.3985pt,draw opacity=0.5] (53.8906,222.5393) .. controls (53.8906,221.8185) and (53.3063,221.2341) ..(52.5855,221.2341) .. controls (51.8647,221.2341) and (51.2804,221.8185) ..(51.2804,222.5393) .. controls (51.2804,223.2601) and (51.8647,223.8444) ..(52.5855,223.8444) .. controls (53.3063,223.8444) and (53.8906,223.2601) ..(53.8906,222.5393);
\path[fill=chfourfig5c0,fill opacity=0.5,draw=chfourfig5c0,line width=0.3985pt,draw opacity=0.5] (63.9132,219.1619) .. controls (63.9132,218.4411) and (63.3288,217.8568) ..(62.6080,217.8568) .. controls (61.8872,217.8568) and (61.3029,218.4411) ..(61.3029,219.1619) .. controls (61.3029,219.8827) and (61.8872,220.4670) ..(62.6080,220.4670) .. controls (63.3288,220.4670) and (63.9132,219.8827) ..(63.9132,219.1619);
\path[fill=chfourfig5c0,fill opacity=0.5,draw=chfourfig5c0,line width=0.3985pt,draw opacity=0.5] (129.3686,186.4641) .. controls (129.3686,185.7433) and (128.7843,185.1590) ..(128.0634,185.1590) .. controls (127.3427,185.1590) and (126.7583,185.7433) ..(126.7583,186.4641) .. controls (126.7583,187.1849) and (127.3427,187.7692) ..(128.0634,187.7692) .. controls (128.7843,187.7692) and (129.3686,187.1849) ..(129.3686,186.4641);
\path[fill=chfourfig5c0,fill opacity=0.5,draw=chfourfig5c0,line width=0.3985pt,draw opacity=0.5] (80.2720,262.3007) .. controls (80.2720,261.5799) and (79.6877,260.9955) ..(78.9669,260.9955) .. controls (78.2461,260.9955) and (77.6618,261.5799) ..(77.6618,262.3007) .. controls (77.6618,263.0215) and (78.2461,263.6058) ..(78.9669,263.6058) .. controls (79.6877,263.6058) and (80.2720,263.0215) ..(80.2720,262.3007);
\path[fill=chfourfig5c0,fill opacity=0.5,draw=chfourfig5c0,line width=0.3985pt,draw opacity=0.5] (165.3939,177.9957) .. controls (165.3939,177.2749) and (164.8096,176.6906) ..(164.0888,176.6906) .. controls (163.3680,176.6906) and (162.7837,177.2749) ..(162.7837,177.9957) .. controls (162.7837,178.7166) and (163.3680,179.3009) ..(164.0888,179.3009) .. controls (164.8096,179.3009) and (165.3939,178.7166) ..(165.3939,177.9957);
\path[fill=chfourfig5c0,fill opacity=0.5,draw=chfourfig5c0,line width=0.3985pt,draw opacity=0.5] (147.9790,230.4896) .. controls (147.9790,229.7687) and (147.3947,229.1844) ..(146.6739,229.1844) .. controls (145.9531,229.1844) and (145.3688,229.7687) ..(145.3688,230.4896) .. controls (145.3688,231.2104) and (145.9531,231.7947) ..(146.6739,231.7947) .. controls (147.3947,231.7947) and (147.9790,231.2104) ..(147.9790,230.4896);
\path[fill=chfourfig5c0,fill opacity=0.5,draw=chfourfig5c0,line width=0.3985pt,draw opacity=0.5] (97.3980,246.7687) .. controls (97.3980,246.0479) and (96.8137,245.4636) ..(96.0929,245.4636) .. controls (95.3721,245.4636) and (94.7878,246.0479) ..(94.7878,246.7687) .. controls (94.7878,247.4895) and (95.3721,248.0739) ..(96.0929,248.0739) .. controls (96.8137,248.0739) and (97.3980,247.4895) ..(97.3980,246.7687);
\path[fill=chfourfig5c0,fill opacity=0.5,draw=chfourfig5c0,line width=0.3985pt,draw opacity=0.5] (84.4863,191.6647) .. controls (84.4863,190.9439) and (83.9020,190.3595) ..(83.1812,190.3595) .. controls (82.4604,190.3595) and (81.8761,190.9439) ..(81.8761,191.6647) .. controls (81.8761,192.3855) and (82.4604,192.9698) ..(83.1812,192.9698) .. controls (83.9020,192.9698) and (84.4863,192.3855) ..(84.4863,191.6647);
\path[fill=chfourfig5c0,fill opacity=0.5,draw=chfourfig5c0,line width=0.3985pt,draw opacity=0.5] (107.1715,243.6304) .. controls (107.1715,242.9097) and (106.5872,242.3253) ..(105.8664,242.3253) .. controls (105.1456,242.3253) and (104.5613,242.9097) ..(104.5613,243.6304) .. controls (104.5613,244.3513) and (105.1456,244.9356) ..(105.8664,244.9356) .. controls (106.5872,244.9356) and (107.1715,244.3513) ..(107.1715,243.6304);
\path[fill=chfourfig5c0,fill opacity=0.5,draw=chfourfig5c0,line width=0.3985pt,draw opacity=0.5] (65.7862,177.2087) .. controls (65.7862,176.4879) and (65.2019,175.9035) ..(64.4810,175.9035) .. controls (63.7603,175.9035) and (63.1759,176.4879) ..(63.1759,177.2087) .. controls (63.1759,177.9295) and (63.7603,178.5138) ..(64.4810,178.5138) .. controls (65.2019,178.5138) and (65.7862,177.9295) ..(65.7862,177.2087);
\path[fill=chfourfig5c0,fill opacity=0.5,draw=chfourfig5c0,line width=0.3985pt,draw opacity=0.5] (74.1947,171.3605) .. controls (74.1947,170.6397) and (73.6104,170.0554) ..(72.8896,170.0554) .. controls (72.1688,170.0554) and (71.5845,170.6397) ..(71.5845,171.3605) .. controls (71.5845,172.0813) and (72.1688,172.6656) ..(72.8896,172.6656) .. controls (73.6104,172.6656) and (74.1947,172.0813) ..(74.1947,171.3605);
\path[fill=chfourfig5c0,fill opacity=0.5,draw=chfourfig5c0,line width=0.3985pt,draw opacity=0.5] (71.1163,181.2038) .. controls (71.1163,180.4830) and (70.5319,179.8986) ..(69.8112,179.8986) .. controls (69.0903,179.8986) and (68.5060,180.4830) ..(68.5060,181.2038) .. controls (68.5060,181.9246) and (69.0903,182.5089) ..(69.8112,182.5089) .. controls (70.5319,182.5089) and (71.1163,181.9246) ..(71.1163,181.2038);
\path[fill=chfourfig5c0,fill opacity=0.5,draw=chfourfig5c0,line width=0.3985pt,draw opacity=0.5] (93.7317,252.6368) .. controls (93.7317,251.9160) and (93.1474,251.3317) ..(92.4266,251.3317) .. controls (91.7058,251.3317) and (91.1215,251.9160) ..(91.1215,252.6368) .. controls (91.1215,253.3576) and (91.7058,253.9419) ..(92.4266,253.9419) .. controls (93.1474,253.9419) and (93.7317,253.3576) ..(93.7317,252.6368);
\path[fill=chfourfig5c0,fill opacity=0.5,draw=chfourfig5c0,line width=0.3985pt,draw opacity=0.5] (63.6442,203.2613) .. controls (63.6442,202.5405) and (63.0599,201.9562) ..(62.3391,201.9562) .. controls (61.6183,201.9562) and (61.0339,202.5405) ..(61.0339,203.2613) .. controls (61.0339,203.9821) and (61.6183,204.5665) ..(62.3391,204.5665) .. controls (63.0599,204.5665) and (63.6442,203.9821) ..(63.6442,203.2613);
\path[fill=chfourfig5c0,fill opacity=0.5,draw=chfourfig5c0,line width=0.3985pt,draw opacity=0.5] (98.5438,265.1301) .. controls (98.5438,264.4093) and (97.9594,263.8250) ..(97.2386,263.8250) .. controls (96.5178,263.8250) and (95.9335,264.4093) ..(95.9335,265.1301) .. controls (95.9335,265.8509) and (96.5178,266.4352) ..(97.2386,266.4352) .. controls (97.9594,266.4352) and (98.5438,265.8509) ..(98.5438,265.1301);
\path[fill=chfourfig5c0,fill opacity=0.5,draw=chfourfig5c0,line width=0.3985pt,draw opacity=0.5] (120.1530,230.3900) .. controls (120.1530,229.6691) and (119.5687,229.0848) ..(118.8479,229.0848) .. controls (118.1271,229.0848) and (117.5428,229.6691) ..(117.5428,230.3900) .. controls (117.5428,231.1107) and (118.1271,231.6951) ..(118.8479,231.6951) .. controls (119.5687,231.6951) and (120.1530,231.1107) ..(120.1530,230.3900);
\path[fill=chfourfig5c0,fill opacity=0.5,draw=chfourfig5c0,line width=0.3985pt,draw opacity=0.5] (54.3987,230.7386) .. controls (54.3987,230.0178) and (53.8144,229.4335) ..(53.0936,229.4335) .. controls (52.3728,229.4335) and (51.7885,230.0178) ..(51.7885,230.7386) .. controls (51.7885,231.4594) and (52.3728,232.0438) ..(53.0936,232.0438) .. controls (53.8144,232.0438) and (54.3987,231.4594) ..(54.3987,230.7386);
\path[fill=chfourfig5c0,fill opacity=0.5,draw=chfourfig5c0,line width=0.3985pt,draw opacity=0.5] (157.8721,206.0310) .. controls (157.8721,205.3102) and (157.2877,204.7259) ..(156.5669,204.7259) .. controls (155.8461,204.7259) and (155.2618,205.3102) ..(155.2618,206.0310) .. controls (155.2618,206.7518) and (155.8461,207.3361) ..(156.5669,207.3361) .. controls (157.2877,207.3361) and (157.8721,206.7518) ..(157.8721,206.0310);
\path[fill=chfourfig5c0,fill opacity=0.5,draw=chfourfig5c0,line width=0.3985pt,draw opacity=0.5] (115.6698,259.7602) .. controls (115.6698,259.0394) and (115.0854,258.4551) ..(114.3646,258.4551) .. controls (113.6438,258.4551) and (113.0595,259.0394) ..(113.0595,259.7602) .. controls (113.0595,260.4810) and (113.6438,261.0653) ..(114.3646,261.0653) .. controls (115.0854,261.0653) and (115.6698,260.4810) ..(115.6698,259.7602);
\path[fill=chfourfig5c0,fill opacity=0.5,draw=chfourfig5c0,line width=0.3985pt,draw opacity=0.5] (72.2620,188.4467) .. controls (72.2620,187.7259) and (71.6777,187.1416) ..(70.9568,187.1416) .. controls (70.2361,187.1416) and (69.6517,187.7259) ..(69.6517,188.4467) .. controls (69.6517,189.1675) and (70.2361,189.7518) ..(70.9568,189.7518) .. controls (71.6777,189.7518) and (72.2620,189.1675) ..(72.2620,188.4467);
\path[fill=chfourfig5c0,fill opacity=0.5,draw=chfourfig5c0,line width=0.3985pt,draw opacity=0.5] (99.4305,201.7968) .. controls (99.4305,201.0760) and (98.8461,200.4917) ..(98.1253,200.4917) .. controls (97.4045,200.4917) and (96.8202,201.0760) ..(96.8202,201.7968) .. controls (96.8202,202.5176) and (97.4045,203.1019) ..(98.1253,203.1019) .. controls (98.8461,203.1019) and (99.4305,202.5176) ..(99.4305,201.7968);
\path[fill=chfourfig5c0,fill opacity=0.5,draw=chfourfig5c0,line width=0.3985pt,draw opacity=0.5] (84.5859,204.5067) .. controls (84.5859,203.7859) and (84.0016,203.2015) ..(83.2808,203.2015) .. controls (82.5600,203.2015) and (81.9757,203.7859) ..(81.9757,204.5067) .. controls (81.9757,205.2275) and (82.5600,205.8118) ..(83.2808,205.8118) .. controls (84.0016,205.8118) and (84.5859,205.2275) ..(84.5859,204.5067);
\path[fill=chfourfig5c0,fill opacity=0.5,draw=chfourfig5c0,line width=0.3985pt,draw opacity=0.5] (71.3155,262.2010) .. controls (71.3155,261.4803) and (70.7312,260.8959) ..(70.0104,260.8959) .. controls (69.2896,260.8959) and (68.7053,261.4803) ..(68.7053,262.2010) .. controls (68.7053,262.9218) and (69.2896,263.5062) ..(70.0104,263.5062) .. controls (70.7312,263.5062) and (71.3155,262.9218) ..(71.3155,262.2010);
\path[fill=chfourfig5c0,fill opacity=0.5,draw=chfourfig5c0,line width=0.3985pt,draw opacity=0.5] (105.9959,182.5089) .. controls (105.9959,181.7881) and (105.4116,181.2038) ..(104.6908,181.2038) .. controls (103.9700,181.2038) and (103.3857,181.7881) ..(103.3857,182.5089) .. controls (103.3857,183.2297) and (103.9700,183.8140) ..(104.6908,183.8140) .. controls (105.4116,183.8140) and (105.9959,183.2297) ..(105.9959,182.5089);
\path[fill=chfourfig5c0,fill opacity=0.5,draw=chfourfig5c0,line width=0.3985pt,draw opacity=0.5] (119.7445,202.6337) .. controls (119.7445,201.9129) and (119.1602,201.3286) ..(118.4394,201.3286) .. controls (117.7186,201.3286) and (117.1343,201.9129) ..(117.1343,202.6337) .. controls (117.1343,203.3545) and (117.7186,203.9388) ..(118.4394,203.9388) .. controls (119.1602,203.9388) and (119.7445,203.3545) ..(119.7445,202.6337);
\path[fill=chfourfig5c0,fill opacity=0.5,draw=chfourfig5c0,line width=0.3985pt,draw opacity=0.5] (52.8844,261.8524) .. controls (52.8844,261.1316) and (52.3001,260.5472) ..(51.5793,260.5472) .. controls (50.8585,260.5472) and (50.2742,261.1316) ..(50.2742,261.8524) .. controls (50.2742,262.5732) and (50.8585,263.1575) ..(51.5793,263.1575) .. controls (52.3001,263.1575) and (52.8844,262.5732) ..(52.8844,261.8524);
\path[fill=chfourfig5c0,fill opacity=0.5,draw=chfourfig5c0,line width=0.3985pt,draw opacity=0.5] (107.3409,186.7032) .. controls (107.3409,185.9824) and (106.7566,185.3981) ..(106.0358,185.3981) .. controls (105.3150,185.3981) and (104.7307,185.9824) ..(104.7307,186.7032) .. controls (104.7307,187.4240) and (105.3150,188.0083) ..(106.0358,188.0083) .. controls (106.7566,188.0083) and (107.3409,187.4240) ..(107.3409,186.7032);
\path[fill=chfourfig5c0,fill opacity=0.5,draw=chfourfig5c0,line width=0.3985pt,draw opacity=0.5] (71.1461,200.2028) .. controls (71.1461,199.4819) and (70.5618,198.8976) ..(69.8410,198.8976) .. controls (69.1202,198.8976) and (68.5359,199.4819) ..(68.5359,200.2028) .. controls (68.5359,200.9236) and (69.1202,201.5079) ..(69.8410,201.5079) .. controls (70.5618,201.5079) and (71.1461,200.9236) ..(71.1461,200.2028);
\path[fill=chfourfig5c0,fill opacity=0.5,draw=chfourfig5c0,line width=0.3985pt,draw opacity=0.5] (92.6259,235.0625) .. controls (92.6259,234.3417) and (92.0415,233.7574) ..(91.3208,233.7574) .. controls (90.5999,233.7574) and (90.0156,234.3417) ..(90.0156,235.0625) .. controls (90.0156,235.7833) and (90.5999,236.3676) ..(91.3208,236.3676) .. controls (92.0415,236.3676) and (92.6259,235.7833) ..(92.6259,235.0625);
\path[fill=chfourfig5c0,fill opacity=0.5,draw=chfourfig5c0,line width=0.3985pt,draw opacity=0.5] (151.1273,240.3228) .. controls (151.1273,239.6020) and (150.5429,239.0177) ..(149.8221,239.0177) .. controls (149.1013,239.0177) and (148.5170,239.6020) ..(148.5170,240.3228) .. controls (148.5170,241.0436) and (149.1013,241.6279) ..(149.8221,241.6279) .. controls (150.5429,241.6279) and (151.1273,241.0436) ..(151.1273,240.3228);
\path[fill=chfourfig5c0,fill opacity=0.5,draw=chfourfig5c0,line width=0.3985pt,draw opacity=0.5] (159.7650,225.8569) .. controls (159.7650,225.1361) and (159.1806,224.5518) ..(158.4599,224.5518) .. controls (157.7390,224.5518) and (157.1547,225.1361) ..(157.1547,225.8569) .. controls (157.1547,226.5777) and (157.7390,227.1620) ..(158.4599,227.1620) .. controls (159.1806,227.1620) and (159.7650,226.5777) ..(159.7650,225.8569);
\path[fill=chfourfig5c0,fill opacity=0.5,draw=chfourfig5c0,line width=0.3985pt,draw opacity=0.5] (153.8670,265.1500) .. controls (153.8670,264.4293) and (153.2827,263.8449) ..(152.5619,263.8449) .. controls (151.8411,263.8449) and (151.2568,264.4293) ..(151.2568,265.1500) .. controls (151.2568,265.8708) and (151.8411,266.4551) ..(152.5619,266.4551) .. controls (153.2827,266.4551) and (153.8670,265.8708) ..(153.8670,265.1500);
\path[fill=chfourfig5c0,fill opacity=0.5,draw=chfourfig5c0,line width=0.3985pt,draw opacity=0.5] (128.1033,194.2849) .. controls (128.1033,193.5641) and (127.5190,192.9797) ..(126.7982,192.9797) .. controls (126.0774,192.9797) and (125.4931,193.5641) ..(125.4931,194.2849) .. controls (125.4931,195.0057) and (126.0774,195.5900) ..(126.7982,195.5900) .. controls (127.5190,195.5900) and (128.1033,195.0057) ..(128.1033,194.2849);
\path[fill=chfourfig5c0,fill opacity=0.5,draw=chfourfig5c0,line width=0.3985pt,draw opacity=0.5] (108.9947,265.5884) .. controls (108.9947,264.8676) and (108.4104,264.2833) ..(107.6896,264.2833) .. controls (106.9688,264.2833) and (106.3845,264.8676) ..(106.3845,265.5884) .. controls (106.3845,266.3092) and (106.9688,266.8935) ..(107.6896,266.8935) .. controls (108.4104,266.8935) and (108.9947,266.3092) ..(108.9947,265.5884);
\path[fill=chfourfig5c0,fill opacity=0.5,draw=chfourfig5c0,line width=0.3985pt,draw opacity=0.5] (71.9133,236.4573) .. controls (71.9133,235.7365) and (71.3289,235.1522) ..(70.6082,235.1522) .. controls (69.8873,235.1522) and (69.3030,235.7365) ..(69.3030,236.4573) .. controls (69.3030,237.1781) and (69.8873,237.7624) ..(70.6082,237.7624) .. controls (71.3289,237.7624) and (71.9133,237.1781) ..(71.9133,236.4573);
\path[fill=chfourfig5c0,fill opacity=0.5,draw=chfourfig5c0,line width=0.3985pt,draw opacity=0.5] (151.0276,227.7697) .. controls (151.0276,227.0489) and (150.4433,226.4646) ..(149.7225,226.4646) .. controls (149.0017,226.4646) and (148.4174,227.0489) ..(148.4174,227.7697) .. controls (148.4174,228.4905) and (149.0017,229.0748) ..(149.7225,229.0748) .. controls (150.4433,229.0748) and (151.0276,228.4905) ..(151.0276,227.7697);
\path[fill=chfourfig5c0,fill opacity=0.5,draw=chfourfig5c0,line width=0.3985pt,draw opacity=0.5] (146.5643,193.6472) .. controls (146.5643,192.9265) and (145.9800,192.3421) ..(145.2592,192.3421) .. controls (144.5384,192.3421) and (143.9541,192.9265) ..(143.9541,193.6472) .. controls (143.9541,194.3680) and (144.5384,194.9524) ..(145.2592,194.9524) .. controls (145.9800,194.9524) and (146.5643,194.3680) ..(146.5643,193.6472);
\path[fill=chfourfig5c0,fill opacity=0.5,draw=chfourfig5c0,line width=0.3985pt,draw opacity=0.5] (119.1268,189.3633) .. controls (119.1268,188.6425) and (118.5425,188.0581) ..(117.8217,188.0581) .. controls (117.1009,188.0581) and (116.5166,188.6425) ..(116.5166,189.3633) .. controls (116.5166,190.0841) and (117.1009,190.6684) ..(117.8217,190.6684) .. controls (118.5425,190.6684) and (119.1268,190.0841) ..(119.1268,189.3633);
\path[fill=chfourfig5c0,fill opacity=0.5,draw=chfourfig5c0,line width=0.3985pt,draw opacity=0.5] (105.6472,233.4385) .. controls (105.6472,232.7177) and (105.0629,232.1334) ..(104.3421,232.1334) .. controls (103.6213,232.1334) and (103.0370,232.7177) ..(103.0370,233.4385) .. controls (103.0370,234.1594) and (103.6213,234.7437) ..(104.3421,234.7437) .. controls (105.0629,234.7437) and (105.6472,234.1594) ..(105.6472,233.4385);
\path[fill=chfourfig5c0,fill opacity=0.5,draw=chfourfig5c0,line width=0.3985pt,draw opacity=0.5] (117.2339,235.2916) .. controls (117.2339,234.5708) and (116.6496,233.9865) ..(115.9288,233.9865) .. controls (115.2080,233.9865) and (114.6237,234.5708) ..(114.6237,235.2916) .. controls (114.6237,236.0124) and (115.2080,236.5967) ..(115.9288,236.5967) .. controls (116.6496,236.5967) and (117.2339,236.0124) ..(117.2339,235.2916);
\path[fill=chfourfig5c0,fill opacity=0.5,draw=chfourfig5c0,line width=0.3985pt,draw opacity=0.5] (128.8405,249.7775) .. controls (128.8405,249.0567) and (128.2562,248.4724) ..(127.5354,248.4724) .. controls (126.8146,248.4724) and (126.2303,249.0567) ..(126.2303,249.7775) .. controls (126.2303,250.4983) and (126.8146,251.0826) ..(127.5354,251.0826) .. controls (128.2562,251.0826) and (128.8405,250.4983) ..(128.8405,249.7775);
\path[fill=chfourfig5c0,fill opacity=0.5,draw=chfourfig5c0,line width=0.3985pt,draw opacity=0.5] (132.4172,189.0245) .. controls (132.4172,188.3037) and (131.8329,187.7194) ..(131.1120,187.7194) .. controls (130.3913,187.7194) and (129.8069,188.3037) ..(129.8069,189.0245) .. controls (129.8069,189.7453) and (130.3913,190.3296) ..(131.1120,190.3296) .. controls (131.8329,190.3296) and (132.4172,189.7453) ..(132.4172,189.0245);
\path[fill=chfourfig5c0,fill opacity=0.5,draw=chfourfig5c0,line width=0.3985pt,draw opacity=0.5] (147.9292,200.7407) .. controls (147.9292,200.0199) and (147.3449,199.4356) ..(146.6241,199.4356) .. controls (145.9033,199.4356) and (145.3190,200.0199) ..(145.3190,200.7407) .. controls (145.3190,201.4615) and (145.9033,202.0458) ..(146.6241,202.0458) .. controls (147.3449,202.0458) and (147.9292,201.4615) ..(147.9292,200.7407);
\path[fill=chfourfig5c0,fill opacity=0.5,draw=chfourfig5c0,line width=0.3985pt,draw opacity=0.5] (56.8894,200.8503) .. controls (56.8894,200.1295) and (56.3051,199.5452) ..(55.5843,199.5452) .. controls (54.8635,199.5452) and (54.2792,200.1295) ..(54.2792,200.8503) .. controls (54.2792,201.5711) and (54.8635,202.1555) ..(55.5843,202.1555) .. controls (56.3051,202.1555) and (56.8894,201.5711) ..(56.8894,200.8503);
\path[fill=chfourfig5c0,fill opacity=0.5,draw=chfourfig5c0,line width=0.3985pt,draw opacity=0.5] (145.2592,189.5625) .. controls (145.2592,188.8417) and (144.6749,188.2574) ..(143.9541,188.2574) .. controls (143.2333,188.2574) and (142.6489,188.8417) ..(142.6489,189.5625) .. controls (142.6489,190.2833) and (143.2333,190.8676) ..(143.9541,190.8676) .. controls (144.6749,190.8676) and (145.2592,190.2833) ..(145.2592,189.5625);
\path[fill=chfourfig5c0,fill opacity=0.5,draw=chfourfig5c0,line width=0.3985pt,draw opacity=0.5] (165.2345,183.0867) .. controls (165.2345,182.3659) and (164.6502,181.7816) ..(163.9294,181.7816) .. controls (163.2086,181.7816) and (162.6243,182.3659) ..(162.6243,183.0867) .. controls (162.6243,183.8075) and (163.2086,184.3918) ..(163.9294,184.3918) .. controls (164.6502,184.3918) and (165.2345,183.8075) ..(165.2345,183.0867);
\path[fill=chfourfig5c0,fill opacity=0.5,draw=chfourfig5c0,line width=0.3985pt,draw opacity=0.5] (91.5399,225.3986) .. controls (91.5399,224.6778) and (90.9556,224.0935) ..(90.2348,224.0935) .. controls (89.5140,224.0935) and (88.9297,224.6778) ..(88.9297,225.3986) .. controls (88.9297,226.1194) and (89.5140,226.7037) ..(90.2348,226.7037) .. controls (90.9556,226.7037) and (91.5399,226.1194) ..(91.5399,225.3986);
\path[fill=chfourfig5c0,fill opacity=0.5,draw=chfourfig5c0,line width=0.3985pt,draw opacity=0.5] (156.9156,247.3067) .. controls (156.9156,246.5859) and (156.3313,246.0016) ..(155.6105,246.0016) .. controls (154.8897,246.0016) and (154.3054,246.5859) ..(154.3054,247.3067) .. controls (154.3054,248.0275) and (154.8897,248.6118) ..(155.6105,248.6118) .. controls (156.3313,248.6118) and (156.9156,248.0275) ..(156.9156,247.3067);
\path[fill=chfourfig5c0,fill opacity=0.5,draw=chfourfig5c0,line width=0.3985pt,draw opacity=0.5] (146.0263,267.2223) .. controls (146.0263,266.5015) and (145.4420,265.9172) ..(144.7212,265.9172) .. controls (144.0004,265.9172) and (143.4161,266.5015) ..(143.4161,267.2223) .. controls (143.4161,267.9431) and (144.0004,268.5274) ..(144.7212,268.5274) .. controls (145.4420,268.5274) and (146.0263,267.9431) ..(146.0263,267.2223);
\path[fill=chfourfig5c0,fill opacity=0.5,draw=chfourfig5c0,line width=0.3985pt,draw opacity=0.5] (155.7998,208.7309) .. controls (155.7998,208.0101) and (155.2155,207.4258) ..(154.4947,207.4258) .. controls (153.7739,207.4258) and (153.1895,208.0101) ..(153.1895,208.7309) .. controls (153.1895,209.4517) and (153.7739,210.0360) ..(154.4947,210.0360) .. controls (155.2155,210.0360) and (155.7998,209.4517) ..(155.7998,208.7309);
\path[fill=chfourfig5c0,fill opacity=0.5,draw=chfourfig5c0,line width=0.3985pt,draw opacity=0.5] (137.7672,219.3213) .. controls (137.7672,218.6005) and (137.1829,218.0162) ..(136.4621,218.0162) .. controls (135.7413,218.0162) and (135.1569,218.6005) ..(135.1569,219.3213) .. controls (135.1569,220.0421) and (135.7413,220.6264) ..(136.4621,220.6264) .. controls (137.1829,220.6264) and (137.7672,220.0421) ..(137.7672,219.3213);
\path[fill=chfourfig5c0,fill opacity=0.5,draw=chfourfig5c0,line width=0.3985pt,draw opacity=0.5] (105.7369,216.0037) .. controls (105.7369,215.2829) and (105.1526,214.6986) ..(104.4318,214.6986) .. controls (103.7110,214.6986) and (103.1266,215.2829) ..(103.1266,216.0037) .. controls (103.1266,216.7245) and (103.7110,217.3088) ..(104.4318,217.3088) .. controls (105.1526,217.3088) and (105.7369,216.7245) ..(105.7369,216.0037);
\path[fill=chfourfig5c0,fill opacity=0.5,draw=chfourfig5c0,line width=0.3985pt,draw opacity=0.5] (126.7085,247.3565) .. controls (126.7085,246.6357) and (126.1242,246.0514) ..(125.4034,246.0514) .. controls (124.6826,246.0514) and (124.0982,246.6357) ..(124.0982,247.3565) .. controls (124.0982,248.0773) and (124.6826,248.6617) ..(125.4034,248.6617) .. controls (126.1242,248.6617) and (126.7085,248.0773) ..(126.7085,247.3565);
\path[fill=chfourfig5c0,fill opacity=0.5,draw=chfourfig5c0,line width=0.3985pt,draw opacity=0.5] (98.0456,239.3465) .. controls (98.0456,238.6257) and (97.4613,238.0413) ..(96.7405,238.0413) .. controls (96.0197,238.0413) and (95.4354,238.6257) ..(95.4354,239.3465) .. controls (95.4354,240.0673) and (96.0197,240.6516) ..(96.7405,240.6516) .. controls (97.4613,240.6516) and (98.0456,240.0673) ..(98.0456,239.3465);
\path[fill=chfourfig5c0,fill opacity=0.5,draw=chfourfig5c0,line width=0.3985pt,draw opacity=0.5] (137.4185,267.5909) .. controls (137.4185,266.8701) and (136.8342,266.2858) ..(136.1134,266.2858) .. controls (135.3925,266.2858) and (134.8082,266.8701) ..(134.8082,267.5909) .. controls (134.8082,268.3117) and (135.3925,268.8961) ..(136.1134,268.8961) .. controls (136.8342,268.8961) and (137.4185,268.3117) ..(137.4185,267.5909);
\path[fill=chfourfig5c0,fill opacity=0.5,draw=chfourfig5c0,line width=0.3985pt,draw opacity=0.5] (114.0757,213.3935) .. controls (114.0757,212.6727) and (113.4914,212.0883) ..(112.7706,212.0883) .. controls (112.0498,212.0883) and (111.4655,212.6727) ..(111.4655,213.3935) .. controls (111.4655,214.1143) and (112.0498,214.6986) ..(112.7706,214.6986) .. controls (113.4914,214.6986) and (114.0757,214.1143) ..(114.0757,213.3935);
\path[fill=chfourfig5c0,fill opacity=0.5,draw=chfourfig5c0,line width=0.3985pt,draw opacity=0.5] (139.6103,265.9570) .. controls (139.6103,265.2362) and (139.0260,264.6519) ..(138.3052,264.6519) .. controls (137.5844,264.6519) and (137.0000,265.2362) ..(137.0000,265.9570) .. controls (137.0000,266.6778) and (137.5844,267.2621) ..(138.3052,267.2621) .. controls (139.0260,267.2621) and (139.6103,266.6778) ..(139.6103,265.9570);
\path[fill=chfourfig5c0,fill opacity=0.5,draw=chfourfig5c0,line width=0.3985pt,draw opacity=0.5] (117.2937,182.8875) .. controls (117.2937,182.1667) and (116.7094,181.5823) ..(115.9886,181.5823) .. controls (115.2678,181.5823) and (114.6834,182.1667) ..(114.6834,182.8875) .. controls (114.6834,183.6083) and (115.2678,184.1926) ..(115.9886,184.1926) .. controls (116.7094,184.1926) and (117.2937,183.6083) ..(117.2937,182.8875);
\path[fill=chfourfig5c0,fill opacity=0.5,draw=chfourfig5c0,line width=0.3985pt,draw opacity=0.5] (97.3582,191.9436) .. controls (97.3582,191.2228) and (96.7739,190.6385) ..(96.0531,190.6385) .. controls (95.3323,190.6385) and (94.7479,191.2228) ..(94.7479,191.9436) .. controls (94.7479,192.6644) and (95.3323,193.2487) ..(96.0531,193.2487) .. controls (96.7739,193.2487) and (97.3582,192.6644) ..(97.3582,191.9436);
\path[fill=chfourfig5c0,fill opacity=0.5,draw=chfourfig5c0,line width=0.3985pt,draw opacity=0.5] (80.9993,236.4174) .. controls (80.9993,235.6966) and (80.4150,235.1123) ..(79.6942,235.1123) .. controls (78.9734,235.1123) and (78.3891,235.6966) ..(78.3891,236.4174) .. controls (78.3891,237.1382) and (78.9734,237.7225) ..(79.6942,237.7225) .. controls (80.4150,237.7225) and (80.9993,237.1382) ..(80.9993,236.4174);
\path[fill=chfourfig5c0,fill opacity=0.5,draw=chfourfig5c0,line width=0.3985pt,draw opacity=0.5] (150.2206,188.0681) .. controls (150.2206,187.3473) and (149.6363,186.7630) ..(148.9155,186.7630) .. controls (148.1947,186.7630) and (147.6104,187.3473) ..(147.6104,188.0681) .. controls (147.6104,188.7889) and (148.1947,189.3732) ..(148.9155,189.3732) .. controls (149.6363,189.3732) and (150.2206,188.7889) ..(150.2206,188.0681);
\path[fill=chfourfig5c0,fill opacity=0.5,draw=chfourfig5c0,line width=0.3985pt,draw opacity=0.5] (95.2162,215.4358) .. controls (95.2162,214.7150) and (94.6319,214.1307) ..(93.9111,214.1307) .. controls (93.1903,214.1307) and (92.6060,214.7150) ..(92.6060,215.4358) .. controls (92.6060,216.1566) and (93.1903,216.7409) ..(93.9111,216.7409) .. controls (94.6319,216.7409) and (95.2162,216.1566) ..(95.2162,215.4358);
\path[fill=chfourfig5c0,fill opacity=0.5,draw=chfourfig5c0,line width=0.3985pt,draw opacity=0.5] (52.4161,175.6047) .. controls (52.4161,174.8839) and (51.8318,174.2995) ..(51.1110,174.2995) .. controls (50.3902,174.2995) and (49.8059,174.8839) ..(49.8059,175.6047) .. controls (49.8059,176.3255) and (50.3902,176.9098) ..(51.1110,176.9098) .. controls (51.8318,176.9098) and (52.4161,176.3255) ..(52.4161,175.6047);
\path[fill=chfourfig5c0,fill opacity=0.5,draw=chfourfig5c0,line width=0.3985pt,draw opacity=0.5] (62.5782,243.4013) .. controls (62.5782,242.6805) and (61.9938,242.0962) ..(61.2730,242.0962) .. controls (60.5522,242.0962) and (59.9679,242.6805) ..(59.9679,243.4013) .. controls (59.9679,244.1221) and (60.5522,244.7064) ..(61.2730,244.7064) .. controls (61.9938,244.7064) and (62.5782,244.1221) ..(62.5782,243.4013);
\path[fill=chfourfig5c0,fill opacity=0.5,draw=chfourfig5c0,line width=0.3985pt,draw opacity=0.5] (149.5930,257.9769) .. controls (149.5930,257.2560) and (149.0087,256.6717) ..(148.2879,256.6717) .. controls (147.5670,256.6717) and (146.9827,257.2560) ..(146.9827,257.9769) .. controls (146.9827,258.6976) and (147.5670,259.2820) ..(148.2879,259.2820) .. controls (149.0087,259.2820) and (149.5930,258.6976) ..(149.5930,257.9769);
\path[fill=chfourfig5c0,fill opacity=0.5,draw=chfourfig5c0,line width=0.3985pt,draw opacity=0.5] (111.5053,225.4285) .. controls (111.5053,224.7077) and (110.9210,224.1234) ..(110.2002,224.1234) .. controls (109.4794,224.1234) and (108.8951,224.7077) ..(108.8951,225.4285) .. controls (108.8951,226.1493) and (109.4794,226.7336) ..(110.2002,226.7336) .. controls (110.9210,226.7336) and (111.5053,226.1493) ..(111.5053,225.4285);
\path[fill=chfourfig5c0,fill opacity=0.5,draw=chfourfig5c0,line width=0.3985pt,draw opacity=0.5] (113.0794,203.7495) .. controls (113.0794,203.0287) and (112.4951,202.4444) ..(111.7743,202.4444) .. controls (111.0535,202.4444) and (110.4692,203.0287) ..(110.4692,203.7495) .. controls (110.4692,204.4703) and (111.0535,205.0546) ..(111.7743,205.0546) .. controls (112.4951,205.0546) and (113.0794,204.4703) ..(113.0794,203.7495);
\path[fill=chfourfig5c0,fill opacity=0.5,draw=chfourfig5c0,line width=0.3985pt,draw opacity=0.5] (57.5071,262.5896) .. controls (57.5071,261.8688) and (56.9228,261.2845) ..(56.2020,261.2845) .. controls (55.4812,261.2845) and (54.8969,261.8688) ..(54.8969,262.5896) .. controls (54.8969,263.3104) and (55.4812,263.8947) ..(56.2020,263.8947) .. controls (56.9228,263.8947) and (57.5071,263.3104) ..(57.5071,262.5896);
\path[fill=chfourfig5c0,fill opacity=0.5,draw=chfourfig5c0,line width=0.3985pt,draw opacity=0.5] (155.3216,195.5701) .. controls (155.3216,194.8492) and (154.7372,194.2650) ..(154.0164,194.2650) .. controls (153.2956,194.2650) and (152.7113,194.8492) ..(152.7113,195.5701) .. controls (152.7113,196.2909) and (153.2956,196.8752) ..(154.0164,196.8752) .. controls (154.7372,196.8752) and (155.3216,196.2909) ..(155.3216,195.5701);
\path[fill=chfourfig5c0,fill opacity=0.5,draw=chfourfig5c0,line width=0.3985pt,draw opacity=0.5] (62.3889,200.4917) .. controls (62.3889,199.7708) and (61.8046,199.1865) ..(61.0837,199.1865) .. controls (60.3629,199.1865) and (59.7786,199.7708) ..(59.7786,200.4917) .. controls (59.7786,201.2125) and (60.3629,201.7968) ..(61.0837,201.7968) .. controls (61.8046,201.7968) and (62.3889,201.2125) ..(62.3889,200.4917);
\path[fill=chfourfig5c0,fill opacity=0.5,draw=chfourfig5c0,line width=0.3985pt,draw opacity=0.5] (67.1212,179.8189) .. controls (67.1212,179.0981) and (66.5369,178.5138) ..(65.8161,178.5138) .. controls (65.0953,178.5138) and (64.5109,179.0981) ..(64.5109,179.8189) .. controls (64.5109,180.5397) and (65.0953,181.1240) ..(65.8161,181.1240) .. controls (66.5369,181.1240) and (67.1212,180.5397) ..(67.1212,179.8189);
\path[fill=chfourfig5c0,fill opacity=0.5,draw=chfourfig5c0,line width=0.3985pt,draw opacity=0.5] (71.2856,191.1566) .. controls (71.2856,190.4357) and (70.7013,189.8514) ..(69.9805,189.8514) .. controls (69.2597,189.8514) and (68.6754,190.4357) ..(68.6754,191.1566) .. controls (68.6754,191.8773) and (69.2597,192.4617) ..(69.9805,192.4617) .. controls (70.7013,192.4617) and (71.2856,191.8773) ..(71.2856,191.1566);
\path[fill=chfourfig5c0,fill opacity=0.5,draw=chfourfig5c0,line width=0.3985pt,draw opacity=0.5] (107.3807,181.7318) .. controls (107.3807,181.0110) and (106.7964,180.4267) ..(106.0756,180.4267) .. controls (105.3548,180.4267) and (104.7705,181.0110) ..(104.7705,181.7318) .. controls (104.7705,182.4526) and (105.3548,183.0369) ..(106.0756,183.0369) .. controls (106.7964,183.0369) and (107.3807,182.4526) ..(107.3807,181.7318);
\path[fill=chfourfig5c0,fill opacity=0.5,draw=chfourfig5c0,line width=0.3985pt,draw opacity=0.5] (91.2809,254.7688) .. controls (91.2809,254.0480) and (90.6966,253.4637) ..(89.9758,253.4637) .. controls (89.2550,253.4637) and (88.6707,254.0480) ..(88.6707,254.7688) .. controls (88.6707,255.4896) and (89.2550,256.0739) ..(89.9758,256.0739) .. controls (90.6966,256.0739) and (91.2809,255.4896) ..(91.2809,254.7688);
\path[fill=chfourfig5c0,fill opacity=0.5,draw=chfourfig5c0,line width=0.3985pt,draw opacity=0.5] (105.6671,246.6193) .. controls (105.6671,245.8985) and (105.0828,245.3141) ..(104.3620,245.3141) .. controls (103.6412,245.3141) and (103.0569,245.8985) ..(103.0569,246.6193) .. controls (103.0569,247.3401) and (103.6412,247.9244) ..(104.3620,247.9244) .. controls (105.0828,247.9244) and (105.6671,247.3401) ..(105.6671,246.6193);
\path[fill=chfourfig5c0,fill opacity=0.5,draw=chfourfig5c0,line width=0.3985pt,draw opacity=0.5] (81.1089,227.1022) .. controls (81.1089,226.3814) and (80.5246,225.7971) ..(79.8038,225.7971) .. controls (79.0830,225.7971) and (78.4987,226.3814) ..(78.4987,227.1022) .. controls (78.4987,227.8230) and (79.0830,228.4073) ..(79.8038,228.4073) .. controls (80.5246,228.4073) and (81.1089,227.8230) ..(81.1089,227.1022);
\path[fill=chfourfig5c0,fill opacity=0.5,draw=chfourfig5c0,line width=0.3985pt,draw opacity=0.5] (56.0027,230.2803) .. controls (56.0027,229.5596) and (55.4184,228.9752) ..(54.6976,228.9752) .. controls (53.9768,228.9752) and (53.3925,229.5596) ..(53.3925,230.2803) .. controls (53.3925,231.0012) and (53.9768,231.5854) ..(54.6976,231.5854) .. controls (55.4184,231.5854) and (56.0027,231.0012) ..(56.0027,230.2803);
\path[fill=chfourfig5c0,fill opacity=0.5,draw=chfourfig5c0,line width=0.3985pt,draw opacity=0.5] (118.5789,249.9568) .. controls (118.5789,249.2360) and (117.9946,248.6517) ..(117.2738,248.6517) .. controls (116.5530,248.6517) and (115.9686,249.2360) ..(115.9686,249.9568) .. controls (115.9686,250.6776) and (116.5530,251.2619) ..(117.2738,251.2619) .. controls (117.9946,251.2619) and (118.5789,250.6776) ..(118.5789,249.9568);
\path[fill=chfourfig5c0,fill opacity=0.5,draw=chfourfig5c0,line width=0.3985pt,draw opacity=0.5] (124.4868,177.1888) .. controls (124.4868,176.4679) and (123.9025,175.8836) ..(123.1817,175.8836) .. controls (122.4609,175.8836) and (121.8766,176.4679) ..(121.8766,177.1888) .. controls (121.8766,177.9095) and (122.4609,178.4939) ..(123.1817,178.4939) .. controls (123.9025,178.4939) and (124.4868,177.9095) ..(124.4868,177.1888);
\path[fill=chfourfig5c0,fill opacity=0.5,draw=chfourfig5c0,line width=0.3985pt,draw opacity=0.5] (110.1105,184.3121) .. controls (110.1105,183.5913) and (109.5262,183.0070) ..(108.8054,183.0070) .. controls (108.0846,183.0070) and (107.5003,183.5913) ..(107.5003,184.3121) .. controls (107.5003,185.0329) and (108.0846,185.6172) ..(108.8054,185.6172) .. controls (109.5262,185.6172) and (110.1105,185.0329) ..(110.1105,184.3121);
\path[fill=chfourfig5c0,fill opacity=0.5,draw=chfourfig5c0,line width=0.3985pt,draw opacity=0.5] (62.5383,239.5457) .. controls (62.5383,238.8249) and (61.9540,238.2406) ..(61.2332,238.2406) .. controls (60.5124,238.2406) and (59.9281,238.8249) ..(59.9281,239.5457) .. controls (59.9281,240.2665) and (60.5124,240.8508) ..(61.2332,240.8508) .. controls (61.9540,240.8508) and (62.5383,240.2665) ..(62.5383,239.5457);
\path[fill=chfourfig5c0,fill opacity=0.5,draw=chfourfig5c0,line width=0.3985pt,draw opacity=0.5] (72.8298,255.9544) .. controls (72.8298,255.2336) and (72.2455,254.6493) ..(71.5247,254.6493) .. controls (70.8039,254.6493) and (70.2196,255.2336) ..(70.2196,255.9544) .. controls (70.2196,256.6752) and (70.8039,257.2595) ..(71.5247,257.2595) .. controls (72.2455,257.2595) and (72.8298,256.6752) ..(72.8298,255.9544);
\path[fill=chfourfig5c0,fill opacity=0.5,draw=chfourfig5c0,line width=0.3985pt,draw opacity=0.5] (116.9848,215.7148) .. controls (116.9848,214.9940) and (116.4005,214.4097) ..(115.6797,214.4097) .. controls (114.9589,214.4097) and (114.3746,214.9940) ..(114.3746,215.7148) .. controls (114.3746,216.4356) and (114.9589,217.0199) ..(115.6797,217.0199) .. controls (116.4005,217.0199) and (116.9848,216.4356) ..(116.9848,215.7148);
\path[fill=chfourfig5c0,fill opacity=0.5,draw=chfourfig5c0,line width=0.3985pt,draw opacity=0.5] (154.3054,241.6479) .. controls (154.3054,240.9271) and (153.7211,240.3428) ..(153.0003,240.3428) .. controls (152.2794,240.3428) and (151.6951,240.9271) ..(151.6951,241.6479) .. controls (151.6951,242.3687) and (152.2794,242.9530) ..(153.0003,242.9530) .. controls (153.7211,242.9530) and (154.3054,242.3687) ..(154.3054,241.6479);
\path[fill=chfourfig5c0,fill opacity=0.5,draw=chfourfig5c0,line width=0.3985pt,draw opacity=0.5] (128.6712,236.9654) .. controls (128.6712,236.2446) and (128.0868,235.6602) ..(127.3660,235.6602) .. controls (126.6452,235.6602) and (126.0609,236.2446) ..(126.0609,236.9654) .. controls (126.0609,237.6862) and (126.6452,238.2705) ..(127.3660,238.2705) .. controls (128.0868,238.2705) and (128.6712,237.6862) ..(128.6712,236.9654);
\path[fill=chfourfig5c0,fill opacity=0.5,draw=chfourfig5c0,line width=0.3985pt,draw opacity=0.5] (122.4444,187.7692) .. controls (122.4444,187.0484) and (121.8601,186.4641) ..(121.1393,186.4641) .. controls (120.4185,186.4641) and (119.8342,187.0484) ..(119.8342,187.7692) .. controls (119.8342,188.4900) and (120.4185,189.0743) ..(121.1393,189.0743) .. controls (121.8601,189.0743) and (122.4444,188.4900) ..(122.4444,187.7692);
\path[fill=chfourfig5c0,fill opacity=0.5,draw=chfourfig5c0,line width=0.3985pt,draw opacity=0.5] (90.0356,264.3829) .. controls (90.0356,263.6621) and (89.4512,263.0778) ..(88.7304,263.0778) .. controls (88.0096,263.0778) and (87.4253,263.6621) ..(87.4253,264.3829) .. controls (87.4253,265.1037) and (88.0096,265.6880) ..(88.7304,265.6880) .. controls (89.4512,265.6880) and (90.0356,265.1037) ..(90.0356,264.3829);
\path[fill=chfourfig5c0,fill opacity=0.5,draw=chfourfig5c0,line width=0.3985pt,draw opacity=0.5] (85.7117,182.9074) .. controls (85.7117,182.1866) and (85.1274,181.6023) ..(84.4066,181.6023) .. controls (83.6858,181.6023) and (83.1015,182.1866) ..(83.1015,182.9074) .. controls (83.1015,183.6282) and (83.6858,184.2125) ..(84.4066,184.2125) .. controls (85.1274,184.2125) and (85.7117,183.6282) ..(85.7117,182.9074);
\path[fill=chfourfig5c0,fill opacity=0.5,draw=chfourfig5c0,line width=0.3985pt,draw opacity=0.5] (114.4941,212.9850) .. controls (114.4941,212.2642) and (113.9098,211.6799) ..(113.1890,211.6799) .. controls (112.4682,211.6799) and (111.8839,212.2642) ..(111.8839,212.9850) .. controls (111.8839,213.7058) and (112.4682,214.2901) ..(113.1890,214.2901) .. controls (113.9098,214.2901) and (114.4941,213.7058) ..(114.4941,212.9850);
\path[fill=chfourfig5c0,fill opacity=0.5,draw=chfourfig5c0,line width=0.3985pt,draw opacity=0.5] (65.3279,230.2405) .. controls (65.3279,229.5197) and (64.7436,228.9354) ..(64.0228,228.9354) .. controls (63.3020,228.9354) and (62.7177,229.5197) ..(62.7177,230.2405) .. controls (62.7177,230.9613) and (63.3020,231.5456) ..(64.0228,231.5456) .. controls (64.7436,231.5456) and (65.3279,230.9613) ..(65.3279,230.2405);
\path[fill=chfourfig5c0,fill opacity=0.5,draw=chfourfig5c0,line width=0.3985pt,draw opacity=0.5] (144.8905,217.9863) .. controls (144.8905,217.2655) and (144.3062,216.6812) ..(143.5854,216.6812) .. controls (142.8646,216.6812) and (142.2803,217.2655) ..(142.2803,217.9863) .. controls (142.2803,218.7071) and (142.8646,219.2914) ..(143.5854,219.2914) .. controls (144.3062,219.2914) and (144.8905,218.7071) ..(144.8905,217.9863);
\path[fill=chfourfig5c0,fill opacity=0.5,draw=chfourfig5c0,line width=0.3985pt,draw opacity=0.5] (58.0949,194.2948) .. controls (58.0949,193.5740) and (57.5106,192.9897) ..(56.7898,192.9897) .. controls (56.0690,192.9897) and (55.4847,193.5740) ..(55.4847,194.2948) .. controls (55.4847,195.0157) and (56.0690,195.6000) ..(56.7898,195.6000) .. controls (57.5106,195.6000) and (58.0949,195.0157) ..(58.0949,194.2948);
\path[fill=chfourfig5c0,fill opacity=0.5,draw=chfourfig5c0,line width=0.3985pt,draw opacity=0.5] (108.3969,178.8027) .. controls (108.3969,178.0819) and (107.8126,177.4976) ..(107.0918,177.4976) .. controls (106.3710,177.4976) and (105.7867,178.0819) ..(105.7867,178.8027) .. controls (105.7867,179.5235) and (106.3710,180.1078) ..(107.0918,180.1078) .. controls (107.8126,180.1078) and (108.3969,179.5235) ..(108.3969,178.8027);
\path[fill=chfourfig5c0,fill opacity=0.5,draw=chfourfig5c0,line width=0.3985pt,draw opacity=0.5] (97.3084,218.6837) .. controls (97.3084,217.9629) and (96.7240,217.3786) ..(96.0033,217.3786) .. controls (95.2825,217.3786) and (94.6981,217.9629) ..(94.6981,218.6837) .. controls (94.6981,219.4045) and (95.2825,219.9888) ..(96.0033,219.9888) .. controls (96.7240,219.9888) and (97.3084,219.4045) ..(97.3084,218.6837);
\path[fill=chfourfig5c0,fill opacity=0.5,draw=chfourfig5c0,line width=0.3985pt,draw opacity=0.5] (138.0163,188.1777) .. controls (138.0163,187.4569) and (137.4319,186.8726) ..(136.7111,186.8726) .. controls (135.9903,186.8726) and (135.4060,187.4569) ..(135.4060,188.1777) .. controls (135.4060,188.8985) and (135.9903,189.4828) ..(136.7111,189.4828) .. controls (137.4319,189.4828) and (138.0163,188.8985) ..(138.0163,188.1777);
\path[fill=chfourfig5c0,fill opacity=0.5,draw=chfourfig5c0,line width=0.3985pt,draw opacity=0.5] (117.6424,218.9427) .. controls (117.6424,218.2219) and (117.0581,217.6376) ..(116.3373,217.6376) .. controls (115.6165,217.6376) and (115.0322,218.2219) ..(115.0322,218.9427) .. controls (115.0322,219.6635) and (115.6165,220.2478) ..(116.3373,220.2478) .. controls (117.0581,220.2478) and (117.6424,219.6635) ..(117.6424,218.9427);
\path[fill=chfourfig5c0,fill opacity=0.5,draw=chfourfig5c0,line width=0.3985pt,draw opacity=0.5] (95.7940,207.3062) .. controls (95.7940,206.5854) and (95.2097,206.0011) ..(94.4889,206.0011) .. controls (93.7681,206.0011) and (93.1838,206.5854) ..(93.1838,207.3062) .. controls (93.1838,208.0270) and (93.7681,208.6113) ..(94.4889,208.6113) .. controls (95.2097,208.6113) and (95.7940,208.0270) ..(95.7940,207.3062);
\path[fill=chfourfig5c0,fill opacity=0.5,draw=chfourfig5c0,line width=0.3985pt,draw opacity=0.5] (123.7296,223.8046) .. controls (123.7296,223.0837) and (123.1453,222.4995) ..(122.4245,222.4995) .. controls (121.7037,222.4995) and (121.1194,223.0837) ..(121.1194,223.8046) .. controls (121.1194,224.5253) and (121.7037,225.1097) ..(122.4245,225.1097) .. controls (123.1453,225.1097) and (123.7296,224.5253) ..(123.7296,223.8046);
\path[fill=chfourfig5c0,fill opacity=0.5,draw=chfourfig5c0,line width=0.3985pt,draw opacity=0.5] (139.7000,210.5939) .. controls (139.7000,209.8731) and (139.1156,209.2888) ..(138.3948,209.2888) .. controls (137.6740,209.2888) and (137.0897,209.8731) ..(137.0897,210.5939) .. controls (137.0897,211.3147) and (137.6740,211.8990) ..(138.3948,211.8990) .. controls (139.1156,211.8990) and (139.7000,211.3147) ..(139.7000,210.5939);
\path[fill=chfourfig5c0,fill opacity=0.5,draw=chfourfig5c0,line width=0.3985pt,draw opacity=0.5] (155.5607,222.0411) .. controls (155.5607,221.3203) and (154.9764,220.7360) ..(154.2556,220.7360) .. controls (153.5348,220.7360) and (152.9504,221.3203) ..(152.9504,222.0411) .. controls (152.9504,222.7619) and (153.5348,223.3463) ..(154.2556,223.3463) .. controls (154.9764,223.3463) and (155.5607,222.7619) ..(155.5607,222.0411);
\path[fill=chfourfig5c0,fill opacity=0.5,draw=chfourfig5c0,line width=0.3985pt,draw opacity=0.5] (127.0970,225.9266) .. controls (127.0970,225.2058) and (126.5127,224.6215) ..(125.7919,224.6215) .. controls (125.0711,224.6215) and (124.4868,225.2058) ..(124.4868,225.9266) .. controls (124.4868,226.6474) and (125.0711,227.2318) ..(125.7919,227.2318) .. controls (126.5127,227.2318) and (127.0970,226.6474) ..(127.0970,225.9266);
\path[fill=chfourfig5c0,fill opacity=0.5,draw=chfourfig5c0,line width=0.3985pt,draw opacity=0.5] (75.8386,207.0671) .. controls (75.8386,206.3463) and (75.2543,205.7620) ..(74.5335,205.7620) .. controls (73.8127,205.7620) and (73.2284,206.3463) ..(73.2284,207.0671) .. controls (73.2284,207.7879) and (73.8127,208.3722) ..(74.5335,208.3722) .. controls (75.2543,208.3722) and (75.8386,207.7879) ..(75.8386,207.0671);
\path[fill=chfourfig5c0,fill opacity=0.5,draw=chfourfig5c0,line width=0.3985pt,draw opacity=0.5] (78.9769,230.3401) .. controls (78.9769,229.6193) and (78.3926,229.0350) ..(77.6718,229.0350) .. controls (76.9510,229.0350) and (76.3666,229.6193) ..(76.3666,230.3401) .. controls (76.3666,231.0609) and (76.9510,231.6453) ..(77.6718,231.6453) .. controls (78.3926,231.6453) and (78.9769,231.0609) ..(78.9769,230.3401);
\path[fill=chfourfig5c0,fill opacity=0.5,draw=chfourfig5c0,line width=0.3985pt,draw opacity=0.5] (94.5786,264.2335) .. controls (94.5786,263.5127) and (93.9942,262.9283) ..(93.2735,262.9283) .. controls (92.5526,262.9283) and (91.9683,263.5127) ..(91.9683,264.2335) .. controls (91.9683,264.9543) and (92.5526,265.5386) ..(93.2735,265.5386) .. controls (93.9942,265.5386) and (94.5786,264.9543) ..(94.5786,264.2335);
\path[fill=chfourfig5c0,fill opacity=0.5,draw=chfourfig5c0,line width=0.3985pt,draw opacity=0.5] (131.0025,234.0961) .. controls (131.0025,233.3753) and (130.4181,232.7910) ..(129.6974,232.7910) .. controls (128.9765,232.7910) and (128.3922,233.3753) ..(128.3922,234.0961) .. controls (128.3922,234.8169) and (128.9765,235.4012) ..(129.6974,235.4012) .. controls (130.4181,235.4012) and (131.0025,234.8169) ..(131.0025,234.0961);
\path[fill=chfourfig5c0,fill opacity=0.5,draw=chfourfig5c0,line width=0.3985pt,draw opacity=0.5] (113.9861,220.9951) .. controls (113.9861,220.2742) and (113.4017,219.6899) ..(112.6809,219.6899) .. controls (111.9601,219.6899) and (111.3758,220.2742) ..(111.3758,220.9951) .. controls (111.3758,221.7159) and (111.9601,222.3002) ..(112.6809,222.3002) .. controls (113.4017,222.3002) and (113.9861,221.7159) ..(113.9861,220.9951);
\path[fill=chfourfig5c0,fill opacity=0.5,draw=chfourfig5c0,line width=0.3985pt,draw opacity=0.5] (141.8918,187.8091) .. controls (141.8918,187.0883) and (141.3074,186.5039) ..(140.5866,186.5039) .. controls (139.8658,186.5039) and (139.2815,187.0883) ..(139.2815,187.8091) .. controls (139.2815,188.5299) and (139.8658,189.1142) ..(140.5866,189.1142) .. controls (141.3074,189.1142) and (141.8918,188.5299) ..(141.8918,187.8091);
\path[fill=chfourfig5c0,fill opacity=0.5,draw=chfourfig5c0,line width=0.3985pt,draw opacity=0.5] (98.2449,175.5051) .. controls (98.2449,174.7843) and (97.6606,174.1999) ..(96.9398,174.1999) .. controls (96.2189,174.1999) and (95.6346,174.7843) ..(95.6346,175.5051) .. controls (95.6346,176.2259) and (96.2189,176.8102) ..(96.9398,176.8102) .. controls (97.6606,176.8102) and (98.2449,176.2259) ..(98.2449,175.5051);
\path[fill=chfourfig5c0,fill opacity=0.5,draw=chfourfig5c0,line width=0.3985pt,draw opacity=0.5] (125.0646,261.7029) .. controls (125.0646,260.9821) and (124.4803,260.3978) ..(123.7595,260.3978) .. controls (123.0387,260.3978) and (122.4544,260.9821) ..(122.4544,261.7029) .. controls (122.4544,262.4237) and (123.0387,263.0081) ..(123.7595,263.0081) .. controls (124.4803,263.0081) and (125.0646,262.4237) ..(125.0646,261.7029);
\path[fill=chfourfig5c0,fill opacity=0.5,draw=chfourfig5c0,line width=0.3985pt,draw opacity=0.5] (132.6164,266.9035) .. controls (132.6164,266.1827) and (132.0321,265.5984) ..(131.3113,265.5984) .. controls (130.5905,265.5984) and (130.0062,266.1827) ..(130.0062,266.9035) .. controls (130.0062,267.6243) and (130.5905,268.2086) ..(131.3113,268.2086) .. controls (132.0321,268.2086) and (132.6164,267.6243) ..(132.6164,266.9035);
\path[fill=chfourfig5c0,fill opacity=0.5,draw=chfourfig5c0,line width=0.3985pt,draw opacity=0.5] (140.9353,243.5508) .. controls (140.9353,242.8300) and (140.3510,242.2456) ..(139.6302,242.2456) .. controls (138.9094,242.2456) and (138.3251,242.8300) ..(138.3251,243.5508) .. controls (138.3251,244.2715) and (138.9094,244.8559) ..(139.6302,244.8559) .. controls (140.3510,244.8559) and (140.9353,244.2715) ..(140.9353,243.5508);
\path[fill=chfourfig5c0,fill opacity=0.5,draw=chfourfig5c0,line width=0.3985pt,draw opacity=0.5] (122.6636,199.6847) .. controls (122.6636,198.9639) and (122.0793,198.3796) ..(121.3585,198.3796) .. controls (120.6377,198.3796) and (120.0534,198.9639) ..(120.0534,199.6847) .. controls (120.0534,200.4055) and (120.6377,200.9898) ..(121.3585,200.9898) .. controls (122.0793,200.9898) and (122.6636,200.4055) ..(122.6636,199.6847);
\path[fill=chfourfig5c0,fill opacity=0.5,draw=chfourfig5c0,line width=0.3985pt,draw opacity=0.5] (73.9756,232.1932) .. controls (73.9756,231.4724) and (73.3913,230.8881) ..(72.6705,230.8881) .. controls (71.9496,230.8881) and (71.3653,231.4724) ..(71.3653,232.1932) .. controls (71.3653,232.9140) and (71.9496,233.4983) ..(72.6705,233.4983) .. controls (73.3913,233.4983) and (73.9756,232.9140) ..(73.9756,232.1932);
\path[fill=chfourfig5c0,fill opacity=0.5,draw=chfourfig5c0,line width=0.3985pt,draw opacity=0.5] (73.2284,236.3178) .. controls (73.2284,235.5970) and (72.6440,235.0127) ..(71.9232,235.0127) .. controls (71.2024,235.0127) and (70.6181,235.5970) ..(70.6181,236.3178) .. controls (70.6181,237.0386) and (71.2024,237.6229) ..(71.9232,237.6229) .. controls (72.6440,237.6229) and (73.2284,237.0386) ..(73.2284,236.3178);
\path[fill=chfourfig5c0,fill opacity=0.5,draw=chfourfig5c0,line width=0.3985pt,draw opacity=0.5] (154.4847,222.2404) .. controls (154.4847,221.5196) and (153.9004,220.9353) ..(153.1796,220.9353) .. controls (152.4588,220.9353) and (151.8745,221.5196) ..(151.8745,222.2404) .. controls (151.8745,222.9612) and (152.4588,223.5455) ..(153.1796,223.5455) .. controls (153.9004,223.5455) and (154.4847,222.9612) ..(154.4847,222.2404);
\path[fill=chfourfig5c0,fill opacity=0.5,draw=chfourfig5c0,line width=0.3985pt,draw opacity=0.5] (48.7598,251.2619) .. controls (48.7598,250.5411) and (48.1755,249.9568) ..(47.4547,249.9568) .. controls (46.7339,249.9568) and (46.1496,250.5411) ..(46.1496,251.2619) .. controls (46.1496,251.9828) and (46.7339,252.5671) ..(47.4547,252.5671) .. controls (48.1755,252.5671) and (48.7598,251.9828) ..(48.7598,251.2619);
\path[fill=chfourfig5c0,fill opacity=0.5,draw=chfourfig5c0,line width=0.3985pt,draw opacity=0.5] (134.1208,222.2703) .. controls (134.1208,221.5495) and (133.5365,220.9651) ..(132.8157,220.9651) .. controls (132.0949,220.9651) and (131.5106,221.5495) ..(131.5106,222.2703) .. controls (131.5106,222.9911) and (132.0949,223.5754) ..(132.8157,223.5754) .. controls (133.5365,223.5754) and (134.1208,222.9911) ..(134.1208,222.2703);
\path[fill=chfourfig5c0,fill opacity=0.5,draw=chfourfig5c0,line width=0.3985pt,draw opacity=0.5] (81.8960,221.9116) .. controls (81.8960,221.1908) and (81.3116,220.6065) ..(80.5909,220.6065) .. controls (79.8700,220.6065) and (79.2857,221.1908) ..(79.2857,221.9116) .. controls (79.2857,222.6324) and (79.8700,223.2167) ..(80.5909,223.2167) .. controls (81.3116,223.2167) and (81.8960,222.6324) ..(81.8960,221.9116);
\path[fill=chfourfig5c0,fill opacity=0.5,draw=chfourfig5c0,line width=0.3985pt,draw opacity=0.5] (100.4965,180.7554) .. controls (100.4965,180.0346) and (99.9121,179.4503) ..(99.1913,179.4503) .. controls (98.4705,179.4503) and (97.8862,180.0346) ..(97.8862,180.7554) .. controls (97.8862,181.4762) and (98.4705,182.0605) ..(99.1913,182.0605) .. controls (99.9121,182.0605) and (100.4965,181.4762) ..(100.4965,180.7554);
\path[fill=chfourfig5c0,fill opacity=0.5,draw=chfourfig5c0,line width=0.3985pt,draw opacity=0.5] (115.2015,206.3298) .. controls (115.2015,205.6090) and (114.6172,205.0247) ..(113.8964,205.0247) .. controls (113.1756,205.0247) and (112.5913,205.6090) ..(112.5913,206.3298) .. controls (112.5913,207.0507) and (113.1756,207.6349) ..(113.8964,207.6349) .. controls (114.6172,207.6349) and (115.2015,207.0507) ..(115.2015,206.3298);
\path[fill=chfourfig5c0,fill opacity=0.5,draw=chfourfig5c0,line width=0.3985pt,draw opacity=0.5] (66.5035,182.3793) .. controls (66.5035,181.6586) and (65.9192,181.0742) ..(65.1984,181.0742) .. controls (64.4776,181.0742) and (63.8932,181.6586) ..(63.8932,182.3793) .. controls (63.8932,183.1002) and (64.4776,183.6845) ..(65.1984,183.6845) .. controls (65.9192,183.6845) and (66.5035,183.1002) ..(66.5035,182.3793);
\path[fill=chfourfig5c0,fill opacity=0.5,draw=chfourfig5c0,line width=0.3985pt,draw opacity=0.5] (117.1243,189.7119) .. controls (117.1243,188.9911) and (116.5400,188.4068) ..(115.8192,188.4068) .. controls (115.0984,188.4068) and (114.5141,188.9911) ..(114.5141,189.7119) .. controls (114.5141,190.4328) and (115.0984,191.0171) ..(115.8192,191.0171) .. controls (116.5400,191.0171) and (117.1243,190.4328) ..(117.1243,189.7119);
\path[fill=chfourfig5c0,fill opacity=0.5,draw=chfourfig5c0,line width=0.3985pt,draw opacity=0.5] (50.4236,244.3876) .. controls (50.4236,243.6668) and (49.8393,243.0825) ..(49.1185,243.0825) .. controls (48.3977,243.0825) and (47.8133,243.6668) ..(47.8133,244.3876) .. controls (47.8133,245.1084) and (48.3977,245.6927) ..(49.1185,245.6927) .. controls (49.8393,245.6927) and (50.4236,245.1084) ..(50.4236,244.3876);
\path[fill=chfourfig5c0,fill opacity=0.5,draw=chfourfig5c0,line width=0.3985pt,draw opacity=0.5] (93.1141,232.4223) .. controls (93.1141,231.7015) and (92.5297,231.1172) ..(91.8089,231.1172) .. controls (91.0881,231.1172) and (90.5038,231.7015) ..(90.5038,232.4223) .. controls (90.5038,233.1431) and (91.0881,233.7275) ..(91.8089,233.7275) .. controls (92.5297,233.7275) and (93.1141,233.1431) ..(93.1141,232.4223);
\path[fill=chfourfig5c0,fill opacity=0.5,draw=chfourfig5c0,line width=0.3985pt,draw opacity=0.5] (66.7227,246.7986) .. controls (66.7227,246.0778) and (66.1384,245.4935) ..(65.4176,245.4935) .. controls (64.6967,245.4935) and (64.1124,246.0778) ..(64.1124,246.7986) .. controls (64.1124,247.5194) and (64.6967,248.1038) ..(65.4176,248.1038) .. controls (66.1384,248.1038) and (66.7227,247.5194) ..(66.7227,246.7986);
\path[fill=chfourfig5c0,fill opacity=0.5,draw=chfourfig5c0,line width=0.3985pt,draw opacity=0.5] (80.4215,210.2651) .. controls (80.4215,209.5443) and (79.8372,208.9600) ..(79.1164,208.9600) .. controls (78.3956,208.9600) and (77.8112,209.5443) ..(77.8112,210.2651) .. controls (77.8112,210.9860) and (78.3956,211.5703) ..(79.1164,211.5703) .. controls (79.8372,211.5703) and (80.4215,210.9860) ..(80.4215,210.2651);
\path[fill=chfourfig5c0,fill opacity=0.5,draw=chfourfig5c0,line width=0.3985pt,draw opacity=0.5] (153.1995,194.1155) .. controls (153.1995,193.3947) and (152.6152,192.8104) ..(151.8944,192.8104) .. controls (151.1736,192.8104) and (150.5893,193.3947) ..(150.5893,194.1155) .. controls (150.5893,194.8363) and (151.1736,195.4206) ..(151.8944,195.4206) .. controls (152.6152,195.4206) and (153.1995,194.8363) ..(153.1995,194.1155);
\path[fill=chfourfig5c0,fill opacity=0.5,draw=chfourfig5c0,line width=0.3985pt,draw opacity=0.5] (157.8222,259.6705) .. controls (157.8222,258.9497) and (157.2379,258.3654) ..(156.5171,258.3654) .. controls (155.7963,258.3654) and (155.2120,258.9497) ..(155.2120,259.6705) .. controls (155.2120,260.3913) and (155.7963,260.9756) ..(156.5171,260.9756) .. controls (157.2379,260.9756) and (157.8222,260.3913) ..(157.8222,259.6705);
\path[fill=chfourfig5c0,fill opacity=0.5,draw=chfourfig5c0,line width=0.3985pt,draw opacity=0.5] (149.4137,236.5270) .. controls (149.4137,235.8062) and (148.8293,235.2219) ..(148.1085,235.2219) .. controls (147.3877,235.2219) and (146.8034,235.8062) ..(146.8034,236.5270) .. controls (146.8034,237.2478) and (147.3877,237.8321) ..(148.1085,237.8321) .. controls (148.8293,237.8321) and (149.4137,237.2478) ..(149.4137,236.5270);
\path[fill=chfourfig5c0,fill opacity=0.5,draw=chfourfig5c0,line width=0.3985pt,draw opacity=0.5] (126.9974,245.2245) .. controls (126.9974,244.5037) and (126.4131,243.9194) ..(125.6923,243.9194) .. controls (124.9715,243.9194) and (124.3872,244.5037) ..(124.3872,245.2245) .. controls (124.3872,245.9453) and (124.9715,246.5296) ..(125.6923,246.5296) .. controls (126.4131,246.5296) and (126.9974,245.9453) ..(126.9974,245.2245);
\path[fill=chfourfig5c0,fill opacity=0.5,draw=chfourfig5c0,line width=0.3985pt,draw opacity=0.5] (44.5555,225.9565) .. controls (44.5555,225.2357) and (43.9712,224.6514) ..(43.2504,224.6514) .. controls (42.5296,224.6514) and (41.9453,225.2357) ..(41.9453,225.9565) .. controls (41.9453,226.6773) and (42.5296,227.2616) ..(43.2504,227.2616) .. controls (43.9712,227.2616) and (44.5555,226.6773) ..(44.5555,225.9565);
\path[fill=chfourfig5c0,fill opacity=0.5,draw=chfourfig5c0,line width=0.3985pt,draw opacity=0.5] (55.9031,225.2890) .. controls (55.9031,224.5682) and (55.3188,223.9839) ..(54.5980,223.9839) .. controls (53.8772,223.9839) and (53.2929,224.5682) ..(53.2929,225.2890) .. controls (53.2929,226.0098) and (53.8772,226.5941) ..(54.5980,226.5941) .. controls (55.3188,226.5941) and (55.9031,226.0098) ..(55.9031,225.2890);
\path[fill=chfourfig5c0,fill opacity=0.5,draw=chfourfig5c0,line width=0.3985pt,draw opacity=0.5] (114.3248,264.1239) .. controls (114.3248,263.4031) and (113.7405,262.8188) ..(113.0197,262.8188) .. controls (112.2989,262.8188) and (111.7145,263.4031) ..(111.7145,264.1239) .. controls (111.7145,264.8447) and (112.2989,265.4290) ..(113.0197,265.4290) .. controls (113.7405,265.4290) and (114.3248,264.8447) ..(114.3248,264.1239);
\path[fill=chfourfig5c0,fill opacity=0.5,draw=chfourfig5c0,line width=0.3985pt,draw opacity=0.5] (50.4236,176.4117) .. controls (50.4236,175.6909) and (49.8393,175.1065) ..(49.1185,175.1065) .. controls (48.3977,175.1065) and (47.8133,175.6909) ..(47.8133,176.4117) .. controls (47.8133,177.1325) and (48.3977,177.7168) ..(49.1185,177.7168) .. controls (49.8393,177.7168) and (50.4236,177.1325) ..(50.4236,176.4117);
\path[fill=chfourfig5c0,fill opacity=0.5,draw=chfourfig5c0,line width=0.3985pt,draw opacity=0.5] (87.8736,259.1724) .. controls (87.8736,258.4516) and (87.2893,257.8672) ..(86.5685,257.8672) .. controls (85.8477,257.8672) and (85.2634,258.4516) ..(85.2634,259.1724) .. controls (85.2634,259.8932) and (85.8477,260.4775) ..(86.5685,260.4775) .. controls (87.2893,260.4775) and (87.8736,259.8932) ..(87.8736,259.1724);
\path[fill=chfourfig5c0,fill opacity=0.5,draw=chfourfig5c0,line width=0.3985pt,draw opacity=0.5] (98.3246,191.7643) .. controls (98.3246,191.0435) and (97.7403,190.4592) ..(97.0195,190.4592) .. controls (96.2987,190.4592) and (95.7143,191.0435) ..(95.7143,191.7643) .. controls (95.7143,192.4851) and (96.2987,193.0694) ..(97.0195,193.0694) .. controls (97.7403,193.0694) and (98.3246,192.4851) ..(98.3246,191.7643);
\path[fill=chfourfig5c0,fill opacity=0.5,draw=chfourfig5c0,line width=0.3985pt,draw opacity=0.5] (115.0620,178.0555) .. controls (115.0620,177.3347) and (114.4777,176.7504) ..(113.7569,176.7504) .. controls (113.0361,176.7504) and (112.4518,177.3347) ..(112.4518,178.0555) .. controls (112.4518,178.7763) and (113.0361,179.3606) ..(113.7569,179.3606) .. controls (114.4777,179.3606) and (115.0620,178.7763) ..(115.0620,178.0555);
\path[fill=chfourfig5c0,fill opacity=0.5,draw=chfourfig5c0,line width=0.3985pt,draw opacity=0.5] (71.8934,219.2715) .. controls (71.8934,218.5507) and (71.3090,217.9664) ..(70.5882,217.9664) .. controls (69.8674,217.9664) and (69.2831,218.5507) ..(69.2831,219.2715) .. controls (69.2831,219.9923) and (69.8674,220.5766) ..(70.5882,220.5766) .. controls (71.3090,220.5766) and (71.8934,219.9923) ..(71.8934,219.2715);
\path[fill=chfourfig5c0,fill opacity=0.5,draw=chfourfig5c0,line width=0.3985pt,draw opacity=0.5] (148.5469,198.8976) .. controls (148.5469,198.1768) and (147.9626,197.5925) ..(147.2418,197.5925) .. controls (146.5210,197.5925) and (145.9367,198.1768) ..(145.9367,198.8976) .. controls (145.9367,199.6184) and (146.5210,200.2028) ..(147.2418,200.2028) .. controls (147.9626,200.2028) and (148.5469,199.6184) ..(148.5469,198.8976);
\path[fill=chfourfig5c0,fill opacity=0.5,draw=chfourfig5c0,line width=0.3985pt,draw opacity=0.5] (141.8419,203.9886) .. controls (141.8419,203.2678) and (141.2576,202.6835) ..(140.5368,202.6835) .. controls (139.8160,202.6835) and (139.2317,203.2678) ..(139.2317,203.9886) .. controls (139.2317,204.7094) and (139.8160,205.2937) ..(140.5368,205.2937) .. controls (141.2576,205.2937) and (141.8419,204.7094) ..(141.8419,203.9886);
\path[fill=chfourfig5c0,fill opacity=0.5,draw=chfourfig5c0,line width=0.3985pt,draw opacity=0.5] (148.0886,226.8432) .. controls (148.0886,226.1224) and (147.5043,225.5381) ..(146.7835,225.5381) .. controls (146.0627,225.5381) and (145.4784,226.1224) ..(145.4784,226.8432) .. controls (145.4784,227.5640) and (146.0627,228.1483) ..(146.7835,228.1483) .. controls (147.5043,228.1483) and (148.0886,227.5640) ..(148.0886,226.8432);
\path[fill=chfourfig5c0,fill opacity=0.5,draw=chfourfig5c0,line width=0.3985pt,draw opacity=0.5] (107.8689,224.7610) .. controls (107.8689,224.0402) and (107.2846,223.4558) ..(106.5638,223.4558) .. controls (105.8430,223.4558) and (105.2587,224.0402) ..(105.2587,224.7610) .. controls (105.2587,225.4818) and (105.8430,226.0661) ..(106.5638,226.0661) .. controls (107.2846,226.0661) and (107.8689,225.4818) ..(107.8689,224.7610);
\path[fill=chfourfig5c0,fill opacity=0.5,draw=chfourfig5c0,line width=0.3985pt,draw opacity=0.5] (119.5353,175.3158) .. controls (119.5353,174.5950) and (118.9510,174.0106) ..(118.2302,174.0106) .. controls (117.5094,174.0106) and (116.9251,174.5950) ..(116.9251,175.3158) .. controls (116.9251,176.0366) and (117.5094,176.6209) ..(118.2302,176.6209) .. controls (118.9510,176.6209) and (119.5353,176.0366) ..(119.5353,175.3158);
\path[fill=chfourfig5c0,fill opacity=0.5,draw=chfourfig5c0,line width=0.3985pt,draw opacity=0.5] (106.9723,205.4531) .. controls (106.9723,204.7323) and (106.3879,204.1480) ..(105.6671,204.1480) .. controls (104.9463,204.1480) and (104.3620,204.7323) ..(104.3620,205.4531) .. controls (104.3620,206.1740) and (104.9463,206.7582) ..(105.6671,206.7582) .. controls (106.3879,206.7582) and (106.9723,206.1740) ..(106.9723,205.4531);
\path[fill=chfourfig5c0,fill opacity=0.5,draw=chfourfig5c0,line width=0.3985pt,draw opacity=0.5] (92.5661,193.4679) .. controls (92.5661,192.7471) and (91.9818,192.1628) ..(91.2610,192.1628) .. controls (90.5402,192.1628) and (89.9559,192.7471) ..(89.9559,193.4679) .. controls (89.9559,194.1887) and (90.5402,194.7730) ..(91.2610,194.7730) .. controls (91.9818,194.7730) and (92.5661,194.1887) ..(92.5661,193.4679);
\path[fill=chfourfig5c0,fill opacity=0.5,draw=chfourfig5c0,line width=0.3985pt,draw opacity=0.5] (105.1491,188.1279) .. controls (105.1491,187.4071) and (104.5647,186.8228) ..(103.8439,186.8228) .. controls (103.1231,186.8228) and (102.5388,187.4071) ..(102.5388,188.1279) .. controls (102.5388,188.8487) and (103.1231,189.4330) ..(103.8439,189.4330) .. controls (104.5647,189.4330) and (105.1491,188.8487) ..(105.1491,188.1279);
\path[fill=chfourfig5c0,fill opacity=0.5,draw=chfourfig5c0,line width=0.3985pt,draw opacity=0.5] (74.0951,232.6316) .. controls (74.0951,231.9108) and (73.5108,231.3264) ..(72.7900,231.3264) .. controls (72.0692,231.3264) and (71.4849,231.9108) ..(71.4849,232.6316) .. controls (71.4849,233.3524) and (72.0692,233.9367) ..(72.7900,233.9367) .. controls (73.5108,233.9367) and (74.0951,233.3524) ..(74.0951,232.6316);
\path[fill=chfourfig5c0,fill opacity=0.5,draw=chfourfig5c0,line width=0.3985pt,draw opacity=0.5] (155.5109,236.3776) .. controls (155.5109,235.6567) and (154.9265,235.0724) ..(154.2057,235.0724) .. controls (153.4849,235.0724) and (152.9006,235.6567) ..(152.9006,236.3776) .. controls (152.9006,237.0984) and (153.4849,237.6827) ..(154.2057,237.6827) .. controls (154.9265,237.6827) and (155.5109,237.0984) ..(155.5109,236.3776);
\path[fill=chfourfig5c0,fill opacity=0.5,draw=chfourfig5c0,line width=0.3985pt,draw opacity=0.5] (152.9305,243.9592) .. controls (152.9305,243.2384) and (152.3462,242.6541) ..(151.6254,242.6541) .. controls (150.9046,242.6541) and (150.3203,243.2384) ..(150.3203,243.9592) .. controls (150.3203,244.6800) and (150.9046,245.2643) ..(151.6254,245.2643) .. controls (152.3462,245.2643) and (152.9305,244.6800) ..(152.9305,243.9592);
\path[fill=chfourfig5c0,fill opacity=0.5,draw=chfourfig5c0,line width=0.3985pt,draw opacity=0.5] (56.4012,221.9116) .. controls (56.4012,221.1908) and (55.8169,220.6065) ..(55.0961,220.6065) .. controls (54.3753,220.6065) and (53.7910,221.1908) ..(53.7910,221.9116) .. controls (53.7910,222.6324) and (54.3753,223.2167) ..(55.0961,223.2167) .. controls (55.8169,223.2167) and (56.4012,222.6324) ..(56.4012,221.9116);
\path[fill=chfourfig5c0,fill opacity=0.5,draw=chfourfig5c0,line width=0.3985pt,draw opacity=0.5] (72.8398,225.0300) .. controls (72.8398,224.3091) and (72.2555,223.7249) ..(71.5347,223.7249) .. controls (70.8139,223.7249) and (70.2296,224.3091) ..(70.2296,225.0300) .. controls (70.2296,225.7508) and (70.8139,226.3351) ..(71.5347,226.3351) .. controls (72.2555,226.3351) and (72.8398,225.7508) ..(72.8398,225.0300);
\path[fill=chfourfig5c0,fill opacity=0.5,draw=chfourfig5c0,line width=0.3985pt,draw opacity=0.5] (45.9702,209.3386) .. controls (45.9702,208.6178) and (45.3859,208.0335) ..(44.6651,208.0335) .. controls (43.9443,208.0335) and (43.3600,208.6178) ..(43.3600,209.3386) .. controls (43.3600,210.0594) and (43.9443,210.6437) ..(44.6651,210.6437) .. controls (45.3859,210.6437) and (45.9702,210.0594) ..(45.9702,209.3386);
\path[fill=chfourfig5c0,fill opacity=0.5,draw=chfourfig5c0,line width=0.3985pt,draw opacity=0.5] (51.2804,212.8555) .. controls (51.2804,212.1346) and (50.6961,211.5503) ..(49.9753,211.5503) .. controls (49.2545,211.5503) and (48.6701,212.1346) ..(48.6701,212.8555) .. controls (48.6701,213.5763) and (49.2545,214.1606) ..(49.9753,214.1606) .. controls (50.6961,214.1606) and (51.2804,213.5763) ..(51.2804,212.8555);
\path[fill=chfourfig5c0,fill opacity=0.5,draw=chfourfig5c0,line width=0.3985pt,draw opacity=0.5] (101.5226,231.5356) .. controls (101.5226,230.8148) and (100.9383,230.2305) ..(100.2175,230.2305) .. controls (99.4967,230.2305) and (98.9124,230.8148) ..(98.9124,231.5356) .. controls (98.9124,232.2565) and (99.4967,232.8408) ..(100.2175,232.8408) .. controls (100.9383,232.8408) and (101.5226,232.2565) ..(101.5226,231.5356);
\path[fill=chfourfig5c0,fill opacity=0.5,draw=chfourfig5c0,line width=0.3985pt,draw opacity=0.5] (57.9754,251.5907) .. controls (57.9754,250.8699) and (57.3911,250.2856) ..(56.6702,250.2856) .. controls (55.9494,250.2856) and (55.3651,250.8699) ..(55.3651,251.5907) .. controls (55.3651,252.3115) and (55.9494,252.8958) ..(56.6702,252.8958) .. controls (57.3911,252.8958) and (57.9754,252.3115) ..(57.9754,251.5907);
\path[fill=chfourfig5c0,fill opacity=0.5,draw=chfourfig5c0,line width=0.3985pt,draw opacity=0.5] (52.8844,224.6215) .. controls (52.8844,223.9007) and (52.3001,223.3164) ..(51.5793,223.3164) .. controls (50.8585,223.3164) and (50.2742,223.9007) ..(50.2742,224.6215) .. controls (50.2742,225.3423) and (50.8585,225.9266) ..(51.5793,225.9266) .. controls (52.3001,225.9266) and (52.8844,225.3423) ..(52.8844,224.6215);
\path[fill=chfourfig5c0,fill opacity=0.5,draw=chfourfig5c0,line width=0.3985pt,draw opacity=0.5] (148.4273,255.6655) .. controls (148.4273,254.9447) and (147.8430,254.3604) ..(147.1222,254.3604) .. controls (146.4014,254.3604) and (145.8171,254.9447) ..(145.8171,255.6655) .. controls (145.8171,256.3863) and (146.4014,256.9706) ..(147.1222,256.9706) .. controls (147.8430,256.9706) and (148.4273,256.3863) ..(148.4273,255.6655);
\path[fill=chfourfig5c0,fill opacity=0.5,draw=chfourfig5c0,line width=0.3985pt,draw opacity=0.5] (104.6908,223.9640) .. controls (104.6908,223.2431) and (104.1065,222.6588) ..(103.3857,222.6588) .. controls (102.6649,222.6588) and (102.0805,223.2431) ..(102.0805,223.9640) .. controls (102.0805,224.6848) and (102.6649,225.2691) ..(103.3857,225.2691) .. controls (104.1065,225.2691) and (104.6908,224.6848) ..(104.6908,223.9640);
\path[fill=chfourfig5c0,fill opacity=0.5,draw=chfourfig5c0,line width=0.3985pt,draw opacity=0.5] (133.8917,214.3499) .. controls (133.8917,213.6291) and (133.3073,213.0447) ..(132.5865,213.0447) .. controls (131.8657,213.0447) and (131.2814,213.6291) ..(131.2814,214.3499) .. controls (131.2814,215.0707) and (131.8657,215.6550) ..(132.5865,215.6550) .. controls (133.3073,215.6550) and (133.8917,215.0707) ..(133.8917,214.3499);
\path[fill=chfourfig5c0,fill opacity=0.5,draw=chfourfig5c0,line width=0.3985pt,draw opacity=0.5] (80.0827,214.8381) .. controls (80.0827,214.1172) and (79.4984,213.5330) ..(78.7776,213.5330) .. controls (78.0568,213.5330) and (77.4725,214.1172) ..(77.4725,214.8381) .. controls (77.4725,215.5588) and (78.0568,216.1432) ..(78.7776,216.1432) .. controls (79.4984,216.1432) and (80.0827,215.5588) ..(80.0827,214.8381);
\path[fill=chfourfig5c0,fill opacity=0.5,draw=chfourfig5c0,line width=0.3985pt,draw opacity=0.5] (94.7380,250.1262) .. controls (94.7380,249.4054) and (94.1537,248.8211) ..(93.4329,248.8211) .. controls (92.7120,248.8211) and (92.1277,249.4054) ..(92.1277,250.1262) .. controls (92.1277,250.8470) and (92.7120,251.4313) ..(93.4329,251.4313) .. controls (94.1537,251.4313) and (94.7380,250.8470) ..(94.7380,250.1262);
\path[fill=chfourfig5c0,fill opacity=0.5,draw=chfourfig5c0,line width=0.3985pt,draw opacity=0.5] (130.7434,229.4236) .. controls (130.7434,228.7028) and (130.1591,228.1184) ..(129.4383,228.1184) .. controls (128.7175,228.1184) and (128.1332,228.7028) ..(128.1332,229.4236) .. controls (128.1332,230.1443) and (128.7175,230.7287) ..(129.4383,230.7287) .. controls (130.1591,230.7287) and (130.7434,230.1443) ..(130.7434,229.4236);
\path[fill=chfourfig5c0,fill opacity=0.5,draw=chfourfig5c0,line width=0.3985pt,draw opacity=0.5] (97.4777,231.2467) .. controls (97.4777,230.5259) and (96.8934,229.9416) ..(96.1726,229.9416) .. controls (95.4518,229.9416) and (94.8675,230.5259) ..(94.8675,231.2467) .. controls (94.8675,231.9675) and (95.4518,232.5519) ..(96.1726,232.5519) .. controls (96.8934,232.5519) and (97.4777,231.9675) ..(97.4777,231.2467);
\path[fill=chfourfig5c0,fill opacity=0.5,draw=chfourfig5c0,line width=0.3985pt,draw opacity=0.5] (86.6482,198.3297) .. controls (86.6482,197.6089) and (86.0639,197.0246) ..(85.3431,197.0246) .. controls (84.6223,197.0246) and (84.0380,197.6089) ..(84.0380,198.3297) .. controls (84.0380,199.0505) and (84.6223,199.6349) ..(85.3431,199.6349) .. controls (86.0639,199.6349) and (86.6482,199.0505) ..(86.6482,198.3297);
\path[fill=chfourfig5c0,fill opacity=0.5,draw=chfourfig5c0,line width=0.3985pt,draw opacity=0.5] (64.1224,244.6267) .. controls (64.1224,243.9059) and (63.5381,243.3216) ..(62.8173,243.3216) .. controls (62.0965,243.3216) and (61.5121,243.9059) ..(61.5121,244.6267) .. controls (61.5121,245.3475) and (62.0965,245.9319) ..(62.8173,245.9319) .. controls (63.5381,245.9319) and (64.1224,245.3475) ..(64.1224,244.6267);
\path[fill=chfourfig5c0,fill opacity=0.5,draw=chfourfig5c0,line width=0.3985pt,draw opacity=0.5] (162.6940,197.0246) .. controls (162.6940,196.3038) and (162.1097,195.7195) ..(161.3889,195.7195) .. controls (160.6681,195.7195) and (160.0838,196.3038) ..(160.0838,197.0246) .. controls (160.0838,197.7454) and (160.6681,198.3297) ..(161.3889,198.3297) .. controls (162.1097,198.3297) and (162.6940,197.7454) ..(162.6940,197.0246);
\path[fill=chfourfig5c0,fill opacity=0.5,draw=chfourfig5c0,line width=0.3985pt,draw opacity=0.5] (69.6318,251.5608) .. controls (69.6318,250.8400) and (69.0475,250.2557) ..(68.3267,250.2557) .. controls (67.6059,250.2557) and (67.0215,250.8400) ..(67.0215,251.5608) .. controls (67.0215,252.2816) and (67.6059,252.8659) ..(68.3267,252.8659) .. controls (69.0475,252.8659) and (69.6318,252.2816) ..(69.6318,251.5608);
\path[fill=chfourfig5c0,fill opacity=0.5,draw=chfourfig5c0,line width=0.3985pt,draw opacity=0.5] (135.5455,239.8147) .. controls (135.5455,239.0939) and (134.9612,238.5096) ..(134.2404,238.5096) .. controls (133.5196,238.5096) and (132.9352,239.0939) ..(132.9352,239.8147) .. controls (132.9352,240.5355) and (133.5196,241.1198) ..(134.2404,241.1198) .. controls (134.9612,241.1198) and (135.5455,240.5355) ..(135.5455,239.8147);
\path[fill=chfourfig5c0,fill opacity=0.5,draw=chfourfig5c0,line width=0.3985pt,draw opacity=0.5] (78.8474,195.4406) .. controls (78.8474,194.7197) and (78.2630,194.1354) ..(77.5422,194.1354) .. controls (76.8214,194.1354) and (76.2371,194.7197) ..(76.2371,195.4406) .. controls (76.2371,196.1613) and (76.8214,196.7457) ..(77.5422,196.7457) .. controls (78.2630,196.7457) and (78.8474,196.1613) ..(78.8474,195.4406);
\path[fill=chfourfig5c0,fill opacity=0.5,draw=chfourfig5c0,line width=0.3985pt,draw opacity=0.5] (125.1045,227.8395) .. controls (125.1045,227.1187) and (124.5202,226.5344) ..(123.7994,226.5344) .. controls (123.0786,226.5344) and (122.4943,227.1187) ..(122.4943,227.8395) .. controls (122.4943,228.5603) and (123.0786,229.1446) ..(123.7994,229.1446) .. controls (124.5202,229.1446) and (125.1045,228.5603) ..(125.1045,227.8395);
\path[fill=chfourfig5c0,fill opacity=0.5,draw=chfourfig5c0,line width=0.3985pt,draw opacity=0.5] (84.8649,223.2068) .. controls (84.8649,222.4860) and (84.2805,221.9017) ..(83.5598,221.9017) .. controls (82.8389,221.9017) and (82.2546,222.4860) ..(82.2546,223.2068) .. controls (82.2546,223.9276) and (82.8389,224.5119) ..(83.5598,224.5119) .. controls (84.2805,224.5119) and (84.8649,223.9276) ..(84.8649,223.2068);
\path[fill=chfourfig5c0,fill opacity=0.5,draw=chfourfig5c0,line width=0.3985pt,draw opacity=0.5] (118.9176,253.4538) .. controls (118.9176,252.7329) and (118.3333,252.1486) ..(117.6125,252.1486) .. controls (116.8917,252.1486) and (116.3074,252.7329) ..(116.3074,253.4538) .. controls (116.3074,254.1746) and (116.8917,254.7589) ..(117.6125,254.7589) .. controls (118.3333,254.7589) and (118.9176,254.1746) ..(118.9176,253.4538);
\path[fill=chfourfig5c0,fill opacity=0.5,draw=chfourfig5c0,line width=0.3985pt,draw opacity=0.5] (60.4461,214.3897) .. controls (60.4461,213.6689) and (59.8618,213.0846) ..(59.1410,213.0846) .. controls (58.4202,213.0846) and (57.8359,213.6689) ..(57.8359,214.3897) .. controls (57.8359,215.1105) and (58.4202,215.6949) ..(59.1410,215.6949) .. controls (59.8618,215.6949) and (60.4461,215.1105) ..(60.4461,214.3897);
\path[fill=chfourfig5c0,fill opacity=0.5,draw=chfourfig5c0,line width=0.3985pt,draw opacity=0.5] (97.8563,244.8957) .. controls (97.8563,244.1749) and (97.2720,243.5906) ..(96.5512,243.5906) .. controls (95.8304,243.5906) and (95.2461,244.1749) ..(95.2461,244.8957) .. controls (95.2461,245.6165) and (95.8304,246.2009) ..(96.5512,246.2009) .. controls (97.2720,246.2009) and (97.8563,245.6165) ..(97.8563,244.8957);
\path[fill=chfourfig5c0,fill opacity=0.5,draw=chfourfig5c0,line width=0.3985pt,draw opacity=0.5] (79.3256,203.8392) .. controls (79.3256,203.1183) and (78.7413,202.5340) ..(78.0204,202.5340) .. controls (77.2997,202.5340) and (76.7153,203.1183) ..(76.7153,203.8392) .. controls (76.7153,204.5599) and (77.2997,205.1443) ..(78.0204,205.1443) .. controls (78.7413,205.1443) and (79.3256,204.5599) ..(79.3256,203.8392);
\path[fill=chfourfig5c0,fill opacity=0.5,draw=chfourfig5c0,line width=0.3985pt,draw opacity=0.5] (80.3318,231.7648) .. controls (80.3318,231.0440) and (79.7475,230.4597) ..(79.0267,230.4597) .. controls (78.3059,230.4597) and (77.7216,231.0440) ..(77.7216,231.7648) .. controls (77.7216,232.4856) and (78.3059,233.0699) ..(79.0267,233.0699) .. controls (79.7475,233.0699) and (80.3318,232.4856) ..(80.3318,231.7648);
\path[fill=chfourfig5c0,fill opacity=0.5,draw=chfourfig5c0,line width=0.3985pt,draw opacity=0.5] (160.8509,233.9467) .. controls (160.8509,233.2258) and (160.2666,232.6415) ..(159.5458,232.6415) .. controls (158.8250,232.6415) and (158.2407,233.2258) ..(158.2407,233.9467) .. controls (158.2407,234.6674) and (158.8250,235.2518) ..(159.5458,235.2518) .. controls (160.2666,235.2518) and (160.8509,234.6674) ..(160.8509,233.9467);
\path[fill=chfourfig5c0,fill opacity=0.5,draw=chfourfig5c0,line width=0.3985pt,draw opacity=0.5] (106.2151,184.9796) .. controls (106.2151,184.2588) and (105.6308,183.6745) ..(104.9100,183.6745) .. controls (104.1892,183.6745) and (103.6048,184.2588) ..(103.6048,184.9796) .. controls (103.6048,185.7004) and (104.1892,186.2848) ..(104.9100,186.2848) .. controls (105.6308,186.2848) and (106.2151,185.7004) ..(106.2151,184.9796);
\path[fill=chfourfig5c0,fill opacity=0.5,draw=chfourfig5c0,line width=0.3985pt,draw opacity=0.5] (125.5130,207.9438) .. controls (125.5130,207.2230) and (124.9286,206.6387) ..(124.2078,206.6387) .. controls (123.4870,206.6387) and (122.9027,207.2230) ..(122.9027,207.9438) .. controls (122.9027,208.6646) and (123.4870,209.2490) ..(124.2078,209.2490) .. controls (124.9286,209.2490) and (125.5130,208.6646) ..(125.5130,207.9438);
\path[fill=chfourfig5c0,fill opacity=0.5,draw=chfourfig5c0,line width=0.3985pt,draw opacity=0.5] (55.6939,171.2609) .. controls (55.6939,170.5401) and (55.1096,169.9558) ..(54.3888,169.9558) .. controls (53.6680,169.9558) and (53.0836,170.5401) ..(53.0836,171.2609) .. controls (53.0836,171.9817) and (53.6680,172.5660) ..(54.3888,172.5660) .. controls (55.1096,172.5660) and (55.6939,171.9817) ..(55.6939,171.2609);
\path[fill=chfourfig5c0,fill opacity=0.5,draw=chfourfig5c0,line width=0.3985pt,draw opacity=0.5] (51.7885,220.4272) .. controls (51.7885,219.7064) and (51.2042,219.1221) ..(50.4834,219.1221) .. controls (49.7626,219.1221) and (49.1782,219.7064) ..(49.1782,220.4272) .. controls (49.1782,221.1480) and (49.7626,221.7323) ..(50.4834,221.7323) .. controls (51.2042,221.7323) and (51.7885,221.1480) ..(51.7885,220.4272);
\path[fill=chfourfig5c0,fill opacity=0.5,draw=chfourfig5c0,line width=0.3985pt,draw opacity=0.5] (48.0524,233.4684) .. controls (48.0524,232.7476) and (47.4681,232.1633) ..(46.7473,232.1633) .. controls (46.0265,232.1633) and (45.4422,232.7476) ..(45.4422,233.4684) .. controls (45.4422,234.1892) and (46.0265,234.7736) ..(46.7473,234.7736) .. controls (47.4681,234.7736) and (48.0524,234.1892) ..(48.0524,233.4684);
\path[fill=chfourfig5c0,fill opacity=0.5,draw=chfourfig5c0,line width=0.3985pt,draw opacity=0.5] (120.3722,212.1282) .. controls (120.3722,211.4074) and (119.7879,210.8231) ..(119.0671,210.8231) .. controls (118.3463,210.8231) and (117.7619,211.4074) ..(117.7619,212.1282) .. controls (117.7619,212.8490) and (118.3463,213.4333) ..(119.0671,213.4333) .. controls (119.7879,213.4333) and (120.3722,212.8490) ..(120.3722,212.1282);
\path[fill=chfourfig5c0,fill opacity=0.5,draw=chfourfig5c0,line width=0.3985pt,draw opacity=0.5] (72.7203,178.1153) .. controls (72.7203,177.3945) and (72.1359,176.8102) ..(71.4151,176.8102) .. controls (70.6943,176.8102) and (70.1100,177.3945) ..(70.1100,178.1153) .. controls (70.1100,178.8361) and (70.6943,179.4204) ..(71.4151,179.4204) .. controls (72.1359,179.4204) and (72.7203,178.8361) ..(72.7203,178.1153);
\path[fill=chfourfig5c0,fill opacity=0.5,draw=chfourfig5c0,line width=0.3985pt,draw opacity=0.5] (79.8735,263.6955) .. controls (79.8735,262.9747) and (79.2892,262.3904) ..(78.5684,262.3904) .. controls (77.8476,262.3904) and (77.2633,262.9747) ..(77.2633,263.6955) .. controls (77.2633,264.4163) and (77.8476,265.0006) ..(78.5684,265.0006) .. controls (79.2892,265.0006) and (79.8735,264.4163) ..(79.8735,263.6955);
\path[fill=none,draw=chfourfig5c0,line width=0.3985pt,draw opacity=1.0] (43.2205,271.4365) --(164.6567,271.4365);
\path[fill=none,draw=chfourfig5c0,line width=0.3985pt,draw opacity=1.0] (43.2205,271.4365) --(43.2205,276.3980);
\path[fill=none,draw=chfourfig5c0,line width=0.3985pt,draw opacity=1.0] (83.6992,271.4365) --(83.6992,276.3980);
\path[fill=none,draw=chfourfig5c0,line width=0.3985pt,draw opacity=1.0] (124.1780,271.4365) --(124.1780,276.3980);
\path[fill=none,draw=chfourfig5c0,line width=0.3985pt,draw opacity=1.0] (164.6567,271.4365) --(164.6567,276.3980);
\path[fill=none,draw=chfourfig5c0,line width=0.3985pt,draw opacity=1.0] (38.4085,256.2433) --(38.4085,181.8115);
\path[fill=none,draw=chfourfig5c0,line width=0.3985pt,draw opacity=1.0] (38.4085,256.2433) --(33.4470,256.2433);
\path[fill=none,draw=chfourfig5c0,line width=0.3985pt,draw opacity=1.0] (38.4085,219.0324) --(33.4470,219.0324);
\path[fill=none,draw=chfourfig5c0,line width=0.3985pt,draw opacity=1.0] (38.4085,181.8115) --(33.4470,181.8115);
\path[fill=none,draw=chfourfig5c0,line width=0.3985pt,draw opacity=1.0] (38.4085,271.4365)--(38.4085,167.4153)--(169.3392,167.4153)--(169.3392,271.4365)--cycle;
\path[fill=chfourfig5c0,fill opacity=0.5,draw=chfourfig5c0,line width=0.3985pt,draw opacity=0.5] (266.8747,200.1729) .. controls (266.8747,199.4521) and (266.2904,198.8677) ..(265.5696,198.8677) .. controls (264.8488,198.8677) and (264.2644,199.4521) ..(264.2644,200.1729) .. controls (264.2644,200.8937) and (264.8488,201.4780) ..(265.5696,201.4780) .. controls (266.2904,201.4780) and (266.8747,200.8937) ..(266.8747,200.1729);
\path[fill=chfourfig5c0,fill opacity=0.5,draw=chfourfig5c0,line width=0.3985pt,draw opacity=0.5] (293.7044,226.4447) .. controls (293.7044,225.7239) and (293.1201,225.1396) ..(292.3993,225.1396) .. controls (291.6785,225.1396) and (291.0942,225.7239) ..(291.0942,226.4447) .. controls (291.0942,227.1655) and (291.6785,227.7498) ..(292.3993,227.7498) .. controls (293.1201,227.7498) and (293.7044,227.1655) ..(293.7044,226.4447);
\path[fill=chfourfig5c0,fill opacity=0.5,draw=chfourfig5c0,line width=0.3985pt,draw opacity=0.5] (256.0053,214.1008) .. controls (256.0053,213.3800) and (255.4210,212.7957) ..(254.7002,212.7957) .. controls (253.9794,212.7957) and (253.3951,213.3800) ..(253.3951,214.1008) .. controls (253.3951,214.8216) and (253.9794,215.4059) ..(254.7002,215.4059) .. controls (255.4210,215.4059) and (256.0053,214.8216) ..(256.0053,214.1008);
\path[fill=chfourfig5c0,fill opacity=0.5,draw=chfourfig5c0,line width=0.3985pt,draw opacity=0.5] (292.3794,262.0815) .. controls (292.3794,261.3607) and (291.7950,260.7764) ..(291.0743,260.7764) .. controls (290.3535,260.7764) and (289.7691,261.3607) ..(289.7691,262.0815) .. controls (289.7691,262.8023) and (290.3535,263.3866) ..(291.0743,263.3866) .. controls (291.7950,263.3866) and (292.3794,262.8023) ..(292.3794,262.0815);
\path[fill=chfourfig5c0,fill opacity=0.5,draw=chfourfig5c0,line width=0.3985pt,draw opacity=0.5] (264.4737,196.8851) .. controls (264.4737,196.1643) and (263.8893,195.5800) ..(263.1685,195.5800) .. controls (262.4478,195.5800) and (261.8634,196.1643) ..(261.8634,196.8851) .. controls (261.8634,197.6060) and (262.4478,198.1902) ..(263.1685,198.1902) .. controls (263.8893,198.1902) and (264.4737,197.6060) ..(264.4737,196.8851);
\path[fill=chfourfig5c0,fill opacity=0.5,draw=chfourfig5c0,line width=0.3985pt,draw opacity=0.5] (290.0880,195.5302) .. controls (290.0880,194.8094) and (289.5036,194.2251) ..(288.7828,194.2251) .. controls (288.0620,194.2251) and (287.4777,194.8094) ..(287.4777,195.5302) .. controls (287.4777,196.2510) and (288.0620,196.8353) ..(288.7828,196.8353) .. controls (289.5036,196.8353) and (290.0880,196.2510) ..(290.0880,195.5302);
\path[fill=chfourfig5c0,fill opacity=0.5,draw=chfourfig5c0,line width=0.3985pt,draw opacity=0.5] (288.9123,261.8623) .. controls (288.9123,261.1415) and (288.3280,260.5572) ..(287.6072,260.5572) .. controls (286.8864,260.5572) and (286.3021,261.1415) ..(286.3021,261.8623) .. controls (286.3021,262.5831) and (286.8864,263.1674) ..(287.6072,263.1674) .. controls (288.3280,263.1674) and (288.9123,262.5831) ..(288.9123,261.8623);
\path[fill=chfourfig5c0,fill opacity=0.5,draw=chfourfig5c0,line width=0.3985pt,draw opacity=0.5] (325.0672,201.4381) .. controls (325.0672,200.7173) and (324.4829,200.1330) ..(323.7621,200.1330) .. controls (323.0413,200.1330) and (322.4570,200.7173) ..(322.4570,201.4381) .. controls (322.4570,202.1589) and (323.0413,202.7433) ..(323.7621,202.7433) .. controls (324.4829,202.7433) and (325.0672,202.1589) ..(325.0672,201.4381);
\path[fill=chfourfig5c0,fill opacity=0.5,draw=chfourfig5c0,line width=0.3985pt,draw opacity=0.5] (283.9509,218.0660) .. controls (283.9509,217.3452) and (283.3666,216.7609) ..(282.6458,216.7609) .. controls (281.9250,216.7609) and (281.3406,217.3452) ..(281.3406,218.0660) .. controls (281.3406,218.7868) and (281.9250,219.3711) ..(282.6458,219.3711) .. controls (283.3666,219.3711) and (283.9509,218.7868) ..(283.9509,218.0660);
\path[fill=chfourfig5c0,fill opacity=0.5,draw=chfourfig5c0,line width=0.3985pt,draw opacity=0.5] (343.3688,210.1855) .. controls (343.3688,209.4647) and (342.7845,208.8803) ..(342.0637,208.8803) .. controls (341.3429,208.8803) and (340.7586,209.4647) ..(340.7586,210.1855) .. controls (340.7586,210.9062) and (341.3429,211.4906) ..(342.0637,211.4906) .. controls (342.7845,211.4906) and (343.3688,210.9062) ..(343.3688,210.1855);
\path[fill=chfourfig5c0,fill opacity=0.5,draw=chfourfig5c0,line width=0.3985pt,draw opacity=0.5] (288.4242,231.7648) .. controls (288.4242,231.0440) and (287.8398,230.4597) ..(287.1190,230.4597) .. controls (286.3982,230.4597) and (285.8139,231.0440) ..(285.8139,231.7648) .. controls (285.8139,232.4856) and (286.3982,233.0699) ..(287.1190,233.0699) .. controls (287.8398,233.0699) and (288.4242,232.4856) ..(288.4242,231.7648);
\path[fill=chfourfig5c0,fill opacity=0.5,draw=chfourfig5c0,line width=0.3985pt,draw opacity=0.5] (255.8160,226.2753) .. controls (255.8160,225.5545) and (255.2317,224.9702) ..(254.5109,224.9702) .. controls (253.7901,224.9702) and (253.2058,225.5545) ..(253.2058,226.2753) .. controls (253.2058,226.9961) and (253.7901,227.5804) ..(254.5109,227.5804) .. controls (255.2317,227.5804) and (255.8160,226.9961) ..(255.8160,226.2753);
\path[fill=chfourfig5c0,fill opacity=0.5,draw=chfourfig5c0,line width=0.3985pt,draw opacity=0.5] (306.8951,243.2021) .. controls (306.8951,242.4813) and (306.3108,241.8969) ..(305.5900,241.8969) .. controls (304.8692,241.8969) and (304.2849,242.4813) ..(304.2849,243.2021) .. controls (304.2849,243.9229) and (304.8692,244.5072) ..(305.5900,244.5072) .. controls (306.3108,244.5072) and (306.8951,243.9229) ..(306.8951,243.2021);
\path[fill=chfourfig5c0,fill opacity=0.5,draw=chfourfig5c0,line width=0.3985pt,draw opacity=0.5] (285.7242,241.7076) .. controls (285.7242,240.9868) and (285.1399,240.4025) ..(284.4191,240.4025) .. controls (283.6983,240.4025) and (283.1140,240.9868) ..(283.1140,241.7076) .. controls (283.1140,242.4285) and (283.6983,243.0128) ..(284.4191,243.0128) .. controls (285.1399,243.0128) and (285.7242,242.4285) ..(285.7242,241.7076);
\path[fill=chfourfig5c0,fill opacity=0.5,draw=chfourfig5c0,line width=0.3985pt,draw opacity=0.5] (253.6840,215.8344) .. controls (253.6840,215.1136) and (253.0997,214.5292) ..(252.3789,214.5292) .. controls (251.6581,214.5292) and (251.0737,215.1136) ..(251.0737,215.8344) .. controls (251.0737,216.5551) and (251.6581,217.1395) ..(252.3789,217.1395) .. controls (253.0997,217.1395) and (253.6840,216.5551) ..(253.6840,215.8344);
\path[fill=chfourfig5c0,fill opacity=0.5,draw=chfourfig5c0,line width=0.3985pt,draw opacity=0.5] (298.5862,243.3116) .. controls (298.5862,242.5909) and (298.0019,242.0065) ..(297.2811,242.0065) .. controls (296.5602,242.0065) and (295.9759,242.5909) ..(295.9759,243.3116) .. controls (295.9759,244.0325) and (296.5602,244.6168) ..(297.2811,244.6168) .. controls (298.0019,244.6168) and (298.5862,244.0325) ..(298.5862,243.3116);
\path[fill=chfourfig5c0,fill opacity=0.5,draw=chfourfig5c0,line width=0.3985pt,draw opacity=0.5] (287.4080,231.0475) .. controls (287.4080,230.3267) and (286.8236,229.7424) ..(286.1028,229.7424) .. controls (285.3820,229.7424) and (284.7977,230.3267) ..(284.7977,231.0475) .. controls (284.7977,231.7683) and (285.3820,232.3526) ..(286.1028,232.3526) .. controls (286.8236,232.3526) and (287.4080,231.7683) ..(287.4080,231.0475);
\path[fill=chfourfig5c0,fill opacity=0.5,draw=chfourfig5c0,line width=0.3985pt,draw opacity=0.5] (302.4119,209.4482) .. controls (302.4119,208.7274) and (301.8275,208.1431) ..(301.1068,208.1431) .. controls (300.3860,208.1431) and (299.8016,208.7274) ..(299.8016,209.4482) .. controls (299.8016,210.1690) and (300.3860,210.7533) ..(301.1068,210.7533) .. controls (301.8275,210.7533) and (302.4119,210.1690) ..(302.4119,209.4482);
\path[fill=chfourfig5c0,fill opacity=0.5,draw=chfourfig5c0,line width=0.3985pt,draw opacity=0.5] (311.7868,213.0447) .. controls (311.7868,212.3239) and (311.2025,211.7396) ..(310.4818,211.7396) .. controls (309.7609,211.7396) and (309.1766,212.3239) ..(309.1766,213.0447) .. controls (309.1766,213.7656) and (309.7609,214.3499) ..(310.4818,214.3499) .. controls (311.2025,214.3499) and (311.7868,213.7656) ..(311.7868,213.0447);
\path[fill=chfourfig5c0,fill opacity=0.5,draw=chfourfig5c0,line width=0.3985pt,draw opacity=0.5] (254.7500,223.7049) .. controls (254.7500,222.9841) and (254.1657,222.3998) ..(253.4449,222.3998) .. controls (252.7241,222.3998) and (252.1398,222.9841) ..(252.1398,223.7049) .. controls (252.1398,224.4257) and (252.7241,225.0100) ..(253.4449,225.0100) .. controls (254.1657,225.0100) and (254.7500,224.4257) ..(254.7500,223.7049);
\path[fill=chfourfig5c0,fill opacity=0.5,draw=chfourfig5c0,line width=0.3985pt,draw opacity=0.5] (259.7314,240.7612) .. controls (259.7314,240.0404) and (259.1471,239.4561) ..(258.4263,239.4561) .. controls (257.7054,239.4561) and (257.1212,240.0404) ..(257.1212,240.7612) .. controls (257.1212,241.4820) and (257.7054,242.0663) ..(258.4263,242.0663) .. controls (259.1471,242.0663) and (259.7314,241.4820) ..(259.7314,240.7612);
\path[fill=chfourfig5c0,fill opacity=0.5,draw=chfourfig5c0,line width=0.3985pt,draw opacity=0.5] (303.6871,259.4812) .. controls (303.6871,258.7604) and (303.1028,258.1761) ..(302.3820,258.1761) .. controls (301.6612,258.1761) and (301.0769,258.7604) ..(301.0769,259.4812) .. controls (301.0769,260.2020) and (301.6612,260.7863) ..(302.3820,260.7863) .. controls (303.1028,260.7863) and (303.6871,260.2020) ..(303.6871,259.4812);
\path[fill=chfourfig5c0,fill opacity=0.5,draw=chfourfig5c0,line width=0.3985pt,draw opacity=0.5] (306.4369,253.6430) .. controls (306.4369,252.9222) and (305.8525,252.3379) ..(305.1317,252.3379) .. controls (304.4109,252.3379) and (303.8266,252.9222) ..(303.8266,253.6430) .. controls (303.8266,254.3638) and (304.4109,254.9482) ..(305.1317,254.9482) .. controls (305.8525,254.9482) and (306.4369,254.3638) ..(306.4369,253.6430);
\path[fill=chfourfig5c0,fill opacity=0.5,draw=chfourfig5c0,line width=0.3985pt,draw opacity=0.5] (318.5117,229.3140) .. controls (318.5117,228.5931) and (317.9274,228.0089) ..(317.2066,228.0089) .. controls (316.4858,228.0089) and (315.9015,228.5931) ..(315.9015,229.3140) .. controls (315.9015,230.0348) and (316.4858,230.6191) ..(317.2066,230.6191) .. controls (317.9274,230.6191) and (318.5117,230.0348) ..(318.5117,229.3140);
\path[fill=chfourfig5c0,fill opacity=0.5,draw=chfourfig5c0,line width=0.3985pt,draw opacity=0.5] (278.6706,221.7323) .. controls (278.6706,221.0115) and (278.0863,220.4272) ..(277.3655,220.4272) .. controls (276.6447,220.4272) and (276.0604,221.0115) ..(276.0604,221.7323) .. controls (276.0604,222.4531) and (276.6447,223.0374) ..(277.3655,223.0374) .. controls (278.0863,223.0374) and (278.6706,222.4531) ..(278.6706,221.7323);
\path[fill=chfourfig5c0,fill opacity=0.5,draw=chfourfig5c0,line width=0.3985pt,draw opacity=0.5] (298.6958,259.6606) .. controls (298.6958,258.9398) and (298.1115,258.3554) ..(297.3907,258.3554) .. controls (296.6699,258.3554) and (296.0855,258.9398) ..(296.0855,259.6606) .. controls (296.0855,260.3814) and (296.6699,260.9657) ..(297.3907,260.9657) .. controls (298.1115,260.9657) and (298.6958,260.3814) ..(298.6958,259.6606);
\path[fill=chfourfig5c0,fill opacity=0.5,draw=chfourfig5c0,line width=0.3985pt,draw opacity=0.5] (222.7297,232.6017) .. controls (222.7297,231.8809) and (222.1453,231.2966) ..(221.4246,231.2966) .. controls (220.7037,231.2966) and (220.1194,231.8809) ..(220.1194,232.6017) .. controls (220.1194,233.3225) and (220.7037,233.9068) ..(221.4246,233.9068) .. controls (222.1453,233.9068) and (222.7297,233.3225) ..(222.7297,232.6017);
\path[fill=chfourfig5c0,fill opacity=0.5,draw=chfourfig5c0,line width=0.3985pt,draw opacity=0.5] (249.4000,241.0601) .. controls (249.4000,240.3392) and (248.8157,239.7549) ..(248.0949,239.7549) .. controls (247.3741,239.7549) and (246.7898,240.3392) ..(246.7898,241.0601) .. controls (246.7898,241.7809) and (247.3741,242.3652) ..(248.0949,242.3652) .. controls (248.8157,242.3652) and (249.4000,241.7809) ..(249.4000,241.0601);
\path[fill=chfourfig5c0,fill opacity=0.5,draw=chfourfig5c0,line width=0.3985pt,draw opacity=0.5] (276.0903,206.1904) .. controls (276.0903,205.4695) and (275.5060,204.8853) ..(274.7851,204.8853) .. controls (274.0643,204.8853) and (273.4800,205.4695) ..(273.4800,206.1904) .. controls (273.4800,206.9112) and (274.0643,207.4955) ..(274.7851,207.4955) .. controls (275.5060,207.4955) and (276.0903,206.9112) ..(276.0903,206.1904);
\path[fill=chfourfig5c0,fill opacity=0.5,draw=chfourfig5c0,line width=0.3985pt,draw opacity=0.5] (318.9401,210.3748) .. controls (318.9401,209.6539) and (318.3558,209.0696) ..(317.6350,209.0696) .. controls (316.9142,209.0696) and (316.3299,209.6539) ..(316.3299,210.3748) .. controls (316.3299,211.0956) and (316.9142,211.6799) ..(317.6350,211.6799) .. controls (318.3558,211.6799) and (318.9401,211.0956) ..(318.9401,210.3748);
\path[fill=chfourfig5c0,fill opacity=0.5,draw=chfourfig5c0,line width=0.3985pt,draw opacity=0.5] (238.3712,230.4099) .. controls (238.3712,229.6891) and (237.7869,229.1047) ..(237.0661,229.1047) .. controls (236.3453,229.1047) and (235.7610,229.6891) ..(235.7610,230.4099) .. controls (235.7610,231.1307) and (236.3453,231.7150) ..(237.0661,231.7150) .. controls (237.7869,231.7150) and (238.3712,231.1307) ..(238.3712,230.4099);
\path[fill=chfourfig5c0,fill opacity=0.5,draw=chfourfig5c0,line width=0.3985pt,draw opacity=0.5] (296.1154,206.6885) .. controls (296.1154,205.9677) and (295.5311,205.3834) ..(294.8103,205.3834) .. controls (294.0895,205.3834) and (293.5052,205.9677) ..(293.5052,206.6885) .. controls (293.5052,207.4093) and (294.0895,207.9936) ..(294.8103,207.9936) .. controls (295.5311,207.9936) and (296.1154,207.4093) ..(296.1154,206.6885);
\path[fill=chfourfig5c0,fill opacity=0.5,draw=chfourfig5c0,line width=0.3985pt,draw opacity=0.5] (292.3296,257.7577) .. controls (292.3296,257.0369) and (291.7452,256.4525) ..(291.0244,256.4525) .. controls (290.3036,256.4525) and (289.7193,257.0369) ..(289.7193,257.7577) .. controls (289.7193,258.4785) and (290.3036,259.0628) ..(291.0244,259.0628) .. controls (291.7452,259.0628) and (292.3296,258.4785) ..(292.3296,257.7577);
\path[fill=chfourfig5c0,fill opacity=0.5,draw=chfourfig5c0,line width=0.3985pt,draw opacity=0.5] (270.4613,201.8765) .. controls (270.4613,201.1557) and (269.8770,200.5714) ..(269.1562,200.5714) .. controls (268.4354,200.5714) and (267.8510,201.1557) ..(267.8510,201.8765) .. controls (267.8510,202.5973) and (268.4354,203.1816) ..(269.1562,203.1816) .. controls (269.8770,203.1816) and (270.4613,202.5973) ..(270.4613,201.8765);
\path[fill=chfourfig5c0,fill opacity=0.5,draw=chfourfig5c0,line width=0.3985pt,draw opacity=0.5] (329.9490,210.4943) .. controls (329.9490,209.7735) and (329.3647,209.1891) ..(328.6439,209.1891) .. controls (327.9230,209.1891) and (327.3387,209.7735) ..(327.3387,210.4943) .. controls (327.3387,211.2151) and (327.9230,211.7994) ..(328.6439,211.7994) .. controls (329.3647,211.7994) and (329.9490,211.2151) ..(329.9490,210.4943);
\path[fill=chfourfig5c0,fill opacity=0.5,draw=chfourfig5c0,line width=0.3985pt,draw opacity=0.5] (291.6023,184.9796) .. controls (291.6023,184.2588) and (291.0180,183.6745) ..(290.2972,183.6745) .. controls (289.5764,183.6745) and (288.9920,184.2588) ..(288.9920,184.9796) .. controls (288.9920,185.7004) and (289.5764,186.2848) ..(290.2972,186.2848) .. controls (291.0180,186.2848) and (291.6023,185.7004) ..(291.6023,184.9796);
\path[fill=chfourfig5c0,fill opacity=0.5,draw=chfourfig5c0,line width=0.3985pt,draw opacity=0.5] (283.2236,243.6603) .. controls (283.2236,242.9395) and (282.6393,242.3552) ..(281.9185,242.3552) .. controls (281.1977,242.3552) and (280.6133,242.9395) ..(280.6133,243.6603) .. controls (280.6133,244.3811) and (281.1977,244.9655) ..(281.9185,244.9655) .. controls (282.6393,244.9655) and (283.2236,244.3811) ..(283.2236,243.6603);
\path[fill=chfourfig5c0,fill opacity=0.5,draw=chfourfig5c0,line width=0.3985pt,draw opacity=0.5] (269.5348,191.1466) .. controls (269.5348,190.4258) and (268.9504,189.8415) ..(268.2296,189.8415) .. controls (267.5088,189.8415) and (266.9245,190.4258) ..(266.9245,191.1466) .. controls (266.9245,191.8674) and (267.5088,192.4517) ..(268.2296,192.4517) .. controls (268.9504,192.4517) and (269.5348,191.8674) ..(269.5348,191.1466);
\path[fill=chfourfig5c0,fill opacity=0.5,draw=chfourfig5c0,line width=0.3985pt,draw opacity=0.5] (316.7583,197.7917) .. controls (316.7583,197.0710) and (316.1740,196.4866) ..(315.4532,196.4866) .. controls (314.7324,196.4866) and (314.1480,197.0710) ..(314.1480,197.7917) .. controls (314.1480,198.5126) and (314.7324,199.0969) ..(315.4532,199.0969) .. controls (316.1740,199.0969) and (316.7583,198.5126) ..(316.7583,197.7917);
\path[fill=chfourfig5c0,fill opacity=0.5,draw=chfourfig5c0,line width=0.3985pt,draw opacity=0.5] (307.4033,218.6339) .. controls (307.4033,217.9131) and (306.8189,217.3288) ..(306.0981,217.3288) .. controls (305.3773,217.3288) and (304.7930,217.9131) ..(304.7930,218.6339) .. controls (304.7930,219.3547) and (305.3773,219.9390) ..(306.0981,219.9390) .. controls (306.8189,219.9390) and (307.4033,219.3547) ..(307.4033,218.6339);
\path[fill=chfourfig5c0,fill opacity=0.5,draw=chfourfig5c0,line width=0.3985pt,draw opacity=0.5] (307.9312,191.8340) .. controls (307.9312,191.1132) and (307.3469,190.5289) ..(306.6262,190.5289) .. controls (305.9053,190.5289) and (305.3210,191.1132) ..(305.3210,191.8340) .. controls (305.3210,192.5548) and (305.9053,193.1392) ..(306.6262,193.1392) .. controls (307.3469,193.1392) and (307.9312,192.5548) ..(307.9312,191.8340);
\path[fill=chfourfig5c0,fill opacity=0.5,draw=chfourfig5c0,line width=0.3985pt,draw opacity=0.5] (281.0219,206.1306) .. controls (281.0219,205.4098) and (280.4375,204.8255) ..(279.7167,204.8255) .. controls (278.9959,204.8255) and (278.4116,205.4098) ..(278.4116,206.1306) .. controls (278.4116,206.8514) and (278.9959,207.4357) ..(279.7167,207.4357) .. controls (280.4375,207.4357) and (281.0219,206.8514) ..(281.0219,206.1306);
\path[fill=chfourfig5c0,fill opacity=0.5,draw=chfourfig5c0,line width=0.3985pt,draw opacity=0.5] (250.1373,235.2717) .. controls (250.1373,234.5509) and (249.5529,233.9666) ..(248.8321,233.9666) .. controls (248.1113,233.9666) and (247.5270,234.5509) ..(247.5270,235.2717) .. controls (247.5270,235.9925) and (248.1113,236.5768) ..(248.8321,236.5768) .. controls (249.5529,236.5768) and (250.1373,235.9925) ..(250.1373,235.2717);
\path[fill=chfourfig5c0,fill opacity=0.5,draw=chfourfig5c0,line width=0.3985pt,draw opacity=0.5] (236.1694,221.4633) .. controls (236.1694,220.7425) and (235.5851,220.1582) ..(234.8643,220.1582) .. controls (234.1435,220.1582) and (233.5592,220.7425) ..(233.5592,221.4633) .. controls (233.5592,222.1841) and (234.1435,222.7684) ..(234.8643,222.7684) .. controls (235.5851,222.7684) and (236.1694,222.1841) ..(236.1694,221.4633);
\path[fill=chfourfig5c0,fill opacity=0.5,draw=chfourfig5c0,line width=0.3985pt,draw opacity=0.5] (262.8597,249.0402) .. controls (262.8597,248.3195) and (262.2754,247.7351) ..(261.5546,247.7351) .. controls (260.8338,247.7351) and (260.2495,248.3195) ..(260.2495,249.0402) .. controls (260.2495,249.7610) and (260.8338,250.3454) ..(261.5546,250.3454) .. controls (262.2754,250.3454) and (262.8597,249.7610) ..(262.8597,249.0402);
\path[fill=chfourfig5c0,fill opacity=0.5,draw=chfourfig5c0,line width=0.3985pt,draw opacity=0.5] (269.7639,221.7522) .. controls (269.7639,221.0314) and (269.1796,220.4471) ..(268.4588,220.4471) .. controls (267.7380,220.4471) and (267.1537,221.0314) ..(267.1537,221.7522) .. controls (267.1537,222.4730) and (267.7380,223.0573) ..(268.4588,223.0573) .. controls (269.1796,223.0573) and (269.7639,222.4730) ..(269.7639,221.7522);
\path[fill=chfourfig5c0,fill opacity=0.5,draw=chfourfig5c0,line width=0.3985pt,draw opacity=0.5] (321.6898,237.7923) .. controls (321.6898,237.0715) and (321.1055,236.4872) ..(320.3847,236.4872) .. controls (319.6639,236.4872) and (319.0796,237.0715) ..(319.0796,237.7923) .. controls (319.0796,238.5131) and (319.6639,239.0974) ..(320.3847,239.0974) .. controls (321.1055,239.0974) and (321.6898,238.5131) ..(321.6898,237.7923);
\path[fill=chfourfig5c0,fill opacity=0.5,draw=chfourfig5c0,line width=0.3985pt,draw opacity=0.5] (233.9677,224.3226) .. controls (233.9677,223.6018) and (233.3834,223.0175) ..(232.6625,223.0175) .. controls (231.9417,223.0175) and (231.3574,223.6018) ..(231.3574,224.3226) .. controls (231.3574,225.0434) and (231.9417,225.6277) ..(232.6625,225.6277) .. controls (233.3834,225.6277) and (233.9677,225.0434) ..(233.9677,224.3226);
\path[fill=chfourfig5c0,fill opacity=0.5,draw=chfourfig5c0,line width=0.3985pt,draw opacity=0.5] (290.2473,210.8031) .. controls (290.2473,210.0823) and (289.6630,209.4980) ..(288.9422,209.4980) .. controls (288.2214,209.4980) and (287.6371,210.0823) ..(287.6371,210.8031) .. controls (287.6371,211.5240) and (288.2214,212.1083) ..(288.9422,212.1083) .. controls (289.6630,212.1083) and (290.2473,211.5240) ..(290.2473,210.8031);
\path[fill=chfourfig5c0,fill opacity=0.5,draw=chfourfig5c0,line width=0.3985pt,draw opacity=0.5] (289.7094,192.0432) .. controls (289.7094,191.3224) and (289.1250,190.7381) ..(288.4042,190.7381) .. controls (287.6834,190.7381) and (287.0991,191.3224) ..(287.0991,192.0432) .. controls (287.0991,192.7640) and (287.6834,193.3484) ..(288.4042,193.3484) .. controls (289.1250,193.3484) and (289.7094,192.7640) ..(289.7094,192.0432);
\path[fill=chfourfig5c0,fill opacity=0.5,draw=chfourfig5c0,line width=0.3985pt,draw opacity=0.5] (261.9930,194.2251) .. controls (261.9930,193.5043) and (261.4086,192.9200) ..(260.6878,192.9200) .. controls (259.9670,192.9200) and (259.3827,193.5043) ..(259.3827,194.2251) .. controls (259.3827,194.9459) and (259.9670,195.5302) ..(260.6878,195.5302) .. controls (261.4086,195.5302) and (261.9930,194.9459) ..(261.9930,194.2251);
\path[fill=chfourfig5c0,fill opacity=0.5,draw=chfourfig5c0,line width=0.3985pt,draw opacity=0.5] (242.9939,242.2456) .. controls (242.9939,241.5248) and (242.4096,240.9405) ..(241.6888,240.9405) .. controls (240.9680,240.9405) and (240.3837,241.5248) ..(240.3837,242.2456) .. controls (240.3837,242.9664) and (240.9680,243.5508) ..(241.6888,243.5508) .. controls (242.4096,243.5508) and (242.9939,242.9664) ..(242.9939,242.2456);
\path[fill=chfourfig5c0,fill opacity=0.5,draw=chfourfig5c0,line width=0.3985pt,draw opacity=0.5] (277.0467,229.4136) .. controls (277.0467,228.6928) and (276.4624,228.1085) ..(275.7416,228.1085) .. controls (275.0208,228.1085) and (274.4364,228.6928) ..(274.4364,229.4136) .. controls (274.4364,230.1344) and (275.0208,230.7187) ..(275.7416,230.7187) .. controls (276.4624,230.7187) and (277.0467,230.1344) ..(277.0467,229.4136);
\path[fill=chfourfig5c0,fill opacity=0.5,draw=chfourfig5c0,line width=0.3985pt,draw opacity=0.5] (255.1983,212.4071) .. controls (255.1983,211.6863) and (254.6140,211.1020) ..(253.8932,211.1020) .. controls (253.1724,211.1020) and (252.5881,211.6863) ..(252.5881,212.4071) .. controls (252.5881,213.1279) and (253.1724,213.7123) ..(253.8932,213.7123) .. controls (254.6140,213.7123) and (255.1983,213.1279) ..(255.1983,212.4071);
\path[fill=chfourfig5c0,fill opacity=0.5,draw=chfourfig5c0,line width=0.3985pt,draw opacity=0.5] (331.5131,225.0997) .. controls (331.5131,224.3789) and (330.9288,223.7946) ..(330.2080,223.7946) .. controls (329.4872,223.7946) and (328.9029,224.3789) ..(328.9029,225.0997) .. controls (328.9029,225.8205) and (329.4872,226.4048) ..(330.2080,226.4048) .. controls (330.9288,226.4048) and (331.5131,225.8205) ..(331.5131,225.0997);
\path[fill=chfourfig5c0,fill opacity=0.5,draw=chfourfig5c0,line width=0.3985pt,draw opacity=0.5] (295.2287,201.4581) .. controls (295.2287,200.7372) and (294.6444,200.1529) ..(293.9236,200.1529) .. controls (293.2028,200.1529) and (292.6185,200.7372) ..(292.6185,201.4581) .. controls (292.6185,202.1788) and (293.2028,202.7632) ..(293.9236,202.7632) .. controls (294.6444,202.7632) and (295.2287,202.1788) ..(295.2287,201.4581);
\path[fill=chfourfig5c0,fill opacity=0.5,draw=chfourfig5c0,line width=0.3985pt,draw opacity=0.5] (297.1615,257.5285) .. controls (297.1615,256.8077) and (296.5772,256.2234) ..(295.8564,256.2234) .. controls (295.1356,256.2234) and (294.5513,256.8077) ..(294.5513,257.5285) .. controls (294.5513,258.2493) and (295.1356,258.8336) ..(295.8564,258.8336) .. controls (296.5772,258.8336) and (297.1615,258.2493) ..(297.1615,257.5285);
\path[fill=chfourfig5c0,fill opacity=0.5,draw=chfourfig5c0,line width=0.3985pt,draw opacity=0.5] (302.1529,231.0873) .. controls (302.1529,230.3665) and (301.5685,229.7822) ..(300.8477,229.7822) .. controls (300.1269,229.7822) and (299.5426,230.3665) ..(299.5426,231.0873) .. controls (299.5426,231.8081) and (300.1269,232.3925) ..(300.8477,232.3925) .. controls (301.5685,232.3925) and (302.1529,231.8081) ..(302.1529,231.0873);
\path[fill=chfourfig5c0,fill opacity=0.5,draw=chfourfig5c0,line width=0.3985pt,draw opacity=0.5] (280.4539,187.3707) .. controls (280.4539,186.6499) and (279.8696,186.0656) ..(279.1488,186.0656) .. controls (278.4280,186.0656) and (277.8437,186.6499) ..(277.8437,187.3707) .. controls (277.8437,188.0915) and (278.4280,188.6758) ..(279.1488,188.6758) .. controls (279.8696,188.6758) and (280.4539,188.0915) ..(280.4539,187.3707);
\path[fill=chfourfig5c0,fill opacity=0.5,draw=chfourfig5c0,line width=0.3985pt,draw opacity=0.5] (241.8383,240.5121) .. controls (241.8383,239.7913) and (241.2539,239.2070) ..(240.5331,239.2070) .. controls (239.8123,239.2070) and (239.2280,239.7913) ..(239.2280,240.5121) .. controls (239.2280,241.2329) and (239.8123,241.8172) ..(240.5331,241.8172) .. controls (241.2539,241.8172) and (241.8383,241.2329) ..(241.8383,240.5121);
\path[fill=chfourfig5c0,fill opacity=0.5,draw=chfourfig5c0,line width=0.3985pt,draw opacity=0.5] (253.5545,209.0696) .. controls (253.5545,208.3488) and (252.9702,207.7645) ..(252.2494,207.7645) .. controls (251.5286,207.7645) and (250.9442,208.3488) ..(250.9442,209.0696) .. controls (250.9442,209.7904) and (251.5286,210.3747) ..(252.2494,210.3747) .. controls (252.9702,210.3747) and (253.5545,209.7904) ..(253.5545,209.0696);
\path[fill=chfourfig5c0,fill opacity=0.5,draw=chfourfig5c0,line width=0.3985pt,draw opacity=0.5] (310.8404,217.8966) .. controls (310.8404,217.1758) and (310.2561,216.5915) ..(309.5353,216.5915) .. controls (308.8145,216.5915) and (308.2301,217.1758) ..(308.2301,217.8966) .. controls (308.2301,218.6174) and (308.8145,219.2017) ..(309.5353,219.2017) .. controls (310.2561,219.2017) and (310.8404,218.6174) ..(310.8404,217.8966);
\path[fill=chfourfig5c0,fill opacity=0.5,draw=chfourfig5c0,line width=0.3985pt,draw opacity=0.5] (304.1554,239.4162) .. controls (304.1554,238.6954) and (303.5710,238.1111) ..(302.8503,238.1111) .. controls (302.1294,238.1111) and (301.5451,238.6954) ..(301.5451,239.4162) .. controls (301.5451,240.1370) and (302.1294,240.7213) ..(302.8503,240.7213) .. controls (303.5710,240.7213) and (304.1554,240.1370) ..(304.1554,239.4162);
\path[fill=chfourfig5c0,fill opacity=0.5,draw=chfourfig5c0,line width=0.3985pt,draw opacity=0.5] (260.1498,236.5668) .. controls (260.1498,235.8461) and (259.5655,235.2617) ..(258.8447,235.2617) .. controls (258.1239,235.2617) and (257.5396,235.8461) ..(257.5396,236.5668) .. controls (257.5396,237.2877) and (258.1239,237.8720) ..(258.8447,237.8720) .. controls (259.5655,237.8720) and (260.1498,237.2877) ..(260.1498,236.5668);
\path[fill=chfourfig5c0,fill opacity=0.5,draw=chfourfig5c0,line width=0.3985pt,draw opacity=0.5] (242.2766,240.1236) .. controls (242.2766,239.4028) and (241.6923,238.8185) ..(240.9715,238.8185) .. controls (240.2507,238.8185) and (239.6664,239.4028) ..(239.6664,240.1236) .. controls (239.6664,240.8444) and (240.2507,241.4287) ..(240.9715,241.4287) .. controls (241.6923,241.4287) and (242.2766,240.8444) ..(242.2766,240.1236);
\path[fill=chfourfig5c0,fill opacity=0.5,draw=chfourfig5c0,line width=0.3985pt,draw opacity=0.5] (287.4777,260.6768) .. controls (287.4777,259.9559) and (286.8934,259.3716) ..(286.1726,259.3716) .. controls (285.4518,259.3716) and (284.8675,259.9559) ..(284.8675,260.6768) .. controls (284.8675,261.3976) and (285.4518,261.9819) ..(286.1726,261.9819) .. controls (286.8934,261.9819) and (287.4777,261.3976) ..(287.4777,260.6768);
\path[fill=chfourfig5c0,fill opacity=0.5,draw=chfourfig5c0,line width=0.3985pt,draw opacity=0.5] (238.5207,220.0884) .. controls (238.5207,219.3676) and (237.9363,218.7833) ..(237.2155,218.7833) .. controls (236.4948,218.7833) and (235.9104,219.3676) ..(235.9104,220.0884) .. controls (235.9104,220.8092) and (236.4948,221.3936) ..(237.2155,221.3936) .. controls (237.9363,221.3936) and (238.5207,220.8092) ..(238.5207,220.0884);
\path[fill=chfourfig5c0,fill opacity=0.5,draw=chfourfig5c0,line width=0.3985pt,draw opacity=0.5] (237.9727,238.7786) .. controls (237.9727,238.0578) and (237.3884,237.4734) ..(236.6676,237.4734) .. controls (235.9468,237.4734) and (235.3625,238.0578) ..(235.3625,238.7786) .. controls (235.3625,239.4994) and (235.9468,240.0837) ..(236.6676,240.0837) .. controls (237.3884,240.0837) and (237.9727,239.4994) ..(237.9727,238.7786);
\path[fill=chfourfig5c0,fill opacity=0.5,draw=chfourfig5c0,line width=0.3985pt,draw opacity=0.5] (302.4119,196.7855) .. controls (302.4119,196.0647) and (301.8275,195.4804) ..(301.1068,195.4804) .. controls (300.3860,195.4804) and (299.8016,196.0647) ..(299.8016,196.7855) .. controls (299.8016,197.5063) and (300.3860,198.0906) ..(301.1068,198.0906) .. controls (301.8275,198.0906) and (302.4119,197.5063) ..(302.4119,196.7855);
\path[fill=chfourfig5c0,fill opacity=0.5,draw=chfourfig5c0,line width=0.3985pt,draw opacity=0.5] (316.7583,194.6336) .. controls (316.7583,193.9128) and (316.1740,193.3284) ..(315.4532,193.3284) .. controls (314.7324,193.3284) and (314.1480,193.9128) ..(314.1480,194.6336) .. controls (314.1480,195.3544) and (314.7324,195.9387) ..(315.4532,195.9387) .. controls (316.1740,195.9387) and (316.7583,195.3544) ..(316.7583,194.6336);
\path[fill=chfourfig5c0,fill opacity=0.5,draw=chfourfig5c0,line width=0.3985pt,draw opacity=0.5] (274.5461,248.3229) .. controls (274.5461,247.6021) and (273.9617,247.0178) ..(273.2409,247.0178) .. controls (272.5201,247.0178) and (271.9358,247.6021) ..(271.9358,248.3229) .. controls (271.9358,249.0437) and (272.5201,249.6281) ..(273.2409,249.6281) .. controls (273.9617,249.6281) and (274.5461,249.0437) ..(274.5461,248.3229);
\path[fill=chfourfig5c0,fill opacity=0.5,draw=chfourfig5c0,line width=0.3985pt,draw opacity=0.5] (298.7356,262.3405) .. controls (298.7356,261.6198) and (298.1513,261.0354) ..(297.4305,261.0354) .. controls (296.7097,261.0354) and (296.1254,261.6198) ..(296.1254,262.3405) .. controls (296.1254,263.0613) and (296.7097,263.6456) ..(297.4305,263.6456) .. controls (298.1513,263.6456) and (298.7356,263.0613) ..(298.7356,262.3405);
\path[fill=chfourfig5c0,fill opacity=0.5,draw=chfourfig5c0,line width=0.3985pt,draw opacity=0.5] (296.5239,211.9688) .. controls (296.5239,211.2480) and (295.9396,210.6637) ..(295.2188,210.6637) .. controls (294.4980,210.6637) and (293.9136,211.2480) ..(293.9136,211.9688) .. controls (293.9136,212.6896) and (294.4980,213.2739) ..(295.2188,213.2739) .. controls (295.9396,213.2739) and (296.5239,212.6896) ..(296.5239,211.9688);
\path[fill=chfourfig5c0,fill opacity=0.5,draw=chfourfig5c0,line width=0.3985pt,draw opacity=0.5] (273.2309,204.3473) .. controls (273.2309,203.6265) and (272.6466,203.0421) ..(271.9258,203.0421) .. controls (271.2050,203.0421) and (270.6207,203.6265) ..(270.6207,204.3473) .. controls (270.6207,205.0681) and (271.2050,205.6524) ..(271.9258,205.6524) .. controls (272.6466,205.6524) and (273.2309,205.0681) ..(273.2309,204.3473);
\path[fill=chfourfig5c0,fill opacity=0.5,draw=chfourfig5c0,line width=0.3985pt,draw opacity=0.5] (275.7814,203.2513) .. controls (275.7814,202.5305) and (275.1971,201.9462) ..(274.4763,201.9462) .. controls (273.7555,201.9462) and (273.1712,202.5305) ..(273.1712,203.2513) .. controls (273.1712,203.9722) and (273.7555,204.5565) ..(274.4763,204.5565) .. controls (275.1971,204.5565) and (275.7814,203.9722) ..(275.7814,203.2513);
\path[fill=chfourfig5c0,fill opacity=0.5,draw=chfourfig5c0,line width=0.3985pt,draw opacity=0.5] (271.7565,199.0570) .. controls (271.7565,198.3362) and (271.1721,197.7519) ..(270.4513,197.7519) .. controls (269.7305,197.7519) and (269.1462,198.3362) ..(269.1462,199.0570) .. controls (269.1462,199.7778) and (269.7305,200.3622) ..(270.4513,200.3622) .. controls (271.1721,200.3622) and (271.7565,199.7778) ..(271.7565,199.0570);
\path[fill=chfourfig5c0,fill opacity=0.5,draw=chfourfig5c0,line width=0.3985pt,draw opacity=0.5] (293.6048,189.4928) .. controls (293.6048,188.7720) and (293.0205,188.1877) ..(292.2997,188.1877) .. controls (291.5789,188.1877) and (290.9946,188.7720) ..(290.9946,189.4928) .. controls (290.9946,190.2136) and (291.5789,190.7979) ..(292.2997,190.7979) .. controls (293.0205,190.7979) and (293.6048,190.2136) ..(293.6048,189.4928);
\path[fill=chfourfig5c0,fill opacity=0.5,draw=chfourfig5c0,line width=0.3985pt,draw opacity=0.5] (265.4102,227.6203) .. controls (265.4102,226.8995) and (264.8258,226.3152) ..(264.1050,226.3152) .. controls (263.3842,226.3152) and (262.7999,226.8995) ..(262.7999,227.6203) .. controls (262.7999,228.3411) and (263.3842,228.9254) ..(264.1050,228.9254) .. controls (264.8258,228.9254) and (265.4102,228.3411) ..(265.4102,227.6203);
\path[fill=chfourfig5c0,fill opacity=0.5,draw=chfourfig5c0,line width=0.3985pt,draw opacity=0.5] (314.5565,232.1135) .. controls (314.5565,231.3927) and (313.9722,230.8084) ..(313.2514,230.8084) .. controls (312.5306,230.8084) and (311.9463,231.3927) ..(311.9463,232.1135) .. controls (311.9463,232.8343) and (312.5306,233.4186) ..(313.2514,233.4186) .. controls (313.9722,233.4186) and (314.5565,232.8343) ..(314.5565,232.1135);
\path[fill=chfourfig5c0,fill opacity=0.5,draw=chfourfig5c0,line width=0.3985pt,draw opacity=0.5] (306.5564,214.0510) .. controls (306.5564,213.3302) and (305.9721,212.7458) ..(305.2513,212.7458) .. controls (304.5305,212.7458) and (303.9461,213.3302) ..(303.9461,214.0510) .. controls (303.9461,214.7718) and (304.5305,215.3561) ..(305.2513,215.3561) .. controls (305.9721,215.3561) and (306.5564,214.7718) ..(306.5564,214.0510);
\path[fill=chfourfig5c0,fill opacity=0.5,draw=chfourfig5c0,line width=0.3985pt,draw opacity=0.5] (298.5363,209.7072) .. controls (298.5363,208.9864) and (297.9520,208.4021) ..(297.2312,208.4021) .. controls (296.5104,208.4021) and (295.9261,208.9864) ..(295.9261,209.7072) .. controls (295.9261,210.4280) and (296.5104,211.0123) ..(297.2312,211.0123) .. controls (297.9520,211.0123) and (298.5363,210.4280) ..(298.5363,209.7072);
\path[fill=chfourfig5c0,fill opacity=0.5,draw=chfourfig5c0,line width=0.3985pt,draw opacity=0.5] (258.9443,244.3079) .. controls (258.9443,243.5871) and (258.3600,243.0028) ..(257.6392,243.0028) .. controls (256.9184,243.0028) and (256.3341,243.5871) ..(256.3341,244.3079) .. controls (256.3341,245.0287) and (256.9184,245.6131) ..(257.6392,245.6131) .. controls (258.3600,245.6131) and (258.9443,245.0287) ..(258.9443,244.3079);
\path[fill=chfourfig5c0,fill opacity=0.5,draw=chfourfig5c0,line width=0.3985pt,draw opacity=0.5] (297.9884,223.7249) .. controls (297.9884,223.0040) and (297.4041,222.4197) ..(296.6833,222.4197) .. controls (295.9625,222.4197) and (295.3782,223.0040) ..(295.3782,223.7249) .. controls (295.3782,224.4456) and (295.9625,225.0300) ..(296.6833,225.0300) .. controls (297.4041,225.0300) and (297.9884,224.4456) ..(297.9884,223.7249);
\path[fill=chfourfig5c0,fill opacity=0.5,draw=chfourfig5c0,line width=0.3985pt,draw opacity=0.5] (258.5657,210.1954) .. controls (258.5657,209.4746) and (257.9814,208.8903) ..(257.2606,208.8903) .. controls (256.5398,208.8903) and (255.9555,209.4746) ..(255.9555,210.1954) .. controls (255.9555,210.9162) and (256.5398,211.5005) ..(257.2606,211.5005) .. controls (257.9814,211.5005) and (258.5657,210.9162) ..(258.5657,210.1954);
\path[fill=chfourfig5c0,fill opacity=0.5,draw=chfourfig5c0,line width=0.3985pt,draw opacity=0.5] (272.2745,180.2374) .. controls (272.2745,179.5166) and (271.6902,178.9323) ..(270.9694,178.9323) .. controls (270.2486,178.9323) and (269.6642,179.5166) ..(269.6642,180.2374) .. controls (269.6642,180.9582) and (270.2486,181.5425) ..(270.9694,181.5425) .. controls (271.6902,181.5425) and (272.2745,180.9582) ..(272.2745,180.2374);
\path[fill=chfourfig5c0,fill opacity=0.5,draw=chfourfig5c0,line width=0.3985pt,draw opacity=0.5] (299.2238,262.0118) .. controls (299.2238,261.2910) and (298.6395,260.7066) ..(297.9187,260.7066) .. controls (297.1979,260.7066) and (296.6136,261.2910) ..(296.6136,262.0118) .. controls (296.6136,262.7325) and (297.1979,263.3169) ..(297.9187,263.3169) .. controls (298.6395,263.3169) and (299.2238,262.7325) ..(299.2238,262.0118);
\path[fill=chfourfig5c0,fill opacity=0.5,draw=chfourfig5c0,line width=0.3985pt,draw opacity=0.5] (312.4046,208.3025) .. controls (312.4046,207.5817) and (311.8203,206.9974) ..(311.0994,206.9974) .. controls (310.3786,206.9974) and (309.7943,207.5817) ..(309.7943,208.3025) .. controls (309.7943,209.0233) and (310.3786,209.6076) ..(311.0994,209.6076) .. controls (311.8203,209.6076) and (312.4046,209.0233) ..(312.4046,208.3025);
\path[fill=chfourfig5c0,fill opacity=0.5,draw=chfourfig5c0,line width=0.3985pt,draw opacity=0.5] (292.1104,248.9506) .. controls (292.1104,248.2298) and (291.5261,247.6455) ..(290.8053,247.6455) .. controls (290.0844,247.6455) and (289.5001,248.2298) ..(289.5001,248.9506) .. controls (289.5001,249.6714) and (290.0844,250.2557) ..(290.8053,250.2557) .. controls (291.5261,250.2557) and (292.1104,249.6714) ..(292.1104,248.9506);
\path[fill=chfourfig5c0,fill opacity=0.5,draw=chfourfig5c0,line width=0.3985pt,draw opacity=0.5] (307.9412,230.8980) .. controls (307.9412,230.1772) and (307.3569,229.5929) ..(306.6361,229.5929) .. controls (305.9153,229.5929) and (305.3310,230.1772) ..(305.3310,230.8980) .. controls (305.3310,231.6188) and (305.9153,232.2032) ..(306.6361,232.2032) .. controls (307.3569,232.2032) and (307.9412,231.6188) ..(307.9412,230.8980);
\path[fill=chfourfig5c0,fill opacity=0.5,draw=chfourfig5c0,line width=0.3985pt,draw opacity=0.5] (283.8513,257.3691) .. controls (283.8513,256.6483) and (283.2669,256.0640) ..(282.5461,256.0640) .. controls (281.8253,256.0640) and (281.2410,256.6483) ..(281.2410,257.3691) .. controls (281.2410,258.0899) and (281.8253,258.6743) ..(282.5461,258.6743) .. controls (283.2669,258.6743) and (283.8513,258.0899) ..(283.8513,257.3691);
\path[fill=chfourfig5c0,fill opacity=0.5,draw=chfourfig5c0,line width=0.3985pt,draw opacity=0.5] (239.4472,211.0920) .. controls (239.4472,210.3713) and (238.8629,209.7869) ..(238.1421,209.7869) .. controls (237.4213,209.7869) and (236.8369,210.3713) ..(236.8369,211.0920) .. controls (236.8369,211.8129) and (237.4213,212.3972) ..(238.1421,212.3972) .. controls (238.8629,212.3972) and (239.4472,211.8129) ..(239.4472,211.0920);
\path[fill=chfourfig5c0,fill opacity=0.5,draw=chfourfig5c0,line width=0.3985pt,draw opacity=0.5] (272.9819,209.1693) .. controls (272.9819,208.4485) and (272.3976,207.8641) ..(271.6768,207.8641) .. controls (270.9559,207.8641) and (270.3716,208.4485) ..(270.3716,209.1693) .. controls (270.3716,209.8900) and (270.9559,210.4744) ..(271.6768,210.4744) .. controls (272.3976,210.4744) and (272.9819,209.8900) ..(272.9819,209.1693);
\path[fill=chfourfig5c0,fill opacity=0.5,draw=chfourfig5c0,line width=0.3985pt,draw opacity=0.5] (270.5211,227.8295) .. controls (270.5211,227.1087) and (269.9368,226.5244) ..(269.2159,226.5244) .. controls (268.4951,226.5244) and (267.9108,227.1087) ..(267.9108,227.8295) .. controls (267.9108,228.5503) and (268.4951,229.1346) ..(269.2159,229.1346) .. controls (269.9368,229.1346) and (270.5211,228.5503) ..(270.5211,227.8295);
\path[fill=chfourfig5c0,fill opacity=0.5,draw=chfourfig5c0,line width=0.3985pt,draw opacity=0.5] (285.3158,234.2356) .. controls (285.3158,233.5148) and (284.7314,232.9305) ..(284.0107,232.9305) .. controls (283.2899,232.9305) and (282.7055,233.5148) ..(282.7055,234.2356) .. controls (282.7055,234.9564) and (283.2899,235.5407) ..(284.0107,235.5407) .. controls (284.7314,235.5407) and (285.3158,234.9564) ..(285.3158,234.2356);
\path[fill=chfourfig5c0,fill opacity=0.5,draw=chfourfig5c0,line width=0.3985pt,draw opacity=0.5] (303.1491,246.5994) .. controls (303.1491,245.8786) and (302.5648,245.2943) ..(301.8440,245.2943) .. controls (301.1232,245.2943) and (300.5389,245.8786) ..(300.5389,246.5994) .. controls (300.5389,247.3202) and (301.1232,247.9045) ..(301.8440,247.9045) .. controls (302.5648,247.9045) and (303.1491,247.3202) ..(303.1491,246.5994);
\path[fill=chfourfig5c0,fill opacity=0.5,draw=chfourfig5c0,line width=0.3985pt,draw opacity=0.5] (293.3757,193.4779) .. controls (293.3757,192.7571) and (292.7914,192.1728) ..(292.0705,192.1728) .. controls (291.3497,192.1728) and (290.7654,192.7571) ..(290.7654,193.4779) .. controls (290.7654,194.1987) and (291.3497,194.7830) ..(292.0705,194.7830) .. controls (292.7914,194.7830) and (293.3757,194.1987) ..(293.3757,193.4779);
\path[fill=chfourfig5c0,fill opacity=0.5,draw=chfourfig5c0,line width=0.3985pt,draw opacity=0.5] (268.7875,208.8504) .. controls (268.7875,208.1296) and (268.2032,207.5453) ..(267.4824,207.5453) .. controls (266.7616,207.5453) and (266.1773,208.1296) ..(266.1773,208.8504) .. controls (266.1773,209.5712) and (266.7616,210.1556) ..(267.4824,210.1556) .. controls (268.2032,210.1556) and (268.7875,209.5712) ..(268.7875,208.8504);
\path[fill=chfourfig5c0,fill opacity=0.5,draw=chfourfig5c0,line width=0.3985pt,draw opacity=0.5] (278.2920,225.2591) .. controls (278.2920,224.5383) and (277.7077,223.9540) ..(276.9869,223.9540) .. controls (276.2661,223.9540) and (275.6818,224.5383) ..(275.6818,225.2591) .. controls (275.6818,225.9799) and (276.2661,226.5642) ..(276.9869,226.5642) .. controls (277.7077,226.5642) and (278.2920,225.9799) ..(278.2920,225.2591);
\path[fill=chfourfig5c0,fill opacity=0.5,draw=chfourfig5c0,line width=0.3985pt,draw opacity=0.5] (264.1250,227.5804) .. controls (264.1250,226.8596) and (263.5406,226.2753) ..(262.8199,226.2753) .. controls (262.0991,226.2753) and (261.5147,226.8596) ..(261.5147,227.5804) .. controls (261.5147,228.3012) and (262.0991,228.8856) ..(262.8199,228.8856) .. controls (263.5406,228.8856) and (264.1250,228.3012) ..(264.1250,227.5804);
\path[fill=chfourfig5c0,fill opacity=0.5,draw=chfourfig5c0,line width=0.3985pt,draw opacity=0.5] (310.6511,196.1479) .. controls (310.6511,195.4271) and (310.0668,194.8428) ..(309.3460,194.8428) .. controls (308.6252,194.8428) and (308.0409,195.4271) ..(308.0409,196.1479) .. controls (308.0409,196.8687) and (308.6252,197.4530) ..(309.3460,197.4530) .. controls (310.0668,197.4530) and (310.6511,196.8687) ..(310.6511,196.1479);
\path[fill=chfourfig5c0,fill opacity=0.5,draw=chfourfig5c0,line width=0.3985pt,draw opacity=0.5] (222.1319,230.3302) .. controls (222.1319,229.6093) and (221.5476,229.0250) ..(220.8268,229.0250) .. controls (220.1060,229.0250) and (219.5217,229.6093) ..(219.5217,230.3302) .. controls (219.5217,231.0510) and (220.1060,231.6353) ..(220.8268,231.6353) .. controls (221.5476,231.6353) and (222.1319,231.0510) ..(222.1319,230.3302);
\path[fill=chfourfig5c0,fill opacity=0.5,draw=chfourfig5c0,line width=0.3985pt,draw opacity=0.5] (260.6679,202.2451) .. controls (260.6679,201.5243) and (260.0836,200.9400) ..(259.3628,200.9400) .. controls (258.6420,200.9400) and (258.0576,201.5243) ..(258.0576,202.2451) .. controls (258.0576,202.9659) and (258.6420,203.5502) ..(259.3628,203.5502) .. controls (260.0836,203.5502) and (260.6679,202.9659) ..(260.6679,202.2451);
\path[fill=chfourfig5c0,fill opacity=0.5,draw=chfourfig5c0,line width=0.3985pt,draw opacity=0.5] (309.5851,241.1597) .. controls (309.5851,240.4389) and (309.0008,239.8546) ..(308.2800,239.8546) .. controls (307.5591,239.8546) and (306.9749,240.4389) ..(306.9749,241.1597) .. controls (306.9749,241.8805) and (307.5591,242.4648) ..(308.2800,242.4648) .. controls (309.0008,242.4648) and (309.5851,241.8805) ..(309.5851,241.1597);
\path[fill=chfourfig5c0,fill opacity=0.5,draw=chfourfig5c0,line width=0.3985pt,draw opacity=0.5] (290.9946,243.0028) .. controls (290.9946,242.2820) and (290.4102,241.6977) ..(289.6894,241.6977) .. controls (288.9687,241.6977) and (288.3843,242.2820) ..(288.3843,243.0028) .. controls (288.3843,243.7236) and (288.9687,244.3079) ..(289.6894,244.3079) .. controls (290.4102,244.3079) and (290.9946,243.7236) ..(290.9946,243.0028);
\path[fill=chfourfig5c0,fill opacity=0.5,draw=chfourfig5c0,line width=0.3985pt,draw opacity=0.5] (319.6176,203.4606) .. controls (319.6176,202.7398) and (319.0333,202.1555) ..(318.3125,202.1555) .. controls (317.5916,202.1555) and (317.0074,202.7398) ..(317.0074,203.4606) .. controls (317.0074,204.1814) and (317.5916,204.7657) ..(318.3125,204.7657) .. controls (319.0333,204.7657) and (319.6176,204.1814) ..(319.6176,203.4606);
\path[fill=chfourfig5c0,fill opacity=0.5,draw=chfourfig5c0,line width=0.3985pt,draw opacity=0.5] (291.6920,189.5526) .. controls (291.6920,188.8317) and (291.1076,188.2474) ..(290.3868,188.2474) .. controls (289.6660,188.2474) and (289.0817,188.8317) ..(289.0817,189.5526) .. controls (289.0817,190.2733) and (289.6660,190.8577) ..(290.3868,190.8577) .. controls (291.1076,190.8577) and (291.6920,190.2733) ..(291.6920,189.5526);
\path[fill=chfourfig5c0,fill opacity=0.5,draw=chfourfig5c0,line width=0.3985pt,draw opacity=0.5] (307.4730,251.7501) .. controls (307.4730,251.0293) and (306.8887,250.4450) ..(306.1678,250.4450) .. controls (305.4471,250.4450) and (304.8627,251.0293) ..(304.8627,251.7501) .. controls (304.8627,252.4709) and (305.4471,253.0552) ..(306.1678,253.0552) .. controls (306.8887,253.0552) and (307.4730,252.4709) ..(307.4730,251.7501);
\path[fill=chfourfig5c0,fill opacity=0.5,draw=chfourfig5c0,line width=0.3985pt,draw opacity=0.5] (236.8967,212.1581) .. controls (236.8967,211.4373) and (236.3124,210.8530) ..(235.5916,210.8530) .. controls (234.8708,210.8530) and (234.2865,211.4373) ..(234.2865,212.1581) .. controls (234.2865,212.8789) and (234.8708,213.4632) ..(235.5916,213.4632) .. controls (236.3124,213.4632) and (236.8967,212.8789) ..(236.8967,212.1581);
\path[fill=chfourfig5c0,fill opacity=0.5,draw=chfourfig5c0,line width=0.3985pt,draw opacity=0.5] (275.0342,187.3807) .. controls (275.0342,186.6599) and (274.4499,186.0755) ..(273.7291,186.0755) .. controls (273.0083,186.0755) and (272.4240,186.6599) ..(272.4240,187.3807) .. controls (272.4240,188.1015) and (273.0083,188.6858) ..(273.7291,188.6858) .. controls (274.4499,188.6858) and (275.0342,188.1015) ..(275.0342,187.3807);
\path[fill=chfourfig5c0,fill opacity=0.5,draw=chfourfig5c0,line width=0.3985pt,draw opacity=0.5] (296.5338,189.9511) .. controls (296.5338,189.2303) and (295.9495,188.6460) ..(295.2287,188.6460) .. controls (294.5079,188.6460) and (293.9236,189.2303) ..(293.9236,189.9511) .. controls (293.9236,190.6719) and (294.5079,191.2562) ..(295.2287,191.2562) .. controls (295.9495,191.2562) and (296.5338,190.6719) ..(296.5338,189.9511);
\path[fill=chfourfig5c0,fill opacity=0.5,draw=chfourfig5c0,line width=0.3985pt,draw opacity=0.5] (328.0660,224.9403) .. controls (328.0660,224.2195) and (327.4817,223.6352) ..(326.7609,223.6352) .. controls (326.0401,223.6352) and (325.4558,224.2195) ..(325.4558,224.9403) .. controls (325.4558,225.6611) and (326.0401,226.2454) ..(326.7609,226.2454) .. controls (327.4817,226.2454) and (328.0660,225.6611) ..(328.0660,224.9403);
\path[fill=chfourfig5c0,fill opacity=0.5,draw=chfourfig5c0,line width=0.3985pt,draw opacity=0.5] (238.8793,218.9527) .. controls (238.8793,218.2319) and (238.2950,217.6476) ..(237.5742,217.6476) .. controls (236.8534,217.6476) and (236.2691,218.2319) ..(236.2691,218.9527) .. controls (236.2691,219.6735) and (236.8534,220.2578) ..(237.5742,220.2578) .. controls (238.2950,220.2578) and (238.8793,219.6735) ..(238.8793,218.9527);
\path[fill=chfourfig5c0,fill opacity=0.5,draw=chfourfig5c0,line width=0.3985pt,draw opacity=0.5] (258.4960,189.8514) .. controls (258.4960,189.1306) and (257.9117,188.5463) ..(257.1909,188.5463) .. controls (256.4701,188.5463) and (255.8858,189.1306) ..(255.8858,189.8514) .. controls (255.8858,190.5722) and (256.4701,191.1566) ..(257.1909,191.1566) .. controls (257.9117,191.1566) and (258.4960,190.5722) ..(258.4960,189.8514);
\path[fill=chfourfig5c0,fill opacity=0.5,draw=chfourfig5c0,line width=0.3985pt,draw opacity=0.5] (334.9603,216.0137) .. controls (334.9603,215.2928) and (334.3759,214.7086) ..(333.6552,214.7086) .. controls (332.9343,214.7086) and (332.3500,215.2928) ..(332.3500,216.0137) .. controls (332.3500,216.7345) and (332.9343,217.3188) ..(333.6552,217.3188) .. controls (334.3759,217.3188) and (334.9603,216.7345) ..(334.9603,216.0137);
\path[fill=chfourfig5c0,fill opacity=0.5,draw=chfourfig5c0,line width=0.3985pt,draw opacity=0.5] (275.4825,244.9854) .. controls (275.4825,244.2646) and (274.8982,243.6803) ..(274.1774,243.6803) .. controls (273.4566,243.6803) and (272.8723,244.2646) ..(272.8723,244.9854) .. controls (272.8723,245.7062) and (273.4566,246.2905) ..(274.1774,246.2905) .. controls (274.8982,246.2905) and (275.4825,245.7062) ..(275.4825,244.9854);
\path[fill=chfourfig5c0,fill opacity=0.5,draw=chfourfig5c0,line width=0.3985pt,draw opacity=0.5] (316.4694,199.4356) .. controls (316.4694,198.7148) and (315.8850,198.1305) ..(315.1642,198.1305) .. controls (314.4434,198.1305) and (313.8591,198.7148) ..(313.8591,199.4356) .. controls (313.8591,200.1564) and (314.4434,200.7407) ..(315.1642,200.7407) .. controls (315.8850,200.7407) and (316.4694,200.1564) ..(316.4694,199.4356);
\path[fill=chfourfig5c0,fill opacity=0.5,draw=chfourfig5c0,line width=0.3985pt,draw opacity=0.5] (297.5302,232.4921) .. controls (297.5302,231.7713) and (296.9458,231.1870) ..(296.2250,231.1870) .. controls (295.5042,231.1870) and (294.9199,231.7713) ..(294.9199,232.4921) .. controls (294.9199,233.2129) and (295.5042,233.7972) ..(296.2250,233.7972) .. controls (296.9458,233.7972) and (297.5302,233.2129) ..(297.5302,232.4921);
\path[fill=chfourfig5c0,fill opacity=0.5,draw=chfourfig5c0,line width=0.3985pt,draw opacity=0.5] (272.8524,176.9795) .. controls (272.8524,176.2587) and (272.2680,175.6744) ..(271.5472,175.6744) .. controls (270.8264,175.6744) and (270.2421,176.2587) ..(270.2421,176.9795) .. controls (270.2421,177.7003) and (270.8264,178.2847) ..(271.5472,178.2847) .. controls (272.2680,178.2847) and (272.8524,177.7003) ..(272.8524,176.9795);
\path[fill=chfourfig5c0,fill opacity=0.5,draw=chfourfig5c0,line width=0.3985pt,draw opacity=0.5] (284.5387,198.7083) .. controls (284.5387,197.9875) and (283.9544,197.4032) ..(283.2336,197.4032) .. controls (282.5128,197.4032) and (281.9284,197.9875) ..(281.9284,198.7083) .. controls (281.9284,199.4291) and (282.5128,200.0135) ..(283.2336,200.0135) .. controls (283.9544,200.0135) and (284.5387,199.4291) ..(284.5387,198.7083);
\path[fill=chfourfig5c0,fill opacity=0.5,draw=chfourfig5c0,line width=0.3985pt,draw opacity=0.5] (267.1138,212.7658) .. controls (267.1138,212.0450) and (266.5295,211.4607) ..(265.8087,211.4607) .. controls (265.0879,211.4607) and (264.5036,212.0450) ..(264.5036,212.7658) .. controls (264.5036,213.4866) and (265.0879,214.0709) ..(265.8087,214.0709) .. controls (266.5295,214.0709) and (267.1138,213.4866) ..(267.1138,212.7658);
\path[fill=chfourfig5c0,fill opacity=0.5,draw=chfourfig5c0,line width=0.3985pt,draw opacity=0.5] (267.8610,227.4011) .. controls (267.8610,226.6803) and (267.2767,226.0960) ..(266.5559,226.0960) .. controls (265.8351,226.0960) and (265.2508,226.6803) ..(265.2508,227.4011) .. controls (265.2508,228.1219) and (265.8351,228.7062) ..(266.5559,228.7062) .. controls (267.2767,228.7062) and (267.8610,228.1219) ..(267.8610,227.4011);
\path[fill=chfourfig5c0,fill opacity=0.5,draw=chfourfig5c0,line width=0.3985pt,draw opacity=0.5] (309.6847,230.0512) .. controls (309.6847,229.3304) and (309.1004,228.7461) ..(308.3796,228.7461) .. controls (307.6588,228.7461) and (307.0745,229.3304) ..(307.0745,230.0512) .. controls (307.0745,230.7720) and (307.6588,231.3563) ..(308.3796,231.3563) .. controls (309.1004,231.3563) and (309.6847,230.7720) ..(309.6847,230.0512);
\path[fill=chfourfig5c0,fill opacity=0.5,draw=chfourfig5c0,line width=0.3985pt,draw opacity=0.5] (242.6054,217.5081) .. controls (242.6054,216.7873) and (242.0211,216.2029) ..(241.3003,216.2029) .. controls (240.5795,216.2029) and (239.9951,216.7873) ..(239.9951,217.5081) .. controls (239.9951,218.2289) and (240.5795,218.8132) ..(241.3003,218.8132) .. controls (242.0211,218.8132) and (242.6054,218.2289) ..(242.6054,217.5081);
\path[fill=chfourfig5c0,fill opacity=0.5,draw=chfourfig5c0,line width=0.3985pt,draw opacity=0.5] (287.8364,228.9553) .. controls (287.8364,228.2345) and (287.2520,227.6502) ..(286.5312,227.6502) .. controls (285.8104,227.6502) and (285.2261,228.2345) ..(285.2261,228.9553) .. controls (285.2261,229.6761) and (285.8104,230.2604) ..(286.5312,230.2604) .. controls (287.2520,230.2604) and (287.8364,229.6761) ..(287.8364,228.9553);
\path[fill=chfourfig5c0,fill opacity=0.5,draw=chfourfig5c0,line width=0.3985pt,draw opacity=0.5] (271.2982,222.7385) .. controls (271.2982,222.0177) and (270.7138,221.4334) ..(269.9930,221.4334) .. controls (269.2722,221.4334) and (268.6879,222.0177) ..(268.6879,222.7385) .. controls (268.6879,223.4594) and (269.2722,224.0436) ..(269.9930,224.0436) .. controls (270.7138,224.0436) and (271.2982,223.4594) ..(271.2982,222.7385);
\path[fill=chfourfig5c0,fill opacity=0.5,draw=chfourfig5c0,line width=0.3985pt,draw opacity=0.5] (237.2355,220.3773) .. controls (237.2355,219.6566) and (236.6511,219.0722) ..(235.9303,219.0722) .. controls (235.2095,219.0722) and (234.6252,219.6566) ..(234.6252,220.3773) .. controls (234.6252,221.0981) and (235.2095,221.6825) ..(235.9303,221.6825) .. controls (236.6511,221.6825) and (237.2355,221.0981) ..(237.2355,220.3773);
\path[fill=chfourfig5c0,fill opacity=0.5,draw=chfourfig5c0,line width=0.3985pt,draw opacity=0.5] (249.5096,196.9350) .. controls (249.5096,196.2142) and (248.9253,195.6298) ..(248.2045,195.6298) .. controls (247.4837,195.6298) and (246.8994,196.2142) ..(246.8994,196.9350) .. controls (246.8994,197.6558) and (247.4837,198.2401) ..(248.2045,198.2401) .. controls (248.9253,198.2401) and (249.5096,197.6558) ..(249.5096,196.9350);
\path[fill=chfourfig5c0,fill opacity=0.5,draw=chfourfig5c0,line width=0.3985pt,draw opacity=0.5] (258.7849,244.4474) .. controls (258.7849,243.7266) and (258.2006,243.1423) ..(257.4798,243.1423) .. controls (256.7590,243.1423) and (256.1747,243.7266) ..(256.1747,244.4474) .. controls (256.1747,245.1682) and (256.7590,245.7525) ..(257.4798,245.7525) .. controls (258.2006,245.7525) and (258.7849,245.1682) ..(258.7849,244.4474);
\path[fill=chfourfig5c0,fill opacity=0.5,draw=chfourfig5c0,line width=0.3985pt,draw opacity=0.5] (270.7104,195.5800) .. controls (270.7104,194.8592) and (270.1260,194.2749) ..(269.4052,194.2749) .. controls (268.6844,194.2749) and (268.1001,194.8592) ..(268.1001,195.5800) .. controls (268.1001,196.3008) and (268.6844,196.8851) ..(269.4052,196.8851) .. controls (270.1260,196.8851) and (270.7104,196.3008) ..(270.7104,195.5800);
\path[fill=chfourfig5c0,fill opacity=0.5,draw=chfourfig5c0,line width=0.3985pt,draw opacity=0.5] (306.2276,250.3155) .. controls (306.2276,249.5947) and (305.6433,249.0104) ..(304.9225,249.0104) .. controls (304.2017,249.0104) and (303.6174,249.5947) ..(303.6174,250.3155) .. controls (303.6174,251.0363) and (304.2017,251.6206) ..(304.9225,251.6206) .. controls (305.6433,251.6206) and (306.2276,251.0363) ..(306.2276,250.3155);
\path[fill=chfourfig5c0,fill opacity=0.5,draw=chfourfig5c0,line width=0.3985pt,draw opacity=0.5] (249.2107,197.9312) .. controls (249.2107,197.2104) and (248.6264,196.6261) ..(247.9056,196.6261) .. controls (247.1848,196.6261) and (246.6005,197.2104) ..(246.6005,197.9312) .. controls (246.6005,198.6520) and (247.1848,199.2364) ..(247.9056,199.2364) .. controls (248.6264,199.2364) and (249.2107,198.6520) ..(249.2107,197.9312);
\path[fill=chfourfig5c0,fill opacity=0.5,draw=chfourfig5c0,line width=0.3985pt,draw opacity=0.5] (293.3159,205.6424) .. controls (293.3159,204.9216) and (292.7316,204.3373) ..(292.0107,204.3373) .. controls (291.2900,204.3373) and (290.7057,204.9216) ..(290.7057,205.6424) .. controls (290.7057,206.3632) and (291.2900,206.9475) ..(292.0107,206.9475) .. controls (292.7316,206.9475) and (293.3159,206.3632) ..(293.3159,205.6424);
\path[fill=chfourfig5c0,fill opacity=0.5,draw=chfourfig5c0,line width=0.3985pt,draw opacity=0.5] (270.5410,196.2177) .. controls (270.5410,195.4968) and (269.9567,194.9125) ..(269.2359,194.9125) .. controls (268.5151,194.9125) and (267.9308,195.4968) ..(267.9308,196.2177) .. controls (267.9308,196.9384) and (268.5151,197.5228) ..(269.2359,197.5228) .. controls (269.9567,197.5228) and (270.5410,196.9384) ..(270.5410,196.2177);
\path[fill=chfourfig5c0,fill opacity=0.5,draw=chfourfig5c0,line width=0.3985pt,draw opacity=0.5] (244.9965,228.9453) .. controls (244.9965,228.2245) and (244.4121,227.6402) ..(243.6913,227.6402) .. controls (242.9705,227.6402) and (242.3862,228.2245) ..(242.3862,228.9453) .. controls (242.3862,229.6661) and (242.9705,230.2505) ..(243.6913,230.2505) .. controls (244.4121,230.2505) and (244.9965,229.6661) ..(244.9965,228.9453);
\path[fill=chfourfig5c0,fill opacity=0.5,draw=chfourfig5c0,line width=0.3985pt,draw opacity=0.5] (236.2093,238.8981) .. controls (236.2093,238.1773) and (235.6250,237.5930) ..(234.9042,237.5930) .. controls (234.1833,237.5930) and (233.5991,238.1773) ..(233.5991,238.8981) .. controls (233.5991,239.6190) and (234.1833,240.2032) ..(234.9042,240.2032) .. controls (235.6250,240.2032) and (236.2093,239.6190) ..(236.2093,238.8981);
\path[fill=chfourfig5c0,fill opacity=0.5,draw=chfourfig5c0,line width=0.3985pt,draw opacity=0.5] (272.8424,204.9948) .. controls (272.8424,204.2740) and (272.2581,203.6897) ..(271.5373,203.6897) .. controls (270.8165,203.6897) and (270.2321,204.2740) ..(270.2321,204.9948) .. controls (270.2321,205.7156) and (270.8165,206.3000) ..(271.5373,206.3000) .. controls (272.2581,206.3000) and (272.8424,205.7156) ..(272.8424,204.9948);
\path[fill=chfourfig5c0,fill opacity=0.5,draw=chfourfig5c0,line width=0.3985pt,draw opacity=0.5] (299.2537,259.3318) .. controls (299.2537,258.6110) and (298.6694,258.0267) ..(297.9486,258.0267) .. controls (297.2278,258.0267) and (296.6434,258.6110) ..(296.6434,259.3318) .. controls (296.6434,260.0526) and (297.2278,260.6369) ..(297.9486,260.6369) .. controls (298.6694,260.6369) and (299.2537,260.0526) ..(299.2537,259.3318);
\path[fill=chfourfig5c0,fill opacity=0.5,draw=chfourfig5c0,line width=0.3985pt,draw opacity=0.5] (250.6354,234.8234) .. controls (250.6354,234.1026) and (250.0511,233.5182) ..(249.3303,233.5182) .. controls (248.6095,233.5182) and (248.0251,234.1026) ..(248.0251,234.8234) .. controls (248.0251,235.5442) and (248.6095,236.1285) ..(249.3303,236.1285) .. controls (250.0511,236.1285) and (250.6354,235.5442) ..(250.6354,234.8234);
\path[fill=chfourfig5c0,fill opacity=0.5,draw=chfourfig5c0,line width=0.3985pt,draw opacity=0.5] (276.1999,206.8679) .. controls (276.1999,206.1470) and (275.6155,205.5627) ..(274.8947,205.5627) .. controls (274.1739,205.5627) and (273.5896,206.1470) ..(273.5896,206.8679) .. controls (273.5896,207.5886) and (274.1739,208.1730) ..(274.8947,208.1730) .. controls (275.6155,208.1730) and (276.1999,207.5886) ..(276.1999,206.8679);
\path[fill=chfourfig5c0,fill opacity=0.5,draw=chfourfig5c0,line width=0.3985pt,draw opacity=0.5] (324.6288,234.5245) .. controls (324.6288,233.8037) and (324.0445,233.2194) ..(323.3238,233.2194) .. controls (322.6029,233.2194) and (322.0186,233.8037) ..(322.0186,234.5245) .. controls (322.0186,235.2453) and (322.6029,235.8296) ..(323.3238,235.8296) .. controls (324.0445,235.8296) and (324.6288,235.2453) ..(324.6288,234.5245);
\path[fill=chfourfig5c0,fill opacity=0.5,draw=chfourfig5c0,line width=0.3985pt,draw opacity=0.5] (304.0558,219.6800) .. controls (304.0558,218.9592) and (303.4714,218.3748) ..(302.7506,218.3748) .. controls (302.0298,218.3748) and (301.4455,218.9592) ..(301.4455,219.6800) .. controls (301.4455,220.4008) and (302.0298,220.9851) ..(302.7506,220.9851) .. controls (303.4714,220.9851) and (304.0558,220.4008) ..(304.0558,219.6800);
\path[fill=chfourfig5c0,fill opacity=0.5,draw=chfourfig5c0,line width=0.3985pt,draw opacity=0.5] (288.0057,184.1228) .. controls (288.0057,183.4020) and (287.4214,182.8177) ..(286.7006,182.8177) .. controls (285.9798,182.8177) and (285.3955,183.4020) ..(285.3955,184.1228) .. controls (285.3955,184.8436) and (285.9798,185.4280) ..(286.7006,185.4280) .. controls (287.4214,185.4280) and (288.0057,184.8436) ..(288.0057,184.1228);
\path[fill=chfourfig5c0,fill opacity=0.5,draw=chfourfig5c0,line width=0.3985pt,draw opacity=0.5] (325.2565,213.3536) .. controls (325.2565,212.6328) and (324.6722,212.0485) ..(323.9514,212.0485) .. controls (323.2306,212.0485) and (322.6463,212.6328) ..(322.6463,213.3536) .. controls (322.6463,214.0744) and (323.2306,214.6587) ..(323.9514,214.6587) .. controls (324.6722,214.6587) and (325.2565,214.0744) ..(325.2565,213.3536);
\path[fill=chfourfig5c0,fill opacity=0.5,draw=chfourfig5c0,line width=0.3985pt,draw opacity=0.5] (299.0843,263.3268) .. controls (299.0843,262.6060) and (298.5000,262.0217) ..(297.7792,262.0217) .. controls (297.0584,262.0217) and (296.4741,262.6060) ..(296.4741,263.3268) .. controls (296.4741,264.0477) and (297.0584,264.6320) ..(297.7792,264.6320) .. controls (298.5000,264.6320) and (299.0843,264.0477) ..(299.0843,263.3268);
\path[fill=chfourfig5c0,fill opacity=0.5,draw=chfourfig5c0,line width=0.3985pt,draw opacity=0.5] (281.1115,249.9767) .. controls (281.1115,249.2560) and (280.5272,248.6716) ..(279.8064,248.6716) .. controls (279.0856,248.6716) and (278.5013,249.2560) ..(278.5013,249.9767) .. controls (278.5013,250.6976) and (279.0856,251.2819) ..(279.8064,251.2819) .. controls (280.5272,251.2819) and (281.1115,250.6976) ..(281.1115,249.9767);
\path[fill=chfourfig5c0,fill opacity=0.5,draw=chfourfig5c0,line width=0.3985pt,draw opacity=0.5] (242.8047,222.9178) .. controls (242.8047,222.1971) and (242.2203,221.6127) ..(241.4995,221.6127) .. controls (240.7787,221.6127) and (240.1944,222.1971) ..(240.1944,222.9178) .. controls (240.1944,223.6387) and (240.7787,224.2230) ..(241.4995,224.2230) .. controls (242.2203,224.2230) and (242.8047,223.6387) ..(242.8047,222.9178);
\path[fill=chfourfig5c0,fill opacity=0.5,draw=chfourfig5c0,line width=0.3985pt,draw opacity=0.5] (285.6744,265.4290) .. controls (285.6744,264.7082) and (285.0901,264.1239) ..(284.3693,264.1239) .. controls (283.6485,264.1239) and (283.0642,264.7082) ..(283.0642,265.4290) .. controls (283.0642,266.1498) and (283.6485,266.7341) ..(284.3693,266.7341) .. controls (285.0901,266.7341) and (285.6744,266.1498) ..(285.6744,265.4290);
\path[fill=chfourfig5c0,fill opacity=0.5,draw=chfourfig5c0,line width=0.3985pt,draw opacity=0.5] (265.5496,190.9075) .. controls (265.5496,190.1867) and (264.9653,189.6024) ..(264.2445,189.6024) .. controls (263.5237,189.6024) and (262.9394,190.1867) ..(262.9394,190.9075) .. controls (262.9394,191.6283) and (263.5237,192.2126) ..(264.2445,192.2126) .. controls (264.9653,192.2126) and (265.5496,191.6283) ..(265.5496,190.9075);
\path[fill=chfourfig5c0,fill opacity=0.5,draw=chfourfig5c0,line width=0.3985pt,draw opacity=0.5] (264.2744,182.7480) .. controls (264.2744,182.0272) and (263.6901,181.4429) ..(262.9693,181.4429) .. controls (262.2485,181.4429) and (261.6642,182.0272) ..(261.6642,182.7480) .. controls (261.6642,183.4688) and (262.2485,184.0531) ..(262.9693,184.0531) .. controls (263.6901,184.0531) and (264.2744,183.4688) ..(264.2744,182.7480);
\path[fill=chfourfig5c0,fill opacity=0.5,draw=chfourfig5c0,line width=0.3985pt,draw opacity=0.5] (307.4730,210.1157) .. controls (307.4730,209.3949) and (306.8887,208.8106) ..(306.1678,208.8106) .. controls (305.4471,208.8106) and (304.8627,209.3949) ..(304.8627,210.1157) .. controls (304.8627,210.8365) and (305.4471,211.4208) ..(306.1678,211.4208) .. controls (306.8887,211.4208) and (307.4730,210.8365) ..(307.4730,210.1157);
\path[fill=chfourfig5c0,fill opacity=0.5,draw=chfourfig5c0,line width=0.3985pt,draw opacity=0.5] (246.3912,231.8545) .. controls (246.3912,231.1337) and (245.8069,230.5493) ..(245.0861,230.5493) .. controls (244.3653,230.5493) and (243.7810,231.1337) ..(243.7810,231.8545) .. controls (243.7810,232.5753) and (244.3653,233.1596) ..(245.0861,233.1596) .. controls (245.8069,233.1596) and (246.3912,232.5753) ..(246.3912,231.8545);
\path[fill=chfourfig5c0,fill opacity=0.5,draw=chfourfig5c0,line width=0.3985pt,draw opacity=0.5] (272.2247,238.7089) .. controls (272.2247,237.9880) and (271.6404,237.4037) ..(270.9196,237.4037) .. controls (270.1988,237.4037) and (269.6145,237.9880) ..(269.6145,238.7089) .. controls (269.6145,239.4296) and (270.1988,240.0140) ..(270.9196,240.0140) .. controls (271.6404,240.0140) and (272.2247,239.4296) ..(272.2247,238.7089);
\path[fill=chfourfig5c0,fill opacity=0.5,draw=chfourfig5c0,line width=0.3985pt,draw opacity=0.5] (259.7214,205.6723) .. controls (259.7214,204.9515) and (259.1371,204.3672) ..(258.4163,204.3672) .. controls (257.6955,204.3672) and (257.1112,204.9515) ..(257.1112,205.6723) .. controls (257.1112,206.3931) and (257.6955,206.9774) ..(258.4163,206.9774) .. controls (259.1371,206.9774) and (259.7214,206.3931) ..(259.7214,205.6723);
\path[fill=chfourfig5c0,fill opacity=0.5,draw=chfourfig5c0,line width=0.3985pt,draw opacity=0.5] (290.1876,220.5567) .. controls (290.1876,219.8359) and (289.6033,219.2516) ..(288.8824,219.2516) .. controls (288.1617,219.2516) and (287.5773,219.8359) ..(287.5773,220.5567) .. controls (287.5773,221.2775) and (288.1617,221.8618) ..(288.8824,221.8618) .. controls (289.6033,221.8618) and (290.1876,221.2775) ..(290.1876,220.5567);
\path[fill=chfourfig5c0,fill opacity=0.5,draw=chfourfig5c0,line width=0.3985pt,draw opacity=0.5] (274.3866,200.7308) .. controls (274.3866,200.0100) and (273.8023,199.4257) ..(273.0815,199.4257) .. controls (272.3607,199.4257) and (271.7764,200.0100) ..(271.7764,200.7308) .. controls (271.7764,201.4516) and (272.3607,202.0359) ..(273.0815,202.0359) .. controls (273.8023,202.0359) and (274.3866,201.4516) ..(274.3866,200.7308);
\path[fill=chfourfig5c0,fill opacity=0.5,draw=chfourfig5c0,line width=0.3985pt,draw opacity=0.5] (284.2099,191.5451) .. controls (284.2099,190.8243) and (283.6256,190.2400) ..(282.9048,190.2400) .. controls (282.1840,190.2400) and (281.5997,190.8243) ..(281.5997,191.5451) .. controls (281.5997,192.2659) and (282.1840,192.8502) ..(282.9048,192.8502) .. controls (283.6256,192.8502) and (284.2099,192.2659) ..(284.2099,191.5451);
\path[fill=chfourfig5c0,fill opacity=0.5,draw=chfourfig5c0,line width=0.3985pt,draw opacity=0.5] (291.1042,215.8542) .. controls (291.1042,215.1335) and (290.5198,214.5491) ..(289.7990,214.5491) .. controls (289.0782,214.5491) and (288.4939,215.1335) ..(288.4939,215.8542) .. controls (288.4939,216.5751) and (289.0782,217.1594) ..(289.7990,217.1594) .. controls (290.5198,217.1594) and (291.1042,216.5751) ..(291.1042,215.8542);
\path[fill=chfourfig5c0,fill opacity=0.5,draw=chfourfig5c0,line width=0.3985pt,draw opacity=0.5] (263.5172,228.8656) .. controls (263.5172,228.1448) and (262.9329,227.5605) ..(262.2121,227.5605) .. controls (261.4913,227.5605) and (260.9070,228.1448) ..(260.9070,228.8656) .. controls (260.9070,229.5865) and (261.4913,230.1708) ..(262.2121,230.1708) .. controls (262.9329,230.1708) and (263.5172,229.5865) ..(263.5172,228.8656);
\path[fill=chfourfig5c0,fill opacity=0.5,draw=chfourfig5c0,line width=0.3985pt,draw opacity=0.5] (236.5381,210.9326) .. controls (236.5381,210.2119) and (235.9537,209.6275) ..(235.2329,209.6275) .. controls (234.5121,209.6275) and (233.9278,210.2119) ..(233.9278,210.9326) .. controls (233.9278,211.6534) and (234.5121,212.2378) ..(235.2329,212.2378) .. controls (235.9537,212.2378) and (236.5381,211.6534) ..(236.5381,210.9326);
\path[fill=chfourfig5c0,fill opacity=0.5,draw=chfourfig5c0,line width=0.3985pt,draw opacity=0.5] (227.9800,233.3788) .. controls (227.9800,232.6580) and (227.3957,232.0737) ..(226.6749,232.0737) .. controls (225.9541,232.0737) and (225.3698,232.6580) ..(225.3698,233.3788) .. controls (225.3698,234.0996) and (225.9541,234.6839) ..(226.6749,234.6839) .. controls (227.3957,234.6839) and (227.9800,234.0996) ..(227.9800,233.3788);
\path[fill=chfourfig5c0,fill opacity=0.5,draw=chfourfig5c0,line width=0.3985pt,draw opacity=0.5] (288.3146,262.2608) .. controls (288.3146,261.5400) and (287.7303,260.9557) ..(287.0095,260.9557) .. controls (286.2886,260.9557) and (285.7043,261.5400) ..(285.7043,262.2608) .. controls (285.7043,262.9816) and (286.2886,263.5659) ..(287.0095,263.5659) .. controls (287.7303,263.5659) and (288.3146,262.9816) ..(288.3146,262.2608);
\path[fill=chfourfig5c0,fill opacity=0.5,draw=chfourfig5c0,line width=0.3985pt,draw opacity=0.5] (272.3343,230.7386) .. controls (272.3343,230.0178) and (271.7500,229.4335) ..(271.0292,229.4335) .. controls (270.3084,229.4335) and (269.7241,230.0178) ..(269.7241,230.7386) .. controls (269.7241,231.4594) and (270.3084,232.0438) ..(271.0292,232.0438) .. controls (271.7500,232.0438) and (272.3343,231.4594) ..(272.3343,230.7386);
\path[fill=chfourfig5c0,fill opacity=0.5,draw=chfourfig5c0,line width=0.3985pt,draw opacity=0.5] (294.5014,190.8178) .. controls (294.5014,190.0970) and (293.9171,189.5127) ..(293.1964,189.5127) .. controls (292.4755,189.5127) and (291.8912,190.0970) ..(291.8912,190.8178) .. controls (291.8912,191.5386) and (292.4755,192.1230) ..(293.1964,192.1230) .. controls (293.9171,192.1230) and (294.5014,191.5386) ..(294.5014,190.8178);
\path[fill=chfourfig5c0,fill opacity=0.5,draw=chfourfig5c0,line width=0.3985pt,draw opacity=0.5] (269.5348,205.4531) .. controls (269.5348,204.7323) and (268.9504,204.1480) ..(268.2296,204.1480) .. controls (267.5088,204.1480) and (266.9245,204.7323) ..(266.9245,205.4531) .. controls (266.9245,206.1740) and (267.5088,206.7582) ..(268.2296,206.7582) .. controls (268.9504,206.7582) and (269.5348,206.1740) ..(269.5348,205.4531);
\path[fill=chfourfig5c0,fill opacity=0.5,draw=chfourfig5c0,line width=0.3985pt,draw opacity=0.5] (306.6062,214.2702) .. controls (306.6062,213.5494) and (306.0219,212.9651) ..(305.3011,212.9651) .. controls (304.5803,212.9651) and (303.9960,213.5494) ..(303.9960,214.2702) .. controls (303.9960,214.9910) and (304.5803,215.5753) ..(305.3011,215.5753) .. controls (306.0219,215.5753) and (306.6062,214.9910) ..(306.6062,214.2702);
\path[fill=chfourfig5c0,fill opacity=0.5,draw=chfourfig5c0,line width=0.3985pt,draw opacity=0.5] (243.5519,235.3614) .. controls (243.5519,234.6406) and (242.9675,234.0562) ..(242.2467,234.0562) .. controls (241.5259,234.0562) and (240.9416,234.6406) ..(240.9416,235.3614) .. controls (240.9416,236.0822) and (241.5259,236.6665) ..(242.2467,236.6665) .. controls (242.9675,236.6665) and (243.5519,236.0822) ..(243.5519,235.3614);
\path[fill=chfourfig5c0,fill opacity=0.5,draw=chfourfig5c0,line width=0.3985pt,draw opacity=0.5] (276.5186,210.1755) .. controls (276.5186,209.4547) and (275.9343,208.8704) ..(275.2135,208.8704) .. controls (274.4927,208.8704) and (273.9084,209.4547) ..(273.9084,210.1755) .. controls (273.9084,210.8963) and (274.4927,211.4806) ..(275.2135,211.4806) .. controls (275.9343,211.4806) and (276.5186,210.8963) ..(276.5186,210.1755);
\path[fill=chfourfig5c0,fill opacity=0.5,draw=chfourfig5c0,line width=0.3985pt,draw opacity=0.5] (315.6325,194.3745) .. controls (315.6325,193.6537) and (315.0482,193.0694) ..(314.3274,193.0694) .. controls (313.6066,193.0694) and (313.0222,193.6537) ..(313.0222,194.3745) .. controls (313.0222,195.0953) and (313.6066,195.6796) ..(314.3274,195.6796) .. controls (315.0482,195.6796) and (315.6325,195.0953) ..(315.6325,194.3745);
\path[fill=chfourfig5c0,fill opacity=0.5,draw=chfourfig5c0,line width=0.3985pt,draw opacity=0.5] (241.5095,220.6862) .. controls (241.5095,219.9654) and (240.9252,219.3811) ..(240.2043,219.3811) .. controls (239.4836,219.3811) and (238.8992,219.9654) ..(238.8992,220.6862) .. controls (238.8992,221.4070) and (239.4836,221.9913) ..(240.2043,221.9913) .. controls (240.9252,221.9913) and (241.5095,221.4070) ..(241.5095,220.6862);
\path[fill=chfourfig5c0,fill opacity=0.5,draw=chfourfig5c0,line width=0.3985pt,draw opacity=0.5] (255.0987,193.7070) .. controls (255.0987,192.9862) and (254.5144,192.4019) ..(253.7936,192.4019) .. controls (253.0728,192.4019) and (252.4885,192.9862) ..(252.4885,193.7070) .. controls (252.4885,194.4278) and (253.0728,195.0121) ..(253.7936,195.0121) .. controls (254.5144,195.0121) and (255.0987,194.4278) ..(255.0987,193.7070);
\path[fill=chfourfig5c0,fill opacity=0.5,draw=chfourfig5c0,line width=0.3985pt,draw opacity=0.5] (306.9549,233.0699) .. controls (306.9549,232.3491) and (306.3706,231.7648) ..(305.6498,231.7648) .. controls (304.9290,231.7648) and (304.3447,232.3491) ..(304.3447,233.0699) .. controls (304.3447,233.7907) and (304.9290,234.3751) ..(305.6498,234.3751) .. controls (306.3706,234.3751) and (306.9549,233.7907) ..(306.9549,233.0699);
\path[fill=chfourfig5c0,fill opacity=0.5,draw=chfourfig5c0,line width=0.3985pt,draw opacity=0.5] (307.1940,211.0124) .. controls (307.1940,210.2915) and (306.6097,209.7072) ..(305.8889,209.7072) .. controls (305.1681,209.7072) and (304.5838,210.2915) ..(304.5838,211.0124) .. controls (304.5838,211.7332) and (305.1681,212.3175) ..(305.8889,212.3175) .. controls (306.6097,212.3175) and (307.1940,211.7332) ..(307.1940,211.0124);
\path[fill=chfourfig5c0,fill opacity=0.5,draw=chfourfig5c0,line width=0.3985pt,draw opacity=0.5] (299.3334,212.2876) .. controls (299.3334,211.5668) and (298.7491,210.9825) ..(298.0283,210.9825) .. controls (297.3075,210.9825) and (296.7231,211.5668) ..(296.7231,212.2876) .. controls (296.7231,213.0084) and (297.3075,213.5927) ..(298.0283,213.5927) .. controls (298.7491,213.5927) and (299.3334,213.0084) ..(299.3334,212.2876);
\path[fill=chfourfig5c0,fill opacity=0.5,draw=chfourfig5c0,line width=0.3985pt,draw opacity=0.5] (318.2527,233.9765) .. controls (318.2527,233.2557) and (317.6684,232.6714) ..(316.9476,232.6714) .. controls (316.2268,232.6714) and (315.6425,233.2557) ..(315.6425,233.9765) .. controls (315.6425,234.6973) and (316.2268,235.2816) ..(316.9476,235.2816) .. controls (317.6684,235.2816) and (318.2527,234.6973) ..(318.2527,233.9765);
\path[fill=chfourfig5c0,fill opacity=0.5,draw=chfourfig5c0,line width=0.3985pt,draw opacity=0.5] (307.8815,213.9713) .. controls (307.8815,213.2505) and (307.2971,212.6662) ..(306.5763,212.6662) .. controls (305.8555,212.6662) and (305.2712,213.2505) ..(305.2712,213.9713) .. controls (305.2712,214.6921) and (305.8555,215.2764) ..(306.5763,215.2764) .. controls (307.2971,215.2764) and (307.8815,214.6921) ..(307.8815,213.9713);
\path[fill=chfourfig5c0,fill opacity=0.5,draw=chfourfig5c0,line width=0.3985pt,draw opacity=0.5] (273.6494,225.0499) .. controls (273.6494,224.3291) and (273.0650,223.7448) ..(272.3443,223.7448) .. controls (271.6234,223.7448) and (271.0391,224.3291) ..(271.0391,225.0499) .. controls (271.0391,225.7707) and (271.6234,226.3550) ..(272.3443,226.3550) .. controls (273.0650,226.3550) and (273.6494,225.7707) ..(273.6494,225.0499);
\path[fill=chfourfig5c0,fill opacity=0.5,draw=chfourfig5c0,line width=0.3985pt,draw opacity=0.5] (267.7713,215.9838) .. controls (267.7713,215.2630) and (267.1870,214.6786) ..(266.4662,214.6786) .. controls (265.7454,214.6786) and (265.1611,215.2630) ..(265.1611,215.9838) .. controls (265.1611,216.7046) and (265.7454,217.2889) ..(266.4662,217.2889) .. controls (267.1870,217.2889) and (267.7713,216.7046) ..(267.7713,215.9838);
\path[fill=chfourfig5c0,fill opacity=0.5,draw=chfourfig5c0,line width=0.3985pt,draw opacity=0.5] (289.2809,200.6311) .. controls (289.2809,199.9103) and (288.6966,199.3260) ..(287.9759,199.3260) .. controls (287.2550,199.3260) and (286.6707,199.9103) ..(286.6707,200.6311) .. controls (286.6707,201.3520) and (287.2550,201.9363) ..(287.9759,201.9363) .. controls (288.6966,201.9363) and (289.2809,201.3520) ..(289.2809,200.6311);
\path[fill=chfourfig5c0,fill opacity=0.5,draw=chfourfig5c0,line width=0.3985pt,draw opacity=0.5] (289.6994,221.6327) .. controls (289.6994,220.9119) and (289.1151,220.3275) ..(288.3943,220.3275) .. controls (287.6735,220.3275) and (287.0891,220.9119) ..(287.0891,221.6327) .. controls (287.0891,222.3535) and (287.6735,222.9378) ..(288.3943,222.9378) .. controls (289.1151,222.9378) and (289.6994,222.3535) ..(289.6994,221.6327);
\path[fill=chfourfig5c0,fill opacity=0.5,draw=chfourfig5c0,line width=0.3985pt,draw opacity=0.5] (285.1364,258.2259) .. controls (285.1364,257.5051) and (284.5521,256.9208) ..(283.8313,256.9208) .. controls (283.1105,256.9208) and (282.5262,257.5051) ..(282.5262,258.2259) .. controls (282.5262,258.9467) and (283.1105,259.5310) ..(283.8313,259.5310) .. controls (284.5521,259.5310) and (285.1364,258.9467) ..(285.1364,258.2259);
\path[fill=chfourfig5c0,fill opacity=0.5,draw=chfourfig5c0,line width=0.3985pt,draw opacity=0.5] (304.5339,237.1248) .. controls (304.5339,236.4040) and (303.9496,235.8196) ..(303.2289,235.8196) .. controls (302.5080,235.8196) and (301.9237,236.4040) ..(301.9237,237.1248) .. controls (301.9237,237.8456) and (302.5080,238.4299) ..(303.2289,238.4299) .. controls (303.9496,238.4299) and (304.5339,237.8456) ..(304.5339,237.1248);
\path[fill=chfourfig5c0,fill opacity=0.5,draw=chfourfig5c0,line width=0.3985pt,draw opacity=0.5] (283.5225,178.4540) .. controls (283.5225,177.7332) and (282.9382,177.1489) ..(282.2173,177.1489) .. controls (281.4966,177.1489) and (280.9123,177.7332) ..(280.9123,178.4540) .. controls (280.9123,179.1748) and (281.4966,179.7592) ..(282.2173,179.7592) .. controls (282.9382,179.7592) and (283.5225,179.1748) ..(283.5225,178.4540);
\path[fill=chfourfig5c0,fill opacity=0.5,draw=chfourfig5c0,line width=0.3985pt,draw opacity=0.5] (301.0271,202.2949) .. controls (301.0271,201.5741) and (300.4427,200.9898) ..(299.7219,200.9898) .. controls (299.0011,200.9898) and (298.4168,201.5741) ..(298.4168,202.2949) .. controls (298.4168,203.0157) and (299.0011,203.6000) ..(299.7219,203.6000) .. controls (300.4427,203.6000) and (301.0271,203.0157) ..(301.0271,202.2949);
\path[fill=chfourfig5c0,fill opacity=0.5,draw=chfourfig5c0,line width=0.3985pt,draw opacity=0.5] (264.0652,182.0307) .. controls (264.0652,181.3098) and (263.4809,180.7256) ..(262.7601,180.7256) .. controls (262.0393,180.7256) and (261.4550,181.3098) ..(261.4550,182.0307) .. controls (261.4550,182.7514) and (262.0393,183.3358) ..(262.7601,183.3358) .. controls (263.4809,183.3358) and (264.0652,182.7514) ..(264.0652,182.0307);
\path[fill=chfourfig5c0,fill opacity=0.5,draw=chfourfig5c0,line width=0.3985pt,draw opacity=0.5] (272.7527,240.2132) .. controls (272.7527,239.4924) and (272.1684,238.9081) ..(271.4476,238.9081) .. controls (270.7268,238.9081) and (270.1425,239.4924) ..(270.1425,240.2132) .. controls (270.1425,240.9340) and (270.7268,241.5184) ..(271.4476,241.5184) .. controls (272.1684,241.5184) and (272.7527,240.9340) ..(272.7527,240.2132);
\path[fill=chfourfig5c0,fill opacity=0.5,draw=chfourfig5c0,line width=0.3985pt,draw opacity=0.5] (312.5440,227.9690) .. controls (312.5440,227.2482) and (311.9597,226.6639) ..(311.2389,226.6639) .. controls (310.5181,226.6639) and (309.9338,227.2482) ..(309.9338,227.9690) .. controls (309.9338,228.6898) and (310.5181,229.2741) ..(311.2389,229.2741) .. controls (311.9597,229.2741) and (312.5440,228.6898) ..(312.5440,227.9690);
\path[fill=chfourfig5c0,fill opacity=0.5,draw=chfourfig5c0,line width=0.3985pt,draw opacity=0.5] (311.3983,205.5727) .. controls (311.3983,204.8519) and (310.8140,204.2676) ..(310.0932,204.2676) .. controls (309.3724,204.2676) and (308.7881,204.8519) ..(308.7881,205.5727) .. controls (308.7881,206.2935) and (309.3724,206.8778) ..(310.0932,206.8778) .. controls (310.8140,206.8778) and (311.3983,206.2935) ..(311.3983,205.5727);
\path[fill=chfourfig5c0,fill opacity=0.5,draw=chfourfig5c0,line width=0.3985pt,draw opacity=0.5] (231.7161,219.0922) .. controls (231.7161,218.3714) and (231.1318,217.7870) ..(230.4109,217.7870) .. controls (229.6902,217.7870) and (229.1058,218.3714) ..(229.1058,219.0922) .. controls (229.1058,219.8130) and (229.6902,220.3973) ..(230.4109,220.3973) .. controls (231.1318,220.3973) and (231.7161,219.8130) ..(231.7161,219.0922);
\path[fill=chfourfig5c0,fill opacity=0.5,draw=chfourfig5c0,line width=0.3985pt,draw opacity=0.5] (292.6284,198.6785) .. controls (292.6284,197.9576) and (292.0441,197.3733) ..(291.3233,197.3733) .. controls (290.6025,197.3733) and (290.0182,197.9576) ..(290.0182,198.6785) .. controls (290.0182,199.3992) and (290.6025,199.9836) ..(291.3233,199.9836) .. controls (292.0441,199.9836) and (292.6284,199.3992) ..(292.6284,198.6785);
\path[fill=chfourfig5c0,fill opacity=0.5,draw=chfourfig5c0,line width=0.3985pt,draw opacity=0.5] (315.9513,209.0098) .. controls (315.9513,208.2890) and (315.3670,207.7047) ..(314.6462,207.7047) .. controls (313.9254,207.7047) and (313.3410,208.2890) ..(313.3410,209.0098) .. controls (313.3410,209.7307) and (313.9254,210.3149) ..(314.6462,210.3149) .. controls (315.3670,210.3149) and (315.9513,209.7307) ..(315.9513,209.0098);
\path[fill=chfourfig5c0,fill opacity=0.5,draw=chfourfig5c0,line width=0.3985pt,draw opacity=0.5] (264.8722,224.7709) .. controls (264.8722,224.0501) and (264.2878,223.4658) ..(263.5671,223.4658) .. controls (262.8463,223.4658) and (262.2619,224.0501) ..(262.2619,224.7709) .. controls (262.2619,225.4918) and (262.8463,226.0760) ..(263.5671,226.0760) .. controls (264.2878,226.0760) and (264.8722,225.4918) ..(264.8722,224.7709);
\path[fill=chfourfig5c0,fill opacity=0.5,draw=chfourfig5c0,line width=0.3985pt,draw opacity=0.5] (228.6675,220.6563) .. controls (228.6675,219.9355) and (228.0832,219.3512) ..(227.3624,219.3512) .. controls (226.6415,219.3512) and (226.0573,219.9355) ..(226.0573,220.6563) .. controls (226.0573,221.3771) and (226.6415,221.9614) ..(227.3624,221.9614) .. controls (228.0832,221.9614) and (228.6675,221.3771) ..(228.6675,220.6563);
\path[fill=chfourfig5c0,fill opacity=0.5,draw=chfourfig5c0,line width=0.3985pt,draw opacity=0.5] (284.2000,248.6617) .. controls (284.2000,247.9409) and (283.6156,247.3565) ..(282.8948,247.3565) .. controls (282.1740,247.3565) and (281.5897,247.9409) ..(281.5897,248.6617) .. controls (281.5897,249.3825) and (282.1740,249.9668) ..(282.8948,249.9668) .. controls (283.6156,249.9668) and (284.2000,249.3825) ..(284.2000,248.6617);
\path[fill=chfourfig5c0,fill opacity=0.5,draw=chfourfig5c0,line width=0.3985pt,draw opacity=0.5] (275.5522,251.5608) .. controls (275.5522,250.8400) and (274.9679,250.2557) ..(274.2472,250.2557) .. controls (273.5264,250.2557) and (272.9420,250.8400) ..(272.9420,251.5608) .. controls (272.9420,252.2816) and (273.5264,252.8659) ..(274.2472,252.8659) .. controls (274.9679,252.8659) and (275.5522,252.2816) ..(275.5522,251.5608);
\path[fill=chfourfig5c0,fill opacity=0.5,draw=chfourfig5c0,line width=0.3985pt,draw opacity=0.5] (310.1430,229.0848) .. controls (310.1430,228.3640) and (309.5587,227.7797) ..(308.8379,227.7797) .. controls (308.1171,227.7797) and (307.5327,228.3640) ..(307.5327,229.0848) .. controls (307.5327,229.8056) and (308.1171,230.3900) ..(308.8379,230.3900) .. controls (309.5587,230.3900) and (310.1430,229.8056) ..(310.1430,229.0848);
\path[fill=chfourfig5c0,fill opacity=0.5,draw=chfourfig5c0,line width=0.3985pt,draw opacity=0.5] (296.3047,192.7506) .. controls (296.3047,192.0298) and (295.7204,191.4455) ..(294.9996,191.4455) .. controls (294.2788,191.4455) and (293.6945,192.0298) ..(293.6945,192.7506) .. controls (293.6945,193.4714) and (294.2788,194.0557) ..(294.9996,194.0557) .. controls (295.7204,194.0557) and (296.3047,193.4714) ..(296.3047,192.7506);
\path[fill=chfourfig5c0,fill opacity=0.5,draw=chfourfig5c0,line width=0.3985pt,draw opacity=0.5] (271.7465,222.6987) .. controls (271.7465,221.9779) and (271.1622,221.3936) ..(270.4413,221.3936) .. controls (269.7206,221.3936) and (269.1363,221.9779) ..(269.1363,222.6987) .. controls (269.1363,223.4195) and (269.7206,224.0038) ..(270.4413,224.0038) .. controls (271.1622,224.0038) and (271.7465,223.4195) ..(271.7465,222.6987);
\path[fill=chfourfig5c0,fill opacity=0.5,draw=chfourfig5c0,line width=0.3985pt,draw opacity=0.5] (313.6599,196.1080) .. controls (313.6599,195.3873) and (313.0755,194.8029) ..(312.3547,194.8029) .. controls (311.6339,194.8029) and (311.0496,195.3873) ..(311.0496,196.1080) .. controls (311.0496,196.8289) and (311.6339,197.4132) ..(312.3547,197.4132) .. controls (313.0755,197.4132) and (313.6599,196.8289) ..(313.6599,196.1080);
\path[fill=chfourfig5c0,fill opacity=0.5,draw=chfourfig5c0,line width=0.3985pt,draw opacity=0.5] (281.3805,207.9139) .. controls (281.3805,207.1931) and (280.7961,206.6088) ..(280.0754,206.6088) .. controls (279.3546,206.6088) and (278.7702,207.1931) ..(278.7702,207.9139) .. controls (278.7702,208.6347) and (279.3546,209.2191) ..(280.0754,209.2191) .. controls (280.7961,209.2191) and (281.3805,208.6347) ..(281.3805,207.9139);
\path[fill=chfourfig5c0,fill opacity=0.5,draw=chfourfig5c0,line width=0.3985pt,draw opacity=0.5] (273.9781,200.7407) .. controls (273.9781,200.0199) and (273.3938,199.4356) ..(272.6730,199.4356) .. controls (271.9522,199.4356) and (271.3679,200.0199) ..(271.3679,200.7407) .. controls (271.3679,201.4615) and (271.9522,202.0458) ..(272.6730,202.0458) .. controls (273.3938,202.0458) and (273.9781,201.4615) ..(273.9781,200.7407);
\path[fill=chfourfig5c0,fill opacity=0.5,draw=chfourfig5c0,line width=0.3985pt,draw opacity=0.5] (289.3208,233.3090) .. controls (289.3208,232.5882) and (288.7365,232.0039) ..(288.0157,232.0039) .. controls (287.2949,232.0039) and (286.7106,232.5882) ..(286.7106,233.3090) .. controls (286.7106,234.0298) and (287.2949,234.6142) ..(288.0157,234.6142) .. controls (288.7365,234.6142) and (289.3208,234.0298) ..(289.3208,233.3090);
\path[fill=chfourfig5c0,fill opacity=0.5,draw=chfourfig5c0,line width=0.3985pt,draw opacity=0.5] (330.9054,211.9389) .. controls (330.9054,211.2181) and (330.3211,210.6338) ..(329.6003,210.6338) .. controls (328.8795,210.6338) and (328.2952,211.2181) ..(328.2952,211.9389) .. controls (328.2952,212.6597) and (328.8795,213.2440) ..(329.6003,213.2440) .. controls (330.3211,213.2440) and (330.9054,212.6597) ..(330.9054,211.9389);
\path[fill=chfourfig5c0,fill opacity=0.5,draw=chfourfig5c0,line width=0.3985pt,draw opacity=0.5] (275.5124,231.1371) .. controls (275.5124,230.4164) and (274.9281,229.8320) ..(274.2073,229.8320) .. controls (273.4865,229.8320) and (272.9022,230.4164) ..(272.9022,231.1371) .. controls (272.9022,231.8580) and (273.4865,232.4423) ..(274.2073,232.4423) .. controls (274.9281,232.4423) and (275.5124,231.8580) ..(275.5124,231.1371);
\path[fill=chfourfig5c0,fill opacity=0.5,draw=chfourfig5c0,line width=0.3985pt,draw opacity=0.5] (301.2064,237.4834) .. controls (301.2064,236.7626) and (300.6221,236.1783) ..(299.9013,236.1783) .. controls (299.1805,236.1783) and (298.5962,236.7626) ..(298.5962,237.4834) .. controls (298.5962,238.2042) and (299.1805,238.7885) ..(299.9013,238.7885) .. controls (300.6221,238.7885) and (301.2064,238.2042) ..(301.2064,237.4834);
\path[fill=chfourfig5c0,fill opacity=0.5,draw=chfourfig5c0,line width=0.3985pt,draw opacity=0.5] (291.4329,183.5251) .. controls (291.4329,182.8043) and (290.8486,182.2200) ..(290.1278,182.2200) .. controls (289.4070,182.2200) and (288.8227,182.8043) ..(288.8227,183.5251) .. controls (288.8227,184.2459) and (289.4070,184.8302) ..(290.1278,184.8302) .. controls (290.8486,184.8302) and (291.4329,184.2459) ..(291.4329,183.5251);
\path[fill=chfourfig5c0,fill opacity=0.5,draw=chfourfig5c0,line width=0.3985pt,draw opacity=0.5] (311.3684,250.2159) .. controls (311.3684,249.4951) and (310.7841,248.9107) ..(310.0633,248.9107) .. controls (309.3425,248.9107) and (308.7582,249.4951) ..(308.7582,250.2159) .. controls (308.7582,250.9366) and (309.3425,251.5210) ..(310.0633,251.5210) .. controls (310.7841,251.5210) and (311.3684,250.9366) ..(311.3684,250.2159);
\path[fill=chfourfig5c0,fill opacity=0.5,draw=chfourfig5c0,line width=0.3985pt,draw opacity=0.5] (257.4998,220.1682) .. controls (257.4998,219.4473) and (256.9154,218.8630) ..(256.1946,218.8630) .. controls (255.4738,218.8630) and (254.8895,219.4473) ..(254.8895,220.1682) .. controls (254.8895,220.8889) and (255.4738,221.4733) ..(256.1946,221.4733) .. controls (256.9154,221.4733) and (257.4998,220.8889) ..(257.4998,220.1682);
\path[fill=chfourfig5c0,fill opacity=0.5,draw=chfourfig5c0,line width=0.3985pt,draw opacity=0.5] (314.8653,216.1631) .. controls (314.8653,215.4423) and (314.2810,214.8580) ..(313.5602,214.8580) .. controls (312.8394,214.8580) and (312.2551,215.4423) ..(312.2551,216.1631) .. controls (312.2551,216.8839) and (312.8394,217.4682) ..(313.5602,217.4682) .. controls (314.2810,217.4682) and (314.8653,216.8839) ..(314.8653,216.1631);
\path[fill=chfourfig5c0,fill opacity=0.5,draw=chfourfig5c0,line width=0.3985pt,draw opacity=0.5] (274.7752,244.5869) .. controls (274.7752,243.8661) and (274.1909,243.2818) ..(273.4700,243.2818) .. controls (272.7493,243.2818) and (272.1649,243.8661) ..(272.1649,244.5869) .. controls (272.1649,245.3077) and (272.7493,245.8920) ..(273.4700,245.8920) .. controls (274.1909,245.8920) and (274.7752,245.3077) ..(274.7752,244.5869);
\path[fill=chfourfig5c0,fill opacity=0.5,draw=chfourfig5c0,line width=0.3985pt,draw opacity=0.5] (275.3231,203.7495) .. controls (275.3231,203.0287) and (274.7388,202.4444) ..(274.0180,202.4444) .. controls (273.2972,202.4444) and (272.7129,203.0287) ..(272.7129,203.7495) .. controls (272.7129,204.4703) and (273.2972,205.0546) ..(274.0180,205.0546) .. controls (274.7388,205.0546) and (275.3231,204.4703) ..(275.3231,203.7495);
\path[fill=chfourfig5c0,fill opacity=0.5,draw=chfourfig5c0,line width=0.3985pt,draw opacity=0.5] (293.7841,246.0315) .. controls (293.7841,245.3107) and (293.1998,244.7264) ..(292.4790,244.7264) .. controls (291.7582,244.7264) and (291.1739,245.3107) ..(291.1739,246.0315) .. controls (291.1739,246.7523) and (291.7582,247.3366) ..(292.4790,247.3366) .. controls (293.1998,247.3366) and (293.7841,246.7523) ..(293.7841,246.0315);
\path[fill=chfourfig5c0,fill opacity=0.5,draw=chfourfig5c0,line width=0.3985pt,draw opacity=0.5] (275.1637,184.2623) .. controls (275.1637,183.5415) and (274.5794,182.9572) ..(273.8586,182.9572) .. controls (273.1378,182.9572) and (272.5535,183.5415) ..(272.5535,184.2623) .. controls (272.5535,184.9831) and (273.1378,185.5674) ..(273.8586,185.5674) .. controls (274.5794,185.5674) and (275.1637,184.9831) ..(275.1637,184.2623);
\path[fill=chfourfig5c0,fill opacity=0.5,draw=chfourfig5c0,line width=0.3985pt,draw opacity=0.5] (243.6614,227.0225) .. controls (243.6614,226.3017) and (243.0771,225.7174) ..(242.3563,225.7174) .. controls (241.6355,225.7174) and (241.0512,226.3017) ..(241.0512,227.0225) .. controls (241.0512,227.7433) and (241.6355,228.3276) ..(242.3563,228.3276) .. controls (243.0771,228.3276) and (243.6614,227.7433) ..(243.6614,227.0225);
\path[fill=chfourfig5c0,fill opacity=0.5,draw=chfourfig5c0,line width=0.3985pt,draw opacity=0.5] (286.0730,224.1931) .. controls (286.0730,223.4723) and (285.4886,222.8880) ..(284.7678,222.8880) .. controls (284.0470,222.8880) and (283.4627,223.4723) ..(283.4627,224.1931) .. controls (283.4627,224.9139) and (284.0470,225.4982) ..(284.7678,225.4982) .. controls (285.4886,225.4982) and (286.0730,224.9139) ..(286.0730,224.1931);
\path[fill=chfourfig5c0,fill opacity=0.5,draw=chfourfig5c0,line width=0.3985pt,draw opacity=0.5] (261.1859,201.9562) .. controls (261.1859,201.2354) and (260.6016,200.6511) ..(259.8808,200.6511) .. controls (259.1600,200.6511) and (258.5757,201.2354) ..(258.5757,201.9562) .. controls (258.5757,202.6770) and (259.1600,203.2613) ..(259.8808,203.2613) .. controls (260.6016,203.2613) and (261.1859,202.6770) ..(261.1859,201.9562);
\path[fill=chfourfig5c0,fill opacity=0.5,draw=chfourfig5c0,line width=0.3985pt,draw opacity=0.5] (229.1158,223.8344) .. controls (229.1158,223.1136) and (228.5315,222.5293) ..(227.8107,222.5293) .. controls (227.0899,222.5293) and (226.5056,223.1136) ..(226.5056,223.8344) .. controls (226.5056,224.5552) and (227.0899,225.1396) ..(227.8107,225.1396) .. controls (228.5315,225.1396) and (229.1158,224.5552) ..(229.1158,223.8344);
\path[fill=chfourfig5c0,fill opacity=0.5,draw=chfourfig5c0,line width=0.3985pt,draw opacity=0.5] (324.3399,203.9487) .. controls (324.3399,203.2279) and (323.7556,202.6436) ..(323.0348,202.6436) .. controls (322.3140,202.6436) and (321.7297,203.2279) ..(321.7297,203.9487) .. controls (321.7297,204.6696) and (322.3140,205.2539) ..(323.0348,205.2539) .. controls (323.7556,205.2539) and (324.3399,204.6696) ..(324.3399,203.9487);
\path[fill=chfourfig5c0,fill opacity=0.5,draw=chfourfig5c0,line width=0.3985pt,draw opacity=0.5] (238.7100,231.0276) .. controls (238.7100,230.3068) and (238.1256,229.7224) ..(237.4048,229.7224) .. controls (236.6840,229.7224) and (236.0997,230.3068) ..(236.0997,231.0276) .. controls (236.0997,231.7484) and (236.6840,232.3327) ..(237.4048,232.3327) .. controls (238.1256,232.3327) and (238.7100,231.7484) ..(238.7100,231.0276);
\path[fill=chfourfig5c0,fill opacity=0.5,draw=chfourfig5c0,line width=0.3985pt,draw opacity=0.5] (280.3045,217.5181) .. controls (280.3045,216.7972) and (279.7202,216.2129) ..(278.9994,216.2129) .. controls (278.2786,216.2129) and (277.6943,216.7972) ..(277.6943,217.5181) .. controls (277.6943,218.2388) and (278.2786,218.8232) ..(278.9994,218.8232) .. controls (279.7202,218.8232) and (280.3045,218.2388) ..(280.3045,217.5181);
\path[fill=chfourfig5c0,fill opacity=0.5,draw=chfourfig5c0,line width=0.3985pt,draw opacity=0.5] (263.7464,228.0387) .. controls (263.7464,227.3179) and (263.1621,226.7336) ..(262.4413,226.7336) .. controls (261.7205,226.7336) and (261.1361,227.3179) ..(261.1361,228.0387) .. controls (261.1361,228.7595) and (261.7205,229.3438) ..(262.4413,229.3438) .. controls (263.1621,229.3438) and (263.7464,228.7595) ..(263.7464,228.0387);
\path[fill=chfourfig5c0,fill opacity=0.5,draw=chfourfig5c0,line width=0.3985pt,draw opacity=0.5] (247.1584,207.9538) .. controls (247.1584,207.2330) and (246.5740,206.6487) ..(245.8533,206.6487) .. controls (245.1324,206.6487) and (244.5482,207.2330) ..(244.5482,207.9538) .. controls (244.5482,208.6746) and (245.1324,209.2589) ..(245.8533,209.2589) .. controls (246.5740,209.2589) and (247.1584,208.6746) ..(247.1584,207.9538);
\path[fill=chfourfig5c0,fill opacity=0.5,draw=chfourfig5c0,line width=0.3985pt,draw opacity=0.5] (258.2469,231.3364) .. controls (258.2469,230.6156) and (257.6626,230.0313) ..(256.9418,230.0313) .. controls (256.2210,230.0313) and (255.6367,230.6156) ..(255.6367,231.3364) .. controls (255.6367,232.0572) and (256.2210,232.6415) ..(256.9418,232.6415) .. controls (257.6626,232.6415) and (258.2469,232.0572) ..(258.2469,231.3364);
\path[fill=chfourfig5c0,fill opacity=0.5,draw=chfourfig5c0,line width=0.3985pt,draw opacity=0.5] (283.8811,231.0774) .. controls (283.8811,230.3566) and (283.2968,229.7722) ..(282.5760,229.7722) .. controls (281.8552,229.7722) and (281.2709,230.3566) ..(281.2709,231.0774) .. controls (281.2709,231.7982) and (281.8552,232.3825) ..(282.5760,232.3825) .. controls (283.2968,232.3825) and (283.8811,231.7982) ..(283.8811,231.0774);
\path[fill=chfourfig5c0,fill opacity=0.5,draw=chfourfig5c0,line width=0.3985pt,draw opacity=0.5] (292.0406,182.4890) .. controls (292.0406,181.7681) and (291.4563,181.1838) ..(290.7355,181.1838) .. controls (290.0147,181.1838) and (289.4304,181.7681) ..(289.4304,182.4890) .. controls (289.4304,183.2097) and (290.0147,183.7941) ..(290.7355,183.7941) .. controls (291.4563,183.7941) and (292.0406,183.2097) ..(292.0406,182.4890);
\path[fill=chfourfig5c0,fill opacity=0.5,draw=chfourfig5c0,line width=0.3985pt,draw opacity=0.5] (256.7625,217.4184) .. controls (256.7625,216.6976) and (256.1782,216.1133) ..(255.4574,216.1133) .. controls (254.7366,216.1133) and (254.1522,216.6976) ..(254.1522,217.4184) .. controls (254.1522,218.1392) and (254.7366,218.7235) ..(255.4574,218.7235) .. controls (256.1782,218.7235) and (256.7625,218.1392) ..(256.7625,217.4184);
\path[fill=chfourfig5c0,fill opacity=0.5,draw=chfourfig5c0,line width=0.3985pt,draw opacity=0.5] (302.5713,220.6165) .. controls (302.5713,219.8957) and (301.9870,219.3113) ..(301.2662,219.3113) .. controls (300.5454,219.3113) and (299.9611,219.8957) ..(299.9611,220.6165) .. controls (299.9611,221.3373) and (300.5454,221.9216) ..(301.2662,221.9216) .. controls (301.9870,221.9216) and (302.5713,221.3373) ..(302.5713,220.6165);
\path[fill=chfourfig5c0,fill opacity=0.5,draw=chfourfig5c0,line width=0.3985pt,draw opacity=0.5] (271.6568,183.3657) .. controls (271.6568,182.6449) and (271.0725,182.0606) ..(270.3517,182.0606) .. controls (269.6309,182.0606) and (269.0466,182.6449) ..(269.0466,183.3657) .. controls (269.0466,184.0865) and (269.6309,184.6708) ..(270.3517,184.6708) .. controls (271.0725,184.6708) and (271.6568,184.0865) ..(271.6568,183.3657);
\path[fill=chfourfig5c0,fill opacity=0.5,draw=chfourfig5c0,line width=0.3985pt,draw opacity=0.5] (304.7531,243.6902) .. controls (304.7531,242.9694) and (304.1688,242.3851) ..(303.4480,242.3851) .. controls (302.7272,242.3851) and (302.1429,242.9694) ..(302.1429,243.6902) .. controls (302.1429,244.4110) and (302.7272,244.9954) ..(303.4480,244.9954) .. controls (304.1688,244.9954) and (304.7531,244.4110) ..(304.7531,243.6902);
\path[fill=chfourfig5c0,fill opacity=0.5,draw=chfourfig5c0,line width=0.3985pt,draw opacity=0.5] (270.4015,216.4122) .. controls (270.4015,215.6914) and (269.8172,215.1071) ..(269.0964,215.1071) .. controls (268.3756,215.1071) and (267.7913,215.6914) ..(267.7913,216.4122) .. controls (267.7913,217.1330) and (268.3756,217.7173) ..(269.0964,217.7173) .. controls (269.8172,217.7173) and (270.4015,217.1330) ..(270.4015,216.4122);
\path[fill=chfourfig5c0,fill opacity=0.5,draw=chfourfig5c0,line width=0.3985pt,draw opacity=0.5] (268.1898,241.3789) .. controls (268.1898,240.6581) and (267.6055,240.0737) ..(266.8846,240.0737) .. controls (266.1639,240.0737) and (265.5796,240.6581) ..(265.5796,241.3789) .. controls (265.5796,242.0997) and (266.1639,242.6840) ..(266.8846,242.6840) .. controls (267.6055,242.6840) and (268.1898,242.0997) ..(268.1898,241.3789);
\path[fill=chfourfig5c0,fill opacity=0.5,draw=chfourfig5c0,line width=0.3985pt,draw opacity=0.5] (239.1184,239.9841) .. controls (239.1184,239.2633) and (238.5341,238.6790) ..(237.8133,238.6790) .. controls (237.0925,238.6790) and (236.5082,239.2633) ..(236.5082,239.9841) .. controls (236.5082,240.7049) and (237.0925,241.2892) ..(237.8133,241.2892) .. controls (238.5341,241.2892) and (239.1184,240.7049) ..(239.1184,239.9841);
\path[fill=chfourfig5c0,fill opacity=0.5,draw=chfourfig5c0,line width=0.3985pt,draw opacity=0.5] (273.6593,229.5232) .. controls (273.6593,228.8024) and (273.0750,228.2181) ..(272.3542,228.2181) .. controls (271.6334,228.2181) and (271.0491,228.8024) ..(271.0491,229.5232) .. controls (271.0491,230.2440) and (271.6334,230.8283) ..(272.3542,230.8283) .. controls (273.0750,230.8283) and (273.6593,230.2440) ..(273.6593,229.5232);
\path[fill=chfourfig5c0,fill opacity=0.5,draw=chfourfig5c0,line width=0.3985pt,draw opacity=0.5] (281.6495,195.0620) .. controls (281.6495,194.3412) and (281.0652,193.7568) ..(280.3444,193.7568) .. controls (279.6235,193.7568) and (279.0392,194.3412) ..(279.0392,195.0620) .. controls (279.0392,195.7828) and (279.6235,196.3671) ..(280.3444,196.3671) .. controls (281.0652,196.3671) and (281.6495,195.7828) ..(281.6495,195.0620);
\path[fill=chfourfig5c0,fill opacity=0.5,draw=chfourfig5c0,line width=0.3985pt,draw opacity=0.5] (292.5886,258.7739) .. controls (292.5886,258.0530) and (292.0043,257.4688) ..(291.2835,257.4688) .. controls (290.5627,257.4688) and (289.9783,258.0530) ..(289.9783,258.7739) .. controls (289.9783,259.4947) and (290.5627,260.0790) ..(291.2835,260.0790) .. controls (292.0043,260.0790) and (292.5886,259.4947) ..(292.5886,258.7739);
\path[fill=chfourfig5c0,fill opacity=0.5,draw=chfourfig5c0,line width=0.3985pt,draw opacity=0.5] (248.8321,229.6726) .. controls (248.8321,228.9518) and (248.2478,228.3675) ..(247.5270,228.3675) .. controls (246.8062,228.3675) and (246.2219,228.9518) ..(246.2219,229.6726) .. controls (246.2219,230.3934) and (246.8062,230.9777) ..(247.5270,230.9777) .. controls (248.2478,230.9777) and (248.8321,230.3934) ..(248.8321,229.6726);
\path[fill=chfourfig5c0,fill opacity=0.5,draw=chfourfig5c0,line width=0.3985pt,draw opacity=0.5] (304.1454,238.0613) .. controls (304.1454,237.3405) and (303.5611,236.7561) ..(302.8403,236.7561) .. controls (302.1195,236.7561) and (301.5352,237.3405) ..(301.5352,238.0613) .. controls (301.5352,238.7821) and (302.1195,239.3664) ..(302.8403,239.3664) .. controls (303.5611,239.3664) and (304.1454,238.7821) ..(304.1454,238.0613);
\path[fill=chfourfig5c0,fill opacity=0.5,draw=chfourfig5c0,line width=0.3985pt,draw opacity=0.5] (281.7093,236.5270) .. controls (281.7093,235.8062) and (281.1249,235.2219) ..(280.4041,235.2219) .. controls (279.6833,235.2219) and (279.0990,235.8062) ..(279.0990,236.5270) .. controls (279.0990,237.2478) and (279.6833,237.8321) ..(280.4041,237.8321) .. controls (281.1249,237.8321) and (281.7093,237.2478) ..(281.7093,236.5270);
\path[fill=chfourfig5c0,fill opacity=0.5,draw=chfourfig5c0,line width=0.3985pt,draw opacity=0.5] (256.5931,244.1585) .. controls (256.5931,243.4377) and (256.0088,242.8534) ..(255.2880,242.8534) .. controls (254.5672,242.8534) and (253.9829,243.4377) ..(253.9829,244.1585) .. controls (253.9829,244.8793) and (254.5672,245.4636) ..(255.2880,245.4636) .. controls (256.0088,245.4636) and (256.5931,244.8793) ..(256.5931,244.1585);
\path[fill=chfourfig5c0,fill opacity=0.5,draw=chfourfig5c0,line width=0.3985pt,draw opacity=0.5] (254.6902,202.8230) .. controls (254.6902,202.1022) and (254.1059,201.5178) ..(253.3851,201.5178) .. controls (252.6643,201.5178) and (252.0800,202.1022) ..(252.0800,202.8230) .. controls (252.0800,203.5438) and (252.6643,204.1281) ..(253.3851,204.1281) .. controls (254.1059,204.1281) and (254.6902,203.5438) ..(254.6902,202.8230);
\path[fill=chfourfig5c0,fill opacity=0.5,draw=chfourfig5c0,line width=0.3985pt,draw opacity=0.5] (289.6097,263.3468) .. controls (289.6097,262.6260) and (289.0254,262.0417) ..(288.3046,262.0417) .. controls (287.5838,262.0417) and (286.9995,262.6260) ..(286.9995,263.3468) .. controls (286.9995,264.0676) and (287.5838,264.6519) ..(288.3046,264.6519) .. controls (289.0254,264.6519) and (289.6097,264.0676) ..(289.6097,263.3468);
\path[fill=chfourfig5c0,fill opacity=0.5,draw=chfourfig5c0,line width=0.3985pt,draw opacity=0.5] (283.2435,225.0499) .. controls (283.2435,224.3291) and (282.6592,223.7448) ..(281.9384,223.7448) .. controls (281.2176,223.7448) and (280.6333,224.3291) ..(280.6333,225.0499) .. controls (280.6333,225.7707) and (281.2176,226.3550) ..(281.9384,226.3550) .. controls (282.6592,226.3550) and (283.2435,225.7707) ..(283.2435,225.0499);
\path[fill=chfourfig5c0,fill opacity=0.5,draw=chfourfig5c0,line width=0.3985pt,draw opacity=0.5] (271.4376,252.1785) .. controls (271.4376,251.4577) and (270.8533,250.8734) ..(270.1325,250.8734) .. controls (269.4117,250.8734) and (268.8274,251.4577) ..(268.8274,252.1785) .. controls (268.8274,252.8993) and (269.4117,253.4836) ..(270.1325,253.4836) .. controls (270.8533,253.4836) and (271.4376,252.8993) ..(271.4376,252.1785);
\path[fill=chfourfig5c0,fill opacity=0.5,draw=chfourfig5c0,line width=0.3985pt,draw opacity=0.5] (277.9334,221.6227) .. controls (277.9334,220.9019) and (277.3491,220.3176) ..(276.6282,220.3176) .. controls (275.9074,220.3176) and (275.3231,220.9019) ..(275.3231,221.6227) .. controls (275.3231,222.3435) and (275.9074,222.9278) ..(276.6282,222.9278) .. controls (277.3491,222.9278) and (277.9334,222.3435) ..(277.9334,221.6227);
\path[fill=chfourfig5c0,fill opacity=0.5,draw=chfourfig5c0,line width=0.3985pt,draw opacity=0.5] (254.1522,239.9542) .. controls (254.1522,239.2334) and (253.5679,238.6490) ..(252.8471,238.6490) .. controls (252.1263,238.6490) and (251.5420,239.2334) ..(251.5420,239.9542) .. controls (251.5420,240.6750) and (252.1263,241.2593) ..(252.8471,241.2593) .. controls (253.5679,241.2593) and (254.1522,240.6750) ..(254.1522,239.9542);
\path[fill=chfourfig5c0,fill opacity=0.5,draw=chfourfig5c0,line width=0.3985pt,draw opacity=0.5] (273.7390,250.9730) .. controls (273.7390,250.2522) and (273.1547,249.6679) ..(272.4339,249.6679) .. controls (271.7131,249.6679) and (271.1288,250.2522) ..(271.1288,250.9730) .. controls (271.1288,251.6938) and (271.7131,252.2781) ..(272.4339,252.2781) .. controls (273.1547,252.2781) and (273.7390,251.6938) ..(273.7390,250.9730);
\path[fill=chfourfig5c0,fill opacity=0.5,draw=chfourfig5c0,line width=0.3985pt,draw opacity=0.5] (338.9553,212.4570) .. controls (338.9553,211.7362) and (338.3710,211.1519) ..(337.6502,211.1519) .. controls (336.9294,211.1519) and (336.3451,211.7362) ..(336.3451,212.4570) .. controls (336.3451,213.1778) and (336.9294,213.7621) ..(337.6502,213.7621) .. controls (338.3710,213.7621) and (338.9553,213.1778) ..(338.9553,212.4570);
\path[fill=chfourfig5c0,fill opacity=0.5,draw=chfourfig5c0,line width=0.3985pt,draw opacity=0.5] (304.1454,194.7133) .. controls (304.1454,193.9925) and (303.5611,193.4081) ..(302.8403,193.4081) .. controls (302.1195,193.4081) and (301.5352,193.9925) ..(301.5352,194.7133) .. controls (301.5352,195.4341) and (302.1195,196.0184) ..(302.8403,196.0184) .. controls (303.5611,196.0184) and (304.1454,195.4341) ..(304.1454,194.7133);
\path[fill=chfourfig5c0,fill opacity=0.5,draw=chfourfig5c0,line width=0.3985pt,draw opacity=0.5] (239.7261,221.6924) .. controls (239.7261,220.9716) and (239.1418,220.3873) ..(238.4210,220.3873) .. controls (237.7002,220.3873) and (237.1159,220.9716) ..(237.1159,221.6924) .. controls (237.1159,222.4133) and (237.7002,222.9976) ..(238.4210,222.9976) .. controls (239.1418,222.9976) and (239.7261,222.4133) ..(239.7261,221.6924);
\path[fill=chfourfig5c0,fill opacity=0.5,draw=chfourfig5c0,line width=0.3985pt,draw opacity=0.5] (257.5097,196.3073) .. controls (257.5097,195.5865) and (256.9254,195.0022) ..(256.2046,195.0022) .. controls (255.4838,195.0022) and (254.8994,195.5865) ..(254.8994,196.3073) .. controls (254.8994,197.0281) and (255.4838,197.6124) ..(256.2046,197.6124) .. controls (256.9254,197.6124) and (257.5097,197.0281) ..(257.5097,196.3073);
\path[fill=chfourfig5c0,fill opacity=0.5,draw=chfourfig5c0,line width=0.3985pt,draw opacity=0.5] (321.7895,203.0820) .. controls (321.7895,202.3612) and (321.2052,201.7769) ..(320.4843,201.7769) .. controls (319.7635,201.7769) and (319.1793,202.3612) ..(319.1793,203.0820) .. controls (319.1793,203.8028) and (319.7635,204.3871) ..(320.4843,204.3871) .. controls (321.2052,204.3871) and (321.7895,203.8028) ..(321.7895,203.0820);
\path[fill=chfourfig5c0,fill opacity=0.5,draw=chfourfig5c0,line width=0.3985pt,draw opacity=0.5] (273.9084,252.0988) .. controls (273.9084,251.3780) and (273.3241,250.7937) ..(272.6033,250.7937) .. controls (271.8825,250.7937) and (271.2982,251.3780) ..(271.2982,252.0988) .. controls (271.2982,252.8196) and (271.8825,253.4039) ..(272.6033,253.4039) .. controls (273.3241,253.4039) and (273.9084,252.8196) ..(273.9084,252.0988);
\path[fill=chfourfig5c0,fill opacity=0.5,draw=chfourfig5c0,line width=0.3985pt,draw opacity=0.5] (287.4877,203.3211) .. controls (287.4877,202.6003) and (286.9034,202.0160) ..(286.1825,202.0160) .. controls (285.4617,202.0160) and (284.8774,202.6003) ..(284.8774,203.3211) .. controls (284.8774,204.0419) and (285.4617,204.6262) ..(286.1825,204.6262) .. controls (286.9034,204.6262) and (287.4877,204.0419) ..(287.4877,203.3211);
\path[fill=chfourfig5c0,fill opacity=0.5,draw=chfourfig5c0,line width=0.3985pt,draw opacity=0.5] (282.2373,238.6690) .. controls (282.2373,237.9482) and (281.6530,237.3639) ..(280.9322,237.3639) .. controls (280.2113,237.3639) and (279.6270,237.9482) ..(279.6270,238.6690) .. controls (279.6270,239.3898) and (280.2113,239.9741) ..(280.9322,239.9741) .. controls (281.6530,239.9741) and (282.2373,239.3898) ..(282.2373,238.6690);
\path[fill=chfourfig5c0,fill opacity=0.5,draw=chfourfig5c0,line width=0.3985pt,draw opacity=0.5] (285.9235,219.0224) .. controls (285.9235,218.3016) and (285.3392,217.7173) ..(284.6184,217.7173) .. controls (283.8976,217.7173) and (283.3133,218.3016) ..(283.3133,219.0224) .. controls (283.3133,219.7432) and (283.8976,220.3275) ..(284.6184,220.3275) .. controls (285.3392,220.3275) and (285.9235,219.7432) ..(285.9235,219.0224);
\path[fill=chfourfig5c0,fill opacity=0.5,draw=chfourfig5c0,line width=0.3985pt,draw opacity=0.5] (286.1825,213.2042) .. controls (286.1825,212.4834) and (285.5982,211.8990) ..(284.8774,211.8990) .. controls (284.1566,211.8990) and (283.5723,212.4834) ..(283.5723,213.2042) .. controls (283.5723,213.9250) and (284.1566,214.5093) ..(284.8774,214.5093) .. controls (285.5982,214.5093) and (286.1825,213.9250) ..(286.1825,213.2042);
\path[fill=chfourfig5c0,fill opacity=0.5,draw=chfourfig5c0,line width=0.3985pt,draw opacity=0.5] (254.7600,219.9589) .. controls (254.7600,219.2381) and (254.1757,218.6538) ..(253.4548,218.6538) .. controls (252.7341,218.6538) and (252.1497,219.2381) ..(252.1497,219.9589) .. controls (252.1497,220.6797) and (252.7341,221.2640) ..(253.4548,221.2640) .. controls (254.1757,221.2640) and (254.7600,220.6797) ..(254.7600,219.9589);
\path[fill=chfourfig5c0,fill opacity=0.5,draw=chfourfig5c0,line width=0.3985pt,draw opacity=0.5] (292.5487,185.9660) .. controls (292.5487,185.2452) and (291.9644,184.6608) ..(291.2436,184.6608) .. controls (290.5228,184.6608) and (289.9385,185.2452) ..(289.9385,185.9660) .. controls (289.9385,186.6868) and (290.5228,187.2711) ..(291.2436,187.2711) .. controls (291.9644,187.2711) and (292.5487,186.6868) ..(292.5487,185.9660);
\path[fill=chfourfig5c0,fill opacity=0.5,draw=chfourfig5c0,line width=0.3985pt,draw opacity=0.5] (257.1411,221.9913) .. controls (257.1411,221.2705) and (256.5568,220.6862) ..(255.8359,220.6862) .. controls (255.1152,220.6862) and (254.5308,221.2705) ..(254.5308,221.9913) .. controls (254.5308,222.7121) and (255.1152,223.2964) ..(255.8359,223.2964) .. controls (256.5568,223.2964) and (257.1411,222.7121) ..(257.1411,221.9913);
\path[fill=chfourfig5c0,fill opacity=0.5,draw=chfourfig5c0,line width=0.3985pt,draw opacity=0.5] (306.0184,191.1566) .. controls (306.0184,190.4357) and (305.4341,189.8514) ..(304.7133,189.8514) .. controls (303.9925,189.8514) and (303.4082,190.4357) ..(303.4082,191.1566) .. controls (303.4082,191.8773) and (303.9925,192.4617) ..(304.7133,192.4617) .. controls (305.4341,192.4617) and (306.0184,191.8773) ..(306.0184,191.1566);
\path[fill=chfourfig5c0,fill opacity=0.5,draw=chfourfig5c0,line width=0.3985pt,draw opacity=0.5] (281.4403,254.6294) .. controls (281.4403,253.9086) and (280.8559,253.3242) ..(280.1351,253.3242) .. controls (279.4143,253.3242) and (278.8300,253.9086) ..(278.8300,254.6294) .. controls (278.8300,255.3502) and (279.4143,255.9345) ..(280.1351,255.9345) .. controls (280.8559,255.9345) and (281.4403,255.3502) ..(281.4403,254.6294);
\path[fill=chfourfig5c0,fill opacity=0.5,draw=chfourfig5c0,line width=0.3985pt,draw opacity=0.5] (293.7642,267.5909) .. controls (293.7642,266.8701) and (293.1799,266.2858) ..(292.4591,266.2858) .. controls (291.7383,266.2858) and (291.1540,266.8701) ..(291.1540,267.5909) .. controls (291.1540,268.3117) and (291.7383,268.8961) ..(292.4591,268.8961) .. controls (293.1799,268.8961) and (293.7642,268.3117) ..(293.7642,267.5909);
\path[fill=chfourfig5c0,fill opacity=0.5,draw=chfourfig5c0,line width=0.3985pt,draw opacity=0.5] (296.0357,221.0847) .. controls (296.0357,220.3639) and (295.4514,219.7796) ..(294.7306,219.7796) .. controls (294.0098,219.7796) and (293.4255,220.3639) ..(293.4255,221.0847) .. controls (293.4255,221.8055) and (294.0098,222.3898) ..(294.7306,222.3898) .. controls (295.4514,222.3898) and (296.0357,221.8055) ..(296.0357,221.0847);
\path[fill=chfourfig5c0,fill opacity=0.5,draw=chfourfig5c0,line width=0.3985pt,draw opacity=0.5] (261.5048,237.4734) .. controls (261.5048,236.7527) and (260.9204,236.1684) ..(260.1996,236.1684) .. controls (259.4789,236.1684) and (258.8945,236.7527) ..(258.8945,237.4734) .. controls (258.8945,238.1943) and (259.4789,238.7786) ..(260.1996,238.7786) .. controls (260.9204,238.7786) and (261.5048,238.1943) ..(261.5048,237.4734);
\path[fill=chfourfig5c0,fill opacity=0.5,draw=chfourfig5c0,line width=0.3985pt,draw opacity=0.5] (294.5712,211.5204) .. controls (294.5712,210.7997) and (293.9869,210.2153) ..(293.2661,210.2153) .. controls (292.5453,210.2153) and (291.9610,210.7997) ..(291.9610,211.5204) .. controls (291.9610,212.2413) and (292.5453,212.8256) ..(293.2661,212.8256) .. controls (293.9869,212.8256) and (294.5712,212.2413) ..(294.5712,211.5204);
\path[fill=chfourfig5c0,fill opacity=0.5,draw=chfourfig5c0,line width=0.3985pt,draw opacity=0.5] (284.0306,224.7909) .. controls (284.0306,224.0701) and (283.4463,223.4857) ..(282.7255,223.4857) .. controls (282.0047,223.4857) and (281.4203,224.0701) ..(281.4203,224.7909) .. controls (281.4203,225.5117) and (282.0047,226.0960) ..(282.7255,226.0960) .. controls (283.4463,226.0960) and (284.0306,225.5117) ..(284.0306,224.7909);
\path[fill=chfourfig5c0,fill opacity=0.5,draw=chfourfig5c0,line width=0.3985pt,draw opacity=0.5] (305.3210,231.1272) .. controls (305.3210,230.4064) and (304.7367,229.8221) ..(304.0159,229.8221) .. controls (303.2951,229.8221) and (302.7108,230.4064) ..(302.7108,231.1272) .. controls (302.7108,231.8480) and (303.2951,232.4323) ..(304.0159,232.4323) .. controls (304.7367,232.4323) and (305.3210,231.8480) ..(305.3210,231.1272);
\path[fill=chfourfig5c0,fill opacity=0.5,draw=chfourfig5c0,line width=0.3985pt,draw opacity=0.5] (289.2311,181.5524) .. controls (289.2311,180.8316) and (288.6468,180.2473) ..(287.9260,180.2473) .. controls (287.2052,180.2473) and (286.6209,180.8316) ..(286.6209,181.5524) .. controls (286.6209,182.2732) and (287.2052,182.8575) ..(287.9260,182.8575) .. controls (288.6468,182.8575) and (289.2311,182.2732) ..(289.2311,181.5524);
\path[fill=chfourfig5c0,fill opacity=0.5,draw=chfourfig5c0,line width=0.3985pt,draw opacity=0.5] (303.8565,208.5615) .. controls (303.8565,207.8407) and (303.2722,207.2564) ..(302.5514,207.2564) .. controls (301.8306,207.2564) and (301.2462,207.8407) ..(301.2462,208.5615) .. controls (301.2462,209.2823) and (301.8306,209.8666) ..(302.5514,209.8666) .. controls (303.2722,209.8666) and (303.8565,209.2823) ..(303.8565,208.5615);
\path[fill=chfourfig5c0,fill opacity=0.5,draw=chfourfig5c0,line width=0.3985pt,draw opacity=0.5] (302.2525,199.2663) .. controls (302.2525,198.5454) and (301.6682,197.9611) ..(300.9474,197.9611) .. controls (300.2266,197.9611) and (299.6422,198.5454) ..(299.6422,199.2663) .. controls (299.6422,199.9871) and (300.2266,200.5714) ..(300.9474,200.5714) .. controls (301.6682,200.5714) and (302.2525,199.9871) ..(302.2525,199.2663);
\path[fill=chfourfig5c0,fill opacity=0.5,draw=chfourfig5c0,line width=0.3985pt,draw opacity=0.5] (266.5858,195.6398) .. controls (266.5858,194.9190) and (266.0014,194.3347) ..(265.2807,194.3347) .. controls (264.5598,194.3347) and (263.9755,194.9190) ..(263.9755,195.6398) .. controls (263.9755,196.3606) and (264.5598,196.9449) ..(265.2807,196.9449) .. controls (266.0014,196.9449) and (266.5858,196.3606) ..(266.5858,195.6398);
\path[fill=chfourfig5c0,fill opacity=0.5,draw=chfourfig5c0,line width=0.3985pt,draw opacity=0.5] (244.3887,213.9514) .. controls (244.3887,213.2306) and (243.8044,212.6462) ..(243.0836,212.6462) .. controls (242.3628,212.6462) and (241.7785,213.2306) ..(241.7785,213.9514) .. controls (241.7785,214.6722) and (242.3628,215.2565) ..(243.0836,215.2565) .. controls (243.8044,215.2565) and (244.3887,214.6722) ..(244.3887,213.9514);
\path[fill=chfourfig5c0,fill opacity=0.5,draw=chfourfig5c0,line width=0.3985pt,draw opacity=0.5] (298.7755,184.8800) .. controls (298.7755,184.1592) and (298.1912,183.5749) ..(297.4704,183.5749) .. controls (296.7495,183.5749) and (296.1652,184.1592) ..(296.1652,184.8800) .. controls (296.1652,185.6008) and (296.7495,186.1851) ..(297.4704,186.1851) .. controls (298.1912,186.1851) and (298.7755,185.6008) ..(298.7755,184.8800);
\path[fill=chfourfig5c0,fill opacity=0.5,draw=chfourfig5c0,line width=0.3985pt,draw opacity=0.5] (285.4054,203.9288) .. controls (285.4054,203.2080) and (284.8211,202.6237) ..(284.1003,202.6237) .. controls (283.3795,202.6237) and (282.7952,203.2080) ..(282.7952,203.9288) .. controls (282.7952,204.6496) and (283.3795,205.2339) ..(284.1003,205.2339) .. controls (284.8211,205.2339) and (285.4054,204.6496) ..(285.4054,203.9288);
\path[fill=chfourfig5c0,fill opacity=0.5,draw=chfourfig5c0,line width=0.3985pt,draw opacity=0.5] (247.3776,219.3213) .. controls (247.3776,218.6005) and (246.7932,218.0162) ..(246.0724,218.0162) .. controls (245.3516,218.0162) and (244.7673,218.6005) ..(244.7673,219.3213) .. controls (244.7673,220.0421) and (245.3516,220.6264) ..(246.0724,220.6264) .. controls (246.7932,220.6264) and (247.3776,220.0421) ..(247.3776,219.3213);
\path[fill=chfourfig5c0,fill opacity=0.5,draw=chfourfig5c0,line width=0.3985pt,draw opacity=0.5] (247.9853,224.0436) .. controls (247.9853,223.3229) and (247.4010,222.7386) ..(246.6802,222.7386) .. controls (245.9594,222.7386) and (245.3751,223.3229) ..(245.3751,224.0436) .. controls (245.3751,224.7645) and (245.9594,225.3488) ..(246.6802,225.3488) .. controls (247.4010,225.3488) and (247.9853,224.7645) ..(247.9853,224.0436);
\path[fill=chfourfig5c0,fill opacity=0.5,draw=chfourfig5c0,line width=0.3985pt,draw opacity=0.5] (300.2599,249.2794) .. controls (300.2599,248.5585) and (299.6756,247.9742) ..(298.9548,247.9742) .. controls (298.2340,247.9742) and (297.6497,248.5585) ..(297.6497,249.2794) .. controls (297.6497,250.0001) and (298.2340,250.5844) ..(298.9548,250.5844) .. controls (299.6756,250.5844) and (300.2599,250.0001) ..(300.2599,249.2794);
\path[fill=chfourfig5c0,fill opacity=0.5,draw=chfourfig5c0,line width=0.3985pt,draw opacity=0.5] (286.2523,244.5470) .. controls (286.2523,243.8262) and (285.6680,243.2419) ..(284.9471,243.2419) .. controls (284.2264,243.2419) and (283.6421,243.8262) ..(283.6421,244.5470) .. controls (283.6421,245.2678) and (284.2264,245.8521) ..(284.9471,245.8521) .. controls (285.6680,245.8521) and (286.2523,245.2678) ..(286.2523,244.5470);
\path[fill=chfourfig5c0,fill opacity=0.5,draw=chfourfig5c0,line width=0.3985pt,draw opacity=0.5] (235.0337,222.8083) .. controls (235.0337,222.0875) and (234.4494,221.5032) ..(233.7286,221.5032) .. controls (233.0078,221.5032) and (232.4235,222.0875) ..(232.4235,222.8083) .. controls (232.4235,223.5291) and (233.0078,224.1134) ..(233.7286,224.1134) .. controls (234.4494,224.1134) and (235.0337,223.5291) ..(235.0337,222.8083);
\path[fill=chfourfig5c0,fill opacity=0.5,draw=chfourfig5c0,line width=0.3985pt,draw opacity=0.5] (292.7779,197.5526) .. controls (292.7779,196.8318) and (292.1936,196.2475) ..(291.4728,196.2475) .. controls (290.7520,196.2475) and (290.1676,196.8318) ..(290.1676,197.5526) .. controls (290.1676,198.2734) and (290.7520,198.8578) ..(291.4728,198.8578) .. controls (292.1936,198.8578) and (292.7779,198.2734) ..(292.7779,197.5526);
\path[fill=chfourfig5c0,fill opacity=0.5,draw=chfourfig5c0,line width=0.3985pt,draw opacity=0.5] (245.2256,242.4947) .. controls (245.2256,241.7739) and (244.6413,241.1896) ..(243.9205,241.1896) .. controls (243.1997,241.1896) and (242.6154,241.7739) ..(242.6154,242.4947) .. controls (242.6154,243.2155) and (243.1997,243.7998) ..(243.9205,243.7998) .. controls (244.6413,243.7998) and (245.2256,243.2155) ..(245.2256,242.4947);
\path[fill=chfourfig5c0,fill opacity=0.5,draw=chfourfig5c0,line width=0.3985pt,draw opacity=0.5] (257.4897,234.9429) .. controls (257.4897,234.2221) and (256.9055,233.6378) ..(256.1846,233.6378) .. controls (255.4638,233.6378) and (254.8795,234.2221) ..(254.8795,234.9429) .. controls (254.8795,235.6637) and (255.4638,236.2480) ..(256.1846,236.2480) .. controls (256.9055,236.2480) and (257.4897,235.6637) ..(257.4897,234.9429);
\path[fill=chfourfig5c0,fill opacity=0.5,draw=chfourfig5c0,line width=0.3985pt,draw opacity=0.5] (274.6855,234.9130) .. controls (274.6855,234.1922) and (274.1012,233.6079) ..(273.3804,233.6079) .. controls (272.6596,233.6079) and (272.0753,234.1922) ..(272.0753,234.9130) .. controls (272.0753,235.6339) and (272.6596,236.2181) ..(273.3804,236.2181) .. controls (274.1012,236.2181) and (274.6855,235.6339) ..(274.6855,234.9130);
\path[fill=chfourfig5c0,fill opacity=0.5,draw=chfourfig5c0,line width=0.3985pt,draw opacity=0.5] (263.9456,253.2844) .. controls (263.9456,252.5636) and (263.3613,251.9792) ..(262.6405,251.9792) .. controls (261.9197,251.9792) and (261.3354,252.5636) ..(261.3354,253.2844) .. controls (261.3354,254.0052) and (261.9197,254.5895) ..(262.6405,254.5895) .. controls (263.3613,254.5895) and (263.9456,254.0052) ..(263.9456,253.2844);
\path[fill=chfourfig5c0,fill opacity=0.5,draw=chfourfig5c0,line width=0.3985pt,draw opacity=0.5] (327.7273,204.1679) .. controls (327.7273,203.4471) and (327.1429,202.8628) ..(326.4222,202.8628) .. controls (325.7014,202.8628) and (325.1170,203.4471) ..(325.1170,204.1679) .. controls (325.1170,204.8887) and (325.7014,205.4731) ..(326.4222,205.4731) .. controls (327.1429,205.4731) and (327.7273,204.8887) ..(327.7273,204.1679);
\path[fill=chfourfig5c0,fill opacity=0.5,draw=chfourfig5c0,line width=0.3985pt,draw opacity=0.5] (256.6130,192.6012) .. controls (256.6130,191.8804) and (256.0287,191.2960) ..(255.3079,191.2960) .. controls (254.5871,191.2960) and (254.0028,191.8804) ..(254.0028,192.6012) .. controls (254.0028,193.3220) and (254.5871,193.9062) ..(255.3079,193.9062) .. controls (256.0287,193.9062) and (256.6130,193.3220) ..(256.6130,192.6012);
\path[fill=chfourfig5c0,fill opacity=0.5,draw=chfourfig5c0,line width=0.3985pt,draw opacity=0.5] (291.7119,208.2228) .. controls (291.7119,207.5020) and (291.1276,206.9177) ..(290.4068,206.9177) .. controls (289.6859,206.9177) and (289.1016,207.5020) ..(289.1016,208.2228) .. controls (289.1016,208.9436) and (289.6859,209.5279) ..(290.4068,209.5279) .. controls (291.1276,209.5279) and (291.7119,208.9436) ..(291.7119,208.2228);
\path[fill=chfourfig5c0,fill opacity=0.5,draw=chfourfig5c0,line width=0.3985pt,draw opacity=0.5] (289.4802,254.0415) .. controls (289.4802,253.3207) and (288.8959,252.7364) ..(288.1751,252.7364) .. controls (287.4543,252.7364) and (286.8700,253.3207) ..(286.8700,254.0415) .. controls (286.8700,254.7623) and (287.4543,255.3467) ..(288.1751,255.3467) .. controls (288.8959,255.3467) and (289.4802,254.7623) ..(289.4802,254.0415);
\path[fill=chfourfig5c0,fill opacity=0.5,draw=chfourfig5c0,line width=0.3985pt,draw opacity=0.5] (320.4246,231.3264) .. controls (320.4246,230.6057) and (319.8403,230.0213) ..(319.1194,230.0213) .. controls (318.3987,230.0213) and (317.8143,230.6057) ..(317.8143,231.3264) .. controls (317.8143,232.0472) and (318.3987,232.6316) ..(319.1194,232.6316) .. controls (319.8403,232.6316) and (320.4246,232.0472) ..(320.4246,231.3264);
\path[fill=chfourfig5c0,fill opacity=0.5,draw=chfourfig5c0,line width=0.3985pt,draw opacity=0.5] (267.6617,244.8061) .. controls (267.6617,244.0853) and (267.0774,243.5009) ..(266.3566,243.5009) .. controls (265.6358,243.5009) and (265.0515,244.0853) ..(265.0515,244.8061) .. controls (265.0515,245.5269) and (265.6358,246.1112) ..(266.3566,246.1112) .. controls (267.0774,246.1112) and (267.6617,245.5269) ..(267.6617,244.8061);
\path[fill=chfourfig5c0,fill opacity=0.5,draw=chfourfig5c0,line width=0.3985pt,draw opacity=0.5] (263.6268,225.5480) .. controls (263.6268,224.8272) and (263.0425,224.2429) ..(262.3217,224.2429) .. controls (261.6009,224.2429) and (261.0166,224.8272) ..(261.0166,225.5480) .. controls (261.0166,226.2688) and (261.6009,226.8531) ..(262.3217,226.8531) .. controls (263.0425,226.8531) and (263.6268,226.2688) ..(263.6268,225.5480);
\path[fill=chfourfig5c0,fill opacity=0.5,draw=chfourfig5c0,line width=0.3985pt,draw opacity=0.5] (260.7874,211.4706) .. controls (260.7874,210.7498) and (260.2031,210.1655) ..(259.4823,210.1655) .. controls (258.7615,210.1655) and (258.1772,210.7498) ..(258.1772,211.4706) .. controls (258.1772,212.1914) and (258.7615,212.7758) ..(259.4823,212.7758) .. controls (260.2031,212.7758) and (260.7874,212.1914) ..(260.7874,211.4706);
\path[fill=chfourfig5c0,fill opacity=0.5,draw=chfourfig5c0,line width=0.3985pt,draw opacity=0.5] (305.1815,245.7924) .. controls (305.1815,245.0716) and (304.5972,244.4872) ..(303.8764,244.4872) .. controls (303.1556,244.4872) and (302.5713,245.0716) ..(302.5713,245.7924) .. controls (302.5713,246.5132) and (303.1556,247.0975) ..(303.8764,247.0975) .. controls (304.5972,247.0975) and (305.1815,246.5132) ..(305.1815,245.7924);
\path[fill=chfourfig5c0,fill opacity=0.5,draw=chfourfig5c0,line width=0.3985pt,draw opacity=0.5] (273.7689,172.2273) .. controls (273.7689,171.5065) and (273.1846,170.9222) ..(272.4638,170.9222) .. controls (271.7430,170.9222) and (271.1587,171.5065) ..(271.1587,172.2273) .. controls (271.1587,172.9481) and (271.7430,173.5324) ..(272.4638,173.5324) .. controls (273.1846,173.5324) and (273.7689,172.9481) ..(273.7689,172.2273);
\path[fill=chfourfig5c0,fill opacity=0.5,draw=chfourfig5c0,line width=0.3985pt,draw opacity=0.5] (327.1594,227.3712) .. controls (327.1594,226.6504) and (326.5751,226.0661) ..(325.8543,226.0661) .. controls (325.1335,226.0661) and (324.5492,226.6504) ..(324.5492,227.3712) .. controls (324.5492,228.0920) and (325.1335,228.6763) ..(325.8543,228.6763) .. controls (326.5751,228.6763) and (327.1594,228.0920) ..(327.1594,227.3712);
\path[fill=chfourfig5c0,fill opacity=0.5,draw=chfourfig5c0,line width=0.3985pt,draw opacity=0.5] (281.0517,260.9955) .. controls (281.0517,260.2748) and (280.4674,259.6904) ..(279.7466,259.6904) .. controls (279.0258,259.6904) and (278.4415,260.2748) ..(278.4415,260.9955) .. controls (278.4415,261.7164) and (279.0258,262.3007) ..(279.7466,262.3007) .. controls (280.4674,262.3007) and (281.0517,261.7164) ..(281.0517,260.9955);
\path[fill=chfourfig5c0,fill opacity=0.5,draw=chfourfig5c0,line width=0.3985pt,draw opacity=0.5] (273.1911,194.5339) .. controls (273.1911,193.8131) and (272.6068,193.2288) ..(271.8860,193.2288) .. controls (271.1652,193.2288) and (270.5809,193.8131) ..(270.5809,194.5339) .. controls (270.5809,195.2547) and (271.1652,195.8391) ..(271.8860,195.8391) .. controls (272.6068,195.8391) and (273.1911,195.2547) ..(273.1911,194.5339);
\path[fill=chfourfig5c0,fill opacity=0.5,draw=chfourfig5c0,line width=0.3985pt,draw opacity=0.5] (316.0609,241.3789) .. controls (316.0609,240.6581) and (315.4766,240.0737) ..(314.7558,240.0737) .. controls (314.0349,240.0737) and (313.4507,240.6581) ..(313.4507,241.3789) .. controls (313.4507,242.0997) and (314.0349,242.6840) ..(314.7558,242.6840) .. controls (315.4766,242.6840) and (316.0609,242.0997) ..(316.0609,241.3789);
\path[fill=chfourfig5c0,fill opacity=0.5,draw=chfourfig5c0,line width=0.3985pt,draw opacity=0.5] (272.9221,217.6476) .. controls (272.9221,216.9268) and (272.3378,216.3424) ..(271.6170,216.3424) .. controls (270.8962,216.3424) and (270.3119,216.9268) ..(270.3119,217.6476) .. controls (270.3119,218.3683) and (270.8962,218.9527) ..(271.6170,218.9527) .. controls (272.3378,218.9527) and (272.9221,218.3683) ..(272.9221,217.6476);
\path[fill=chfourfig5c0,fill opacity=0.5,draw=chfourfig5c0,line width=0.3985pt,draw opacity=0.5] (261.9830,243.9792) .. controls (261.9830,243.2584) and (261.3987,242.6740) ..(260.6779,242.6740) .. controls (259.9571,242.6740) and (259.3727,243.2584) ..(259.3727,243.9792) .. controls (259.3727,244.7000) and (259.9571,245.2843) ..(260.6779,245.2843) .. controls (261.3987,245.2843) and (261.9830,244.7000) ..(261.9830,243.9792);
\path[fill=chfourfig5c0,fill opacity=0.5,draw=chfourfig5c0,line width=0.3985pt,draw opacity=0.5] (309.0272,244.3777) .. controls (309.0272,243.6569) and (308.4428,243.0725) ..(307.7220,243.0725) .. controls (307.0013,243.0725) and (306.4169,243.6569) ..(306.4169,244.3777) .. controls (306.4169,245.0984) and (307.0013,245.6828) ..(307.7220,245.6828) .. controls (308.4428,245.6828) and (309.0272,245.0984) ..(309.0272,244.3777);
\path[fill=chfourfig5c0,fill opacity=0.5,draw=chfourfig5c0,line width=0.3985pt,draw opacity=0.5] (289.0518,231.9242) .. controls (289.0518,231.2034) and (288.4675,230.6191) ..(287.7467,230.6191) .. controls (287.0259,230.6191) and (286.4416,231.2034) ..(286.4416,231.9242) .. controls (286.4416,232.6450) and (287.0259,233.2293) ..(287.7467,233.2293) .. controls (288.4675,233.2293) and (289.0518,232.6450) ..(289.0518,231.9242);
\path[fill=chfourfig5c0,fill opacity=0.5,draw=chfourfig5c0,line width=0.3985pt,draw opacity=0.5] (271.4376,238.9978) .. controls (271.4376,238.2770) and (270.8533,237.6927) ..(270.1325,237.6927) .. controls (269.4117,237.6927) and (268.8274,238.2770) ..(268.8274,238.9978) .. controls (268.8274,239.7186) and (269.4117,240.3029) ..(270.1325,240.3029) .. controls (270.8533,240.3029) and (271.4376,239.7186) ..(271.4376,238.9978);
\path[fill=chfourfig5c0,fill opacity=0.5,draw=chfourfig5c0,line width=0.3985pt,draw opacity=0.5] (293.1067,256.6020) .. controls (293.1067,255.8812) and (292.5223,255.2968) ..(291.8015,255.2968) .. controls (291.0807,255.2968) and (290.4964,255.8812) ..(290.4964,256.6020) .. controls (290.4964,257.3228) and (291.0807,257.9071) ..(291.8015,257.9071) .. controls (292.5223,257.9071) and (293.1067,257.3228) ..(293.1067,256.6020);
\path[fill=chfourfig5c0,fill opacity=0.5,draw=chfourfig5c0,line width=0.3985pt,draw opacity=0.5] (314.5366,236.4473) .. controls (314.5366,235.7265) and (313.9523,235.1422) ..(313.2314,235.1422) .. controls (312.5107,235.1422) and (311.9264,235.7265) ..(311.9264,236.4473) .. controls (311.9264,237.1681) and (312.5107,237.7524) ..(313.2314,237.7524) .. controls (313.9523,237.7524) and (314.5366,237.1681) ..(314.5366,236.4473);
\path[fill=chfourfig5c0,fill opacity=0.5,draw=chfourfig5c0,line width=0.3985pt,draw opacity=0.5] (286.1028,233.1696) .. controls (286.1028,232.4487) and (285.5185,231.8644) ..(284.7977,231.8644) .. controls (284.0769,231.8644) and (283.4926,232.4487) ..(283.4926,233.1696) .. controls (283.4926,233.8904) and (284.0769,234.4747) ..(284.7977,234.4747) .. controls (285.5185,234.4747) and (286.1028,233.8904) ..(286.1028,233.1696);
\path[fill=chfourfig5c0,fill opacity=0.5,draw=chfourfig5c0,line width=0.3985pt,draw opacity=0.5] (295.2487,185.3981) .. controls (295.2487,184.6773) and (294.6644,184.0930) ..(293.9435,184.0930) .. controls (293.2227,184.0930) and (292.6384,184.6773) ..(292.6384,185.3981) .. controls (292.6384,186.1189) and (293.2227,186.7032) ..(293.9435,186.7032) .. controls (294.6644,186.7032) and (295.2487,186.1189) ..(295.2487,185.3981);
\path[fill=chfourfig5c0,fill opacity=0.5,draw=chfourfig5c0,line width=0.3985pt,draw opacity=0.5] (341.9840,208.0135) .. controls (341.9840,207.2928) and (341.3997,206.7084) ..(340.6789,206.7084) .. controls (339.9581,206.7084) and (339.3738,207.2928) ..(339.3738,208.0135) .. controls (339.3738,208.7343) and (339.9581,209.3187) ..(340.6789,209.3187) .. controls (341.3997,209.3187) and (341.9840,208.7343) ..(341.9840,208.0135);
\path[fill=chfourfig5c0,fill opacity=0.5,draw=chfourfig5c0,line width=0.3985pt,draw opacity=0.5] (273.8088,218.4047) .. controls (273.8088,217.6839) and (273.2245,217.0996) ..(272.5037,217.0996) .. controls (271.7829,217.0996) and (271.1985,217.6839) ..(271.1985,218.4047) .. controls (271.1985,219.1255) and (271.7829,219.7098) ..(272.5037,219.7098) .. controls (273.2245,219.7098) and (273.8088,219.1255) ..(273.8088,218.4047);
\path[fill=chfourfig5c0,fill opacity=0.5,draw=chfourfig5c0,line width=0.3985pt,draw opacity=0.5] (246.8994,208.5615) .. controls (246.8994,207.8407) and (246.3150,207.2564) ..(245.5942,207.2564) .. controls (244.8734,207.2564) and (244.2891,207.8407) ..(244.2891,208.5615) .. controls (244.2891,209.2823) and (244.8734,209.8666) ..(245.5942,209.8666) .. controls (246.3150,209.8666) and (246.8994,209.2823) ..(246.8994,208.5615);
\path[fill=chfourfig5c0,fill opacity=0.5,draw=chfourfig5c0,line width=0.3985pt,draw opacity=0.5] (259.8908,243.7201) .. controls (259.8908,242.9993) and (259.3065,242.4150) ..(258.5857,242.4150) .. controls (257.8649,242.4150) and (257.2805,242.9993) ..(257.2805,243.7201) .. controls (257.2805,244.4409) and (257.8649,245.0253) ..(258.5857,245.0253) .. controls (259.3065,245.0253) and (259.8908,244.4409) ..(259.8908,243.7201);
\path[fill=chfourfig5c0,fill opacity=0.5,draw=chfourfig5c0,line width=0.3985pt,draw opacity=0.5] (295.9162,186.5836) .. controls (295.9162,185.8628) and (295.3318,185.2785) ..(294.6111,185.2785) .. controls (293.8903,185.2785) and (293.3059,185.8628) ..(293.3059,186.5836) .. controls (293.3059,187.3044) and (293.8903,187.8888) ..(294.6111,187.8888) .. controls (295.3318,187.8888) and (295.9162,187.3044) ..(295.9162,186.5836);
\path[fill=chfourfig5c0,fill opacity=0.5,draw=chfourfig5c0,line width=0.3985pt,draw opacity=0.5] (276.7578,254.7589) .. controls (276.7578,254.0381) and (276.1734,253.4538) ..(275.4526,253.4538) .. controls (274.7318,253.4538) and (274.1475,254.0381) ..(274.1475,254.7589) .. controls (274.1475,255.4797) and (274.7318,256.0640) ..(275.4526,256.0640) .. controls (276.1734,256.0640) and (276.7578,255.4797) ..(276.7578,254.7589);
\path[fill=chfourfig5c0,fill opacity=0.5,draw=chfourfig5c0,line width=0.3985pt,draw opacity=0.5] (253.5046,193.4579) .. controls (253.5046,192.7372) and (252.9203,192.1528) ..(252.1995,192.1528) .. controls (251.4787,192.1528) and (250.8944,192.7372) ..(250.8944,193.4579) .. controls (250.8944,194.1788) and (251.4787,194.7631) ..(252.1995,194.7631) .. controls (252.9203,194.7631) and (253.5046,194.1788) ..(253.5046,193.4579);
\path[fill=chfourfig5c0,fill opacity=0.5,draw=chfourfig5c0,line width=0.3985pt,draw opacity=0.5] (274.4464,225.2591) .. controls (274.4464,224.5383) and (273.8621,223.9540) ..(273.1413,223.9540) .. controls (272.4205,223.9540) and (271.8362,224.5383) ..(271.8362,225.2591) .. controls (271.8362,225.9799) and (272.4205,226.5642) ..(273.1413,226.5642) .. controls (273.8621,226.5642) and (274.4464,225.9799) ..(274.4464,225.2591);
\path[fill=chfourfig5c0,fill opacity=0.5,draw=chfourfig5c0,line width=0.3985pt,draw opacity=0.5] (272.7926,194.0557) .. controls (272.7926,193.3349) and (272.2083,192.7506) ..(271.4875,192.7506) .. controls (270.7666,192.7506) and (270.1823,193.3349) ..(270.1823,194.0557) .. controls (270.1823,194.7765) and (270.7666,195.3608) ..(271.4875,195.3608) .. controls (272.2083,195.3608) and (272.7926,194.7765) ..(272.7926,194.0557);
\path[fill=chfourfig5c0,fill opacity=0.5,draw=chfourfig5c0,line width=0.3985pt,draw opacity=0.5] (265.2508,248.4823) .. controls (265.2508,247.7615) and (264.6664,247.1772) ..(263.9456,247.1772) .. controls (263.2248,247.1772) and (262.6405,247.7615) ..(262.6405,248.4823) .. controls (262.6405,249.2032) and (263.2248,249.7874) ..(263.9456,249.7874) .. controls (264.6664,249.7874) and (265.2508,249.2032) ..(265.2508,248.4823);
\path[fill=chfourfig5c0,fill opacity=0.5,draw=chfourfig5c0,line width=0.3985pt,draw opacity=0.5] (266.3865,241.8073) .. controls (266.3865,241.0865) and (265.8022,240.5021) ..(265.0814,240.5021) .. controls (264.3606,240.5021) and (263.7763,241.0865) ..(263.7763,241.8073) .. controls (263.7763,242.5281) and (264.3606,243.1124) ..(265.0814,243.1124) .. controls (265.8022,243.1124) and (266.3865,242.5281) ..(266.3865,241.8073);
\path[fill=chfourfig5c0,fill opacity=0.5,draw=chfourfig5c0,line width=0.3985pt,draw opacity=0.5] (264.9718,193.6672) .. controls (264.9718,192.9464) and (264.3875,192.3621) ..(263.6667,192.3621) .. controls (262.9459,192.3621) and (262.3616,192.9464) ..(262.3616,193.6672) .. controls (262.3616,194.3880) and (262.9459,194.9723) ..(263.6667,194.9723) .. controls (264.3875,194.9723) and (264.9718,194.3880) ..(264.9718,193.6672);
\path[fill=chfourfig5c0,fill opacity=0.5,draw=chfourfig5c0,line width=0.3985pt,draw opacity=0.5] (308.4593,251.8796) .. controls (308.4593,251.1588) and (307.8750,250.5745) ..(307.1542,250.5745) .. controls (306.4334,250.5745) and (305.8491,251.1588) ..(305.8491,251.8796) .. controls (305.8491,252.6004) and (306.4334,253.1847) ..(307.1542,253.1847) .. controls (307.8750,253.1847) and (308.4593,252.6004) ..(308.4593,251.8796);
\path[fill=chfourfig5c0,fill opacity=0.5,draw=chfourfig5c0,line width=0.3985pt,draw opacity=0.5] (272.0852,171.2609) .. controls (272.0852,170.5401) and (271.5009,169.9558) ..(270.7801,169.9558) .. controls (270.0593,169.9558) and (269.4750,170.5401) ..(269.4750,171.2609) .. controls (269.4750,171.9817) and (270.0593,172.5660) ..(270.7801,172.5660) .. controls (271.5009,172.5660) and (272.0852,171.9817) ..(272.0852,171.2609);
\path[fill=chfourfig5c0,fill opacity=0.5,draw=chfourfig5c0,line width=0.3985pt,draw opacity=0.5] (308.6685,202.9425) .. controls (308.6685,202.2217) and (308.0842,201.6374) ..(307.3634,201.6374) .. controls (306.6426,201.6374) and (306.0583,202.2217) ..(306.0583,202.9425) .. controls (306.0583,203.6633) and (306.6426,204.2476) ..(307.3634,204.2476) .. controls (308.0842,204.2476) and (308.6685,203.6633) ..(308.6685,202.9425);
\path[fill=chfourfig5c0,fill opacity=0.5,draw=chfourfig5c0,line width=0.3985pt,draw opacity=0.5] (302.3222,236.2680) .. controls (302.3222,235.5472) and (301.7379,234.9628) ..(301.0171,234.9628) .. controls (300.2963,234.9628) and (299.7120,235.5472) ..(299.7120,236.2680) .. controls (299.7120,236.9888) and (300.2963,237.5731) ..(301.0171,237.5731) .. controls (301.7379,237.5731) and (302.3222,236.9888) ..(302.3222,236.2680);
\path[fill=chfourfig5c0,fill opacity=0.5,draw=chfourfig5c0,line width=0.3985pt,draw opacity=0.5] (251.7014,219.2516) .. controls (251.7014,218.5308) and (251.1171,217.9464) ..(250.3963,217.9464) .. controls (249.6755,217.9464) and (249.0912,218.5308) ..(249.0912,219.2516) .. controls (249.0912,219.9724) and (249.6755,220.5567) ..(250.3963,220.5567) .. controls (251.1171,220.5567) and (251.7014,219.9724) ..(251.7014,219.2516);
\path[fill=chfourfig5c0,fill opacity=0.5,draw=chfourfig5c0,line width=0.3985pt,draw opacity=0.5] (303.5576,229.8420) .. controls (303.5576,229.1212) and (302.9733,228.5369) ..(302.2525,228.5369) .. controls (301.5317,228.5369) and (300.9474,229.1212) ..(300.9474,229.8420) .. controls (300.9474,230.5628) and (301.5317,231.1471) ..(302.2525,231.1471) .. controls (302.9733,231.1471) and (303.5576,230.5628) ..(303.5576,229.8420);
\path[fill=chfourfig5c0,fill opacity=0.5,draw=chfourfig5c0,line width=0.3985pt,draw opacity=0.5] (232.4234,222.7684) .. controls (232.4234,222.0476) and (231.8391,221.4633) ..(231.1183,221.4633) .. controls (230.3975,221.4633) and (229.8132,222.0476) ..(229.8132,222.7684) .. controls (229.8132,223.4892) and (230.3975,224.0735) ..(231.1183,224.0735) .. controls (231.8391,224.0735) and (232.4234,223.4892) ..(232.4234,222.7684);
\path[fill=chfourfig5c0,fill opacity=0.5,draw=chfourfig5c0,line width=0.3985pt,draw opacity=0.5] (230.8792,215.8742) .. controls (230.8792,215.1534) and (230.2949,214.5691) ..(229.5741,214.5691) .. controls (228.8533,214.5691) and (228.2690,215.1534) ..(228.2690,215.8742) .. controls (228.2690,216.5950) and (228.8533,217.1793) ..(229.5741,217.1793) .. controls (230.2949,217.1793) and (230.8792,216.5950) ..(230.8792,215.8742);
\path[fill=chfourfig5c0,fill opacity=0.5,draw=chfourfig5c0,line width=0.3985pt,draw opacity=0.5] (237.8133,221.7123) .. controls (237.8133,220.9916) and (237.2290,220.4072) ..(236.5082,220.4072) .. controls (235.7874,220.4072) and (235.2031,220.9916) ..(235.2031,221.7123) .. controls (235.2031,222.4332) and (235.7874,223.0175) ..(236.5082,223.0175) .. controls (237.2290,223.0175) and (237.8133,222.4332) ..(237.8133,221.7123);
\path[fill=chfourfig5c0,fill opacity=0.5,draw=chfourfig5c0,line width=0.3985pt,draw opacity=0.5] (305.7395,240.7114) .. controls (305.7395,239.9906) and (305.1552,239.4062) ..(304.4343,239.4062) .. controls (303.7135,239.4062) and (303.1292,239.9906) ..(303.1292,240.7114) .. controls (303.1292,241.4322) and (303.7135,242.0165) ..(304.4343,242.0165) .. controls (305.1552,242.0165) and (305.7395,241.4322) ..(305.7395,240.7114);
\path[fill=chfourfig5c0,fill opacity=0.5,draw=chfourfig5c0,line width=0.3985pt,draw opacity=0.5] (253.1261,235.1023) .. controls (253.1261,234.3815) and (252.5417,233.7972) ..(251.8210,233.7972) .. controls (251.1002,233.7972) and (250.5158,234.3815) ..(250.5158,235.1023) .. controls (250.5158,235.8231) and (251.1002,236.4074) ..(251.8210,236.4074) .. controls (252.5417,236.4074) and (253.1261,235.8231) ..(253.1261,235.1023);
\path[fill=chfourfig5c0,fill opacity=0.5,draw=chfourfig5c0,line width=0.3985pt,draw opacity=0.5] (317.9538,243.5807) .. controls (317.9538,242.8598) and (317.3695,242.2755) ..(316.6487,242.2755) .. controls (315.9279,242.2755) and (315.3436,242.8598) ..(315.3436,243.5807) .. controls (315.3436,244.3014) and (315.9279,244.8857) ..(316.6487,244.8857) .. controls (317.3695,244.8857) and (317.9538,244.3014) ..(317.9538,243.5807);
\path[fill=chfourfig5c0,fill opacity=0.5,draw=chfourfig5c0,line width=0.3985pt,draw opacity=0.5] (297.6895,217.3786) .. controls (297.6895,216.6578) and (297.1052,216.0735) ..(296.3844,216.0735) .. controls (295.6636,216.0735) and (295.0793,216.6578) ..(295.0793,217.3786) .. controls (295.0793,218.0994) and (295.6636,218.6837) ..(296.3844,218.6837) .. controls (297.1052,218.6837) and (297.6895,218.0994) ..(297.6895,217.3786);
\path[fill=chfourfig5c0,fill opacity=0.5,draw=chfourfig5c0,line width=0.3985pt,draw opacity=0.5] (272.2347,251.7501) .. controls (272.2347,251.0293) and (271.6504,250.4450) ..(270.9295,250.4450) .. controls (270.2087,250.4450) and (269.6244,251.0293) ..(269.6244,251.7501) .. controls (269.6244,252.4709) and (270.2087,253.0552) ..(270.9295,253.0552) .. controls (271.6504,253.0552) and (272.2347,252.4709) ..(272.2347,251.7501);
\path[fill=chfourfig5c0,fill opacity=0.5,draw=chfourfig5c0,line width=0.3985pt,draw opacity=0.5] (292.3096,187.2511) .. controls (292.3096,186.5303) and (291.7253,185.9460) ..(291.0045,185.9460) .. controls (290.2837,185.9460) and (289.6994,186.5303) ..(289.6994,187.2511) .. controls (289.6994,187.9720) and (290.2837,188.5563) ..(291.0045,188.5563) .. controls (291.7253,188.5563) and (292.3096,187.9720) ..(292.3096,187.2511);
\path[fill=chfourfig5c0,fill opacity=0.5,draw=chfourfig5c0,line width=0.3985pt,draw opacity=0.5] (320.2154,232.3625) .. controls (320.2154,231.6418) and (319.6310,231.0574) ..(318.9102,231.0574) .. controls (318.1894,231.0574) and (317.6051,231.6418) ..(317.6051,232.3625) .. controls (317.6051,233.0833) and (318.1894,233.6677) ..(318.9102,233.6677) .. controls (319.6310,233.6677) and (320.2154,233.0833) ..(320.2154,232.3625);
\path[fill=chfourfig5c0,fill opacity=0.5,draw=chfourfig5c0,line width=0.3985pt,draw opacity=0.5] (244.2493,224.2329) .. controls (244.2493,223.5121) and (243.6649,222.9278) ..(242.9441,222.9278) .. controls (242.2233,222.9278) and (241.6390,223.5121) ..(241.6390,224.2329) .. controls (241.6390,224.9538) and (242.2233,225.5381) ..(242.9441,225.5381) .. controls (243.6649,225.5381) and (244.2493,224.9538) ..(244.2493,224.2329);
\path[fill=chfourfig5c0,fill opacity=0.5,draw=chfourfig5c0,line width=0.3985pt,draw opacity=0.5] (265.9282,215.8443) .. controls (265.9282,215.1235) and (265.3439,214.5392) ..(264.6231,214.5392) .. controls (263.9023,214.5392) and (263.3180,215.1235) ..(263.3180,215.8443) .. controls (263.3180,216.5651) and (263.9023,217.1494) ..(264.6231,217.1494) .. controls (265.3439,217.1494) and (265.9282,216.5651) ..(265.9282,215.8443);
\path[fill=chfourfig5c0,fill opacity=0.5,draw=chfourfig5c0,line width=0.3985pt,draw opacity=0.5] (262.3615,224.7411) .. controls (262.3615,224.0203) and (261.7772,223.4359) ..(261.0565,223.4359) .. controls (260.3356,223.4359) and (259.7513,224.0203) ..(259.7513,224.7411) .. controls (259.7513,225.4619) and (260.3356,226.0462) ..(261.0565,226.0462) .. controls (261.7772,226.0462) and (262.3615,225.4619) ..(262.3615,224.7411);
\path[fill=chfourfig5c0,fill opacity=0.5,draw=chfourfig5c0,line width=0.3985pt,draw opacity=0.5] (304.8627,256.4824) .. controls (304.8627,255.7616) and (304.2784,255.1773) ..(303.5576,255.1773) .. controls (302.8368,255.1773) and (302.2525,255.7616) ..(302.2525,256.4824) .. controls (302.2525,257.2032) and (302.8368,257.7876) ..(303.5576,257.7876) .. controls (304.2784,257.7876) and (304.8627,257.2032) ..(304.8627,256.4824);
\path[fill=chfourfig5c0,fill opacity=0.5,draw=chfourfig5c0,line width=0.3985pt,draw opacity=0.5] (252.5482,204.1181) .. controls (252.5482,203.3973) and (251.9639,202.8130) ..(251.2431,202.8130) .. controls (250.5223,202.8130) and (249.9380,203.3973) ..(249.9380,204.1181) .. controls (249.9380,204.8389) and (250.5223,205.4232) ..(251.2431,205.4232) .. controls (251.9639,205.4232) and (252.5482,204.8389) ..(252.5482,204.1181);
\path[fill=chfourfig5c0,fill opacity=0.5,draw=chfourfig5c0,line width=0.3985pt,draw opacity=0.5] (274.5759,205.6324) .. controls (274.5759,204.9117) and (273.9916,204.3273) ..(273.2708,204.3273) .. controls (272.5500,204.3273) and (271.9657,204.9117) ..(271.9657,205.6324) .. controls (271.9657,206.3533) and (272.5500,206.9376) ..(273.2708,206.9376) .. controls (273.9916,206.9376) and (274.5759,206.3533) ..(274.5759,205.6324);
\path[fill=chfourfig5c0,fill opacity=0.5,draw=chfourfig5c0,line width=0.3985pt,draw opacity=0.5] (297.3608,265.9371) .. controls (297.3608,265.2163) and (296.7765,264.6320) ..(296.0556,264.6320) .. controls (295.3348,264.6320) and (294.7505,265.2163) ..(294.7505,265.9371) .. controls (294.7505,266.6579) and (295.3348,267.2422) ..(296.0556,267.2422) .. controls (296.7765,267.2422) and (297.3608,266.6579) ..(297.3608,265.9371);
\path[fill=chfourfig5c0,fill opacity=0.5,draw=chfourfig5c0,line width=0.3985pt,draw opacity=0.5] (256.5433,205.2240) .. controls (256.5433,204.5032) and (255.9590,203.9189) ..(255.2382,203.9189) .. controls (254.5174,203.9189) and (253.9330,204.5032) ..(253.9330,205.2240) .. controls (253.9330,205.9448) and (254.5174,206.5291) ..(255.2382,206.5291) .. controls (255.9590,206.5291) and (256.5433,205.9448) ..(256.5433,205.2240);
\path[fill=chfourfig5c0,fill opacity=0.5,draw=chfourfig5c0,line width=0.3985pt,draw opacity=0.5] (253.5345,229.8719) .. controls (253.5345,229.1511) and (252.9502,228.5668) ..(252.2294,228.5668) .. controls (251.5086,228.5668) and (250.9243,229.1511) ..(250.9243,229.8719) .. controls (250.9243,230.5927) and (251.5086,231.1770) ..(252.2294,231.1770) .. controls (252.9502,231.1770) and (253.5345,230.5927) ..(253.5345,229.8719);
\path[fill=chfourfig5c0,fill opacity=0.5,draw=chfourfig5c0,line width=0.3985pt,draw opacity=0.5] (290.7754,233.7175) .. controls (290.7754,232.9967) and (290.1911,232.4124) ..(289.4702,232.4124) .. controls (288.7495,232.4124) and (288.1651,232.9967) ..(288.1651,233.7175) .. controls (288.1651,234.4383) and (288.7495,235.0226) ..(289.4702,235.0226) .. controls (290.1911,235.0226) and (290.7754,234.4383) ..(290.7754,233.7175);
\path[fill=chfourfig5c0,fill opacity=0.5,draw=chfourfig5c0,line width=0.3985pt,draw opacity=0.5] (318.0335,205.5328) .. controls (318.0335,204.8120) and (317.4492,204.2277) ..(316.7284,204.2277) .. controls (316.0076,204.2277) and (315.4233,204.8120) ..(315.4233,205.5328) .. controls (315.4233,206.2536) and (316.0076,206.8380) ..(316.7284,206.8380) .. controls (317.4492,206.8380) and (318.0335,206.2536) ..(318.0335,205.5328);
\path[fill=chfourfig5c0,fill opacity=0.5,draw=chfourfig5c0,line width=0.3985pt,draw opacity=0.5] (309.4556,194.8129) .. controls (309.4556,194.0921) and (308.8712,193.5078) ..(308.1505,193.5078) .. controls (307.4296,193.5078) and (306.8453,194.0921) ..(306.8453,194.8129) .. controls (306.8453,195.5337) and (307.4296,196.1180) ..(308.1505,196.1180) .. controls (308.8712,196.1180) and (309.4556,195.5337) ..(309.4556,194.8129);
\path[fill=chfourfig5c0,fill opacity=0.5,draw=chfourfig5c0,line width=0.3985pt,draw opacity=0.5] (339.6328,214.7982) .. controls (339.6328,214.0774) and (339.0485,213.4931) ..(338.3277,213.4931) .. controls (337.6068,213.4931) and (337.0226,214.0774) ..(337.0226,214.7982) .. controls (337.0226,215.5190) and (337.6068,216.1033) ..(338.3277,216.1033) .. controls (339.0485,216.1033) and (339.6328,215.5190) ..(339.6328,214.7982);
\path[fill=chfourfig5c0,fill opacity=0.5,draw=chfourfig5c0,line width=0.3985pt,draw opacity=0.5] (270.9395,197.6921) .. controls (270.9395,196.9713) and (270.3552,196.3870) ..(269.6344,196.3870) .. controls (268.9136,196.3870) and (268.3293,196.9713) ..(268.3293,197.6921) .. controls (268.3293,198.4129) and (268.9136,198.9973) ..(269.6344,198.9973) .. controls (270.3552,198.9973) and (270.9395,198.4129) ..(270.9395,197.6921);
\path[fill=chfourfig5c0,fill opacity=0.5,draw=chfourfig5c0,line width=0.3985pt,draw opacity=0.5] (322.4271,238.3402) .. controls (322.4271,237.6194) and (321.8428,237.0351) ..(321.1220,237.0351) .. controls (320.4012,237.0351) and (319.8168,237.6194) ..(319.8168,238.3402) .. controls (319.8168,239.0610) and (320.4012,239.6454) ..(321.1220,239.6454) .. controls (321.8428,239.6454) and (322.4271,239.0610) ..(322.4271,238.3402);
\path[fill=chfourfig5c0,fill opacity=0.5,draw=chfourfig5c0,line width=0.3985pt,draw opacity=0.5] (283.8114,245.0950) .. controls (283.8114,244.3742) and (283.2271,243.7899) ..(282.5063,243.7899) .. controls (281.7855,243.7899) and (281.2012,244.3742) ..(281.2012,245.0950) .. controls (281.2012,245.8158) and (281.7855,246.4001) ..(282.5063,246.4001) .. controls (283.2271,246.4001) and (283.8114,245.8158) ..(283.8114,245.0950);
\path[fill=chfourfig5c0,fill opacity=0.5,draw=chfourfig5c0,line width=0.3985pt,draw opacity=0.5] (307.6125,200.1828) .. controls (307.6125,199.4620) and (307.0281,198.8777) ..(306.3073,198.8777) .. controls (305.5865,198.8777) and (305.0022,199.4620) ..(305.0022,200.1828) .. controls (305.0022,200.9036) and (305.5865,201.4879) ..(306.3073,201.4879) .. controls (307.0281,201.4879) and (307.6125,200.9036) ..(307.6125,200.1828);
\path[fill=chfourfig5c0,fill opacity=0.5,draw=chfourfig5c0,line width=0.3985pt,draw opacity=0.5] (277.6544,187.8091) .. controls (277.6544,187.0883) and (277.0701,186.5039) ..(276.3493,186.5039) .. controls (275.6285,186.5039) and (275.0442,187.0883) ..(275.0442,187.8091) .. controls (275.0442,188.5299) and (275.6285,189.1142) ..(276.3493,189.1142) .. controls (277.0701,189.1142) and (277.6544,188.5299) ..(277.6544,187.8091);
\path[fill=chfourfig5c0,fill opacity=0.5,draw=chfourfig5c0,line width=0.3985pt,draw opacity=0.5] (263.3678,200.2426) .. controls (263.3678,199.5218) and (262.7835,198.9375) ..(262.0627,198.9375) .. controls (261.3419,198.9375) and (260.7576,199.5218) ..(260.7576,200.2426) .. controls (260.7576,200.9634) and (261.3419,201.5477) ..(262.0627,201.5477) .. controls (262.7835,201.5477) and (263.3678,200.9634) ..(263.3678,200.2426);
\path[fill=chfourfig5c0,fill opacity=0.5,draw=chfourfig5c0,line width=0.3985pt,draw opacity=0.5] (294.5214,226.2355) .. controls (294.5214,225.5146) and (293.9370,224.9304) ..(293.2162,224.9304) .. controls (292.4955,224.9304) and (291.9111,225.5146) ..(291.9111,226.2355) .. controls (291.9111,226.9563) and (292.4955,227.5406) ..(293.2162,227.5406) .. controls (293.9370,227.5406) and (294.5214,226.9563) ..(294.5214,226.2355);
\path[fill=chfourfig5c0,fill opacity=0.5,draw=chfourfig5c0,line width=0.3985pt,draw opacity=0.5] (300.5987,221.0050) .. controls (300.5987,220.2842) and (300.0143,219.6999) ..(299.2935,219.6999) .. controls (298.5727,219.6999) and (297.9884,220.2842) ..(297.9884,221.0050) .. controls (297.9884,221.7258) and (298.5727,222.3101) ..(299.2935,222.3101) .. controls (300.0143,222.3101) and (300.5987,221.7258) ..(300.5987,221.0050);
\path[fill=chfourfig5c0,fill opacity=0.5,draw=chfourfig5c0,line width=0.3985pt,draw opacity=0.5] (317.1269,221.2043) .. controls (317.1269,220.4835) and (316.5426,219.8991) ..(315.8218,219.8991) .. controls (315.1010,219.8991) and (314.5167,220.4835) ..(314.5167,221.2043) .. controls (314.5167,221.9250) and (315.1010,222.5094) ..(315.8218,222.5094) .. controls (316.5426,222.5094) and (317.1269,221.9250) ..(317.1269,221.2043);
\path[fill=chfourfig5c0,fill opacity=0.5,draw=chfourfig5c0,line width=0.3985pt,draw opacity=0.5] (268.2794,193.1790) .. controls (268.2794,192.4582) and (267.6951,191.8739) ..(266.9743,191.8739) .. controls (266.2535,191.8739) and (265.6692,192.4582) ..(265.6692,193.1790) .. controls (265.6692,193.8998) and (266.2535,194.4841) ..(266.9743,194.4841) .. controls (267.6951,194.4841) and (268.2794,193.8998) ..(268.2794,193.1790);
\path[fill=chfourfig5c0,fill opacity=0.5,draw=chfourfig5c0,line width=0.3985pt,draw opacity=0.5] (284.0406,190.1503) .. controls (284.0406,189.4295) and (283.4562,188.8452) ..(282.7354,188.8452) .. controls (282.0146,188.8452) and (281.4303,189.4295) ..(281.4303,190.1503) .. controls (281.4303,190.8711) and (282.0146,191.4554) ..(282.7354,191.4554) .. controls (283.4562,191.4554) and (284.0406,190.8711) ..(284.0406,190.1503);
\path[fill=chfourfig5c0,fill opacity=0.5,draw=chfourfig5c0,line width=0.3985pt,draw opacity=0.5] (248.4834,237.5731) .. controls (248.4834,236.8523) and (247.8991,236.2680) ..(247.1783,236.2680) .. controls (246.4575,236.2680) and (245.8732,236.8523) ..(245.8732,237.5731) .. controls (245.8732,238.2939) and (246.4575,238.8782) ..(247.1783,238.8782) .. controls (247.8991,238.8782) and (248.4834,238.2939) ..(248.4834,237.5731);
\path[fill=chfourfig5c0,fill opacity=0.5,draw=chfourfig5c0,line width=0.3985pt,draw opacity=0.5] (273.7590,186.7231) .. controls (273.7590,186.0023) and (273.1747,185.4180) ..(272.4539,185.4180) .. controls (271.7330,185.4180) and (271.1487,186.0023) ..(271.1487,186.7231) .. controls (271.1487,187.4439) and (271.7330,188.0282) ..(272.4539,188.0282) .. controls (273.1747,188.0282) and (273.7590,187.4439) ..(273.7590,186.7231);
\path[fill=chfourfig5c0,fill opacity=0.5,draw=chfourfig5c0,line width=0.3985pt,draw opacity=0.5] (276.2297,173.5424) .. controls (276.2297,172.8216) and (275.6454,172.2373) ..(274.9246,172.2373) .. controls (274.2038,172.2373) and (273.6195,172.8216) ..(273.6195,173.5424) .. controls (273.6195,174.2632) and (274.2038,174.8475) ..(274.9246,174.8475) .. controls (275.6454,174.8475) and (276.2297,174.2632) ..(276.2297,173.5424);
\path[fill=chfourfig5c0,fill opacity=0.5,draw=chfourfig5c0,line width=0.3985pt,draw opacity=0.5] (282.3568,230.1110) .. controls (282.3568,229.3902) and (281.7725,228.8058) ..(281.0517,228.8058) .. controls (280.3309,228.8058) and (279.7466,229.3902) ..(279.7466,230.1110) .. controls (279.7466,230.8318) and (280.3309,231.4161) ..(281.0517,231.4161) .. controls (281.7725,231.4161) and (282.3568,230.8318) ..(282.3568,230.1110);
\path[fill=chfourfig5c0,fill opacity=0.5,draw=chfourfig5c0,line width=0.3985pt,draw opacity=0.5] (326.0735,205.5328) .. controls (326.0735,204.8120) and (325.4891,204.2277) ..(324.7683,204.2277) .. controls (324.0475,204.2277) and (323.4632,204.8120) ..(323.4632,205.5328) .. controls (323.4632,206.2536) and (324.0475,206.8380) ..(324.7683,206.8380) .. controls (325.4891,206.8380) and (326.0735,206.2536) ..(326.0735,205.5328);
\path[fill=chfourfig5c0,fill opacity=0.5,draw=chfourfig5c0,line width=0.3985pt,draw opacity=0.5] (338.2778,219.7796) .. controls (338.2778,219.0588) and (337.6935,218.4745) ..(336.9727,218.4745) .. controls (336.2519,218.4745) and (335.6676,219.0588) ..(335.6676,219.7796) .. controls (335.6676,220.5004) and (336.2519,221.0847) ..(336.9727,221.0847) .. controls (337.6935,221.0847) and (338.2778,220.5004) ..(338.2778,219.7796);
\path[fill=chfourfig5c0,fill opacity=0.5,draw=chfourfig5c0,line width=0.3985pt,draw opacity=0.5] (293.7343,189.4928) .. controls (293.7343,188.7720) and (293.1500,188.1877) ..(292.4292,188.1877) .. controls (291.7084,188.1877) and (291.1241,188.7720) ..(291.1241,189.4928) .. controls (291.1241,190.2136) and (291.7084,190.7979) ..(292.4292,190.7979) .. controls (293.1500,190.7979) and (293.7343,190.2136) ..(293.7343,189.4928);
\path[fill=chfourfig5c0,fill opacity=0.5,draw=chfourfig5c0,line width=0.3985pt,draw opacity=0.5] (295.4380,203.4406) .. controls (295.4380,202.7198) and (294.8536,202.1355) ..(294.1328,202.1355) .. controls (293.4120,202.1355) and (292.8277,202.7198) ..(292.8277,203.4406) .. controls (292.8277,204.1614) and (293.4120,204.7458) ..(294.1328,204.7458) .. controls (294.8536,204.7458) and (295.4380,204.1614) ..(295.4380,203.4406);
\path[fill=chfourfig5c0,fill opacity=0.5,draw=chfourfig5c0,line width=0.3985pt,draw opacity=0.5] (280.1451,218.6837) .. controls (280.1451,217.9629) and (279.5608,217.3786) ..(278.8400,217.3786) .. controls (278.1192,217.3786) and (277.5349,217.9629) ..(277.5349,218.6837) .. controls (277.5349,219.4045) and (278.1192,219.9888) ..(278.8400,219.9888) .. controls (279.5608,219.9888) and (280.1451,219.4045) ..(280.1451,218.6837);
\path[fill=chfourfig5c0,fill opacity=0.5,draw=chfourfig5c0,line width=0.3985pt,draw opacity=0.5] (314.2875,221.2740) .. controls (314.2875,220.5532) and (313.7032,219.9689) ..(312.9824,219.9689) .. controls (312.2616,219.9689) and (311.6773,220.5532) ..(311.6773,221.2740) .. controls (311.6773,221.9948) and (312.2616,222.5791) ..(312.9824,222.5791) .. controls (313.7032,222.5791) and (314.2875,221.9948) ..(314.2875,221.2740);
\path[fill=chfourfig5c0,fill opacity=0.5,draw=chfourfig5c0,line width=0.3985pt,draw opacity=0.5] (296.4442,232.7312) .. controls (296.4442,232.0104) and (295.8599,231.4261) ..(295.1391,231.4261) .. controls (294.4183,231.4261) and (293.8340,232.0104) ..(293.8340,232.7312) .. controls (293.8340,233.4520) and (294.4183,234.0363) ..(295.1391,234.0363) .. controls (295.8599,234.0363) and (296.4442,233.4520) ..(296.4442,232.7312);
\path[fill=chfourfig5c0,fill opacity=0.5,draw=chfourfig5c0,line width=0.3985pt,draw opacity=0.5] (335.2094,224.4123) .. controls (335.2094,223.6915) and (334.6250,223.1072) ..(333.9042,223.1072) .. controls (333.1834,223.1072) and (332.5991,223.6915) ..(332.5991,224.4123) .. controls (332.5991,225.1331) and (333.1834,225.7174) ..(333.9042,225.7174) .. controls (334.6250,225.7174) and (335.2094,225.1331) ..(335.2094,224.4123);
\path[fill=chfourfig5c0,fill opacity=0.5,draw=chfourfig5c0,line width=0.3985pt,draw opacity=0.5] (281.2609,213.2241) .. controls (281.2609,212.5033) and (280.6766,211.9190) ..(279.9558,211.9190) .. controls (279.2350,211.9190) and (278.6507,212.5033) ..(278.6507,213.2241) .. controls (278.6507,213.9449) and (279.2350,214.5292) ..(279.9558,214.5292) .. controls (280.6766,214.5292) and (281.2609,213.9449) ..(281.2609,213.2241);
\path[fill=chfourfig5c0,fill opacity=0.5,draw=chfourfig5c0,line width=0.3985pt,draw opacity=0.5] (334.7212,222.5692) .. controls (334.7212,221.8484) and (334.1368,221.2640) ..(333.4160,221.2640) .. controls (332.6952,221.2640) and (332.1109,221.8484) ..(332.1109,222.5692) .. controls (332.1109,223.2900) and (332.6952,223.8743) ..(333.4160,223.8743) .. controls (334.1368,223.8743) and (334.7212,223.2900) ..(334.7212,222.5692);
\path[fill=chfourfig5c0,fill opacity=0.5,draw=chfourfig5c0,line width=0.3985pt,draw opacity=0.5] (257.2805,198.4095) .. controls (257.2805,197.6887) and (256.6962,197.1043) ..(255.9754,197.1043) .. controls (255.2546,197.1043) and (254.6703,197.6887) ..(254.6703,198.4095) .. controls (254.6703,199.1302) and (255.2546,199.7146) ..(255.9754,199.7146) .. controls (256.6962,199.7146) and (257.2805,199.1302) ..(257.2805,198.4095);
\path[fill=chfourfig5c0,fill opacity=0.5,draw=chfourfig5c0,line width=0.3985pt,draw opacity=0.5] (256.9717,212.6761) .. controls (256.9717,211.9553) and (256.3874,211.3710) ..(255.6666,211.3710) .. controls (254.9458,211.3710) and (254.3615,211.9553) ..(254.3615,212.6761) .. controls (254.3615,213.3969) and (254.9458,213.9813) ..(255.6666,213.9813) .. controls (256.3874,213.9813) and (256.9717,213.3969) ..(256.9717,212.6761);
\path[fill=chfourfig5c0,fill opacity=0.5,draw=chfourfig5c0,line width=0.3985pt,draw opacity=0.5] (287.3382,240.3527) .. controls (287.3382,239.6319) and (286.7539,239.0476) ..(286.0331,239.0476) .. controls (285.3123,239.0476) and (284.7280,239.6319) ..(284.7280,240.3527) .. controls (284.7280,241.0735) and (285.3123,241.6578) ..(286.0331,241.6578) .. controls (286.7539,241.6578) and (287.3382,241.0735) ..(287.3382,240.3527);
\path[fill=chfourfig5c0,fill opacity=0.5,draw=chfourfig5c0,line width=0.3985pt,draw opacity=0.5] (274.4663,183.5052) .. controls (274.4663,182.7843) and (273.8820,182.2000) ..(273.1612,182.2000) .. controls (272.4404,182.2000) and (271.8561,182.7843) ..(271.8561,183.5052) .. controls (271.8561,184.2260) and (272.4404,184.8103) ..(273.1612,184.8103) .. controls (273.8820,184.8103) and (274.4663,184.2260) ..(274.4663,183.5052);
\path[fill=chfourfig5c0,fill opacity=0.5,draw=chfourfig5c0,line width=0.3985pt,draw opacity=0.5] (275.5522,223.9141) .. controls (275.5522,223.1934) and (274.9679,222.6090) ..(274.2472,222.6090) .. controls (273.5264,222.6090) and (272.9420,223.1934) ..(272.9420,223.9141) .. controls (272.9420,224.6349) and (273.5264,225.2192) ..(274.2472,225.2192) .. controls (274.9679,225.2192) and (275.5522,224.6349) ..(275.5522,223.9141);
\path[fill=chfourfig5c0,fill opacity=0.5,draw=chfourfig5c0,line width=0.3985pt,draw opacity=0.5] (225.8580,229.0350) .. controls (225.8580,228.3142) and (225.2737,227.7299) ..(224.5529,227.7299) .. controls (223.8320,227.7299) and (223.2477,228.3142) ..(223.2477,229.0350) .. controls (223.2477,229.7558) and (223.8320,230.3401) ..(224.5529,230.3401) .. controls (225.2737,230.3401) and (225.8580,229.7558) ..(225.8580,229.0350);
\path[fill=chfourfig5c0,fill opacity=0.5,draw=chfourfig5c0,line width=0.3985pt,draw opacity=0.5] (285.8936,252.9357) .. controls (285.8936,252.2149) and (285.3093,251.6306) ..(284.5885,251.6306) .. controls (283.8677,251.6306) and (283.2834,252.2149) ..(283.2834,252.9357) .. controls (283.2834,253.6565) and (283.8677,254.2408) ..(284.5885,254.2408) .. controls (285.3093,254.2408) and (285.8936,253.6565) ..(285.8936,252.9357);
\path[fill=chfourfig5c0,fill opacity=0.5,draw=chfourfig5c0,line width=0.3985pt,draw opacity=0.5] (332.0312,213.9414) .. controls (332.0312,213.2206) and (331.4469,212.6363) ..(330.7261,212.6363) .. controls (330.0053,212.6363) and (329.4210,213.2206) ..(329.4210,213.9414) .. controls (329.4210,214.6622) and (330.0053,215.2465) ..(330.7261,215.2465) .. controls (331.4469,215.2465) and (332.0312,214.6622) ..(332.0312,213.9414);
\path[fill=chfourfig5c0,fill opacity=0.5,draw=chfourfig5c0,line width=0.3985pt,draw opacity=0.5] (290.1975,219.7397) .. controls (290.1975,219.0189) and (289.6132,218.4346) ..(288.8924,218.4346) .. controls (288.1716,218.4346) and (287.5873,219.0189) ..(287.5873,219.7397) .. controls (287.5873,220.4605) and (288.1716,221.0449) ..(288.8924,221.0449) .. controls (289.6132,221.0449) and (290.1975,220.4605) ..(290.1975,219.7397);
\path[fill=chfourfig5c0,fill opacity=0.5,draw=chfourfig5c0,line width=0.3985pt,draw opacity=0.5] (272.8922,209.5877) .. controls (272.8922,208.8669) and (272.3079,208.2825) ..(271.5871,208.2825) .. controls (270.8663,208.2825) and (270.2820,208.8669) ..(270.2820,209.5877) .. controls (270.2820,210.3085) and (270.8663,210.8928) ..(271.5871,210.8928) .. controls (272.3079,210.8928) and (272.8922,210.3085) ..(272.8922,209.5877);
\path[fill=chfourfig5c0,fill opacity=0.5,draw=chfourfig5c0,line width=0.3985pt,draw opacity=0.5] (299.7817,263.8449) .. controls (299.7817,263.1241) and (299.1974,262.5398) ..(298.4766,262.5398) .. controls (297.7558,262.5398) and (297.1715,263.1241) ..(297.1715,263.8449) .. controls (297.1715,264.5657) and (297.7558,265.1500) ..(298.4766,265.1500) .. controls (299.1974,265.1500) and (299.7817,264.5657) ..(299.7817,263.8449);
\path[fill=chfourfig5c0,fill opacity=0.5,draw=chfourfig5c0,line width=0.3985pt,draw opacity=0.5] (282.6657,184.0830) .. controls (282.6657,183.3622) and (282.0814,182.7779) ..(281.3606,182.7779) .. controls (280.6398,182.7779) and (280.0555,183.3622) ..(280.0555,184.0830) .. controls (280.0555,184.8038) and (280.6398,185.3881) ..(281.3606,185.3881) .. controls (282.0814,185.3881) and (282.6657,184.8038) ..(282.6657,184.0830);
\path[fill=chfourfig5c0,fill opacity=0.5,draw=chfourfig5c0,line width=0.3985pt,draw opacity=0.5] (250.3465,234.5544) .. controls (250.3465,233.8336) and (249.7621,233.2493) ..(249.0413,233.2493) .. controls (248.3206,233.2493) and (247.7362,233.8336) ..(247.7362,234.5544) .. controls (247.7362,235.2752) and (248.3206,235.8595) ..(249.0413,235.8595) .. controls (249.7621,235.8595) and (250.3465,235.2752) ..(250.3465,234.5544);
\path[fill=chfourfig5c0,fill opacity=0.5,draw=chfourfig5c0,line width=0.3985pt,draw opacity=0.5] (235.1034,223.1869) .. controls (235.1034,222.4661) and (234.5191,221.8817) ..(233.7983,221.8817) .. controls (233.0775,221.8817) and (232.4932,222.4661) ..(232.4932,223.1869) .. controls (232.4932,223.9077) and (233.0775,224.4920) ..(233.7983,224.4920) .. controls (234.5191,224.4920) and (235.1034,223.9077) ..(235.1034,223.1869);
\path[fill=chfourfig5c0,fill opacity=0.5,draw=chfourfig5c0,line width=0.3985pt,draw opacity=0.5] (246.1123,225.9067) .. controls (246.1123,225.1859) and (245.5280,224.6016) ..(244.8072,224.6016) .. controls (244.0864,224.6016) and (243.5020,225.1859) ..(243.5020,225.9067) .. controls (243.5020,226.6275) and (244.0864,227.2118) ..(244.8072,227.2118) .. controls (245.5280,227.2118) and (246.1123,226.6275) ..(246.1123,225.9067);
\path[fill=chfourfig5c0,fill opacity=0.5,draw=chfourfig5c0,line width=0.3985pt,draw opacity=0.5] (252.4486,203.0621) .. controls (252.4486,202.3412) and (251.8643,201.7570) ..(251.1435,201.7570) .. controls (250.4227,201.7570) and (249.8384,202.3412) ..(249.8384,203.0621) .. controls (249.8384,203.7829) and (250.4227,204.3672) ..(251.1435,204.3672) .. controls (251.8643,204.3672) and (252.4486,203.7829) ..(252.4486,203.0621);
\path[fill=chfourfig5c0,fill opacity=0.5,draw=chfourfig5c0,line width=0.3985pt,draw opacity=0.5] (306.5365,242.9032) .. controls (306.5365,242.1824) and (305.9521,241.5981) ..(305.2314,241.5981) .. controls (304.5106,241.5981) and (303.9262,242.1824) ..(303.9262,242.9032) .. controls (303.9262,243.6240) and (304.5106,244.2083) ..(305.2314,244.2083) .. controls (305.9521,244.2083) and (306.5365,243.6240) ..(306.5365,242.9032);
\path[fill=chfourfig5c0,fill opacity=0.5,draw=chfourfig5c0,line width=0.3985pt,draw opacity=0.5] (305.4306,231.9043) .. controls (305.4306,231.1835) and (304.8463,230.5992) ..(304.1255,230.5992) .. controls (303.4047,230.5992) and (302.8204,231.1835) ..(302.8204,231.9043) .. controls (302.8204,232.6251) and (303.4047,233.2094) ..(304.1255,233.2094) .. controls (304.8463,233.2094) and (305.4306,232.6251) ..(305.4306,231.9043);
\path[fill=chfourfig5c0,fill opacity=0.5,draw=chfourfig5c0,line width=0.3985pt,draw opacity=0.5] (279.6768,236.2779) .. controls (279.6768,235.5571) and (279.0925,234.9728) ..(278.3717,234.9728) .. controls (277.6509,234.9728) and (277.0666,235.5571) ..(277.0666,236.2779) .. controls (277.0666,236.9987) and (277.6509,237.5831) ..(278.3717,237.5831) .. controls (279.0925,237.5831) and (279.6768,236.9987) ..(279.6768,236.2779);
\path[fill=chfourfig5c0,fill opacity=0.5,draw=chfourfig5c0,line width=0.3985pt,draw opacity=0.5] (272.4738,250.7140) .. controls (272.4738,249.9932) and (271.8894,249.4088) ..(271.1687,249.4088) .. controls (270.4478,249.4088) and (269.8635,249.9932) ..(269.8635,250.7140) .. controls (269.8635,251.4348) and (270.4478,252.0191) ..(271.1687,252.0191) .. controls (271.8894,252.0191) and (272.4738,251.4348) ..(272.4738,250.7140);
\path[fill=chfourfig5c0,fill opacity=0.5,draw=chfourfig5c0,line width=0.3985pt,draw opacity=0.5] (313.2514,226.6240) .. controls (313.2514,225.9032) and (312.6671,225.3189) ..(311.9463,225.3189) .. controls (311.2255,225.3189) and (310.6411,225.9032) ..(310.6411,226.6240) .. controls (310.6411,227.3448) and (311.2255,227.9291) ..(311.9463,227.9291) .. controls (312.6671,227.9291) and (313.2514,227.3448) ..(313.2514,226.6240);
\path[fill=chfourfig5c0,fill opacity=0.5,draw=chfourfig5c0,line width=0.3985pt,draw opacity=0.5] (255.3876,192.2126) .. controls (255.3876,191.4918) and (254.8033,190.9075) ..(254.0825,190.9075) .. controls (253.3617,190.9075) and (252.7774,191.4918) ..(252.7774,192.2126) .. controls (252.7774,192.9334) and (253.3617,193.5177) ..(254.0825,193.5177) .. controls (254.8033,193.5177) and (255.3876,192.9334) ..(255.3876,192.2126);
\path[fill=chfourfig5c0,fill opacity=0.5,draw=chfourfig5c0,line width=0.3985pt,draw opacity=0.5] (255.6566,202.7632) .. controls (255.6566,202.0424) and (255.0723,201.4581) ..(254.3515,201.4581) .. controls (253.6307,201.4581) and (253.0464,202.0424) ..(253.0464,202.7632) .. controls (253.0464,203.4840) and (253.6307,204.0683) ..(254.3515,204.0683) .. controls (255.0723,204.0683) and (255.6566,203.4840) ..(255.6566,202.7632);
\path[fill=chfourfig5c0,fill opacity=0.5,draw=chfourfig5c0,line width=0.3985pt,draw opacity=0.5] (282.6956,251.3018) .. controls (282.6956,250.5810) and (282.1112,249.9966) ..(281.3904,249.9966) .. controls (280.6696,249.9966) and (280.0853,250.5810) ..(280.0853,251.3018) .. controls (280.0853,252.0226) and (280.6696,252.6069) ..(281.3904,252.6069) .. controls (282.1112,252.6069) and (282.6956,252.0226) ..(282.6956,251.3018);
\path[fill=chfourfig5c0,fill opacity=0.5,draw=chfourfig5c0,line width=0.3985pt,draw opacity=0.5] (300.2998,253.0154) .. controls (300.2998,252.2946) and (299.7155,251.7103) ..(298.9947,251.7103) .. controls (298.2739,251.7103) and (297.6895,252.2946) ..(297.6895,253.0154) .. controls (297.6895,253.7362) and (298.2739,254.3205) ..(298.9947,254.3205) .. controls (299.7155,254.3205) and (300.2998,253.7362) ..(300.2998,253.0154);
\path[fill=chfourfig5c0,fill opacity=0.5,draw=chfourfig5c0,line width=0.3985pt,draw opacity=0.5] (284.3195,212.7060) .. controls (284.3195,211.9852) and (283.7352,211.4009) ..(283.0144,211.4009) .. controls (282.2936,211.4009) and (281.7093,211.9852) ..(281.7093,212.7060) .. controls (281.7093,213.4268) and (282.2936,214.0111) ..(283.0144,214.0111) .. controls (283.7352,214.0111) and (284.3195,213.4268) ..(284.3195,212.7060);
\path[fill=chfourfig5c0,fill opacity=0.5,draw=chfourfig5c0,line width=0.3985pt,draw opacity=0.5] (320.3748,204.4469) .. controls (320.3748,203.7261) and (319.7904,203.1418) ..(319.0696,203.1418) .. controls (318.3488,203.1418) and (317.7645,203.7261) ..(317.7645,204.4469) .. controls (317.7645,205.1677) and (318.3488,205.7520) ..(319.0696,205.7520) .. controls (319.7904,205.7520) and (320.3748,205.1677) ..(320.3748,204.4469);
\path[fill=chfourfig5c0,fill opacity=0.5,draw=chfourfig5c0,line width=0.3985pt,draw opacity=0.5] (306.4667,215.7746) .. controls (306.4667,215.0538) and (305.8824,214.4694) ..(305.1616,214.4694) .. controls (304.4408,214.4694) and (303.8565,215.0538) ..(303.8565,215.7746) .. controls (303.8565,216.4954) and (304.4408,217.0797) ..(305.1616,217.0797) .. controls (305.8824,217.0797) and (306.4667,216.4954) ..(306.4667,215.7746);
\path[fill=chfourfig5c0,fill opacity=0.5,draw=chfourfig5c0,line width=0.3985pt,draw opacity=0.5] (263.3379,197.8217) .. controls (263.3379,197.1008) and (262.7536,196.5165) ..(262.0328,196.5165) .. controls (261.3120,196.5165) and (260.7277,197.1008) ..(260.7277,197.8217) .. controls (260.7277,198.5424) and (261.3120,199.1268) ..(262.0328,199.1268) .. controls (262.7536,199.1268) and (263.3379,198.5424) ..(263.3379,197.8217);
\path[fill=chfourfig5c0,fill opacity=0.5,draw=chfourfig5c0,line width=0.3985pt,draw opacity=0.5] (313.9986,247.6953) .. controls (313.9986,246.9745) and (313.4142,246.3901) ..(312.6935,246.3901) .. controls (311.9727,246.3901) and (311.3883,246.9745) ..(311.3883,247.6953) .. controls (311.3883,248.4161) and (311.9727,249.0004) ..(312.6935,249.0004) .. controls (313.4142,249.0004) and (313.9986,248.4161) ..(313.9986,247.6953);
\path[fill=chfourfig5c0,fill opacity=0.5,draw=chfourfig5c0,line width=0.3985pt,draw opacity=0.5] (244.9367,214.8480) .. controls (244.9367,214.1272) and (244.3524,213.5429) ..(243.6316,213.5429) .. controls (242.9108,213.5429) and (242.3264,214.1272) ..(242.3264,214.8480) .. controls (242.3264,215.5688) and (242.9108,216.1531) ..(243.6316,216.1531) .. controls (244.3524,216.1531) and (244.9367,215.5688) ..(244.9367,214.8480);
\path[fill=chfourfig5c0,fill opacity=0.5,draw=chfourfig5c0,line width=0.3985pt,draw opacity=0.5] (281.0816,212.8256) .. controls (281.0816,212.1048) and (280.4973,211.5204) ..(279.7765,211.5204) .. controls (279.0557,211.5204) and (278.4713,212.1048) ..(278.4713,212.8256) .. controls (278.4713,213.5464) and (279.0557,214.1307) ..(279.7765,214.1307) .. controls (280.4973,214.1307) and (281.0816,213.5464) ..(281.0816,212.8256);
\path[fill=chfourfig5c0,fill opacity=0.5,draw=chfourfig5c0,line width=0.3985pt,draw opacity=0.5] (276.1002,245.8422) .. controls (276.1002,245.1214) and (275.5159,244.5370) ..(274.7951,244.5370) .. controls (274.0743,244.5370) and (273.4900,245.1214) ..(273.4900,245.8422) .. controls (273.4900,246.5630) and (274.0743,247.1473) ..(274.7951,247.1473) .. controls (275.5159,247.1473) and (276.1002,246.5630) ..(276.1002,245.8422);
\path[fill=chfourfig5c0,fill opacity=0.5,draw=chfourfig5c0,line width=0.3985pt,draw opacity=0.5] (297.1216,199.1567) .. controls (297.1216,198.4359) and (296.5373,197.8515) ..(295.8165,197.8515) .. controls (295.0957,197.8515) and (294.5114,198.4359) ..(294.5114,199.1567) .. controls (294.5114,199.8774) and (295.0957,200.4618) ..(295.8165,200.4618) .. controls (296.5373,200.4618) and (297.1216,199.8774) ..(297.1216,199.1567);
\path[fill=chfourfig5c0,fill opacity=0.5,draw=chfourfig5c0,line width=0.3985pt,draw opacity=0.5] (243.5319,234.1260) .. controls (243.5319,233.4052) and (242.9476,232.8209) ..(242.2268,232.8209) .. controls (241.5060,232.8209) and (240.9217,233.4052) ..(240.9217,234.1260) .. controls (240.9217,234.8468) and (241.5060,235.4311) ..(242.2268,235.4311) .. controls (242.9476,235.4311) and (243.5319,234.8468) ..(243.5319,234.1260);
\path[fill=chfourfig5c0,fill opacity=0.5,draw=chfourfig5c0,line width=0.3985pt,draw opacity=0.5] (250.4261,201.2787) .. controls (250.4261,200.5579) and (249.8419,199.9736) ..(249.1210,199.9736) .. controls (248.4002,199.9736) and (247.8159,200.5579) ..(247.8159,201.2787) .. controls (247.8159,201.9995) and (248.4002,202.5839) ..(249.1210,202.5839) .. controls (249.8419,202.5839) and (250.4261,201.9995) ..(250.4261,201.2787);
\path[fill=chfourfig5c0,fill opacity=0.5,draw=chfourfig5c0,line width=0.3985pt,draw opacity=0.5] (279.0691,224.2230) .. controls (279.0691,223.5022) and (278.4848,222.9178) ..(277.7640,222.9178) .. controls (277.0432,222.9178) and (276.4589,223.5022) ..(276.4589,224.2230) .. controls (276.4589,224.9438) and (277.0432,225.5281) ..(277.7640,225.5281) .. controls (278.4848,225.5281) and (279.0691,224.9438) ..(279.0691,224.2230);
\path[fill=chfourfig5c0,fill opacity=0.5,draw=chfourfig5c0,line width=0.3985pt,draw opacity=0.5] (269.7739,189.9012) .. controls (269.7739,189.1805) and (269.1895,188.5961) ..(268.4688,188.5961) .. controls (267.7479,188.5961) and (267.1636,189.1805) ..(267.1636,189.9012) .. controls (267.1636,190.6221) and (267.7479,191.2064) ..(268.4688,191.2064) .. controls (269.1895,191.2064) and (269.7739,190.6221) ..(269.7739,189.9012);
\path[fill=chfourfig5c0,fill opacity=0.5,draw=chfourfig5c0,line width=0.3985pt,draw opacity=0.5] (287.2386,213.7521) .. controls (287.2386,213.0313) and (286.6543,212.4470) ..(285.9335,212.4470) .. controls (285.2126,212.4470) and (284.6284,213.0313) ..(284.6284,213.7521) .. controls (284.6284,214.4729) and (285.2126,215.0573) ..(285.9335,215.0573) .. controls (286.6543,215.0573) and (287.2386,214.4729) ..(287.2386,213.7521);
\path[fill=chfourfig5c0,fill opacity=0.5,draw=chfourfig5c0,line width=0.3985pt,draw opacity=0.5] (269.0565,220.1083) .. controls (269.0565,219.3876) and (268.4722,218.8032) ..(267.7514,218.8032) .. controls (267.0306,218.8032) and (266.4463,219.3876) ..(266.4463,220.1083) .. controls (266.4463,220.8292) and (267.0306,221.4135) ..(267.7514,221.4135) .. controls (268.4722,221.4135) and (269.0565,220.8292) ..(269.0565,220.1083);
\path[fill=chfourfig5c0,fill opacity=0.5,draw=chfourfig5c0,line width=0.3985pt,draw opacity=0.5] (293.6447,212.6761) .. controls (293.6447,211.9553) and (293.0603,211.3710) ..(292.3395,211.3710) .. controls (291.6187,211.3710) and (291.0344,211.9553) ..(291.0344,212.6761) .. controls (291.0344,213.3969) and (291.6187,213.9813) ..(292.3395,213.9813) .. controls (293.0603,213.9813) and (293.6447,213.3969) ..(293.6447,212.6761);
\path[fill=chfourfig5c0,fill opacity=0.5,draw=chfourfig5c0,line width=0.3985pt,draw opacity=0.5] (288.9721,198.6785) .. controls (288.9721,197.9576) and (288.3878,197.3733) ..(287.6670,197.3733) .. controls (286.9462,197.3733) and (286.3619,197.9576) ..(286.3619,198.6785) .. controls (286.3619,199.3992) and (286.9462,199.9836) ..(287.6670,199.9836) .. controls (288.3878,199.9836) and (288.9721,199.3992) ..(288.9721,198.6785);
\path[fill=chfourfig5c0,fill opacity=0.5,draw=chfourfig5c0,line width=0.3985pt,draw opacity=0.5] (304.6535,195.3509) .. controls (304.6535,194.6301) and (304.0692,194.0458) ..(303.3484,194.0458) .. controls (302.6276,194.0458) and (302.0432,194.6301) ..(302.0432,195.3509) .. controls (302.0432,196.0717) and (302.6276,196.6560) ..(303.3484,196.6560) .. controls (304.0692,196.6560) and (304.6535,196.0717) ..(304.6535,195.3509);
\path[fill=chfourfig5c0,fill opacity=0.5,draw=chfourfig5c0,line width=0.3985pt,draw opacity=0.5] (296.7231,211.8393) .. controls (296.7231,211.1185) and (296.1388,210.5341) ..(295.4180,210.5341) .. controls (294.6972,210.5341) and (294.1129,211.1185) ..(294.1129,211.8393) .. controls (294.1129,212.5601) and (294.6972,213.1444) ..(295.4180,213.1444) .. controls (296.1388,213.1444) and (296.7231,212.5601) ..(296.7231,211.8393);
\path[fill=chfourfig5c0,fill opacity=0.5,draw=chfourfig5c0,line width=0.3985pt,draw opacity=0.5] (261.0465,230.3900) .. controls (261.0465,229.6691) and (260.4622,229.0848) ..(259.7414,229.0848) .. controls (259.0206,229.0848) and (258.4362,229.6691) ..(258.4362,230.3900) .. controls (258.4362,231.1107) and (259.0206,231.6951) ..(259.7414,231.6951) .. controls (260.4622,231.6951) and (261.0465,231.1107) ..(261.0465,230.3900);
\path[fill=chfourfig5c0,fill opacity=0.5,draw=chfourfig5c0,line width=0.3985pt,draw opacity=0.5] (281.5299,238.7786) .. controls (281.5299,238.0578) and (280.9456,237.4734) ..(280.2248,237.4734) .. controls (279.5040,237.4734) and (278.9197,238.0578) ..(278.9197,238.7786) .. controls (278.9197,239.4994) and (279.5040,240.0837) ..(280.2248,240.0837) .. controls (280.9456,240.0837) and (281.5299,239.4994) ..(281.5299,238.7786);
\path[fill=chfourfig5c0,fill opacity=0.5,draw=chfourfig5c0,line width=0.3985pt,draw opacity=0.5] (315.6624,245.2643) .. controls (315.6624,244.5435) and (315.0780,243.9592) ..(314.3573,243.9592) .. controls (313.6364,243.9592) and (313.0521,244.5435) ..(313.0521,245.2643) .. controls (313.0521,245.9852) and (313.6364,246.5695) ..(314.3573,246.5695) .. controls (315.0780,246.5695) and (315.6624,245.9852) ..(315.6624,245.2643);
\path[fill=chfourfig5c0,fill opacity=0.5,draw=chfourfig5c0,line width=0.3985pt,draw opacity=0.5] (305.0022,213.3237) .. controls (305.0022,212.6029) and (304.4179,212.0186) ..(303.6971,212.0186) .. controls (302.9763,212.0186) and (302.3920,212.6029) ..(302.3920,213.3237) .. controls (302.3920,214.0445) and (302.9763,214.6288) ..(303.6971,214.6288) .. controls (304.4179,214.6288) and (305.0022,214.0445) ..(305.0022,213.3237);
\path[fill=chfourfig5c0,fill opacity=0.5,draw=chfourfig5c0,line width=0.3985pt,draw opacity=0.5] (287.5076,216.5317) .. controls (287.5076,215.8109) and (286.9233,215.2266) ..(286.2025,215.2266) .. controls (285.4817,215.2266) and (284.8974,215.8109) ..(284.8974,216.5317) .. controls (284.8974,217.2525) and (285.4817,217.8369) ..(286.2025,217.8369) .. controls (286.9233,217.8369) and (287.5076,217.2525) ..(287.5076,216.5317);
\path[fill=chfourfig5c0,fill opacity=0.5,draw=chfourfig5c0,line width=0.3985pt,draw opacity=0.5] (270.9893,187.7294) .. controls (270.9893,187.0086) and (270.4050,186.4243) ..(269.6842,186.4243) .. controls (268.9634,186.4243) and (268.3791,187.0086) ..(268.3791,187.7294) .. controls (268.3791,188.4502) and (268.9634,189.0345) ..(269.6842,189.0345) .. controls (270.4050,189.0345) and (270.9893,188.4502) ..(270.9893,187.7294);
\path[fill=chfourfig5c0,fill opacity=0.5,draw=chfourfig5c0,line width=0.3985pt,draw opacity=0.5] (243.7212,205.1343) .. controls (243.7212,204.4135) and (243.1369,203.8292) ..(242.4161,203.8292) .. controls (241.6953,203.8292) and (241.1110,204.4135) ..(241.1110,205.1343) .. controls (241.1110,205.8551) and (241.6953,206.4395) ..(242.4161,206.4395) .. controls (243.1369,206.4395) and (243.7212,205.8551) ..(243.7212,205.1343);
\path[fill=chfourfig5c0,fill opacity=0.5,draw=chfourfig5c0,line width=0.3985pt,draw opacity=0.5] (325.5056,228.3077) .. controls (325.5056,227.5869) and (324.9213,227.0026) ..(324.2005,227.0026) .. controls (323.4796,227.0026) and (322.8953,227.5869) ..(322.8953,228.3077) .. controls (322.8953,229.0285) and (323.4796,229.6128) ..(324.2005,229.6128) .. controls (324.9213,229.6128) and (325.5056,229.0285) ..(325.5056,228.3077);
\path[fill=chfourfig5c0,fill opacity=0.5,draw=chfourfig5c0,line width=0.3985pt,draw opacity=0.5] (332.7585,226.6140) .. controls (332.7585,225.8932) and (332.1742,225.3089) ..(331.4534,225.3089) .. controls (330.7325,225.3089) and (330.1482,225.8932) ..(330.1482,226.6140) .. controls (330.1482,227.3348) and (330.7325,227.9192) ..(331.4534,227.9192) .. controls (332.1742,227.9192) and (332.7585,227.3348) ..(332.7585,226.6140);
\path[fill=chfourfig5c0,fill opacity=0.5,draw=chfourfig5c0,line width=0.3985pt,draw opacity=0.5] (316.7085,212.2278) .. controls (316.7085,211.5070) and (316.1241,210.9227) ..(315.4034,210.9227) .. controls (314.6825,210.9227) and (314.0982,211.5070) ..(314.0982,212.2278) .. controls (314.0982,212.9486) and (314.6825,213.5330) ..(315.4034,213.5330) .. controls (316.1241,213.5330) and (316.7085,212.9486) ..(316.7085,212.2278);
\path[fill=chfourfig5c0,fill opacity=0.5,draw=chfourfig5c0,line width=0.3985pt,draw opacity=0.5] (273.2808,202.8429) .. controls (273.2808,202.1221) and (272.6964,201.5378) ..(271.9756,201.5378) .. controls (271.2548,201.5378) and (270.6705,202.1221) ..(270.6705,202.8429) .. controls (270.6705,203.5637) and (271.2548,204.1480) ..(271.9756,204.1480) .. controls (272.6964,204.1480) and (273.2808,203.5637) ..(273.2808,202.8429);
\path[fill=chfourfig5c0,fill opacity=0.5,draw=chfourfig5c0,line width=0.3985pt,draw opacity=0.5] (281.0916,242.1859) .. controls (281.0916,241.4651) and (280.5072,240.8807) ..(279.7864,240.8807) .. controls (279.0656,240.8807) and (278.4813,241.4651) ..(278.4813,242.1859) .. controls (278.4813,242.9067) and (279.0656,243.4910) ..(279.7864,243.4910) .. controls (280.5072,243.4910) and (281.0916,242.9067) ..(281.0916,242.1859);
\path[fill=chfourfig5c0,fill opacity=0.5,draw=chfourfig5c0,line width=0.3985pt,draw opacity=0.5] (284.2199,244.3478) .. controls (284.2199,243.6270) and (283.6356,243.0427) ..(282.9147,243.0427) .. controls (282.1940,243.0427) and (281.6096,243.6270) ..(281.6096,244.3478) .. controls (281.6096,245.0686) and (282.1940,245.6529) ..(282.9147,245.6529) .. controls (283.6356,245.6529) and (284.2199,245.0686) ..(284.2199,244.3478);
\path[fill=chfourfig5c0,fill opacity=0.5,draw=chfourfig5c0,line width=0.3985pt,draw opacity=0.5] (304.3945,195.9984) .. controls (304.3945,195.2776) and (303.8102,194.6934) ..(303.0894,194.6934) .. controls (302.3686,194.6934) and (301.7842,195.2776) ..(301.7842,195.9984) .. controls (301.7842,196.7192) and (302.3686,197.3036) ..(303.0894,197.3036) .. controls (303.8102,197.3036) and (304.3945,196.7192) ..(304.3945,195.9984);
\path[fill=chfourfig5c0,fill opacity=0.5,draw=chfourfig5c0,line width=0.3985pt,draw opacity=0.5] (286.9895,263.5161) .. controls (286.9895,262.7953) and (286.4052,262.2110) ..(285.6844,262.2110) .. controls (284.9636,262.2110) and (284.3793,262.7953) ..(284.3793,263.5161) .. controls (284.3793,264.2369) and (284.9636,264.8213) ..(285.6844,264.8213) .. controls (286.4052,264.8213) and (286.9895,264.2369) ..(286.9895,263.5161);
\path[fill=chfourfig5c0,fill opacity=0.5,draw=chfourfig5c0,line width=0.3985pt,draw opacity=0.5] (296.4442,206.6088) .. controls (296.4442,205.8880) and (295.8599,205.3037) ..(295.1391,205.3037) .. controls (294.4183,205.3037) and (293.8340,205.8880) ..(293.8340,206.6088) .. controls (293.8340,207.3296) and (294.4183,207.9139) ..(295.1391,207.9139) .. controls (295.8599,207.9139) and (296.4442,207.3296) ..(296.4442,206.6088);
\path[fill=chfourfig5c0,fill opacity=0.5,draw=chfourfig5c0,line width=0.3985pt,draw opacity=0.5] (275.6917,233.6278) .. controls (275.6917,232.9070) and (275.1074,232.3227) ..(274.3866,232.3227) .. controls (273.6658,232.3227) and (273.0815,232.9070) ..(273.0815,233.6278) .. controls (273.0815,234.3486) and (273.6658,234.9330) ..(274.3866,234.9330) .. controls (275.1074,234.9330) and (275.6917,234.3486) ..(275.6917,233.6278);
\path[fill=chfourfig5c0,fill opacity=0.5,draw=chfourfig5c0,line width=0.3985pt,draw opacity=0.5] (248.9517,206.2302) .. controls (248.9517,205.5094) and (248.3674,204.9251) ..(247.6465,204.9251) .. controls (246.9258,204.9251) and (246.3414,205.5094) ..(246.3414,206.2302) .. controls (246.3414,206.9510) and (246.9258,207.5353) ..(247.6465,207.5353) .. controls (248.3674,207.5353) and (248.9517,206.9510) ..(248.9517,206.2302);
\path[fill=chfourfig5c0,fill opacity=0.5,draw=chfourfig5c0,line width=0.3985pt,draw opacity=0.5] (275.8611,209.5977) .. controls (275.8611,208.8768) and (275.2768,208.2925) ..(274.5560,208.2925) .. controls (273.8352,208.2925) and (273.2509,208.8768) ..(273.2509,209.5977) .. controls (273.2509,210.3185) and (273.8352,210.9028) ..(274.5560,210.9028) .. controls (275.2768,210.9028) and (275.8611,210.3185) ..(275.8611,209.5977);
\path[fill=chfourfig5c0,fill opacity=0.5,draw=chfourfig5c0,line width=0.3985pt,draw opacity=0.5] (236.9764,224.6215) .. controls (236.9764,223.9007) and (236.3921,223.3164) ..(235.6713,223.3164) .. controls (234.9505,223.3164) and (234.3662,223.9007) ..(234.3662,224.6215) .. controls (234.3662,225.3423) and (234.9505,225.9266) ..(235.6713,225.9266) .. controls (236.3921,225.9266) and (236.9764,225.3423) ..(236.9764,224.6215);
\path[fill=chfourfig5c0,fill opacity=0.5,draw=chfourfig5c0,line width=0.3985pt,draw opacity=0.5] (262.8597,201.4282) .. controls (262.8597,200.7074) and (262.2754,200.1230) ..(261.5546,200.1230) .. controls (260.8338,200.1230) and (260.2495,200.7074) ..(260.2495,201.4282) .. controls (260.2495,202.1490) and (260.8338,202.7333) ..(261.5546,202.7333) .. controls (262.2754,202.7333) and (262.8597,202.1490) ..(262.8597,201.4282);
\path[fill=chfourfig5c0,fill opacity=0.5,draw=chfourfig5c0,line width=0.3985pt,draw opacity=0.5] (281.9583,259.6904) .. controls (281.9583,258.9697) and (281.3740,258.3853) ..(280.6532,258.3853) .. controls (279.9324,258.3853) and (279.3481,258.9697) ..(279.3481,259.6904) .. controls (279.3481,260.4113) and (279.9324,260.9955) ..(280.6532,260.9955) .. controls (281.3740,260.9955) and (281.9583,260.4113) ..(281.9583,259.6904);
\path[fill=chfourfig5c0,fill opacity=0.5,draw=chfourfig5c0,line width=0.3985pt,draw opacity=0.5] (288.7729,232.3227) .. controls (288.7729,231.6019) and (288.1886,231.0176) ..(287.4677,231.0176) .. controls (286.7469,231.0176) and (286.1626,231.6019) ..(286.1626,232.3227) .. controls (286.1626,233.0435) and (286.7469,233.6278) ..(287.4677,233.6278) .. controls (288.1886,233.6278) and (288.7729,233.0435) ..(288.7729,232.3227);
\path[fill=chfourfig5c0,fill opacity=0.5,draw=chfourfig5c0,line width=0.3985pt,draw opacity=0.5] (290.3370,252.2483) .. controls (290.3370,251.5275) and (289.7527,250.9431) ..(289.0319,250.9431) .. controls (288.3111,250.9431) and (287.7268,251.5275) ..(287.7268,252.2483) .. controls (287.7268,252.9691) and (288.3111,253.5534) ..(289.0319,253.5534) .. controls (289.7527,253.5534) and (290.3370,252.9691) ..(290.3370,252.2483);
\path[fill=chfourfig5c0,fill opacity=0.5,draw=chfourfig5c0,line width=0.3985pt,draw opacity=0.5] (265.4899,229.3837) .. controls (265.4899,228.6629) and (264.9055,228.0786) ..(264.1847,228.0786) .. controls (263.4640,228.0786) and (262.8796,228.6629) ..(262.8796,229.3837) .. controls (262.8796,230.1045) and (263.4640,230.6888) ..(264.1847,230.6888) .. controls (264.9055,230.6888) and (265.4899,230.1045) ..(265.4899,229.3837);
\path[fill=chfourfig5c0,fill opacity=0.5,draw=chfourfig5c0,line width=0.3985pt,draw opacity=0.5] (280.6333,184.5512) .. controls (280.6333,183.8304) and (280.0490,183.2461) ..(279.3282,183.2461) .. controls (278.6074,183.2461) and (278.0230,183.8304) ..(278.0230,184.5512) .. controls (278.0230,185.2720) and (278.6074,185.8564) ..(279.3282,185.8564) .. controls (280.0490,185.8564) and (280.6333,185.2720) ..(280.6333,184.5512);
\path[fill=chfourfig5c0,fill opacity=0.5,draw=chfourfig5c0,line width=0.3985pt,draw opacity=0.5] (336.6539,210.3847) .. controls (336.6539,209.6639) and (336.0696,209.0796) ..(335.3488,209.0796) .. controls (334.6280,209.0796) and (334.0437,209.6639) ..(334.0437,210.3847) .. controls (334.0437,211.1055) and (334.6280,211.6898) ..(335.3488,211.6898) .. controls (336.0696,211.6898) and (336.6539,211.1055) ..(336.6539,210.3847);
\path[fill=chfourfig5c0,fill opacity=0.5,draw=chfourfig5c0,line width=0.3985pt,draw opacity=0.5] (314.2178,204.7856) .. controls (314.2178,204.0648) and (313.6335,203.4805) ..(312.9127,203.4805) .. controls (312.1918,203.4805) and (311.6075,204.0648) ..(311.6075,204.7856) .. controls (311.6075,205.5064) and (312.1918,206.0908) ..(312.9127,206.0908) .. controls (313.6335,206.0908) and (314.2178,205.5064) ..(314.2178,204.7856);
\path[fill=chfourfig5c0,fill opacity=0.5,draw=chfourfig5c0,line width=0.3985pt,draw opacity=0.5] (312.6436,220.2080) .. controls (312.6436,219.4872) and (312.0593,218.9029) ..(311.3385,218.9029) .. controls (310.6177,218.9029) and (310.0334,219.4872) ..(310.0334,220.2080) .. controls (310.0334,220.9288) and (310.6177,221.5131) ..(311.3385,221.5131) .. controls (312.0593,221.5131) and (312.6436,220.9288) ..(312.6436,220.2080);
\path[fill=chfourfig5c0,fill opacity=0.5,draw=chfourfig5c0,line width=0.3985pt,draw opacity=0.5] (264.4239,254.8087) .. controls (264.4239,254.0879) and (263.8395,253.5035) ..(263.1187,253.5035) .. controls (262.3979,253.5035) and (261.8136,254.0879) ..(261.8136,254.8087) .. controls (261.8136,255.5295) and (262.3979,256.1138) ..(263.1187,256.1138) .. controls (263.8395,256.1138) and (264.4239,255.5295) ..(264.4239,254.8087);
\path[fill=chfourfig5c0,fill opacity=0.5,draw=chfourfig5c0,line width=0.3985pt,draw opacity=0.5] (268.3093,248.6118) .. controls (268.3093,247.8910) and (267.7250,247.3067) ..(267.0042,247.3067) .. controls (266.2834,247.3067) and (265.6991,247.8910) ..(265.6991,248.6118) .. controls (265.6991,249.3326) and (266.2834,249.9170) ..(267.0042,249.9170) .. controls (267.7250,249.9170) and (268.3093,249.3326) ..(268.3093,248.6118);
\path[fill=chfourfig5c0,fill opacity=0.5,draw=chfourfig5c0,line width=0.3985pt,draw opacity=0.5] (323.3038,234.9330) .. controls (323.3038,234.2122) and (322.7195,233.6278) ..(321.9987,233.6278) .. controls (321.2779,233.6278) and (320.6936,234.2122) ..(320.6936,234.9330) .. controls (320.6936,235.6537) and (321.2779,236.2381) ..(321.9987,236.2381) .. controls (322.7195,236.2381) and (323.3038,235.6537) ..(323.3038,234.9330);
\path[fill=chfourfig5c0,fill opacity=0.5,draw=chfourfig5c0,line width=0.3985pt,draw opacity=0.5] (225.7484,230.4198) .. controls (225.7484,229.6990) and (225.1641,229.1147) ..(224.4433,229.1147) .. controls (223.7225,229.1147) and (223.1381,229.6990) ..(223.1381,230.4198) .. controls (223.1381,231.1406) and (223.7225,231.7249) ..(224.4433,231.7249) .. controls (225.1641,231.7249) and (225.7484,231.1406) ..(225.7484,230.4198);
\path[fill=chfourfig5c0,fill opacity=0.5,draw=chfourfig5c0,line width=0.3985pt,draw opacity=0.5] (308.8578,246.5695) .. controls (308.8578,245.8487) and (308.2735,245.2643) ..(307.5527,245.2643) .. controls (306.8319,245.2643) and (306.2476,245.8487) ..(306.2476,246.5695) .. controls (306.2476,247.2903) and (306.8319,247.8746) ..(307.5527,247.8746) .. controls (308.2735,247.8746) and (308.8578,247.2903) ..(308.8578,246.5695);
\path[fill=chfourfig5c0,fill opacity=0.5,draw=chfourfig5c0,line width=0.3985pt,draw opacity=0.5] (257.2009,212.0983) .. controls (257.2009,211.3775) and (256.6165,210.7932) ..(255.8957,210.7932) .. controls (255.1749,210.7932) and (254.5906,211.3775) ..(254.5906,212.0983) .. controls (254.5906,212.8191) and (255.1749,213.4034) ..(255.8957,213.4034) .. controls (256.6165,213.4034) and (257.2009,212.8191) ..(257.2009,212.0983);
\path[fill=chfourfig5c0,fill opacity=0.5,draw=chfourfig5c0,line width=0.3985pt,draw opacity=0.5] (252.4187,197.4829) .. controls (252.4187,196.7621) and (251.8344,196.1778) ..(251.1136,196.1778) .. controls (250.3928,196.1778) and (249.8085,196.7621) ..(249.8085,197.4829) .. controls (249.8085,198.2037) and (250.3928,198.7881) ..(251.1136,198.7881) .. controls (251.8344,198.7881) and (252.4187,198.2037) ..(252.4187,197.4829);
\path[fill=chfourfig5c0,fill opacity=0.5,draw=chfourfig5c0,line width=0.3985pt,draw opacity=0.5] (269.5945,237.7026) .. controls (269.5945,236.9818) and (269.0102,236.3975) ..(268.2894,236.3975) .. controls (267.5686,236.3975) and (266.9843,236.9818) ..(266.9843,237.7026) .. controls (266.9843,238.4234) and (267.5686,239.0078) ..(268.2894,239.0078) .. controls (269.0102,239.0078) and (269.5945,238.4234) ..(269.5945,237.7026);
\path[fill=chfourfig5c0,fill opacity=0.5,draw=chfourfig5c0,line width=0.3985pt,draw opacity=0.5] (282.7653,189.0345) .. controls (282.7653,188.3137) and (282.1810,187.7294) ..(281.4602,187.7294) .. controls (280.7394,187.7294) and (280.1551,188.3137) ..(280.1551,189.0345) .. controls (280.1551,189.7553) and (280.7394,190.3396) ..(281.4602,190.3396) .. controls (282.1810,190.3396) and (282.7653,189.7553) ..(282.7653,189.0345);
\path[fill=chfourfig5c0,fill opacity=0.5,draw=chfourfig5c0,line width=0.3985pt,draw opacity=0.5] (284.3494,194.7232) .. controls (284.3494,194.0024) and (283.7651,193.4181) ..(283.0443,193.4181) .. controls (282.3235,193.4181) and (281.7391,194.0024) ..(281.7391,194.7232) .. controls (281.7391,195.4440) and (282.3235,196.0284) ..(283.0443,196.0284) .. controls (283.7651,196.0284) and (284.3494,195.4440) ..(284.3494,194.7232);
\path[fill=chfourfig5c0,fill opacity=0.5,draw=chfourfig5c0,line width=0.3985pt,draw opacity=0.5] (305.6896,201.3086) .. controls (305.6896,200.5878) and (305.1053,200.0035) ..(304.3845,200.0035) .. controls (303.6637,200.0035) and (303.0794,200.5878) ..(303.0794,201.3086) .. controls (303.0794,202.0294) and (303.6637,202.6137) ..(304.3845,202.6137) .. controls (305.1053,202.6137) and (305.6896,202.0294) ..(305.6896,201.3086);
\path[fill=chfourfig5c0,fill opacity=0.5,draw=chfourfig5c0,line width=0.3985pt,draw opacity=0.5] (288.2249,221.3438) .. controls (288.2249,220.6230) and (287.6406,220.0386) ..(286.9198,220.0386) .. controls (286.1990,220.0386) and (285.6147,220.6230) ..(285.6147,221.3438) .. controls (285.6147,222.0645) and (286.1990,222.6489) ..(286.9198,222.6489) .. controls (287.6406,222.6489) and (288.2249,222.0645) ..(288.2249,221.3438);
\path[fill=chfourfig5c0,fill opacity=0.5,draw=chfourfig5c0,line width=0.3985pt,draw opacity=0.5] (251.9106,193.9760) .. controls (251.9106,193.2552) and (251.3263,192.6709) ..(250.6055,192.6709) .. controls (249.8847,192.6709) and (249.3004,193.2552) ..(249.3004,193.9760) .. controls (249.3004,194.6968) and (249.8847,195.2811) ..(250.6055,195.2811) .. controls (251.3263,195.2811) and (251.9106,194.6968) ..(251.9106,193.9760);
\path[fill=chfourfig5c0,fill opacity=0.5,draw=chfourfig5c0,line width=0.3985pt,draw opacity=0.5] (271.9059,213.4931) .. controls (271.9059,212.7723) and (271.3216,212.1880) ..(270.6008,212.1880) .. controls (269.8800,212.1880) and (269.2957,212.7723) ..(269.2957,213.4931) .. controls (269.2957,214.2139) and (269.8800,214.7982) ..(270.6008,214.7982) .. controls (271.3216,214.7982) and (271.9059,214.2139) ..(271.9059,213.4931);
\path[fill=chfourfig5c0,fill opacity=0.5,draw=chfourfig5c0,line width=0.3985pt,draw opacity=0.5] (256.3540,215.8344) .. controls (256.3540,215.1136) and (255.7697,214.5292) ..(255.0489,214.5292) .. controls (254.3281,214.5292) and (253.7438,215.1136) ..(253.7438,215.8344) .. controls (253.7438,216.5551) and (254.3281,217.1395) ..(255.0489,217.1395) .. controls (255.7697,217.1395) and (256.3540,216.5551) ..(256.3540,215.8344);
\path[fill=chfourfig5c0,fill opacity=0.5,draw=chfourfig5c0,line width=0.3985pt,draw opacity=0.5] (256.8621,206.9774) .. controls (256.8621,206.2567) and (256.2778,205.6723) ..(255.5570,205.6723) .. controls (254.8362,205.6723) and (254.2519,206.2567) ..(254.2519,206.9774) .. controls (254.2519,207.6982) and (254.8362,208.2825) ..(255.5570,208.2825) .. controls (256.2778,208.2825) and (256.8621,207.6982) ..(256.8621,206.9774);
\path[fill=chfourfig5c0,fill opacity=0.5,draw=chfourfig5c0,line width=0.3985pt,draw opacity=0.5] (281.5001,242.3054) .. controls (281.5001,241.5846) and (280.9157,241.0003) ..(280.1949,241.0003) .. controls (279.4741,241.0003) and (278.8898,241.5846) ..(278.8898,242.3054) .. controls (278.8898,243.0262) and (279.4741,243.6105) ..(280.1949,243.6105) .. controls (280.9157,243.6105) and (281.5001,243.0262) ..(281.5001,242.3054);
\path[fill=chfourfig5c0,fill opacity=0.5,draw=chfourfig5c0,line width=0.3985pt,draw opacity=0.5] (316.4893,201.5577) .. controls (316.4893,200.8369) and (315.9050,200.2526) ..(315.1841,200.2526) .. controls (314.4634,200.2526) and (313.8791,200.8369) ..(313.8791,201.5577) .. controls (313.8791,202.2785) and (314.4634,202.8628) ..(315.1841,202.8628) .. controls (315.9050,202.8628) and (316.4893,202.2785) ..(316.4893,201.5577);
\path[fill=chfourfig5c0,fill opacity=0.5,draw=chfourfig5c0,line width=0.3985pt,draw opacity=0.5] (321.7397,206.1605) .. controls (321.7397,205.4397) and (321.1553,204.8553) ..(320.4345,204.8553) .. controls (319.7137,204.8553) and (319.1294,205.4397) ..(319.1294,206.1605) .. controls (319.1294,206.8813) and (319.7137,207.4656) ..(320.4345,207.4656) .. controls (321.1553,207.4656) and (321.7397,206.8813) ..(321.7397,206.1605);
\path[fill=chfourfig5c0,fill opacity=0.5,draw=chfourfig5c0,line width=0.3985pt,draw opacity=0.5] (265.7090,246.8982) .. controls (265.7090,246.1774) and (265.1247,245.5931) ..(264.4039,245.5931) .. controls (263.6831,245.5931) and (263.0988,246.1774) ..(263.0988,246.8982) .. controls (263.0988,247.6190) and (263.6831,248.2034) ..(264.4039,248.2034) .. controls (265.1247,248.2034) and (265.7090,247.6190) ..(265.7090,246.8982);
\path[fill=chfourfig5c0,fill opacity=0.5,draw=chfourfig5c0,line width=0.3985pt,draw opacity=0.5] (274.7353,239.6852) .. controls (274.7353,238.9644) and (274.1510,238.3801) ..(273.4302,238.3801) .. controls (272.7094,238.3801) and (272.1251,238.9644) ..(272.1251,239.6852) .. controls (272.1251,240.4060) and (272.7094,240.9903) ..(273.4302,240.9903) .. controls (274.1510,240.9903) and (274.7353,240.4060) ..(274.7353,239.6852);
\path[fill=chfourfig5c0,fill opacity=0.5,draw=chfourfig5c0,line width=0.3985pt,draw opacity=0.5] (251.2332,246.9082) .. controls (251.2332,246.1874) and (250.6488,245.6031) ..(249.9280,245.6031) .. controls (249.2072,245.6031) and (248.6229,246.1874) ..(248.6229,246.9082) .. controls (248.6229,247.6290) and (249.2072,248.2133) ..(249.9280,248.2133) .. controls (250.6488,248.2133) and (251.2332,247.6290) ..(251.2332,246.9082);
\path[fill=chfourfig5c0,fill opacity=0.5,draw=chfourfig5c0,line width=0.3985pt,draw opacity=0.5] (256.2145,245.6629) .. controls (256.2145,244.9421) and (255.6302,244.3578) ..(254.9094,244.3578) .. controls (254.1886,244.3578) and (253.6043,244.9421) ..(253.6043,245.6629) .. controls (253.6043,246.3837) and (254.1886,246.9680) ..(254.9094,246.9680) .. controls (255.6302,246.9680) and (256.2145,246.3837) ..(256.2145,245.6629);
\path[fill=chfourfig5c0,fill opacity=0.5,draw=chfourfig5c0,line width=0.3985pt,draw opacity=0.5] (291.3333,227.5605) .. controls (291.3333,226.8397) and (290.7490,226.2554) ..(290.0282,226.2554) .. controls (289.3074,226.2554) and (288.7231,226.8397) ..(288.7231,227.5605) .. controls (288.7231,228.2813) and (289.3074,228.8656) ..(290.0282,228.8656) .. controls (290.7490,228.8656) and (291.3333,228.2813) ..(291.3333,227.5605);
\path[fill=chfourfig5c0,fill opacity=0.5,draw=chfourfig5c0,line width=0.3985pt,draw opacity=0.5] (290.8750,258.8636) .. controls (290.8750,258.1427) and (290.2907,257.5584) ..(289.5699,257.5584) .. controls (288.8491,257.5584) and (288.2648,258.1427) ..(288.2648,258.8636) .. controls (288.2648,259.5844) and (288.8491,260.1686) ..(289.5699,260.1686) .. controls (290.2907,260.1686) and (290.8750,259.5844) ..(290.8750,258.8636);
\path[fill=chfourfig5c0,fill opacity=0.5,draw=chfourfig5c0,line width=0.3985pt,draw opacity=0.5] (266.5758,249.8970) .. controls (266.5758,249.1762) and (265.9915,248.5919) ..(265.2707,248.5919) .. controls (264.5499,248.5919) and (263.9656,249.1762) ..(263.9656,249.8970) .. controls (263.9656,250.6179) and (264.5499,251.2021) ..(265.2707,251.2021) .. controls (265.9915,251.2021) and (266.5758,250.6179) ..(266.5758,249.8970);
\path[fill=chfourfig5c0,fill opacity=0.5,draw=chfourfig5c0,line width=0.3985pt,draw opacity=0.5] (329.6601,213.5429) .. controls (329.6601,212.8221) and (329.0757,212.2378) ..(328.3549,212.2378) .. controls (327.6341,212.2378) and (327.0498,212.8221) ..(327.0498,213.5429) .. controls (327.0498,214.2637) and (327.6341,214.8480) ..(328.3549,214.8480) .. controls (329.0757,214.8480) and (329.6601,214.2637) ..(329.6601,213.5429);
\path[fill=chfourfig5c0,fill opacity=0.5,draw=chfourfig5c0,line width=0.3985pt,draw opacity=0.5] (286.3220,222.6489) .. controls (286.3220,221.9281) and (285.7377,221.3438) ..(285.0169,221.3438) .. controls (284.2961,221.3438) and (283.7118,221.9281) ..(283.7118,222.6489) .. controls (283.7118,223.3697) and (284.2961,223.9540) ..(285.0169,223.9540) .. controls (285.7377,223.9540) and (286.3220,223.3697) ..(286.3220,222.6489);
\path[fill=chfourfig5c0,fill opacity=0.5,draw=chfourfig5c0,line width=0.3985pt,draw opacity=0.5] (289.8090,203.3211) .. controls (289.8090,202.6003) and (289.2247,202.0160) ..(288.5038,202.0160) .. controls (287.7831,202.0160) and (287.1988,202.6003) ..(287.1988,203.3211) .. controls (287.1988,204.0419) and (287.7831,204.6262) ..(288.5038,204.6262) .. controls (289.2247,204.6262) and (289.8090,204.0419) ..(289.8090,203.3211);
\path[fill=chfourfig5c0,fill opacity=0.5,draw=chfourfig5c0,line width=0.3985pt,draw opacity=0.5] (269.1363,231.5257) .. controls (269.1363,230.8049) and (268.5519,230.2206) ..(267.8311,230.2206) .. controls (267.1103,230.2206) and (266.5260,230.8049) ..(266.5260,231.5257) .. controls (266.5260,232.2465) and (267.1103,232.8308) ..(267.8311,232.8308) .. controls (268.5519,232.8308) and (269.1363,232.2465) ..(269.1363,231.5257);
\path[fill=chfourfig5c0,fill opacity=0.5,draw=chfourfig5c0,line width=0.3985pt,draw opacity=0.5] (304.0657,239.0974) .. controls (304.0657,238.3766) and (303.4814,237.7923) ..(302.7606,237.7923) .. controls (302.0398,237.7923) and (301.4554,238.3766) ..(301.4554,239.0974) .. controls (301.4554,239.8182) and (302.0398,240.4025) ..(302.7606,240.4025) .. controls (303.4814,240.4025) and (304.0657,239.8182) ..(304.0657,239.0974);
\path[fill=chfourfig5c0,fill opacity=0.5,draw=chfourfig5c0,line width=0.3985pt,draw opacity=0.5] (301.0370,211.4507) .. controls (301.0370,210.7299) and (300.4527,210.1456) ..(299.7319,210.1456) .. controls (299.0111,210.1456) and (298.4268,210.7299) ..(298.4268,211.4507) .. controls (298.4268,212.1715) and (299.0111,212.7558) ..(299.7319,212.7558) .. controls (300.4527,212.7558) and (301.0370,212.1715) ..(301.0370,211.4507);
\path[fill=chfourfig5c0,fill opacity=0.5,draw=chfourfig5c0,line width=0.3985pt,draw opacity=0.5] (289.5101,229.5431) .. controls (289.5101,228.8223) and (288.9258,228.2380) ..(288.2050,228.2380) .. controls (287.4842,228.2380) and (286.8999,228.8223) ..(286.8999,229.5431) .. controls (286.8999,230.2639) and (287.4842,230.8482) ..(288.2050,230.8482) .. controls (288.9258,230.8482) and (289.5101,230.2639) ..(289.5101,229.5431);
\path[fill=chfourfig5c0,fill opacity=0.5,draw=chfourfig5c0,line width=0.3985pt,draw opacity=0.5] (258.0576,221.0050) .. controls (258.0576,220.2842) and (257.4733,219.6999) ..(256.7525,219.6999) .. controls (256.0317,219.6999) and (255.4474,220.2842) ..(255.4474,221.0050) .. controls (255.4474,221.7258) and (256.0317,222.3101) ..(256.7525,222.3101) .. controls (257.4733,222.3101) and (258.0576,221.7258) ..(258.0576,221.0050);
\path[fill=chfourfig5c0,fill opacity=0.5,draw=chfourfig5c0,line width=0.3985pt,draw opacity=0.5] (287.5175,252.6667) .. controls (287.5175,251.9459) and (286.9332,251.3616) ..(286.2124,251.3616) .. controls (285.4916,251.3616) and (284.9073,251.9459) ..(284.9073,252.6667) .. controls (284.9073,253.3875) and (285.4916,253.9718) ..(286.2124,253.9718) .. controls (286.9332,253.9718) and (287.5175,253.3875) ..(287.5175,252.6667);
\path[fill=chfourfig5c0,fill opacity=0.5,draw=chfourfig5c0,line width=0.3985pt,draw opacity=0.5] (286.7604,180.8750) .. controls (286.7604,180.1542) and (286.1761,179.5699) ..(285.4553,179.5699) .. controls (284.7345,179.5699) and (284.1501,180.1542) ..(284.1501,180.8750) .. controls (284.1501,181.5958) and (284.7345,182.1801) ..(285.4553,182.1801) .. controls (286.1761,182.1801) and (286.7604,181.5958) ..(286.7604,180.8750);
\path[fill=chfourfig5c0,fill opacity=0.5,draw=chfourfig5c0,line width=0.3985pt,draw opacity=0.5] (295.4081,252.7862) .. controls (295.4081,252.0654) and (294.8238,251.4811) ..(294.1029,251.4811) .. controls (293.3821,251.4811) and (292.7978,252.0654) ..(292.7978,252.7862) .. controls (292.7978,253.5070) and (293.3821,254.0914) ..(294.1029,254.0914) .. controls (294.8238,254.0914) and (295.4081,253.5070) ..(295.4081,252.7862);
\path[fill=chfourfig5c0,fill opacity=0.5,draw=chfourfig5c0,line width=0.3985pt,draw opacity=0.5] (311.5079,213.4233) .. controls (311.5079,212.7025) and (310.9236,212.1182) ..(310.2028,212.1182) .. controls (309.4820,212.1182) and (308.8976,212.7025) ..(308.8976,213.4233) .. controls (308.8976,214.1441) and (309.4820,214.7285) ..(310.2028,214.7285) .. controls (310.9236,214.7285) and (311.5079,214.1441) ..(311.5079,213.4233);
\path[fill=chfourfig5c0,fill opacity=0.5,draw=chfourfig5c0,line width=0.3985pt,draw opacity=0.5] (252.6180,223.8245) .. controls (252.6180,223.1037) and (252.0336,222.5193) ..(251.3129,222.5193) .. controls (250.5920,222.5193) and (250.0078,223.1037) ..(250.0078,223.8245) .. controls (250.0078,224.5453) and (250.5920,225.1296) ..(251.3129,225.1296) .. controls (252.0336,225.1296) and (252.6180,224.5453) ..(252.6180,223.8245);
\path[fill=chfourfig5c0,fill opacity=0.5,draw=chfourfig5c0,line width=0.3985pt,draw opacity=0.5] (297.5202,213.7023) .. controls (297.5202,212.9815) and (296.9359,212.3972) ..(296.2151,212.3972) .. controls (295.4942,212.3972) and (294.9099,212.9815) ..(294.9099,213.7023) .. controls (294.9099,214.4231) and (295.4942,215.0074) ..(296.2151,215.0074) .. controls (296.9359,215.0074) and (297.5202,214.4231) ..(297.5202,213.7023);
\path[fill=chfourfig5c0,fill opacity=0.5,draw=chfourfig5c0,line width=0.3985pt,draw opacity=0.5] (277.9234,232.6415) .. controls (277.9234,231.9207) and (277.3391,231.3364) ..(276.6183,231.3364) .. controls (275.8975,231.3364) and (275.3132,231.9207) ..(275.3132,232.6415) .. controls (275.3132,233.3623) and (275.8975,233.9467) ..(276.6183,233.9467) .. controls (277.3391,233.9467) and (277.9234,233.3623) ..(277.9234,232.6415);
\path[fill=chfourfig5c0,fill opacity=0.5,draw=chfourfig5c0,line width=0.3985pt,draw opacity=0.5] (316.2801,227.9491) .. controls (316.2801,227.2283) and (315.6958,226.6440) ..(314.9749,226.6440) .. controls (314.2542,226.6440) and (313.6698,227.2283) ..(313.6698,227.9491) .. controls (313.6698,228.6699) and (314.2542,229.2542) ..(314.9749,229.2542) .. controls (315.6958,229.2542) and (316.2801,228.6699) ..(316.2801,227.9491);
\path[fill=chfourfig5c0,fill opacity=0.5,draw=chfourfig5c0,line width=0.3985pt,draw opacity=0.5] (261.0764,241.5482) .. controls (261.0764,240.8275) and (260.4920,240.2431) ..(259.7712,240.2431) .. controls (259.0504,240.2431) and (258.4661,240.8275) ..(258.4661,241.5482) .. controls (258.4661,242.2690) and (259.0504,242.8534) ..(259.7712,242.8534) .. controls (260.4920,242.8534) and (261.0764,242.2690) ..(261.0764,241.5482);
\path[fill=chfourfig5c0,fill opacity=0.5,draw=chfourfig5c0,line width=0.3985pt,draw opacity=0.5] (300.9573,235.2318) .. controls (300.9573,234.5110) and (300.3730,233.9267) ..(299.6522,233.9267) .. controls (298.9314,233.9267) and (298.3471,234.5110) ..(298.3471,235.2318) .. controls (298.3471,235.9526) and (298.9314,236.5370) ..(299.6522,236.5370) .. controls (300.3730,236.5370) and (300.9573,235.9526) ..(300.9573,235.2318);
\path[fill=chfourfig5c0,fill opacity=0.5,draw=chfourfig5c0,line width=0.3985pt,draw opacity=0.5] (259.7414,227.1819) .. controls (259.7414,226.4611) and (259.1570,225.8768) ..(258.4362,225.8768) .. controls (257.7155,225.8768) and (257.1311,226.4611) ..(257.1311,227.1819) .. controls (257.1311,227.9027) and (257.7155,228.4871) ..(258.4362,228.4871) .. controls (259.1570,228.4871) and (259.7414,227.9027) ..(259.7414,227.1819);
\path[fill=chfourfig5c0,fill opacity=0.5,draw=chfourfig5c0,line width=0.3985pt,draw opacity=0.5] (283.2336,238.6889) .. controls (283.2336,237.9681) and (282.6492,237.3838) ..(281.9284,237.3838) .. controls (281.2076,237.3838) and (280.6233,237.9681) ..(280.6233,238.6889) .. controls (280.6233,239.4097) and (281.2076,239.9940) ..(281.9284,239.9940) .. controls (282.6492,239.9940) and (283.2336,239.4097) ..(283.2336,238.6889);
\path[fill=chfourfig5c0,fill opacity=0.5,draw=chfourfig5c0,line width=0.3985pt,draw opacity=0.5] (316.5690,197.7220) .. controls (316.5690,197.0012) and (315.9847,196.4169) ..(315.2639,196.4169) .. controls (314.5431,196.4169) and (313.9587,197.0012) ..(313.9587,197.7220) .. controls (313.9587,198.4428) and (314.5431,199.0271) ..(315.2639,199.0271) .. controls (315.9847,199.0271) and (316.5690,198.4428) ..(316.5690,197.7220);
\path[fill=chfourfig5c0,fill opacity=0.5,draw=chfourfig5c0,line width=0.3985pt,draw opacity=0.5] (254.6803,205.0646) .. controls (254.6803,204.3438) and (254.0959,203.7595) ..(253.3752,203.7595) .. controls (252.6544,203.7595) and (252.0700,204.3438) ..(252.0700,205.0646) .. controls (252.0700,205.7854) and (252.6544,206.3697) ..(253.3752,206.3697) .. controls (254.0959,206.3697) and (254.6803,205.7854) ..(254.6803,205.0646);
\path[fill=chfourfig5c0,fill opacity=0.5,draw=chfourfig5c0,line width=0.3985pt,draw opacity=0.5] (281.2310,204.9151) .. controls (281.2310,204.1943) and (280.6467,203.6100) ..(279.9259,203.6100) .. controls (279.2051,203.6100) and (278.6208,204.1943) ..(278.6208,204.9151) .. controls (278.6208,205.6360) and (279.2051,206.2203) ..(279.9259,206.2203) .. controls (280.6467,206.2203) and (281.2310,205.6360) ..(281.2310,204.9151);
\path[fill=chfourfig5c0,fill opacity=0.5,draw=chfourfig5c0,line width=0.3985pt,draw opacity=0.5] (223.5566,225.4584) .. controls (223.5566,224.7375) and (222.9723,224.1533) ..(222.2515,224.1533) .. controls (221.5306,224.1533) and (220.9464,224.7375) ..(220.9464,225.4584) .. controls (220.9464,226.1792) and (221.5306,226.7635) ..(222.2515,226.7635) .. controls (222.9723,226.7635) and (223.5566,226.1792) ..(223.5566,225.4584);
\path[fill=chfourfig5c0,fill opacity=0.5,draw=chfourfig5c0,line width=0.3985pt,draw opacity=0.5] (262.6804,248.6617) .. controls (262.6804,247.9409) and (262.0961,247.3565) ..(261.3752,247.3565) .. controls (260.6544,247.3565) and (260.0701,247.9409) ..(260.0701,248.6617) .. controls (260.0701,249.3825) and (260.6544,249.9668) ..(261.3752,249.9668) .. controls (262.0961,249.9668) and (262.6804,249.3825) ..(262.6804,248.6617);
\path[fill=chfourfig5c0,fill opacity=0.5,draw=chfourfig5c0,line width=0.3985pt,draw opacity=0.5] (271.9955,256.2234) .. controls (271.9955,255.5026) and (271.4112,254.9183) ..(270.6904,254.9183) .. controls (269.9696,254.9183) and (269.3853,255.5026) ..(269.3853,256.2234) .. controls (269.3853,256.9442) and (269.9696,257.5285) ..(270.6904,257.5285) .. controls (271.4112,257.5285) and (271.9955,256.9442) ..(271.9955,256.2234);
\path[fill=chfourfig5c0,fill opacity=0.5,draw=chfourfig5c0,line width=0.3985pt,draw opacity=0.5] (282.6757,209.3984) .. controls (282.6757,208.6776) and (282.0913,208.0933) ..(281.3705,208.0933) .. controls (280.6497,208.0933) and (280.0654,208.6776) ..(280.0654,209.3984) .. controls (280.0654,210.1192) and (280.6497,210.7035) ..(281.3705,210.7035) .. controls (282.0913,210.7035) and (282.6757,210.1192) ..(282.6757,209.3984);
\path[fill=chfourfig5c0,fill opacity=0.5,draw=chfourfig5c0,line width=0.3985pt,draw opacity=0.5] (235.8805,219.5504) .. controls (235.8805,218.8297) and (235.2962,218.2453) ..(234.5754,218.2453) .. controls (233.8546,218.2453) and (233.2703,218.8297) ..(233.2703,219.5504) .. controls (233.2703,220.2712) and (233.8546,220.8556) ..(234.5754,220.8556) .. controls (235.2962,220.8556) and (235.8805,220.2712) ..(235.8805,219.5504);
\path[fill=chfourfig5c0,fill opacity=0.5,draw=chfourfig5c0,line width=0.3985pt,draw opacity=0.5] (309.4556,252.6866) .. controls (309.4556,251.9658) and (308.8712,251.3815) ..(308.1505,251.3815) .. controls (307.4296,251.3815) and (306.8453,251.9658) ..(306.8453,252.6866) .. controls (306.8453,253.4074) and (307.4296,253.9918) ..(308.1505,253.9918) .. controls (308.8712,253.9918) and (309.4556,253.4074) ..(309.4556,252.6866);
\path[fill=none,draw=chfourfig5c0,line width=0.3985pt,draw opacity=1.0] (224.1743,271.4365) --(336.7635,271.4365);
\path[fill=none,draw=chfourfig5c0,line width=0.3985pt,draw opacity=1.0] (224.1743,271.4365) --(224.1743,276.3980);
\path[fill=none,draw=chfourfig5c0,line width=0.3985pt,draw opacity=1.0] (280.4739,271.4365) --(280.4739,276.3980);
\path[fill=none,draw=chfourfig5c0,line width=0.3985pt,draw opacity=1.0] (336.7635,271.4365) --(336.7635,276.3980);
\path[fill=none,draw=chfourfig5c0,line width=0.3985pt,draw opacity=1.0] (215.9749,270.4403) --(215.9749,190.2898);
\path[fill=none,draw=chfourfig5c0,line width=0.3985pt,draw opacity=1.0] (215.9749,270.4403) --(211.0135,270.4403);
\path[fill=none,draw=chfourfig5c0,line width=0.3985pt,draw opacity=1.0] (215.9749,230.3600) --(211.0135,230.3600);
\path[fill=none,draw=chfourfig5c0,line width=0.3985pt,draw opacity=1.0] (215.9749,190.2898) --(211.0135,190.2898);
\path[fill=none,draw=chfourfig5c0,line width=0.3985pt,draw opacity=1.0] (215.9749,271.4365)--(215.9749,167.4153)--(346.9056,167.4153)--(346.9056,271.4365)--cycle;
\node[font=\small,inner sep=0pt,rotate=0] (ch04f5t0) at (50.3832,136.2954) {$-1$};
\node[font=\small,inner sep=0pt,rotate=0] (ch04f5t1) at (86.0396,136.2954) {$0$};
\node[font=\small,inner sep=0pt,rotate=0] (ch04f5t2) at (121.6956,136.2954) {$1$};
\node[font=\small,inner sep=0pt,rotate=0] (ch04f5t3) at (157.3516,136.2954) {$2$};
\node[font=\small,inner sep=0pt,rotate=0] (ch04f5t4) at (24.1534,119.7748) {$-1$};
\node[font=\small,inner sep=0pt,rotate=0] (ch04f5t5) at (24.1534,79.6953) {$0$};
\node[font=\small,inner sep=0pt,rotate=0] (ch04f5t6) at (24.1534,39.6253) {$1$};
\node[font=\small,inner sep=0pt,rotate=0] (ch04f5t7) at (103.4472,149.0384) {$X$};
\node[font=\small,inner sep=0pt,rotate=90] (ch04f5t8) at (11.1034,69.5680) {$Y$};
\node[font=\small,inner sep=0pt,rotate=0] (ch04f5t9) at (227.9572,136.2954) {$-1$};
\node[font=\small,inner sep=0pt,rotate=0] (ch04f5t10) at (263.6136,136.2954) {$0$};
\node[font=\small,inner sep=0pt,rotate=0] (ch04f5t11) at (299.2696,136.2954) {$1$};
\node[font=\small,inner sep=0pt,rotate=0] (ch04f5t12) at (334.9256,136.2954) {$2$};
\node[font=\small,inner sep=0pt,rotate=0] (ch04f5t13) at (201.7274,119.7748) {$-1$};
\node[font=\small,inner sep=0pt,rotate=0] (ch04f5t14) at (201.7274,79.6954) {$0$};
\node[font=\small,inner sep=0pt,rotate=0] (ch04f5t15) at (201.7274,39.6254) {$1$};
\node[font=\small,inner sep=0pt,rotate=0] (ch04f5t16) at (281.0112,149.0384) {$X$};
\node[font=\small,inner sep=0pt,rotate=90] (ch04f5t17) at (188.6674,69.5672) {$Y$};
\node[font=\small,inner sep=0pt,rotate=0] (ch04f5t18) at (43.2202,286.9604) {$-1$};
\node[font=\small,inner sep=0pt,rotate=0] (ch04f5t19) at (83.6986,286.9604) {$0$};
\node[font=\small,inner sep=0pt,rotate=0] (ch04f5t20) at (124.1766,286.9604) {$1$};
\node[font=\small,inner sep=0pt,rotate=0] (ch04f5t21) at (164.6546,286.9604) {$2$};
\node[font=\small,inner sep=0pt,rotate=0] (ch04f5t22) at (24.1534,256.2433) {$-0.4$};
\node[font=\small,inner sep=0pt,rotate=0] (ch04f5t23) at (24.1534,219.0329) {$0.0$};
\node[font=\small,inner sep=0pt,rotate=0] (ch04f5t24) at (24.1534,181.8129) {$0.4$};
\node[font=\small,inner sep=0pt,rotate=0] (ch04f5t25) at (103.4472,299.7034) {$X$};
\node[font=\small,inner sep=0pt,rotate=90] (ch04f5t26) at (11.1824,219.8571) {$Y$ 关于 $X$ 回归的残差};
\node[font=\small,inner sep=0pt,rotate=0] (ch04f5t27) at (224.1712,286.9604) {$-1$};
\node[font=\small,inner sep=0pt,rotate=0] (ch04f5t28) at (280.4706,286.9604) {$0$};
\node[font=\small,inner sep=0pt,rotate=0] (ch04f5t29) at (336.7596,286.9604) {$1$};
\node[font=\small,inner sep=0pt,rotate=0] (ch04f5t30) at (201.7274,270.4398) {$-1$};
\node[font=\small,inner sep=0pt,rotate=0] (ch04f5t31) at (201.7274,230.3604) {$0$};
\node[font=\small,inner sep=0pt,rotate=0] (ch04f5t32) at (201.7274,190.2904) {$1$};
\node[font=\small,inner sep=0pt,rotate=0] (ch04f5t33) at (280.8615,299.7824) {$X$ 关于 $Y$ 回归的残差};
\node[font=\small,inner sep=0pt,rotate=90] (ch04f5t34) at (188.6674,220.2322) {$Y$};
\end{tikzpicture}
\caption{给定从底层分布抽取的样本，分别在 $X\to Y$（左）和 $Y\to X$（右）两个方向上进行线性回归。上排显示拟合的函数，下排显示对应残差。只有 $X\to Y$ 方向产生独立残差，另见图4.1。}
\label{fig:ch04-5}
\end{figure}


图4.5展示了在模拟数据集上实施这一过程的结果，底层分布见图4.1。至少在连续情形下，前两步都是机器学习和统计学的标准问题，见附录A.1及A.2；但这里还有一个额外挑战：这两步相互耦合。$\widehat f_Y$ 对 $f_Y$ 的偏离，可能掩盖或者制造噪声与输入变量之间的相依关系。一般而言，任何基于估计残差的检验，都可能无法继续控制第一类错误。样本分割是一种可能的解决办法［Kpotufe et al., 2014］。此外，必须选择能够考虑高阶统计特征的独立性检验，而不能仅检验相关性。只要回归方法包含线性项和截距项，并以最小化平方误差为目标，它就会产生不相关的噪声。\footnote{用下面这个标准几何图景就很容易看清这一点：$\Cov[\cdot,\cdot]$ 在中心化、有限方差的随机变量空间上定义了一个内积。因此，当 $Y-\alpha X$ 与 $X$ 正交时，向量 $Y-\alpha X$ 的长度最小。}

在实际中，例如可以使用\textbf{Hilbert--Schmidt 独立性准则（Hilbert--Schmidt independence criterion，HSIC）}［Gretton et al., 2008］；附录A.2将对此作简要介绍。Mooij et al.［2016，定理20］利用 HSIC 的连续性证明：即使不进行样本分割，当数据量趋于无穷时，也能获得正确的 HSIC 值；不过，该结果并未对检验的 $p$ 值作出结论。

最后一步尤其值得注意，因为它涉及概率与因果性之间的联系。根据用于拒绝和接受独立性的显著性水平，可能在两个方向上都得到加性噪声模型，也可能两个方向都得不到，还可能只在一个方向上得到。如果必须作出判断，就可以直接选择拒绝独立性时 $p$ 值较大的方向，并将其推断为因果方向。

近来的研究提供了一些证据，表明这一过程在真实数据上的成功率高于随机猜测［Mooij et al., 2016］。图4.6给出了一个真实数据的散点图；\footnote{这是原因—结果数据对数据库 \url{https://webdav.tuebingen.mpg.de/cause-effect/} 中的 pair001，另见 Mooij et al.［2016］。}
对于该数据，加性噪声模型仅在因果方向上拟合得较为合理。关于离散数据所需的修改，参见相应文献［Peters et al., 2011a］。还应注意，后非线性模型 (4.9) 在实际中比更标准的非线性回归模型 (4.3) 难拟合得多。

% 中文 TikZ 重排：坐标及曲线提取自原 PDF，未重新模拟数据。
\begin{figure}[htbp]
\centering
\definecolor{chfourfig6c0}{rgb}{0.13669,0.12195,0.12529}
\definecolor{chfourfig6c1}{rgb}{0.93042,0.1109,0.14145}
\definecolor{chfourfig6c2}{rgb}{0.84295,0.84816,0.85371}
\begin{tikzpicture}[x=0.039cm,y=-0.039cm]
\useasboundingbox (0,0) rectangle (360,229);
\path[fill=chfourfig6c0,fill opacity=0.5,draw=chfourfig6c0,line width=0.3985pt,draw opacity=0.5] (57.0890,50.2038) .. controls (57.0890,49.3124) and (56.3664,48.5898) ..(55.4750,48.5898) .. controls (54.5836,48.5898) and (53.8610,49.3124) ..(53.8610,50.2038) .. controls (53.8610,51.0951) and (54.5836,51.8177) ..(55.4750,51.8177) .. controls (56.3664,51.8177) and (57.0890,51.0951) ..(57.0890,50.2038);
\path[fill=chfourfig6c0,fill opacity=0.5,draw=chfourfig6c0,line width=0.3985pt,draw opacity=0.5] (129.1895,68.3659) .. controls (129.1895,67.4745) and (128.4669,66.7519) ..(127.5755,66.7519) .. controls (126.6842,66.7519) and (125.9616,67.4745) ..(125.9616,68.3659) .. controls (125.9616,69.2572) and (126.6842,69.9798) ..(127.5755,69.9798) .. controls (128.4669,69.9798) and (129.1895,69.2572) ..(129.1895,68.3659);
\path[fill=chfourfig6c0,fill opacity=0.5,draw=chfourfig6c0,line width=0.3985pt,draw opacity=0.5] (83.9785,51.2100) .. controls (83.9785,50.3186) and (83.2559,49.5960) ..(82.3645,49.5960) .. controls (81.4731,49.5960) and (80.7505,50.3186) ..(80.7505,51.2100) .. controls (80.7505,52.1014) and (81.4731,52.8240) ..(82.3645,52.8240) .. controls (83.2559,52.8240) and (83.9785,52.1014) ..(83.9785,51.2100);
\path[fill=chfourfig6c0,fill opacity=0.5,draw=chfourfig6c0,line width=0.3985pt,draw opacity=0.5] (100.8455,70.3883) .. controls (100.8455,69.4970) and (100.1229,68.7744) ..(99.2315,68.7744) .. controls (98.3401,68.7744) and (97.6175,69.4970) ..(97.6175,70.3883) .. controls (97.6175,71.2797) and (98.3401,72.0023) ..(99.2315,72.0023) .. controls (100.1229,72.0023) and (100.8455,71.2797) ..(100.8455,70.3883);
\path[fill=chfourfig6c0,fill opacity=0.5,draw=chfourfig6c0,line width=0.3985pt,draw opacity=0.5] (73.3781,38.0890) .. controls (73.3781,37.1976) and (72.6555,36.4751) ..(71.7641,36.4751) .. controls (70.8728,36.4751) and (70.1502,37.1976) ..(70.1502,38.0890) .. controls (70.1502,38.9804) and (70.8728,39.7030) ..(71.7641,39.7030) .. controls (72.6555,39.7030) and (73.3781,38.9804) ..(73.3781,38.0890);
\path[fill=chfourfig6c0,fill opacity=0.5,draw=chfourfig6c0,line width=0.3985pt,draw opacity=0.5] (89.5676,54.2387) .. controls (89.5676,53.3473) and (88.8450,52.6247) ..(87.9536,52.6247) .. controls (87.0623,52.6247) and (86.3397,53.3473) ..(86.3397,54.2387) .. controls (86.3397,55.1300) and (87.0623,55.8526) ..(87.9536,55.8526) .. controls (88.8450,55.8526) and (89.5676,55.1300) ..(89.5676,54.2387);
\path[fill=chfourfig6c0,fill opacity=0.5,draw=chfourfig6c0,line width=0.3985pt,draw opacity=0.5] (57.8561,49.1875) .. controls (57.8561,48.2962) and (57.1335,47.5736) ..(56.2421,47.5736) .. controls (55.3508,47.5736) and (54.6282,48.2962) ..(54.6282,49.1875) .. controls (54.6282,50.0789) and (55.3508,50.8015) ..(56.2421,50.8015) .. controls (57.1335,50.8015) and (57.8561,50.0789) ..(57.8561,49.1875);
\path[fill=chfourfig6c0,fill opacity=0.5,draw=chfourfig6c0,line width=0.3985pt,draw opacity=0.5] (113.3786,66.3534) .. controls (113.3786,65.4620) and (112.6560,64.7394) ..(111.7646,64.7394) .. controls (110.8733,64.7394) and (110.1507,65.4620) ..(110.1507,66.3534) .. controls (110.1507,67.2448) and (110.8733,67.9674) ..(111.7646,67.9674) .. controls (112.6560,67.9674) and (113.3786,67.2448) ..(113.3786,66.3534);
\path[fill=chfourfig6c0,fill opacity=0.5,draw=chfourfig6c0,line width=0.3985pt,draw opacity=0.5] (92.4668,49.1875) .. controls (92.4668,48.2962) and (91.7442,47.5736) ..(90.8528,47.5736) .. controls (89.9614,47.5736) and (89.2388,48.2962) ..(89.2388,49.1875) .. controls (89.2388,50.0789) and (89.9614,50.8015) ..(90.8528,50.8015) .. controls (91.7442,50.8015) and (92.4668,50.0789) ..(92.4668,49.1875);
\path[fill=chfourfig6c0,fill opacity=0.5,draw=chfourfig6c0,line width=0.3985pt,draw opacity=0.5] (99.0223,57.2673) .. controls (99.0223,56.3760) and (98.2997,55.6534) ..(97.4083,55.6534) .. controls (96.5169,55.6534) and (95.7943,56.3760) ..(95.7943,57.2673) .. controls (95.7943,58.1587) and (96.5169,58.8813) ..(97.4083,58.8813) .. controls (98.2997,58.8813) and (99.0223,58.1587) ..(99.0223,57.2673);
\path[fill=chfourfig6c0,fill opacity=0.5,draw=chfourfig6c0,line width=0.3985pt,draw opacity=0.5] (56.7004,53.2324) .. controls (56.7004,52.3411) and (55.9778,51.6185) ..(55.0864,51.6185) .. controls (54.1951,51.6185) and (53.4725,52.3411) ..(53.4725,53.2324) .. controls (53.4725,54.1238) and (54.1951,54.8464) ..(55.0864,54.8464) .. controls (55.9778,54.8464) and (56.7004,54.1238) ..(56.7004,53.2324);
\path[fill=chfourfig6c0,fill opacity=0.5,draw=chfourfig6c0,line width=0.3985pt,draw opacity=0.5] (75.9784,48.1813) .. controls (75.9784,47.2899) and (75.2558,46.5673) ..(74.3644,46.5673) .. controls (73.4731,46.5673) and (72.7505,47.2899) ..(72.7505,48.1813) .. controls (72.7505,49.0727) and (73.4731,49.7953) ..(74.3644,49.7953) .. controls (75.2558,49.7953) and (75.9784,49.0727) ..(75.9784,48.1813);
\path[fill=chfourfig6c0,fill opacity=0.5,draw=chfourfig6c0,line width=0.3985pt,draw opacity=0.5] (68.4565,47.1751) .. controls (68.4565,46.2837) and (67.7339,45.5611) ..(66.8425,45.5611) .. controls (65.9511,45.5611) and (65.2285,46.2837) ..(65.2285,47.1751) .. controls (65.2285,48.0664) and (65.9511,48.7890) ..(66.8425,48.7890) .. controls (67.7339,48.7890) and (68.4565,48.0664) ..(68.4565,47.1751);
\path[fill=chfourfig6c0,fill opacity=0.5,draw=chfourfig6c0,line width=0.3985pt,draw opacity=0.5] (97.1891,52.2162) .. controls (97.1891,51.3249) and (96.4665,50.6023) ..(95.5751,50.6023) .. controls (94.6838,50.6023) and (93.9612,51.3249) ..(93.9612,52.2162) .. controls (93.9612,53.1076) and (94.6838,53.8302) ..(95.5751,53.8302) .. controls (96.4665,53.8302) and (97.1891,53.1076) ..(97.1891,52.2162);
\path[fill=chfourfig6c0,fill opacity=0.5,draw=chfourfig6c0,line width=0.3985pt,draw opacity=0.5] (106.6338,56.2611) .. controls (106.6338,55.3697) and (105.9112,54.6471) ..(105.0199,54.6471) .. controls (104.1285,54.6471) and (103.4059,55.3697) ..(103.4059,56.2611) .. controls (103.4059,57.1525) and (104.1285,57.8751) ..(105.0199,57.8751) .. controls (105.9112,57.8751) and (106.6338,57.1525) ..(106.6338,56.2611);
\path[fill=chfourfig6c0,fill opacity=0.5,draw=chfourfig6c0,line width=0.3985pt,draw opacity=0.5] (53.0341,48.1813) .. controls (53.0341,47.2899) and (52.3115,46.5673) ..(51.4202,46.5673) .. controls (50.5288,46.5673) and (49.8062,47.2899) ..(49.8062,48.1813) .. controls (49.8062,49.0727) and (50.5288,49.7953) ..(51.4202,49.7953) .. controls (52.3115,49.7953) and (53.0341,49.0727) ..(53.0341,48.1813);
\path[fill=chfourfig6c0,fill opacity=0.5,draw=chfourfig6c0,line width=0.3985pt,draw opacity=0.5] (112.6115,60.2960) .. controls (112.6115,59.4047) and (111.8889,58.6821) ..(110.9975,58.6821) .. controls (110.1061,58.6821) and (109.3835,59.4047) ..(109.3835,60.2960) .. controls (109.3835,61.1874) and (110.1061,61.9100) ..(110.9975,61.9100) .. controls (111.8889,61.9100) and (112.6115,61.1874) ..(112.6115,60.2960);
\path[fill=chfourfig6c0,fill opacity=0.5,draw=chfourfig6c0,line width=0.3985pt,draw opacity=0.5] (75.6895,47.1751) .. controls (75.6895,46.2837) and (74.9669,45.5611) ..(74.0755,45.5611) .. controls (73.1841,45.5611) and (72.4615,46.2837) ..(72.4615,47.1751) .. controls (72.4615,48.0664) and (73.1841,48.7890) ..(74.0755,48.7890) .. controls (74.9669,48.7890) and (75.6895,48.0664) ..(75.6895,47.1751);
\path[fill=chfourfig6c0,fill opacity=0.5,draw=chfourfig6c0,line width=0.3985pt,draw opacity=0.5] (56.5111,46.1589) .. controls (56.5111,45.2675) and (55.7885,44.5449) ..(54.8972,44.5449) .. controls (54.0058,44.5449) and (53.2832,45.2675) ..(53.2832,46.1589) .. controls (53.2832,47.0502) and (54.0058,47.7728) ..(54.8972,47.7728) .. controls (55.7885,47.7728) and (56.5111,47.0502) ..(56.5111,46.1589);
\path[fill=chfourfig6c0,fill opacity=0.5,draw=chfourfig6c0,line width=0.3985pt,draw opacity=0.5] (84.4666,55.2449) .. controls (84.4666,54.3535) and (83.7440,53.6309) ..(82.8527,53.6309) .. controls (81.9613,53.6309) and (81.2387,54.3535) ..(81.2387,55.2449) .. controls (81.2387,56.1363) and (81.9613,56.8589) ..(82.8527,56.8589) .. controls (83.7440,56.8589) and (84.4666,56.1363) ..(84.4666,55.2449);
\path[fill=chfourfig6c0,fill opacity=0.5,draw=chfourfig6c0,line width=0.3985pt,draw opacity=0.5] (128.7113,66.3534) .. controls (128.7113,65.4620) and (127.9887,64.7394) ..(127.0974,64.7394) .. controls (126.2060,64.7394) and (125.4834,65.4620) ..(125.4834,66.3534) .. controls (125.4834,67.2448) and (126.2060,67.9674) ..(127.0974,67.9674) .. controls (127.9887,67.9674) and (128.7113,67.2448) ..(128.7113,66.3534);
\path[fill=chfourfig6c0,fill opacity=0.5,draw=chfourfig6c0,line width=0.3985pt,draw opacity=0.5] (96.0334,52.2162) .. controls (96.0334,51.3249) and (95.3108,50.6023) ..(94.4195,50.6023) .. controls (93.5281,50.6023) and (92.8055,51.3249) ..(92.8055,52.2162) .. controls (92.8055,53.1076) and (93.5281,53.8302) ..(94.4195,53.8302) .. controls (95.3108,53.8302) and (96.0334,53.1076) ..(96.0334,52.2162);
\path[fill=chfourfig6c0,fill opacity=0.5,draw=chfourfig6c0,line width=0.3985pt,draw opacity=0.5] (56.5111,44.1464) .. controls (56.5111,43.2550) and (55.7885,42.5324) ..(54.8972,42.5324) .. controls (54.0058,42.5324) and (53.2832,43.2550) ..(53.2832,44.1464) .. controls (53.2832,45.0378) and (54.0058,45.7603) ..(54.8972,45.7603) .. controls (55.7885,45.7603) and (56.5111,45.0378) ..(56.5111,44.1464);
\path[fill=chfourfig6c0,fill opacity=0.5,draw=chfourfig6c0,line width=0.3985pt,draw opacity=0.5] (76.9448,34.0541) .. controls (76.9448,33.1627) and (76.2222,32.4401) ..(75.3308,32.4401) .. controls (74.4394,32.4401) and (73.7168,33.1627) ..(73.7168,34.0541) .. controls (73.7168,34.9455) and (74.4394,35.6681) ..(75.3308,35.6681) .. controls (76.2222,35.6681) and (76.9448,34.9455) ..(76.9448,34.0541);
\path[fill=chfourfig6c0,fill opacity=0.5,draw=chfourfig6c0,line width=0.3985pt,draw opacity=0.5] (95.5453,62.3085) .. controls (95.5453,61.4171) and (94.8226,60.6945) ..(93.9313,60.6945) .. controls (93.0399,60.6945) and (92.3173,61.4171) ..(92.3173,62.3085) .. controls (92.3173,63.1999) and (93.0399,63.9225) ..(93.9313,63.9225) .. controls (94.8226,63.9225) and (95.5453,63.1999) ..(95.5453,62.3085);
\path[fill=chfourfig6c0,fill opacity=0.5,draw=chfourfig6c0,line width=0.3985pt,draw opacity=0.5] (58.4339,44.1464) .. controls (58.4339,43.2550) and (57.7113,42.5324) ..(56.8200,42.5324) .. controls (55.9286,42.5324) and (55.2060,43.2550) ..(55.2060,44.1464) .. controls (55.2060,45.0378) and (55.9286,45.7603) ..(56.8200,45.7603) .. controls (57.7113,45.7603) and (58.4339,45.0378) ..(58.4339,44.1464);
\path[fill=chfourfig6c0,fill opacity=0.5,draw=chfourfig6c0,line width=0.3985pt,draw opacity=0.5] (58.2446,43.1402) .. controls (58.2446,42.2488) and (57.5220,41.5262) ..(56.6307,41.5262) .. controls (55.7393,41.5262) and (55.0167,42.2488) ..(55.0167,43.1402) .. controls (55.0167,44.0315) and (55.7393,44.7541) ..(56.6307,44.7541) .. controls (57.5220,44.7541) and (58.2446,44.0315) ..(58.2446,43.1402);
\path[fill=chfourfig6c0,fill opacity=0.5,draw=chfourfig6c0,line width=0.3985pt,draw opacity=0.5] (56.1226,39.0953) .. controls (56.1226,38.2039) and (55.4000,37.4813) ..(54.5086,37.4813) .. controls (53.6173,37.4813) and (52.8947,38.2039) ..(52.8947,39.0953) .. controls (52.8947,39.9866) and (53.6173,40.7092) ..(54.5086,40.7092) .. controls (55.4000,40.7092) and (56.1226,39.9866) ..(56.1226,39.0953);
\path[fill=chfourfig6c0,fill opacity=0.5,draw=chfourfig6c0,line width=0.3985pt,draw opacity=0.5] (57.3779,38.0890) .. controls (57.3779,37.1976) and (56.6553,36.4751) ..(55.7639,36.4751) .. controls (54.8726,36.4751) and (54.1500,37.1976) ..(54.1500,38.0890) .. controls (54.1500,38.9804) and (54.8726,39.7030) ..(55.7639,39.7030) .. controls (56.6553,39.7030) and (57.3779,38.9804) ..(57.3779,38.0890);
\path[fill=chfourfig6c0,fill opacity=0.5,draw=chfourfig6c0,line width=0.3985pt,draw opacity=0.5] (60.7453,40.1115) .. controls (60.7453,39.2201) and (60.0227,38.4975) ..(59.1313,38.4975) .. controls (58.2400,38.4975) and (57.5174,39.2201) ..(57.5174,40.1115) .. controls (57.5174,41.0028) and (58.2400,41.7254) ..(59.1313,41.7254) .. controls (60.0227,41.7254) and (60.7453,41.0028) ..(60.7453,40.1115);
\path[fill=chfourfig6c0,fill opacity=0.5,draw=chfourfig6c0,line width=0.3985pt,draw opacity=0.5] (64.2223,31.0254) .. controls (64.2223,30.1340) and (63.4997,29.4115) ..(62.6084,29.4115) .. controls (61.7170,29.4115) and (60.9944,30.1340) ..(60.9944,31.0254) .. controls (60.9944,31.9168) and (61.7170,32.6394) ..(62.6084,32.6394) .. controls (63.4997,32.6394) and (64.2223,31.9168) ..(64.2223,31.0254);
\path[fill=chfourfig6c0,fill opacity=0.5,draw=chfourfig6c0,line width=0.3985pt,draw opacity=0.5] (87.6448,44.1464) .. controls (87.6448,43.2550) and (86.9222,42.5324) ..(86.0308,42.5324) .. controls (85.1395,42.5324) and (84.4169,43.2550) ..(84.4169,44.1464) .. controls (84.4169,45.0378) and (85.1395,45.7603) ..(86.0308,45.7603) .. controls (86.9222,45.7603) and (87.6448,45.0378) ..(87.6448,44.1464);
\path[fill=chfourfig6c0,fill opacity=0.5,draw=chfourfig6c0,line width=0.3985pt,draw opacity=0.5] (80.6111,55.2449) .. controls (80.6111,54.3535) and (79.8885,53.6309) ..(78.9971,53.6309) .. controls (78.1057,53.6309) and (77.3831,54.3535) ..(77.3831,55.2449) .. controls (77.3831,56.1363) and (78.1057,56.8589) ..(78.9971,56.8589) .. controls (79.8885,56.8589) and (80.6111,56.1363) ..(80.6111,55.2449);
\path[fill=chfourfig6c0,fill opacity=0.5,draw=chfourfig6c0,line width=0.3985pt,draw opacity=0.5] (92.3671,57.2673) .. controls (92.3671,56.3760) and (91.6445,55.6534) ..(90.7532,55.6534) .. controls (89.8618,55.6534) and (89.1392,56.3760) ..(89.1392,57.2673) .. controls (89.1392,58.1587) and (89.8618,58.8813) ..(90.7532,58.8813) .. controls (91.6445,58.8813) and (92.3671,58.1587) ..(92.3671,57.2673);
\path[fill=chfourfig6c0,fill opacity=0.5,draw=chfourfig6c0,line width=0.3985pt,draw opacity=0.5] (60.3667,32.0317) .. controls (60.3667,31.1403) and (59.6441,30.4177) ..(58.7527,30.4177) .. controls (57.8614,30.4177) and (57.1388,31.1403) ..(57.1388,32.0317) .. controls (57.1388,32.9230) and (57.8614,33.6456) ..(58.7527,33.6456) .. controls (59.6441,33.6456) and (60.3667,32.9230) ..(60.3667,32.0317);
\path[fill=chfourfig6c0,fill opacity=0.5,draw=chfourfig6c0,line width=0.3985pt,draw opacity=0.5] (54.6780,38.0890) .. controls (54.6780,37.1976) and (53.9554,36.4751) ..(53.0640,36.4751) .. controls (52.1727,36.4751) and (51.4501,37.1976) ..(51.4501,38.0890) .. controls (51.4501,38.9804) and (52.1727,39.7030) ..(53.0640,39.7030) .. controls (53.9554,39.7030) and (54.6780,38.9804) ..(54.6780,38.0890);
\path[fill=chfourfig6c0,fill opacity=0.5,draw=chfourfig6c0,line width=0.3985pt,draw opacity=0.5] (56.9893,49.1875) .. controls (56.9893,48.2962) and (56.2668,47.5736) ..(55.3754,47.5736) .. controls (54.4840,47.5736) and (53.7614,48.2962) ..(53.7614,49.1875) .. controls (53.7614,50.0789) and (54.4840,50.8015) ..(55.3754,50.8015) .. controls (56.2668,50.8015) and (56.9893,50.0789) ..(56.9893,49.1875);
\path[fill=chfourfig6c0,fill opacity=0.5,draw=chfourfig6c0,line width=0.3985pt,draw opacity=0.5] (54.1001,48.1813) .. controls (54.1001,47.2899) and (53.3775,46.5673) ..(52.4862,46.5673) .. controls (51.5948,46.5673) and (50.8722,47.2899) ..(50.8722,48.1813) .. controls (50.8722,49.0727) and (51.5948,49.7953) ..(52.4862,49.7953) .. controls (53.3775,49.7953) and (54.1001,49.0727) ..(54.1001,48.1813);
\path[fill=chfourfig6c0,fill opacity=0.5,draw=chfourfig6c0,line width=0.3985pt,draw opacity=0.5] (58.6232,30.0192) .. controls (58.6232,29.1278) and (57.9006,28.4052) ..(57.0093,28.4052) .. controls (56.1179,28.4052) and (55.3953,29.1278) ..(55.3953,30.0192) .. controls (55.3953,30.9106) and (56.1179,31.6331) ..(57.0093,31.6331) .. controls (57.9006,31.6331) and (58.6232,30.9106) ..(58.6232,30.0192);
\path[fill=chfourfig6c0,fill opacity=0.5,draw=chfourfig6c0,line width=0.3985pt,draw opacity=0.5] (137.0003,64.3310) .. controls (137.0003,63.4396) and (136.2777,62.7170) ..(135.3864,62.7170) .. controls (134.4950,62.7170) and (133.7724,63.4396) ..(133.7724,64.3310) .. controls (133.7724,65.2223) and (134.4950,65.9449) ..(135.3864,65.9449) .. controls (136.2777,65.9449) and (137.0003,65.2223) ..(137.0003,64.3310);
\path[fill=chfourfig6c0,fill opacity=0.5,draw=chfourfig6c0,line width=0.3985pt,draw opacity=0.5] (53.1338,41.1177) .. controls (53.1338,40.2263) and (52.4112,39.5037) ..(51.5198,39.5037) .. controls (50.6284,39.5037) and (49.9058,40.2263) ..(49.9058,41.1177) .. controls (49.9058,42.0091) and (50.6284,42.7317) ..(51.5198,42.7317) .. controls (52.4112,42.7317) and (53.1338,42.0091) ..(53.1338,41.1177);
\path[fill=chfourfig6c0,fill opacity=0.5,draw=chfourfig6c0,line width=0.3985pt,draw opacity=0.5] (88.6111,49.1875) .. controls (88.6111,48.2962) and (87.8885,47.5736) ..(86.9972,47.5736) .. controls (86.1058,47.5736) and (85.3832,48.2962) ..(85.3832,49.1875) .. controls (85.3832,50.0789) and (86.1058,50.8015) ..(86.9972,50.8015) .. controls (87.8885,50.8015) and (88.6111,50.0789) ..(88.6111,49.1875);
\path[fill=chfourfig6c0,fill opacity=0.5,draw=chfourfig6c0,line width=0.3985pt,draw opacity=0.5] (52.7452,44.1464) .. controls (52.7452,43.2550) and (52.0226,42.5324) ..(51.1312,42.5324) .. controls (50.2399,42.5324) and (49.5173,43.2550) ..(49.5173,44.1464) .. controls (49.5173,45.0378) and (50.2399,45.7603) ..(51.1312,45.7603) .. controls (52.0226,45.7603) and (52.7452,45.0378) ..(52.7452,44.1464);
\path[fill=chfourfig6c0,fill opacity=0.5,draw=chfourfig6c0,line width=0.3985pt,draw opacity=0.5] (53.5223,51.2100) .. controls (53.5223,50.3186) and (52.7997,49.5960) ..(51.9083,49.5960) .. controls (51.0170,49.5960) and (50.2944,50.3186) ..(50.2944,51.2100) .. controls (50.2944,52.1014) and (51.0170,52.8240) ..(51.9083,52.8240) .. controls (52.7997,52.8240) and (53.5223,52.1014) ..(53.5223,51.2100);
\path[fill=chfourfig6c0,fill opacity=0.5,draw=chfourfig6c0,line width=0.3985pt,draw opacity=0.5] (111.1669,73.4170) .. controls (111.1669,72.5256) and (110.4443,71.8030) ..(109.5529,71.8030) .. controls (108.6616,71.8030) and (107.9390,72.5256) ..(107.9390,73.4170) .. controls (107.9390,74.3084) and (108.6616,75.0310) ..(109.5529,75.0310) .. controls (110.4443,75.0310) and (111.1669,74.3084) ..(111.1669,73.4170);
\path[fill=chfourfig6c0,fill opacity=0.5,draw=chfourfig6c0,line width=0.3985pt,draw opacity=0.5] (61.1339,44.1464) .. controls (61.1339,43.2550) and (60.4113,42.5324) ..(59.5199,42.5324) .. controls (58.6285,42.5324) and (57.9059,43.2550) ..(57.9059,44.1464) .. controls (57.9059,45.0378) and (58.6285,45.7603) ..(59.5199,45.7603) .. controls (60.4113,45.7603) and (61.1339,45.0378) ..(61.1339,44.1464);
\path[fill=chfourfig6c0,fill opacity=0.5,draw=chfourfig6c0,line width=0.3985pt,draw opacity=0.5] (53.1338,44.1464) .. controls (53.1338,43.2550) and (52.4112,42.5324) ..(51.5198,42.5324) .. controls (50.6284,42.5324) and (49.9058,43.2550) ..(49.9058,44.1464) .. controls (49.9058,45.0378) and (50.6284,45.7603) ..(51.5198,45.7603) .. controls (52.4112,45.7603) and (53.1338,45.0378) ..(53.1338,44.1464);
\path[fill=chfourfig6c0,fill opacity=0.5,draw=chfourfig6c0,line width=0.3985pt,draw opacity=0.5] (53.3230,42.1239) .. controls (53.3230,41.2326) and (52.6004,40.5100) ..(51.7091,40.5100) .. controls (50.8177,40.5100) and (50.0951,41.2326) ..(50.0951,42.1239) .. controls (50.0951,43.0153) and (50.8177,43.7379) ..(51.7091,43.7379) .. controls (52.6004,43.7379) and (53.3230,43.0153) ..(53.3230,42.1239);
\path[fill=chfourfig6c0,fill opacity=0.5,draw=chfourfig6c0,line width=0.3985pt,draw opacity=0.5] (162.7342,103.6938) .. controls (162.7342,102.8025) and (162.0116,102.0799) ..(161.1202,102.0799) .. controls (160.2288,102.0799) and (159.5062,102.8025) ..(159.5062,103.6938) .. controls (159.5062,104.5852) and (160.2288,105.3078) ..(161.1202,105.3078) .. controls (162.0116,105.3078) and (162.7342,104.5852) ..(162.7342,103.6938);
\path[fill=chfourfig6c0,fill opacity=0.5,draw=chfourfig6c0,line width=0.3985pt,draw opacity=0.5] (105.5778,54.2387) .. controls (105.5778,53.3473) and (104.8552,52.6247) ..(103.9638,52.6247) .. controls (103.0724,52.6247) and (102.3499,53.3473) ..(102.3499,54.2387) .. controls (102.3499,55.1300) and (103.0724,55.8526) ..(103.9638,55.8526) .. controls (104.8552,55.8526) and (105.5778,55.1300) ..(105.5778,54.2387);
\path[fill=chfourfig6c0,fill opacity=0.5,draw=chfourfig6c0,line width=0.3985pt,draw opacity=0.5] (82.5339,54.2387) .. controls (82.5339,53.3473) and (81.8113,52.6247) ..(80.9199,52.6247) .. controls (80.0285,52.6247) and (79.3059,53.3473) ..(79.3059,54.2387) .. controls (79.3059,55.1300) and (80.0285,55.8526) ..(80.9199,55.8526) .. controls (81.8113,55.8526) and (82.5339,55.1300) ..(82.5339,54.2387);
\path[fill=chfourfig6c0,fill opacity=0.5,draw=chfourfig6c0,line width=0.3985pt,draw opacity=0.5] (85.4231,51.2100) .. controls (85.4231,50.3186) and (84.7005,49.5960) ..(83.8091,49.5960) .. controls (82.9178,49.5960) and (82.1952,50.3186) ..(82.1952,51.2100) .. controls (82.1952,52.1014) and (82.9178,52.8240) ..(83.8091,52.8240) .. controls (84.7005,52.8240) and (85.4231,52.1014) ..(85.4231,51.2100);
\path[fill=chfourfig6c0,fill opacity=0.5,draw=chfourfig6c0,line width=0.3985pt,draw opacity=0.5] (60.0778,54.2387) .. controls (60.0778,53.3473) and (59.3552,52.6247) ..(58.4638,52.6247) .. controls (57.5725,52.6247) and (56.8499,53.3473) ..(56.8499,54.2387) .. controls (56.8499,55.1300) and (57.5725,55.8526) ..(58.4638,55.8526) .. controls (59.3552,55.8526) and (60.0778,55.1300) ..(60.0778,54.2387);
\path[fill=chfourfig6c0,fill opacity=0.5,draw=chfourfig6c0,line width=0.3985pt,draw opacity=0.5] (59.2110,40.1115) .. controls (59.2110,39.2201) and (58.4884,38.4975) ..(57.5971,38.4975) .. controls (56.7057,38.4975) and (55.9831,39.2201) ..(55.9831,40.1115) .. controls (55.9831,41.0028) and (56.7057,41.7254) ..(57.5971,41.7254) .. controls (58.4884,41.7254) and (59.2110,41.0028) ..(59.2110,40.1115);
\path[fill=chfourfig6c0,fill opacity=0.5,draw=chfourfig6c0,line width=0.3985pt,draw opacity=0.5] (56.4115,43.1402) .. controls (56.4115,42.2488) and (55.6889,41.5262) ..(54.7975,41.5262) .. controls (53.9062,41.5262) and (53.1836,42.2488) ..(53.1836,43.1402) .. controls (53.1836,44.0315) and (53.9062,44.7541) ..(54.7975,44.7541) .. controls (55.6889,44.7541) and (56.4115,44.0315) ..(56.4115,43.1402);
\path[fill=chfourfig6c0,fill opacity=0.5,draw=chfourfig6c0,line width=0.3985pt,draw opacity=0.5] (90.8229,53.2324) .. controls (90.8229,52.3411) and (90.1003,51.6185) ..(89.2089,51.6185) .. controls (88.3176,51.6185) and (87.5950,52.3411) ..(87.5950,53.2324) .. controls (87.5950,54.1238) and (88.3176,54.8464) ..(89.2089,54.8464) .. controls (90.1003,54.8464) and (90.8229,54.1238) ..(90.8229,53.2324);
\path[fill=chfourfig6c0,fill opacity=0.5,draw=chfourfig6c0,line width=0.3985pt,draw opacity=0.5] (92.9450,53.2324) .. controls (92.9450,52.3411) and (92.2224,51.6185) ..(91.3310,51.6185) .. controls (90.4396,51.6185) and (89.7170,52.3411) ..(89.7170,53.2324) .. controls (89.7170,54.1238) and (90.4396,54.8464) ..(91.3310,54.8464) .. controls (92.2224,54.8464) and (92.9450,54.1238) ..(92.9450,53.2324);
\path[fill=chfourfig6c0,fill opacity=0.5,draw=chfourfig6c0,line width=0.3985pt,draw opacity=0.5] (107.2117,70.3883) .. controls (107.2117,69.4970) and (106.4891,68.7744) ..(105.5977,68.7744) .. controls (104.7063,68.7744) and (103.9837,69.4970) ..(103.9837,70.3883) .. controls (103.9837,71.2797) and (104.7063,72.0023) ..(105.5977,72.0023) .. controls (106.4891,72.0023) and (107.2117,71.2797) ..(107.2117,70.3883);
\path[fill=chfourfig6c0,fill opacity=0.5,draw=chfourfig6c0,line width=0.3985pt,draw opacity=0.5] (59.3007,43.1402) .. controls (59.3007,42.2488) and (58.5781,41.5262) ..(57.6867,41.5262) .. controls (56.7954,41.5262) and (56.0728,42.2488) ..(56.0728,43.1402) .. controls (56.0728,44.0315) and (56.7954,44.7541) ..(57.6867,44.7541) .. controls (58.5781,44.7541) and (59.3007,44.0315) ..(59.3007,43.1402);
\path[fill=chfourfig6c0,fill opacity=0.5,draw=chfourfig6c0,line width=0.3985pt,draw opacity=0.5] (53.1338,43.1402) .. controls (53.1338,42.2488) and (52.4112,41.5262) ..(51.5198,41.5262) .. controls (50.6284,41.5262) and (49.9058,42.2488) ..(49.9058,43.1402) .. controls (49.9058,44.0315) and (50.6284,44.7541) ..(51.5198,44.7541) .. controls (52.4112,44.7541) and (53.1338,44.0315) ..(53.1338,43.1402);
\path[fill=chfourfig6c0,fill opacity=0.5,draw=chfourfig6c0,line width=0.3985pt,draw opacity=0.5] (65.4677,32.0317) .. controls (65.4677,31.1403) and (64.7451,30.4177) ..(63.8537,30.4177) .. controls (62.9623,30.4177) and (62.2397,31.1403) ..(62.2397,32.0317) .. controls (62.2397,32.9230) and (62.9623,33.6456) ..(63.8537,33.6456) .. controls (64.7451,33.6456) and (65.4677,32.9230) ..(65.4677,32.0317);
\path[fill=chfourfig6c0,fill opacity=0.5,draw=chfourfig6c0,line width=0.3985pt,draw opacity=0.5] (98.9226,53.2324) .. controls (98.9226,52.3411) and (98.2000,51.6185) ..(97.3087,51.6185) .. controls (96.4173,51.6185) and (95.6947,52.3411) ..(95.6947,53.2324) .. controls (95.6947,54.1238) and (96.4173,54.8464) ..(97.3087,54.8464) .. controls (98.2000,54.8464) and (98.9226,54.1238) ..(98.9226,53.2324);
\path[fill=chfourfig6c0,fill opacity=0.5,draw=chfourfig6c0,line width=0.3985pt,draw opacity=0.5] (119.7448,58.2736) .. controls (119.7448,57.3822) and (119.0222,56.6596) ..(118.1309,56.6596) .. controls (117.2395,56.6596) and (116.5169,57.3822) ..(116.5169,58.2736) .. controls (116.5169,59.1650) and (117.2395,59.8876) ..(118.1309,59.8876) .. controls (119.0222,59.8876) and (119.7448,59.1650) ..(119.7448,58.2736);
\path[fill=chfourfig6c0,fill opacity=0.5,draw=chfourfig6c0,line width=0.3985pt,draw opacity=0.5] (62.0006,47.1751) .. controls (62.0006,46.2837) and (61.2780,45.5611) ..(60.3867,45.5611) .. controls (59.4953,45.5611) and (58.7727,46.2837) ..(58.7727,47.1751) .. controls (58.7727,48.0664) and (59.4953,48.7890) ..(60.3867,48.7890) .. controls (61.2780,48.7890) and (62.0006,48.0664) ..(62.0006,47.1751);
\path[fill=chfourfig6c0,fill opacity=0.5,draw=chfourfig6c0,line width=0.3985pt,draw opacity=0.5] (117.9117,62.3085) .. controls (117.9117,61.4171) and (117.1891,60.6945) ..(116.2977,60.6945) .. controls (115.4063,60.6945) and (114.6837,61.4171) ..(114.6837,62.3085) .. controls (114.6837,63.1999) and (115.4063,63.9225) ..(116.2977,63.9225) .. controls (117.1891,63.9225) and (117.9117,63.1999) ..(117.9117,62.3085);
\path[fill=chfourfig6c0,fill opacity=0.5,draw=chfourfig6c0,line width=0.3985pt,draw opacity=0.5] (74.5338,43.1402) .. controls (74.5338,42.2488) and (73.8112,41.5262) ..(72.9198,41.5262) .. controls (72.0284,41.5262) and (71.3058,42.2488) ..(71.3058,43.1402) .. controls (71.3058,44.0315) and (72.0284,44.7541) ..(72.9198,44.7541) .. controls (73.8112,44.7541) and (74.5338,44.0315) ..(74.5338,43.1402);
\path[fill=chfourfig6c0,fill opacity=0.5,draw=chfourfig6c0,line width=0.3985pt,draw opacity=0.5] (71.1564,49.1875) .. controls (71.1564,48.2962) and (70.4338,47.5736) ..(69.5424,47.5736) .. controls (68.6511,47.5736) and (67.9285,48.2962) ..(67.9285,49.1875) .. controls (67.9285,50.0789) and (68.6511,50.8015) ..(69.5424,50.8015) .. controls (70.4338,50.8015) and (71.1564,50.0789) ..(71.1564,49.1875);
\path[fill=chfourfig6c0,fill opacity=0.5,draw=chfourfig6c0,line width=0.3985pt,draw opacity=0.5] (119.3562,63.3247) .. controls (119.3562,62.4333) and (118.6337,61.7107) ..(117.7423,61.7107) .. controls (116.8509,61.7107) and (116.1283,62.4333) ..(116.1283,63.3247) .. controls (116.1283,64.2161) and (116.8509,64.9387) ..(117.7423,64.9387) .. controls (118.6337,64.9387) and (119.3562,64.2161) ..(119.3562,63.3247);
\path[fill=chfourfig6c0,fill opacity=0.5,draw=chfourfig6c0,line width=0.3985pt,draw opacity=0.5] (56.2222,29.0030) .. controls (56.2222,28.1116) and (55.4996,27.3890) ..(54.6082,27.3890) .. controls (53.7169,27.3890) and (52.9943,28.1116) ..(52.9943,29.0030) .. controls (52.9943,29.8943) and (53.7169,30.6169) ..(54.6082,30.6169) .. controls (55.4996,30.6169) and (56.2222,29.8943) ..(56.2222,29.0030);
\path[fill=chfourfig6c0,fill opacity=0.5,draw=chfourfig6c0,line width=0.3985pt,draw opacity=0.5] (56.2222,29.0030) .. controls (56.2222,28.1116) and (55.4996,27.3890) ..(54.6082,27.3890) .. controls (53.7169,27.3890) and (52.9943,28.1116) ..(52.9943,29.0030) .. controls (52.9943,29.8943) and (53.7169,30.6169) ..(54.6082,30.6169) .. controls (55.4996,30.6169) and (56.2222,29.8943) ..(56.2222,29.0030);
\path[fill=chfourfig6c0,fill opacity=0.5,draw=chfourfig6c0,line width=0.3985pt,draw opacity=0.5] (64.5112,44.1464) .. controls (64.5112,43.2550) and (63.7886,42.5324) ..(62.8973,42.5324) .. controls (62.0059,42.5324) and (61.2833,43.2550) ..(61.2833,44.1464) .. controls (61.2833,45.0378) and (62.0059,45.7603) ..(62.8973,45.7603) .. controls (63.7886,45.7603) and (64.5112,45.0378) ..(64.5112,44.1464);
\path[fill=chfourfig6c0,fill opacity=0.5,draw=chfourfig6c0,line width=0.3985pt,draw opacity=0.5] (66.7230,47.1751) .. controls (66.7230,46.2837) and (66.0004,45.5611) ..(65.1090,45.5611) .. controls (64.2176,45.5611) and (63.4951,46.2837) ..(63.4951,47.1751) .. controls (63.4951,48.0664) and (64.2176,48.7890) ..(65.1090,48.7890) .. controls (66.0004,48.7890) and (66.7230,48.0664) ..(66.7230,47.1751);
\path[fill=chfourfig6c0,fill opacity=0.5,draw=chfourfig6c0,line width=0.3985pt,draw opacity=0.5] (96.9899,55.2449) .. controls (96.9899,54.3535) and (96.2672,53.6309) ..(95.3759,53.6309) .. controls (94.4845,53.6309) and (93.7619,54.3535) ..(93.7619,55.2449) .. controls (93.7619,56.1363) and (94.4845,56.8589) ..(95.3759,56.8589) .. controls (96.2672,56.8589) and (96.9899,56.1363) ..(96.9899,55.2449);
\path[fill=chfourfig6c0,fill opacity=0.5,draw=chfourfig6c0,line width=0.3985pt,draw opacity=0.5] (60.8449,34.0541) .. controls (60.8449,33.1627) and (60.1223,32.4401) ..(59.2310,32.4401) .. controls (58.3396,32.4401) and (57.6170,33.1627) ..(57.6170,34.0541) .. controls (57.6170,34.9455) and (58.3396,35.6681) ..(59.2310,35.6681) .. controls (60.1223,35.6681) and (60.8449,34.9455) ..(60.8449,34.0541);
\path[fill=chfourfig6c0,fill opacity=0.5,draw=chfourfig6c0,line width=0.3985pt,draw opacity=0.5] (53.8112,41.1177) .. controls (53.8112,40.2263) and (53.0886,39.5037) ..(52.1973,39.5037) .. controls (51.3059,39.5037) and (50.5833,40.2263) ..(50.5833,41.1177) .. controls (50.5833,42.0091) and (51.3059,42.7317) ..(52.1973,42.7317) .. controls (53.0886,42.7317) and (53.8112,42.0091) ..(53.8112,41.1177);
\path[fill=chfourfig6c0,fill opacity=0.5,draw=chfourfig6c0,line width=0.3985pt,draw opacity=0.5] (52.6555,44.1464) .. controls (52.6555,43.2550) and (51.9329,42.5324) ..(51.0416,42.5324) .. controls (50.1502,42.5324) and (49.4276,43.2550) ..(49.4276,44.1464) .. controls (49.4276,45.0378) and (50.1502,45.7603) ..(51.0416,45.7603) .. controls (51.9329,45.7603) and (52.6555,45.0378) ..(52.6555,44.1464);
\path[fill=chfourfig6c0,fill opacity=0.5,draw=chfourfig6c0,line width=0.3985pt,draw opacity=0.5] (66.5337,46.1589) .. controls (66.5337,45.2675) and (65.8111,44.5449) ..(64.9197,44.5449) .. controls (64.0283,44.5449) and (63.3058,45.2675) ..(63.3058,46.1589) .. controls (63.3058,47.0502) and (64.0283,47.7728) ..(64.9197,47.7728) .. controls (65.8111,47.7728) and (66.5337,47.0502) ..(66.5337,46.1589);
\path[fill=chfourfig6c0,fill opacity=0.5,draw=chfourfig6c0,line width=0.3985pt,draw opacity=0.5] (120.6116,64.3310) .. controls (120.6116,63.4396) and (119.8890,62.7170) ..(118.9976,62.7170) .. controls (118.1062,62.7170) and (117.3837,63.4396) ..(117.3837,64.3310) .. controls (117.3837,65.2223) and (118.1062,65.9449) ..(118.9976,65.9449) .. controls (119.8890,65.9449) and (120.6116,65.2223) ..(120.6116,64.3310);
\path[fill=chfourfig6c0,fill opacity=0.5,draw=chfourfig6c0,line width=0.3985pt,draw opacity=0.5] (110.4894,64.3310) .. controls (110.4894,63.4396) and (109.7668,62.7170) ..(108.8754,62.7170) .. controls (107.9841,62.7170) and (107.2615,63.4396) ..(107.2615,64.3310) .. controls (107.2615,65.2223) and (107.9841,65.9449) ..(108.8754,65.9449) .. controls (109.7668,65.9449) and (110.4894,65.2223) ..(110.4894,64.3310);
\path[fill=chfourfig6c0,fill opacity=0.5,draw=chfourfig6c0,line width=0.3985pt,draw opacity=0.5] (95.4157,54.2387) .. controls (95.4157,53.3473) and (94.6931,52.6247) ..(93.8018,52.6247) .. controls (92.9104,52.6247) and (92.1878,53.3473) ..(92.1878,54.2387) .. controls (92.1878,55.1300) and (92.9104,55.8526) ..(93.8018,55.8526) .. controls (94.6931,55.8526) and (95.4157,55.1300) ..(95.4157,54.2387);
\path[fill=chfourfig6c0,fill opacity=0.5,draw=chfourfig6c0,line width=0.3985pt,draw opacity=0.5] (54.3891,49.1875) .. controls (54.3891,48.2962) and (53.6665,47.5736) ..(52.7751,47.5736) .. controls (51.8837,47.5736) and (51.1611,48.2962) ..(51.1611,49.1875) .. controls (51.1611,50.0789) and (51.8837,50.8015) ..(52.7751,50.8015) .. controls (53.6665,50.8015) and (54.3891,50.0789) ..(54.3891,49.1875);
\path[fill=chfourfig6c0,fill opacity=0.5,draw=chfourfig6c0,line width=0.3985pt,draw opacity=0.5] (83.1117,53.2324) .. controls (83.1117,52.3411) and (82.3891,51.6185) ..(81.4977,51.6185) .. controls (80.6064,51.6185) and (79.8838,52.3411) ..(79.8838,53.2324) .. controls (79.8838,54.1238) and (80.6064,54.8464) ..(81.4977,54.8464) .. controls (82.3891,54.8464) and (83.1117,54.1238) ..(83.1117,53.2324);
\path[fill=chfourfig6c0,fill opacity=0.5,draw=chfourfig6c0,line width=0.3985pt,draw opacity=0.5] (91.9786,44.1464) .. controls (91.9786,43.2550) and (91.2560,42.5324) ..(90.3646,42.5324) .. controls (89.4732,42.5324) and (88.7506,43.2550) ..(88.7506,44.1464) .. controls (88.7506,45.0378) and (89.4732,45.7603) ..(90.3646,45.7603) .. controls (91.2560,45.7603) and (91.9786,45.0378) ..(91.9786,44.1464);
\path[fill=chfourfig6c0,fill opacity=0.5,draw=chfourfig6c0,line width=0.3985pt,draw opacity=0.5] (78.6783,46.1589) .. controls (78.6783,45.2675) and (77.9557,44.5449) ..(77.0643,44.5449) .. controls (76.1730,44.5449) and (75.4504,45.2675) ..(75.4504,46.1589) .. controls (75.4504,47.0502) and (76.1730,47.7728) ..(77.0643,47.7728) .. controls (77.9557,47.7728) and (78.6783,47.0502) ..(78.6783,46.1589);
\path[fill=chfourfig6c0,fill opacity=0.5,draw=chfourfig6c0,line width=0.3985pt,draw opacity=0.5] (75.2112,31.0254) .. controls (75.2112,30.1340) and (74.4886,29.4115) ..(73.5973,29.4115) .. controls (72.7059,29.4115) and (71.9833,30.1340) ..(71.9833,31.0254) .. controls (71.9833,31.9168) and (72.7059,32.6394) ..(73.5973,32.6394) .. controls (74.4886,32.6394) and (75.2112,31.9168) ..(75.2112,31.0254);
\path[fill=chfourfig6c0,fill opacity=0.5,draw=chfourfig6c0,line width=0.3985pt,draw opacity=0.5] (67.1115,36.0666) .. controls (67.1115,35.1752) and (66.3889,34.4526) ..(65.4976,34.4526) .. controls (64.6062,34.4526) and (63.8836,35.1752) ..(63.8836,36.0666) .. controls (63.8836,36.9579) and (64.6062,37.6805) ..(65.4976,37.6805) .. controls (66.3889,37.6805) and (67.1115,36.9579) ..(67.1115,36.0666);
\path[fill=chfourfig6c0,fill opacity=0.5,draw=chfourfig6c0,line width=0.3985pt,draw opacity=0.5] (195.9002,99.6589) .. controls (195.9002,98.7676) and (195.1776,98.0450) ..(194.2863,98.0450) .. controls (193.3949,98.0450) and (192.6723,98.7676) ..(192.6723,99.6589) .. controls (192.6723,100.5503) and (193.3949,101.2729) ..(194.2863,101.2729) .. controls (195.1776,101.2729) and (195.9002,100.5503) ..(195.9002,99.6589);
\path[fill=chfourfig6c0,fill opacity=0.5,draw=chfourfig6c0,line width=0.3985pt,draw opacity=0.5] (169.5786,103.6938) .. controls (169.5786,102.8025) and (168.8560,102.0799) ..(167.9646,102.0799) .. controls (167.0732,102.0799) and (166.3506,102.8025) ..(166.3506,103.6938) .. controls (166.3506,104.5852) and (167.0732,105.3078) ..(167.9646,105.3078) .. controls (168.8560,105.3078) and (169.5786,104.5852) ..(169.5786,103.6938);
\path[fill=chfourfig6c0,fill opacity=0.5,draw=chfourfig6c0,line width=0.3985pt,draw opacity=0.5] (56.9893,50.2038) .. controls (56.9893,49.3124) and (56.2668,48.5898) ..(55.3754,48.5898) .. controls (54.4840,48.5898) and (53.7614,49.3124) ..(53.7614,50.2038) .. controls (53.7614,51.0951) and (54.4840,51.8177) ..(55.3754,51.8177) .. controls (56.2668,51.8177) and (56.9893,51.0951) ..(56.9893,50.2038);
\path[fill=chfourfig6c0,fill opacity=0.5,draw=chfourfig6c0,line width=0.3985pt,draw opacity=0.5] (63.4452,35.0603) .. controls (63.4452,34.1690) and (62.7226,33.4464) ..(61.8313,33.4464) .. controls (60.9399,33.4464) and (60.2173,34.1690) ..(60.2173,35.0603) .. controls (60.2173,35.9517) and (60.9399,36.6743) ..(61.8313,36.6743) .. controls (62.7226,36.6743) and (63.4452,35.9517) ..(63.4452,35.0603);
\path[fill=chfourfig6c0,fill opacity=0.5,draw=chfourfig6c0,line width=0.3985pt,draw opacity=0.5] (64.7005,35.0603) .. controls (64.7005,34.1690) and (63.9779,33.4464) ..(63.0866,33.4464) .. controls (62.1952,33.4464) and (61.4726,34.1690) ..(61.4726,35.0603) .. controls (61.4726,35.9517) and (62.1952,36.6743) ..(63.0866,36.6743) .. controls (63.9779,36.6743) and (64.7005,35.9517) ..(64.7005,35.0603);
\path[fill=chfourfig6c0,fill opacity=0.5,draw=chfourfig6c0,line width=0.3985pt,draw opacity=0.5] (57.2783,45.1526) .. controls (57.2783,44.2613) and (56.5557,43.5387) ..(55.6643,43.5387) .. controls (54.7729,43.5387) and (54.0503,44.2613) ..(54.0503,45.1526) .. controls (54.0503,46.0440) and (54.7729,46.7666) ..(55.6643,46.7666) .. controls (56.5557,46.7666) and (57.2783,46.0440) ..(57.2783,45.1526);
\path[fill=chfourfig6c0,fill opacity=0.5,draw=chfourfig6c0,line width=0.3985pt,draw opacity=0.5] (99.0223,41.1177) .. controls (99.0223,40.2263) and (98.2997,39.5037) ..(97.4083,39.5037) .. controls (96.5169,39.5037) and (95.7943,40.2263) ..(95.7943,41.1177) .. controls (95.7943,42.0091) and (96.5169,42.7317) ..(97.4083,42.7317) .. controls (98.2997,42.7317) and (99.0223,42.0091) ..(99.0223,41.1177);
\path[fill=chfourfig6c0,fill opacity=0.5,draw=chfourfig6c0,line width=0.3985pt,draw opacity=0.5] (52.8448,48.1813) .. controls (52.8448,47.2899) and (52.1222,46.5673) ..(51.2309,46.5673) .. controls (50.3395,46.5673) and (49.6169,47.2899) ..(49.6169,48.1813) .. controls (49.6169,49.0727) and (50.3395,49.7953) ..(51.2309,49.7953) .. controls (52.1222,49.7953) and (52.8448,49.0727) ..(52.8448,48.1813);
\path[fill=chfourfig6c0,fill opacity=0.5,draw=chfourfig6c0,line width=0.3985pt,draw opacity=0.5] (129.4784,66.3534) .. controls (129.4784,65.4620) and (128.7558,64.7394) ..(127.8645,64.7394) .. controls (126.9731,64.7394) and (126.2505,65.4620) ..(126.2505,66.3534) .. controls (126.2505,67.2448) and (126.9731,67.9674) ..(127.8645,67.9674) .. controls (128.7558,67.9674) and (129.4784,67.2448) ..(129.4784,66.3534);
\path[fill=chfourfig6c0,fill opacity=0.5,draw=chfourfig6c0,line width=0.3985pt,draw opacity=0.5] (142.1112,73.4170) .. controls (142.1112,72.5256) and (141.3886,71.8030) ..(140.4973,71.8030) .. controls (139.6059,71.8030) and (138.8833,72.5256) ..(138.8833,73.4170) .. controls (138.8833,74.3084) and (139.6059,75.0310) ..(140.4973,75.0310) .. controls (141.3886,75.0310) and (142.1112,74.3084) ..(142.1112,73.4170);
\path[fill=chfourfig6c0,fill opacity=0.5,draw=chfourfig6c0,line width=0.3985pt,draw opacity=0.5] (103.0671,65.3372) .. controls (103.0671,64.4458) and (102.3445,63.7232) ..(101.4532,63.7232) .. controls (100.5618,63.7232) and (99.8392,64.4458) ..(99.8392,65.3372) .. controls (99.8392,66.2286) and (100.5618,66.9512) ..(101.4532,66.9512) .. controls (102.3445,66.9512) and (103.0671,66.2286) ..(103.0671,65.3372);
\path[fill=chfourfig6c0,fill opacity=0.5,draw=chfourfig6c0,line width=0.3985pt,draw opacity=0.5] (90.6336,41.1177) .. controls (90.6336,40.2263) and (89.9110,39.5037) ..(89.0197,39.5037) .. controls (88.1283,39.5037) and (87.4057,40.2263) ..(87.4057,41.1177) .. controls (87.4057,42.0091) and (88.1283,42.7317) ..(89.0197,42.7317) .. controls (89.9110,42.7317) and (90.6336,42.0091) ..(90.6336,41.1177);
\path[fill=chfourfig6c0,fill opacity=0.5,draw=chfourfig6c0,line width=0.3985pt,draw opacity=0.5] (53.4227,45.1526) .. controls (53.4227,44.2613) and (52.7001,43.5387) ..(51.8087,43.5387) .. controls (50.9173,43.5387) and (50.1947,44.2613) ..(50.1947,45.1526) .. controls (50.1947,46.0440) and (50.9173,46.7666) ..(51.8087,46.7666) .. controls (52.7001,46.7666) and (53.4227,46.0440) ..(53.4227,45.1526);
\path[fill=chfourfig6c0,fill opacity=0.5,draw=chfourfig6c0,line width=0.3985pt,draw opacity=0.5] (77.2337,52.2162) .. controls (77.2337,51.3249) and (76.5111,50.6023) ..(75.6197,50.6023) .. controls (74.7283,50.6023) and (74.0057,51.3249) ..(74.0057,52.2162) .. controls (74.0057,53.1076) and (74.7283,53.8302) ..(75.6197,53.8302) .. controls (76.5111,53.8302) and (77.2337,53.1076) ..(77.2337,52.2162);
\path[fill=chfourfig6c0,fill opacity=0.5,draw=chfourfig6c0,line width=0.3985pt,draw opacity=0.5] (57.1786,47.1751) .. controls (57.1786,46.2837) and (56.4560,45.5611) ..(55.5647,45.5611) .. controls (54.6733,45.5611) and (53.9507,46.2837) ..(53.9507,47.1751) .. controls (53.9507,48.0664) and (54.6733,48.7890) ..(55.5647,48.7890) .. controls (56.4560,48.7890) and (57.1786,48.0664) ..(57.1786,47.1751);
\path[fill=chfourfig6c0,fill opacity=0.5,draw=chfourfig6c0,line width=0.3985pt,draw opacity=0.5] (121.9565,67.3596) .. controls (121.9565,66.4683) and (121.2340,65.7457) ..(120.3426,65.7457) .. controls (119.4512,65.7457) and (118.7286,66.4683) ..(118.7286,67.3596) .. controls (118.7286,68.2510) and (119.4512,68.9736) ..(120.3426,68.9736) .. controls (121.2340,68.9736) and (121.9565,68.2510) ..(121.9565,67.3596);
\path[fill=chfourfig6c0,fill opacity=0.5,draw=chfourfig6c0,line width=0.3985pt,draw opacity=0.5] (64.0230,33.0379) .. controls (64.0230,32.1465) and (63.3004,31.4239) ..(62.4091,31.4239) .. controls (61.5177,31.4239) and (60.7951,32.1465) ..(60.7951,33.0379) .. controls (60.7951,33.9293) and (61.5177,34.6519) ..(62.4091,34.6519) .. controls (63.3004,34.6519) and (64.0230,33.9293) ..(64.0230,33.0379);
\path[fill=chfourfig6c0,fill opacity=0.5,draw=chfourfig6c0,line width=0.3985pt,draw opacity=0.5] (56.0230,46.1589) .. controls (56.0230,45.2675) and (55.3004,44.5449) ..(54.4090,44.5449) .. controls (53.5176,44.5449) and (52.7950,45.2675) ..(52.7950,46.1589) .. controls (52.7950,47.0502) and (53.5176,47.7728) ..(54.4090,47.7728) .. controls (55.3004,47.7728) and (56.0230,47.0502) ..(56.0230,46.1589);
\path[fill=chfourfig6c0,fill opacity=0.5,draw=chfourfig6c0,line width=0.3985pt,draw opacity=0.5] (82.6335,54.2387) .. controls (82.6335,53.3473) and (81.9109,52.6247) ..(81.0195,52.6247) .. controls (80.1282,52.6247) and (79.4056,53.3473) ..(79.4056,54.2387) .. controls (79.4056,55.1300) and (80.1282,55.8526) ..(81.0195,55.8526) .. controls (81.9109,55.8526) and (82.6335,55.1300) ..(82.6335,54.2387);
\path[fill=chfourfig6c0,fill opacity=0.5,draw=chfourfig6c0,line width=0.3985pt,draw opacity=0.5] (61.3231,32.0317) .. controls (61.3231,31.1403) and (60.6005,30.4177) ..(59.7092,30.4177) .. controls (58.8178,30.4177) and (58.0952,31.1403) ..(58.0952,32.0317) .. controls (58.0952,32.9230) and (58.8178,33.6456) ..(59.7092,33.6456) .. controls (60.6005,33.6456) and (61.3231,32.9230) ..(61.3231,32.0317);
\path[fill=chfourfig6c0,fill opacity=0.5,draw=chfourfig6c0,line width=0.3985pt,draw opacity=0.5] (70.5786,41.1177) .. controls (70.5786,40.2263) and (69.8560,39.5037) ..(68.9646,39.5037) .. controls (68.0732,39.5037) and (67.3506,40.2263) ..(67.3506,41.1177) .. controls (67.3506,42.0091) and (68.0732,42.7317) ..(68.9646,42.7317) .. controls (69.8560,42.7317) and (70.5786,42.0091) ..(70.5786,41.1177);
\path[fill=chfourfig6c0,fill opacity=0.5,draw=chfourfig6c0,line width=0.3985pt,draw opacity=0.5] (52.8448,47.1751) .. controls (52.8448,46.2837) and (52.1222,45.5611) ..(51.2309,45.5611) .. controls (50.3395,45.5611) and (49.6169,46.2837) ..(49.6169,47.1751) .. controls (49.6169,48.0664) and (50.3395,48.7890) ..(51.2309,48.7890) .. controls (52.1222,48.7890) and (52.8448,48.0664) ..(52.8448,47.1751);
\path[fill=chfourfig6c0,fill opacity=0.5,draw=chfourfig6c0,line width=0.3985pt,draw opacity=0.5] (75.5898,50.2038) .. controls (75.5898,49.3124) and (74.8672,48.5898) ..(73.9759,48.5898) .. controls (73.0845,48.5898) and (72.3619,49.3124) ..(72.3619,50.2038) .. controls (72.3619,51.0951) and (73.0845,51.8177) ..(73.9759,51.8177) .. controls (74.8672,51.8177) and (75.5898,51.0951) ..(75.5898,50.2038);
\path[fill=chfourfig6c0,fill opacity=0.5,draw=chfourfig6c0,line width=0.3985pt,draw opacity=0.5] (69.3232,45.1526) .. controls (69.3232,44.2613) and (68.6006,43.5387) ..(67.7093,43.5387) .. controls (66.8179,43.5387) and (66.0953,44.2613) ..(66.0953,45.1526) .. controls (66.0953,46.0440) and (66.8179,46.7666) ..(67.7093,46.7666) .. controls (68.6006,46.7666) and (69.3232,46.0440) ..(69.3232,45.1526);
\path[fill=chfourfig6c0,fill opacity=0.5,draw=chfourfig6c0,line width=0.3985pt,draw opacity=0.5] (58.2446,51.2100) .. controls (58.2446,50.3186) and (57.5220,49.5960) ..(56.6307,49.5960) .. controls (55.7393,49.5960) and (55.0167,50.3186) ..(55.0167,51.2100) .. controls (55.0167,52.1014) and (55.7393,52.8240) ..(56.6307,52.8240) .. controls (57.5220,52.8240) and (58.2446,52.1014) ..(58.2446,51.2100);
\path[fill=chfourfig6c0,fill opacity=0.5,draw=chfourfig6c0,line width=0.3985pt,draw opacity=0.5] (55.2558,49.1875) .. controls (55.2558,48.2962) and (54.5332,47.5736) ..(53.6418,47.5736) .. controls (52.7505,47.5736) and (52.0279,48.2962) ..(52.0279,49.1875) .. controls (52.0279,50.0789) and (52.7505,50.8015) ..(53.6418,50.8015) .. controls (54.5332,50.8015) and (55.2558,50.0789) ..(55.2558,49.1875);
\path[fill=chfourfig6c0,fill opacity=0.5,draw=chfourfig6c0,line width=0.3985pt,draw opacity=0.5] (53.2334,51.2100) .. controls (53.2334,50.3186) and (52.5108,49.5960) ..(51.6194,49.5960) .. controls (50.7280,49.5960) and (50.0054,50.3186) ..(50.0054,51.2100) .. controls (50.0054,52.1014) and (50.7280,52.8240) ..(51.6194,52.8240) .. controls (52.5108,52.8240) and (53.2334,52.1014) ..(53.2334,51.2100);
\path[fill=chfourfig6c0,fill opacity=0.5,draw=chfourfig6c0,line width=0.3985pt,draw opacity=0.5] (56.5111,54.2387) .. controls (56.5111,53.3473) and (55.7885,52.6247) ..(54.8972,52.6247) .. controls (54.0058,52.6247) and (53.2832,53.3473) ..(53.2832,54.2387) .. controls (53.2832,55.1300) and (54.0058,55.8526) ..(54.8972,55.8526) .. controls (55.7885,55.8526) and (56.5111,55.1300) ..(56.5111,54.2387);
\path[fill=chfourfig6c0,fill opacity=0.5,draw=chfourfig6c0,line width=0.3985pt,draw opacity=0.5] (56.7004,45.1526) .. controls (56.7004,44.2613) and (55.9778,43.5387) ..(55.0864,43.5387) .. controls (54.1951,43.5387) and (53.4725,44.2613) ..(53.4725,45.1526) .. controls (53.4725,46.0440) and (54.1951,46.7666) ..(55.0864,46.7666) .. controls (55.9778,46.7666) and (56.7004,46.0440) ..(56.7004,45.1526);
\path[fill=chfourfig6c0,fill opacity=0.5,draw=chfourfig6c0,line width=0.3985pt,draw opacity=0.5] (191.0782,105.7163) .. controls (191.0782,104.8249) and (190.3557,104.1023) ..(189.4643,104.1023) .. controls (188.5729,104.1023) and (187.8503,104.8249) ..(187.8503,105.7163) .. controls (187.8503,106.6077) and (188.5729,107.3303) ..(189.4643,107.3303) .. controls (190.3557,107.3303) and (191.0782,106.6077) ..(191.0782,105.7163);
\path[fill=chfourfig6c0,fill opacity=0.5,draw=chfourfig6c0,line width=0.3985pt,draw opacity=0.5] (57.9557,53.2324) .. controls (57.9557,52.3411) and (57.2331,51.6185) ..(56.3418,51.6185) .. controls (55.4504,51.6185) and (54.7278,52.3411) ..(54.7278,53.2324) .. controls (54.7278,54.1238) and (55.4504,54.8464) ..(56.3418,54.8464) .. controls (57.2331,54.8464) and (57.9557,54.1238) ..(57.9557,53.2324);
\path[fill=chfourfig6c0,fill opacity=0.5,draw=chfourfig6c0,line width=0.3985pt,draw opacity=0.5] (81.5675,58.2736) .. controls (81.5675,57.3822) and (80.8449,56.6596) ..(79.9535,56.6596) .. controls (79.0622,56.6596) and (78.3396,57.3822) ..(78.3396,58.2736) .. controls (78.3396,59.1650) and (79.0622,59.8876) ..(79.9535,59.8876) .. controls (80.8449,59.8876) and (81.5675,59.1650) ..(81.5675,58.2736);
\path[fill=chfourfig6c0,fill opacity=0.5,draw=chfourfig6c0,line width=0.3985pt,draw opacity=0.5] (107.3113,71.3946) .. controls (107.3113,70.5032) and (106.5887,69.7806) ..(105.6973,69.7806) .. controls (104.8059,69.7806) and (104.0833,70.5032) ..(104.0833,71.3946) .. controls (104.0833,72.2859) and (104.8059,73.0085) ..(105.6973,73.0085) .. controls (106.5887,73.0085) and (107.3113,72.2859) ..(107.3113,71.3946);
\path[fill=chfourfig6c0,fill opacity=0.5,draw=chfourfig6c0,line width=0.3985pt,draw opacity=0.5] (61.6121,41.1177) .. controls (61.6121,40.2263) and (60.8895,39.5037) ..(59.9981,39.5037) .. controls (59.1067,39.5037) and (58.3841,40.2263) ..(58.3841,41.1177) .. controls (58.3841,42.0091) and (59.1067,42.7317) ..(59.9981,42.7317) .. controls (60.8895,42.7317) and (61.6121,42.0091) ..(61.6121,41.1177);
\path[fill=chfourfig6c0,fill opacity=0.5,draw=chfourfig6c0,line width=0.3985pt,draw opacity=0.5] (53.7116,46.1589) .. controls (53.7116,45.2675) and (52.9890,44.5449) ..(52.0976,44.5449) .. controls (51.2063,44.5449) and (50.4837,45.2675) ..(50.4837,46.1589) .. controls (50.4837,47.0502) and (51.2063,47.7728) ..(52.0976,47.7728) .. controls (52.9890,47.7728) and (53.7116,47.0502) ..(53.7116,46.1589);
\path[fill=chfourfig6c0,fill opacity=0.5,draw=chfourfig6c0,line width=0.3985pt,draw opacity=0.5] (54.7677,38.0890) .. controls (54.7677,37.1976) and (54.0451,36.4751) ..(53.1537,36.4751) .. controls (52.2623,36.4751) and (51.5397,37.1976) ..(51.5397,38.0890) .. controls (51.5397,38.9804) and (52.2623,39.7030) ..(53.1537,39.7030) .. controls (54.0451,39.7030) and (54.7677,38.9804) ..(54.7677,38.0890);
\path[fill=chfourfig6c0,fill opacity=0.5,draw=chfourfig6c0,line width=0.3985pt,draw opacity=0.5] (54.3891,42.1239) .. controls (54.3891,41.2326) and (53.6665,40.5100) ..(52.7751,40.5100) .. controls (51.8837,40.5100) and (51.1611,41.2326) ..(51.1611,42.1239) .. controls (51.1611,43.0153) and (51.8837,43.7379) ..(52.7751,43.7379) .. controls (53.6665,43.7379) and (54.3891,43.0153) ..(54.3891,42.1239);
\path[fill=chfourfig6c0,fill opacity=0.5,draw=chfourfig6c0,line width=0.3985pt,draw opacity=0.5] (58.3343,43.1402) .. controls (58.3343,42.2488) and (57.6117,41.5262) ..(56.7203,41.5262) .. controls (55.8290,41.5262) and (55.1064,42.2488) ..(55.1064,43.1402) .. controls (55.1064,44.0315) and (55.8290,44.7541) ..(56.7203,44.7541) .. controls (57.6117,44.7541) and (58.3343,44.0315) ..(58.3343,43.1402);
\path[fill=chfourfig6c0,fill opacity=0.5,draw=chfourfig6c0,line width=0.3985pt,draw opacity=0.5] (78.5786,47.1751) .. controls (78.5786,46.2837) and (77.8560,45.5611) ..(76.9647,45.5611) .. controls (76.0733,45.5611) and (75.3507,46.2837) ..(75.3507,47.1751) .. controls (75.3507,48.0664) and (76.0733,48.7890) ..(76.9647,48.7890) .. controls (77.8560,48.7890) and (78.5786,48.0664) ..(78.5786,47.1751);
\path[fill=chfourfig6c0,fill opacity=0.5,draw=chfourfig6c0,line width=0.3985pt,draw opacity=0.5] (91.6000,64.3310) .. controls (91.6000,63.4396) and (90.8774,62.7170) ..(89.9860,62.7170) .. controls (89.0947,62.7170) and (88.3721,63.4396) ..(88.3721,64.3310) .. controls (88.3721,65.2223) and (89.0947,65.9449) ..(89.9860,65.9449) .. controls (90.8774,65.9449) and (91.6000,65.2223) ..(91.6000,64.3310);
\path[fill=chfourfig6c0,fill opacity=0.5,draw=chfourfig6c0,line width=0.3985pt,draw opacity=0.5] (63.3555,25.9743) .. controls (63.3555,25.0829) and (62.6330,24.3603) ..(61.7416,24.3603) .. controls (60.8502,24.3603) and (60.1276,25.0829) ..(60.1276,25.9743) .. controls (60.1276,26.8657) and (60.8502,27.5883) ..(61.7416,27.5883) .. controls (62.6330,27.5883) and (63.3555,26.8657) ..(63.3555,25.9743);
\path[fill=chfourfig6c0,fill opacity=0.5,draw=chfourfig6c0,line width=0.3985pt,draw opacity=0.5] (100.8455,64.3310) .. controls (100.8455,63.4396) and (100.1229,62.7170) ..(99.2315,62.7170) .. controls (98.3401,62.7170) and (97.6175,63.4396) ..(97.6175,64.3310) .. controls (97.6175,65.2223) and (98.3401,65.9449) ..(99.2315,65.9449) .. controls (100.1229,65.9449) and (100.8455,65.2223) ..(100.8455,64.3310);
\path[fill=chfourfig6c0,fill opacity=0.5,draw=chfourfig6c0,line width=0.3985pt,draw opacity=0.5] (100.4669,55.2449) .. controls (100.4669,54.3535) and (99.7443,53.6309) ..(98.8529,53.6309) .. controls (97.9615,53.6309) and (97.2389,54.3535) ..(97.2389,55.2449) .. controls (97.2389,56.1363) and (97.9615,56.8589) ..(98.8529,56.8589) .. controls (99.7443,56.8589) and (100.4669,56.1363) ..(100.4669,55.2449);
\path[fill=chfourfig6c0,fill opacity=0.5,draw=chfourfig6c0,line width=0.3985pt,draw opacity=0.5] (53.0341,41.1177) .. controls (53.0341,40.2263) and (52.3115,39.5037) ..(51.4202,39.5037) .. controls (50.5288,39.5037) and (49.8062,40.2263) ..(49.8062,41.1177) .. controls (49.8062,42.0091) and (50.5288,42.7317) ..(51.4202,42.7317) .. controls (52.3115,42.7317) and (53.0341,42.0091) ..(53.0341,41.1177);
\path[fill=chfourfig6c0,fill opacity=0.5,draw=chfourfig6c0,line width=0.3985pt,draw opacity=0.5] (66.1451,43.1402) .. controls (66.1451,42.2488) and (65.4225,41.5262) ..(64.5312,41.5262) .. controls (63.6398,41.5262) and (62.9172,42.2488) ..(62.9172,43.1402) .. controls (62.9172,44.0315) and (63.6398,44.7541) ..(64.5312,44.7541) .. controls (65.4225,44.7541) and (66.1451,44.0315) ..(66.1451,43.1402);
\path[fill=chfourfig6c0,fill opacity=0.5,draw=chfourfig6c0,line width=0.3985pt,draw opacity=0.5] (52.8448,50.2038) .. controls (52.8448,49.3124) and (52.1222,48.5898) ..(51.2309,48.5898) .. controls (50.3395,48.5898) and (49.6169,49.3124) ..(49.6169,50.2038) .. controls (49.6169,51.0951) and (50.3395,51.8177) ..(51.2309,51.8177) .. controls (52.1222,51.8177) and (52.8448,51.0951) ..(52.8448,50.2038);
\path[fill=chfourfig6c0,fill opacity=0.5,draw=chfourfig6c0,line width=0.3985pt,draw opacity=0.5] (60.0778,39.0953) .. controls (60.0778,38.2039) and (59.3552,37.4813) ..(58.4638,37.4813) .. controls (57.5725,37.4813) and (56.8499,38.2039) ..(56.8499,39.0953) .. controls (56.8499,39.9866) and (57.5725,40.7092) ..(58.4638,40.7092) .. controls (59.3552,40.7092) and (60.0778,39.9866) ..(60.0778,39.0953);
\path[fill=chfourfig6c0,fill opacity=0.5,draw=chfourfig6c0,line width=0.3985pt,draw opacity=0.5] (86.3895,47.1751) .. controls (86.3895,46.2837) and (85.6669,45.5611) ..(84.7755,45.5611) .. controls (83.8841,45.5611) and (83.1615,46.2837) ..(83.1615,47.1751) .. controls (83.1615,48.0664) and (83.8841,48.7890) ..(84.7755,48.7890) .. controls (85.6669,48.7890) and (86.3895,48.0664) ..(86.3895,47.1751);
\path[fill=chfourfig6c0,fill opacity=0.5,draw=chfourfig6c0,line width=0.3985pt,draw opacity=0.5] (103.7446,63.3247) .. controls (103.7446,62.4333) and (103.0220,61.7107) ..(102.1306,61.7107) .. controls (101.2393,61.7107) and (100.5167,62.4333) ..(100.5167,63.3247) .. controls (100.5167,64.2161) and (101.2393,64.9387) ..(102.1306,64.9387) .. controls (103.0220,64.9387) and (103.7446,64.2161) ..(103.7446,63.3247);
\path[fill=chfourfig6c0,fill opacity=0.5,draw=chfourfig6c0,line width=0.3985pt,draw opacity=0.5] (78.8676,45.1526) .. controls (78.8676,44.2613) and (78.1450,43.5387) ..(77.2536,43.5387) .. controls (76.3622,43.5387) and (75.6396,44.2613) ..(75.6396,45.1526) .. controls (75.6396,46.0440) and (76.3622,46.7666) ..(77.2536,46.7666) .. controls (78.1450,46.7666) and (78.8676,46.0440) ..(78.8676,45.1526);
\path[fill=chfourfig6c0,fill opacity=0.5,draw=chfourfig6c0,line width=0.3985pt,draw opacity=0.5] (154.1562,73.4170) .. controls (154.1562,72.5256) and (153.4336,71.8030) ..(152.5423,71.8030) .. controls (151.6509,71.8030) and (150.9283,72.5256) ..(150.9283,73.4170) .. controls (150.9283,74.3084) and (151.6509,75.0310) ..(152.5423,75.0310) .. controls (153.4336,75.0310) and (154.1562,74.3084) ..(154.1562,73.4170);
\path[fill=chfourfig6c0,fill opacity=0.5,draw=chfourfig6c0,line width=0.3985pt,draw opacity=0.5] (137.7675,76.4457) .. controls (137.7675,75.5543) and (137.0449,74.8317) ..(136.1535,74.8317) .. controls (135.2621,74.8317) and (134.5396,75.5543) ..(134.5396,76.4457) .. controls (134.5396,77.3370) and (135.2621,78.0596) ..(136.1535,78.0596) .. controls (137.0449,78.0596) and (137.7675,77.3370) ..(137.7675,76.4457);
\path[fill=chfourfig6c0,fill opacity=0.5,draw=chfourfig6c0,line width=0.3985pt,draw opacity=0.5] (149.8224,71.3946) .. controls (149.8224,70.5032) and (149.0998,69.7806) ..(148.2084,69.7806) .. controls (147.3171,69.7806) and (146.5945,70.5032) ..(146.5945,71.3946) .. controls (146.5945,72.2859) and (147.3171,73.0085) ..(148.2084,73.0085) .. controls (149.0998,73.0085) and (149.8224,72.2859) ..(149.8224,71.3946);
\path[fill=chfourfig6c0,fill opacity=0.5,draw=chfourfig6c0,line width=0.3985pt,draw opacity=0.5] (107.3113,68.3659) .. controls (107.3113,67.4745) and (106.5887,66.7519) ..(105.6973,66.7519) .. controls (104.8059,66.7519) and (104.0833,67.4745) ..(104.0833,68.3659) .. controls (104.0833,69.2572) and (104.8059,69.9798) ..(105.6973,69.9798) .. controls (106.5887,69.9798) and (107.3113,69.2572) ..(107.3113,68.3659);
\path[fill=chfourfig6c0,fill opacity=0.5,draw=chfourfig6c0,line width=0.3985pt,draw opacity=0.5] (75.9784,55.2449) .. controls (75.9784,54.3535) and (75.2558,53.6309) ..(74.3644,53.6309) .. controls (73.4731,53.6309) and (72.7505,54.3535) ..(72.7505,55.2449) .. controls (72.7505,56.1363) and (73.4731,56.8589) ..(74.3644,56.8589) .. controls (75.2558,56.8589) and (75.9784,56.1363) ..(75.9784,55.2449);
\path[fill=chfourfig6c0,fill opacity=0.5,draw=chfourfig6c0,line width=0.3985pt,draw opacity=0.5] (146.8336,67.3596) .. controls (146.8336,66.4683) and (146.1110,65.7457) ..(145.2196,65.7457) .. controls (144.3282,65.7457) and (143.6057,66.4683) ..(143.6057,67.3596) .. controls (143.6057,68.2510) and (144.3282,68.9736) ..(145.2196,68.9736) .. controls (146.1110,68.9736) and (146.8336,68.2510) ..(146.8336,67.3596);
\path[fill=chfourfig6c0,fill opacity=0.5,draw=chfourfig6c0,line width=0.3985pt,draw opacity=0.5] (54.2894,52.2162) .. controls (54.2894,51.3249) and (53.5668,50.6023) ..(52.6754,50.6023) .. controls (51.7841,50.6023) and (51.0615,51.3249) ..(51.0615,52.2162) .. controls (51.0615,53.1076) and (51.7841,53.8302) ..(52.6754,53.8302) .. controls (53.5668,53.8302) and (54.2894,53.1076) ..(54.2894,52.2162);
\path[fill=chfourfig6c0,fill opacity=0.5,draw=chfourfig6c0,line width=0.3985pt,draw opacity=0.5] (64.9894,43.1402) .. controls (64.9894,42.2488) and (64.2668,41.5262) ..(63.3755,41.5262) .. controls (62.4841,41.5262) and (61.7615,42.2488) ..(61.7615,43.1402) .. controls (61.7615,44.0315) and (62.4841,44.7541) ..(63.3755,44.7541) .. controls (64.2668,44.7541) and (64.9894,44.0315) ..(64.9894,43.1402);
\path[fill=chfourfig6c0,fill opacity=0.5,draw=chfourfig6c0,line width=0.3985pt,draw opacity=0.5] (95.0670,63.3247) .. controls (95.0670,62.4333) and (94.3444,61.7107) ..(93.4531,61.7107) .. controls (92.5617,61.7107) and (91.8391,62.4333) ..(91.8391,63.3247) .. controls (91.8391,64.2161) and (92.5617,64.9387) ..(93.4531,64.9387) .. controls (94.3444,64.9387) and (95.0670,64.2161) ..(95.0670,63.3247);
\path[fill=chfourfig6c0,fill opacity=0.5,draw=chfourfig6c0,line width=0.3985pt,draw opacity=0.5] (161.1003,84.5155) .. controls (161.1003,83.6242) and (160.3777,82.9016) ..(159.4863,82.9016) .. controls (158.5949,82.9016) and (157.8723,83.6242) ..(157.8723,84.5155) .. controls (157.8723,85.4069) and (158.5949,86.1295) ..(159.4863,86.1295) .. controls (160.3777,86.1295) and (161.1003,85.4069) ..(161.1003,84.5155);
\path[fill=chfourfig6c0,fill opacity=0.5,draw=chfourfig6c0,line width=0.3985pt,draw opacity=0.5] (94.8777,56.2611) .. controls (94.8777,55.3697) and (94.1552,54.6471) ..(93.2638,54.6471) .. controls (92.3724,54.6471) and (91.6498,55.3697) ..(91.6498,56.2611) .. controls (91.6498,57.1525) and (92.3724,57.8751) ..(93.2638,57.8751) .. controls (94.1552,57.8751) and (94.8777,57.1525) ..(94.8777,56.2611);
\path[fill=chfourfig6c0,fill opacity=0.5,draw=chfourfig6c0,line width=0.3985pt,draw opacity=0.5] (52.9445,50.2038) .. controls (52.9445,49.3124) and (52.2219,48.5898) ..(51.3305,48.5898) .. controls (50.4391,48.5898) and (49.7165,49.3124) ..(49.7165,50.2038) .. controls (49.7165,51.0951) and (50.4391,51.8177) ..(51.3305,51.8177) .. controls (52.2219,51.8177) and (52.9445,51.0951) ..(52.9445,50.2038);
\path[fill=chfourfig6c0,fill opacity=0.5,draw=chfourfig6c0,line width=0.3985pt,draw opacity=0.5] (56.2222,54.2387) .. controls (56.2222,53.3473) and (55.4996,52.6247) ..(54.6082,52.6247) .. controls (53.7169,52.6247) and (52.9943,53.3473) ..(52.9943,54.2387) .. controls (52.9943,55.1300) and (53.7169,55.8526) ..(54.6082,55.8526) .. controls (55.4996,55.8526) and (56.2222,55.1300) ..(56.2222,54.2387);
\path[fill=chfourfig6c0,fill opacity=0.5,draw=chfourfig6c0,line width=0.3985pt,draw opacity=0.5] (121.2890,64.3310) .. controls (121.2890,63.4396) and (120.5664,62.7170) ..(119.6751,62.7170) .. controls (118.7837,62.7170) and (118.0611,63.4396) ..(118.0611,64.3310) .. controls (118.0611,65.2223) and (118.7837,65.9449) ..(119.6751,65.9449) .. controls (120.5664,65.9449) and (121.2890,65.2223) ..(121.2890,64.3310);
\path[fill=chfourfig6c0,fill opacity=0.5,draw=chfourfig6c0,line width=0.3985pt,draw opacity=0.5] (67.5897,39.0953) .. controls (67.5897,38.2039) and (66.8671,37.4813) ..(65.9758,37.4813) .. controls (65.0844,37.4813) and (64.3618,38.2039) ..(64.3618,39.0953) .. controls (64.3618,39.9866) and (65.0844,40.7092) ..(65.9758,40.7092) .. controls (66.8671,40.7092) and (67.5897,39.9866) ..(67.5897,39.0953);
\path[fill=chfourfig6c0,fill opacity=0.5,draw=chfourfig6c0,line width=0.3985pt,draw opacity=0.5] (53.3230,46.1589) .. controls (53.3230,45.2675) and (52.6004,44.5449) ..(51.7091,44.5449) .. controls (50.8177,44.5449) and (50.0951,45.2675) ..(50.0951,46.1589) .. controls (50.0951,47.0502) and (50.8177,47.7728) ..(51.7091,47.7728) .. controls (52.6004,47.7728) and (53.3230,47.0502) ..(53.3230,46.1589);
\path[fill=chfourfig6c0,fill opacity=0.5,draw=chfourfig6c0,line width=0.3985pt,draw opacity=0.5] (61.4228,36.0666) .. controls (61.4228,35.1752) and (60.7002,34.4526) ..(59.8088,34.4526) .. controls (58.9174,34.4526) and (58.1949,35.1752) ..(58.1949,36.0666) .. controls (58.1949,36.9579) and (58.9174,37.6805) ..(59.8088,37.6805) .. controls (60.7002,37.6805) and (61.4228,36.9579) ..(61.4228,36.0666);
\path[fill=chfourfig6c0,fill opacity=0.5,draw=chfourfig6c0,line width=0.3985pt,draw opacity=0.5] (59.8785,45.1526) .. controls (59.8785,44.2613) and (59.1559,43.5387) ..(58.2646,43.5387) .. controls (57.3732,43.5387) and (56.6506,44.2613) ..(56.6506,45.1526) .. controls (56.6506,46.0440) and (57.3732,46.7666) ..(58.2646,46.7666) .. controls (59.1559,46.7666) and (59.8785,46.0440) ..(59.8785,45.1526);
\path[fill=chfourfig6c0,fill opacity=0.5,draw=chfourfig6c0,line width=0.3985pt,draw opacity=0.5] (133.5233,83.5093) .. controls (133.5233,82.6179) and (132.8007,81.8953) ..(131.9094,81.8953) .. controls (131.0180,81.8953) and (130.2954,82.6179) ..(130.2954,83.5093) .. controls (130.2954,84.4007) and (131.0180,85.1232) ..(131.9094,85.1232) .. controls (132.8007,85.1232) and (133.5233,84.4007) ..(133.5233,83.5093);
\path[fill=chfourfig6c0,fill opacity=0.5,draw=chfourfig6c0,line width=0.3985pt,draw opacity=0.5] (78.6783,38.0890) .. controls (78.6783,37.1976) and (77.9557,36.4751) ..(77.0643,36.4751) .. controls (76.1730,36.4751) and (75.4504,37.1976) ..(75.4504,38.0890) .. controls (75.4504,38.9804) and (76.1730,39.7030) ..(77.0643,39.7030) .. controls (77.9557,39.7030) and (78.6783,38.9804) ..(78.6783,38.0890);
\path[fill=chfourfig6c0,fill opacity=0.5,draw=chfourfig6c0,line width=0.3985pt,draw opacity=0.5] (101.3336,59.2898) .. controls (101.3336,58.3984) and (100.6110,57.6758) ..(99.7197,57.6758) .. controls (98.8283,57.6758) and (98.1057,58.3984) ..(98.1057,59.2898) .. controls (98.1057,60.1812) and (98.8283,60.9038) ..(99.7197,60.9038) .. controls (100.6110,60.9038) and (101.3336,60.1812) ..(101.3336,59.2898);
\path[fill=chfourfig6c0,fill opacity=0.5,draw=chfourfig6c0,line width=0.3985pt,draw opacity=0.5] (99.6001,65.3372) .. controls (99.6001,64.4458) and (98.8775,63.7232) ..(97.9861,63.7232) .. controls (97.0948,63.7232) and (96.3722,64.4458) ..(96.3722,65.3372) .. controls (96.3722,66.2286) and (97.0948,66.9512) ..(97.9861,66.9512) .. controls (98.8775,66.9512) and (99.6001,66.2286) ..(99.6001,65.3372);
\path[fill=chfourfig6c0,fill opacity=0.5,draw=chfourfig6c0,line width=0.3985pt,draw opacity=0.5] (88.7008,54.2387) .. controls (88.7008,53.3473) and (87.9782,52.6247) ..(87.0869,52.6247) .. controls (86.1955,52.6247) and (85.4729,53.3473) ..(85.4729,54.2387) .. controls (85.4729,55.1300) and (86.1955,55.8526) ..(87.0869,55.8526) .. controls (87.9782,55.8526) and (88.7008,55.1300) ..(88.7008,54.2387);
\path[fill=chfourfig6c0,fill opacity=0.5,draw=chfourfig6c0,line width=0.3985pt,draw opacity=0.5] (63.4452,29.0030) .. controls (63.4452,28.1116) and (62.7226,27.3890) ..(61.8313,27.3890) .. controls (60.9399,27.3890) and (60.2173,28.1116) ..(60.2173,29.0030) .. controls (60.2173,29.8943) and (60.9399,30.6169) ..(61.8313,30.6169) .. controls (62.7226,30.6169) and (63.4452,29.8943) ..(63.4452,29.0030);
\path[fill=chfourfig6c0,fill opacity=0.5,draw=chfourfig6c0,line width=0.3985pt,draw opacity=0.5] (74.9223,47.1751) .. controls (74.9223,46.2837) and (74.1997,45.5611) ..(73.3083,45.5611) .. controls (72.4170,45.5611) and (71.6944,46.2837) ..(71.6944,47.1751) .. controls (71.6944,48.0664) and (72.4170,48.7890) ..(73.3083,48.7890) .. controls (74.1997,48.7890) and (74.9223,48.0664) ..(74.9223,47.1751);
\path[fill=chfourfig6c0,fill opacity=0.5,draw=chfourfig6c0,line width=0.3985pt,draw opacity=0.5] (109.0448,58.2736) .. controls (109.0448,57.3822) and (108.3222,56.6596) ..(107.4308,56.6596) .. controls (106.5395,56.6596) and (105.8169,57.3822) ..(105.8169,58.2736) .. controls (105.8169,59.1650) and (106.5395,59.8876) ..(107.4308,59.8876) .. controls (108.3222,59.8876) and (109.0448,59.1650) ..(109.0448,58.2736);
\path[fill=chfourfig6c0,fill opacity=0.5,draw=chfourfig6c0,line width=0.3985pt,draw opacity=0.5] (120.6116,63.3247) .. controls (120.6116,62.4333) and (119.8890,61.7107) ..(118.9976,61.7107) .. controls (118.1062,61.7107) and (117.3837,62.4333) ..(117.3837,63.3247) .. controls (117.3837,64.2161) and (118.1062,64.9387) ..(118.9976,64.9387) .. controls (119.8890,64.9387) and (120.6116,64.2161) ..(120.6116,63.3247);
\path[fill=chfourfig6c0,fill opacity=0.5,draw=chfourfig6c0,line width=0.3985pt,draw opacity=0.5] (55.2558,47.1751) .. controls (55.2558,46.2837) and (54.5332,45.5611) ..(53.6418,45.5611) .. controls (52.7505,45.5611) and (52.0279,46.2837) ..(52.0279,47.1751) .. controls (52.0279,48.0664) and (52.7505,48.7890) ..(53.6418,48.7890) .. controls (54.5332,48.7890) and (55.2558,48.0664) ..(55.2558,47.1751);
\path[fill=chfourfig6c0,fill opacity=0.5,draw=chfourfig6c0,line width=0.3985pt,draw opacity=0.5] (53.0341,49.1875) .. controls (53.0341,48.2962) and (52.3115,47.5736) ..(51.4202,47.5736) .. controls (50.5288,47.5736) and (49.8062,48.2962) ..(49.8062,49.1875) .. controls (49.8062,50.0789) and (50.5288,50.8015) ..(51.4202,50.8015) .. controls (52.3115,50.8015) and (53.0341,50.0789) ..(53.0341,49.1875);
\path[fill=chfourfig6c0,fill opacity=0.5,draw=chfourfig6c0,line width=0.3985pt,draw opacity=0.5] (79.8340,46.1589) .. controls (79.8340,45.2675) and (79.1114,44.5449) ..(78.2200,44.5449) .. controls (77.3286,44.5449) and (76.6060,45.2675) ..(76.6060,46.1589) .. controls (76.6060,47.0502) and (77.3286,47.7728) ..(78.2200,47.7728) .. controls (79.1114,47.7728) and (79.8340,47.0502) ..(79.8340,46.1589);
\path[fill=chfourfig6c0,fill opacity=0.5,draw=chfourfig6c0,line width=0.3985pt,draw opacity=0.5] (130.2456,76.4457) .. controls (130.2456,75.5543) and (129.5230,74.8317) ..(128.6316,74.8317) .. controls (127.7402,74.8317) and (127.0176,75.5543) ..(127.0176,76.4457) .. controls (127.0176,77.3370) and (127.7402,78.0596) ..(128.6316,78.0596) .. controls (129.5230,78.0596) and (130.2456,77.3370) ..(130.2456,76.4457);
\path[fill=chfourfig6c0,fill opacity=0.5,draw=chfourfig6c0,line width=0.3985pt,draw opacity=0.5] (57.0890,36.0666) .. controls (57.0890,35.1752) and (56.3664,34.4526) ..(55.4750,34.4526) .. controls (54.5836,34.4526) and (53.8610,35.1752) ..(53.8610,36.0666) .. controls (53.8610,36.9579) and (54.5836,37.6805) ..(55.4750,37.6805) .. controls (56.3664,37.6805) and (57.0890,36.9579) ..(57.0890,36.0666);
\path[fill=chfourfig6c0,fill opacity=0.5,draw=chfourfig6c0,line width=0.3985pt,draw opacity=0.5] (146.4450,70.3883) .. controls (146.4450,69.4970) and (145.7224,68.7744) ..(144.8311,68.7744) .. controls (143.9397,68.7744) and (143.2171,69.4970) ..(143.2171,70.3883) .. controls (143.2171,71.2797) and (143.9397,72.0023) ..(144.8311,72.0023) .. controls (145.7224,72.0023) and (146.4450,71.2797) ..(146.4450,70.3883);
\path[fill=chfourfig6c0,fill opacity=0.5,draw=chfourfig6c0,line width=0.3985pt,draw opacity=0.5] (61.5224,36.0666) .. controls (61.5224,35.1752) and (60.7998,34.4526) ..(59.9084,34.4526) .. controls (59.0171,34.4526) and (58.2945,35.1752) ..(58.2945,36.0666) .. controls (58.2945,36.9579) and (59.0171,37.6805) ..(59.9084,37.6805) .. controls (60.7998,37.6805) and (61.5224,36.9579) ..(61.5224,36.0666);
\path[fill=chfourfig6c0,fill opacity=0.5,draw=chfourfig6c0,line width=0.3985pt,draw opacity=0.5] (125.3339,66.3534) .. controls (125.3339,65.4620) and (124.6113,64.7394) ..(123.7200,64.7394) .. controls (122.8286,64.7394) and (122.1060,65.4620) ..(122.1060,66.3534) .. controls (122.1060,67.2448) and (122.8286,67.9674) ..(123.7200,67.9674) .. controls (124.6113,67.9674) and (125.3339,67.2448) ..(125.3339,66.3534);
\path[fill=chfourfig6c0,fill opacity=0.5,draw=chfourfig6c0,line width=0.3985pt,draw opacity=0.5] (92.8453,54.2387) .. controls (92.8453,53.3473) and (92.1227,52.6247) ..(91.2314,52.6247) .. controls (90.3400,52.6247) and (89.6174,53.3473) ..(89.6174,54.2387) .. controls (89.6174,55.1300) and (90.3400,55.8526) ..(91.2314,55.8526) .. controls (92.1227,55.8526) and (92.8453,55.1300) ..(92.8453,54.2387);
\path[fill=chfourfig6c0,fill opacity=0.5,draw=chfourfig6c0,line width=0.3985pt,draw opacity=0.5] (134.5894,63.3247) .. controls (134.5894,62.4333) and (133.8668,61.7107) ..(132.9754,61.7107) .. controls (132.0840,61.7107) and (131.3614,62.4333) ..(131.3614,63.3247) .. controls (131.3614,64.2161) and (132.0840,64.9387) ..(132.9754,64.9387) .. controls (133.8668,64.9387) and (134.5894,64.2161) ..(134.5894,63.3247);
\path[fill=chfourfig6c0,fill opacity=0.5,draw=chfourfig6c0,line width=0.3985pt,draw opacity=0.5] (95.3560,40.1115) .. controls (95.3560,39.2201) and (94.6333,38.4975) ..(93.7420,38.4975) .. controls (92.8506,38.4975) and (92.1280,39.2201) ..(92.1280,40.1115) .. controls (92.1280,41.0028) and (92.8506,41.7254) ..(93.7420,41.7254) .. controls (94.6333,41.7254) and (95.3560,41.0028) ..(95.3560,40.1115);
\path[fill=chfourfig6c0,fill opacity=0.5,draw=chfourfig6c0,line width=0.3985pt,draw opacity=0.5] (127.4560,67.3596) .. controls (127.4560,66.4683) and (126.7334,65.7457) ..(125.8420,65.7457) .. controls (124.9507,65.7457) and (124.2281,66.4683) ..(124.2281,67.3596) .. controls (124.2281,68.2510) and (124.9507,68.9736) ..(125.8420,68.9736) .. controls (126.7334,68.9736) and (127.4560,68.2510) ..(127.4560,67.3596);
\path[fill=chfourfig6c0,fill opacity=0.5,draw=chfourfig6c0,line width=0.3985pt,draw opacity=0.5] (62.4888,37.0828) .. controls (62.4888,36.1914) and (61.7662,35.4688) ..(60.8748,35.4688) .. controls (59.9835,35.4688) and (59.2609,36.1914) ..(59.2609,37.0828) .. controls (59.2609,37.9742) and (59.9835,38.6967) ..(60.8748,38.6967) .. controls (61.7662,38.6967) and (62.4888,37.9742) ..(62.4888,37.0828);
\path[fill=chfourfig6c0,fill opacity=0.5,draw=chfourfig6c0,line width=0.3985pt,draw opacity=0.5] (102.7782,55.2449) .. controls (102.7782,54.3535) and (102.0556,53.6309) ..(101.1642,53.6309) .. controls (100.2729,53.6309) and (99.5503,54.3535) ..(99.5503,55.2449) .. controls (99.5503,56.1363) and (100.2729,56.8589) ..(101.1642,56.8589) .. controls (102.0556,56.8589) and (102.7782,56.1363) ..(102.7782,55.2449);
\path[fill=chfourfig6c0,fill opacity=0.5,draw=chfourfig6c0,line width=0.3985pt,draw opacity=0.5] (94.6785,51.2100) .. controls (94.6785,50.3186) and (93.9559,49.5960) ..(93.0645,49.5960) .. controls (92.1732,49.5960) and (91.4506,50.3186) ..(91.4506,51.2100) .. controls (91.4506,52.1014) and (92.1732,52.8240) ..(93.0645,52.8240) .. controls (93.9559,52.8240) and (94.6785,52.1014) ..(94.6785,51.2100);
\path[fill=chfourfig6c0,fill opacity=0.5,draw=chfourfig6c0,line width=0.3985pt,draw opacity=0.5] (124.6565,65.3372) .. controls (124.6565,64.4458) and (123.9339,63.7232) ..(123.0425,63.7232) .. controls (122.1511,63.7232) and (121.4285,64.4458) ..(121.4285,65.3372) .. controls (121.4285,66.2286) and (122.1511,66.9512) ..(123.0425,66.9512) .. controls (123.9339,66.9512) and (124.6565,66.2286) ..(124.6565,65.3372);
\path[fill=chfourfig6c0,fill opacity=0.5,draw=chfourfig6c0,line width=0.3985pt,draw opacity=0.5] (53.1338,45.1526) .. controls (53.1338,44.2613) and (52.4112,43.5387) ..(51.5198,43.5387) .. controls (50.6284,43.5387) and (49.9058,44.2613) ..(49.9058,45.1526) .. controls (49.9058,46.0440) and (50.6284,46.7666) ..(51.5198,46.7666) .. controls (52.4112,46.7666) and (53.1338,46.0440) ..(53.1338,45.1526);
\path[fill=chfourfig6c0,fill opacity=0.5,draw=chfourfig6c0,line width=0.3985pt,draw opacity=0.5] (64.1227,43.1402) .. controls (64.1227,42.2488) and (63.4001,41.5262) ..(62.5087,41.5262) .. controls (61.6174,41.5262) and (60.8947,42.2488) ..(60.8947,43.1402) .. controls (60.8947,44.0315) and (61.6174,44.7541) ..(62.5087,44.7541) .. controls (63.4001,44.7541) and (64.1227,44.0315) ..(64.1227,43.1402);
\path[fill=chfourfig6c0,fill opacity=0.5,draw=chfourfig6c0,line width=0.3985pt,draw opacity=0.5] (81.5675,57.2673) .. controls (81.5675,56.3760) and (80.8449,55.6534) ..(79.9535,55.6534) .. controls (79.0622,55.6534) and (78.3396,56.3760) ..(78.3396,57.2673) .. controls (78.3396,58.1587) and (79.0622,58.8813) ..(79.9535,58.8813) .. controls (80.8449,58.8813) and (81.5675,58.1587) ..(81.5675,57.2673);
\path[fill=chfourfig6c0,fill opacity=0.5,draw=chfourfig6c0,line width=0.3985pt,draw opacity=0.5] (85.8615,51.2100) .. controls (85.8615,50.3186) and (85.1389,49.5960) ..(84.2475,49.5960) .. controls (83.3561,49.5960) and (82.6335,50.3186) ..(82.6335,51.2100) .. controls (82.6335,52.1014) and (83.3561,52.8240) ..(84.2475,52.8240) .. controls (85.1389,52.8240) and (85.8615,52.1014) ..(85.8615,51.2100);
\path[fill=chfourfig6c0,fill opacity=0.5,draw=chfourfig6c0,line width=0.3985pt,draw opacity=0.5] (53.3230,53.2324) .. controls (53.3230,52.3411) and (52.6004,51.6185) ..(51.7091,51.6185) .. controls (50.8177,51.6185) and (50.0951,52.3411) ..(50.0951,53.2324) .. controls (50.0951,54.1238) and (50.8177,54.8464) ..(51.7091,54.8464) .. controls (52.6004,54.8464) and (53.3230,54.1238) ..(53.3230,53.2324);
\path[fill=chfourfig6c0,fill opacity=0.5,draw=chfourfig6c0,line width=0.3985pt,draw opacity=0.5] (53.5223,50.2038) .. controls (53.5223,49.3124) and (52.7997,48.5898) ..(51.9083,48.5898) .. controls (51.0170,48.5898) and (50.2944,49.3124) ..(50.2944,50.2038) .. controls (50.2944,51.0951) and (51.0170,51.8177) ..(51.9083,51.8177) .. controls (52.7997,51.8177) and (53.5223,51.0951) ..(53.5223,50.2038);
\path[fill=chfourfig6c0,fill opacity=0.5,draw=chfourfig6c0,line width=0.3985pt,draw opacity=0.5] (87.8341,54.2387) .. controls (87.8341,53.3473) and (87.1115,52.6247) ..(86.2201,52.6247) .. controls (85.3287,52.6247) and (84.6061,53.3473) ..(84.6061,54.2387) .. controls (84.6061,55.1300) and (85.3287,55.8526) ..(86.2201,55.8526) .. controls (87.1115,55.8526) and (87.8341,55.1300) ..(87.8341,54.2387);
\path[fill=chfourfig6c0,fill opacity=0.5,draw=chfourfig6c0,line width=0.3985pt,draw opacity=0.5] (86.9673,57.2673) .. controls (86.9673,56.3760) and (86.2447,55.6534) ..(85.3533,55.6534) .. controls (84.4620,55.6534) and (83.7394,56.3760) ..(83.7394,57.2673) .. controls (83.7394,58.1587) and (84.4620,58.8813) ..(85.3533,58.8813) .. controls (86.2447,58.8813) and (86.9673,58.1587) ..(86.9673,57.2673);
\path[fill=chfourfig6c0,fill opacity=0.5,draw=chfourfig6c0,line width=0.3985pt,draw opacity=0.5] (65.9558,41.1177) .. controls (65.9558,40.2263) and (65.2332,39.5037) ..(64.3419,39.5037) .. controls (63.4505,39.5037) and (62.7279,40.2263) ..(62.7279,41.1177) .. controls (62.7279,42.0091) and (63.4505,42.7317) ..(64.3419,42.7317) .. controls (65.2332,42.7317) and (65.9558,42.0091) ..(65.9558,41.1177);
\path[fill=chfourfig6c0,fill opacity=0.5,draw=chfourfig6c0,line width=0.3985pt,draw opacity=0.5] (66.2448,44.1464) .. controls (66.2448,43.2550) and (65.5222,42.5324) ..(64.6308,42.5324) .. controls (63.7394,42.5324) and (63.0168,43.2550) ..(63.0168,44.1464) .. controls (63.0168,45.0378) and (63.7394,45.7603) ..(64.6308,45.7603) .. controls (65.5222,45.7603) and (66.2448,45.0378) ..(66.2448,44.1464);
\path[fill=chfourfig6c0,fill opacity=0.5,draw=chfourfig6c0,line width=0.3985pt,draw opacity=0.5] (125.7225,60.2960) .. controls (125.7225,59.4047) and (124.9999,58.6821) ..(124.1085,58.6821) .. controls (123.2171,58.6821) and (122.4946,59.4047) ..(122.4946,60.2960) .. controls (122.4946,61.1874) and (123.2171,61.9100) ..(124.1085,61.9100) .. controls (124.9999,61.9100) and (125.7225,61.1874) ..(125.7225,60.2960);
\path[fill=chfourfig6c0,fill opacity=0.5,draw=chfourfig6c0,line width=0.3985pt,draw opacity=0.5] (131.5009,71.3946) .. controls (131.5009,70.5032) and (130.7783,69.7806) ..(129.8869,69.7806) .. controls (128.9955,69.7806) and (128.2729,70.5032) ..(128.2729,71.3946) .. controls (128.2729,72.2859) and (128.9955,73.0085) ..(129.8869,73.0085) .. controls (130.7783,73.0085) and (131.5009,72.2859) ..(131.5009,71.3946);
\path[fill=chfourfig6c0,fill opacity=0.5,draw=chfourfig6c0,line width=0.3985pt,draw opacity=0.5] (56.8897,23.9618) .. controls (56.8897,23.0704) and (56.1671,22.3479) ..(55.2757,22.3479) .. controls (54.3844,22.3479) and (53.6618,23.0704) ..(53.6618,23.9618) .. controls (53.6618,24.8532) and (54.3844,25.5758) ..(55.2757,25.5758) .. controls (56.1671,25.5758) and (56.8897,24.8532) ..(56.8897,23.9618);
\path[fill=chfourfig6c0,fill opacity=0.5,draw=chfourfig6c0,line width=0.3985pt,draw opacity=0.5] (80.0233,43.1402) .. controls (80.0233,42.2488) and (79.3007,41.5262) ..(78.4093,41.5262) .. controls (77.5179,41.5262) and (76.7953,42.2488) ..(76.7953,43.1402) .. controls (76.7953,44.0315) and (77.5179,44.7541) ..(78.4093,44.7541) .. controls (79.3007,44.7541) and (80.0233,44.0315) ..(80.0233,43.1402);
\path[fill=chfourfig6c0,fill opacity=0.5,draw=chfourfig6c0,line width=0.3985pt,draw opacity=0.5] (70.4889,39.0953) .. controls (70.4889,38.2039) and (69.7663,37.4813) ..(68.8749,37.4813) .. controls (67.9836,37.4813) and (67.2610,38.2039) ..(67.2610,39.0953) .. controls (67.2610,39.9866) and (67.9836,40.7092) ..(68.8749,40.7092) .. controls (69.7663,40.7092) and (70.4889,39.9866) ..(70.4889,39.0953);
\path[fill=chfourfig6c0,fill opacity=0.5,draw=chfourfig6c0,line width=0.3985pt,draw opacity=0.5] (62.1002,46.1589) .. controls (62.1002,45.2675) and (61.3777,44.5449) ..(60.4863,44.5449) .. controls (59.5949,44.5449) and (58.8723,45.2675) ..(58.8723,46.1589) .. controls (58.8723,47.0502) and (59.5949,47.7728) ..(60.4863,47.7728) .. controls (61.3777,47.7728) and (62.1002,47.0502) ..(62.1002,46.1589);
\path[fill=chfourfig6c0,fill opacity=0.5,draw=chfourfig6c0,line width=0.3985pt,draw opacity=0.5] (85.5227,38.0890) .. controls (85.5227,37.1976) and (84.8001,36.4751) ..(83.9087,36.4751) .. controls (83.0174,36.4751) and (82.2948,37.1976) ..(82.2948,38.0890) .. controls (82.2948,38.9804) and (83.0174,39.7030) ..(83.9087,39.7030) .. controls (84.8001,39.7030) and (85.5227,38.9804) ..(85.5227,38.0890);
\path[fill=chfourfig6c0,fill opacity=0.5,draw=chfourfig6c0,line width=0.3985pt,draw opacity=0.5] (54.7677,38.0890) .. controls (54.7677,37.1976) and (54.0451,36.4751) ..(53.1537,36.4751) .. controls (52.2623,36.4751) and (51.5397,37.1976) ..(51.5397,38.0890) .. controls (51.5397,38.9804) and (52.2623,39.7030) ..(53.1537,39.7030) .. controls (54.0451,39.7030) and (54.7677,38.9804) ..(54.7677,38.0890);
\path[fill=chfourfig6c0,fill opacity=0.5,draw=chfourfig6c0,line width=0.3985pt,draw opacity=0.5] (67.7890,43.1402) .. controls (67.7890,42.2488) and (67.0664,41.5262) ..(66.1750,41.5262) .. controls (65.2836,41.5262) and (64.5610,42.2488) ..(64.5610,43.1402) .. controls (64.5610,44.0315) and (65.2836,44.7541) ..(66.1750,44.7541) .. controls (67.0664,44.7541) and (67.7890,44.0315) ..(67.7890,43.1402);
\path[fill=chfourfig6c0,fill opacity=0.5,draw=chfourfig6c0,line width=0.3985pt,draw opacity=0.5] (61.5224,39.0953) .. controls (61.5224,38.2039) and (60.7998,37.4813) ..(59.9084,37.4813) .. controls (59.0171,37.4813) and (58.2945,38.2039) ..(58.2945,39.0953) .. controls (58.2945,39.9866) and (59.0171,40.7092) ..(59.9084,40.7092) .. controls (60.7998,40.7092) and (61.5224,39.9866) ..(61.5224,39.0953);
\path[fill=chfourfig6c0,fill opacity=0.5,draw=chfourfig6c0,line width=0.3985pt,draw opacity=0.5] (55.1562,48.1813) .. controls (55.1562,47.2899) and (54.4336,46.5673) ..(53.5422,46.5673) .. controls (52.6509,46.5673) and (51.9283,47.2899) ..(51.9283,48.1813) .. controls (51.9283,49.0727) and (52.6509,49.7953) ..(53.5422,49.7953) .. controls (54.4336,49.7953) and (55.1562,49.0727) ..(55.1562,48.1813);
\path[fill=chfourfig6c0,fill opacity=0.5,draw=chfourfig6c0,line width=0.3985pt,draw opacity=0.5] (54.7677,42.1239) .. controls (54.7677,41.2326) and (54.0451,40.5100) ..(53.1537,40.5100) .. controls (52.2623,40.5100) and (51.5397,41.2326) ..(51.5397,42.1239) .. controls (51.5397,43.0153) and (52.2623,43.7379) ..(53.1537,43.7379) .. controls (54.0451,43.7379) and (54.7677,43.0153) ..(54.7677,42.1239);
\path[fill=chfourfig6c0,fill opacity=0.5,draw=chfourfig6c0,line width=0.3985pt,draw opacity=0.5] (68.1676,46.1589) .. controls (68.1676,45.2675) and (67.4450,44.5449) ..(66.5536,44.5449) .. controls (65.6622,44.5449) and (64.9396,45.2675) ..(64.9396,46.1589) .. controls (64.9396,47.0502) and (65.6622,47.7728) ..(66.5536,47.7728) .. controls (67.4450,47.7728) and (68.1676,47.0502) ..(68.1676,46.1589);
\path[fill=chfourfig6c0,fill opacity=0.5,draw=chfourfig6c0,line width=0.3985pt,draw opacity=0.5] (53.1338,51.2100) .. controls (53.1338,50.3186) and (52.4112,49.5960) ..(51.5198,49.5960) .. controls (50.6284,49.5960) and (49.9058,50.3186) ..(49.9058,51.2100) .. controls (49.9058,52.1014) and (50.6284,52.8240) ..(51.5198,52.8240) .. controls (52.4112,52.8240) and (53.1338,52.1014) ..(53.1338,51.2100);
\path[fill=chfourfig6c0,fill opacity=0.5,draw=chfourfig6c0,line width=0.3985pt,draw opacity=0.5] (54.2894,46.1589) .. controls (54.2894,45.2675) and (53.5668,44.5449) ..(52.6754,44.5449) .. controls (51.7841,44.5449) and (51.0615,45.2675) ..(51.0615,46.1589) .. controls (51.0615,47.0502) and (51.7841,47.7728) ..(52.6754,47.7728) .. controls (53.5668,47.7728) and (54.2894,47.0502) ..(54.2894,46.1589);
\path[fill=chfourfig6c0,fill opacity=0.5,draw=chfourfig6c0,line width=0.3985pt,draw opacity=0.5] (95.4556,53.2324) .. controls (95.4556,52.3411) and (94.7330,51.6185) ..(93.8416,51.6185) .. controls (92.9503,51.6185) and (92.2277,52.3411) ..(92.2277,53.2324) .. controls (92.2277,54.1238) and (92.9503,54.8464) ..(93.8416,54.8464) .. controls (94.7330,54.8464) and (95.4556,54.1238) ..(95.4556,53.2324);
\path[fill=chfourfig6c0,fill opacity=0.5,draw=chfourfig6c0,line width=0.3985pt,draw opacity=0.5] (58.2446,38.0890) .. controls (58.2446,37.1976) and (57.5220,36.4751) ..(56.6307,36.4751) .. controls (55.7393,36.4751) and (55.0167,37.1976) ..(55.0167,38.0890) .. controls (55.0167,38.9804) and (55.7393,39.7030) ..(56.6307,39.7030) .. controls (57.5220,39.7030) and (58.2446,38.9804) ..(58.2446,38.0890);
\path[fill=chfourfig6c0,fill opacity=0.5,draw=chfourfig6c0,line width=0.3985pt,draw opacity=0.5] (53.7116,43.1402) .. controls (53.7116,42.2488) and (52.9890,41.5262) ..(52.0976,41.5262) .. controls (51.2063,41.5262) and (50.4837,42.2488) ..(50.4837,43.1402) .. controls (50.4837,44.0315) and (51.2063,44.7541) ..(52.0976,44.7541) .. controls (52.9890,44.7541) and (53.7116,44.0315) ..(53.7116,43.1402);
\path[fill=chfourfig6c0,fill opacity=0.5,draw=chfourfig6c0,line width=0.3985pt,draw opacity=0.5] (59.9782,45.1526) .. controls (59.9782,44.2613) and (59.2556,43.5387) ..(58.3642,43.5387) .. controls (57.4728,43.5387) and (56.7502,44.2613) ..(56.7502,45.1526) .. controls (56.7502,46.0440) and (57.4728,46.7666) ..(58.3642,46.7666) .. controls (59.2556,46.7666) and (59.9782,46.0440) ..(59.9782,45.1526);
\path[fill=chfourfig6c0,fill opacity=0.5,draw=chfourfig6c0,line width=0.3985pt,draw opacity=0.5] (93.0446,54.2387) .. controls (93.0446,53.3473) and (92.3220,52.6247) ..(91.4306,52.6247) .. controls (90.5393,52.6247) and (89.8167,53.3473) ..(89.8167,54.2387) .. controls (89.8167,55.1300) and (90.5393,55.8526) ..(91.4306,55.8526) .. controls (92.3220,55.8526) and (93.0446,55.1300) ..(93.0446,54.2387);
\path[fill=chfourfig6c0,fill opacity=0.5,draw=chfourfig6c0,line width=0.3985pt,draw opacity=0.5] (61.9010,30.0192) .. controls (61.9010,29.1278) and (61.1784,28.4052) ..(60.2870,28.4052) .. controls (59.3956,28.4052) and (58.6731,29.1278) ..(58.6731,30.0192) .. controls (58.6731,30.9106) and (59.3956,31.6331) ..(60.2870,31.6331) .. controls (61.1784,31.6331) and (61.9010,30.9106) ..(61.9010,30.0192);
\path[fill=chfourfig6c0,fill opacity=0.5,draw=chfourfig6c0,line width=0.3985pt,draw opacity=0.5] (70.5786,43.1402) .. controls (70.5786,42.2488) and (69.8560,41.5262) ..(68.9646,41.5262) .. controls (68.0732,41.5262) and (67.3506,42.2488) ..(67.3506,43.1402) .. controls (67.3506,44.0315) and (68.0732,44.7541) ..(68.9646,44.7541) .. controls (69.8560,44.7541) and (70.5786,44.0315) ..(70.5786,43.1402);
\path[fill=chfourfig6c0,fill opacity=0.5,draw=chfourfig6c0,line width=0.3985pt,draw opacity=0.5] (105.3785,64.3310) .. controls (105.3785,63.4396) and (104.6559,62.7170) ..(103.7645,62.7170) .. controls (102.8732,62.7170) and (102.1506,63.4396) ..(102.1506,64.3310) .. controls (102.1506,65.2223) and (102.8732,65.9449) ..(103.7645,65.9449) .. controls (104.6559,65.9449) and (105.3785,65.2223) ..(105.3785,64.3310);
\path[fill=chfourfig6c0,fill opacity=0.5,draw=chfourfig6c0,line width=0.3985pt,draw opacity=0.5] (96.0334,63.3247) .. controls (96.0334,62.4333) and (95.3108,61.7107) ..(94.4195,61.7107) .. controls (93.5281,61.7107) and (92.8055,62.4333) ..(92.8055,63.3247) .. controls (92.8055,64.2161) and (93.5281,64.9387) ..(94.4195,64.9387) .. controls (95.3108,64.9387) and (96.0334,64.2161) ..(96.0334,63.3247);
\path[fill=chfourfig6c0,fill opacity=0.5,draw=chfourfig6c0,line width=0.3985pt,draw opacity=0.5] (60.0778,52.2162) .. controls (60.0778,51.3249) and (59.3552,50.6023) ..(58.4638,50.6023) .. controls (57.5725,50.6023) and (56.8499,51.3249) ..(56.8499,52.2162) .. controls (56.8499,53.1076) and (57.5725,53.8302) ..(58.4638,53.8302) .. controls (59.3552,53.8302) and (60.0778,53.1076) ..(60.0778,52.2162);
\path[fill=chfourfig6c0,fill opacity=0.5,draw=chfourfig6c0,line width=0.3985pt,draw opacity=0.5] (82.8228,55.2449) .. controls (82.8228,54.3535) and (82.1002,53.6309) ..(81.2088,53.6309) .. controls (80.3175,53.6309) and (79.5949,54.3535) ..(79.5949,55.2449) .. controls (79.5949,56.1363) and (80.3175,56.8589) ..(81.2088,56.8589) .. controls (82.1002,56.8589) and (82.8228,56.1363) ..(82.8228,55.2449);
\path[fill=chfourfig6c0,fill opacity=0.5,draw=chfourfig6c0,line width=0.3985pt,draw opacity=0.5] (112.5118,57.2673) .. controls (112.5118,56.3760) and (111.7892,55.6534) ..(110.8979,55.6534) .. controls (110.0065,55.6534) and (109.2839,56.3760) ..(109.2839,57.2673) .. controls (109.2839,58.1587) and (110.0065,58.8813) ..(110.8979,58.8813) .. controls (111.7892,58.8813) and (112.5118,58.1587) ..(112.5118,57.2673);
\path[fill=chfourfig6c0,fill opacity=0.5,draw=chfourfig6c0,line width=0.3985pt,draw opacity=0.5] (56.7004,48.1813) .. controls (56.7004,47.2899) and (55.9778,46.5673) ..(55.0864,46.5673) .. controls (54.1951,46.5673) and (53.4725,47.2899) ..(53.4725,48.1813) .. controls (53.4725,49.0727) and (54.1951,49.7953) ..(55.0864,49.7953) .. controls (55.9778,49.7953) and (56.7004,49.0727) ..(56.7004,48.1813);
\path[fill=chfourfig6c0,fill opacity=0.5,draw=chfourfig6c0,line width=0.3985pt,draw opacity=0.5] (91.6897,52.2162) .. controls (91.6897,51.3249) and (90.9671,50.6023) ..(90.0757,50.6023) .. controls (89.1843,50.6023) and (88.4617,51.3249) ..(88.4617,52.2162) .. controls (88.4617,53.1076) and (89.1843,53.8302) ..(90.0757,53.8302) .. controls (90.9671,53.8302) and (91.6897,53.1076) ..(91.6897,52.2162);
\path[fill=chfourfig6c0,fill opacity=0.5,draw=chfourfig6c0,line width=0.3985pt,draw opacity=0.5] (78.9672,35.0603) .. controls (78.9672,34.1690) and (78.2446,33.4464) ..(77.3532,33.4464) .. controls (76.4619,33.4464) and (75.7393,34.1690) ..(75.7393,35.0603) .. controls (75.7393,35.9517) and (76.4619,36.6743) ..(77.3532,36.6743) .. controls (78.2446,36.6743) and (78.9672,35.9517) ..(78.9672,35.0603);
\path[fill=chfourfig6c0,fill opacity=0.5,draw=chfourfig6c0,line width=0.3985pt,draw opacity=0.5] (58.7229,49.1875) .. controls (58.7229,48.2962) and (58.0002,47.5736) ..(57.1089,47.5736) .. controls (56.2175,47.5736) and (55.4949,48.2962) ..(55.4949,49.1875) .. controls (55.4949,50.0789) and (56.2175,50.8015) ..(57.1089,50.8015) .. controls (58.0002,50.8015) and (58.7229,50.0789) ..(58.7229,49.1875);
\path[fill=chfourfig6c0,fill opacity=0.5,draw=chfourfig6c0,line width=0.3985pt,draw opacity=0.5] (102.3000,51.2100) .. controls (102.3000,50.3186) and (101.5774,49.5960) ..(100.6860,49.5960) .. controls (99.7947,49.5960) and (99.0721,50.3186) ..(99.0721,51.2100) .. controls (99.0721,52.1014) and (99.7947,52.8240) ..(100.6860,52.8240) .. controls (101.5774,52.8240) and (102.3000,52.1014) ..(102.3000,51.2100);
\path[fill=chfourfig6c0,fill opacity=0.5,draw=chfourfig6c0,line width=0.3985pt,draw opacity=0.5] (103.5453,54.2387) .. controls (103.5453,53.3473) and (102.8228,52.6247) ..(101.9314,52.6247) .. controls (101.0400,52.6247) and (100.3174,53.3473) ..(100.3174,54.2387) .. controls (100.3174,55.1300) and (101.0400,55.8526) ..(101.9314,55.8526) .. controls (102.8228,55.8526) and (103.5453,55.1300) ..(103.5453,54.2387);
\path[fill=chfourfig6c0,fill opacity=0.5,draw=chfourfig6c0,line width=0.3985pt,draw opacity=0.5] (102.3000,41.1177) .. controls (102.3000,40.2263) and (101.5774,39.5037) ..(100.6860,39.5037) .. controls (99.7947,39.5037) and (99.0721,40.2263) ..(99.0721,41.1177) .. controls (99.0721,42.0091) and (99.7947,42.7317) ..(100.6860,42.7317) .. controls (101.5774,42.7317) and (102.3000,42.0091) ..(102.3000,41.1177);
\path[fill=chfourfig6c0,fill opacity=0.5,draw=chfourfig6c0,line width=0.3985pt,draw opacity=0.5] (58.6232,40.1115) .. controls (58.6232,39.2201) and (57.9006,38.4975) ..(57.0093,38.4975) .. controls (56.1179,38.4975) and (55.3953,39.2201) ..(55.3953,40.1115) .. controls (55.3953,41.0028) and (56.1179,41.7254) ..(57.0093,41.7254) .. controls (57.9006,41.7254) and (58.6232,41.0028) ..(58.6232,40.1115);
\path[fill=chfourfig6c0,fill opacity=0.5,draw=chfourfig6c0,line width=0.3985pt,draw opacity=0.5] (60.4564,53.2324) .. controls (60.4564,52.3411) and (59.7338,51.6185) ..(58.8424,51.6185) .. controls (57.9510,51.6185) and (57.2285,52.3411) ..(57.2285,53.2324) .. controls (57.2285,54.1238) and (57.9510,54.8464) ..(58.8424,54.8464) .. controls (59.7338,54.8464) and (60.4564,54.1238) ..(60.4564,53.2324);
\path[fill=chfourfig6c0,fill opacity=0.5,draw=chfourfig6c0,line width=0.3985pt,draw opacity=0.5] (63.3555,34.0541) .. controls (63.3555,33.1627) and (62.6330,32.4401) ..(61.7416,32.4401) .. controls (60.8502,32.4401) and (60.1276,33.1627) ..(60.1276,34.0541) .. controls (60.1276,34.9455) and (60.8502,35.6681) ..(61.7416,35.6681) .. controls (62.6330,35.6681) and (63.3555,34.9455) ..(63.3555,34.0541);
\path[fill=chfourfig6c0,fill opacity=0.5,draw=chfourfig6c0,line width=0.3985pt,draw opacity=0.5] (100.8455,60.2960) .. controls (100.8455,59.4047) and (100.1229,58.6821) ..(99.2315,58.6821) .. controls (98.3401,58.6821) and (97.6175,59.4047) ..(97.6175,60.2960) .. controls (97.6175,61.1874) and (98.3401,61.9100) ..(99.2315,61.9100) .. controls (100.1229,61.9100) and (100.8455,61.1874) ..(100.8455,60.2960);
\path[fill=chfourfig6c0,fill opacity=0.5,draw=chfourfig6c0,line width=0.3985pt,draw opacity=0.5] (55.1562,48.1813) .. controls (55.1562,47.2899) and (54.4336,46.5673) ..(53.5422,46.5673) .. controls (52.6509,46.5673) and (51.9283,47.2899) ..(51.9283,48.1813) .. controls (51.9283,49.0727) and (52.6509,49.7953) ..(53.5422,49.7953) .. controls (54.4336,49.7953) and (55.1562,49.0727) ..(55.1562,48.1813);
\path[fill=chfourfig6c0,fill opacity=0.5,draw=chfourfig6c0,line width=0.3985pt,draw opacity=0.5] (56.3119,47.1751) .. controls (56.3119,46.2837) and (55.5893,45.5611) ..(54.6979,45.5611) .. controls (53.8065,45.5611) and (53.0839,46.2837) ..(53.0839,47.1751) .. controls (53.0839,48.0664) and (53.8065,48.7890) ..(54.6979,48.7890) .. controls (55.5893,48.7890) and (56.3119,48.0664) ..(56.3119,47.1751);
\path[fill=chfourfig6c0,fill opacity=0.5,draw=chfourfig6c0,line width=0.3985pt,draw opacity=0.5] (64.7005,31.0254) .. controls (64.7005,30.1340) and (63.9779,29.4115) ..(63.0866,29.4115) .. controls (62.1952,29.4115) and (61.4726,30.1340) ..(61.4726,31.0254) .. controls (61.4726,31.9168) and (62.1952,32.6394) ..(63.0866,32.6394) .. controls (63.9779,32.6394) and (64.7005,31.9168) ..(64.7005,31.0254);
\path[fill=chfourfig6c0,fill opacity=0.5,draw=chfourfig6c0,line width=0.3985pt,draw opacity=0.5] (58.8225,52.2162) .. controls (58.8225,51.3249) and (58.0999,50.6023) ..(57.2085,50.6023) .. controls (56.3172,50.6023) and (55.5946,51.3249) ..(55.5946,52.2162) .. controls (55.5946,53.1076) and (56.3172,53.8302) ..(57.2085,53.8302) .. controls (58.0999,53.8302) and (58.8225,53.1076) ..(58.8225,52.2162);
\path[fill=chfourfig6c0,fill opacity=0.5,draw=chfourfig6c0,line width=0.3985pt,draw opacity=0.5] (55.1562,42.1239) .. controls (55.1562,41.2326) and (54.4336,40.5100) ..(53.5422,40.5100) .. controls (52.6509,40.5100) and (51.9283,41.2326) ..(51.9283,42.1239) .. controls (51.9283,43.0153) and (52.6509,43.7379) ..(53.5422,43.7379) .. controls (54.4336,43.7379) and (55.1562,43.0153) ..(55.1562,42.1239);
\path[fill=chfourfig6c0,fill opacity=0.5,draw=chfourfig6c0,line width=0.3985pt,draw opacity=0.5] (93.6224,52.2162) .. controls (93.6224,51.3249) and (92.8998,50.6023) ..(92.0085,50.6023) .. controls (91.1171,50.6023) and (90.3945,51.3249) ..(90.3945,52.2162) .. controls (90.3945,53.1076) and (91.1171,53.8302) ..(92.0085,53.8302) .. controls (92.8998,53.8302) and (93.6224,53.1076) ..(93.6224,52.2162);
\path[fill=chfourfig6c0,fill opacity=0.5,draw=chfourfig6c0,line width=0.3985pt,draw opacity=0.5] (53.7116,42.1239) .. controls (53.7116,41.2326) and (52.9890,40.5100) ..(52.0976,40.5100) .. controls (51.2063,40.5100) and (50.4837,41.2326) ..(50.4837,42.1239) .. controls (50.4837,43.0153) and (51.2063,43.7379) ..(52.0976,43.7379) .. controls (52.9890,43.7379) and (53.7116,43.0153) ..(53.7116,42.1239);
\path[fill=chfourfig6c0,fill opacity=0.5,draw=chfourfig6c0,line width=0.3985pt,draw opacity=0.5] (113.0897,65.3372) .. controls (113.0897,64.4458) and (112.3671,63.7232) ..(111.4757,63.7232) .. controls (110.5844,63.7232) and (109.8618,64.4458) ..(109.8618,65.3372) .. controls (109.8618,66.2286) and (110.5844,66.9512) ..(111.4757,66.9512) .. controls (112.3671,66.9512) and (113.0897,66.2286) ..(113.0897,65.3372);
\path[fill=chfourfig6c0,fill opacity=0.5,draw=chfourfig6c0,line width=0.3985pt,draw opacity=0.5] (82.9224,44.1464) .. controls (82.9224,43.2550) and (82.1998,42.5324) ..(81.3085,42.5324) .. controls (80.4171,42.5324) and (79.6945,43.2550) ..(79.6945,44.1464) .. controls (79.6945,45.0378) and (80.4171,45.7603) ..(81.3085,45.7603) .. controls (82.1998,45.7603) and (82.9224,45.0378) ..(82.9224,44.1464);
\path[fill=chfourfig6c0,fill opacity=0.5,draw=chfourfig6c0,line width=0.3985pt,draw opacity=0.5] (79.6447,44.1464) .. controls (79.6447,43.2550) and (78.9221,42.5324) ..(78.0307,42.5324) .. controls (77.1393,42.5324) and (76.4167,43.2550) ..(76.4167,44.1464) .. controls (76.4167,45.0378) and (77.1393,45.7603) ..(78.0307,45.7603) .. controls (78.9221,45.7603) and (79.6447,45.0378) ..(79.6447,44.1464);
\path[fill=chfourfig6c0,fill opacity=0.5,draw=chfourfig6c0,line width=0.3985pt,draw opacity=0.5] (99.8890,55.2449) .. controls (99.8890,54.3535) and (99.1664,53.6309) ..(98.2751,53.6309) .. controls (97.3837,53.6309) and (96.6611,54.3535) ..(96.6611,55.2449) .. controls (96.6611,56.1363) and (97.3837,56.8589) ..(98.2751,56.8589) .. controls (99.1664,56.8589) and (99.8890,56.1363) ..(99.8890,55.2449);
\path[fill=chfourfig6c0,fill opacity=0.5,draw=chfourfig6c0,line width=0.3985pt,draw opacity=0.5] (130.3452,71.3946) .. controls (130.3452,70.5032) and (129.6226,69.7806) ..(128.7312,69.7806) .. controls (127.8399,69.7806) and (127.1173,70.5032) ..(127.1173,71.3946) .. controls (127.1173,72.2859) and (127.8399,73.0085) ..(128.7312,73.0085) .. controls (129.6226,73.0085) and (130.3452,72.2859) ..(130.3452,71.3946);
\path[fill=chfourfig6c0,fill opacity=0.5,draw=chfourfig6c0,line width=0.3985pt,draw opacity=0.5] (79.2561,41.1177) .. controls (79.2561,40.2263) and (78.5335,39.5037) ..(77.6422,39.5037) .. controls (76.7508,39.5037) and (76.0282,40.2263) ..(76.0282,41.1177) .. controls (76.0282,42.0091) and (76.7508,42.7317) ..(77.6422,42.7317) .. controls (78.5335,42.7317) and (79.2561,42.0091) ..(79.2561,41.1177);
\path[fill=chfourfig6c0,fill opacity=0.5,draw=chfourfig6c0,line width=0.3985pt,draw opacity=0.5] (63.2559,29.0030) .. controls (63.2559,28.1116) and (62.5333,27.3890) ..(61.6419,27.3890) .. controls (60.7506,27.3890) and (60.0280,28.1116) ..(60.0280,29.0030) .. controls (60.0280,29.8943) and (60.7506,30.6169) ..(61.6419,30.6169) .. controls (62.5333,30.6169) and (63.2559,29.8943) ..(63.2559,29.0030);
\path[fill=chfourfig6c0,fill opacity=0.5,draw=chfourfig6c0,line width=0.3985pt,draw opacity=0.5] (67.1115,46.1589) .. controls (67.1115,45.2675) and (66.3889,44.5449) ..(65.4976,44.5449) .. controls (64.6062,44.5449) and (63.8836,45.2675) ..(63.8836,46.1589) .. controls (63.8836,47.0502) and (64.6062,47.7728) ..(65.4976,47.7728) .. controls (66.3889,47.7728) and (67.1115,47.0502) ..(67.1115,46.1589);
\path[fill=chfourfig6c0,fill opacity=0.5,draw=chfourfig6c0,line width=0.3985pt,draw opacity=0.5] (61.8113,41.1177) .. controls (61.8113,40.2263) and (61.0887,39.5037) ..(60.1974,39.5037) .. controls (59.3060,39.5037) and (58.5834,40.2263) ..(58.5834,41.1177) .. controls (58.5834,42.0091) and (59.3060,42.7317) ..(60.1974,42.7317) .. controls (61.0887,42.7317) and (61.8113,42.0091) ..(61.8113,41.1177);
\path[fill=chfourfig6c0,fill opacity=0.5,draw=chfourfig6c0,line width=0.3985pt,draw opacity=0.5] (92.0782,52.2162) .. controls (92.0782,51.3249) and (91.3556,50.6023) ..(90.4642,50.6023) .. controls (89.5729,50.6023) and (88.8503,51.3249) ..(88.8503,52.2162) .. controls (88.8503,53.1076) and (89.5729,53.8302) ..(90.4642,53.8302) .. controls (91.3556,53.8302) and (92.0782,53.1076) ..(92.0782,52.2162);
\path[fill=chfourfig6c0,fill opacity=0.5,draw=chfourfig6c0,line width=0.3985pt,draw opacity=0.5] (94.5789,56.2611) .. controls (94.5789,55.3697) and (93.8563,54.6471) ..(92.9649,54.6471) .. controls (92.0735,54.6471) and (91.3509,55.3697) ..(91.3509,56.2611) .. controls (91.3509,57.1525) and (92.0735,57.8751) ..(92.9649,57.8751) .. controls (93.8563,57.8751) and (94.5789,57.1525) ..(94.5789,56.2611);
\path[fill=chfourfig6c0,fill opacity=0.5,draw=chfourfig6c0,line width=0.3985pt,draw opacity=0.5] (76.3669,41.1177) .. controls (76.3669,40.2263) and (75.6443,39.5037) ..(74.7530,39.5037) .. controls (73.8616,39.5037) and (73.1390,40.2263) ..(73.1390,41.1177) .. controls (73.1390,42.0091) and (73.8616,42.7317) ..(74.7530,42.7317) .. controls (75.6443,42.7317) and (76.3669,42.0091) ..(76.3669,41.1177);
\path[fill=chfourfig6c0,fill opacity=0.5,draw=chfourfig6c0,line width=0.3985pt,draw opacity=0.5] (89.8565,57.2673) .. controls (89.8565,56.3760) and (89.1339,55.6534) ..(88.2426,55.6534) .. controls (87.3512,55.6534) and (86.6286,56.3760) ..(86.6286,57.2673) .. controls (86.6286,58.1587) and (87.3512,58.8813) ..(88.2426,58.8813) .. controls (89.1339,58.8813) and (89.8565,58.1587) ..(89.8565,57.2673);
\path[fill=chfourfig6c0,fill opacity=0.5,draw=chfourfig6c0,line width=0.3985pt,draw opacity=0.5] (54.9669,48.1813) .. controls (54.9669,47.2899) and (54.2443,46.5673) ..(53.3529,46.5673) .. controls (52.4616,46.5673) and (51.7390,47.2899) ..(51.7390,48.1813) .. controls (51.7390,49.0727) and (52.4616,49.7953) ..(53.3529,49.7953) .. controls (54.2443,49.7953) and (54.9669,49.0727) ..(54.9669,48.1813);
\path[fill=chfourfig6c0,fill opacity=0.5,draw=chfourfig6c0,line width=0.3985pt,draw opacity=0.5] (60.4564,45.1526) .. controls (60.4564,44.2613) and (59.7338,43.5387) ..(58.8424,43.5387) .. controls (57.9510,43.5387) and (57.2285,44.2613) ..(57.2285,45.1526) .. controls (57.2285,46.0440) and (57.9510,46.7666) ..(58.8424,46.7666) .. controls (59.7338,46.7666) and (60.4564,46.0440) ..(60.4564,45.1526);
\path[fill=chfourfig6c0,fill opacity=0.5,draw=chfourfig6c0,line width=0.3985pt,draw opacity=0.5] (56.5111,53.2324) .. controls (56.5111,52.3411) and (55.7885,51.6185) ..(54.8972,51.6185) .. controls (54.0058,51.6185) and (53.2832,52.3411) ..(53.2832,53.2324) .. controls (53.2832,54.1238) and (54.0058,54.8464) ..(54.8972,54.8464) .. controls (55.7885,54.8464) and (56.5111,54.1238) ..(56.5111,53.2324);
\path[fill=chfourfig6c0,fill opacity=0.5,draw=chfourfig6c0,line width=0.3985pt,draw opacity=0.5] (64.5112,44.1464) .. controls (64.5112,43.2550) and (63.7886,42.5324) ..(62.8973,42.5324) .. controls (62.0059,42.5324) and (61.2833,43.2550) ..(61.2833,44.1464) .. controls (61.2833,45.0378) and (62.0059,45.7603) ..(62.8973,45.7603) .. controls (63.7886,45.7603) and (64.5112,45.0378) ..(64.5112,44.1464);
\path[fill=chfourfig6c0,fill opacity=0.5,draw=chfourfig6c0,line width=0.3985pt,draw opacity=0.5] (210.7448,98.6527) .. controls (210.7448,97.7613) and (210.0222,97.0387) ..(209.1308,97.0387) .. controls (208.2394,97.0387) and (207.5168,97.7613) ..(207.5168,98.6527) .. controls (207.5168,99.5441) and (208.2394,100.2667) ..(209.1308,100.2667) .. controls (210.0222,100.2667) and (210.7448,99.5441) ..(210.7448,98.6527);
\path[fill=chfourfig6c0,fill opacity=0.5,draw=chfourfig6c0,line width=0.3985pt,draw opacity=0.5] (63.1563,36.0666) .. controls (63.1563,35.1752) and (62.4337,34.4526) ..(61.5423,34.4526) .. controls (60.6510,34.4526) and (59.9284,35.1752) ..(59.9284,36.0666) .. controls (59.9284,36.9579) and (60.6510,37.6805) ..(61.5423,37.6805) .. controls (62.4337,37.6805) and (63.1563,36.9579) ..(63.1563,36.0666);
\path[fill=chfourfig6c0,fill opacity=0.5,draw=chfourfig6c0,line width=0.3985pt,draw opacity=0.5] (87.9337,50.2038) .. controls (87.9337,49.3124) and (87.2111,48.5898) ..(86.3197,48.5898) .. controls (85.4284,48.5898) and (84.7058,49.3124) ..(84.7058,50.2038) .. controls (84.7058,51.0951) and (85.4284,51.8177) ..(86.3197,51.8177) .. controls (87.2111,51.8177) and (87.9337,51.0951) ..(87.9337,50.2038);
\path[fill=chfourfig6c0,fill opacity=0.5,draw=chfourfig6c0,line width=0.3985pt,draw opacity=0.5] (97.9563,52.2162) .. controls (97.9563,51.3249) and (97.2336,50.6023) ..(96.3423,50.6023) .. controls (95.4509,50.6023) and (94.7283,51.3249) ..(94.7283,52.2162) .. controls (94.7283,53.1076) and (95.4509,53.8302) ..(96.3423,53.8302) .. controls (97.2336,53.8302) and (97.9563,53.1076) ..(97.9563,52.2162);
\path[fill=chfourfig6c0,fill opacity=0.5,draw=chfourfig6c0,line width=0.3985pt,draw opacity=0.5] (119.1670,70.3883) .. controls (119.1670,69.4970) and (118.4444,68.7744) ..(117.5530,68.7744) .. controls (116.6616,68.7744) and (115.9391,69.4970) ..(115.9391,70.3883) .. controls (115.9391,71.2797) and (116.6616,72.0023) ..(117.5530,72.0023) .. controls (118.4444,72.0023) and (119.1670,71.2797) ..(119.1670,70.3883);
\path[fill=chfourfig6c0,fill opacity=0.5,draw=chfourfig6c0,line width=0.3985pt,draw opacity=0.5] (65.2784,35.0603) .. controls (65.2784,34.1690) and (64.5558,33.4464) ..(63.6644,33.4464) .. controls (62.7730,33.4464) and (62.0504,34.1690) ..(62.0504,35.0603) .. controls (62.0504,35.9517) and (62.7730,36.6743) ..(63.6644,36.6743) .. controls (64.5558,36.6743) and (65.2784,35.9517) ..(65.2784,35.0603);
\path[fill=chfourfig6c0,fill opacity=0.5,draw=chfourfig6c0,line width=0.3985pt,draw opacity=0.5] (80.3221,32.0317) .. controls (80.3221,31.1403) and (79.5995,30.4177) ..(78.7082,30.4177) .. controls (77.8168,30.4177) and (77.0942,31.1403) ..(77.0942,32.0317) .. controls (77.0942,32.9230) and (77.8168,33.6456) ..(78.7082,33.6456) .. controls (79.5995,33.6456) and (80.3221,32.9230) ..(80.3221,32.0317);
\path[fill=chfourfig6c0,fill opacity=0.5,draw=chfourfig6c0,line width=0.3985pt,draw opacity=0.5] (63.8338,34.0541) .. controls (63.8338,33.1627) and (63.1112,32.4401) ..(62.2198,32.4401) .. controls (61.3284,32.4401) and (60.6058,33.1627) ..(60.6058,34.0541) .. controls (60.6058,34.9455) and (61.3284,35.6681) ..(62.2198,35.6681) .. controls (63.1112,35.6681) and (63.8338,34.9455) ..(63.8338,34.0541);
\path[fill=chfourfig6c0,fill opacity=0.5,draw=chfourfig6c0,line width=0.3985pt,draw opacity=0.5] (95.4556,49.1875) .. controls (95.4556,48.2962) and (94.7330,47.5736) ..(93.8416,47.5736) .. controls (92.9503,47.5736) and (92.2277,48.2962) ..(92.2277,49.1875) .. controls (92.2277,50.0789) and (92.9503,50.8015) ..(93.8416,50.8015) .. controls (94.7330,50.8015) and (95.4556,50.0789) ..(95.4556,49.1875);
\path[fill=chfourfig6c0,fill opacity=0.5,draw=chfourfig6c0,line width=0.3985pt,draw opacity=0.5] (53.0341,48.1813) .. controls (53.0341,47.2899) and (52.3115,46.5673) ..(51.4202,46.5673) .. controls (50.5288,46.5673) and (49.8062,47.2899) ..(49.8062,48.1813) .. controls (49.8062,49.0727) and (50.5288,49.7953) ..(51.4202,49.7953) .. controls (52.3115,49.7953) and (53.0341,49.0727) ..(53.0341,48.1813);
\path[fill=chfourfig6c0,fill opacity=0.5,draw=chfourfig6c0,line width=0.3985pt,draw opacity=0.5] (61.0342,42.1239) .. controls (61.0342,41.2326) and (60.3116,40.5100) ..(59.4203,40.5100) .. controls (58.5289,40.5100) and (57.8063,41.2326) ..(57.8063,42.1239) .. controls (57.8063,43.0153) and (58.5289,43.7379) ..(59.4203,43.7379) .. controls (60.3116,43.7379) and (61.0342,43.0153) ..(61.0342,42.1239);
\path[fill=chfourfig6c0,fill opacity=0.5,draw=chfourfig6c0,line width=0.3985pt,draw opacity=0.5] (124.6565,64.3310) .. controls (124.6565,63.4396) and (123.9339,62.7170) ..(123.0425,62.7170) .. controls (122.1511,62.7170) and (121.4285,63.4396) ..(121.4285,64.3310) .. controls (121.4285,65.2223) and (122.1511,65.9449) ..(123.0425,65.9449) .. controls (123.9339,65.9449) and (124.6565,65.2223) ..(124.6565,64.3310);
\path[fill=chfourfig6c0,fill opacity=0.5,draw=chfourfig6c0,line width=0.3985pt,draw opacity=0.5] (109.3337,59.2898) .. controls (109.3337,58.3984) and (108.6111,57.6758) ..(107.7198,57.6758) .. controls (106.8284,57.6758) and (106.1058,58.3984) ..(106.1058,59.2898) .. controls (106.1058,60.1812) and (106.8284,60.9038) ..(107.7198,60.9038) .. controls (108.6111,60.9038) and (109.3337,60.1812) ..(109.3337,59.2898);
\path[fill=chfourfig6c0,fill opacity=0.5,draw=chfourfig6c0,line width=0.3985pt,draw opacity=0.5] (119.3562,72.4108) .. controls (119.3562,71.5194) and (118.6337,70.7968) ..(117.7423,70.7968) .. controls (116.8509,70.7968) and (116.1283,71.5194) ..(116.1283,72.4108) .. controls (116.1283,73.3021) and (116.8509,74.0247) ..(117.7423,74.0247) .. controls (118.6337,74.0247) and (119.3562,73.3021) ..(119.3562,72.4108);
\path[fill=chfourfig6c0,fill opacity=0.5,draw=chfourfig6c0,line width=0.3985pt,draw opacity=0.5] (83.5003,43.1402) .. controls (83.5003,42.2488) and (82.7777,41.5262) ..(81.8863,41.5262) .. controls (80.9949,41.5262) and (80.2723,42.2488) ..(80.2723,43.1402) .. controls (80.2723,44.0315) and (80.9949,44.7541) ..(81.8863,44.7541) .. controls (82.7777,44.7541) and (83.5003,44.0315) ..(83.5003,43.1402);
\path[fill=chfourfig6c0,fill opacity=0.5,draw=chfourfig6c0,line width=0.3985pt,draw opacity=0.5] (71.2560,27.9967) .. controls (71.2560,27.1054) and (70.5334,26.3828) ..(69.6421,26.3828) .. controls (68.7507,26.3828) and (68.0281,27.1054) ..(68.0281,27.9967) .. controls (68.0281,28.8881) and (68.7507,29.6107) ..(69.6421,29.6107) .. controls (70.5334,29.6107) and (71.2560,28.8881) ..(71.2560,27.9967);
\path[fill=chfourfig6c0,fill opacity=0.5,draw=chfourfig6c0,line width=0.3985pt,draw opacity=0.5] (84.7556,60.2960) .. controls (84.7556,59.4047) and (84.0330,58.6821) ..(83.1416,58.6821) .. controls (82.2502,58.6821) and (81.5276,59.4047) ..(81.5276,60.2960) .. controls (81.5276,61.1874) and (82.2502,61.9100) ..(83.1416,61.9100) .. controls (84.0330,61.9100) and (84.7556,61.1874) ..(84.7556,60.2960);
\path[fill=chfourfig6c0,fill opacity=0.5,draw=chfourfig6c0,line width=0.3985pt,draw opacity=0.5] (65.6669,40.1115) .. controls (65.6669,39.2201) and (64.9443,38.4975) ..(64.0529,38.4975) .. controls (63.1616,38.4975) and (62.4390,39.2201) ..(62.4390,40.1115) .. controls (62.4390,41.0028) and (63.1616,41.7254) ..(64.0529,41.7254) .. controls (64.9443,41.7254) and (65.6669,41.0028) ..(65.6669,40.1115);
\path[fill=chfourfig6c0,fill opacity=0.5,draw=chfourfig6c0,line width=0.3985pt,draw opacity=0.5] (129.2892,70.3883) .. controls (129.2892,69.4970) and (128.5666,68.7744) ..(127.6752,68.7744) .. controls (126.7838,68.7744) and (126.0612,69.4970) ..(126.0612,70.3883) .. controls (126.0612,71.2797) and (126.7838,72.0023) ..(127.6752,72.0023) .. controls (128.5666,72.0023) and (129.2892,71.2797) ..(129.2892,70.3883);
\path[fill=chfourfig6c0,fill opacity=0.5,draw=chfourfig6c0,line width=0.3985pt,draw opacity=0.5] (53.1338,47.1751) .. controls (53.1338,46.2837) and (52.4112,45.5611) ..(51.5198,45.5611) .. controls (50.6284,45.5611) and (49.9058,46.2837) ..(49.9058,47.1751) .. controls (49.9058,48.0664) and (50.6284,48.7890) ..(51.5198,48.7890) .. controls (52.4112,48.7890) and (53.1338,48.0664) ..(53.1338,47.1751);
\path[fill=chfourfig6c0,fill opacity=0.5,draw=chfourfig6c0,line width=0.3985pt,draw opacity=0.5] (67.5897,40.1115) .. controls (67.5897,39.2201) and (66.8671,38.4975) ..(65.9758,38.4975) .. controls (65.0844,38.4975) and (64.3618,39.2201) ..(64.3618,40.1115) .. controls (64.3618,41.0028) and (65.0844,41.7254) ..(65.9758,41.7254) .. controls (66.8671,41.7254) and (67.5897,41.0028) ..(67.5897,40.1115);
\path[fill=chfourfig6c0,fill opacity=0.5,draw=chfourfig6c0,line width=0.3985pt,draw opacity=0.5] (67.5897,46.1589) .. controls (67.5897,45.2675) and (66.8671,44.5449) ..(65.9758,44.5449) .. controls (65.0844,44.5449) and (64.3618,45.2675) ..(64.3618,46.1589) .. controls (64.3618,47.0502) and (65.0844,47.7728) ..(65.9758,47.7728) .. controls (66.8671,47.7728) and (67.5897,47.0502) ..(67.5897,46.1589);
\path[fill=chfourfig6c0,fill opacity=0.5,draw=chfourfig6c0,line width=0.3985pt,draw opacity=0.5] (111.3562,79.4744) .. controls (111.3562,78.5830) and (110.6336,77.8604) ..(109.7422,77.8604) .. controls (108.8508,77.8604) and (108.1282,78.5830) ..(108.1282,79.4744) .. controls (108.1282,80.3657) and (108.8508,81.0883) ..(109.7422,81.0883) .. controls (110.6336,81.0883) and (111.3562,80.3657) ..(111.3562,79.4744);
\path[fill=chfourfig6c0,fill opacity=0.5,draw=chfourfig6c0,line width=0.3985pt,draw opacity=0.5] (56.8000,52.2162) .. controls (56.8000,51.3249) and (56.0775,50.6023) ..(55.1861,50.6023) .. controls (54.2947,50.6023) and (53.5721,51.3249) ..(53.5721,52.2162) .. controls (53.5721,53.1076) and (54.2947,53.8302) ..(55.1861,53.8302) .. controls (56.0775,53.8302) and (56.8000,53.1076) ..(56.8000,52.2162);
\path[fill=chfourfig6c0,fill opacity=0.5,draw=chfourfig6c0,line width=0.3985pt,draw opacity=0.5] (69.0343,40.1115) .. controls (69.0343,39.2201) and (68.3117,38.4975) ..(67.4203,38.4975) .. controls (66.5290,38.4975) and (65.8064,39.2201) ..(65.8064,40.1115) .. controls (65.8064,41.0028) and (66.5290,41.7254) ..(67.4203,41.7254) .. controls (68.3117,41.7254) and (69.0343,41.0028) ..(69.0343,40.1115);
\path[fill=chfourfig6c0,fill opacity=0.5,draw=chfourfig6c0,line width=0.3985pt,draw opacity=0.5] (148.3778,76.4457) .. controls (148.3778,75.5543) and (147.6552,74.8317) ..(146.7639,74.8317) .. controls (145.8725,74.8317) and (145.1499,75.5543) ..(145.1499,76.4457) .. controls (145.1499,77.3370) and (145.8725,78.0596) ..(146.7639,78.0596) .. controls (147.6552,78.0596) and (148.3778,77.3370) ..(148.3778,76.4457);
\path[fill=chfourfig6c0,fill opacity=0.5,draw=chfourfig6c0,line width=0.3985pt,draw opacity=0.5] (142.9780,88.5504) .. controls (142.9780,87.6591) and (142.2554,86.9365) ..(141.3640,86.9365) .. controls (140.4727,86.9365) and (139.7501,87.6591) ..(139.7501,88.5504) .. controls (139.7501,89.4418) and (140.4727,90.1644) ..(141.3640,90.1644) .. controls (142.2554,90.1644) and (142.9780,89.4418) ..(142.9780,88.5504);
\path[fill=chfourfig6c0,fill opacity=0.5,draw=chfourfig6c0,line width=0.3985pt,draw opacity=0.5] (113.8668,58.2736) .. controls (113.8668,57.3822) and (113.1442,56.6596) ..(112.2528,56.6596) .. controls (111.3615,56.6596) and (110.6389,57.3822) ..(110.6389,58.2736) .. controls (110.6389,59.1650) and (111.3615,59.8876) ..(112.2528,59.8876) .. controls (113.1442,59.8876) and (113.8668,59.1650) ..(113.8668,58.2736);
\path[fill=chfourfig6c0,fill opacity=0.5,draw=chfourfig6c0,line width=0.3985pt,draw opacity=0.5] (112.0336,60.2960) .. controls (112.0336,59.4047) and (111.3110,58.6821) ..(110.4197,58.6821) .. controls (109.5283,58.6821) and (108.8057,59.4047) ..(108.8057,60.2960) .. controls (108.8057,61.1874) and (109.5283,61.9100) ..(110.4197,61.9100) .. controls (111.3110,61.9100) and (112.0336,61.1874) ..(112.0336,60.2960);
\path[fill=chfourfig6c0,fill opacity=0.5,draw=chfourfig6c0,line width=0.3985pt,draw opacity=0.5] (78.2001,50.2038) .. controls (78.2001,49.3124) and (77.4775,48.5898) ..(76.5861,48.5898) .. controls (75.6947,48.5898) and (74.9722,49.3124) ..(74.9722,50.2038) .. controls (74.9722,51.0951) and (75.6947,51.8177) ..(76.5861,51.8177) .. controls (77.4775,51.8177) and (78.2001,51.0951) ..(78.2001,50.2038);
\path[fill=chfourfig6c0,fill opacity=0.5,draw=chfourfig6c0,line width=0.3985pt,draw opacity=0.5] (101.2340,70.3883) .. controls (101.2340,69.4970) and (100.5114,68.7744) ..(99.6200,68.7744) .. controls (98.7287,68.7744) and (98.0061,69.4970) ..(98.0061,70.3883) .. controls (98.0061,71.2797) and (98.7287,72.0023) ..(99.6200,72.0023) .. controls (100.5114,72.0023) and (101.2340,71.2797) ..(101.2340,70.3883);
\path[fill=chfourfig6c0,fill opacity=0.5,draw=chfourfig6c0,line width=0.3985pt,draw opacity=0.5] (87.8341,46.1589) .. controls (87.8341,45.2675) and (87.1115,44.5449) ..(86.2201,44.5449) .. controls (85.3287,44.5449) and (84.6061,45.2675) ..(84.6061,46.1589) .. controls (84.6061,47.0502) and (85.3287,47.7728) ..(86.2201,47.7728) .. controls (87.1115,47.7728) and (87.8341,47.0502) ..(87.8341,46.1589);
\path[fill=chfourfig6c0,fill opacity=0.5,draw=chfourfig6c0,line width=0.3985pt,draw opacity=0.5] (95.4556,55.2449) .. controls (95.4556,54.3535) and (94.7330,53.6309) ..(93.8416,53.6309) .. controls (92.9503,53.6309) and (92.2277,54.3535) ..(92.2277,55.2449) .. controls (92.2277,56.1363) and (92.9503,56.8589) ..(93.8416,56.8589) .. controls (94.7330,56.8589) and (95.4556,56.1363) ..(95.4556,55.2449);
\path[fill=chfourfig6c0,fill opacity=0.5,draw=chfourfig6c0,line width=0.3985pt,draw opacity=0.5] (129.0002,64.3310) .. controls (129.0002,63.4396) and (128.2776,62.7170) ..(127.3863,62.7170) .. controls (126.4949,62.7170) and (125.7723,63.4396) ..(125.7723,64.3310) .. controls (125.7723,65.2223) and (126.4949,65.9449) ..(127.3863,65.9449) .. controls (128.2776,65.9449) and (129.0002,65.2223) ..(129.0002,64.3310);
\path[fill=chfourfig6c0,fill opacity=0.5,draw=chfourfig6c0,line width=0.3985pt,draw opacity=0.5] (58.3343,48.1813) .. controls (58.3343,47.2899) and (57.6117,46.5673) ..(56.7203,46.5673) .. controls (55.8290,46.5673) and (55.1064,47.2899) ..(55.1064,48.1813) .. controls (55.1064,49.0727) and (55.8290,49.7953) ..(56.7203,49.7953) .. controls (57.6117,49.7953) and (58.3343,49.0727) ..(58.3343,48.1813);
\path[fill=chfourfig6c0,fill opacity=0.5,draw=chfourfig6c0,line width=0.3985pt,draw opacity=0.5] (54.6780,47.1751) .. controls (54.6780,46.2837) and (53.9554,45.5611) ..(53.0640,45.5611) .. controls (52.1727,45.5611) and (51.4501,46.2837) ..(51.4501,47.1751) .. controls (51.4501,48.0664) and (52.1727,48.7890) ..(53.0640,48.7890) .. controls (53.9554,48.7890) and (54.6780,48.0664) ..(54.6780,47.1751);
\path[fill=chfourfig6c0,fill opacity=0.5,draw=chfourfig6c0,line width=0.3985pt,draw opacity=0.5] (70.5786,47.1751) .. controls (70.5786,46.2837) and (69.8560,45.5611) ..(68.9646,45.5611) .. controls (68.0732,45.5611) and (67.3506,46.2837) ..(67.3506,47.1751) .. controls (67.3506,48.0664) and (68.0732,48.7890) ..(68.9646,48.7890) .. controls (69.8560,48.7890) and (70.5786,48.0664) ..(70.5786,47.1751);
\path[fill=chfourfig6c0,fill opacity=0.5,draw=chfourfig6c0,line width=0.3985pt,draw opacity=0.5] (111.5454,70.3883) .. controls (111.5454,69.4970) and (110.8228,68.7744) ..(109.9315,68.7744) .. controls (109.0401,68.7744) and (108.3175,69.4970) ..(108.3175,70.3883) .. controls (108.3175,71.2797) and (109.0401,72.0023) ..(109.9315,72.0023) .. controls (110.8228,72.0023) and (111.5454,71.2797) ..(111.5454,70.3883);
\path[fill=chfourfig6c0,fill opacity=0.5,draw=chfourfig6c0,line width=0.3985pt,draw opacity=0.5] (87.3559,52.2162) .. controls (87.3559,51.3249) and (86.6333,50.6023) ..(85.7419,50.6023) .. controls (84.8505,50.6023) and (84.1279,51.3249) ..(84.1279,52.2162) .. controls (84.1279,53.1076) and (84.8505,53.8302) ..(85.7419,53.8302) .. controls (86.6333,53.8302) and (87.3559,53.1076) ..(87.3559,52.2162);
\path[fill=chfourfig6c0,fill opacity=0.5,draw=chfourfig6c0,line width=0.3985pt,draw opacity=0.5] (60.0778,48.1813) .. controls (60.0778,47.2899) and (59.3552,46.5673) ..(58.4638,46.5673) .. controls (57.5725,46.5673) and (56.8499,47.2899) ..(56.8499,48.1813) .. controls (56.8499,49.0727) and (57.5725,49.7953) ..(58.4638,49.7953) .. controls (59.3552,49.7953) and (60.0778,49.0727) ..(60.0778,48.1813);
\path[fill=chfourfig6c0,fill opacity=0.5,draw=chfourfig6c0,line width=0.3985pt,draw opacity=0.5] (87.8341,58.2736) .. controls (87.8341,57.3822) and (87.1115,56.6596) ..(86.2201,56.6596) .. controls (85.3287,56.6596) and (84.6061,57.3822) ..(84.6061,58.2736) .. controls (84.6061,59.1650) and (85.3287,59.8876) ..(86.2201,59.8876) .. controls (87.1115,59.8876) and (87.8341,59.1650) ..(87.8341,58.2736);
\path[fill=chfourfig6c0,fill opacity=0.5,draw=chfourfig6c0,line width=0.3985pt,draw opacity=0.5] (66.8226,44.1464) .. controls (66.8226,43.2550) and (66.1000,42.5324) ..(65.2086,42.5324) .. controls (64.3172,42.5324) and (63.5947,43.2550) ..(63.5947,44.1464) .. controls (63.5947,45.0378) and (64.3172,45.7603) ..(65.2086,45.7603) .. controls (66.1000,45.7603) and (66.8226,45.0378) ..(66.8226,44.1464);
\path[fill=chfourfig6c0,fill opacity=0.5,draw=chfourfig6c0,line width=0.3985pt,draw opacity=0.5] (61.0342,43.1402) .. controls (61.0342,42.2488) and (60.3116,41.5262) ..(59.4203,41.5262) .. controls (58.5289,41.5262) and (57.8063,42.2488) ..(57.8063,43.1402) .. controls (57.8063,44.0315) and (58.5289,44.7541) ..(59.4203,44.7541) .. controls (60.3116,44.7541) and (61.0342,44.0315) ..(61.0342,43.1402);
\path[fill=chfourfig6c0,fill opacity=0.5,draw=chfourfig6c0,line width=0.3985pt,draw opacity=0.5] (123.4111,64.3310) .. controls (123.4111,63.4396) and (122.6885,62.7170) ..(121.7971,62.7170) .. controls (120.9058,62.7170) and (120.1832,63.4396) ..(120.1832,64.3310) .. controls (120.1832,65.2223) and (120.9058,65.9449) ..(121.7971,65.9449) .. controls (122.6885,65.9449) and (123.4111,65.2223) ..(123.4111,64.3310);
\path[fill=chfourfig6c0,fill opacity=0.5,draw=chfourfig6c0,line width=0.3985pt,draw opacity=0.5] (86.3895,49.1875) .. controls (86.3895,48.2962) and (85.6669,47.5736) ..(84.7755,47.5736) .. controls (83.8841,47.5736) and (83.1615,48.2962) ..(83.1615,49.1875) .. controls (83.1615,50.0789) and (83.8841,50.8015) ..(84.7755,50.8015) .. controls (85.6669,50.8015) and (86.3895,50.0789) ..(86.3895,49.1875);
\path[fill=chfourfig6c0,fill opacity=0.5,draw=chfourfig6c0,line width=0.3985pt,draw opacity=0.5] (80.2225,31.0254) .. controls (80.2225,30.1340) and (79.4999,29.4115) ..(78.6086,29.4115) .. controls (77.7172,29.4115) and (76.9946,30.1340) ..(76.9946,31.0254) .. controls (76.9946,31.9168) and (77.7172,32.6394) ..(78.6086,32.6394) .. controls (79.4999,32.6394) and (80.2225,31.9168) ..(80.2225,31.0254);
\path[fill=chfourfig6c0,fill opacity=0.5,draw=chfourfig6c0,line width=0.3985pt,draw opacity=0.5] (82.6335,37.0828) .. controls (82.6335,36.1914) and (81.9109,35.4688) ..(81.0195,35.4688) .. controls (80.1282,35.4688) and (79.4056,36.1914) ..(79.4056,37.0828) .. controls (79.4056,37.9742) and (80.1282,38.6967) ..(81.0195,38.6967) .. controls (81.9109,38.6967) and (82.6335,37.9742) ..(82.6335,37.0828);
\path[fill=chfourfig6c0,fill opacity=0.5,draw=chfourfig6c0,line width=0.3985pt,draw opacity=0.5] (91.3111,45.1526) .. controls (91.3111,44.2613) and (90.5885,43.5387) ..(89.6971,43.5387) .. controls (88.8058,43.5387) and (88.0832,44.2613) ..(88.0832,45.1526) .. controls (88.0832,46.0440) and (88.8058,46.7666) ..(89.6971,46.7666) .. controls (90.5885,46.7666) and (91.3111,46.0440) ..(91.3111,45.1526);
\path[fill=chfourfig6c0,fill opacity=0.5,draw=chfourfig6c0,line width=0.3985pt,draw opacity=0.5] (133.3340,58.2736) .. controls (133.3340,57.3822) and (132.6115,56.6596) ..(131.7201,56.6596) .. controls (130.8287,56.6596) and (130.1061,57.3822) ..(130.1061,58.2736) .. controls (130.1061,59.1650) and (130.8287,59.8876) ..(131.7201,59.8876) .. controls (132.6115,59.8876) and (133.3340,59.1650) ..(133.3340,58.2736);
\path[fill=chfourfig6c0,fill opacity=0.5,draw=chfourfig6c0,line width=0.3985pt,draw opacity=0.5] (57.0890,51.2100) .. controls (57.0890,50.3186) and (56.3664,49.5960) ..(55.4750,49.5960) .. controls (54.5836,49.5960) and (53.8610,50.3186) ..(53.8610,51.2100) .. controls (53.8610,52.1014) and (54.5836,52.8240) ..(55.4750,52.8240) .. controls (56.3664,52.8240) and (57.0890,52.1014) ..(57.0890,51.2100);
\path[fill=chfourfig6c0,fill opacity=0.5,draw=chfourfig6c0,line width=0.3985pt,draw opacity=0.5] (58.5336,44.1464) .. controls (58.5336,43.2550) and (57.8110,42.5324) ..(56.9196,42.5324) .. controls (56.0282,42.5324) and (55.3056,43.2550) ..(55.3056,44.1464) .. controls (55.3056,45.0378) and (56.0282,45.7603) ..(56.9196,45.7603) .. controls (57.8110,45.7603) and (58.5336,45.0378) ..(58.5336,44.1464);
\path[fill=chfourfig6c0,fill opacity=0.5,draw=chfourfig6c0,line width=0.3985pt,draw opacity=0.5] (90.8229,48.1813) .. controls (90.8229,47.2899) and (90.1003,46.5673) ..(89.2089,46.5673) .. controls (88.3176,46.5673) and (87.5950,47.2899) ..(87.5950,48.1813) .. controls (87.5950,49.0727) and (88.3176,49.7953) ..(89.2089,49.7953) .. controls (90.1003,49.7953) and (90.8229,49.0727) ..(90.8229,48.1813);
\path[fill=chfourfig6c0,fill opacity=0.5,draw=chfourfig6c0,line width=0.3985pt,draw opacity=0.5] (136.5121,75.4295) .. controls (136.5121,74.5381) and (135.7896,73.8155) ..(134.8982,73.8155) .. controls (134.0068,73.8155) and (133.2842,74.5381) ..(133.2842,75.4295) .. controls (133.2842,76.3208) and (134.0068,77.0434) ..(134.8982,77.0434) .. controls (135.7896,77.0434) and (136.5121,76.3208) ..(136.5121,75.4295);
\path[fill=chfourfig6c0,fill opacity=0.5,draw=chfourfig6c0,line width=0.3985pt,draw opacity=0.5] (134.2008,76.4457) .. controls (134.2008,75.5543) and (133.4782,74.8317) ..(132.5868,74.8317) .. controls (131.6955,74.8317) and (130.9729,75.5543) ..(130.9729,76.4457) .. controls (130.9729,77.3370) and (131.6955,78.0596) ..(132.5868,78.0596) .. controls (133.4782,78.0596) and (134.2008,77.3370) ..(134.2008,76.4457);
\path[fill=chfourfig6c0,fill opacity=0.5,draw=chfourfig6c0,line width=0.3985pt,draw opacity=0.5] (127.9342,70.3883) .. controls (127.9342,69.4970) and (127.2116,68.7744) ..(126.3202,68.7744) .. controls (125.4289,68.7744) and (124.7063,69.4970) ..(124.7063,70.3883) .. controls (124.7063,71.2797) and (125.4289,72.0023) ..(126.3202,72.0023) .. controls (127.2116,72.0023) and (127.9342,71.2797) ..(127.9342,70.3883);
\path[fill=chfourfig6c0,fill opacity=0.5,draw=chfourfig6c0,line width=0.3985pt,draw opacity=0.5] (114.3450,60.2960) .. controls (114.3450,59.4047) and (113.6224,58.6821) ..(112.7310,58.6821) .. controls (111.8397,58.6821) and (111.1171,59.4047) ..(111.1171,60.2960) .. controls (111.1171,61.1874) and (111.8397,61.9100) ..(112.7310,61.9100) .. controls (113.6224,61.9100) and (114.3450,61.1874) ..(114.3450,60.2960);
\path[fill=chfourfig6c0,fill opacity=0.5,draw=chfourfig6c0,line width=0.3985pt,draw opacity=0.5] (60.3667,45.1526) .. controls (60.3667,44.2613) and (59.6441,43.5387) ..(58.7527,43.5387) .. controls (57.8614,43.5387) and (57.1388,44.2613) ..(57.1388,45.1526) .. controls (57.1388,46.0440) and (57.8614,46.7666) ..(58.7527,46.7666) .. controls (59.6441,46.7666) and (60.3667,46.0440) ..(60.3667,45.1526);
\path[fill=chfourfig6c0,fill opacity=0.5,draw=chfourfig6c0,line width=0.3985pt,draw opacity=0.5] (53.5223,46.1589) .. controls (53.5223,45.2675) and (52.7997,44.5449) ..(51.9083,44.5449) .. controls (51.0170,44.5449) and (50.2944,45.2675) ..(50.2944,46.1589) .. controls (50.2944,47.0502) and (51.0170,47.7728) ..(51.9083,47.7728) .. controls (52.7997,47.7728) and (53.5223,47.0502) ..(53.5223,46.1589);
\path[fill=chfourfig6c0,fill opacity=0.5,draw=chfourfig6c0,line width=0.3985pt,draw opacity=0.5] (120.1334,61.3023) .. controls (120.1334,60.4109) and (119.4108,59.6883) ..(118.5194,59.6883) .. controls (117.6280,59.6883) and (116.9055,60.4109) ..(116.9055,61.3023) .. controls (116.9055,62.1936) and (117.6280,62.9162) ..(118.5194,62.9162) .. controls (119.4108,62.9162) and (120.1334,62.1936) ..(120.1334,61.3023);
\path[fill=chfourfig6c0,fill opacity=0.5,draw=chfourfig6c0,line width=0.3985pt,draw opacity=0.5] (66.5337,33.0379) .. controls (66.5337,32.1465) and (65.8111,31.4239) ..(64.9197,31.4239) .. controls (64.0283,31.4239) and (63.3058,32.1465) ..(63.3058,33.0379) .. controls (63.3058,33.9293) and (64.0283,34.6519) ..(64.9197,34.6519) .. controls (65.8111,34.6519) and (66.5337,33.9293) ..(66.5337,33.0379);
\path[fill=chfourfig6c0,fill opacity=0.5,draw=chfourfig6c0,line width=0.3985pt,draw opacity=0.5] (78.2001,41.1177) .. controls (78.2001,40.2263) and (77.4775,39.5037) ..(76.5861,39.5037) .. controls (75.6947,39.5037) and (74.9722,40.2263) ..(74.9722,41.1177) .. controls (74.9722,42.0091) and (75.6947,42.7317) ..(76.5861,42.7317) .. controls (77.4775,42.7317) and (78.2001,42.0091) ..(78.2001,41.1177);
\path[fill=chfourfig6c0,fill opacity=0.5,draw=chfourfig6c0,line width=0.3985pt,draw opacity=0.5] (99.8890,45.1526) .. controls (99.8890,44.2613) and (99.1664,43.5387) ..(98.2751,43.5387) .. controls (97.3837,43.5387) and (96.6611,44.2613) ..(96.6611,45.1526) .. controls (96.6611,46.0440) and (97.3837,46.7666) ..(98.2751,46.7666) .. controls (99.1664,46.7666) and (99.8890,46.0440) ..(99.8890,45.1526);
\path[fill=chfourfig6c0,fill opacity=0.5,draw=chfourfig6c0,line width=0.3985pt,draw opacity=0.5] (57.4675,47.1751) .. controls (57.4675,46.2837) and (56.7450,45.5611) ..(55.8536,45.5611) .. controls (54.9622,45.5611) and (54.2396,46.2837) ..(54.2396,47.1751) .. controls (54.2396,48.0664) and (54.9622,48.7890) ..(55.8536,48.7890) .. controls (56.7450,48.7890) and (57.4675,48.0664) ..(57.4675,47.1751);
\path[fill=chfourfig6c0,fill opacity=0.5,draw=chfourfig6c0,line width=0.3985pt,draw opacity=0.5] (107.6899,53.2324) .. controls (107.6899,52.3411) and (106.9673,51.6185) ..(106.0759,51.6185) .. controls (105.1845,51.6185) and (104.4619,52.3411) ..(104.4619,53.2324) .. controls (104.4619,54.1238) and (105.1845,54.8464) ..(106.0759,54.8464) .. controls (106.9673,54.8464) and (107.6899,54.1238) ..(107.6899,53.2324);
\path[fill=chfourfig6c0,fill opacity=0.5,draw=chfourfig6c0,line width=0.3985pt,draw opacity=0.5] (98.0559,52.2162) .. controls (98.0559,51.3249) and (97.3333,50.6023) ..(96.4419,50.6023) .. controls (95.5505,50.6023) and (94.8279,51.3249) ..(94.8279,52.2162) .. controls (94.8279,53.1076) and (95.5505,53.8302) ..(96.4419,53.8302) .. controls (97.3333,53.8302) and (98.0559,53.1076) ..(98.0559,52.2162);
\path[fill=chfourfig6c0,fill opacity=0.5,draw=chfourfig6c0,line width=0.3985pt,draw opacity=0.5] (80.7007,50.2038) .. controls (80.7007,49.3124) and (79.9781,48.5898) ..(79.0868,48.5898) .. controls (78.1954,48.5898) and (77.4728,49.3124) ..(77.4728,50.2038) .. controls (77.4728,51.0951) and (78.1954,51.8177) ..(79.0868,51.8177) .. controls (79.9781,51.8177) and (80.7007,51.0951) ..(80.7007,50.2038);
\path[fill=chfourfig6c0,fill opacity=0.5,draw=chfourfig6c0,line width=0.3985pt,draw opacity=0.5] (75.7891,50.2038) .. controls (75.7891,49.3124) and (75.0665,48.5898) ..(74.1751,48.5898) .. controls (73.2837,48.5898) and (72.5612,49.3124) ..(72.5612,50.2038) .. controls (72.5612,51.0951) and (73.2837,51.8177) ..(74.1751,51.8177) .. controls (75.0665,51.8177) and (75.7891,51.0951) ..(75.7891,50.2038);
\path[fill=chfourfig6c0,fill opacity=0.5,draw=chfourfig6c0,line width=0.3985pt,draw opacity=0.5] (57.5672,44.1464) .. controls (57.5672,43.2550) and (56.8446,42.5324) ..(55.9532,42.5324) .. controls (55.0619,42.5324) and (54.3393,43.2550) ..(54.3393,44.1464) .. controls (54.3393,45.0378) and (55.0619,45.7603) ..(55.9532,45.7603) .. controls (56.8446,45.7603) and (57.5672,45.0378) ..(57.5672,44.1464);
\path[fill=chfourfig6c0,fill opacity=0.5,draw=chfourfig6c0,line width=0.3985pt,draw opacity=0.5] (74.1452,31.0254) .. controls (74.1452,30.1340) and (73.4226,29.4115) ..(72.5312,29.4115) .. controls (71.6399,29.4115) and (70.9173,30.1340) ..(70.9173,31.0254) .. controls (70.9173,31.9168) and (71.6399,32.6394) ..(72.5312,32.6394) .. controls (73.4226,32.6394) and (74.1452,31.9168) ..(74.1452,31.0254);
\path[fill=chfourfig6c0,fill opacity=0.5,draw=chfourfig6c0,line width=0.3985pt,draw opacity=0.5] (69.6221,48.1813) .. controls (69.6221,47.2899) and (68.8995,46.5673) ..(68.0082,46.5673) .. controls (67.1168,46.5673) and (66.3942,47.2899) ..(66.3942,48.1813) .. controls (66.3942,49.0727) and (67.1168,49.7953) ..(68.0082,49.7953) .. controls (68.8995,49.7953) and (69.6221,49.0727) ..(69.6221,48.1813);
\path[fill=chfourfig6c0,fill opacity=0.5,draw=chfourfig6c0,line width=0.3985pt,draw opacity=0.5] (89.2787,51.2100) .. controls (89.2787,50.3186) and (88.5561,49.5960) ..(87.6647,49.5960) .. controls (86.7733,49.5960) and (86.0508,50.3186) ..(86.0508,51.2100) .. controls (86.0508,52.1014) and (86.7733,52.8240) ..(87.6647,52.8240) .. controls (88.5561,52.8240) and (89.2787,52.1014) ..(89.2787,51.2100);
\path[fill=chfourfig6c0,fill opacity=0.5,draw=chfourfig6c0,line width=0.3985pt,draw opacity=0.5] (143.2669,76.4457) .. controls (143.2669,75.5543) and (142.5443,74.8317) ..(141.6530,74.8317) .. controls (140.7616,74.8317) and (140.0390,75.5543) ..(140.0390,76.4457) .. controls (140.0390,77.3370) and (140.7616,78.0596) ..(141.6530,78.0596) .. controls (142.5443,78.0596) and (143.2669,77.3370) ..(143.2669,76.4457);
\path[fill=chfourfig6c0,fill opacity=0.5,draw=chfourfig6c0,line width=0.3985pt,draw opacity=0.5] (59.4003,51.2100) .. controls (59.4003,50.3186) and (58.6777,49.5960) ..(57.7864,49.5960) .. controls (56.8950,49.5960) and (56.1724,50.3186) ..(56.1724,51.2100) .. controls (56.1724,52.1014) and (56.8950,52.8240) ..(57.7864,52.8240) .. controls (58.6777,52.8240) and (59.4003,52.1014) ..(59.4003,51.2100);
\path[fill=chfourfig6c0,fill opacity=0.5,draw=chfourfig6c0,line width=0.3985pt,draw opacity=0.5] (141.4337,84.5155) .. controls (141.4337,83.6242) and (140.7112,82.9016) ..(139.8198,82.9016) .. controls (138.9284,82.9016) and (138.2058,83.6242) ..(138.2058,84.5155) .. controls (138.2058,85.4069) and (138.9284,86.1295) ..(139.8198,86.1295) .. controls (140.7112,86.1295) and (141.4337,85.4069) ..(141.4337,84.5155);
\path[fill=chfourfig6c0,fill opacity=0.5,draw=chfourfig6c0,line width=0.3985pt,draw opacity=0.5] (94.8777,58.2736) .. controls (94.8777,57.3822) and (94.1552,56.6596) ..(93.2638,56.6596) .. controls (92.3724,56.6596) and (91.6498,57.3822) ..(91.6498,58.2736) .. controls (91.6498,59.1650) and (92.3724,59.8876) ..(93.2638,59.8876) .. controls (94.1552,59.8876) and (94.8777,59.1650) ..(94.8777,58.2736);
\path[fill=chfourfig6c0,fill opacity=0.5,draw=chfourfig6c0,line width=0.3985pt,draw opacity=0.5] (97.9563,57.2673) .. controls (97.9563,56.3760) and (97.2336,55.6534) ..(96.3423,55.6534) .. controls (95.4509,55.6534) and (94.7283,56.3760) ..(94.7283,57.2673) .. controls (94.7283,58.1587) and (95.4509,58.8813) ..(96.3423,58.8813) .. controls (97.2336,58.8813) and (97.9563,58.1587) ..(97.9563,57.2673);
\path[fill=chfourfig6c0,fill opacity=0.5,draw=chfourfig6c0,line width=0.3985pt,draw opacity=0.5] (66.7230,36.0666) .. controls (66.7230,35.1752) and (66.0004,34.4526) ..(65.1090,34.4526) .. controls (64.2176,34.4526) and (63.4951,35.1752) ..(63.4951,36.0666) .. controls (63.4951,36.9579) and (64.2176,37.6805) ..(65.1090,37.6805) .. controls (66.0004,37.6805) and (66.7230,36.9579) ..(66.7230,36.0666);
\path[fill=chfourfig6c0,fill opacity=0.5,draw=chfourfig6c0,line width=0.3985pt,draw opacity=0.5] (105.9563,54.2387) .. controls (105.9563,53.3473) and (105.2337,52.6247) ..(104.3424,52.6247) .. controls (103.4510,52.6247) and (102.7284,53.3473) ..(102.7284,54.2387) .. controls (102.7284,55.1300) and (103.4510,55.8526) ..(104.3424,55.8526) .. controls (105.2337,55.8526) and (105.9563,55.1300) ..(105.9563,54.2387);
\path[fill=chfourfig6c0,fill opacity=0.5,draw=chfourfig6c0,line width=0.3985pt,draw opacity=0.5] (95.0670,45.1526) .. controls (95.0670,44.2613) and (94.3444,43.5387) ..(93.4531,43.5387) .. controls (92.5617,43.5387) and (91.8391,44.2613) ..(91.8391,45.1526) .. controls (91.8391,46.0440) and (92.5617,46.7666) ..(93.4531,46.7666) .. controls (94.3444,46.7666) and (95.0670,46.0440) ..(95.0670,45.1526);
\path[fill=chfourfig6c0,fill opacity=0.5,draw=chfourfig6c0,line width=0.3985pt,draw opacity=0.5] (73.7666,36.0666) .. controls (73.7666,35.1752) and (73.0440,34.4526) ..(72.1527,34.4526) .. controls (71.2613,34.4526) and (70.5387,35.1752) ..(70.5387,36.0666) .. controls (70.5387,36.9579) and (71.2613,37.6805) ..(72.1527,37.6805) .. controls (73.0440,37.6805) and (73.7666,36.9579) ..(73.7666,36.0666);
\path[fill=chfourfig6c0,fill opacity=0.5,draw=chfourfig6c0,line width=0.3985pt,draw opacity=0.5] (93.3335,51.2100) .. controls (93.3335,50.3186) and (92.6109,49.5960) ..(91.7195,49.5960) .. controls (90.8282,49.5960) and (90.1056,50.3186) ..(90.1056,51.2100) .. controls (90.1056,52.1014) and (90.8282,52.8240) ..(91.7195,52.8240) .. controls (92.6109,52.8240) and (93.3335,52.1014) ..(93.3335,51.2100);
\path[fill=chfourfig6c0,fill opacity=0.5,draw=chfourfig6c0,line width=0.3985pt,draw opacity=0.5] (229.2456,111.7737) .. controls (229.2456,110.8823) and (228.5230,110.1597) ..(227.6317,110.1597) .. controls (226.7403,110.1597) and (226.0177,110.8823) ..(226.0177,111.7737) .. controls (226.0177,112.6650) and (226.7403,113.3876) ..(227.6317,113.3876) .. controls (228.5230,113.3876) and (229.2456,112.6650) ..(229.2456,111.7737);
\path[fill=chfourfig6c0,fill opacity=0.5,draw=chfourfig6c0,line width=0.3985pt,draw opacity=0.5] (66.1451,41.1177) .. controls (66.1451,40.2263) and (65.4225,39.5037) ..(64.5312,39.5037) .. controls (63.6398,39.5037) and (62.9172,40.2263) ..(62.9172,41.1177) .. controls (62.9172,42.0091) and (63.6398,42.7317) ..(64.5312,42.7317) .. controls (65.4225,42.7317) and (66.1451,42.0091) ..(66.1451,41.1177);
\path[fill=chfourfig6c0,fill opacity=0.5,draw=chfourfig6c0,line width=0.3985pt,draw opacity=0.5] (52.9445,49.1875) .. controls (52.9445,48.2962) and (52.2219,47.5736) ..(51.3305,47.5736) .. controls (50.4391,47.5736) and (49.7165,48.2962) ..(49.7165,49.1875) .. controls (49.7165,50.0789) and (50.4391,50.8015) ..(51.3305,50.8015) .. controls (52.2219,50.8015) and (52.9445,50.0789) ..(52.9445,49.1875);
\path[fill=chfourfig6c0,fill opacity=0.5,draw=chfourfig6c0,line width=0.3985pt,draw opacity=0.5] (66.3444,34.0541) .. controls (66.3444,33.1627) and (65.6218,32.4401) ..(64.7304,32.4401) .. controls (63.8390,32.4401) and (63.1165,33.1627) ..(63.1165,34.0541) .. controls (63.1165,34.9455) and (63.8390,35.6681) ..(64.7304,35.6681) .. controls (65.6218,35.6681) and (66.3444,34.9455) ..(66.3444,34.0541);
\path[fill=chfourfig6c0,fill opacity=0.5,draw=chfourfig6c0,line width=0.3985pt,draw opacity=0.5] (123.9890,60.2960) .. controls (123.9890,59.4047) and (123.2664,58.6821) ..(122.3750,58.6821) .. controls (121.4836,58.6821) and (120.7610,59.4047) ..(120.7610,60.2960) .. controls (120.7610,61.1874) and (121.4836,61.9100) ..(122.3750,61.9100) .. controls (123.2664,61.9100) and (123.9890,61.1874) ..(123.9890,60.2960);
\path[fill=chfourfig6c0,fill opacity=0.5,draw=chfourfig6c0,line width=0.3985pt,draw opacity=0.5] (108.5566,68.3659) .. controls (108.5566,67.4745) and (107.8340,66.7519) ..(106.9427,66.7519) .. controls (106.0513,66.7519) and (105.3287,67.4745) ..(105.3287,68.3659) .. controls (105.3287,69.2572) and (106.0513,69.9798) ..(106.9427,69.9798) .. controls (107.8340,69.9798) and (108.5566,69.2572) ..(108.5566,68.3659);
\path[fill=chfourfig6c0,fill opacity=0.5,draw=chfourfig6c0,line width=0.3985pt,draw opacity=0.5] (118.3002,73.4170) .. controls (118.3002,72.5256) and (117.5776,71.8030) ..(116.6862,71.8030) .. controls (115.7949,71.8030) and (115.0723,72.5256) ..(115.0723,73.4170) .. controls (115.0723,74.3084) and (115.7949,75.0310) ..(116.6862,75.0310) .. controls (117.5776,75.0310) and (118.3002,74.3084) ..(118.3002,73.4170);
\path[fill=chfourfig6c0,fill opacity=0.5,draw=chfourfig6c0,line width=0.3985pt,draw opacity=0.5] (62.7777,45.1526) .. controls (62.7777,44.2613) and (62.0551,43.5387) ..(61.1638,43.5387) .. controls (60.2724,43.5387) and (59.5498,44.2613) ..(59.5498,45.1526) .. controls (59.5498,46.0440) and (60.2724,46.7666) ..(61.1638,46.7666) .. controls (62.0551,46.7666) and (62.7777,46.0440) ..(62.7777,45.1526);
\path[fill=chfourfig6c0,fill opacity=0.5,draw=chfourfig6c0,line width=0.3985pt,draw opacity=0.5] (65.4677,45.1526) .. controls (65.4677,44.2613) and (64.7451,43.5387) ..(63.8537,43.5387) .. controls (62.9623,43.5387) and (62.2397,44.2613) ..(62.2397,45.1526) .. controls (62.2397,46.0440) and (62.9623,46.7666) ..(63.8537,46.7666) .. controls (64.7451,46.7666) and (65.4677,46.0440) ..(65.4677,45.1526);
\path[fill=chfourfig6c0,fill opacity=0.5,draw=chfourfig6c0,line width=0.3985pt,draw opacity=0.5] (97.3784,57.2673) .. controls (97.3784,56.3760) and (96.6558,55.6534) ..(95.7644,55.6534) .. controls (94.8731,55.6534) and (94.1505,56.3760) ..(94.1505,57.2673) .. controls (94.1505,58.1587) and (94.8731,58.8813) ..(95.7644,58.8813) .. controls (96.6558,58.8813) and (97.3784,58.1587) ..(97.3784,57.2673);
\path[fill=chfourfig6c0,fill opacity=0.5,draw=chfourfig6c0,line width=0.3985pt,draw opacity=0.5] (80.7007,38.0890) .. controls (80.7007,37.1976) and (79.9781,36.4751) ..(79.0868,36.4751) .. controls (78.1954,36.4751) and (77.4728,37.1976) ..(77.4728,38.0890) .. controls (77.4728,38.9804) and (78.1954,39.7030) ..(79.0868,39.7030) .. controls (79.9781,39.7030) and (80.7007,38.9804) ..(80.7007,38.0890);
\path[fill=chfourfig6c0,fill opacity=0.5,draw=chfourfig6c0,line width=0.3985pt,draw opacity=0.5] (52.7452,48.1813) .. controls (52.7452,47.2899) and (52.0226,46.5673) ..(51.1312,46.5673) .. controls (50.2399,46.5673) and (49.5173,47.2899) ..(49.5173,48.1813) .. controls (49.5173,49.0727) and (50.2399,49.7953) ..(51.1312,49.7953) .. controls (52.0226,49.7953) and (52.7452,49.0727) ..(52.7452,48.1813);
\path[fill=chfourfig6c0,fill opacity=0.5,draw=chfourfig6c0,line width=0.3985pt,draw opacity=0.5] (78.4890,41.1177) .. controls (78.4890,40.2263) and (77.7664,39.5037) ..(76.8750,39.5037) .. controls (75.9837,39.5037) and (75.2611,40.2263) ..(75.2611,41.1177) .. controls (75.2611,42.0091) and (75.9837,42.7317) ..(76.8750,42.7317) .. controls (77.7664,42.7317) and (78.4890,42.0091) ..(78.4890,41.1177);
\path[fill=chfourfig6c0,fill opacity=0.5,draw=chfourfig6c0,line width=0.3985pt,draw opacity=0.5] (137.1896,89.5667) .. controls (137.1896,88.6753) and (136.4670,87.9527) ..(135.5757,87.9527) .. controls (134.6843,87.9527) and (133.9617,88.6753) ..(133.9617,89.5667) .. controls (133.9617,90.4580) and (134.6843,91.1806) ..(135.5757,91.1806) .. controls (136.4670,91.1806) and (137.1896,90.4580) ..(137.1896,89.5667);
\path[fill=chfourfig6c0,fill opacity=0.5,draw=chfourfig6c0,line width=0.3985pt,draw opacity=0.5] (337.9793,181.4134) .. controls (337.9793,180.5220) and (337.2567,179.7994) ..(336.3653,179.7994) .. controls (335.4739,179.7994) and (334.7513,180.5220) ..(334.7513,181.4134) .. controls (334.7513,182.3048) and (335.4739,183.0274) ..(336.3653,183.0274) .. controls (337.2567,183.0274) and (337.9793,182.3048) ..(337.9793,181.4134);
\path[fill=chfourfig6c0,fill opacity=0.5,draw=chfourfig6c0,line width=0.3985pt,draw opacity=0.5] (111.9340,70.3883) .. controls (111.9340,69.4970) and (111.2114,68.7744) ..(110.3201,68.7744) .. controls (109.4287,68.7744) and (108.7061,69.4970) ..(108.7061,70.3883) .. controls (108.7061,71.2797) and (109.4287,72.0023) ..(110.3201,72.0023) .. controls (111.2114,72.0023) and (111.9340,71.2797) ..(111.9340,70.3883);
\path[fill=none,draw=chfourfig6c0,line width=0.3985pt,draw opacity=1.0] (51.0416,187.7099) --(340.2309,187.7099);
\path[fill=none,draw=chfourfig6c0,line width=0.3985pt,draw opacity=1.0] (51.0416,187.7099) --(51.0416,193.0898);
\path[fill=none,draw=chfourfig6c0,line width=0.3985pt,draw opacity=1.0] (99.2315,187.7099) --(99.2315,193.0898);
\path[fill=none,draw=chfourfig6c0,line width=0.3985pt,draw opacity=1.0] (147.4314,187.7099) --(147.4314,193.0898);
\path[fill=none,draw=chfourfig6c0,line width=0.3985pt,draw opacity=1.0] (195.6312,187.7099) --(195.6312,193.0898);
\path[fill=none,draw=chfourfig6c0,line width=0.3985pt,draw opacity=1.0] (243.8311,187.7099) --(243.8311,193.0898);
\path[fill=none,draw=chfourfig6c0,line width=0.3985pt,draw opacity=1.0] (292.0310,187.7099) --(292.0310,193.0898);
\path[fill=none,draw=chfourfig6c0,line width=0.3985pt,draw opacity=1.0] (340.2309,187.7099) --(340.2309,193.0898);
\path[fill=none,draw=chfourfig6c0,line width=0.3985pt,draw opacity=1.0] (39.6242,183.4259) --(39.6242,32.0317);
\path[fill=none,draw=chfourfig6c0,line width=0.3985pt,draw opacity=1.0] (39.6242,183.4259) --(34.2443,183.4259);
\path[fill=none,draw=chfourfig6c0,line width=0.3985pt,draw opacity=1.0] (39.6242,132.9645) --(34.2443,132.9645);
\path[fill=none,draw=chfourfig6c0,line width=0.3985pt,draw opacity=1.0] (39.6242,82.5030) --(34.2443,82.5030);
\path[fill=none,draw=chfourfig6c0,line width=0.3985pt,draw opacity=1.0] (39.6242,32.0317) --(34.2443,32.0317);
\path[fill=none,draw=chfourfig6c0,line width=0.3985pt,draw opacity=1.0] (39.6242,187.7099)--(39.6242,17.6653)--(347.7826,17.6653)--(347.7826,187.7099)--cycle;
\path[fill=none,draw=chfourfig6c1,line width=1.5940pt,draw opacity=1.0] (51.0415,47.1253) --(51.1312,47.0555) --(51.2308,46.9858) --(51.3305,46.9160) --(51.4201,46.8463) --(51.5198,46.7766) --(51.6194,46.7068) --(51.7090,46.6371) --(51.8087,46.5673) --(51.9083,46.4976) --(51.9980,46.4279) --(52.0976,46.3581) --(52.1972,46.2884) --(52.2869,46.2186) --(52.3865,46.1489) --(52.4861,46.0792) --(52.5758,46.0094) --(52.6754,45.9397) --(52.7751,45.8699) --(52.8647,45.7903) --(52.9644,45.7205) --(53.0640,45.6508) --(53.1537,45.5910) --(53.2533,45.5212) --(53.3529,45.4515) --(53.4426,45.3818) --(53.5422,45.3120) --(53.6418,45.2423) --(53.7415,45.1725) --(53.8311,45.1028) --(53.9308,45.0331) --(54.0304,44.9633) --(54.1200,44.8936) --(54.2197,44.8338) --(54.3193,44.7641) --(54.4090,44.6943) --(54.5086,44.6246) --(54.6082,44.5549) --(54.6979,44.4951) --(54.7975,44.4253) --(54.8971,44.3556) --(54.9868,44.2958) --(55.0864,44.2261) --(55.1861,44.1564) --(55.2757,44.0966) --(55.3754,44.0268) --(55.4750,43.9671) --(55.5647,43.8973) --(55.6643,43.8375) --(55.7639,43.7678) --(55.8536,43.7080) --(55.9532,43.6482) --(56.0528,43.5785) --(56.1425,43.5187) --(56.2421,43.4590) --(56.3417,43.3992) --(56.4314,43.3294) --(56.5310,43.2697) --(56.6307,43.2099) --(56.7203,43.1501) --(56.8199,43.0903) --(56.9196,43.0305) --(57.0092,42.9708) --(57.1089,42.9110) --(57.2085,42.8612) --(57.3081,42.8014) --(57.3978,42.7416) --(57.4974,42.6819) --(57.5970,42.6320) --(57.6867,42.5723) --(57.7863,42.5225) --(57.8860,42.4627) --(57.9756,42.4129) --(58.0753,42.3531) --(58.1749,42.3033) --(58.2646,42.2535) --(58.3642,42.2036) --(58.4638,42.1538) --(58.5535,42.0941) --(58.6531,42.0442) --(58.7527,41.9944) --(58.8424,41.9546) --(58.9420,41.9048) --(59.0416,41.8549) --(59.1313,41.8051) --(59.2309,41.7653) --(59.3306,41.7155) --(59.4202,41.6756) --(59.5199,41.6258) --(59.6195,41.5860) --(59.7092,41.5361) --(59.8088,41.4963) --(59.9084,41.4564) --(59.9981,41.4166) --(60.0977,41.3767) --(60.1973,41.3369) --(60.2870,41.2970) --(60.3866,41.2572) --(60.4863,41.2173) --(60.5759,41.1874) --(60.6755,41.1476) --(60.7752,41.1177) --(60.8748,41.0779) --(60.9645,41.0480) --(61.0641,41.0181) --(61.1637,40.9782) --(61.2534,40.9483) --(61.3530,40.9185) --(61.4526,40.8886) --(61.5423,40.8587) --(61.6419,40.8288) --(61.7416,40.8089) --(61.8312,40.7790) --(61.9309,40.7491) --(62.0305,40.7292) --(62.1201,40.6993) --(62.2198,40.6793) --(62.3194,40.6594) --(62.4091,40.6295) --(62.5087,40.6096) --(62.6083,40.5897) --(62.6980,40.5697) --(62.7976,40.5498) --(62.8972,40.5399) --(62.9869,40.5199) --(63.0865,40.5000) --(63.1862,40.4901) --(63.2758,40.4701) --(63.3755,40.4602) --(63.4751,40.4502) --(63.5647,40.4303) --(63.6644,40.4203) --(63.7640,40.4103) --(63.8537,40.4004) --(63.9533,40.3904) --(64.0529,40.3805) --(64.1426,40.3805) --(64.2422,40.3705) --(64.3418,40.3605) --(64.4415,40.3605) --(64.5311,40.3506) --(64.6308,40.3506) --(64.7304,40.3506) --(64.8201,40.3506) --(64.9197,40.3506) --(65.0193,40.3506) --(65.1090,40.3506) --(65.2086,40.3506) --(65.3082,40.3506) --(65.3979,40.3506) --(65.4975,40.3605) --(65.5972,40.3605) --(65.6868,40.3705) --(65.7864,40.3705) --(65.8861,40.3805) --(65.9757,40.3904) --(66.0754,40.3904) --(66.1750,40.4004) --(66.2646,40.4103) --(66.3643,40.4203) --(66.4639,40.4303) --(66.5536,40.4502) --(66.6532,40.4602) --(66.7528,40.4701) --(66.8425,40.4901) --(66.9421,40.5000) --(67.0417,40.5199) --(67.1314,40.5299) --(67.2310,40.5498) --(67.3307,40.5697) --(67.4203,40.5897) --(67.5200,40.6096) --(67.6196,40.6295) --(67.7093,40.6495) --(67.8089,40.6694) --(67.9085,40.6893) --(68.0081,40.7092) --(68.0978,40.7391) --(68.1974,40.7590) --(68.2971,40.7790) --(68.3867,40.8089) --(68.4864,40.8288) --(68.5860,40.8587) --(68.6756,40.8886) --(68.7753,40.9185) --(68.8749,40.9384) --(68.9646,40.9683) --(69.0642,40.9982) --(69.1638,41.0280) --(69.2535,41.0579) --(69.3531,41.0878) --(69.4527,41.1177) --(69.5424,41.1575) --(69.6420,41.1874) --(69.7417,41.2173) --(69.8313,41.2572) --(69.9310,41.2871) --(70.0306,41.3269) --(70.1203,41.3568) --(70.2199,41.3967) --(70.3195,41.4365) --(70.4092,41.4664) --(70.5088,41.5063) --(70.6084,41.5461) --(70.6981,41.5860) --(70.7977,41.6258) --(70.8974,41.6657) --(70.9870,41.7055) --(71.0866,41.7454) --(71.1863,41.7852) --(71.2759,41.8251) --(71.3756,41.8649) --(71.4752,41.9048) --(71.5748,41.9546) --(71.6645,41.9944) --(71.7641,42.0343) --(71.8637,42.0841) --(71.9534,42.1239) --(72.0530,42.1738) --(72.1526,42.2136) --(72.2423,42.2634) --(72.3419,42.3033) --(72.4416,42.3531) --(72.5312,42.4029) --(72.6309,42.4427) --(72.7305,42.4926) --(72.8202,42.5424) --(72.9198,42.5922) --(73.0194,42.6420) --(73.1091,42.6819) --(73.2087,42.7317) --(73.3083,42.7815) --(73.3980,42.8313) --(73.4976,42.8811) --(73.5972,42.9309) --(73.6869,42.9807) --(73.7865,43.0305) --(73.8862,43.0903) --(73.9758,43.1402) --(74.0755,43.1900) --(74.1751,43.2398) --(74.2648,43.2896) --(74.3644,43.3494) --(74.4640,43.3992) --(74.5537,43.4490) --(74.6533,43.5088) --(74.7529,43.5586) --(74.8426,43.6084) --(74.9422,43.6682) --(75.0419,43.7180) --(75.1415,43.7778) --(75.2311,43.8276) --(75.3308,43.8874) --(75.4304,43.9372) --(75.5201,43.9969) --(75.6197,44.0468) --(75.7193,44.1065) --(75.8090,44.1564) --(75.9086,44.2161) --(76.0082,44.2659) --(76.0979,44.3257) --(76.1975,44.3855) --(76.2972,44.4353) --(76.3868,44.4951) --(76.4865,44.5549) --(76.5861,44.6047) --(76.6758,44.6644) --(76.7754,44.7242) --(76.8750,44.7740) --(76.9647,44.8338) --(77.0643,44.8936) --(77.1639,44.9434) --(77.2536,45.0032) --(77.3532,45.0630) --(77.4529,45.1227) --(77.5425,45.1825) --(77.6421,45.2323) --(77.7418,45.2921) --(77.8314,45.3519) --(77.9311,45.4117) --(78.0307,45.4615) --(78.1204,45.5212) --(78.2200,45.5810) --(78.3196,45.6408) --(78.4093,45.7006) --(78.5089,45.7504) --(78.6085,45.8102) --(78.7082,45.8699) --(78.7978,45.9297) --(78.8974,45.9895) --(78.9971,46.0493) --(79.0867,46.0991) --(79.1864,46.1589) --(79.2860,46.2186) --(79.3757,46.2784) --(79.4753,46.3382) --(79.5749,46.3980) --(79.6646,46.4478) --(79.7642,46.5076) --(79.8638,46.5673) --(79.9535,46.6271) --(80.0531,46.6869) --(80.1528,46.7367) --(80.2424,46.7965) --(80.3421,46.8563) --(80.4417,46.9160) --(80.5313,46.9758) --(80.6310,47.0356) --(80.7306,47.0854) --(80.8203,47.1452) --(80.9199,47.2050) --(81.0195,47.2647) --(81.1092,47.3145) --(81.2088,47.3743) --(81.3084,47.4341) --(81.3981,47.4939) --(81.4977,47.5536) --(81.5974,47.6035) --(81.6870,47.6632) --(81.7867,47.7230) --(81.8863,47.7828) --(81.9759,47.8326) --(82.0756,47.8924) --(82.1752,47.9522) --(82.2748,48.0119) --(82.3645,48.0618) --(82.4641,48.1215) --(82.5637,48.1813) --(82.6534,48.2311) --(82.7530,48.2909) --(82.8527,48.3507) --(82.9423,48.4005) --(83.0420,48.4603) --(83.1416,48.5200) --(83.2312,48.5699) --(83.3309,48.6296) --(83.4305,48.6894) --(83.5202,48.7392) --(83.6198,48.7990) --(83.7194,48.8588) --(83.8091,48.9086) --(83.9087,48.9684) --(84.0083,49.0282) --(84.0980,49.0780) --(84.1976,49.1377) --(84.2973,49.1875) --(84.3869,49.2473) --(84.4866,49.3071) --(84.5862,49.3569) --(84.6759,49.4167) --(84.7755,49.4665) --(84.8751,49.5263) --(84.9648,49.5761) --(85.0644,49.6359) --(85.1640,49.6857) --(85.2537,49.7455) --(85.3533,49.8052) --(85.4530,49.8551) --(85.5426,49.9148) --(85.6422,49.9646) --(85.7419,50.0244) --(85.8415,50.0742) --(85.9312,50.1340) --(86.0308,50.1838) --(86.1304,50.2436) --(86.2201,50.2934) --(86.3197,50.3532) --(86.4193,50.4030) --(86.5090,50.4528) --(86.6086,50.5126) --(86.7083,50.5624) --(86.7979,50.6222) --(86.8976,50.6720) --(86.9972,50.7318) --(87.0869,50.7816) --(87.1865,50.8414) --(87.2861,50.8912) --(87.3758,50.9410) --(87.4754,51.0008) --(87.5750,51.0506) --(87.6647,51.1104) --(87.7643,51.1602) --(87.8639,51.2100) --(87.9536,51.2698) --(88.0532,51.3196) --(88.1529,51.3794) --(88.2425,51.4292) --(88.3421,51.4790) --(88.4418,51.5388) --(88.5314,51.5886) --(88.6311,51.6483) --(88.7307,51.6982) --(88.8204,51.7480) --(88.9200,51.8078) --(89.0196,51.8576) --(89.1093,51.9074) --(89.2089,51.9672) --(89.3085,52.0170) --(89.3982,52.0768) --(89.4978,52.1266) --(89.5975,52.1764) --(89.6971,52.2362) --(89.7867,52.2860) --(89.8864,52.3358) --(89.9860,52.3956) --(90.0757,52.4454) --(90.1753,52.4952) --(90.2749,52.5550) --(90.3646,52.6048) --(90.4642,52.6646) --(90.5638,52.7144) --(90.6535,52.7642) --(90.7531,52.8240) --(90.8528,52.8738) --(90.9424,52.9236) --(91.0421,52.9834) --(91.1417,53.0332) --(91.2314,53.0830) --(91.3310,53.1428) --(91.4306,53.1926) --(91.5203,53.2523) --(91.6199,53.3022) --(91.7195,53.3520) --(91.8092,53.4118) --(91.9088,53.4616) --(92.0085,53.5114) --(92.0981,53.5712) --(92.1977,53.6210) --(92.2974,53.6808) --(92.3870,53.7306) --(92.4867,53.7804) --(92.5863,53.8402) --(92.6760,53.8900) --(92.7756,53.9498) --(92.8752,53.9996) --(92.9649,54.0494) --(93.0645,54.1091) --(93.1641,54.1590) --(93.2638,54.2187) --(93.3534,54.2686) --(93.4530,54.3184) --(93.5527,54.3781) --(93.6423,54.4280) --(93.7420,54.4877) --(93.8416,54.5376) --(93.9313,54.5973) --(94.0309,54.6471) --(94.1305,54.6970) --(94.2202,54.7567) --(94.3198,54.8065) --(94.4194,54.8663) --(94.5091,54.9161) --(94.6087,54.9759) --(94.7084,55.0257) --(94.7980,55.0855) --(94.8977,55.1353) --(94.9973,55.1851) --(95.0869,55.2449) --(95.1866,55.2947) --(95.2862,55.3545) --(95.3759,55.4043) --(95.4755,55.4641) --(95.5751,55.5139) --(95.6648,55.5737) --(95.7644,55.6235) --(95.8640,55.6833) --(95.9537,55.7331) --(96.0533,55.7929) --(96.1530,55.8427) --(96.2426,55.9025) --(96.3423,55.9523) --(96.4419,56.0120) --(96.5316,56.0619) --(96.6312,56.1216) --(96.7308,56.1714) --(96.8304,56.2312) --(96.9201,56.2810) --(97.0197,56.3408) --(97.1194,56.3906) --(97.2090,56.4504) --(97.3086,56.5002) --(97.4083,56.5600) --(97.4979,56.6098) --(97.5976,56.6696) --(97.6972,56.7194) --(97.7868,56.7792) --(97.8865,56.8290) --(97.9861,56.8888) --(98.0758,56.9386) --(98.1754,56.9984) --(98.2750,57.0482) --(98.3647,57.1080) --(98.4643,57.1578) --(98.5639,57.2175) --(98.6536,57.2673) --(98.7532,57.3271) --(98.8529,57.3769) --(98.9425,57.4367) --(99.0422,57.4865) --(99.1418,57.5463) --(99.2315,57.5961) --(99.3311,57.6559) --(99.4307,57.7157) --(99.5204,57.7655) --(99.6200,57.8153) --(99.7196,57.8751) --(99.8093,57.9249) --(99.9089,57.9847) --(100.0086,58.0345) --(100.0982,58.0943) --(100.1978,58.1441) --(100.2975,58.2038) --(100.3971,58.2537) --(100.4868,58.3134) --(100.5864,58.3633) --(100.6860,58.4230) --(100.7757,58.4728) --(100.8753,58.5326) --(100.9749,58.5824) --(101.0646,58.6422) --(101.1642,58.6920) --(101.2639,58.7419) --(101.3535,58.8016) --(101.4532,58.8514) --(101.5528,58.9112) --(101.6425,58.9610) --(101.7421,59.0108) --(101.8417,59.0706) --(101.9314,59.1204) --(102.0310,59.1802) --(102.1306,59.2300) --(102.2203,59.2798) --(102.3199,59.3396) --(102.4195,59.3894) --(102.5092,59.4392) --(102.6088,59.4990) --(102.7085,59.5488) --(102.7981,59.5986) --(102.8978,59.6584) --(102.9974,59.7082) --(103.0870,59.7580) --(103.1867,59.8178) --(103.2863,59.8676) --(103.3759,59.9174) --(103.4756,59.9673) --(103.5752,60.0270) --(103.6649,60.0769) --(103.7645,60.1267) --(103.8641,60.1765) --(103.9638,60.2263) --(104.0534,60.2861) --(104.1530,60.3359) --(104.2527,60.3857) --(104.3423,60.4355) --(104.4420,60.4853) --(104.5416,60.5351) --(104.6313,60.5849) --(104.7309,60.6348) --(104.8305,60.6945) --(104.9202,60.7444) --(105.0198,60.7942) --(105.1194,60.8440) --(105.2091,60.8938) --(105.3087,60.9436) --(105.4084,60.9934) --(105.4980,61.0432) --(105.5977,61.0931) --(105.6973,61.1429) --(105.7870,61.1927) --(105.8866,61.2425) --(105.9862,61.2923) --(106.0759,61.3322) --(106.1755,61.3820) --(106.2751,61.4318) --(106.3648,61.4816) --(106.4644,61.5314) --(106.5641,61.5812) --(106.6537,61.6211) --(106.7533,61.6709) --(106.8530,61.7207) --(106.9427,61.7705) --(107.0423,61.8203) --(107.1419,61.8602) --(107.2316,61.9100) --(107.3312,61.9598) --(107.4308,61.9997) --(107.5304,62.0495) --(107.6201,62.0993) --(107.7197,62.1391) --(107.8194,62.1889) --(107.9090,62.2288) --(108.0087,62.2786) --(108.1083,62.3185) --(108.1980,62.3683) --(108.2976,62.4181) --(108.3972,62.4579) --(108.4869,62.4978) --(108.5865,62.5476) --(108.6861,62.5875) --(108.7758,62.6373) --(108.8754,62.6771) --(108.9750,62.7270) --(109.0647,62.7668) --(109.1643,62.8067) --(109.2639,62.8565) --(109.3536,62.8963) --(109.4533,62.9362) --(109.5529,62.9760) --(109.6425,63.0258) --(109.7422,63.0657) --(109.8418,63.1055) --(109.9315,63.1454) --(110.0311,63.1852) --(110.1307,63.2251) --(110.2204,63.2749) --(110.3200,63.3148) --(110.4196,63.3546) --(110.5093,63.3945) --(110.6089,63.4343) --(110.7086,63.4742) --(110.7982,63.5140) --(110.8979,63.5539) --(110.9975,63.5937) --(111.0971,63.6336) --(111.1868,63.6734) --(111.2864,63.7133) --(111.3860,63.7431) --(111.4757,63.7830) --(111.5753,63.8228) --(111.6750,63.8627) --(111.7646,63.9026) --(111.8643,63.9324) --(111.9639,63.9723) --(112.0536,64.0121) --(112.1532,64.0520) --(112.2528,64.0819) --(112.3425,64.1217) --(112.4421,64.1616) --(112.5417,64.1915) --(112.6314,64.2313) --(112.7310,64.2712) --(112.8307,64.3011) --(112.9203,64.3409) --(113.0199,64.3708) --(113.1196,64.4107) --(113.2092,64.4406) --(113.3089,64.4804) --(113.4085,64.5103) --(113.4981,64.5501) --(113.5978,64.5800) --(113.6974,64.6199) --(113.7871,64.6498) --(113.8867,64.6796) --(113.9863,64.7195) --(114.0760,64.7494) --(114.1756,64.7892) --(114.2752,64.8191) --(114.3649,64.8490) --(114.4645,64.8789) --(114.5642,64.9188) --(114.6638,64.9486) --(114.7534,64.9785) --(114.8531,65.0184) --(114.9527,65.0483) --(115.0424,65.0782) --(115.1420,65.1080) --(115.2416,65.1379) --(115.3313,65.1678) --(115.4309,65.2077) --(115.5305,65.2376) --(115.6202,65.2675) --(115.7198,65.2973) --(115.8195,65.3272) --(115.9091,65.3571) --(116.0088,65.3870) --(116.1084,65.4169) --(116.1981,65.4468) --(116.2977,65.4767) --(116.3973,65.5066) --(116.4870,65.5365) --(116.5866,65.5663) --(116.6862,65.5962) --(116.7759,65.6261) --(116.8755,65.6560) --(116.9752,65.6859) --(117.0648,65.7158) --(117.1644,65.7457) --(117.2641,65.7756) --(117.3537,65.8055) --(117.4534,65.8353) --(117.5530,65.8553) --(117.6427,65.8851) --(117.7423,65.9150) --(117.8419,65.9449) --(117.9316,65.9748) --(118.0312,66.0047) --(118.1308,66.0246) --(118.2305,66.0545) --(118.3201,66.0844) --(118.4198,66.1143) --(118.5194,66.1442) --(118.6090,66.1641) --(118.7087,66.1940) --(118.8083,66.2239) --(118.8980,66.2538) --(118.9976,66.2737) --(119.0972,66.3036) --(119.1869,66.3335) --(119.2865,66.3634) --(119.3861,66.3833) --(119.4758,66.4132) --(119.5754,66.4431) --(119.6750,66.4630) --(119.7647,66.4929) --(119.8643,66.5228) --(119.9640,66.5427) --(120.0536,66.5726) --(120.1533,66.6025) --(120.2529,66.6224) --(120.3426,66.6523) --(120.4422,66.6822) --(120.5418,66.7021) --(120.6315,66.7320) --(120.7311,66.7619) --(120.8307,66.7818) --(120.9204,66.8117) --(121.0200,66.8416) --(121.1197,66.8615) --(121.2093,66.8914) --(121.3089,66.9213) --(121.4086,66.9412) --(121.4982,66.9711) --(121.5979,67.0010) --(121.6975,67.0209) --(121.7971,67.0508) --(121.8868,67.0807) --(121.9864,67.1006) --(122.0860,67.1305) --(122.1757,67.1604) --(122.2753,67.1803) --(122.3750,67.2102) --(122.4646,67.2401) --(122.5643,67.2600) --(122.6639,67.2899) --(122.7536,67.3198) --(122.8532,67.3397) --(122.9528,67.3696) --(123.0425,67.3995) --(123.1421,67.4194) --(123.2417,67.4493) --(123.3314,67.4792) --(123.4310,67.5091) --(123.5307,67.5290) --(123.6203,67.5589) --(123.7199,67.5888) --(123.8196,67.6187) --(123.9092,67.6386) --(124.0089,67.6685) --(124.1085,67.6984) --(124.1982,67.7283) --(124.2978,67.7482) --(124.3974,67.7781) --(124.4871,67.8080) --(124.5867,67.8378) --(124.6863,67.8677) --(124.7760,67.8976) --(124.8756,67.9175) --(124.9753,67.9474) --(125.0649,67.9773) --(125.1646,68.0072) --(125.2642,68.0371) --(125.3638,68.0670) --(125.4535,68.0969) --(125.5531,68.1268) --(125.6527,68.1567) --(125.7424,68.1865) --(125.8420,68.2164) --(125.9417,68.2463) --(126.0313,68.2762) --(126.1309,68.3061) --(126.2306,68.3360) --(126.3202,68.3659) --(126.4199,68.3958) --(126.5195,68.4257) --(126.6091,68.4556) --(126.7088,68.4854) --(126.8084,68.5153) --(126.8981,68.5452) --(126.9977,68.5751) --(127.0973,68.6149) --(127.1870,68.6448) --(127.2866,68.6747) --(127.3862,68.7046) --(127.4759,68.7445) --(127.5755,68.7744) --(127.6752,68.8042) --(127.7648,68.8341) --(127.8645,68.8740) --(127.9641,68.9039) --(128.0537,68.9337) --(128.1534,68.9736) --(128.2530,69.0035) --(128.3427,69.0434) --(128.4423,69.0732) --(128.5419,69.1131) --(128.6316,69.1430) --(128.7312,69.1828) --(128.8308,69.2127) --(128.9305,69.2526) --(129.0201,69.2825) --(129.1198,69.3223) --(129.2194,69.3622) --(129.3091,69.3920) --(129.4087,69.4319) --(129.5083,69.4618) --(129.5980,69.5016) --(129.6976,69.5415) --(129.7972,69.5813) --(129.8869,69.6112) --(129.9865,69.6511) --(130.0862,69.6909) --(130.1758,69.7308) --(130.2754,69.7706) --(130.3751,69.8105) --(130.4647,69.8503) --(130.5644,69.8802) --(130.6640,69.9201) --(130.7537,69.9599) --(130.8533,69.9998) --(130.9529,70.0396) --(131.0426,70.0894) --(131.1422,70.1293) --(131.2418,70.1691) --(131.3315,70.2090) --(131.4311,70.2488) --(131.5308,70.2887) --(131.6204,70.3285) --(131.7200,70.3784) --(131.8197,70.4182) --(131.9093,70.4581) --(132.0090,70.4979) --(132.1086,70.5477) --(132.1983,70.5876) --(132.2979,70.6274) --(132.3975,70.6772) --(132.4971,70.7171) --(132.5868,70.7669) --(132.6864,70.8068) --(132.7861,70.8566) --(132.8757,70.8964) --(132.9754,70.9462) --(133.0750,70.9861) --(133.1647,71.0359) --(133.2643,71.0857) --(133.3639,71.1256) --(133.4536,71.1754) --(133.5532,71.2252) --(133.6528,71.2650) --(133.7425,71.3148) --(133.8421,71.3647) --(133.9417,71.4145) --(134.0314,71.4643) --(134.1310,71.5042) --(134.2307,71.5540) --(134.3203,71.6038) --(134.4200,71.6536) --(134.5196,71.7034) --(134.6093,71.7532) --(134.7089,71.8030) --(134.8085,71.8528) --(134.8982,71.9026) --(134.9978,71.9525) --(135.0974,72.0023) --(135.1871,72.0621) --(135.2867,72.1119) --(135.3864,72.1617) --(135.4760,72.2115) --(135.5756,72.2613) --(135.6753,72.3211) --(135.7649,72.3709) --(135.8645,72.4207) --(135.9642,72.4805) --(136.0638,72.5303) --(136.1535,72.5801) --(136.2531,72.6399) --(136.3527,72.6897) --(136.4424,72.7495) --(136.5420,72.7993) --(136.6416,72.8591) --(136.7313,72.9089) --(136.8309,72.9687) --(136.9306,73.0185) --(137.0202,73.0783) --(137.1199,73.1380) --(137.2195,73.1879) --(137.3092,73.2476) --(137.4088,73.3074) --(137.5084,73.3672) --(137.5981,73.4170) --(137.6977,73.4768) --(137.7973,73.5365) --(137.8870,73.5963) --(137.9866,73.6561) --(138.0863,73.7059) --(138.1759,73.7657) --(138.2755,73.8255) --(138.3752,73.8853) --(138.4648,73.9450) --(138.5645,74.0048) --(138.6641,74.0646) --(138.7538,74.1244) --(138.8534,74.1841) --(138.9530,74.2439) --(139.0427,74.3037) --(139.1423,74.3635) --(139.2419,74.4332) --(139.3316,74.4930) --(139.4312,74.5527) --(139.5309,74.6125) --(139.6305,74.6723) --(139.7202,74.7421) --(139.8198,74.8018) --(139.9194,74.8616) --(140.0091,74.9214) --(140.1087,74.9911) --(140.2083,75.0509) --(140.2980,75.1107) --(140.3976,75.1804) --(140.4973,75.2402) --(140.5869,75.3099) --(140.6865,75.3697) --(140.7862,75.4295) --(140.8758,75.4992) --(140.9755,75.5590) --(141.0751,75.6287) --(141.1647,75.6885) --(141.2644,75.7583) --(141.3640,75.8280) --(141.4537,75.8878) --(141.5533,75.9575) --(141.6529,76.0173) --(141.7426,76.0870) --(141.8422,76.1568) --(141.9418,76.2165) --(142.0315,76.2863) --(142.1311,76.3560) --(142.2308,76.4158) --(142.3204,76.4855) --(142.4201,76.5553) --(142.5197,76.6250) --(142.6094,76.6848) --(142.7090,76.7545) --(142.8086,76.8243) --(142.8983,76.8940) --(142.9979,76.9637) --(143.0975,77.0335) --(143.1972,77.0933) --(143.2868,77.1630) --(143.3864,77.2327) --(143.4861,77.3025) --(143.5757,77.3722) --(143.6754,77.4420) --(143.7750,77.5117) --(143.8647,77.5814) --(143.9643,77.6512) --(144.0639,77.7209) --(144.1536,77.7906) --(144.2532,77.8604) --(144.3528,77.9301) --(144.4425,77.9999) --(144.5421,78.0696) --(144.6418,78.1394) --(144.7314,78.2091) --(144.8310,78.2788) --(144.9307,78.3486) --(145.0203,78.4183) --(145.1200,78.4881) --(145.2196,78.5677) --(145.3093,78.6375) --(145.4089,78.7072) --(145.5085,78.7770) --(145.5982,78.8467) --(145.6978,78.9165) --(145.7974,78.9962) --(145.8871,79.0659) --(145.9867,79.1356) --(146.0864,79.2054) --(146.1760,79.2751) --(146.2757,79.3548) --(146.3753,79.4245) --(146.4649,79.4943) --(146.5646,79.5640) --(146.6642,79.6437) --(146.7638,79.7135) --(146.8535,79.7832) --(146.9531,79.8530) --(147.0527,79.9326) --(147.1424,80.0024) --(147.2420,80.0721) --(147.3417,80.1518) --(147.4313,80.2216) --(147.5310,80.2913) --(147.6306,80.3710) --(147.7203,80.4408) --(147.8199,80.5105) --(147.9195,80.5802) --(148.0092,80.6599) --(148.1088,80.7297) --(148.2084,80.8094) --(148.2981,80.8791) --(148.3977,80.9489) --(148.4973,81.0285) --(148.5870,81.0983) --(148.6866,81.1680) --(148.7863,81.2477) --(148.8759,81.3175) --(148.9756,81.3872) --(149.0752,81.4669) --(149.1649,81.5367) --(149.2645,81.6163) --(149.3641,81.6861) --(149.4538,81.7558) --(149.5534,81.8355) --(149.6530,81.9053) --(149.7427,81.9750) --(149.8423,82.0547) --(149.9420,82.1245) --(150.0316,82.2042) --(150.1312,82.2739) --(150.2309,82.3436) --(150.3305,82.4233) --(150.4202,82.4931) --(150.5198,82.5728) --(150.6194,82.6425) --(150.7091,82.7123) --(150.8087,82.7920) --(150.9083,82.8617) --(150.9980,82.9314) --(151.0976,83.0112) --(151.1973,83.0809) --(151.2869,83.1606) --(151.3866,83.2303) --(151.4862,83.3001) --(151.5759,83.3798) --(151.6755,83.4495) --(151.7751,83.5192) --(151.8648,83.5990) --(151.9644,83.6687) --(152.0640,83.7384) --(152.1537,83.8181) --(152.2533,83.8879) --(152.3529,83.9676) --(152.4426,84.0373) --(152.5422,84.1070) --(152.6419,84.1868) --(152.7315,84.2565) --(152.8311,84.3262) --(152.9308,84.4059) --(153.0204,84.4757) --(153.1201,84.5454) --(153.2197,84.6152) --(153.3094,84.6949) --(153.4090,84.7646) --(153.5086,84.8343) --(153.5983,84.9140) --(153.6979,84.9838) --(153.7975,85.0535) --(153.8972,85.1232) --(153.9868,85.2030) --(154.0865,85.2727) --(154.1861,85.3424) --(154.2758,85.4122) --(154.3754,85.4919) --(154.4750,85.5616) --(154.5647,85.6313) --(154.6643,85.7011) --(154.7639,85.7708) --(154.8536,85.8505) --(154.9532,85.9203) --(155.0529,85.9900) --(155.1425,86.0598) --(155.2421,86.1295) --(155.3418,86.1992) --(155.4314,86.2789) --(155.5311,86.3487) --(155.6307,86.4184) --(155.7204,86.4881) --(155.8200,86.5579) --(155.9196,86.6276) --(156.0093,86.6974) --(156.1089,86.7671) --(156.2085,86.8369) --(156.2982,86.9066) --(156.3978,86.9763) --(156.4975,87.0461) --(156.5871,87.1158) --(156.6868,87.1855) --(156.7864,87.2553) --(156.8760,87.3250) --(156.9757,87.3948) --(157.0753,87.4645) --(157.1650,87.5342) --(157.2646,87.6040) --(157.3642,87.6737) --(157.4638,87.7435) --(157.5535,87.8132) --(157.6531,87.8829) --(157.7528,87.9527) --(157.8424,88.0224) --(157.9420,88.0921) --(158.0417,88.1519) --(158.1313,88.2217) --(158.2310,88.2914) --(158.3306,88.3611) --(158.4203,88.4309) --(158.5199,88.4907) --(158.6195,88.5604) --(158.7092,88.6302) --(158.8088,88.6999) --(158.9084,88.7597) --(158.9981,88.8294) --(159.0977,88.8991) --(159.1974,88.9589) --(159.2870,89.0287) --(159.3867,89.0984) --(159.4863,89.1582) --(159.5759,89.2279) --(159.6756,89.2977) --(159.7752,89.3574) --(159.8649,89.4272) --(159.9645,89.4870) --(160.0641,89.5567) --(160.1538,89.6165) --(160.2534,89.6862) --(160.3530,89.7460) --(160.4427,89.8157) --(160.5423,89.8755) --(160.6420,89.9452) --(160.7316,90.0050) --(160.8313,90.0648) --(160.9309,90.1345) --(161.0205,90.1943) --(161.1202,90.2640) --(161.2198,90.3238) --(161.3194,90.3836) --(161.4091,90.4533) --(161.5087,90.5131) --(161.6084,90.5729) --(161.6980,90.6327) --(161.7976,90.7024) --(161.8973,90.7622) --(161.9869,90.8220) --(162.0866,90.8817) --(162.1862,90.9415) --(162.2759,91.0112) --(162.3755,91.0710) --(162.4751,91.1308) --(162.5648,91.1906) --(162.6644,91.2504) --(162.7640,91.3101) --(162.8537,91.3699) --(162.9533,91.4297) --(163.0530,91.4895) --(163.1426,91.5492) --(163.2422,91.6090) --(163.3419,91.6688) --(163.4315,91.7286) --(163.5312,91.7883) --(163.6308,91.8481) --(163.7205,91.9079) --(163.8201,91.9677) --(163.9197,92.0275) --(164.0094,92.0773) --(164.1090,92.1370) --(164.2086,92.1968) --(164.2983,92.2566) --(164.3979,92.3164) --(164.4976,92.3662) --(164.5872,92.4260) --(164.6868,92.4857) --(164.7865,92.5356) --(164.8861,92.5953) --(164.9758,92.6551) --(165.0754,92.7049) --(165.1750,92.7647) --(165.2647,92.8245) --(165.3643,92.8743) --(165.4639,92.9341) --(165.5536,92.9839) --(165.6532,93.0437) --(165.7529,93.0935) --(165.8425,93.1532) --(165.9422,93.2031) --(166.0418,93.2628) --(166.1315,93.3127) --(166.2311,93.3724) --(166.3307,93.4222) --(166.4204,93.4820) --(166.5200,93.5318) --(166.6196,93.5816) --(166.7093,93.6414) --(166.8089,93.6912) --(166.9085,93.7410) --(166.9982,93.7909) --(167.0978,93.8506) --(167.1975,93.9005) --(167.2871,93.9503) --(167.3867,94.0001) --(167.4864,94.0599) --(167.5760,94.1097) --(167.6757,94.1595) --(167.7753,94.2093) --(167.8650,94.2591) --(167.9646,94.3089) --(168.0642,94.3587) --(168.1539,94.4086) --(168.2535,94.4584) --(168.3531,94.5082) --(168.4528,94.5580) --(168.5424,94.6078) --(168.6421,94.6576) --(168.7417,94.7074) --(168.8314,94.7572) --(168.9310,94.8071) --(169.0306,94.8569) --(169.1203,94.9067) --(169.2199,94.9565) --(169.3195,95.0063) --(169.4092,95.0462) --(169.5088,95.0960) --(169.6085,95.1458) --(169.6981,95.1956) --(169.7977,95.2355) --(169.8974,95.2853) --(169.9870,95.3351) --(170.0867,95.3849) --(170.1863,95.4248) --(170.2760,95.4746) --(170.3756,95.5244) --(170.4752,95.5642) --(170.5649,95.6140) --(170.6645,95.6539) --(170.7641,95.7037) --(170.8538,95.7535) --(170.9534,95.7934) --(171.0531,95.8432) --(171.1427,95.8830) --(171.2424,95.9329) --(171.3420,95.9727) --(171.4316,96.0126) --(171.5313,96.0624) --(171.6309,96.1022) --(171.7206,96.1520) --(171.8202,96.1919) --(171.9198,96.2317) --(172.0195,96.2816) --(172.1091,96.3214) --(172.2087,96.3613) --(172.3084,96.4111) --(172.3980,96.4509) --(172.4977,96.4908) --(172.5973,96.5406) --(172.6870,96.5804) --(172.7866,96.6203) --(172.8862,96.6601) --(172.9759,96.7000) --(173.0755,96.7498) --(173.1751,96.7897) --(173.2648,96.8295) --(173.3644,96.8694) --(173.4641,96.9092) --(173.5537,96.9491) --(173.6533,96.9889) --(173.7530,97.0288) --(173.8426,97.0686) --(173.9423,97.1085) --(174.0419,97.1483) --(174.1315,97.1882) --(174.2312,97.2280) --(174.3308,97.2679) --(174.4205,97.3077) --(174.5201,97.3476) --(174.6197,97.3874) --(174.7094,97.4273) --(174.8090,97.4671) --(174.9086,97.5070) --(174.9983,97.5369) --(175.0979,97.5767) --(175.1976,97.6166) --(175.2872,97.6564) --(175.3869,97.6963) --(175.4865,97.7262) --(175.5861,97.7660) --(175.6758,97.8058) --(175.7754,97.8457) --(175.8750,97.8756) --(175.9647,97.9155) --(176.0643,97.9553) --(176.1640,97.9852) --(176.2536,98.0250) --(176.3532,98.0649) --(176.4529,98.0948) --(176.5425,98.1346) --(176.6422,98.1745) --(176.7418,98.2044) --(176.8315,98.2442) --(176.9311,98.2741) --(177.0307,98.3140) --(177.1204,98.3439) --(177.2200,98.3837) --(177.3196,98.4136) --(177.4093,98.4534) --(177.5089,98.4833) --(177.6086,98.5232) --(177.6982,98.5531) --(177.7979,98.5929) --(177.8975,98.6228) --(177.9871,98.6527) --(178.0868,98.6925) --(178.1864,98.7224) --(178.2761,98.7523) --(178.3757,98.7922) --(178.4753,98.8221) --(178.5650,98.8519) --(178.6646,98.8918) --(178.7642,98.9217) --(178.8539,98.9516) --(178.9535,98.9815) --(179.0532,99.0213) --(179.1528,99.0512) --(179.2425,99.0811) --(179.3421,99.1110) --(179.4417,99.1508) --(179.5314,99.1807) --(179.6310,99.2106) --(179.7306,99.2405) --(179.8203,99.2704) --(179.9199,99.3003) --(180.0195,99.3302) --(180.1092,99.3600) --(180.2088,99.3999) --(180.3085,99.4298) --(180.3981,99.4597) --(180.4978,99.4896) --(180.5974,99.5195) --(180.6871,99.5493) --(180.7867,99.5792) --(180.8863,99.6091) --(180.9760,99.6390) --(181.0756,99.6689) --(181.1752,99.6988) --(181.2649,99.7287) --(181.3645,99.7586) --(181.4642,99.7785) --(181.5538,99.8084) --(181.6534,99.8383) --(181.7531,99.8681) --(181.8427,99.8980) --(181.9424,99.9279) --(182.0420,99.9578) --(182.1317,99.9877) --(182.2313,100.0076) --(182.3309,100.0375) --(182.4206,100.0674) --(182.5202,100.0973) --(182.6198,100.1272) --(182.7195,100.1471) --(182.8091,100.1770) --(182.9088,100.2069) --(183.0084,100.2368) --(183.0980,100.2567) --(183.1977,100.2866) --(183.2973,100.3165) --(183.3870,100.3364) --(183.4866,100.3663) --(183.5862,100.3962) --(183.6759,100.4161) --(183.7755,100.4460) --(183.8751,100.4759) --(183.9648,100.4958) --(184.0645,100.5257) --(184.1641,100.5456) --(184.2537,100.5755) --(184.3533,100.5954) --(184.4530,100.6253) --(184.5427,100.6552) --(184.6423,100.6751) --(184.7419,100.7050) --(184.8316,100.7249) --(184.9312,100.7548) --(185.0308,100.7748) --(185.1205,100.8047) --(185.2201,100.8246) --(185.3197,100.8545) --(185.4094,100.8744) --(185.5090,100.8943) --(185.6086,100.9242) --(185.6983,100.9441) --(185.7980,100.9740) --(185.8976,100.9939) --(185.9872,101.0238) --(186.0869,101.0437) --(186.1865,101.0637) --(186.2861,101.0936) --(186.3758,101.1135) --(186.4754,101.1334) --(186.5751,101.1633) --(186.6647,101.1832) --(186.7643,101.2032) --(186.8640,101.2330) --(186.9536,101.2530) --(187.0533,101.2729) --(187.1529,101.2928) --(187.2426,101.3227) --(187.3422,101.3426) --(187.4418,101.3626) --(187.5315,101.3825) --(187.6311,101.4124) --(187.7307,101.4323) --(187.8204,101.4522) --(187.9200,101.4722) --(188.0197,101.4921) --(188.1093,101.5220) --(188.2090,101.5419) --(188.3086,101.5618) --(188.3983,101.5818) --(188.4979,101.6017) --(188.5975,101.6216) --(188.6872,101.6415) --(188.7868,101.6714) --(188.8864,101.6913) --(188.9761,101.7113) --(189.0757,101.7312) --(189.1754,101.7511) --(189.2650,101.7710) --(189.3646,101.7910) --(189.4642,101.8109) --(189.5539,101.8308) --(189.6536,101.8507) --(189.7532,101.8707) --(189.8528,101.8906) --(189.9424,101.9105) --(190.0421,101.9304) --(190.1417,101.9504) --(190.2314,101.9703) --(190.3310,101.9902) --(190.4307,102.0101) --(190.5203,102.0301) --(190.6199,102.0500) --(190.7195,102.0699) --(190.8092,102.0898) --(190.9088,102.1098) --(191.0085,102.1297) --(191.0981,102.1496) --(191.1978,102.1696) --(191.2974,102.1895) --(191.3871,102.1994) --(191.4867,102.2194) --(191.5863,102.2393) --(191.6760,102.2592) --(191.7756,102.2791) --(191.8752,102.2991) --(191.9649,102.3190) --(192.0645,102.3389) --(192.1642,102.3489) --(192.2538,102.3688) --(192.3535,102.3887) --(192.4531,102.4087) --(192.5427,102.4286) --(192.6424,102.4386) --(192.7420,102.4585) --(192.8317,102.4784) --(192.9313,102.4983) --(193.0309,102.5182) --(193.1206,102.5282) --(193.2202,102.5481) --(193.3198,102.5681) --(193.4195,102.5880) --(193.5092,102.5979) --(193.6088,102.6179) --(193.7084,102.6378) --(193.7981,102.6577) --(193.8977,102.6677) --(193.9973,102.6876) --(194.0870,102.7075) --(194.1866,102.7175) --(194.2862,102.7374) --(194.3759,102.7574) --(194.4755,102.7773) --(194.5751,102.7872) --(194.6648,102.8072) --(194.7645,102.8271) --(194.8641,102.8371) --(194.9537,102.8570) --(195.0533,102.8769) --(195.1530,102.8869) --(195.2427,102.9068) --(195.3423,102.9267) --(195.4419,102.9367) --(195.5316,102.9566) --(195.6312,102.9666) --(195.7308,102.9865) --(195.8205,103.0064) --(195.9201,103.0164) --(196.0198,103.0363) --(196.1094,103.0562) --(196.2090,103.0662) --(196.3087,103.0861) --(196.3983,103.0961) --(196.4980,103.1160) --(196.5976,103.1359) --(196.6873,103.1459) --(196.7869,103.1658) --(196.8865,103.1758) --(196.9861,103.1957) --(197.0758,103.2057) --(197.1754,103.2256) --(197.2751,103.2455) --(197.3647,103.2555) --(197.4644,103.2754) --(197.5640,103.2854) --(197.6537,103.3053) --(197.7533,103.3153) --(197.8529,103.3352) --(197.9426,103.3452) --(198.0422,103.3651) --(198.1418,103.3750) --(198.2315,103.3950) --(198.3311,103.4049) --(198.4307,103.4249) --(198.5204,103.4448) --(198.6201,103.4547) --(198.7197,103.4747) --(198.8093,103.4846) --(198.9089,103.5046) --(199.0086,103.5145) --(199.0983,103.5344) --(199.1979,103.5444) --(199.2975,103.5643) --(199.3872,103.5743) --(199.4868,103.5942) --(199.5864,103.6042) --(199.6761,103.6241) --(199.7757,103.6341) --(199.8753,103.6540) --(199.9650,103.6640) --(200.0646,103.6739) --(200.1643,103.6938) --(200.2539,103.7038) --(200.3536,103.7237) --(200.4532,103.7337) --(200.5528,103.7536) --(200.6425,103.7636) --(200.7421,103.7835) --(200.8417,103.7935) --(200.9314,103.8134) --(201.0310,103.8234) --(201.1307,103.8433) --(201.2203,103.8532) --(201.3199,103.8732) --(201.4196,103.8831) --(201.5092,103.9031) --(201.6089,103.9130) --(201.7085,103.9230) --(201.7982,103.9429) --(201.8978,103.9529) --(201.9974,103.9728) --(202.0871,103.9828) --(202.1867,104.0027) --(202.2863,104.0127) --(202.3760,104.0326) --(202.4756,104.0425) --(202.5753,104.0625) --(202.6649,104.0724) --(202.7646,104.0824) --(202.8642,104.1023) --(202.9539,104.1123) --(203.0535,104.1322) --(203.1531,104.1422) --(203.2428,104.1621) --(203.3424,104.1721) --(203.4420,104.1920) --(203.5317,104.2020) --(203.6313,104.2119) --(203.7310,104.2318) --(203.8206,104.2418) --(203.9202,104.2617) --(204.0198,104.2717) --(204.1195,104.2916) --(204.2092,104.3016) --(204.3088,104.3215) --(204.4084,104.3315) --(204.4980,104.3514) --(204.5977,104.3614) --(204.6973,104.3713) --(204.7870,104.3913) --(204.8866,104.4012) --(204.9863,104.4211) --(205.0759,104.4311) --(205.1755,104.4510) --(205.2751,104.4610) --(205.3648,104.4809) --(205.4645,104.4909) --(205.5641,104.5108) --(205.6537,104.5208) --(205.7534,104.5407) --(205.8530,104.5506) --(205.9427,104.5706) --(206.0423,104.5805) --(206.1419,104.6005) --(206.2316,104.6104) --(206.3312,104.6304) --(206.4308,104.6403) --(206.5205,104.6503) --(206.6201,104.6702) --(206.7198,104.6802) --(206.8094,104.7001) --(206.9091,104.7101) --(207.0087,104.7300) --(207.0983,104.7399) --(207.1980,104.7599) --(207.2976,104.7698) --(207.3873,104.7898) --(207.4869,104.8097) --(207.5865,104.8196) --(207.6862,104.8396) --(207.7758,104.8495) --(207.8754,104.8695) --(207.9751,104.8794) --(208.0648,104.8994) --(208.1644,104.9093) --(208.2640,104.9292) --(208.3537,104.9392) --(208.4533,104.9591) --(208.5529,104.9691) --(208.6426,104.9890) --(208.7422,104.9990) --(208.8418,105.0189) --(208.9315,105.0388) --(209.0311,105.0488) --(209.1308,105.0687) --(209.2204,105.0787) --(209.3201,105.0986) --(209.4197,105.1086) --(209.5093,105.1285) --(209.6090,105.1484) --(209.7086,105.1584) --(209.7983,105.1783) --(209.8979,105.1883) --(209.9975,105.2082) --(210.0872,105.2281) --(210.1868,105.2381) --(210.2864,105.2580) --(210.3761,105.2779) --(210.4757,105.2879) --(210.5754,105.3078) --(210.6650,105.3178) --(210.7646,105.3377) --(210.8643,105.3576) --(210.9539,105.3676) --(211.0536,105.3875) --(211.1532,105.4075) --(211.2528,105.4174) --(211.3425,105.4373) --(211.4421,105.4573) --(211.5417,105.4672) --(211.6314,105.4872) --(211.7310,105.5071) --(211.8307,105.5170) --(211.9203,105.5370) --(212.0200,105.5569) --(212.1196,105.5669) --(212.2093,105.5868) --(212.3089,105.6067) --(212.4085,105.6266) --(212.4982,105.6366) --(212.5978,105.6565) --(212.6974,105.6765) --(212.7871,105.6864) --(212.8867,105.7063) --(212.9863,105.7263) --(213.0760,105.7462) --(213.1757,105.7561) --(213.2753,105.7761) --(213.3649,105.7960) --(213.4645,105.8159) --(213.5642,105.8358) --(213.6539,105.8458) --(213.7535,105.8658) --(213.8531,105.8857) --(213.9428,105.9056) --(214.0424,105.9255) --(214.1420,105.9355) --(214.2317,105.9554) --(214.3313,105.9753) --(214.4310,105.9953) --(214.5206,106.0152) --(214.6202,106.0251) --(214.7199,106.0451) --(214.8195,106.0650) --(214.9092,106.0849) --(215.0088,106.1048) --(215.1084,106.1248) --(215.1981,106.1447) --(215.2977,106.1547) --(215.3973,106.1746) --(215.4870,106.1945) --(215.5866,106.2144) --(215.6863,106.2344) --(215.7759,106.2543) --(215.8755,106.2742) --(215.9752,106.2941) --(216.0648,106.3141) --(216.1645,106.3340) --(216.2641,106.3440) --(216.3538,106.3639) --(216.4534,106.3838) --(216.5530,106.4037) --(216.6427,106.4237) --(216.7423,106.4436) --(216.8419,106.4635) --(216.9316,106.4834) --(217.0312,106.5034) --(217.1309,106.5233) --(217.2205,106.5432) --(217.3202,106.5631) --(217.4198,106.5831) --(217.5095,106.6030) --(217.6091,106.6229) --(217.7087,106.6429) --(217.7984,106.6628) --(217.8980,106.6827) --(217.9976,106.7026) --(218.0873,106.7225) --(218.1869,106.7425) --(218.2866,106.7624) --(218.3862,106.7923) --(218.4758,106.8122) --(218.5755,106.8321) --(218.6751,106.8521) --(218.7648,106.8720) --(218.8644,106.8919) --(218.9640,106.9118) --(219.0536,106.9318) --(219.1533,106.9517) --(219.2529,106.9716) --(219.3426,107.0015) --(219.4422,107.0214) --(219.5419,107.0413) --(219.6315,107.0613) --(219.7311,107.0812) --(219.8307,107.1011) --(219.9204,107.1310) --(220.0201,107.1509) --(220.1197,107.1709) --(220.2094,107.1908) --(220.3090,107.2107) --(220.4086,107.2406) --(220.4983,107.2605) --(220.5979,107.2805) --(220.6975,107.3004) --(220.7872,107.3303) --(220.8868,107.3502) --(220.9865,107.3701) --(221.0761,107.3900) --(221.1758,107.4199) --(221.2754,107.4398) --(221.3651,107.4598) --(221.4647,107.4797) --(221.5643,107.5096) --(221.6540,107.5295) --(221.7536,107.5494) --(221.8532,107.5793) --(221.9528,107.5993) --(222.0425,107.6192) --(222.1422,107.6491) --(222.2418,107.6690) --(222.3314,107.6989) --(222.4310,107.7188) --(222.5307,107.7387) --(222.6204,107.7686) --(222.7200,107.7885) --(222.8196,107.8085) --(222.9093,107.8384) --(223.0089,107.8583) --(223.1085,107.8882) --(223.1982,107.9081) --(223.2978,107.9380) --(223.3975,107.9579) --(223.4871,107.9878) --(223.5867,108.0077) --(223.6864,108.0376) --(223.7760,108.0576) --(223.8757,108.0775) --(223.9753,108.1074) --(224.0649,108.1373) --(224.1646,108.1572) --(224.2642,108.1871) --(224.3539,108.2070) --(224.4535,108.2369) --(224.5531,108.2568) --(224.6428,108.2867) --(224.7424,108.3066) --(224.8420,108.3365) --(224.9317,108.3664) --(225.0313,108.3863) --(225.1310,108.4162) --(225.2206,108.4361) --(225.3203,108.4660) --(225.4199,108.4959) --(225.5195,108.5158) --(225.6092,108.5457) --(225.7088,108.5756) --(225.8084,108.5955) --(225.8981,108.6254) --(225.9977,108.6553) --(226.0974,108.6752) --(226.1870,108.7051) --(226.2867,108.7350) --(226.3863,108.7549) --(226.4760,108.7848) --(226.5756,108.8147) --(226.6752,108.8446) --(226.7649,108.8645) --(226.8645,108.8944) --(226.9641,108.9243) --(227.0538,108.9542) --(227.1534,108.9841) --(227.2531,109.0040) --(227.3427,109.0339) --(227.4423,109.0638) --(227.5420,109.0937) --(227.6316,109.1236) --(227.7313,109.1534) --(227.8309,109.1734) --(227.9205,109.2033) --(228.0201,109.2332) --(228.1198,109.2630) --(228.2095,109.2929) --(228.3091,109.3228) --(228.4087,109.3527) --(228.4984,109.3826) --(228.5980,109.4125) --(228.6976,109.4424) --(228.7873,109.4723) --(228.8869,109.5022) --(228.9866,109.5320) --(229.0862,109.5619) --(229.1758,109.5918) --(229.2755,109.6217) --(229.3751,109.6516) --(229.4648,109.6815) --(229.5644,109.7114) --(229.6640,109.7412) --(229.7537,109.7711) --(229.8533,109.8010) --(229.9529,109.8309) --(230.0426,109.8608) --(230.1422,109.8907) --(230.2419,109.9206) --(230.3315,109.9505) --(230.4312,109.9804) --(230.5308,110.0102) --(230.6205,110.0501) --(230.7201,110.0800) --(230.8197,110.1099) --(230.9094,110.1398) --(231.0090,110.1696) --(231.1086,110.1995) --(231.1983,110.2394) --(231.2979,110.2693) --(231.3975,110.2992) --(231.4872,110.3291) --(231.5869,110.3589) --(231.6865,110.3988) --(231.7761,110.4287) --(231.8758,110.4586) --(231.9754,110.4885) --(232.0651,110.5283) --(232.1647,110.5582) --(232.2643,110.5881) --(232.3540,110.6279) --(232.4536,110.6578) --(232.5532,110.6877) --(232.6429,110.7276) --(232.7425,110.7574) --(232.8422,110.7874) --(232.9418,110.8272) --(233.0314,110.8571) --(233.1311,110.8870) --(233.2307,110.9268) --(233.3204,110.9567) --(233.4200,110.9966) --(233.5196,111.0264) --(233.6093,111.0563) --(233.7089,111.0962) --(233.8085,111.1261) --(233.8982,111.1659) --(233.9978,111.1958) --(234.0975,111.2357) --(234.1871,111.2656) --(234.2867,111.3054) --(234.3864,111.3353) --(234.4760,111.3752) --(234.5757,111.4050) --(234.6753,111.4449) --(234.7650,111.4748) --(234.8646,111.5146) --(234.9642,111.5445) --(235.0539,111.5844) --(235.1535,111.6143) --(235.2531,111.6541) --(235.3428,111.6840) --(235.4424,111.7238) --(235.5421,111.7637) --(235.6317,111.7936) --(235.7314,111.8334) --(235.8310,111.8633) --(235.9207,111.9032) --(236.0203,111.9430) --(236.1199,111.9729) --(236.2096,112.0128) --(236.3092,112.0526) --(236.4088,112.0825) --(236.5085,112.1224) --(236.5981,112.1622) --(236.6978,112.2021) --(236.7974,112.2319) --(236.8870,112.2718) --(236.9866,112.3116) --(237.0863,112.3416) --(237.1760,112.3814) --(237.2756,112.4212) --(237.3752,112.4611) --(237.4649,112.4910) --(237.5645,112.5308) --(237.6641,112.5707) --(237.7538,112.6105) --(237.8534,112.6504) --(237.9531,112.6803) --(238.0427,112.7201) --(238.1423,112.7600) --(238.2420,112.7998) --(238.3316,112.8397) --(238.4313,112.8795) --(238.5309,112.9194) --(238.6205,112.9493) --(238.7202,112.9891) --(238.8198,113.0290) --(238.9095,113.0688) --(239.0091,113.1087) --(239.1087,113.1485) --(239.1984,113.1884) --(239.2980,113.2282) --(239.3976,113.2681) --(239.4873,113.3079) --(239.5869,113.3478) --(239.6866,113.3876) --(239.7762,113.4275) --(239.8759,113.4673) --(239.9755,113.5072) --(240.0751,113.5470) --(240.1648,113.5869) --(240.2644,113.6267) --(240.3640,113.6666) --(240.4537,113.7064) --(240.5533,113.7463) --(240.6530,113.7862) --(240.7426,113.8260) --(240.8423,113.8658) --(240.9419,113.9057) --(241.0316,113.9455) --(241.1312,113.9854) --(241.2308,114.0252) --(241.3205,114.0651) --(241.4201,114.1049) --(241.5197,114.1548) --(241.6094,114.1946) --(241.7090,114.2345) --(241.8087,114.2743) --(241.8983,114.3142) --(241.9979,114.3540) --(242.0976,114.4038) --(242.1872,114.4437) --(242.2869,114.4835) --(242.3865,114.5234) --(242.4761,114.5632) --(242.5758,114.6131) --(242.6754,114.6529) --(242.7651,114.6927) --(242.8647,114.7326) --(242.9643,114.7824) --(243.0540,114.8223) --(243.1536,114.8621) --(243.2532,114.9020) --(243.3429,114.9518) --(243.4425,114.9916) --(243.5422,115.0315) --(243.6418,115.0813) --(243.7314,115.1212) --(243.8311,115.1610) --(243.9307,115.2009) --(244.0204,115.2507) --(244.1200,115.2905) --(244.2196,115.3403) --(244.3093,115.3802) --(244.4089,115.4200) --(244.5085,115.4698) --(244.5982,115.5097) --(244.6978,115.5495) --(244.7975,115.5994) --(244.8871,115.6392) --(244.9868,115.6890) --(245.0864,115.7289) --(245.1761,115.7787) --(245.2757,115.8185) --(245.3753,115.8584) --(245.4650,115.9082) --(245.5646,115.9481) --(245.6642,115.9979) --(245.7539,116.0377) --(245.8535,116.0875) --(245.9531,116.1274) --(246.0428,116.1772) --(246.1425,116.2170) --(246.2421,116.2669) --(246.3317,116.3067) --(246.4314,116.3565) --(246.5310,116.4063) --(246.6207,116.4462) --(246.7203,116.4960) --(246.8199,116.5359) --(246.9096,116.5857) --(247.0092,116.6255) --(247.1088,116.6753) --(247.2085,116.7252) --(247.2981,116.7650) --(247.3978,116.8148) --(247.4974,116.8547) --(247.5870,116.9045) --(247.6867,116.9543) --(247.7863,116.9942) --(247.8760,117.0440) --(247.9756,117.0938) --(248.0752,117.1336) --(248.1649,117.1834) --(248.2645,117.2332) --(248.3641,117.2731) --(248.4538,117.3229) --(248.5534,117.3727) --(248.6531,117.4225) --(248.7427,117.4624) --(248.8423,117.5122) --(248.9420,117.5620) --(249.0316,117.6118) --(249.1313,117.6517) --(249.2309,117.7015) --(249.3206,117.7513) --(249.4202,117.8011) --(249.5198,117.8410) --(249.6095,117.8908) --(249.7091,117.9406) --(249.8087,117.9904) --(249.8984,118.0403) --(249.9980,118.0801) --(250.0977,118.1299) --(250.1873,118.1797) --(250.2870,118.2295) --(250.3866,118.2793) --(250.4763,118.3292) --(250.5759,118.3790) --(250.6755,118.4188) --(250.7751,118.4686) --(250.8648,118.5185) --(250.9644,118.5683) --(251.0641,118.6181) --(251.1537,118.6679) --(251.2534,118.7177) --(251.3530,118.7675) --(251.4426,118.8173) --(251.5423,118.8671) --(251.6419,118.9170) --(251.7316,118.9668) --(251.8312,119.0166) --(251.9308,119.0664) --(252.0205,119.1162) --(252.1201,119.1660) --(252.2197,119.2159) --(252.3094,119.2657) --(252.4090,119.3155) --(252.5087,119.3653) --(252.5983,119.4151) --(252.6979,119.4649) --(252.7976,119.5147) --(252.8872,119.5645) --(252.9869,119.6143) --(253.0865,119.6642) --(253.1761,119.7140) --(253.2758,119.7638) --(253.3754,119.8136) --(253.4651,119.8634) --(253.5647,119.9132) --(253.6643,119.9631) --(253.7540,120.0129) --(253.8536,120.0726) --(253.9532,120.1225) --(254.0429,120.1723) --(254.1425,120.2221) --(254.2422,120.2719) --(254.3418,120.3217) --(254.4315,120.3715) --(254.5311,120.4313) --(254.6307,120.4811) --(254.7204,120.5309) --(254.8200,120.5807) --(254.9196,120.6306) --(255.0093,120.6903) --(255.1089,120.7401) --(255.2086,120.7900) --(255.2982,120.8398) --(255.3979,120.8996) --(255.4975,120.9494) --(255.5872,120.9992) --(255.6868,121.0490) --(255.7864,121.1088) --(255.8761,121.1586) --(255.9757,121.2084) --(256.0753,121.2582) --(256.1650,121.3180) --(256.2646,121.3678) --(256.3643,121.4176) --(256.4539,121.4774) --(256.5535,121.5272) --(256.6532,121.5770) --(256.7428,121.6368) --(256.8425,121.6866) --(256.9421,121.7364) --(257.0318,121.7962) --(257.1314,121.8460) --(257.2310,121.8958) --(257.3207,121.9556) --(257.4203,122.0054) --(257.5199,122.0652) --(257.6096,122.1150) --(257.7092,122.1648) --(257.8089,122.2246) --(257.9085,122.2744) --(257.9981,122.3342) --(258.0978,122.3840) --(258.1974,122.4338) --(258.2871,122.4936) --(258.3867,122.5434) --(258.4863,122.6032) --(258.5760,122.6530) --(258.6756,122.7128) --(258.7752,122.7626) --(258.8649,122.8224) --(258.9645,122.8722) --(259.0641,122.9320) --(259.1538,122.9818) --(259.2534,123.0415) --(259.3531,123.0914) --(259.4427,123.1512) --(259.5424,123.2010) --(259.6420,123.2607) --(259.7317,123.3105) --(259.8313,123.3703) --(259.9309,123.4301) --(260.0206,123.4799) --(260.1202,123.5397) --(260.2198,123.5895) --(260.3095,123.6493) --(260.4091,123.6991) --(260.5088,123.7589) --(260.5984,123.8186) --(260.6981,123.8685) --(260.7977,123.9283) --(260.8874,123.9780) --(260.9870,124.0378) --(261.0866,124.0976) --(261.1763,124.1474) --(261.2759,124.2072) --(261.3755,124.2670) --(261.4752,124.3168) --(261.5648,124.3766) --(261.6645,124.4363) --(261.7641,124.4861) --(261.8537,124.5459) --(261.9534,124.6057) --(262.0530,124.6655) --(262.1426,124.7153) --(262.2423,124.7751) --(262.3419,124.8349) --(262.4316,124.8847) --(262.5312,124.9444) --(262.6308,125.0042) --(262.7205,125.0640) --(262.8201,125.1138) --(262.9197,125.1736) --(263.0094,125.2334) --(263.1090,125.2931) --(263.2087,125.3429) --(263.2983,125.4027) --(263.3980,125.4625) --(263.4976,125.5223) --(263.5872,125.5821) --(263.6869,125.6418) --(263.7865,125.6917) --(263.8762,125.7514) --(263.9758,125.8112) --(264.0754,125.8710) --(264.1651,125.9308) --(264.2647,125.9905) --(264.3643,126.0403) --(264.4540,126.1001) --(264.5536,126.1599) --(264.6533,126.2197) --(264.7430,126.2794) --(264.8426,126.3392) --(264.9422,126.3990) --(265.0418,126.4588) --(265.1315,126.5086) --(265.2311,126.5684) --(265.3307,126.6281) --(265.4204,126.6879) --(265.5200,126.7477) --(265.6197,126.8075) --(265.7093,126.8672) --(265.8090,126.9270) --(265.9086,126.9868) --(265.9983,127.0466) --(266.0979,127.1064) --(266.1975,127.1661) --(266.2872,127.2259) --(266.3868,127.2857) --(266.4864,127.3455) --(266.5761,127.4052) --(266.6757,127.4650) --(266.7754,127.5248) --(266.8650,127.5846) --(266.9646,127.6443) --(267.0643,127.7041) --(267.1539,127.7639) --(267.2536,127.8237) --(267.3532,127.8835) --(267.4428,127.9432) --(267.5425,128.0030) --(267.6421,128.0628) --(267.7318,128.1226) --(267.8314,128.1823) --(267.9310,128.2521) --(268.0207,128.3119) --(268.1203,128.3716) --(268.2199,128.4314) --(268.3096,128.4912) --(268.4092,128.5510) --(268.5089,128.6107) --(268.6085,128.6705) --(268.6982,128.7303) --(268.7978,128.8000) --(268.8974,128.8598) --(268.9871,128.9196) --(269.0867,128.9794) --(269.1863,129.0391) --(269.2760,129.0989) --(269.3756,129.1687) --(269.4753,129.2284) --(269.5649,129.2882) --(269.6646,129.3480) --(269.7642,129.4078) --(269.8539,129.4775) --(269.9535,129.5373) --(270.0531,129.5970) --(270.1428,129.6568) --(270.2424,129.7166) --(270.3420,129.7863) --(270.4317,129.8461) --(270.5313,129.9059) --(270.6310,129.9657) --(270.7206,130.0354) --(270.8202,130.0952) --(270.9199,130.1550) --(271.0095,130.2147) --(271.1091,130.2845) --(271.2088,130.3443) --(271.2984,130.4041) --(271.3981,130.4738) --(271.4977,130.5336) --(271.5874,130.5933) --(271.6870,130.6631) --(271.7866,130.7229) --(271.8763,130.7826) --(271.9759,130.8424) --(272.0755,130.9121) --(272.1752,130.9719) --(272.2648,131.0417) --(272.3645,131.1014) --(272.4641,131.1612) --(272.5537,131.2309) --(272.6534,131.2907) --(272.7530,131.3505) --(272.8427,131.4202) --(272.9423,131.4800) --(273.0419,131.5498) --(273.1316,131.6095) --(273.2312,131.6693) --(273.3308,131.7390) --(273.4205,131.7988) --(273.5201,131.8686) --(273.6198,131.9283) --(273.7094,131.9881) --(273.8091,132.0579) --(273.9087,132.1176) --(273.9984,132.1874) --(274.0980,132.2471) --(274.1976,132.3169) --(274.2873,132.3767) --(274.3869,132.4464) --(274.4865,132.5062) --(274.5762,132.5759) --(274.6758,132.6357) --(274.7755,132.7054) --(274.8651,132.7652) --(274.9648,132.8349) --(275.0644,132.8947) --(275.1540,132.9645) --(275.2537,133.0242) --(275.3533,133.0940) --(275.4429,133.1538) --(275.5426,133.2235) --(275.6422,133.2833) --(275.7419,133.3530) --(275.8315,133.4128) --(275.9311,133.4825) --(276.0308,133.5423) --(276.1204,133.6120) --(276.2201,133.6818) --(276.3197,133.7416) --(276.4093,133.8113) --(276.5090,133.8711) --(276.6086,133.9408) --(276.6982,134.0006) --(276.7979,134.0703) --(276.8975,134.1401) --(276.9872,134.1998) --(277.0868,134.2696) --(277.1864,134.3393) --(277.2761,134.3991) --(277.3757,134.4688) --(277.4753,134.5286) --(277.5650,134.5984) --(277.6646,134.6681) --(277.7643,134.7279) --(277.8539,134.7976) --(277.9536,134.8674) --(278.0532,134.9271) --(278.1429,134.9969) --(278.2425,135.0666) --(278.3421,135.1264) --(278.4318,135.1961) --(278.5314,135.2659) --(278.6310,135.3257) --(278.7207,135.3954) --(278.8203,135.4651) --(278.9200,135.5349) --(279.0096,135.5946) --(279.1092,135.6644) --(279.2089,135.7341) --(279.3085,135.8039) --(279.3982,135.8636) --(279.4978,135.9334) --(279.5974,136.0031) --(279.6871,136.0728) --(279.7867,136.1326) --(279.8863,136.2024) --(279.9760,136.2721) --(280.0756,136.3418) --(280.1753,136.4016) --(280.2649,136.4714) --(280.3646,136.5411) --(280.4642,136.6108) --(280.5538,136.6806) --(280.6535,136.7404) --(280.7531,136.8101) --(280.8428,136.8798) --(280.9424,136.9496) --(281.0420,137.0193) --(281.1317,137.0791) --(281.2313,137.1488) --(281.3309,137.2186) --(281.4206,137.2883) --(281.5202,137.3581) --(281.6198,137.4278) --(281.7095,137.4975) --(281.8091,137.5573) --(281.9088,137.6270) --(281.9984,137.6968) --(282.0981,137.7665) --(282.1977,137.8363) --(282.2874,137.9060) --(282.3870,137.9757) --(282.4866,138.0455) --(282.5763,138.1053) --(282.6759,138.1750) --(282.7755,138.2447) --(282.8752,138.3145) --(282.9648,138.3842) --(283.0645,138.4539) --(283.1641,138.5237) --(283.2538,138.5934) --(283.3534,138.6632) --(283.4530,138.7329) --(283.5427,138.8027) --(283.6423,138.8724) --(283.7419,138.9421) --(283.8316,139.0119) --(283.9312,139.0816) --(284.0309,139.1513) --(284.1205,139.2111) --(284.2201,139.2809) --(284.3198,139.3506) --(284.4094,139.4203) --(284.5091,139.4901) --(284.6087,139.5598) --(284.6983,139.6296) --(284.7980,139.6993) --(284.8976,139.7690) --(284.9873,139.8388) --(285.0869,139.9085) --(285.1865,139.9782) --(285.2762,140.0580) --(285.3758,140.1277) --(285.4754,140.1974) --(285.5651,140.2672) --(285.6647,140.3369) --(285.7644,140.4067) --(285.8540,140.4764) --(285.9537,140.5461) --(286.0533,140.6159) --(286.1429,140.6856) --(286.2426,140.7554) --(286.3422,140.8251) --(286.4418,140.8948) --(286.5315,140.9646) --(286.6311,141.0343) --(286.7307,141.1040) --(286.8204,141.1838) --(286.9200,141.2535) --(287.0197,141.3232) --(287.1093,141.3930) --(287.2090,141.4627) --(287.3086,141.5324) --(287.3983,141.6022) --(287.4979,141.6719) --(287.5975,141.7516) --(287.6872,141.8214) --(287.7868,141.8911) --(287.8864,141.9608) --(287.9761,142.0306) --(288.0757,142.1003) --(288.1754,142.1701) --(288.2650,142.2498) --(288.3647,142.3195) --(288.4643,142.3892) --(288.5540,142.4590) --(288.6536,142.5287) --(288.7532,142.5985) --(288.8429,142.6782) --(288.9425,142.7479) --(289.0421,142.8176) --(289.1318,142.8874) --(289.2314,142.9571) --(289.3311,143.0368) --(289.4207,143.1066) --(289.5203,143.1763) --(289.6200,143.2461) --(289.7096,143.3257) --(289.8093,143.3955) --(289.9089,143.4652) --(290.0085,143.5350) --(290.0982,143.6047) --(290.1978,143.6844) --(290.2974,143.7542) --(290.3871,143.8239) --(290.4867,143.8936) --(290.5863,143.9733) --(290.6760,144.0431) --(290.7756,144.1128) --(290.8753,144.1925) --(290.9649,144.2623) --(291.0645,144.3320) --(291.1642,144.4017) --(291.2538,144.4814) --(291.3535,144.5512) --(291.4531,144.6209) --(291.5428,144.7006) --(291.6424,144.7703) --(291.7420,144.8401) --(291.8317,144.9198) --(291.9313,144.9895) --(292.0309,145.0593) --(292.1206,145.1290) --(292.2202,145.2087) --(292.3199,145.2785) --(292.4095,145.3482) --(292.5092,145.4279) --(292.6088,145.4976) --(292.6985,145.5773) --(292.7981,145.6471) --(292.8977,145.7168) --(292.9874,145.7965) --(293.0870,145.8662) --(293.1866,145.9360) --(293.2763,146.0157) --(293.3759,146.0854) --(293.4756,146.1552) --(293.5752,146.2349) --(293.6648,146.3046) --(293.7645,146.3843) --(293.8641,146.4541) --(293.9538,146.5238) --(294.0534,146.6035) --(294.1530,146.6732) --(294.2427,146.7529) --(294.3423,146.8227) --(294.4419,146.8924) --(294.5316,146.9721) --(294.6312,147.0419) --(294.7309,147.1216) --(294.8205,147.1913) --(294.9202,147.2710) --(295.0198,147.3407) --(295.1094,147.4105) --(295.2091,147.4902) --(295.3087,147.5599) --(295.3984,147.6396) --(295.4980,147.7094) --(295.5976,147.7891) --(295.6873,147.8588) --(295.7869,147.9385) --(295.8865,148.0082) --(295.9762,148.0880) --(296.0758,148.1577) --(296.1754,148.2374) --(296.2651,148.3071) --(296.3647,148.3868) --(296.4644,148.4566) --(296.5540,148.5263) --(296.6537,148.6060) --(296.7533,148.6758) --(296.8430,148.7555) --(296.9426,148.8352) --(297.0422,148.9049) --(297.1418,148.9846) --(297.2315,149.0543) --(297.3311,149.1340) --(297.4308,149.2038) --(297.5204,149.2835) --(297.6201,149.3532) --(297.7197,149.4329) --(297.8094,149.5027) --(297.9090,149.5824) --(298.0086,149.6521) --(298.0983,149.7318) --(298.1979,149.8015) --(298.2975,149.8812) --(298.3872,149.9610) --(298.4868,150.0307) --(298.5865,150.1104) --(298.6761,150.1801) --(298.7758,150.2598) --(298.8754,150.3296) --(298.9651,150.4093) --(299.0647,150.4890) --(299.1643,150.5587) --(299.2539,150.6384) --(299.3536,150.7082) --(299.4532,150.7878) --(299.5429,150.8676) --(299.6425,150.9373) --(299.7421,151.0170) --(299.8318,151.0867) --(299.9314,151.1664) --(300.0310,151.2461) --(300.1207,151.3159) --(300.2204,151.3956) --(300.3200,151.4653) --(300.4096,151.5450) --(300.5093,151.6247) --(300.6089,151.6945) --(300.7085,151.7742) --(300.7982,151.8539) --(300.8978,151.9236) --(300.9974,152.0033) --(301.0871,152.0830) --(301.1867,152.1528) --(301.2864,152.2325) --(301.3760,152.3022) --(301.4757,152.3819) --(301.5753,152.4616) --(301.6649,152.5313) --(301.7646,152.6111) --(301.8642,152.6908) --(301.9539,152.7704) --(302.0535,152.8402) --(302.1531,152.9199) --(302.2428,152.9996) --(302.3424,153.0693) --(302.4420,153.1490) --(302.5317,153.2287) --(302.6313,153.2985) --(302.7310,153.3782) --(302.8206,153.4579) --(302.9203,153.5276) --(303.0199,153.6073) --(303.1096,153.6870) --(303.2092,153.7667) --(303.3088,153.8365) --(303.3985,153.9162) --(303.4981,153.9959) --(303.5977,154.0656) --(303.6874,154.1453) --(303.7870,154.2250) --(303.8867,154.3047) --(303.9763,154.3745) --(304.0760,154.4542) --(304.1756,154.5339) --(304.2652,154.6136) --(304.3649,154.6833) --(304.4645,154.7630) --(304.5641,154.8427) --(304.6538,154.9224) --(304.7534,154.9922) --(304.8531,155.0719) --(304.9427,155.1516) --(305.0423,155.2312) --(305.1420,155.3010) --(305.2316,155.3807) --(305.3313,155.4604) --(305.4309,155.5401) --(305.5206,155.6198) --(305.6202,155.6895) --(305.7198,155.7692) --(305.8095,155.8490) --(305.9091,155.9286) --(306.0087,156.0083) --(306.0984,156.0781) --(306.1980,156.1578) --(306.2976,156.2375) --(306.3873,156.3172) --(306.4869,156.3969) --(306.5866,156.4666) --(306.6762,156.5463) --(306.7758,156.6261) --(306.8755,156.7057) --(306.9651,156.7854) --(307.0648,156.8552) --(307.1644,156.9349) --(307.2541,157.0146) --(307.3537,157.0943) --(307.4533,157.1740) --(307.5430,157.2537) --(307.6426,157.3234) --(307.7422,157.4031) --(307.8319,157.4828) --(307.9315,157.5625) --(308.0312,157.6422) --(308.1308,157.7219) --(308.2205,157.8017) --(308.3201,157.8714) --(308.4197,157.9511) --(308.5094,158.0308) --(308.6090,158.1105) --(308.7086,158.1902) --(308.7983,158.2699) --(308.8979,158.3496) --(308.9976,158.4193) --(309.0872,158.4990) --(309.1869,158.5787) --(309.2865,158.6584) --(309.3761,158.7381) --(309.4758,158.8179) --(309.5754,158.8976) --(309.6651,158.9773) --(309.7647,159.0570) --(309.8643,159.1267) --(309.9540,159.2064) --(310.0536,159.2861) --(310.1532,159.3658) --(310.2429,159.4455) --(310.3425,159.5252) --(310.4422,159.6049) --(310.5318,159.6846) --(310.6314,159.7643) --(310.7311,159.8440) --(310.8207,159.9237) --(310.9204,159.9935) --(311.0200,160.0732) --(311.1097,160.1529) --(311.2093,160.2326) --(311.3089,160.3123) --(311.3986,160.3920) --(311.4982,160.4717) --(311.5978,160.5514) --(311.6975,160.6311) --(311.7871,160.7108) --(311.8867,160.7905) --(311.9864,160.8702) --(312.0760,160.9499) --(312.1757,161.0296) --(312.2753,161.1093) --(312.3650,161.1890) --(312.4646,161.2587) --(312.5642,161.3384) --(312.6539,161.4181) --(312.7535,161.4978) --(312.8531,161.5775) --(312.9428,161.6572) --(313.0424,161.7369) --(313.1421,161.8166) --(313.2317,161.8963) --(313.3314,161.9760) --(313.4310,162.0558) --(313.5207,162.1355) --(313.6203,162.2152) --(313.7199,162.2949) --(313.8096,162.3746) --(313.9092,162.4543) --(314.0088,162.5340) --(314.0985,162.6137) --(314.1981,162.6934) --(314.2978,162.7731) --(314.3874,162.8528) --(314.4871,162.9325) --(314.5867,163.0122) --(314.6763,163.0919) --(314.7760,163.1716) --(314.8756,163.2513) --(314.9653,163.3310) --(315.0649,163.4107) --(315.1645,163.4904) --(315.2641,163.5701) --(315.3538,163.6498) --(315.4534,163.7295) --(315.5531,163.8092) --(315.6427,163.8889) --(315.7423,163.9686) --(315.8420,164.0483) --(315.9316,164.1280) --(316.0313,164.2077) --(316.1309,164.2874) --(316.2206,164.3671) --(316.3202,164.4468) --(316.4198,164.5265) --(316.5095,164.6062) --(316.6091,164.6859) --(316.7087,164.7656) --(316.7984,164.8453) --(316.8980,164.9350) --(316.9977,165.0147) --(317.0873,165.0944) --(317.1870,165.1741) --(317.2866,165.2538) --(317.3763,165.3335) --(317.4759,165.4132) --(317.5755,165.4929) --(317.6652,165.5726) --(317.7648,165.6523) --(317.8644,165.7320) --(317.9541,165.8117) --(318.0537,165.8914) --(318.1534,165.9711) --(318.2430,166.0508) --(318.3426,166.1305) --(318.4423,166.2102) --(318.5319,166.2999) --(318.6316,166.3796) --(318.7312,166.4593) --(318.8308,166.5390) --(318.9205,166.6187) --(319.0201,166.6984) --(319.1197,166.7781) --(319.2094,166.8578) --(319.3090,166.9375) --(319.4087,167.0172) --(319.4983,167.0969) --(319.5979,167.1766) --(319.6976,167.2563) --(319.7872,167.3460) --(319.8869,167.4257) --(319.9865,167.5054) --(320.0762,167.5851) --(320.1758,167.6648) --(320.2754,167.7445) --(320.3651,167.8242) --(320.4647,167.9039) --(320.5643,167.9836) --(320.6540,168.0633) --(320.7536,168.1530) --(320.8532,168.2327) --(320.9429,168.3124) --(321.0425,168.3921) --(321.1422,168.4718) --(321.2318,168.5515) --(321.3315,168.6312) --(321.4311,168.7109) --(321.5208,168.7906) --(321.6204,168.8802) --(321.7200,168.9599) --(321.8097,169.0396) --(321.9093,169.1194) --(322.0089,169.1991) --(322.0986,169.2788) --(322.1982,169.3585) --(322.2979,169.4382) --(322.3975,169.5278) --(322.4872,169.6075) --(322.5868,169.6872) --(322.6864,169.7669) --(322.7761,169.8466) --(322.8757,169.9263) --(322.9753,170.0060) --(323.0650,170.0857) --(323.1646,170.1754) --(323.2643,170.2551) --(323.3539,170.3348) --(323.4536,170.4145) --(323.5532,170.4942) --(323.6428,170.5739) --(323.7425,170.6536) --(323.8421,170.7433) --(323.9317,170.8230) --(324.0314,170.9027) --(324.1310,170.9824) --(324.2206,171.0621) --(324.3203,171.1418) --(324.4199,171.2215) --(324.5096,171.3112) --(324.6092,171.3909) --(324.7088,171.4706) --(324.7985,171.5503) --(324.8981,171.6300) --(324.9978,171.7097) --(325.0874,171.7993) --(325.1870,171.8790) --(325.2867,171.9587) --(325.3763,172.0385) --(325.4760,172.1181) --(325.5756,172.1978) --(325.6653,172.2775) --(325.7649,172.3672) --(325.8645,172.4469) --(325.9641,172.5266) --(326.0538,172.6063) --(326.1534,172.6860) --(326.2531,172.7657) --(326.3427,172.8554) --(326.4424,172.9351) --(326.5420,173.0148) --(326.6317,173.0945) --(326.7313,173.1742) --(326.8309,173.2639) --(326.9206,173.3436) --(327.0202,173.4233) --(327.1198,173.5030) --(327.2095,173.5827) --(327.3091,173.6624) --(327.4088,173.7520) --(327.4984,173.8317) --(327.5980,173.9114) --(327.6977,173.9911) --(327.7873,174.0708) --(327.8870,174.1506) --(327.9866,174.2402) --(328.0762,174.3199) --(328.1759,174.3996) --(328.2755,174.4793) --(328.3652,174.5590) --(328.4648,174.6487) --(328.5644,174.7284) --(328.6541,174.8081) --(328.7537,174.8878) --(328.8533,174.9675) --(328.9430,175.0572) --(329.0426,175.1369) --(329.1422,175.2166) --(329.2319,175.2963) --(329.3315,175.3760) --(329.4312,175.4557) --(329.5308,175.5453) --(329.6205,175.6250) --(329.7201,175.7047) --(329.8197,175.7844) --(329.9094,175.8642) --(330.0090,175.9538) --(330.1086,176.0335) --(330.1983,176.1132) --(330.2979,176.1929) --(330.3976,176.2726) --(330.4872,176.3623) --(330.5869,176.4420) --(330.6865,176.5217) --(330.7762,176.6014) --(330.8758,176.6811) --(330.9754,176.7708) --(331.0651,176.8505) --(331.1647,176.9302) --(331.2643,177.0099) --(331.3540,177.0896) --(331.4536,177.1792) --(331.5533,177.2589) --(331.6429,177.3386) --(331.7426,177.4183) --(331.8422,177.4980) --(331.9318,177.5877) --(332.0315,177.6674) --(332.1311,177.7471) --(332.2208,177.8268) --(332.3204,177.9165) --(332.4200,177.9962) --(332.5097,178.0759) --(332.6093,178.1556) --(332.7089,178.2353) --(332.7986,178.3250) --(332.8982,178.4047) --(332.9979,178.4844) --(333.0975,178.5641) --(333.1871,178.6438) --(333.2868,178.7334) --(333.3864,178.8131) --(333.4761,178.8928) --(333.5757,178.9725) --(333.6753,179.0522) --(333.7650,179.1419) --(333.8646,179.2216) --(333.9642,179.3013) --(334.0539,179.3810) --(334.1535,179.4707) --(334.2532,179.5504) --(334.3428,179.6301) --(334.4424,179.7098) --(334.5421,179.7895) --(334.6317,179.8792) --(334.7314,179.9588) --(334.8310,180.0386) --(334.9207,180.1183) --(335.0203,180.1980) --(335.1199,180.2876) --(335.2096,180.3673) --(335.3092,180.4470) --(335.4088,180.5267) --(335.4985,180.6164) --(335.5981,180.6961) --(335.6978,180.7758) --(335.7874,180.8555) --(335.8871,180.9352) --(335.9867,181.0249) --(336.0764,181.1046) --(336.1760,181.1843) --(336.2756,181.2640) --(336.3653,181.3437) --(336.4649,181.4333) --(336.5645,181.5130) --(336.6642,181.5927) --(336.7538,181.6724) --(336.8535,181.7621) --(336.9531,181.8418) --(337.0427,181.9215) --(337.1424,182.0012) --(337.2420,182.0809) --(337.3317,182.1706) --(337.4313,182.2503) --(337.5309,182.3300) --(337.6206,182.4097) --(337.7202,182.4894) --(337.8198,182.5791) --(337.9095,182.6588) --(338.0091,182.7384) --(338.1087,182.8182) --(338.1984,182.9078) --(338.2980,182.9875) --(338.3977,183.0672) --(338.4873,183.1469) --(338.5870,183.2266) --(338.6866,183.3163) --(338.7762,183.3960) --(338.8759,183.4757) --(338.9755,183.5554) --(339.0652,183.6451) --(339.1648,183.7248) --(339.2644,183.8045) --(339.3541,183.8842) --(339.4537,183.9639) --(339.5533,184.0535) --(339.6430,184.1333) --(339.7426,184.2130) --(339.8423,184.2926) --(339.9319,184.3723) --(340.0316,184.4620) --(340.1312,184.5417) --(340.2308,184.6214);
\path[fill=none,draw=chfourfig6c2,line width=0.3985pt,draw opacity=1.0,dash pattern=on .99628pt off 2.98883pt] (51.0415,187.7099) --(51.0415,17.6653);
\path[fill=none,draw=chfourfig6c2,line width=0.3985pt,draw opacity=1.0,dash pattern=on .99628pt off 2.98883pt] (99.2315,187.7099) --(99.2315,17.6653);
\path[fill=none,draw=chfourfig6c2,line width=0.3985pt,draw opacity=1.0,dash pattern=on .99628pt off 2.98883pt] (147.4313,187.7099) --(147.4313,17.6653);
\path[fill=none,draw=chfourfig6c2,line width=0.3985pt,draw opacity=1.0,dash pattern=on .99628pt off 2.98883pt] (195.6312,187.7099) --(195.6312,17.6653);
\path[fill=none,draw=chfourfig6c2,line width=0.3985pt,draw opacity=1.0,dash pattern=on .99628pt off 2.98883pt] (243.8311,187.7099) --(243.8311,17.6653);
\path[fill=none,draw=chfourfig6c2,line width=0.3985pt,draw opacity=1.0,dash pattern=on .99628pt off 2.98883pt] (292.0309,187.7099) --(292.0309,17.6653);
\path[fill=none,draw=chfourfig6c2,line width=0.3985pt,draw opacity=1.0,dash pattern=on .99628pt off 2.98883pt] (340.2308,187.7099) --(340.2308,17.6653);
\path[fill=none,draw=chfourfig6c2,line width=0.3985pt,draw opacity=1.0,dash pattern=on .99628pt off 2.98883pt] (39.6242,183.4259) --(347.7826,183.4259);
\path[fill=none,draw=chfourfig6c2,line width=0.3985pt,draw opacity=1.0,dash pattern=on .99628pt off 2.98883pt] (39.6242,132.9645) --(347.7826,132.9645);
\path[fill=none,draw=chfourfig6c2,line width=0.3985pt,draw opacity=1.0,dash pattern=on .99628pt off 2.98883pt] (39.6242,82.5030) --(347.7826,82.5030);
\path[fill=none,draw=chfourfig6c2,line width=0.3985pt,draw opacity=1.0,dash pattern=on .99628pt off 2.98883pt] (39.6242,32.0317) --(347.7826,32.0317);
\node[font=\small,inner sep=0pt,rotate=0] (ch04f6t0) at (51.0405,204.5411) {$0$};
\node[font=\small,inner sep=0pt,rotate=0] (ch04f6t1) at (99.2296,204.5411) {$500$};
\node[font=\small,inner sep=0pt,rotate=0] (ch04f6t2) at (147.4292,204.5411) {$1000$};
\node[font=\small,inner sep=0pt,rotate=0] (ch04f6t3) at (195.6282,204.5411) {$1500$};
\node[font=\small,inner sep=0pt,rotate=0] (ch04f6t4) at (243.8272,204.5411) {$2000$};
\node[font=\small,inner sep=0pt,rotate=0] (ch04f6t5) at (292.0272,204.5411) {$2500$};
\node[font=\small,inner sep=0pt,rotate=0] (ch04f6t6) at (340.2262,204.5411) {$3000$};
\node[font=\small,inner sep=0pt,rotate=0] (ch04f6t7) at (24.1751,183.4268) {$-5$};
\node[font=\small,inner sep=0pt,rotate=0] (ch04f6t8) at (24.1751,132.9655) {$0$};
\node[font=\small,inner sep=0pt,rotate=0] (ch04f6t9) at (24.1751,82.5045) {$5$};
\node[font=\small,inner sep=0pt,rotate=0] (ch04f6t10) at (24.1751,32.0339) {$10$};
\node[font=\small,inner sep=0pt,rotate=0] (ch04f6t11) at (193.7060,218.5281) {海拔（米）};
\node[font=\small,inner sep=0pt,rotate=90] (ch04f6t12) at (10.1871,102.6891) {平均气温（摄氏度）};
\end{tikzpicture}
\caption{德国各地海拔（米，$X$）与平均气温（摄氏度，$Y$）之间的关系。数据来自德国气象局（Deutscher Wetterdienst），另见 Mooij et al.［2016］。带加性噪声的非线性函数能够较合理地拟合这些经验观测；在远离海平面的范围内，该函数接近线性。}
\label{fig:ch04-6}
\end{figure}


作为上述方法的替代，也可以采用基于\textbf{最大似然（maximum likelihood）}的方法。例如，考虑具有高斯加性误差项的非线性结构因果模型。此时，可以比较两个模型的似然得分，以区分 $X\to Y$ 和 $X\leftarrow Y$。为此，先对 $Y$ 关于 $X$ 进行非线性回归，得到残差 $R_Y:=Y-\widehat f_Y(X)$。然后，将
\begin{equation}
 L_{X\to Y}
 =-\log\widehat{\operatorname{var}}[X]
  -\log\widehat{\operatorname{var}}[R_Y]
 \tag{4.18}
\end{equation}
与交换 $X,Y$ 的角色后得到的对应版本
\begin{equation}
 L_{X\leftarrow Y}
 =-\log\widehat{\operatorname{var}}[R_X]
  -\log\widehat{\operatorname{var}}[Y]
 \tag{4.19}
\end{equation}
进行比较。不难证明，见习题4.16，如果不实施回归，而是使用真实条件均值 $\widehat f_Y(x)=\EE[Y\mid X=x]$，并类似地处理 $\widehat f_X$，那么这的确对应于比较似然。不过，与前面一样，先回归、再计算样本方差的两步过程需要理论论证。Bühlmann et al.［2014］利用经验过程理论［van de Geer, 2009］证明了其一致性。如果噪声不一定服从高斯分布，就必须调整得分函数，将残差经验方差的对数替换为误差项微分熵的估计［Nowzohour and Bühlmann, 2016］。

\Needspace{14cm}
\begin{codesnippet}
下面的代码给出一个有限数据集的例子。它使用了 \texttt{dHSIC}［Pfister et al., 2017］和 \texttt{mgcv}［Wood, 2006］两个软件包。前者包含函数 \texttt{dhsic.test}，实现了 Gretton et al.［2008］提出的独立性检验；后者包含函数 \texttt{gam}，在第10、11行中被用作非线性回归方法，见第A.1节。只有在反方向上，残差与输入的独立性才被拒绝，见第15、17行。第21、23行还显示，高斯似然得分同样支持正向模型，另见式 (4.18)、(4.19)。
\begin{lstlisting}[language=R]
library(dHSIC)
library(mgcv)
#
# generate data set
set.seed(1)
X <- rnorm(200)
Y <- X^3 + rnorm(200)
#
# fit models
modelforw <- gam(Y ~ s(X))
modelbackw <- gam(X ~ s(Y))
#
# independence tests
dhsic.test(modelforw$residuals, X)$p.value
# [1] 0.7628932
dhsic.test(modelbackw$residuals, Y)$p.value
# [1] 0.004221031
#
# computing likelihoods
- log(var(X)) - log(var(modelforw$residuals))
# [1] 0.1420063
- log(var(modelbackw$residuals)) - log(var(Y))
# [1] -1.014013
\end{lstlisting}
\end{codesnippet}

\subsection{信息几何因果推断}

这里简要概述 IGCI 的实现，详情参见 Mooij et al.［2016］。其理论基础是定理4.10的可识别性结果，以及由此得到的一些简单推论。可以证明，如果定义
\[
 C_{X\to Y}:=\int_0^1\log f'(x)p(x)\,\mathrm{d}x,
\]
并类似地定义 $C_{Y\to X}$，那么独立性条件 (4.10) 蕴含
\[
 C_{X\to Y}\leq C_{Y\to X}.
\]
这里使用下面这个直接的估计量：
\[
 \widehat C_{X\to Y}:=
 \frac{1}{N-1}\sum_{j=1}^{N-1}
 \log\frac{|y_{j+1}-y_j|}{|x_{j+1}-x_j|},
\]
其中，$x_1<x_2<\cdots<x_N$ 是按升序排列的观测 $x$ 值。如果 $Y$ 是 $X$ 的递增函数，相应的 $y$ 值也按顺序排列；但对于真实数据，通常并非如此。类似地定义估计量 $\widehat C_{Y\to X}$，并在 $\widehat C_{X\to Y}<\widehat C_{Y\to X}$ 时推断 $X\to Y$。

除了这种所谓基于斜率的方法，还存在一种基于熵的方法。可以证明，式 (4.10) 还蕴含\chfourtnote{原书所印 $H(X)\leq H(Y)$ 与紧随其后的解释不一致。对于本节的递增微分同胚，有 $H(Y)=H(X)+\EE[\log f^\prime(X)]$；由独立性条件及 Jensen 不等式可得 $\EE[\log f^\prime(X)]\leq0$，因而应为 $H(Y)\leq H(X)$，与“熵较大的变量是原因”一致。正文保留原书不等号。}
\[
 H(X)\leq H(Y),
\]
其中，$H$ 表示微分 Shannon 熵：
\[
 H(X):=-\int_0^1p(x)\log p(x)\,\mathrm{d}x.
\]
% 原书此处不等号与下文“较大熵者为原因”及前述积分推论不一致，见报告。
直观地说，将非线性函数 $f$ 作用于 $p_X$，会引入额外的不规则性，从而使 $p_Y$ 比 $p_X$ 更不均匀；只有当 $f$ 的非线性相对于 $p_X$ 作了调谐时才会例外。因此，假定熵较大的变量是原因。为估计 $H$，可以采用文献中任意一种标准的熵估计量。

\subsection{迹方法}

回顾一下，这种方法依赖高维变量 $X$ 和 $Y$ 之间的线性关系。先假设样本量相对于 $X,Y$ 的维度足够大，能够通过标准线性回归估计协方差矩阵 $\Sigma_{XX},\Sigma_{YY}$ 以及结构矩阵 $A_Y,A_X$。为了利用定理4.12的可识别性结果，可以计算\textbf{迹相依比（tracial dependency ratio）}
\[
 r_{X\to Y}:=
 \frac{\tau(A_Y\Sigma_{XX}A_Y^{\mathsf T})}
 {\tau(A_YA_Y^{\mathsf T})\tau(\Sigma_{XX})},
\]
并通过交换 $X,Y$ 的角色，类似地计算 $r_{Y\to X}$。然后，将更接近1的比值所对应的方向推断为因果方向［Janzing et al., 2010］。

Zscheischler et al.［2011］给出了一种判断这种对1的偏离是否显著的方法。其生成模型通过一个随机正交映射，使两个矩阵 $A$ 与 $\Sigma_{XX}$ 发生相对旋转，以模拟二者的独立性。\textbf{自由概率论（free probability theory）}［Voiculescu, 1997］是描述大型随机矩阵渐近行为的一套数学框架。利用其中的思想，Zscheischler et al.［2011］构造了维度大于样本量时迹条件的实现方法。他们证明，在无噪声情形下，只要额外假定 $A$ 与 $X$ 的经验协方差矩阵独立，即使没有足够数据估计 $\Sigma_{XX}$ 和 $A$，仍然能够估计 $r_{X\to Y}$。因此，当估计量显著偏离1时，就可以拒绝假设 $X\to Y$。此时，要么额外的独立性假设不成立，要么 $r_{X\to Y}$ 的确显著偏离1。

\subsection{监督学习方法}

最后，介绍一种更接近机器学习视角的因果学习方法。原则上，这种方法既能利用受限函数类，也能利用独立性条件。假设给定带标签的训练数据
\[
 (D_1,A_1),\ldots,(D_n,A_n).
\]
其中，每个 $D_i$ 都是一个数据集，
\[
 D_i=\{(X_1,Y_1),\ldots,(X_{n_i},Y_{n_i})\},
\]
包含实现值 $(X_1,Y_1),\ldots,(X_{n_i},Y_{n_i})\overset{\mathrm{i.i.d.}}{\sim}P^i_{X,Y}$；每个标签 $A_i\in\{\to,\leftarrow\}$ 说明数据集 $D_i$ 对应 $X\to Y$ 还是 $X\leftarrow Y$。这样，因果学习就变成了一个经典的预测问题。我们可以训练分类器，并希望它们能够从因果真相已知的数据集，良好地泛化到未见过的测试数据集。

据我们所知，Guyon［2013］最早以竞赛形式系统研究了这一方法，提供了由模拟数据和真实数据混合构成的、因果真相已知的数据集。显然，若只利用相关、协方差等对称特征，这种方法就不可能成功。竞赛中许多有竞争力的分类器都基于人工构造的特征。例如，边缘分布的熵估计，以及对 $X$ 关于 $Y$ 或 $Y$ 关于 $X$ 进行回归后，所得残差分布的熵估计。

有意思的是，这些特征可以与加性噪声模型的概念联系起来。例如，对于高斯分布变量，熵是方差对数的线性函数。因此，这些特征具有足够的表达能力，可以重构式 (4.18)、(4.19) 的得分。用熵代替方差的对数，则相当于放宽高斯性假设［Nowzohour and Bühlmann, 2016］。

Lopez-Paz et al.［2015］致力于自动构造这类特征。其思想是，将联合分布 $P^i_{X,Y}$（$i=1,\ldots,n$）映射到一个\textbf{再生核 Hilbert 空间（reproducing kernel Hilbert space）}，见附录A.2，再在这个空间中分类。在实际中，无法获得完整分布 $P^i_{X,Y}$，而是以经验分布近似它。类似的方法也曾用于区分时间逆转后的时间序列与其原始版本［Peters et al., 2009a］。由于原因与结果的分类似乎依赖于联合分布相当复杂的性质，因此，训练集需要较大的样本量 $n$。为了加入有用的模拟数据集，必须从具有可识别性的情形中生成这些数据。例如，Lopez-Paz et al.［2015］就加入了来自加性噪声模型的额外样本。

监督学习方法目前还不能作为独立的因果学习方法使用。不过，作为能够有效利用已知可识别性性质及其组合的统计工具，它们也许会很有价值。
